%%%%%%%%%%%%%%%%%%%%%%%%%%%%%%%%%%%%%%%%%%%%%%%%%%%%%%%%%%%%%%%%%%%%%%%%%%%%%%%%
% PREAMBOLO
%%%%%%%%%%%%%%%%%%%%%%%%%%%%%%%%%%%%%%%%%%%%%%%%%%%%%%%%%%%%%%%%%%%%%%%%%%%%%%%%
\documentclass[12pt, a4paper, openany]{book} % <-- Aggiunto "openany" per rimuovere le pagine bianche

% --- Lingua e Codifica ---
\usepackage[utf8]{inputenc}
\usepackage[T1]{fontenc}
\usepackage[italian]{babel}

% --- Matematica ---
\usepackage{amsmath}
\usepackage{amssymb}
\usepackage{amsfonts}

% --- Grafica e Layout ---
\usepackage{graphicx}
\usepackage[a4paper, top=2.5cm, bottom=2.5cm, left=2.5cm, right=2.5cm]{geometry}
\usepackage{fancyhdr} % Per personalizzare header e footer
\setlength{\headheight}{14.5pt}


\usepackage{wrapfig} 
\usepackage{caption} 

\captionsetup{labelfont=bf} % Aggiunge il grassetto all'etichetta della didascalia

\usepackage{draftwatermark}
\usepackage{xcolor}
\usepackage[most]{tcolorbox}
\usepackage{pgfplots}
\usetikzlibrary{arrows.meta, decorations.markings, shapes.geometric, calc, positioning}
\usetikzlibrary{decorations.pathmorphing, arrows.meta, chains, shapes.geometric, shapes.misc}

\usepackage{tabularx}
\usepackage{fontawesome5} % Per inserire le icone (es. \faGithub)

\usepackage{lmodern} % Font moderno opzionale
\SetWatermarkText{\sffamily\bfseries\textcolor{gray!10}{github.com / elisabettagnello}}
\SetWatermarkScale{0.3}
\SetWatermarkColor{gray!8}
\SetWatermarkAngle{45}

% Definizione della palette sequenziale (Viridis-like)
\definecolor{colorA}{HTML}{003f5c} % 1 - Blu Profondo
\definecolor{colorB}{HTML}{2f4b7c} % 2 - Indaco
\definecolor{colorC}{HTML}{665191} % 3 - Viola Bluastro
\definecolor{colorD}{HTML}{a05195} % 4 - Magenta Scuro
\definecolor{colorE}{HTML}{d45087} % 5 - Rosso Porpora
\definecolor{colorF}{HTML}{f95d6a} % 6 - Arancio-Rosso
\definecolor{colorG}{HTML}{ff7c43} % 7 - Arancio Brillante

\usepackage{hyperref} % Per link interni (es. tabella dei contenuti, riferimenti)
\hypersetup{
    colorlinks=true,
    linkcolor=colorA,
    filecolor=colorE,      
    urlcolor=colorB,
}

%%%%%%%%%%%%%%%%%%%%%%%%%%%%%%%%%%%%%%%%%%%%%%%%%%%%%%%%%%%%%%%%%%%%%%%%%%%%%%%%
% CONFIGURAZIONE STILE PAGINA (FANCYHDR)
%%%%%%%%%%%%%%%%%%%%%%%%%%%%%%%%%%%%%%%%%%%%%%%%%%%%%%%%%%%%%%%%%%%%%%%%%%%%%%%%
\pagestyle{fancy}
\fancyhf{} % Pulisce tutti i campi di header e footer
\fancyhead[LE,RO]{\nouppercase{\leftmark}} % Titolo del capitolo in alto a sinistra (pagine pari) e destra (pagine dispari)
\fancyfoot[C]{\thepage} % Numero di pagina al centro del footer
\renewcommand{\headrulewidth}{0.4pt} % Linea sotto l'header
\renewcommand{\footrulewidth}{0.4pt} % Linea sopra il footer

\renewcommand{\chaptermark}[1]{\markboth{#1}{}}

\usepackage{titlesec}

\titleformat{\chapter}[display]
  {\normalfont\huge\bfseries} % stile del titolo
  {}                
  {0pt}                      % spazio
  {\Huge}                     % stile applicato al titolo vero e proprio


%%%%%%%%%%%%%%%%%%%%%%%%%%%%%%%%%%%%%%%%%%%%%%%%%%%%%%%%%%%%%%%%%%%%%%%%%%%%%%%%
% INIZIO DEL DOCUMENTO
%%%%%%%%%%%%%%%%%%%%%%%%%%%%%%%%%%%%%%%%%%%%%%%%%%%%%%%%%%%%%%%%%%%%%%%%%%%%%%%%
\begin{document}

% --- Pagina del Titolo (da personalizzare) ---
\title{\textbf{\Huge Neural Networks}}
\author{Autore: Agnello Elisabetta}
\date{anno: 2023/2024}
\maketitle

% --- Indice ---
\frontmatter % Numerazione romana per le pagine iniziali (indice)

% --- Disclaimer ---
\chapter*{Disclaimer}
\thispagestyle{empty} % Rimuove header e footer solo per questa pagina

Il presente documento costituisce una raccolta di appunti personali presi durante le lezioni del corso di \textit{Neural Networks}. 

Si prega di notare che questo materiale \textbf{non è stato sottoposto a una revisione sistematica} né è stato verificato dal docente. Pertanto, è possibile che siano presenti:
\begin{itemize}
    \item Errori concettuali o imprecisioni tecniche;
    \item Refusi o errori di battitura;
    \item Parti incomplete o sezioni mancanti.
\end{itemize}

Questi appunti sono destinati esclusivamente come supporto allo studio e non sostituiscono in alcun modo il materiale didattico ufficiale. L'autore non si assume alcuna responsabilità per l'uso delle informazioni qui contenute.


% =============================================================================
\vspace{3cm} % Spazio verticale per staccare il box dal testo sopra

\begin{tcolorbox}[
    colback=colorA!5,      
    colframe=colorA,       
    title={Contatti e Segnalazioni},
    fonttitle=\bfseries\sffamily,
    arc=3mm,               
    boxrule=1.2pt,         
    halign title=flush left
]
Per consultare l'archivio aggiornato, segnalare refusi o proporre integrazioni agli appunti, puoi utilizzare i seguenti canali:
\begin{description}
    \item[\textsf{\faGlobe\ Archivio Appunti:}] \href{https://elisabettagnello.github.io/elisabetta-physics-notes/}{\texttt{elisabettagnello.github.io/elisabetta-physics-notes/}}
    \item[\textsf{\faGithub\ GitHub:}] \href{https://github.com/elisabettagnello}{\texttt{github.com/elisabettagnello}}
    \item[\textsf{\faLinkedin\ LinkedIn:}] \href{https://www.linkedin.com/in/elisabetta-agnello}{\texttt{linkedin.com/in/elisabetta-agnello}}
    \item[\textsf{\faEnvelope\ Email:}] \href{mailto:agnelloe24@gmail.com}{\texttt{agnelloe24@gmail.com}}
\end{description}
\end{tcolorbox}
% =============================================================================


\clearpage % Forza il passaggio alla pagina successiva

\tableofcontents

\mainmatter 

% --- INCLUSIONE DEI CAPITOLI ---

\chapter{Lezione 1}

\textbf{Docente:} Maurizio Matti (Istituto Nazionale Italiano di Salute, ricercatore senior; fisico con PhD in neurofisiologia)\\
\textbf{Argomenti:} Presentazione del corso, modello di Hodgkin-Huxley, equazione di Fokker-Planck, population firing rates, dinamica non lineare, neuroanatomia, struttura del neurone, organizzazione corticale, aree di Brodmann, storia (Broca, Wernicke, Cajal), elaborazione visiva, scaling cerebrale\\
\textbf{Testo di Riferimento:} Wolfram Gerstner, \textit{Neuronal Dynamics}


\vspace{0.5cm}
\noindent\rule{\textwidth}{0.4pt}


\subsubsection{Argomenti Principali del Corso}

Il corso si concentra sulle \textbf{reti neurali biologiche}, non sulle reti neurali artificiali del machine learning. I temi principali includono:

\begin{enumerate}
    \item \textbf{Composizione e organizzazione del cervello}
    \item \textbf{Elaborazione dell'informazione} e i vincoli fisici che la governano
    \item \textbf{Il singolo neurone}: struttura, dinamica elettrochimica
    \item \textbf{Modello di Hodgkin-Huxley}: descrizione quantitativa del potenziale d'azione
    \item \textbf{Neuroni spiking semplificati}: modelli a riduzione di dimensionalità
    \item \textbf{Equazione di Fokker-Planck} e approcci della fisica statistica
    \item \textbf{Population firing rates} (tassi di scarica delle popolazioni neuronali)
    \item \textbf{Dinamica non lineare} alla base delle funzioni cognitive: decision-making, planning, elaborazione sensoriale, memoria
\end{enumerate}


\subsubsection{Perché Studiare le Reti Neurali Biologiche}

Le reti neurali artificiali hanno prodotto risultati straordinari: \textbf{ChatGPT}, \textbf{LaMDA} e modelli analoghi dimostrano capacità emergenti di elaborazione del linguaggio, ragionamento e generalizzazione. Questi sistemi utilizzano \textbf{reti profonde} (\textit{deep networks}) e \textbf{reinforcement learning}.

\begin{tcolorbox}[colback=gray!5, colframe=gray!40, arc=2mm, boxrule=0.5pt]
I neuroni nelle reti biologiche si auto-organizzano in sequenze simboliche a livelli progressivi: sillabe, parole, frasi. Le reti biologiche dimostrano che sistemi fisici composti da molti elementi semplici possono apprendere, gestire linguaggio, matematica e musica, e generalizzare a situazioni nuove.
\end{tcolorbox}



Dal punto di vista fisico, il cervello è particolarmente sfidante per la sua:

\begin{itemize}
    \item \textbf{Eterogeneità}: diversi tipi di neuroni, connessioni, trasmettitori
    \item \textbf{Gerarchia spaziale e temporale}: organizzazione su scale molto diverse
\end{itemize}

Il cervello opera su \textbf{molteplici scale temporali}:

\begin{table}[htbp]
    \centering
    \renewcommand{\arraystretch}{1.3}
    \begin{tabularx}{\textwidth}{|l|X|}
    \hline
    \textbf{Scala temporale} & \textbf{Processo} \\
    \hline
    Millisecondi & Dinamica dei singoli neuroni (spike) \\
    Minuti & Elaborazione del parlato, comprensione del testo \\
    Ore–giorni & Consolidamento della memoria a breve termine \\
    Anni–decenni & Memorie a lungo termine, apprendimento motorio \\
    \hline
    \end{tabularx}
\end{table}


La \textbf{neuroscienza computazionale} è un campo relativamente giovane (iniziato negli anni '80), caratterizzato da una crescente interazione tra sperimentatori e teorici. I fisici contribuiscono sempre di più alla progettazione degli esperimenti, degli strumenti di analisi e delle simulazioni computazionali.


\newpage

\section{Struttura del Cervello}

\subsubsection{Elementi del Cervello}

Il cervello è un organo complesso composto da diversi elementi funzionali. Il corso si concentra principalmente sulla \textbf{corteccia} (la parte esterna), ma è fondamentale comprendere anche i principali \textbf{elementi subcorticali}:

\begin{itemize}
    \item \textbf{Gangli della base}: regioni ricche di neuroni \textbf{dopaminergici}, rilevanti per l'apprendimento, la ricompensa e la motivazione. Hanno un ruolo cruciale nei meccanismi di reinforcement learning biologico.
    \item \textbf{Tronco cerebrale} (\textit{brainstem}): regola i \textbf{ritmi circadiani} e il controllo di base del corpo (respirazione, battito cardiaco, vigilanza).
    \item \textbf{Cervelletto}: spesso chiamato "secondo cervello", è specializzato per la \textbf{coordinazione motoria} e la \textbf{predizione} dei movimenti. Riveste un ruolo importante nell'apprendimento motorio.
\end{itemize}

\subsubsection{I Quattro Lobi della Corteccia}

La corteccia cerebrale è anatomicamente suddivisa in \textbf{quattro lobi principali}:

\begin{enumerate}
    \item \textbf{Lobo frontale}: funzioni esecutive, pianificazione, produzione del linguaggio (area di Broca), controllo motorio
    \item \textbf{Lobo parietale}: elaborazione sensoriale somatica, integrazione spaziale, via visiva dorsale
    \item \textbf{Lobo occipitale}: elaborazione visiva primaria (V1) e secondaria
    \item \textbf{Lobo temporale}: elaborazione uditiva, comprensione del linguaggio (area di Wernicke), via visiva ventrale, memoria
\end{enumerate}

\noindent Il \textbf{solco centrale} (\textit{central sulcus}) è il solco che separa le cortecce sensoriali (parte posteriore) dalle cortecce motorie (parte anteriore).




\noindent Il cervello è descritto secondo tre assi principali:

\begin{table}[htbp]
    \centering
    \renewcommand{\arraystretch}{1.3}
    \begin{tabularx}{\textwidth}{|l|X|}
    \hline
    \textbf{Asse} & \textbf{Terminologia} \\
    \hline
    Anteriore–Posteriore & \textbf{Rostrale} (anteriore) — \textbf{Caudale} (posteriore) \\
    Interno–Esterno & \textbf{Mediale} (interno) — \textbf{Laterale} (esterno) \\
    Alto–Basso & \textbf{Dorsale} (superiore) — \textbf{Ventrale} (inferiore) \\
    \hline
    \end{tabularx}
\end{table}







\subsection{Il Cervelletto}

Il \textbf{cervelletto} (\textit{cerebellum}, letteralmente "piccolo cervello") merita una trattazione separata per la sua importanza funzionale e strutturale.



\begin{itemize}
    \item \textbf{Coordinazione motoria}: integra informazioni propriocettive, vestibolari e visive per produrre movimenti fluidi e coordinati
    \item \textbf{Predizione}: genera modelli interni predittivi del comportamento motorio — anticipa le conseguenze sensoriali dei movimenti
    \item \textbf{Apprendimento motorio}: è il substrato principale per l'apprendimento di sequenze motorie complesse (strumento musicale, sport, scrittura)
\end{itemize}

\begin{tcolorbox}[colback=gray!5, colframe=gray!40, arc=2mm, boxrule=0.5pt]
Il cervelletto è come un "secondo cervello" : un sistema computazionale autonomo con propri circuiti e propria logica elaborativa, distinto dalla corteccia cerebrale ma in costante comunicazione con essa.
\end{tcolorbox}


\section{Scoperte Storiche}

\subsection{Paul Broca e la Produzione del Linguaggio}

Nel XIX secolo, il neurologo francese \textbf{Paul Broca} studiò un paziente affetto da sifilide celebrale. Questo paziente presentava un deficit molto specifico: era in grado di \textbf{comprendere perfettamente il linguaggio} degli altri, ma era in grado di pronunciare \textbf{una sola parola}: "Tan". Per questa ragione nella letteratura medica il paziente è passato alla storia come "\textbf{Tan}".

Alla morte del paziente, l'\textbf{autopsia} rivelò un danno circoscritto a una specifica area del \textbf{lobo frontale sinistro}: il \textbf{giro inferofrontale}. Questa regione, ora nota come \textbf{area di Broca}, è responsabile della \textbf{produzione del linguaggio}, la capacità di formare ed articolare parole.

\subsection{Carl Wernicke e la Comprensione del Linguaggio}

Il neurologo tedesco \textbf{Carl Wernicke} studiò un paziente che aveva subito un ictus e presentava un deficit complementare: era in grado di parlare fluentemente, ma le sue frasi erano prive di senso, non era in grado di \textbf{comprendere} il linguaggio altrui.

L'analisi anatomica rivelò un'area danneggiata \textbf{diversa} da quella di Broca, localizzata nel \textbf{lobo temporale sinistro}, in prossimità della \textbf{scissura silviana} (\textit{lateral sulcus}). Questa regione, ora nota come \textbf{area di Wernicke}, è responsabile della \textbf{comprensione del linguaggio}.

\begin{tcolorbox}[colback=gray!5, colframe=gray!40, arc=2mm, boxrule=0.5pt]
 Il caso di Broca e Wernicke costituisce una delle prime e più solide evidenze che la corteccia cerebrale non è un substrato omogeneo, bensì è \textbf{modulare}: aree diverse sono specializzate per funzioni diverse. La doppia dissociazione (produzione vs. comprensione) è particolarmente convincente perché dimostra che le due funzioni sono \textbf{indipendenti} dal punto di vista anatomico.
\end{tcolorbox}



\subsection{Korbinian Brodmann e la Mappa Corticale}

Il neuroanatomista tedesco \textbf{Korbinian Brodmann} (inizio del XX secolo) applicò sistematicamente la \textbf{colorazione di Golgi}, tecnica scoperta dall'italiano \textbf{Camillo Golgi}, per colorare selettivamente i neuroni corticali e rivelarne l'organizzazione microscopica.

Dall'analisi sistematica di sezioni corticali colorate, Brodmann scoprì che la corteccia non è istologicamente omogenea: identificò \textbf{52 aree distinte} (numerate da 1 a 52) sulla base delle differenze nell'\textbf{architettura neuronale}: diversa densità cellulare, diverso spessore degli strati, diversa morfologia dei neuroni.

Le principali aree di Brodmann includono:

\begin{table}[htbp]
    \centering
    \renewcommand{\arraystretch}{1.3}
    \begin{tabularx}{\textwidth}{|l|X|}
    \hline
    \textbf{Numero} & \textbf{Nome/Funzione} \\
    \hline
    1, 2, 3 & Corteccia somatosensoriale primaria \\
    4 & Corteccia motoria primaria \\
    6 & Corteccia premotoria e supplementare motoria \\
    17 & Corteccia visiva primaria (V1) \\
    18, 19 & Cortecce visive secondarie (V2, V3) \\
    22 & Corteccia uditiva secondaria (area di Wernicke) \\
    41, 42 & Corteccia uditiva primaria \\
    44, 45 & Corteccia del linguaggio (area di Broca) \\
    \hline
    \end{tabularx}
\end{table}

\begin{tcolorbox}[colback=gray!5, colframe=gray!40, arc=2mm, boxrule=0.5pt]
La scoperta di Brodmann è fondamentale perché stabilisce che la \textbf{modularità funzionale} (emersa dai casi clinici di Broca e Wernicke) ha un \textbf{substrato anatomico-istologico} preciso: aree diverse non solo svolgono funzioni diverse, ma hanno una struttura microscopica diversa.
\end{tcolorbox}



\subsection{Santiago Ramon y Cajal e la Dottrina Neuronale}

\textbf{Santiago Ramon y Cajal} e \textbf{Camillo Golgi} condivisero il \textbf{Premio Nobel per la Fisiologia o la Medicina nel 1906} per aver scoperto che la corteccia cerebrale è composta da semplici \textbf{elementi microscopici}, i neuroni, che comunicano tra loro.

I celebri disegni di Ramon y Cajal, realizzati con straordinaria precisione, mostrarono per la prima volta la struttura dettagliata dei neuroni: corpi cellulari colorati in nero, \textbf{dendriti} (fibre ramificate con \textbf{bottoni sinaptici}) e \textbf{assoni} che si diffondono nello spazio tridimensionale del tessuto.

Questo lavoro fondò la cosiddetta "\textbf{dottrina neuronale}": il principio secondo cui il cervello è composto di elementi fondamentali distinti (i neuroni) e che la cognizione risulta dalle loro interazioni. Si tratta del fondamento concettuale dell'intera neuroscienzas moderna.

\newpage
\section{Neurone e Comunicazione Sinaptica}

\subsubsection{Anatomia del Neurone}

Il neurone è l'elemento fondamentale del circuito cerebrale. È organizzato in \textbf{tre parti principali}:

\begin{enumerate}
    \item \textbf{Soma} (corpo cellulare): contiene il nucleo e il macchinario metabolico della cellula. È il "centro computazionale" che integra gli input in entrata.
    \item \textbf{Albero dendritico}: un sistema di ramificazioni che si estende dal soma. I dendriti \textbf{ricevono informazioni} da altri neuroni tramite i \textbf{bottoni sinaptici}, ovvero le strutture specializzate al termine degli assoni dei neuroni presinapici.
    \item \textbf{Assone}: una singola fibra lunga e sottile che si estende dal soma. Attraverso l'assone, il neurone \textbf{trasmette informazioni} ai neuroni postsinaptici. L'assone può ramificarsi terminalmente, contattando molti neuroni diversi.
\end{enumerate}


\subsubsection{Il Neurone come Macchina Elettrochimica}

I neuroni sono essenzialmente \textbf{macchine elettrochimiche}. Il \textbf{doppio strato lipidico} della membrana cellulare è impermeabile agli ioni e separa l'interno del neurone (con alta concentrazione di $K^+$ e bassa di $Na^+$) dall'esterno (con alta concentrazione di $Na^+$ e bassa di $K^+$).

Questa differenza di concentrazione ionica genera una differenza di \textbf{potenziale elettrico} attraverso la membrana: il cosiddetto \textbf{potenziale di membrana a riposo}, tipicamente intorno a $V_{\text{rest}} \approx -70$ mV (interno negativo rispetto all'esterno).

Il neurone funziona quindi come un \textbf{circuito elettrico} biologico: la membrana è il condensatore, i canali ionici sono le resistenze variabili, e le correnti ioniche sono le sorgenti di corrente.

\subsubsection{Il Potenziale d'Azione (Spike)}

Il meccanismo fondamentale della comunicazione neuronale è il \textbf{potenziale d'azione} (o \textit{spike}):

\begin{enumerate}
    \item Il neurone \textbf{integra} le correnti dendritiche attraverso il soma
    \item Quando il potenziale di membrana supera una soglia critica $V_{\text{thresh}}$, avviene una brusca \textbf{depolarizzazione stereotipata}: lo spike
    \item Questa depolarizzazione è dovuta all'apertura rapida di \textbf{canali ionici voltaggio-dipendenti} (principalmente canali al sodio)
    \item Il potenziale d'azione è un segnale \textbf{binario} e \textbf{stereotipato}: la sua ampiezza e forma sono sempre le stesse, l'informazione è codificata nella \textbf{frequenza di scarica}, non nell'ampiezza
\end{enumerate}

I potenziali d'azione viaggiano lungo gli assoni e vengono consegnati ai neuroni postsinaptici attraverso le sinapsi.

\subsubsection{La Sinapsi e l'EPSP}

La comunicazione tra neuroni avviene nelle \textbf{sinapsi}. Il meccanismo dettagliato è il seguente:

\begin{enumerate}
    \item Il potenziale d'azione del neurone \textbf{presinaptico} raggiunge il terminale assonale
    \item L'arrivo dello spike apre \textbf{canali ionici} nel terminale assonale
    \item Vescicole contenenti \textbf{neurotrasmettitori} si fondono con la membrana presinaptica e rilasciano i neurotrasmettitori nella \textbf{fessura sinaptica} (\textit{synaptic cleft})
    \item I neurotrasmettitori diffondono attraverso la fessura e si legano a recettori sulla membrana \textbf{postsinaptica}
    \item Questo legame apre canali ionici nel neurone postsinaptico, generando un afflusso di corrente
    \item La corrente produce una variazione del potenziale di membrana postsinaptico: il \textbf{potenziale postsinaptico eccitatorio} (\textbf{EPSP}, \textit{Excitatory Post-Synaptic Potential})
\end{enumerate}

La registrazione sperimentale mostra: potenziale d'azione presinaptico $\rightarrow$ propagazione lungo l'assone $\rightarrow$ rilascio di neurotrasmettitori $\rightarrow$ corrente nel neurone postsinaptico $\rightarrow$ EPSP.



\section{Cellule Gliali e Mielinizzazione}

\subsubsection{Le Cellule Gliali}

Le \textbf{cellule gliali} (\textit{glial cells}, dal greco "glia" = colla) costituiscono circa l'\textbf{80\% delle cellule cerebrali}. Per lungo tempo sono state considerate semplice infrastruttura di supporto ai neuroni, da cui il nome "colla". Tuttavia, la ricerca recente ha rivelato che svolgono funzioni importanti e specifiche.

I principali sottotipi di cellule gliali includono:

\begin{itemize}
    \item \textbf{Astrociti}: regolano la concentrazione ionica nell'ambiente extracellulare (fondamentale per il corretto funzionamento dei neuroni); stabiliscono e mantengono la \textbf{barriera emato-encefalica} (\textit{blood-brain barrier}), che protegge il cervello da sostanze tossiche presenti nel sangue
    \item \textbf{Cellule di Schwann} (nel sistema nervoso periferico) / \textbf{Oligodendrociti} (nel sistema nervoso centrale): producono la \textbf{guaina mielinica} che avvolge e isola gli assoni
\end{itemize}


\subsubsection{La Mielina}

La \textbf{mielina} è una sostanza lipidica bianca che avvolge gli assoni a intervalli regolari, lasciando liberi piccoli segmenti detti \textbf{nodi di Ranvier}. La mielina svolge due funzioni essenziali:

\begin{enumerate}
    \item \textbf{Isolamento elettrico}: impedisce la dissipazione della corrente lungo l'assone, evitando che il segnale si disperda
    \item \textbf{Conduzione saltatoria}: il potenziale d'azione "salta" da un nodo di Ranvier al successivo (\textit{conduzione saltatoria}), rendendo la propagazione molto più rapida rispetto a un assone non mielinizzato
\end{enumerate}

Grazie alla mielinizzazione, i potenziali d'azione si propagano lungo l'assone a velocità di decine di m/s, percorrendo anche distanze di centimetri in pochi millisecondi.


\section{Organizzazione Corticale}

\subsubsection{La Struttura a Strati della Corteccia}

La corteccia cerebrale è organizzata in uno strato sottile di \textbf{materia grigia superficiale}, sotto cui si trova la \textbf{materia bianca profonda} contenente gli assoni mielinizzati. 
La corteccia cerebrale è sottile (pochi millimetri di spessore) ma copre un'area molto estesa grazie alle pieghe (\textit{giri} e \textit{solchi}).

La colorazione di Golgi (che colora solo una piccola percentuale di neuroni in modo casuale, ma li rende completamente visibili) mostra che i neuroni sono organizzati in modo non uniforme nello spessore della corteccia: si osservano \textbf{diversi tipi e dimensioni cellulari a diverse profondità}.

La colorazione di \textbf{Nissl} (che colora i corpi cellulari di tutti i neuroni) ha confermato e precisato questa osservazione, permettendo il conteggio sistematico delle cellule per strato.

\subsubsection{I Sei Strati Corticali}

La corteccia cerebrale (neocorteccia) è universalmente organizzata in \textbf{sei strati}, numerati dall'esterno verso l'interno:

\begin{table}[htbp]
    \centering
    \renewcommand{\arraystretch}{1.3}
    \begin{tabularx}{\textwidth}{|l|l|X|}
    \hline
    \textbf{Strato} & \textbf{Profondità} & \textbf{Caratteristiche principali} \\
    \hline
    \textbf{I} & Superficiale (molecolare) & Pochi corpi cellulari, principalmente assoni e dendriti \\
    \textbf{II} & Granulare esterno & Piccole cellule granulari e piramidali \\
    \textbf{III} & Piramidale esterno & Cellule piramidali di medie dimensioni \\
    \textbf{IV} & Granulare interno & Piccole cellule stellate; principale destinazione degli \textbf{input talamici} \\
    \textbf{V} & Piramidale interno & Grandi \textbf{cellule piramidali di Betz} (corteccia motoria); principale output subcorticale \\
    \textbf{VI} & Fusiforme & Cellule fusiformi; connessioni con il talamo \\
    \hline
    \end{tabularx}
\end{table}



\subsection{L'Organizzazione a Minicolonne}

L'organizzazione della mielina nella corteccia rivela una struttura verticale caratteristica: \textbf{fasci verticali di fibre di output} che si proiettano dalla corteccia nella materia bianca sottostante. Questi fasci verticali corrispondono alle \textbf{minicolonne}: l'unità funzionale di base della corteccia. Una minicolonna è un gruppo di neuroni distribuiti attraverso tutti e sei gli strati corticali, orientati verticalmente, che:

\begin{itemize}
    \item Condividono lo stesso ruolo funzionale (rispondono agli stessi stimoli)
    \item Ricevono input dallo stesso territorio
    \item Proiettano verso le stesse aree target
\end{itemize}

Le minicolonne hanno un diametro di circa 50 $\mu$m e contengono tipicamente $\sim$80–100 neuroni in tutte le specie di mammiferi esaminate.

\subsection{La Cellula Piramidale come Prototipo}

La \textbf{ricostruzione tridimensionale della cellula piramidale} mostra un'organizzazione che riflette perfettamente l'architettura a minicolonne:

\begin{itemize}
    \item \textbf{Dendriti apicali}: si estendono verticalmente verso gli strati superficiali (I–III), ricevendo input dalle cortecce associative
    \item \textbf{Dendriti basali}: si estendono orizzontalmente negli strati profondi (IV–V), ricevendo input locali
    \item \textbf{Corpo cellulare} (soma): localizzato nello strato V (cellule piramidali grandi) o III (cellule piramidali medie)
    \item \textbf{Assone}: si proietta verticalmente verso il basso, attraversa la materia bianca e raggiunge i target remoti
\end{itemize}

Il flusso di informazioni nella minicolonna è quindi prevalentemente \textbf{verticale}: input dagli strati superficiali $\rightarrow$ integrazione nel soma $\rightarrow$ output attraverso la materia bianca verso gli strati profondi.

\subsection{Tipi Cellulari Principali della Corteccia}

La corteccia è composta da due grandi classi di neuroni:

\begin{enumerate}
    \item \textbf{Neuroni eccitatori} (circa 80\% della popolazione): principalmente \textbf{cellule piramidali}, usano il \textbf{glutammato} come neurotrasmettitore. Hanno corpi cellulari a forma piramidale con un lungo dendrite apicale. Proiettano a lunga distanza (anche in aree lontane e nell'emisfero controlaterale attraverso il corpo calloso).
    \item \textbf{Neuroni inibitori} (circa 20\% della popolazione): cellule \textbf{GABAergiche} (\textit{gamma-aminobutyric acid}), che usano il GABA come neurotrasmettitore. Esistono molti sottotipi, tra i più importanti:
    \begin{itemize}
        \item \textbf{Cellule a candelabro}: innervano specificamente il segmento iniziale dell'assone delle piramidali, il punto dove si generano gli spike, esercitando un controllo inibitorio molto efficace
        \item \textbf{Cellule a cesto}: innervano i corpi cellulari delle cellule piramidali
    \end{itemize}
\end{enumerate}

I due tipi di neuroni corticali producono effetti opposti sul neurone postsinaptico:

\begin{itemize}
    \item \textbf{Neuroni eccitatori} (glutammatergici): il legame del glutammato con i recettori AMPA/NMDA genera un afflusso di ioni positivi ($Na^+$, $Ca^{2+}$), producendo una \textbf{deflessione depolarizzante} (positiva) del potenziale di membrana. Questo segnale è detto \textbf{EPSP} (\textit{Excitatory Post-Synaptic Potential}).
    \item \textbf{Neuroni inibitori} (GABAergici): il legame del GABA con i recettori GABA-A apre canali per $Cl^-$ o $K^+$, producendo una \textbf{deflessione iperpolarizzante} (negativa) del potenziale di membrana. Questo segnale è detto \textbf{IPSP} (\textit{Inhibitory Post-Synaptic Potential}).
\end{itemize}

La combinazione di EPSP e IPSP integrata nel soma determina se il neurone raggiungerà o meno la soglia di scarica.

\subsection{La Legge di Dale}

Un principio fondamentale della trasmissione sinaptica è la \textbf{Legge di Dale}:

\begin{tcolorbox}[colback=colorE!40, colframe=colorE, arc=2mm, boxrule=0.5pt]
\textit{Se un neurone presinaptico è eccitatorio, tutti i suoi target postsinaptici ricevono depolarizzazione. Se è inibitorio, tutti i suoi target ricevono iperpolarizzazione.}
\end{tcolorbox}

In altre parole, ogni neurone libera \textbf{un solo tipo} di neurotrasmettitore principale in tutte le sue terminazioni sinaptiche. Un neurone non può essere eccitatorio per alcuni target e inibitorio per altri.

 La Legge di Dale è una semplificazione utile a livello di modellistica. In realtà, molti neuroni co-rilasciano più di un neurotrasmettitore, ma il principio di base rimane valido per la grande maggioranza dei casi.


\newpage

\section{Organizzazione Sensoriale e Motoria della Corteccia}

L'organizzazione corticale rispetta un preciso \textbf{vincolo fisico}: le aree di input sensoriale e di output motorio si trovano in prossimità del \textbf{solco centrale}, mentre l'elaborazione associativa complessa avviene nelle regioni più distanti.

Le principali aree sensoriali primarie:

\begin{table}[htbp]
    \centering
    \renewcommand{\arraystretch}{1.3}
    \begin{tabularx}{\textwidth}{|l|X|}
    \hline
    \textbf{Modalità sensoriale} & \textbf{Localizzazione corticale} \\
    \hline
    \textbf{Somatosensazione} (tatto, propriocezione) & Corteccia somatosensoriale primaria (postcentrale, strati 1-2-3 di Brodmann) \\
    \textbf{Visione} & Corteccia visiva primaria V1, lobo occipitale (caudale) \\
    \textbf{Udito} & Corteccia uditiva primaria, lobo temporale superiore (vicino all'area di Wernicke) \\
    \textbf{Olfatto} & Bulbo olfattivo (regione filogeneticamente \textbf{antica}, solo tre strati — non neocorteccia) \\
    \hline
    \end{tabularx}
\end{table}

 Il bulbo olfattivo è particolare perché è una delle strutture cerebrali più antiche dal punto di vista evolutivo (\textit{allocorteccia}), con soli tre strati invece dei sei della neocorteccia. L'olfatto ha una connessione diretta con le strutture limbiche (amigdala, ippocampo) che processano memoria ed emozioni, da qui il potere evocativo degli odori.


\subsubsection{Cortecce Associative}

Le \textbf{cortecce associative} si trovano tra le aree sensoriali primarie e le aree motorie. Integrano informazioni provenienti da più modalità sensoriali:

\begin{tcolorbox}[colback=gray!5, colframe=gray!40, arc=2mm, boxrule=0.5pt]
\textbf{Esempio:} Quando percepiamo una mela, la corteccia visiva elabora forma e colore, la corteccia somatosensoriale elabora la consistenza (tattile), la corteccia olfattiva elabora il profumo, la corteccia gustativa il sapore. Le cortecce associative integrano tutte queste informazioni in un'unica rappresentazione coerente dell'oggetto "mela".
\end{tcolorbox}

\subsubsection{Cortecce Motorie}

Sul lato opposto del solco centrale rispetto alle cortecce sensoriali si trovano le \textbf{cortecce motorie}:

\begin{itemize}
    \item \textbf{Corteccia motoria primaria (M1)}: traduce i comandi motori in attivazione muscolare; movimento, tatto, parlato
    \item \textbf{Corteccia premotoria}: pianificazione dei movimenti
    \item \textbf{Area motoria supplementare (SMA)}: coordinazione di sequenze motorie complesse
\end{itemize}




\subsection{La Via Visiva}

L'elaborazione visiva segue un percorso ben definito dal recettore alla corteccia:

$$ \text{Retina} \rightarrow \text{Nervo ottico} \rightarrow \text{Nucleo Genicolato Laterale} \rightarrow \text{Corteccia Visiva Primaria} $$

\begin{itemize}
    \item \textbf{Retina}: i fotorecettori (coni e bastoncelli) traducono la luce in segnali elettrici. Le cellule gangliari della retina proiettano attraverso il nervo ottico.
    \item \textbf{Nucleo Genicolato Laterale}: struttura talamìca che funge da stazione di trasmissione intermedia. Riceve l'input dalla retina e proietta alla corteccia visiva.
    \item \textbf{Corteccia Visiva Primaria V1} (area 17 di Brodmann, \textit{corteccia striata}): il primo stadio di elaborazione corticale della visione. Il nome "striata" deriva dall'organizzazione a \textbf{strisce} visibile a occhio nudo nelle sezioni colorate.
\end{itemize}


\subsubsection{L'Organizzazione Retinotopica di V1}

Attraverso tecniche di \textbf{imaging ottico} (che misurano le variazioni di riflettanza della corteccia legate all'attività neuronale), è possibile mappare l'organizzazione di V1 in risposta a stimoli visivi.
I risultati mostrano che V1 è organizzata in modo \textbf{retinotopico}: esiste una corrispondenza biunivoca tra la posizione nello spazio visivo (retina) e la posizione sulla corteccia. Stimoli in posizioni diverse del campo visivo attivano strisce diverse di corteccia.

Il sistema di riferimento retinotopico usa \textbf{coordinate polari}:
\begin{itemize}
    \item \textbf{Distanza radiale dalla fovea}: codificata lungo un asse della corteccia
    \item \textbf{Posizione angolare} (0°–360°): codificata lungo l'altro asse
\end{itemize}

La \textbf{fovea} (centro del campo visivo, massima acuità) è rappresentata da un'area corticale sproporzionatamente grande rispetto alla periferia retinica (\textit{magnification factor}).

\subsubsection{Dominanza Oculare}

V1 mostra anche una organizzazione legata alla \textbf{dominanza oculare}: strisce adiacenti sono dominate alternativamente dall'occhio destro e dall'occhio sinistro. Questo arrangement è fondamentale per la \textbf{visione stereoscopica} (percezione della profondità), poiché permette il confronto delle informazioni provenienti dai due occhi.

\begin{tcolorbox}[colback=gray!5, colframe=gray!40, arc=2mm, boxrule=0.5pt]
Gli occhiali da cinema 3D sfruttano proprio questo meccanismo: proiettando immagini leggermente diverse ai due occhi, si genera artificialmente la disparità binoculare che il sistema visivo interpreta come profondità.
\end{tcolorbox}

\subsubsection{Selettività di Orientamento}

Oltre alla retinotopia e alla dominanza oculare, i neuroni di V1 mostrano \textbf{selettività di orientamento}: ogni neurone risponde preferenzialmente a barre o bordi orientati in una direzione specifica (es. verticale, orizzontale, 45°, 135°).

Ruotando la barra stimolo, si osserva che diverse regioni corticali rispondono a diversi orientamenti. Attraverso l'imaging ottico, è possibile mappare la distribuzione degli orientamenti preferiti sulla superficie di V1.

La mappa degli orientamenti in V1 rivela una struttura affascinante: i neuroni selettivi per diversi orientamenti sono disposti in \textbf{spirali} (\textit{pinwheels}) attorno a \textbf{punti di singolarità} al centro di ciascun pinwheel.

In corrispondenza del punto centrale di ogni pinwheel, tutti gli orientamenti (da 0° a 180°) sono rappresentati da neuroni collocati nelle sue immediate vicinanze. Procedendo radialmente verso l'esterno, l'orientamento preferito ruota progressivamente.

Un risultato notevole: la \textbf{densità dei centri di pinwheel} sulla superficie corticale è proporzionale a \textbf{$\pi$ (pi greco)}, una corrispondenza matematica sorprendente che ha attirato l'attenzione della comunità neuroscientifica e che suggerisce un principio organizzativo profondo nell'architettura corticale.

\begin{tcolorbox}[colback=gray!5, colframe=gray!40, arc=2mm, boxrule=0.5pt]
La relazione con $\pi$ non è un'astrazione matematica ma una proprietà misurabile sperimentalmente. Questo tipo di regolarità matematica in un sistema biologico è raro e suggerisce che l'organizzazione della corteccia visiva segua principi fisici precisi.
\end{tcolorbox}



\subsubsection{Gerarchia dell'Elaborazione Visiva}

L'informazione visiva è elaborata in modo \textbf{gerarchico}. V1 codifica le caratteristiche visive di \textbf{basso livello}:
\begin{itemize}
    \item Posizione nello spazio visivo
    \item Orientamento locale
    \item Occhio di provenienza
    \item Colore
\end{itemize}

Le aree visive successive (V2, V3, V4, V5/MT) e le aree associative elaborano caratteristiche progressivamente più complesse:

$$ V1 \xrightarrow{\text{bordi, orientamenti}} V2 \xrightarrow{\text{angoli, forme semplici}} V4 \xrightarrow{\text{colore, forme}} IT \xrightarrow{\text{oggetti, volti}} $$

\begin{tcolorbox}[colback=gray!5, colframe=gray!40, arc=2mm, boxrule=0.5pt]
Questa organizzazione gerarchica è notevolmente simile all'architettura delle \textbf{reti neurali convoluzionali} (\textit{Convolutional Neural Networks, CNN}) nel machine learning, un esempio di come i principi computazionali biologici abbiano ispirato l'ingegneria dell'intelligenza artificiale.
\end{tcolorbox}


Dopo V1, l'informazione visiva si biforca in due principali \textbf{vie di elaborazione} che proiettano verso regioni diverse della corteccia:

\begin{enumerate}
    \item \textbf{Via Ventrale} (\textit{"What" pathway}, "cosa"): proietta verso il \textbf{lobo temporale} inferiore. È responsabile del \textbf{riconoscimento degli oggetti} — identificare cosa si sta guardando (forma, colore, identità).
    \item \textbf{Via Dorsale} (\textit{"Where/How" pathway}, "dove/come"): proietta verso il \textbf{lobo parietale}. È responsabile della \textbf{localizzazione spaziale degli oggetti} e del controllo visuomotorio — dove si trova l'oggetto e come interagirci.
\end{enumerate}

\subsubsection{La Corteccia Infratemporale (IT) e il Riconoscimento degli Oggetti}

La \textbf{corteccia infratemporale} (\textit{inferotemporal cortex, IT}) si trova alla fine della via ventrale e rappresenta il livello più alto dell'elaborazione visiva per il riconoscimento degli oggetti.

Un esperimento condotto da un gruppo di ricerca giapponese ha dimostrato una proprietà notevole della corteccia IT: quando gli animali osservano un oggetto specifico (ad esempio un \textbf{estintore}), si attivano \textbf{regioni segregate} e specifiche della corteccia IT — visibili come cerchi rossi nelle mappe di attivazione.

Aspetto cruciale: queste stesse regioni si attivano in risposta a variazioni dell'oggetto:
\begin{itemize}
    \item Cambio di colore (rosso vs. nero)
    \item Presenza di ombre
    \item Variazioni di illuminazione
\end{itemize}

La risposta è quindi \textbf{invariante} rispetto a queste trasformazioni superficiali. Questo rivela che la corteccia IT codifica una \textbf{rappresentazione astratta} dell'oggetto, non le sue proprietà fisiche specifiche.

\subsubsection{Codifica Distribuita}

La rappresentazione nella corteccia IT è \textbf{distribuita}: non è un singolo neurone a codificare "estintore", ma una \textbf{popolazione di neuroni} il cui pattern di attività collettivo rappresenta l'oggetto.

La \textbf{modularità} della risposta (cerchi diversi per oggetti diversi) riflette il meccanismo di \textbf{astrazione}: diverse categorie di oggetti attivano popolazioni neuronali diverse, ma ogni singola popolazione risponde a variazioni dello stesso oggetto. Questo è il meccanismo biologico che permette di riconoscere una mela sia rossa che verde, sia in piena luce che in ombra.

\begin{tcolorbox}[colback=gray!5, colframe=gray!40, arc=2mm, boxrule=0.5pt]
Un esperimento aggiuntivo ha mostrato che oggetti simili ma privi di manico attivano popolazioni neuronali diverse (contorni verdi nelle mappe) rispetto agli oggetti completi. Questo dimostra che anche dettagli funzionalmente rilevanti (il manico permette di "afferrare") sono codificati nella rappresentazione neuronale.
\end{tcolorbox}


\subsection{Principi di Astrazione nelle Cortecce Associative}

Le cortecce associative costruiscono \textbf{rappresentazioni astratte} delle informazioni sensoriali. Il processo è:

\begin{enumerate}
    \item Le cortecce sensoriali primarie codificano caratteristiche fisiche concrete (bordi, frequenze sonore, pressioni tattili)
    \item Ogni livello gerarchico combina e astrae le caratteristiche del livello inferiore
    \item Nelle cortecce associative di ordine superiore, i neuroni rispondono a concetti astratti (oggetti, categorie, relazioni) indipendentemente dalle proprietà fisiche superficiali dello stimolo
\end{enumerate}

Questa capacità di \textbf{generalizzazione} è essenziale per la sopravvivenza: permette di riconoscere un predatore indipendentemente dalla sua posizione, illuminazione, distanza; di identificare il cibo anche in condizioni parziali.



I calcoli cognitivi rilevanti per la sopravvivenza sono notevolmente \textbf{simili tra le specie}:

\begin{itemize}
    \item Trovare cibo (riconoscimento, localizzazione, perseguimento)
    \item Evitare pericoli (riconoscimento dei predatori, risposta di fuga)
    \item Riproduzione (riconoscimento del partner, comportamenti sociali)
\end{itemize}

Queste funzioni sono condivise da mammiferi, uccelli e probabilmente dai dinosauri. Le differenze cognitive tra le specie riguardano principalmente le \textbf{capacità aggiuntive}: matematica, musica, linguaggio, pensiero astratto — capacità che gli esseri umani hanno sviluppato in misura molto maggiore rispetto ad altri animali.



\section{Scaling del Cervello tra Specie e Proprietà Small-World}



Le dimensioni del cervello variano enormemente tra le specie di mammiferi, da pochi grammi nel topo a diversi chilogrammi nella balena, eppure l'\textbf{organizzazione di base} è notevolmente conservata: stessi lobi, stessa struttura corticale a sei strati, stesso tipo di neuroni, stessa divisione in materia grigia e bianca.

\begin{tcolorbox}[colback=gray!5, colframe=gray!40, arc=2mm, boxrule=0.5pt]
Questa conservazione evolutiva suggerisce che l'architettura corticale a sei strati organizzata in minicolonne rappresenta una soluzione computazionale ottimale per l'elaborazione dell'informazione nei vertebrati. L'evoluzione ha "scoperto" questa architettura una volta sola e l'ha mantenuta per centinaia di milioni di anni.
\end{tcolorbox}

Osservando la variazione del volume cerebrale tra le specie di mammiferi, dal topo all'elefante, si osservano variazioni di \textbf{diversi ordini di grandezza}:

\begin{itemize}
    \item Il \textbf{volume della materia grigia} (corpi neuronali) aumenta enormemente
    \item Il \textbf{volume della materia bianca} (fibre mielinizzate) aumenta in modo analogo
\end{itemize}

Questo pone un importante \textbf{problema ingegneristico}: se tutti i neuroni fossero connessi con tutti gli altri (\textit{connettività all-to-all}), la quantità di materia bianca necessaria crescerebbe in modo \textbf{quadratico} con il numero di neuroni $N$:

$$ V_{\text{bianca}} \propto N^2 $$

Tuttavia, sperimentalmente si osserva che la materia bianca scala in modo sublineare o lineare con la materia grigia, non quadratico. Questo implica che la \textbf{connettività deve essere selettiva}: non tutti i neuroni sono connessi con tutti gli altri.

\subsubsection{Il Numero di Sinapsi per Neurone è Conservato}

Neuroanatomisti spagnoli (studenti, ratti, topi) hanno eseguito conteggi sistematici dei \textbf{bottoni sinaptici} sui neuroni in diversi strati corticali e in diverse specie.

Il risultato è sorprendente: il numero di sinapsi per neurone è \textbf{approssimativamente costante} tra le specie:

$$ N_{\text{sinapsi/neurone}} \approx 10.000 - 20.000 $$

Questo numero rimane costante nonostante l'enorme variazione nel \textbf{numero totale di neuroni}:

\begin{table}[htbp]
    \centering
    \renewcommand{\arraystretch}{1.3}
    \begin{tabularx}{\textwidth}{|l|X|}
    \hline
    \textbf{Specie} & \textbf{Neuroni corticali approssimativi} \\
    \hline
    Topo & $\sim$30 milioni \\
    Ratto & $\sim$70 milioni \\
    Gatto & $\sim$300 milioni \\
    Uomo & $\sim$10–15 miliardi \\
    \hline
    \end{tabularx}
\end{table}

\begin{tcolorbox}[colback=gray!5, colframe=gray!40, arc=2mm, boxrule=0.5pt]
Il vincolo evolutivo è chiaro: minimizzare la distanza di connettività (per ridurre i costi metabolici e la quantità di materia bianca), preservare gli hub di comunicazione importanti, e sacrificare le connessioni periferiche meno critiche.
\end{tcolorbox}

\subsubsection{Proprietà Small-World della Connettività}

La connettività sinaptica nel cervello non è né completamente casuale né completamente regolare. Presenta invece le caratteristiche di una \textbf{rete small-world}:

\begin{itemize}
    \item Alcuni neuroni (gli \textbf{hub}) sono connessi con un numero molto elevato di altri neuroni
    \item La maggior parte dei neuroni (periferici) ha poche connessioni
    \item Nonostante la bassa connettività media, l'informazione si diffonde in modo efficiente attraverso gli hub
\end{itemize}

Questa architettura permette la \textbf{diffusione efficiente delle informazioni} senza richiedere connettività completa (e quindi senza richiedere una quantità quadratica di materia bianca).

Il numero chiave di $\sim$10.000 sinapsi per neurone è conservato tra le specie: è il compromesso evolutivo ottimale tra connettività (per la ricchezza dell'elaborazione) e costo metabolico/anatomico.

\begin{tcolorbox}[colback=gray!5, colframe=gray!40, arc=2mm, boxrule=0.5pt]
I roditori dimostrano che cognizione complessa è possibile anche con circa 100 volte meno neuroni degli esseri umani. Questo suggerisce che la \textbf{qualità} dell'architettura connettiva (la struttura della rete) conta più della pura quantità di neuroni. In qualche modo, anche reti molto più piccole riescono a implementare le funzioni cognitive essenziali.
\end{tcolorbox}

\newpage
\section{Informazioni sull'Esame}

Il docente ha fornito informazioni dettagliate sul formato dell'esame al termine della lezione.

L'esame è \textbf{orale} e strutturato nel modo seguente:

\begin{enumerate}
    \item \textbf{Un argomento a scelta dello studente}: lo studente sceglie liberamente un tema del corso su cui essere esaminato
    \item \textbf{Due argomenti a scelta del docente}: il docente seleziona due temi aggiuntivi
\end{enumerate}

\begin{tcolorbox}[colback=gray!5, colframe=gray!40, arc=2mm, boxrule=0.5pt]
Non tutto il materiale del corso sarà oggetto d'esame. In particolare, i contenuti di questa prima lezione, l'\textbf{introduzione e l'anatomia cerebrale di base}, non appariranno nelle domande a scelta del docente, in quanto considerati conoscenze di prerequisito o di contesto generale.
\end{tcolorbox}

Il focus principale dell'esame sarà sul \textbf{corpo principale del corso}, con particolare attenzione al:

\begin{itemize}
    \item \textbf{Modello di Hodgkin-Huxley} (descrizione quantitativa del potenziale d'azione)
    \item Neuroni spiking semplificati
    \item Population firing rates e dinamica di popolazione
    \item Dinamica non lineare e funzioni cognitive
\end{itemize}

\newpage


\section*{Riepilogo dei Concetti Fondamentali}
\addcontentsline{toc}{section}{Riepilogo dei Concetti Fondamentali}

\begin{table}[htbp]
    \centering
    \renewcommand{\arraystretch}{1.3}
    \begin{tabularx}{\textwidth}{|l|X|}
    \hline
    \textbf{Concetto} & \textbf{Descrizione sintetica} \\
    \hline
    \textbf{Dottrina neuronale} & Il cervello è composto di elementi fondamentali (neuroni); la cognizione emerge dalle loro interazioni \\
    \textbf{Potenziale d'azione} & Segnale binario stereotipato generato quando il potenziale di membrana supera la soglia; informazione codificata nella frequenza \\
    \textbf{EPSP / IPSP} & Potenziali postsinaptici eccitatori (glutammato) e inibitori (GABA); la loro somma determina la scarica del neurone \\
    \textbf{Legge di Dale} & Ogni neurone rilascia un solo tipo di neurotrasmettitore principale in tutte le sue terminazioni \\
    \textbf{Aree di Brodmann} & 52 aree corticali distinte per architettura istologica, corrispondenti a funzioni diverse \\
    \textbf{Minicolonna} & Unità funzionale verticale della corteccia: $\sim$80–100 neuroni attraverso tutti e sei gli strati con la stessa funzione \\
    \textbf{Retinotopia} & Corrispondenza punto-a-punto tra posizione nello spazio visivo e posizione sulla corteccia visiva \\
    \textbf{Pinwheel} & Organizzazione a spirale dei neuroni selettivi per l'orientamento in V1; densità proporzionale a $\pi$ \\
    \textbf{Via ventrale / dorsale} & Due vie visive: "cosa" (lobo temporale) e "dove/come" (lobo parietale) \\
    \textbf{Codifica distribuita} & Gli oggetti sono rappresentati da popolazioni di neuroni, non da singoli neuroni \\
    \textbf{Rete small-world} & Architettura connettiva con hub ad alta connettività: diffusione efficiente senza connettività completa \\
    \textbf{$\sim$10.000 sinapsi/neurone} & Numero conservato tra tutte le specie di mammiferi; compromesso tra connettività e costo metabolico \\
    \hline
    \end{tabularx}
\end{table}

\chapter{Lezione 2}

\textbf{Argomenti:} Metodi sperimentali in neuroscienze, EEG, ERP, MEG, TMS, PET, fMRI, segnale BOLD, imaging ottico VSD, calcium imaging, ECoG, patch clamp, optogenetica, campo elettrico extracellulare, local field potential, SUA, MUA

\vspace{0.5cm}
\noindent\rule{\textwidth}{0.4pt}

\section{Metodi Sperimentali in Neuroscienze}


La seconda parte dell'introduzione al corso è dedicata ai \textbf{metodi sperimentali} attualmente disponibili nei laboratori di neuroscienze per investigare il comportamento delle reti cerebrali. Il punto di partenza è un grafico che mostra la \textbf{risoluzione temporale e spaziale} dei diversi approcci strumentali.

Sull'asse temporale, la gamma di copertura si estende dai \textbf{millisecondi} fino a giorni e mesi: è quindi possibile seguire l'attività neuronale a una risoluzione paragonabile alla durata di un potenziale d'azione, ma anche monitorare l'evoluzione di processi plastici su lunghe scale temporali. Sull'asse spaziale, invece, la finestra è altrettanto vasta: si va dalla \textbf{sinapsi singola}, con una risoluzione dell'ordine di frazioni di micron, fino all'intero cervello esplorato nella sua globalità.

Va notato che il grafico originale non include il \textbf{calcium imaging} con microscopi avanzati, poiché si tratta di una tecnologia più recente. Con questa tecnica, utilizzando microscopi a due fotoni, è oggi possibile seguire in vivo l'attività di singole sinapsi, raggiungendo risoluzioni spaziali nell'ordine di frazioni di micron.


Ogni tecnica occupa una posizione specifica nel piano risoluzione spaziale–risoluzione temporale. Poiché nessun metodo copre simultaneamente tutta la gamma di scale, l'indagine neuroscientifica richiede la combinazione di approcci complementari, selezionati in funzione della domanda scientifica. Le tecniche non invasive — che non richiedono l'apertura del cranio — tendono a collocarsi nella parte superiore del grafico (bassa risoluzione spaziale, buona risoluzione temporale o viceversa). Le tecniche invasive permettono, al prezzo di una procedura chirurgica, di guadagnare in risoluzione spaziale e di accedere al segnale neuronale direttamente alla sua sorgente.


\section{Elettroencefalografia (EEG) e Potenziali Evento-Correlati (ERP)}

\subsubsection{L'Esperienza di Berger e il Principio Fisico}

Per investigare il cervello nella sua interezza, a bassa risoluzione spaziale, l'approccio più noto è la registrazione dell'attività elettrica spontaneamente prodotta dal \textbf{sistema nervoso centrale}. I neuroni generano messaggi sotto forma di \textbf{potenziali d'azione} (spike), che perturbano il campo elettrico — e in misura minore quello magnetico — dello spazio circostante. Quando l'attività è coordinata tra una grande popolazione di neuroni appartenenti alla stessa popolazione funzionale (ad esempio, neuroni che codificano la stessa informazione), queste perturbazioni del campo elettrico si sovrappongono in modo coerente e possono essere rilevate da \textbf{elettrodi posizionati sul cranio}.

Questo principio fu sfruttato per la prima volta negli anni '30 del XX secolo da \textbf{Hans Berger}, il cui lavoro portò alla registrazione della prima \textbf{elettroencefalografia (EEG)}: la registrazione della attività elettrica del cervello attraverso il cuoio capelluto.

\subsubsection{I Potenziali Evento-Correlati (ERP)}

Il termine \textbf{ERP} (Event-Related Potentials, potenziali evento-correlati) descrive una variante dell'EEG in cui si studia la risposta elettrica cerebrale evocata da una stimolazione esogena — tipicamente sensoriale. La stimolazione può essere visiva, uditiva o somatosensoriale.



\section{Magnetoencefalografia (MEG)}

\subsubsection{Il Principio SQUID e la Rilevazione del Campo Magnetico}

Un approccio più sofisticato è la \textbf{magnetoencefalografia (MEG)}, che sfrutta la misurazione delle variazioni del \textbf{campo magnetico} prodotto dall'attività elettrica cerebrale. Quando fluiscono correnti elettriche nei neuroni, si genera anche un campo magnetico associato. Se questa generazione di campo magnetico è coerente tra molti neuroni attivi simultaneamente, il segnale risultante può essere rilevato da sensori estremamente sensibili.

Il dispositivo di rilevazione alla base della MEG è lo \textbf{SQUID} (Superconducting QUantum Interference Device, interferometro quantistico superconduttore), che sfrutta il fenomeno dell'interferenza quantistica nei circuiti superconduttori per misurare variazioni di flusso magnetico straordinariamente piccole. 

\subsubsection{Necessità di Isolamento Ambientale}

A differenza del campo elettrico, il campo magnetico prodotto dall'attività neuronale è estremamente tenue. Per poterlo rilevare è necessario operare in un ambiente di misura fortemente isolato, dove qualsiasi perturbazione magnetica esterna (traffico, elettrodomestici, attrezzature elettriche) sia eliminata. La misura viene quindi eseguita in \textbf{camere schermature magneticamente} (camere di Faraday per il campo magnetico), che isolano il soggetto dal rumore di fondo.

\subsubsection{Caratteristiche della MEG}

Come l'EEG, anche la MEG è una tecnica \textbf{non invasiva}: non richiede alcuna apertura del cranio né l'impianto di elettrodi. La \textbf{risoluzione temporale} è eccellente, raggiungendo la scala dei millisecondi. La \textbf{risoluzione spaziale}, invece, è bassa, poiché il segnale registrato è la sovrapposizione dei contributi di molte sorgenti distribuite nel cervello — una sorta di grande media del generatore distribuito. Per questa ragione la MEG si colloca nella parte superiore dell'asse verticale del grafico di risoluzione.

\section{Stimolazione Magnetica Transcranica (TMS)}

Le tecniche discusse finora sono tutte di tipo \textbf{registrativo}: misurano passivamente l'attività cerebrale. Una prospettiva tipica del fisico è invece quella di \textbf{perturbare}: si dà un "kick" al sistema e si osserva la risposta, per comprendere le proprietà della rete. In neuroscienze, questo approccio è incarnato dalla \textbf{stimolazione magnetica transcranica (TMS)}, che consente di perturbare le reti corticali in modo non invasivo.



Il dispositivo TMS consiste in un \textbf{solenoide} posizionato esternamente al cranio. Quando si fa scorrere un impulso di corrente nel solenoide, questo genera un campo magnetico variabile nel tempo. Per la \textbf{legge di induzione di Faraday}, un campo magnetico variabile induce una forza elettromotrice e una corrente nel tessuto conduttivo sottostante — in questo caso, la corteccia cerebrale. Il flusso di corrente indotta aumenta l'\textbf{eccitabilità corticale} e porta all'attivazione dei neuroni corticali (firing spontaneo).

Se invece di un singolo impulso si eroga un \textbf{burst di impulsi} (sequenza rapida), il meccanismo è opposto: il network corticale viene \textbf{inibito} (down-regulation dell'eccitabilità), e l'attività spontanea viene soppressa.

\
La TMS può essere utilizzata in combinazione con l'EEG: si perturba la rete con la TMS e si osserva con l'EEG l'effetto di questa perturbazione sull'attività elettrica corticale in corso. Questo schema \textbf{perturbazione + registrazione} è uno strumento potente per caratterizzare le proprietà dinamiche delle reti corticali.

\section{Tomografia a Emissione di Positroni (PET)}

Un'ulteriore modalità non invasiva per investigare l'attività cerebrale si basa sulla misurazione dei \textbf{prodotti del metabolismo neuronale}. Le reti neurali più attive consumano più energia: questa correlazione costituisce la base della \textbf{tomografia a emissione di positroni (PET)}.

Nella PET si somministra al soggetto un \textbf{tracciante radioattivo} (solitamente glucosio marcato con un isotopo emettitore di positroni, come il $^{18}$F-fluorodeossiglucosio). I neuroni più attivi metabolicamente lo captano in misura maggiore. Il decadimento positronico del tracciante produce fotoni gamma rilevabili dall'esterno, permettendo di ricostruire una mappa tridimensionale dell'attività metabolica cerebrale.


La PET presenta una \textbf{risoluzione spaziale e temporale molto bassa}: nel grafico risoluzione spazio-temporale, essa occupa l'angolo in alto a destra, indicando scarsa capacità di localizzare sorgenti con precisione e di seguire variazioni rapide. Nonostante queste limitazioni, la PET è ampiamente utilizzata sia in ambito clinico (neurologia, oncologia) sia nella ricerca di base per lo studio del metabolismo cerebrale.



\section{Risonanza Magnetica Funzionale (fMRI)}

\subsubsection{I Protoni dell'Acqua come Sonde}

La tecnica non invasiva di maggiore impatto nella neuroscienze cognitive è la \textbf{risonanza magnetica funzionale (fMRI)}. Il suo funzionamento sfrutta le proprietà quantistiche dei \textbf{protoni dell'acqua} ($^{1}\mathrm{H}$), che costituiscono i candidati ideali per due ragioni: l'acqua rappresenta circa i tre quarti della massa corporea ed è pervasiva nel tessuto cerebrale.

I protoni dell'acqua possiedono un momento angolare intrinseco, lo \textbf{spin nucleare}, che in assenza di campo magnetico esterno è orientato casualmente. Questa casualità degli orientamenti produce un momento magnetico netto nullo e nessun segnale rilevabile.

\subsubsection{Allineamento degli Spin nel Campo Forte}

La fMRI sfrutta un \textbf{magnete principale} estremamente potente che genera un campo magnetico statico e uniforme $\mathbf{B}_0$. Quando il campo viene attivato, gli spin dei protoni tendono ad allinearsi lungo la direzione di $\mathbf{B}_0$ (allineamento parallelo o antiparallelo, con una lieve prevalenza del parallelo per ragioni termodinamiche). Più forte e più omogeneo è il campo $\mathbf{B}_0$, più coerente è l'orientamento degli spin e migliore è la qualità dell'immagine ricostruibile.

In ambito clinico si utilizzano tipicamente magneti da \textbf{1.5 Tesla}. I laboratori di ricerca in neuroscienze dispongono di magneti da \textbf{7 Tesla}, che consentono una risoluzione spaziale superiore al millimetro cubico (il \textbf{voxel} — pixel tridimensionale). Per la sperimentazione animale esistono magneti ancora più potenti.


\subsubsection{Il Polso a Radiofrequenza e la Precessione}

Dopo aver allineato gli spin lungo $\mathbf{B}_0$, la sequenza fMRI prevede l'applicazione di un secondo solenoide che eroga per un breve intervallo di tempo una corrente alternata a \textbf{radiofrequenza (RF)}. Questo genera un secondo campo magnetico $\mathbf{B}_1$, orientato \textbf{ortogonalmente} a $\mathbf{B}_0$ e variabile nel tempo alla frequenza di Larmor:

$$ \omega_L = \gamma B_0 $$

dove $\gamma$ è il rapporto giromagnetico del protone. La durata del polso RF è calibrata affinché gli spin inizino una \textbf{precessione coerente} attorno alla direzione di $\mathbf{B}_1$: tutti gli spin ruotano in sincronia alla frequenza del polso RF, spostandosi dal piano parallelo a $\mathbf{B}_0$ verso un piano ortogonale (piano trasversale).

\subsubsection{Rilassamento Trasversale}

Quando il polso RF viene rimosso, gli spin continuano a ruotare nel piano trasversale. Questo moto coerente dei protoni genera un campo magnetico oscillante rilevabile dall'esterno: per la legge di Faraday, la variazione del flusso magnetico induce una corrente misurabile nei circuiti ricevitori. Si osserva quindi un segnale oscillante — la cosiddetta \textbf{Free Induction Decay (FID)} — che decade in ampiezza nel tempo.

Il decadimento è causato dalla progressiva \textbf{perdita di coerenza di fase} (defasamento) tra gli spin individuali. Rumore termico, disomogeneità locali del campo magnetico, e — come si vedrà tra breve — la presenza di molecole paramagnetiche, accelerano questo dephasing. Il segnale rilevato è:

$$ S(t) = S_0 \, e^{-t / T_2^*} $$

La costante di tempo \textbf{$T_2^*$} è la misura della velocità di defasamento degli spin in un dato volume cerebrale (voxel). Il valore tipico di $T_2^*$ nella corteccia cerebrale è dell'ordine di \textbf{decine di millisecondi}: un tempo sufficientemente breve da permettere l'acquisizione di molti polsi RF in rapida successione, coprendo in tal modo tutto il volume cerebrale.

\subsubsection{Rilassamento Longitudinale}

Dopo il defasamento trasversale, gli spin rientrano progressivamente nell'allineamento lungo $\mathbf{B}_0$ — il cosiddetto \textbf{rilassamento longitudinale}. Questo processo avviene su scale temporali molto più lunghe, caratterizzate dalla costante \textbf{$T_1$}, dell'ordine di \textbf{frazioni di secondo} (tipicamente $\sim 0.5$--$1$ s per il tessuto cerebrale a 3 T).

Il ciclo sperimentale dell'fMRI è quindi:

\begin{enumerate}
    \item Polso RF $\rightarrow$ precessione coerente degli spin
    \item Misura del decadimento $T_2^*$ in tutti i voxel
    \item Attesa del rilassamento longitudinale $T_1$ ($\sim 0.5$ s)
    \item Nuovo polso RF $\rightarrow$ inizio del ciclo successivo
\end{enumerate}

Questo schema impone una \textbf{risoluzione temporale} dell'fMRI intrinsecamente limitata: non è possibile acquisire immagini a frequenza superiore a circa \textbf{2 Hz} (una immagine ogni 0.5 s). Tuttavia, questo permette comunque di campionare l'attività cerebrale in tutto il volume con un tempo di ripetizione (TR) dell'ordine di mezzo secondo.



\subsection{Il Segnale BOLD — Emoglobina e Attività Neuronale}

\subsubsection{L'Emoglobina come Marcatore dell'Attività}

Il legame tra il segnale $T_2^*$ e l'attività neuronale passa attraverso l'\textbf{emoglobina} — la proteina trasportatrice dell'ossigeno presente nei globuli rossi. L'ossigeno è la fonte di energia sfruttata dal metabolismo neuronale per mantenere l'equilibrio ionico trans-membrana (pompe ioniche) e per supportare la sinaptica. La regola generale è: \textbf{più una rete neurale è attiva, più ossigeno consuma}.

\subsubsection{Deossiemoglobina e Paramagnetismo}

Nelle condizioni di bassa attività neuronale (firing rate basale $\approx 1$ Hz), la rete non richiede grandi quantità di ossigeno. Di conseguenza, nei vasi sanguigni del cervello è presente una frazione significativa di \textbf{deossiemoglobina} (emoglobina che ha rilasciato l'ossigeno al tessuto).

La deossiemoglobina contiene un nucleo di ferro che, in assenza di ossigeno legato, si trova in uno stato ad alto spin (\textbf{paramagnetico}). Il ferro paramagnetico è una sorgente di campo magnetico locale disordinato che \textbf{perturba fortemente} l'omogeneità del campo magnetico locale. Questa perturbazione accelerara il defasamento degli spin dell'acqua:

$$ T_2^* \text{ piccolo} \Longleftrightarrow \text{alta [deossiHb]} \Longleftrightarrow \text{bassa attività neuronale} $$

\subsubsection{Ossiemoglobina e Riduzione del Disturbo Paramagnetico}

Quando invece la rete neuronale è attiva e richiede ossigeno, i vasi sanguigni si dilatano e il flusso ematico aumenta, portando una maggiore quantità di \textbf{ossiemoglobina} (emoglobina con l'ossigeno legato al ferro). In questo stato, il nucleo di ferro è \textbf{diamagnetico}: l'ossigeno ne maschera le proprietà paramagnetiche, riducendo drasticamente la perturbazione del campo magnetico locale.

Il campo locale è quindi più omogeneo, gli spin dell'acqua defasano più lentamente, e $T_2^*$ risulta più grande:

$$ T_2^* \text{ grande} \Longleftrightarrow \text{alta [ossiHb]} \Longleftrightarrow \text{alta attività neuronale} $$

\subsubsection{Il Segnale BOLD}

Il segnale \textbf{BOLD} (Blood Oxygen Level Dependent) è esattamente la mappa di $T_2^*$ acquisita in ogni voxel ad ogni intervallo di acquisizione ($\sim 0.5$ s). La codifica a colori convenzionale è:

\begin{itemize}
    \item \textbf{Blu} $\rightarrow$ $T_2^*$ piccolo $\rightarrow$ alta [deossiHb] $\rightarrow$ bassa attività
    \item \textbf{Rosso} $\rightarrow$ $T_2^*$ grande $\rightarrow$ alta [ossiHb] $\rightarrow$ alta attività
\end{itemize}

Il BOLD costituisce pertanto una \textbf{misura indiretta} dell'attività neuronale, mediata dal metabolismo ossidativo e dalla vascolarizzazione locale. Questa misura è accoppiata all'attività attraverso la risposta emodinamica.



\subsection{Limiti del BOLD}

Il segnale BOLD non è una misura diretta dell'attività elettrica neuronale, ma è mediato dal metabolismo: la risposta emodinamica (vasodilatazione, aumento del flusso ematico, variazione del rapporto ossi/deossiemoglobina) risponde con un \textbf{ritardo} di diversi secondi alla variazione di attività. Inoltre, la relazione tra attività elettrica e risposta metabolica potrebbe essere non lineare, soprattutto per stimolazioni prolungate.

\subsubsection{L'Esperimento di Logothetis: Registrazione Simultanea}

Per valutare quantitativamente la relazione tra BOLD e attività elettrica, il gruppo di \textbf{Nikos Logothetis} (Max Planck Institute for Biological Cybernetics, Tübingen, Germania) ha condotto all'inizio degli anni 2000 un esperimento fondamentale: la registrazione simultanea del segnale BOLD con fMRI a 1.5 T e dell'attività elettrica con un \textbf{elettrodo di tungsteno} impiantato nella \textbf{corteccia visiva} di un macaco.

L'esperimento non era banale dal punto di vista tecnico: un elettrodo metallico inserito in un magnete da 1.5 T è soggetto a forze meccaniche considerevoli e a una serie di artefatti elettromagnetici. Il gruppo di Logothetis riuscì tuttavia a dominare questi problemi e a ottenere registrazioni simultanee attendibili.

Lo stimolo visivo era una \textbf{scacchiera in movimento}, che elicita una forte risposta nella corteccia visiva primaria (V1).

\subsubsection{Multi-Unit Activity (MUA) e Risposta BOLD}

L'attività elettrica registrata dall'elettrodo di tungsteno era la \textbf{Multi-Unit Activity (MUA)}: l'attività di scarica (firing) di numerosi neuroni vicini alla punta dell'elettrodo. La MUA era rappresentata da una curva temporale a colore ciano. Il segnale BOLD del voxel corrispondente alla posizione dell'elettrodo era tracciato in rosso.

I risultati mostrarono:

\begin{enumerate}
    \item \textbf{Ritardo del BOLD rispetto alla MUA}: la risposta BOLD arriva con un ritardo di \textbf{secondi} rispetto all'inizio dell'attività neuronale.
    \item \textbf{Adattamento della MUA}: le unità neurali mostrano un adattamento del firing rate durante la stimolazione sostenuta — il tasso di scarica si riduce progressivamente (le reti cercano di risparmiare energia evitando un eccessivo consumo di spike).
    \item \textbf{Rebound post-stimolazione}: alla cessazione della scacchiera, si osserva un transitorio rimbalzo dell'attività.
\end{enumerate}

\subsubsection{Il Filtro di Kalman e la Trasformazione Lineare}

Per stimolazioni brevi e acute, Logothetis et al. dimostrarono che esiste una \textbf{trasformazione lineare} dalla MUA al segnale BOLD. Costruendo un \textbf{filtro di Kalman} (o filtro lineare ottimale), a partire dalla MUA è possibile predire con alta fedeltà il corso temporale del BOLD. Questa convoluzione lineare è rappresentata dalla curva nera sovrapposta al BOLD nel lavoro originale.

La conclusione importante è:

$$ \text{BOLD}(t) \approx (h * \text{MUA})(t) $$

dove $h(t)$ è la \textbf{funzione di risposta emodinamica} (HRF, Hemodynamic Response Function), un kernel lineare con picco intorno a 5–6 s. Questo significa che il BOLD è un segnale filtrato passa-basso della MUA — il che è una buona notizia, perché conoscendo il kernel si può in linea di principio invertire la trasformazione e stimare la MUA dal BOLD.

\subsubsection{Limite Fondamentale: Divergenza per Stimolazioni Prolungate}

Il problema emerge chiaramente quando la stimolazione è mantenuta per lungo tempo (condizione tipica dell'attività cerebrale in stato naturale, in cui il cervello non è mai silente). In questi casi, si osserva una \textbf{divergenza netta} tra la predizione del filtro di Kalman e il segnale BOLD effettivo:

\begin{itemize}
    \item La MUA può essere nulla o quasi nulla (la rete si è adattata e non produce spike)
    \item Eppure il segnale BOLD rimane \textbf{positivo} — suggerendo falsamente che la rete sia attiva
\end{itemize}

In questo scenario, il BOLD non riflette più l'attività di scarica della rete. Il segnale BOLD è un ottimo indicatore dell'attività neuronale in risposta a stimoli brevi e acuti. Per stati prolungati — che sono la norma nel cervello umano in condizioni naturali — il BOLD può fornire indicazioni fuorvianti, mostrando segnale positivo in assenza di attività di scarica. 



\section{Imaging Ottico con Coloranti Voltaggio-Sensibili}

\subsubsection{Principio Fisico dei Coloranti VSD}

L'attività neuronale può essere registrata con alta risoluzione spaziale anche attraverso l'\textbf{imaging ottico} con \textbf{coloranti voltaggio-sensibili} (Voltage-Sensitive Dyes, VSD). Questi coloranti sono molecole che si legano alla membrana plasmatica dei neuroni — inserendosi specificamente nel \textbf{doppio strato fosfolipidico} (bilipid layer) che separa il compartimento intracellulare da quello extracellulare.

La proprietà fondamentale dei VSD è che le loro caratteristiche di \textbf{fluorescenza} dipendono dal \textbf{potenziale di membrana} trans-bilayer. In particolare, illuminando il tessuto con una lunghezza d'onda appropriata (solitamente nel range dell'arancione o del rosso), i VSD assorbono fotoni e li riemettono per fluorescenza con un'intensità e/o uno spettro che variano in funzione della differenza di potenziale tra i due lati della membrana.

Quando la membrana è a riposo (potenziale di membrana di equilibrio, tipicamente intorno a $-70$~mV), la fluorescenza ha una certa intensità basale. Quando la membrana si depolarizza — a causa di correnti sinaptiche o di uno spike — la variazione del campo elettrico trans-membrana modifica la configurazione elettronica dei VSD, producendo una variazione misurabile della fluorescenza.

\subsubsection{Applicazione Sperimentale}

La procedura sperimentale prevede di applicare i VSD direttamente alla superficie corticale esposta chirurgicamente (ad esempio, posando in polvere o in soluzione il colorante sulla corteccia). Il tessuto viene quindi illuminato e la fluorescenza riemessa viene raccolta da un \textbf{microscopio} dotato di una \textbf{camera ad alta velocità}, che acquisisce una sequenza rapida di immagini della corteccia.

Il risultato è un \textbf{film dell'attività corticale}: ogni pixel dell'immagine corrisponde a una piccola area del tessuto corticale, e il suo valore di intensità (codificato in colori) rappresenta il potenziale di membrana medio dei neuroni in quella posizione. Per convenzione:

\begin{itemize}
    \item \textbf{Blu} $\rightarrow$ potenziale di membrana a riposo $\rightarrow$ bassa attività
    \item \textbf{Rosso} $\rightarrow$ marcata depolarizzazione $\rightarrow$ alta attività
\end{itemize}

Va sottolineato che i VSD registrano principalmente l'attività delle \textbf{strutture dendritiche superficiali}: il segnale riflette il potenziale medio dei \textbf{dendriti apicali} che si proiettano verso la superficie del tessuto, non direttamente il potenziale somatico.


\subsection{Esperimenti di Percezione Visiva con VSD}

\subsubsection{Esperimento 1: Risposta a un Punto Rosso}

In una serie di esperimenti paradigmatici condotti su macaco, i VSD sono stati utilizzati per registrare l'attività della \textbf{corteccia visiva primaria (V1)} in risposta a diversi tipi di stimoli visivi. La corteccia visiva è una struttura ideale per questi esperimenti perché presenta una rappresentazione \textbf{topografica} del campo visivo: ogni posizione retinica corrisponde a un'area specifica della corteccia (mappa retinotopica).

Nel primo esperimento, si presenta all'animale un \textbf{piccolo punto rosso} per 40 millisecondi, poi lo si spegne. Ciò che si osserva nel film VSD è:

\begin{enumerate}
    \item Dopo un breve ritardo, nella posizione corticale corrispondente alla mappa retinotopica del punto appare un'area di attivazione (colori caldi).
    \item Questo spot di attività è il risultato dell'attivazione selettiva dei neuroni selettivi a quella posizione nel campo visivo.
    \item L'attività si diffonde lateralmente nelle aree corticali vicine — perché i neuroni comunicano con i neuroni circostanti attraverso connessioni ricorrenti.
    \item L'intera sequenza si svolge nell'arco di circa 100 ms.
\end{enumerate}

\subsubsection{Esperimento 2: Risposta a una Barra Verticale}

Nel secondo esperimento si presenta all'animale una \textbf{barra verticale}, ovvero un oggetto esteso nel campo visivo. Si osserva che la risposta corticale ha un \textbf{ritardo maggiore} rispetto al caso del punto. Perché?

La risposta sta nella struttura funzionale della corteccia visiva: i neuroni in V1 sono \textbf{selettivi per l'orientamento} (ciascun neurone risponde preferenzialmente a stimoli con una certa orientazione e posizione). Quando si presenta una barra verticale estesa, si attivano selettivamente i neuroni selettivi all'orientazione verticale nella zona corrispondente. Contemporaneamente, però, si attivano anche i circuiti \textbf{inibitori}: i neuroni non selettivi per quella orientazione vengono soppressi dall'inibizione laterale.
Nasce così una competizione tra inibizione ed eccitazione che ritarda la risposta corticale.
La risposta è più lenta perché la \textbf{competizione} tra eccitazione e inibizione deve risolversi prima che emerga un pattern di attività stabile e selettivo. Con uno stimolo piccolo e coerente (come il punto), si coinvolgono solo i neuroni che rispondono a quello stimolo, senza attivare circuiti inibitori estesi — e per questo la risposta è più rapida.

\subsubsection{Esperimento 3: Il Punto in Movimento}

Nel terzo esperimento si presenta all'animale un \textbf{punto in movimento}: un piccolo disco che si sposta verticalmente verso il basso per tutta la durata della finestra di osservazione (circa 125 ms). I VSD mostrano ora una \textbf{onda di attivazione} che scorre lungo la corteccia visiva, seguendo la traiettoria topografica del punto nel campo visivo — prima si attiva la regione corrispondente alla posizione iniziale del punto, poi progressivamente le regioni più in basso.
L'attivazione finale, dopo che il punto ha percorso tutto il suo tragitto, è simile a quella prodotta dalla barra verticale.

\subsubsection{L'Illusione di Movimento e il Cinema}

La sequenza degli esperimenti porta a un risultato notevole per la percezione visiva: se si presenta all'animale prima il punto fermo per 50 ms (esperimento 1), e poi la barra verticale per i successivi 100 ms (esperimento 2), l'attivazione della corteccia visiva è \textbf{indistinguibile} da quella prodotta dal punto in movimento continuo (esperimento 3).

I film VSD delle due condizioni sono, letteralmente, indistinguibili. Questo significa che il cervello — sulla base dell'attività corticale — non distingue tra:
\begin{itemize}
    \item un punto che si muove continuamente verso il basso, e
    \item una sequenza di due snapshot: un punto in alto, poi una barra verticale.
\end{itemize}

Questo è proprio il principio del cinema e dello stop-motion. Il cervello interpola tra snapshot discreti e percepisce un continuum. La ragione evolutiva è che i sistemi visivi biologici sono principalmente ottimizzati per rilevare pattern in movimento. Questo è il motivo per cui percepiamo la sequenza di fotogrammi di un film come un'immagine fluida, anche se viene presentata come una serie di snapshot distinti.



\section{Calcium Imaging}

\subsubsection{Principio e Coloranti Calcio-Sensibili}

Un'alternativa ai VSD per il monitoraggio ottico dell'attività neuronale è il \textbf{calcium imaging} (imaging del calcio intracellulare). I \textbf{coloranti calcio-sensibili} — o, nelle versioni geneticamente codificate, le proteine indicatori del calcio come GCaMP — sono molecole la cui fluorescenza varia in funzione della \textbf{concentrazione di ioni Ca$^{2+}$} nel citoplasma neuronale.

Il razionale biologico è il seguente: ogni volta che un neurone produce uno \textbf{spike} (potenziale d'azione), si verifica un significativo \textbf{influsso di Ca$^{2+}$} attraverso i canali del calcio voltaggio-dipendenti presenti nella membrana somatica. Questo picco transitorio di calcio intracellulare determina un aumento della fluorescenza rilevabile con il microscopio a fluorescenza.

\subsubsection{Risoluzione Spaziale a Livello di Singola Cellula}

La risoluzione spaziale del calcium imaging è superiore a quella dei VSD: mentre i VSD registrano principalmente il potenziale medio dei dendriti apicali superficiali, il calcium imaging permette di identificare il \textbf{corpo cellulare (soma) di singole cellule}, visibili come "nuvole bianche" nel campo microscopico. Con questo approccio è possibile costruire \textbf{tracce di fluorescenza per singole cellule}  utilizzando la maschera spaziale fornita dal microscopio.

Grazie ai microscopi a due fotoni e agli indicatori del calcio di ultima generazione (come GCaMP8), è oggi possibile spingersi ulteriormente nella risoluzione spaziale, raggiungendo le \textbf{sinapsi individuali}. Poiché le sinapsi sono i siti fisici dove avviene la \textbf{plasticità sinaptica} (e quindi l'apprendimento), il calcium imaging offre uno strumento straordinario per osservare in tempo reale come l'informazione viene codificata e consolidata nelle reti corticali.

\subsubsection{Limitazione: Lenta Cinetica di Decadimento}

Il principale limite del calcium imaging è la \textbf{lenta cinetica di decadimento} della fluorescenza. Dopo uno spike, la concentrazione di calcio intracellulare aumenta rapidamente e poi decade con una costante di tempo relativamente lunga (circa \textbf{5 secondi})

\begin{itemize}
    \item Se il neurone produce pochi spike ben separati nel tempo, la tecnica li risolve singolarmente.
    \item Se il neurone produce un burst di spike ad alta frequenza, i segnali si sovrappongono e non si possono distinguere.
\end{itemize}

La capacità di riconoscere singoli eventi di firing è quindi limitata alle basse frequenze di scarica.



\section{Elettrofisiologia Multi-scala}

\subsubsection{Registrazione Simultanea a Più Scale}

La tecnica di riferimento per eccellenza nell'indagine della dinamica neuronale rimane l'\textbf{elettrofisiologia}, che misura direttamente il campo elettrico prodotto dall'attività delle reti neurali. Un esempio emblematico di registrazione multi-scala è fornito dal laboratorio di \textbf{György Buzsáki} (New York University), che ha sviluppato sistemi per la registrazione simultanea di segnali a diversi livelli di profondità nel cervello di pazienti epilettici con elettrodi intracranici:

\begin{enumerate}
    \item \textbf{EEG sullo scalpo} (fronto-parietale e occipitale)
    \item \textbf{ECoG} (elettrodi epidurali sulla superficie corticale)
    \item \textbf{Elettrodi intraparenchimali} (all'interno del tessuto cerebrale, a vari livelli)
\end{enumerate}

L'osservazione più immediata di questa registrazione è che l'\textbf{ampiezza delle fluttuazioni del potenziale elettrico} aumenta drasticamente all'avvicinarsi della sorgente: il segnale EEG sullo scalpo è dell'ordine di decine di microvolt, mentre il segnale registrato dagli elettrodi intra-parenchimali vicini ai neuroni è dell'ordine di millivolt — tre ordini di grandezza maggiore.

\subsubsection{La Slow-Wave Sleep come Esempio Paradigmatico}

Un esempio paradigmatico mostrato nel corso è la registrazione durante il \textbf{sonno a onde lente (Slow Wave Sleep, SWS)} — la prima fase del sonno profondo. Durante il SWS, l'EEG registra grandi onde di ampiezza (le onde lente o \textbf{delta waves}), caratterizzate da una sequenza ritmica di depolarizzazione e iperpolarizzazione.

L'origine di queste onde è stato oggetto di decenni di ricerca. La lettura multi-scala rivela che:

\begin{itemize}
    \item \textbf{All'EEG}: si registra una grande onda (polarizzazione $\rightarrow$ onda positiva, poi negativa)
    \item \textbf{Ai singoli neuroni}: si osserva alternanza di periodi di \textbf{firing intenso} (onda UP) e di \textbf{silenzio completo} (onda DOWN)
\end{itemize}

Berger negli anni '30 registrava questi grandi segnali sul cranio senza conoscere l'attività sottostante a livello cellulare. Le moderne registrazioni multi-scala mostrano che la silenziosa "onda DOWN" corrisponde, paradossalmente, a una depolarizzazione osservata sullo scalpo.

Le onde lente del sonno sono conservate attraverso l'evoluzione — sono presenti nei rettili, negli uccelli e nei mammiferi. Questo suggerisce che il sonno lento serva a funzioni biologiche fondamentali: il \textbf{consolidamento della memoria} (attraverso la replay di sequenze di attività) e la \textbf{pulizia dai prodotti tossici del metabolismo} cerebrale, poiché le onde ioniche prodotte dall'alternanza UP/DOWN creano gradienti che favoriscono il movimento del fluido interstiziale e la rimozione delle scorie metaboliche.

\section{ECoG — Corticografia Epidurale}

\subsubsection{Indicazione Clinica: Epilessia Farmacoresistente}

La \textbf{corticografia epidurale (ECoG)} — acronimo di Electrocorticography — è una tecnica di registrazione invasiva dell'attività corticale che trova la sua principale applicazione clinica nei \textbf{pazienti con epilessia farmacoresistente}. Si tratta di pazienti affetti da epilessia che non risponde ai farmaci antiepilettici disponibili, con crisi molto frequenti che compromettono severamente la qualità della vita.

In questi casi, una delle opzioni terapeutiche è la \textbf{resezione chirurgica} della zona epilettogena: il tessuto corticale da cui originano le crisi viene asportato.  Il tessuto che genera le crisi epilettiche ripetute è già un tessuto patologico. Le crisi ripetute inducono una plasticità aberrante: le connessioni tra neuroni si rafforzano enormemente, rendendo la rete sempre più ipereccitabile. Questo tessuto è già, in senso funzionale, inutilizzabile. Rimuoverlo non riduce la capacità cognitiva del resto del cervello; al contrario, la frequenza delle crisi si riduce drasticamente — in molti casi, esse scompaiono del tutto. Nei pazienti pediatrici, intervenire precocemente è ancora più importante: ogni crisi aggiuntiva aggrava il danno alla rete.


\subsubsection{Procedura Chirurgica e Mappa Pre-operatoria}

Prima della resezione, i neurochirurghi devono costruire una \textbf{mappa dettagliata} di due informazioni fondamentali:

\begin{enumerate}
    \item \textbf{Dove originano le crisi} (zona epilettogena)
    \item \textbf{Dove si trovano le aree funzionali essenziali} (linguaggio, motricità, etc.) che non devono essere toccate
\end{enumerate}

A questo scopo si impiantano temporaneamente al paziente:
\begin{itemize}
    \item Una \textbf{griglia di elettrodi epidurali} (sulla superficie corticale)
    \item \textbf{Elettrodi profondi} (sEEG, stereo-EEG) nel parenchima cerebrale
\end{itemize}

L'impianto è eseguito in modo estremamente preciso, con mappatura preoperatoria dei vasi cerebrali, per evitare qualsiasi danno da compressione o emorragia.

\subsubsection{Registrazione e Microstimolazione}

La griglia ECoG permette di:

\begin{itemize}
    \item \textbf{Registrare} l'attività corticale durante le crisi spontanee per identificare il punto di insorgenza
    \item \textbf{Micro-stimolare} selettivamente diverse aree per mappare le funzioni: stimolando l'area del linguaggio si provoca una transitoria incapacità di parlare, confermando la localizzazione; stimolando l'area motoria si osserva un movimento involontario dell'arto controlaterale.
\end{itemize}

Questo schema di registrazione + stimolazione consente al neurochirurgo di definire un'\textbf{area minima di resezione} che massimizza l'eliminazione delle crisi e minimizza il deficit funzionale post-operatorio: il \textbf{rapporto rischio/beneficio} risulta così ottimizzato.


\section{Patch Clamp}

Il \textbf{patch clamp} è la tecnica che offre la \textbf{massima risoluzione} possibile per lo studio dell'attività elettrica di un singolo neurone: registra direttamente il \textbf{potenziale di membrana del soma} in condizioni in vivo o in vitro.

La procedura utilizza una \textbf{pipetta di vetro} con un'apertura submicrometrica (solitamente dell'ordine di 1–2 $\mu$m), riempita con una soluzione elettrolitica. Le fasi dell'approccio sono:

\begin{enumerate}
    \item La pipetta viene avvicinata alla membrana del corpo cellulare (soma) del neurone bersaglio con un micromanipola micrometrica.
    \item Applicando una leggera aspirazione, si crea un sigillo ad alta resistenza (gigaohm seal, o "giga-seal") tra la punta della pipetta e la membrana plasmatica.
    \item Continuando l'aspirazione, si \textbf{rompe} il patch di membrana sottostante la pipetta.
    \item Si stabilisce così un contatto diretto tra il lume della pipetta e il \textbf{fluido citoplasmatico} intracellulare.
\end{enumerate}

In questo modo, la pipetta e il soma del neurone sono in contatto diretto e hanno lo \textbf{stesso potenziale elettrico} (configurazione di whole-cell patch clamp). Un amplificatore ad alta impedenza collegato alla pipetta registra in tempo reale le variazioni di questo potenziale: si ha così la \textbf{registrazione diretta del potenziale di membrana} con una risoluzione temporale di frazioni di millisecondo.

ì

Tramite la stessa pipetta, è possibile anche \textbf{iniettare ioni o corrente} nel citoplasma del neurone, modificando così la sua eccitabilità. Questa capacità permette di:

\begin{itemize}
    \item Iniettare una corrente depolarizzante per forzare il neurone a produrre spike (studio del pattern di scarica)
    \item Iniettare corrente iperpolarizzante per sopprimere l'attività
    \item Introdurre composti chimici specifici per studiare il ruolo di recettori o canali
\end{itemize}


La procedura di patch clamp, pur essendo di massima risoluzione, ha un costo biologico: la rottura della membrana \textbf{danneggia la cellula} in modo progressivo e irreversibile. La registrazione stabile è tipicamente mantenuta per \textbf{poche ore}, talvolta meno in condizioni in vivo. Non è quindi possibile seguire lo stesso neurone per lunghi periodi temporali.


\section{Optogenetica}

\subsubsection{Channelrhodopsine}

L'\textbf{optogenetica} è una tecnica che permette di \textbf{perturbare e registrare} l'attività neuronale utilizzando la luce, con una selettività cellulare impossibile da ottenere con gli elettrodi. Il principio fisico sfrutta una classe di proteine di membrana — le \textbf{channelrhodopsine} (ChR) — che sono canali ionici la cui apertura è controllata dall'assorbimento di fotoni.

Le channelrhodopsine sono naturalmente presenti nella retina delle alghe verdi e, in forma omologa, nella retina di tutti i mammiferi: nella retina dei vertebrati, la rodopsina nei fotorecettori (bastoncelli e coni) è il corrispondente evolutivo che trasduce la luce in segnale elettrico. L'idea portante dell'optogenetica è di trapiantare questi canali fotosensibili in neuroni corticali, in modo da poterli controllare con la luce.

\subsubsection{Consegna Virale e Ingegneria Genetica}

Poiché le channelrhodopsine non sono presenti nei neuroni corticali, è necessario introdurre il loro gene mediante \textbf{vettori virali} (di solito adeno-associated virus, AAV) \textbf{geneticamente modificati}. Una volta iniettato il virus nella zona corticale di interesse, il virus infetta i neuroni bersaglio e inserisce nel loro DNA (o ne esprime transitoriamente) il gene della channelrhodopsina. I neuroni così trattati iniziano a produrre la proteina e a inserirla nella membrana plasmatica.

Questa procedura è applicabile esclusivamente in \textbf{modelli animali} (topi, macachi, etc.): l'inserimento di materiale genetico estraneo in un organismo umano sano per finalità di ricerca non è eticamente consentito.

\subsubsection{Perturbazione e Registrazione Ottica}

Una volta espressa la channelrhodopsina, i neuroni transgenici possono essere attivati o inibiti mediante impulsi luminosi di lunghezze d'onda appropriate:

\begin{itemize}
    \item \textbf{Luce blu} ($\sim 470$ nm) $\rightarrow$ apertura di ChR2 (channelrhodopsina-2) $\rightarrow$ afflusso di cationi (Na$^+$, Ca$^{2+}$) $\rightarrow$ depolarizzazione $\rightarrow$ potenziale d'azione
    \item \textbf{Luce gialla/arancione} ($\sim 590$ nm) $\rightarrow$ apertura di NpHR (halorhodopsina) $\rightarrow$ afflusso di Cl$^-$ $\rightarrow$ iperpolarizzazione $\rightarrow$ soppressione dell'attività
\end{itemize}

Combinando optogenetica con l'imaging ottico (VSD o calcium imaging), è possibile sia \textbf{stimolare} (con la luce) sia \textbf{registrare} (con la fluorescenza) l'attività neuronale nello stesso preparato — un approccio totalmente all-optical che offre selettività cellulare e risoluzione spaziale straordinarie.

\newpage

\section{Teoria del Campo Elettrico Extracellulare}

\subsubsection{Modello dell'Acqua Salata}

Per comprendere il meccanismo alla base della generazione del \textbf{campo elettrico extracellulare} — il segnale rilevato dagli elettrodi — si parte da un esperimento modello. Si considera un contenitore di \textbf{acqua salata} (soluzione di NaCl) in cui sono immersi un anodo e un catodo collegati a un generatore di corrente. L'acqua salata è un modello eccellente del mezzo extracellulare cerebrale, che è anch'esso una soluzione acquosa di ioni (Na$^+$, K$^+$, Cl$^-$, Ca$^{2+}$).

In questa configurazione si osservano:
\begin{itemize}
    \item \textbf{Linee isoequipotenziali}: curve chiuse attorno all'anodo (potenziale positivo) e al catodo (potenziale negativo)
    \item \textbf{Linee di flusso di corrente}: ortogonali alle isoequipotenziali, seguono il gradiente del potenziale
\end{itemize}

Un elettrodo di misura immerso in diverse posizioni nel contenitore registrerà potenziali diversi, in funzione della sua distanza relativa dall'anodo e dal catodo. Questo è esattamente il principio che governa la misurazione del campo elettrico cerebrale.

\subsubsection{L'Approssimazione Quasi-Statica di Maxwell}

Il punto di partenza fisico è la \textbf{legge di Faraday-Maxwell}:

$$ \nabla \times \mathbf{E} = -\frac{\partial \mathbf{B}}{\partial t} $$

In biologia, è lecito adottare l'\textbf{approssimazione quasi-statica}: i neuroni sparano tipicamente a frequenze dell'ordine di 1 Hz (a riposo), con potenziali d'azione a $\sim 1$ kHz al massimo. Queste frequenze sono immensamente inferiori a quelle associate alla propagazione elettromagnetica (frequenze dell'ordine di $c/L$, dove $c$ è la velocità della luce e $L$ la dimensione tipica del sistema). Le variazioni temporali del campo magnetico indotto da queste correnti lente sono quindi trascurabili:

$$ \frac{\partial \mathbf{B}}{\partial t} \approx 0 $$

Ne consegue che il campo elettrico extracellulare è \textbf{conservativo} (irrotazionale):

$$ \nabla \times \mathbf{E} = 0 \implies \mathbf{E} = -\nabla \varphi $$

dove $\varphi$ è il \textbf{potenziale elettrico extracellulare} — la grandezza misurata dagli elettrodi.

\subsubsection{Conservazione delle Cariche e Legge di Ohm con Sorgenti Neurali}

La seconda equazione fondamentale è la \textbf{conservazione delle cariche}. A scala mesoscopica (dell'ordine del millimetro), il mezzo extracellulare è elettricamente neutro: la densità di carica netta $\rho$ è nulla. Poiché $\rho$ non varia nel tempo a questa scala:

$$ \frac{\partial \rho}{\partial t} = 0 \implies \nabla \cdot \mathbf{J} = 0 $$

dove \textbf{J} è la densità di corrente totale. Tuttavia, la densità di corrente non è solo quella di deriva nel mezzo (corrente ohmica), ma include anche le \textbf{sorgenti di corrente neurali} $\mathbf{J}_s$, associate all'apertura dei canali ionici durante l'attività sinaptica o la generazione di spike:

$$ \mathbf{J} = \sigma \mathbf{E} + \mathbf{J}_s = -\sigma \nabla \varphi + \mathbf{J}_s $$

dove $\sigma$ è la \textbf{conducibilità elettrica} del mezzo extracellulare, assunta omogenea nello spazio con buona approssimazione. Sostituendo nella condizione di continuità:

$$ \nabla \cdot \left( -\sigma \nabla \varphi + \mathbf{J}_s \right) = 0 $$

che porta all'\textbf{equazione di Poisson generalizzata}:

$$ \sigma \nabla^2 \varphi = \nabla \cdot \mathbf{J}_s $$

Questa è l'equazione fondamentale del campo potenziale extracellulare: il \textbf{Laplaciano del potenziale} è determinato dalla \textbf{divergenza delle sorgenti di corrente neurali}.

\subsubsection{Contributo di una Sorgente Puntiforme}

Si considera ora una singola sorgente di corrente puntiforme di intensità $I_n$, localizzata in posizione $\mathbf{r}_n$. Per \textbf{simmetria sferica}, le correnti fluiscono radialmente verso l'esterno (o verso l'interno, se è un sink). La \textbf{densità di corrente} alla distanza $\delta = |\mathbf{r}_e - \mathbf{r}_n|$ è:

$$ J(\delta) = \frac{I_n}{4 \pi \delta^2} $$

dove $4\pi\delta^2$ è l'area della sfera di raggio $\delta$ centrata sulla sorgente. Poiché $\mathbf{E} = -\nabla\varphi$ e $J = \sigma E$, si ha:

$$ \frac{\partial \varphi}{\partial \delta} = -\frac{J}{\sigma} = -\frac{I_n}{4\pi\sigma\delta^2} $$

Integrando dal punto di misura (posizione dell'elettrodo a distanza $\delta$ dalla sorgente) fino all'infinito (dove $\varphi = 0$):

$$ \varphi(\delta) = \int_{\delta}^{\infty} \frac{I_n}{4\pi\sigma\delta'^2} \, d\delta' = \frac{I_n}{4\pi\sigma\delta} $$

Per \textbf{linearità del campo elettrico}, il contributo di $N$ sorgenti si somma:

$$ \varphi(\mathbf{r}_e) = \sum_{n=1}^{N} \frac{I_n}{4\pi\sigma |\mathbf{r}_e - \mathbf{r}_n|} $$

Questa è la \textbf{legge di Coulomb} per il potenziale extracellulare: ogni sorgente contribuisce con un termine inversamente proporzionale alla distanza $|\mathbf{r}_e - \mathbf{r}_n|$. Le sorgenti vicine all'elettrodo dominano il segnale; quelle lontane contribuiscono con un fattore trascurabile.


\subsection{Il Dipolo Neuronale: Sinapsi e Soma}

Le sorgenti di corrente neurali $\mathbf{J}_s$ corrispondono a specifici eventi elettrici nella membrana neuronale:

\begin{itemize}
    \item \textbf{Correnti sinaptiche}: l'apertura di recettori ionici postsinaptici (AMPA, NMDA per l'eccitazione; GABA$_A$ per l'inibizione) introduce un flusso di ioni attraverso la membrana in una posizione specifica (la sinapsi, tipicamente sui dendriti)
    \item \textbf{Potenziali d'azione}: l'apertura sequenziale di canali Na$^+$ e K$^+$ durante lo spike produce correnti trans-membrana molto più intense ma localizzate al soma e al segmento iniziale dell'assone
\end{itemize}

In una simulazione computazionale del gruppo Hines-Carnevale, su un neurone ricostruito in silico, si inietta una corrente iperpolarizzante nella sinapsi apicale e si calcola il campo di potenziale risultante. I risultati mostrano:

\begin{itemize}
    \item \textbf{Zona sinaptica (apicale)}: il flusso di corrente entra nel dendrite $\rightarrow$ questa zona è un \textbf{sink} (sorgente negativa) $\rightarrow$ equipotenziali negative
    \item \textbf{Soma}: il soma è soggetto alla corrente di leakage della membrana, che tende a portare il potenziale di membrana verso l'equilibrio $\rightarrow$ flusso di corrente uscente $\rightarrow$ il soma funziona come un \textbf{catodo} (sorgente positiva) $\rightarrow$ equipotenziali positive
\end{itemize}

Si forma quindi un \textbf{dipolo biologico} con il polo positivo somatico e il polo negativo dendritico. Le linee di campo del potenziale seguono la geometria di un dipolo classico, con le isoequipotenziali positive attorno al soma e negative attorno alla sinapsi apicale.

\subsection{Celle Piramidali e Generazione del Segnale EEG}

La ragione per cui l'EEG è registrabile dallo scalpo è legata alla geometria delle \textbf{cellule piramidali}: queste cellule di grandi dimensioni, che costituiscono il principale tipo di neurone eccitatorio della corteccia cerebrale, sono orientate in modo \textbf{perpendicolare alla superficie corticale}, con i dendriti apicali puntati verso la pia mater e il soma più in profondità.

Quando molte cellule piramidali ricevono contemporaneamente correnti sinaptiche eccitatoriche ai loro dendriti apicali, ciascuna produce un dipolo orientato nella stessa direzione. I dipoli individuali si \textbf{sommano coerentemente}:

$$ \varphi_{\text{EEG}} = \sum_{k=1}^{N_{\text{piramidali}}} \frac{I_k}{4\pi\sigma |\mathbf{r}_{\text{scalpo}} - \mathbf{r}_k|} $$

Il risultato è un \textbf{macrodipolo} sufficientemente intenso da generare variazioni di potenziale rilevabili sullo scalpo.

Al contrario, le \textbf{basket cells} (interneuroni inibitori) hanno geometrie più compatte e orientamenti variabili a seconda della sinapsi attivata. I loro dipoli si orientano in direzioni diverse e si cancellano statisticamente, contribuendo in modo trascurabile al segnale EEG. Ciò implica che \textbf{l'EEG registra principalmente l'attività delle cellule eccitatorie piramidali}, non quella degli interneuroni inibitori.



\section{Local Field Potential (LFP)}

\subsubsection{Il Local Field Potential come Filtro Passa-Basso}

Il segnale registrato da un elettrodo immerso nel tessuto corticale — denominato \textbf{Local Field Potential (LFP)} o potenziale di campo locale — è la sovrapposizione delle correnti di tutte le sorgenti neurali nel volume di tessuto intorno alla punta dell'elettrodo.

Una proprietà fondamentale dell'LFP è che esso è una \textbf{versione filtrata passa-basso} delle correnti sinaptiche che alimentano la rete. Questo si capisce con il seguente argomento: la corrente sinaptica iniettata in un dendrite deve propagarsi attraverso la membrana fino al punto di misura dell'elettrodo. La membrana neuronale è un sistema con proprietà capacitive e resistive, e agisce come un \textbf{filtro passa-basso} per i segnali elettrici che si propagano lungo i suoi compartimenti. Le componenti ad alta frequenza della corrente sinaptica vengono attenuate più rapidamente nel propagarsi dal punto di iniezione al punto di misura:



\subsubsection{Spikes e Localizzazione Spaziale}

Oltre alle correnti sinaptiche, l'LFP include anche i \textbf{potenziali d'azione (spikes)} prodotti dai soma vicini all'elettrodo. Uno spike è un evento enormemente più intenso di una corrente sinaptica (ampiezza del potenziale d'azione $\sim 100$ mV, rispetto a potenziali postsinaptici di pochi mV), ma la sua influenza sull'LFP decade molto rapidamente con la distanza: a \textbf{poche centinaia di micron} dalla sorgente, la deflection associata allo spike si confonde con il rumore di fondo.

Questo decadimento rapido è dovuto a due fattori concorrenti: la legge di Coulomb ($1/\delta$) e l'attenuazione da parte del mezzo stesso. In pratica, un elettrodo immerso nel tessuto corticale "vede" gli spike solo dei neuroni entro un cilindro di raggio di circa \textbf{50 $\mu$m} attorno alla punta.

\subsubsection{Registrazione Multi-informazione con un Singolo Elettrodo}

Un aspetto notevole della registrazione elettrofisiologica è che un \textbf{singolo elettrodo di tungsteno} fornisce simultaneamente più tipi di informazione:

\begin{enumerate}
    \item \textbf{Componente a bassa frequenza dell'LFP} ($<300$ Hz): riflette le correnti sinaptiche della rete nel volume intorno all'elettrodo $\rightarrow$ attività dei circuiti locali
    \item \textbf{Componente ad alta frequenza (dopo high-pass filter)} ($>300$ Hz): contiene i transienti veloci degli spikes $\rightarrow$ attività di scarica dei neuroni vicini
\end{enumerate}

Il procedimento pratico è il seguente: si registra il segnale grezzo (raw LFP), che contiene entrambe le componenti. Applicando un \textbf{filtro passa-basso} si isola la componente sinaptica. Applicando un \textbf{filtro passa-alto} si eliminano le basse frequenze e si ottiene un segnale in cui sono visibili le deflections veloci degli spike.

\subsubsection{Single Unit Activity (SUA)}

Quando la deflection di uno spike nell'LFP filtrato passa-alto è sufficientemente ampia e ha una forma caratteristica riconoscibile, essa appartiene a un neurone specifico e ben localizzato — entro il raggio di $\sim 50$ $\mu$m dall'elettrodo. Questi eventi riconoscibili individualmente costituiscono la \textbf{Single Unit Activity (SUA)}.

La forma della SUA è \textbf{stereotipata per ciascun neurone} (firma neuronale): ha un profilo temporale peculiare che dipende dalla geometria del neurone, dalla distribuzione dei suoi canali ionici e dalla distanza relativa dal punto di misura. Riconoscendo questa forma attraverso algoritmi di \textbf{spike sorting}, è possibile attribuire ogni spike a un neurone specifico, ottenendo la serie temporale di firing di quel neurone.

\subsubsection{Multi-Unit Activity (MUA)}

Oltre alle SUA, nell'LFP filtrato passa-alto sono visibili molte altre deflections più piccole, non attribuibili singolarmente a neuroni specifici ma che rappresentano l'attività di tutti i neuroni nel volume di influenza dell'elettrodo (raggio $> 50$ $\mu$m). Queste deflections collettive costituiscono la \textbf{Multi-Unit Activity (MUA)}: l'attività di scarica non classificata dell'intera popolazione neuronale intorno alla punta.


\subsubsection{Discriminazione Eccitatori/Inibitori dalla Forma dello Spike}

La forma temporale della SUA non è solo una firma individuale: contiene informazione sul \textbf{tipo cellulare} del neurone che ha generato lo spike. Questo si basa sul seguente principio fisico:

\begin{itemize}
    \item \textbf{Neuroni eccitatori (cellule piramidali)}: sono cellule di grandi dimensioni con un albero dendritico esteso. L'apertura dei canali Na$^+$ e K$^+$ durante lo spike genera correnti trans-membrana distribuite su molti compartimenti della cellula, e l'integrazione di questi contributi produce una deflection nell'LFP che è relativamente \textbf{ampia nel tempo} (durata $\sim 0.5$--$1$ ms; spike "broad").
    \item \textbf{Neuroni inibitori (interneuroni fast-spiking, basket cells)}: sono cellule più piccole con dendriti meno estesi. Il loro spike è più compatto nel tempo: la deflection nell'LFP è \textbf{più stretta} (durata $\sim 0.2$--$0.3$ ms; spike "narrow").
\end{itemize}

Questa distinzione permette, almeno in parte, di classificare le unità registrate come \textbf{eccitatorie} (spike larghi) o \textbf{inibitorie} (spike stretti) direttamente dalla forma temporale della SUA, senza necessità di identificare morfologicamente la cellula.



\chapter{Lezione 3}

\textbf{Argomenti:} Membrana bilipidica, canali ionici, specie ioniche, equazione di Nernst-Planck, equazione di Nernst, potenziale di equilibrio/reversione, pompe ioniche, modello Goldman-Hodgkin-Katz, potenziale di riposo

\vspace{0.5cm}
\noindent\rule{\textwidth}{0.4pt}

\section{Il Neurone come Dispositivo Elettrochimico}


Il punto di partenza di questa lezione è apparentemente paradossale: il mezzo cerebrale è, a livello macroscopico, privo di cariche nette. Eppure nei neuroni si misura con grande precisione una differenza di potenziale tra il compartimento intracellulare e quello extracellulare. Come è possibile conciliare queste due osservazioni? La risposta risiede nella struttura microscopica della \textbf{membrana neuronale} e nella distribuzione spazialmente non casuale degli ioni ai suoi bordi. Questa lezione si propone di costruire, passo per passo, il quadro teorico che spiega l'origine di tale differenza di potenziale.

\subsubsection{Anatomia Funzionale del Neurone}

Prima di concentrarsi sulla membrana, è utile richiamare l'architettura del neurone nei suoi compartimenti principali, così come introdotta nella lezione precedente.

Il neurone è organizzato attorno a tre grandi strutture funzionali:

\begin{itemize}
    \item \textbf{Albero dendritico:} costituisce il compartimento di ingresso (\textit{input}) del neurone. Riceve i segnali provenienti da altri neuroni attraverso i contatti sinaptici localizzati sulle spine dendritiche. Le correnti ioniche generate a livello sinaptico si propagano lungo i dendriti verso il soma.
    \item \textbf{Soma:} è il corpo cellulare, sede dell'integrazione delle correnti provenienti dall'albero dendritico. Le proprietà elettriche passive e attive della membrana del soma determinano se e quando il neurone genera un segnale di uscita. Come vedremo nella prossima lezione, il soma è dotato di \textbf{componenti attive della membrana} capaci di amplificare le correnti e produrre il \textbf{potenziale d'azione}.
    \item \textbf{Assone:} è il compartimento di uscita (\textit{output}) del neurone. Conduce i potenziali d'azione dal soma verso i \textbf{bottoni sinaptici presinaptici} (\textit{synaptic boutons}), dove viene trasmesso il segnale al neurone successivo. L'assone è tipicamente avvolto dalla \textbf{mielina}, una guaina isolante prodotta dalle cellule di Schwann (nel sistema nervoso periferico) o dagli oligodendrociti (nel sistema nervoso centrale). La mielina consente una propagazione più efficiente del segnale elettrico, riducendo la dispersione e il consumo energetico.
\end{itemize}

Il flusso d'informazione è quindi: dendriti $\rightarrow$ soma $\rightarrow$ assone $\rightarrow$ bottoni sinaptici.

\subsubsection{La Domanda Centrale della Lezione}

Il focus di questa lezione è sul soma e, in particolare, su una piccola sezione della sua membrana. La domanda che ci poniamo è: \textbf{come si genera una differenza di potenziale tra l'interno e l'esterno della cellula, senza che il mezzo sia macroscopicamente carico?}

Per rispondere, isoliamo concettualmente un piccolo \textit{patch} della membrana del soma e analizziamo cosa accade alla distribuzione degli ioni nelle immediate vicinanze.

\section{La Membrana Bilipidica come Capacitore}

\subsubsection{Struttura del Doppio Strato Lipidico}

La membrana cellulare è un \textbf{doppio strato lipidico} (\textit{bilipidic layer}): due fogli di molecole fosfolipidiche disposti con le code idrofobiche rivolte verso l'interno e le teste idrofiliche verso i mezzi acquosi intracellulare ed extracellulare. Questa struttura la rende un ottimo \textbf{isolante elettrico}: in linea di principio, nessun ion può attraversarla spontaneamente.

Le dimensioni caratteristiche sono fondamentali per comprendere le approssimazioni che useremo nel seguito:

\begin{itemize}
    \item \textbf{Spessore della membrana:} $A \approx 5\,\text{nm}$
    \item \textbf{Raggio tipico di un assone:} $r \approx 5\,\mu\text{m}$
\end{itemize}

Il rapporto $A/r \approx 10^{-3}$, cioè lo spessore della membrana è circa \textbf{1000 volte più piccolo} del raggio dell'assone. Questa disparità di scale ha una conseguenza geometrica fondamentale: su qualsiasi \textit{patch} di membrana di dimensioni non trascurabili, la superficie può essere considerata \textbf{piana}, esattamente come le placche di un condensatore piano.

\subsubsection{La Membrana come Condensatore}

Il fatto che la membrana sia un isolante con due conduttori (il liquido intracellulare e quello extracellulare) ai suoi lati la rende funzionalmente equivalente a un \textbf{condensatore}. La capacità per unità di superficie della membrana è un parametro ben misurato sperimentalmente:

\begin{equation}
C_m \approx 1\,\mu\text{F/cm}^2
\end{equation}

Questa grandezza entra direttamente nel calcolo del numero di cariche che devono spostarsi per modificare il potenziale di membrana.

\subsubsection{Neutralità Macroscopica e Distribuzione Microscopica}

Un punto concettualmente sottile, ma cruciale, è la seguente osservazione: la neutralità macroscopica del mezzo è compatibile con la presenza di un campo elettrico locale nelle immediate vicinanze della membrana.

Consideriamo tre volumi distinti:

\begin{enumerate}
    \item \textbf{Volume intracellulare (lontano dalla membrana):} il conteggio degli ioni positivi e negativi dà zero carica netta.
    \item \textbf{Volume extracellulare (lontano dalla membrana):} ugualmente, la carica netta è zero.
    \item \textbf{Volume "mesoscopico" attorno alla membrana:} anche questo volume contiene lo stesso numero di cariche positive e negative. Tuttavia, la loro distribuzione \textbf{non è casuale}: la maggior parte delle cariche positive si accumula sulla superficie esterna della membrana, e un numero non trascurabile di cariche negative si addensa sulla superficie interna.
\end{enumerate}

Questo è il punto chiave: la neutralità locale è preservata a scala macroscopica, ma la distribuzione spazialmente non uniforme delle cariche nelle immediate vicinanze della membrana genera un \textbf{campo elettrico locale} orientato dall'esterno verso l'interno. 

\section{I Canali Ionici}

\subsubsection{L'Eccezione all'Impermeabilità della Membrana}

Se la membrana fosse un isolante perfetto, non ci sarebbe alcun flusso ionico e il potenziale di membrana sarebbe determinato esclusivamente dalla distribuzione iniziale degli ioni. In realtà, la membrana è percorsa da un'alta densità di \textbf{canali ionici}: macromolecole proteiche transmembrana che attraversano l'intero spessore del doppio strato lipidico e formano un poro acquoso attraverso il quale specifici ioni possono diffondere.

\subsubsection{Specificità dei Canali Ionici}

I canali ionici sono \textbf{selettivi}: ogni tipo di canale è permeabile a una sola specie ionica (o a un gruppo ristretto di specie chimicamente simili). Esistono pertanto:

\begin{itemize}
    \item Canali per il \textbf{potassio} (K$^+$)
    \item Canali per il \textbf{sodio} (Na$^+$)
    \item Canali per il \textbf{cloruro} (Cl$^-$)
    \item Canali per il \textbf{calcio} (Ca$^{2+}$)
\end{itemize}

La \textbf{densità di canali} per unità di superficie della membrana determina la \textbf{mobilità} degli ioni di quella specie attraverso la membrana, e quindi la \textbf{resistenza} che la membrana oppone al loro flusso.

\subsubsection{Scoperta Recente: i Canali Ionici si Muovono}

\begin{tcolorbox}[colback=gray!5, colframe=gray!40, arc=2mm, boxrule=0.5pt]
\textbf{Nota del docente:} "Circa sette anni fa ho scoperto — con mia stessa sorpresa, dato che si trattava di una verità consolidata della biologia strutturale — che i canali ionici non sono fissi nella membrana. Questi macromolecole si trovano in continuo movimento laterale all'interno del doppio strato lipidico. Questo complica notevolmente il quadro teorico classico, ma per gli scopi di questa lezione non vogliamo rendere la storia troppo difficile."
\end{tcolorbox}

Per oggi, facciamo quindi l'assunzione semplificativa che la \textbf{mobilità degli ioni attraverso la membrana sia costante}, ovvero che non dipenda né dalla concentrazione degli ioni né dal potenziale di membrana. Questo significa che trattiamo esclusivamente le \textbf{proprietà elettriche passive} della membrana. La dipendenza della mobilità dal potenziale di membrana — fondamentale per la generazione del potenziale d'azione — sarà discussa nella prossima lezione.


\section{Le Specie Ioniche Principali}

Nel contesto delle proprietà elettriche del soma neuronale, le specie ioniche rilevanti sono essenzialmente quattro. Descriviamo ciascuna con le sue caratteristiche principali.

\textbf{Potassio (K$^+$):}\\
La valenza è $Z_K = +1$. Il potassio è la specie ionica \textbf{più abbondante all'interno} della cellula. Nella corteccia cerebrale dei mammiferi, la concentrazione intracellulare è circa $C_K^{in} = 140\,\text{mM}$, mentre quella extracellulare è circa $C_K^{out} = 5\,\text{mM}$. Il gradiente di concentrazione spinge gli ioni K$^+$ verso l'esterno.

\textbf{Sodio (Na$^+$):}\\
La valenza è $Z_{Na} = +1$. Il sodio ha una situazione opposta al potassio: è molto più abbondante \textbf{fuori} dalla cellula ($C_{Na}^{out} \approx 142\,\text{mM}$) rispetto all'interno ($C_{Na}^{in} \approx 10\,\text{mM}$). Il gradiente di concentrazione spinge gli ioni Na$^+$ verso l'interno.

\textbf{Cloruro (Cl$^-$):}\\
La valenza è $Z_{Cl} = -1$. Il cloruro è uno ione negativo con una distribuzione analoga al sodio: più abbondante all'esterno che all'interno. Il suo ruolo nelle proprietà passive del soma è considerato minore rispetto a K$^+$ e Na$^+$, ma contribuisce in modo non trascurabile al potenziale di riposo.

\textbf{Calcio (Ca$^{2+}$):}\\
La valenza è $Z_{Ca} = +2$. Il calcio ha una concentrazione intracellulare estremamente bassa, dell'ordine del nanomolare ($C_{Ca}^{in} \approx 10^{-4}\,\text{mM}$), mentre la concentrazione extracellulare è di circa $C_{Ca}^{out} \approx 2.5\,\text{mM}$. Il calcio non verrà trattato in questa lezione, ma riveste un ruolo fondamentale in due contesti che affronteremo in seguito:

\begin{enumerate}
    \item La \textbf{trasmissione sinaptica}: il Ca$^{2+}$ è il segnale che innesca il rilascio dei neurotrasmettitori nei bottoni sinaptici.
    \item La \textbf{regolazione della frequenza di scarica} (\textit{firing rate}) del neurone: un afflusso di Ca$^{2+}$ nel soma può modificare la permeabilità della membrana ad altri ioni.
\end{enumerate}

\begin{tcolorbox}[colback=gray!5, colframe=gray!40, arc=2mm, boxrule=0.5pt]
Senza complicare eccessivamente il quadro, ci limitiamo oggi a K$^+$, Na$^+$ e Cl$^-$ come protagonisti principali. Il Ca$^{2+}$ sarà discusso nelle prossime lezioni.
\end{tcolorbox}



\section{Flusso Molare e Origine del Campo Elettrico}

\subsubsection{Il Flusso Molare}

Per quantificare il movimento degli ioni attraverso la membrana, introduciamo il \textbf{flusso molare} $\mathbf{J}^{(s)}$ della specie $s$: un vettore che rappresenta il numero di moli di ioni che attraversano l'unità di superficie nell'unità di tempo, espresso in $[\text{mol}/\text{m}^2\text{s}]$.

I flussi molari dipendono in generale dalla posizione nello spazio (sono \textbf{campi vettoriali}) e dal tempo, poiché il sistema biologico è tipicamente fuori dall'equilibrio.

\subsubsection{Flusso Diffusivo e Squilibrio di Cariche}

Consideriamo il caso del potassio: la sua concentrazione è molto maggiore all'interno ($140\,\text{mM}$) che all'esterno ($5\,\text{mM}$). Questa differenza di concentrazione costituisce un \textbf{gradiente di concentrazione} che, termodinamicamente, tende a equalizzarsi per aumentare l'entropia del sistema. Il risultato è un \textbf{flusso diffusivo} $J_D$ di ioni K$^+$ dall'interno verso l'esterno.

Questo flusso di ioni K$^+$ verso l'esterno crea uno \textbf{squilibrio di cariche}: si accumulano cariche positive sulla superficie esterna della membrana, mentre si crea un deficit di cariche positive (equivalente a un eccesso di cariche negative) sulla superficie interna. La membrana si comporta così come un \textbf{condensatore carico}, con:

\begin{itemize}
    \item Superficie esterna: $\rho_+ > 0$ (eccesso di cariche positive)
    \item Superficie interna: $\rho_- < 0$ (eccesso di cariche negative)
\end{itemize}

Questa distribuzione genera un \textbf{campo elettrico} $\mathbf{E}$ orientato dall'esterno verso l'interno, che tende a rallentare il flusso di K$^+$ verso l'esterno e ad attrarre le cariche positive verso l'interno. Nasce così un \textbf{flusso di deriva} $J_E$ in direzione opposta al flusso diffusivo $J_D$.

\subsubsection{Verso la Competizione tra i Due Flussi}

Il quadro fisico che emerge è quindi quello di una competizione tra due contributi al flusso molare totale:

\begin{equation}
J_{tot}^{(s)} = J_D^{(s)} + J_E^{(s)}
\end{equation}

Il sistema evolve verso uno stato di equilibrio in cui questi due contributi si compensano esattamente, ovvero in cui il flusso netto di ioni è zero. Questo stato di equilibrio è descritto dall'\textbf{equazione di Nernst}.


\section{Il Potenziale di Membrana}

\subsubsection{Asse di Riferimento}

Prima di procedere con la derivazione formale, fissiamo un sistema di riferimento che useremo nel seguito. Consideriamo un \textit{patch} di membrana e orientiamo l'asse $x$ \textbf{perpendicolare} alla membrana:

\begin{itemize}
    \item $x = 0$: parete interna della membrana (a contatto con il citoplasma)
    \item $x = A$: parete esterna della membrana (a contatto con il mezzo extracellulare)
    \item $A \approx 5\,\text{nm}$: spessore della membrana
\end{itemize}

All'interno della membrana ($0 < x < A$) il campo elettrico è non nullo. Al di fuori (nel mezzo intracellulare o extracellulare), poiché la neutralità macroscopica è preservata, il campo elettrico è asintoticamente nullo.

\subsubsection{Definizione del Potenziale di Membrana}

Nell'esempio del potassio che fluisce verso l'esterno, le cariche positive si accumulano all'esterno e le negative all'interno. Il campo elettrico $\mathbf{E}$ è orientato dall'esterno verso l'interno (nella direzione $-x$). Poiché il campo è il gradiente negativo del potenziale, questo significa che il potenziale \textbf{decresce} andando dall'esterno verso l'interno:

\begin{equation}
\phi_{out} > \phi_{in}
\end{equation}

La \textbf{convenzione standard} per il potenziale di membrana è:

\begin{equation}
\boxed{V_m = V_{in} - V_{out} = \phi(0) - \phi(A)}
\end{equation}

In questa convenzione, per la situazione descritta sopra, $V_m < 0$. Questa è esattamente la situazione fisiologica: il potenziale di riposo dei neuroni è tipicamente compreso tra $-60\,\text{mV}$ e $-90\,\text{mV}$.

\subsubsection{Interpretazione in Termini di Correnti}

Con la convenzione stabilita:

\begin{itemize}
    \item Un \textbf{flusso di ioni positivi} nella direzione $+x$ (dall'interno verso l'esterno) corrisponde a una \textbf{corrente positiva}.
    \item Un \textbf{flusso di ioni positivi} nella direzione $-x$ (dall'esterno verso l'interno) corrisponde a una \textbf{corrente negativa}.
\end{itemize}

Per il potassio nel regime di equilibrio:
\begin{itemize}
    \item Il flusso diffusivo $J_D > 0$ (da dentro verso fuori): corrente positiva, guidata dal gradiente di concentrazione.
    \item Il flusso di deriva $J_E < 0$ (da fuori verso dentro): corrente negativa, guidata dal campo elettrico che tende a riportare i K$^+$ all'interno.
\end{itemize}

All'equilibrio, questi due contributi si annullano.

\section{Il Flusso di Deriva: La Legge di Ohm Riformulata}

\subsubsection{Velocità degli Ioni in un Campo Elettrico}

Consideriamo una specie ionica $s$ con valenza $Z_s$ e mobilità $\mu_s$. In presenza di un campo elettrico $\mathbf{E}$, gli ioni acquistano una velocità di deriva proporzionale al campo e alla loro mobilità:

\begin{equation}
\mathbf{v}_s = Z_s \mu_s \mathbf{E}
\end{equation}

Il fattore $Z_s$ tiene conto del segno della carica: uno ione positivo si muove nella direzione del campo, uno negativo nella direzione opposta. La mobilità $\mu_s$ è inversamente proporzionale alla resistenza del mezzo al moto degli ioni.

Poiché, come discusso nella lezione precedente, le variazioni del campo magnetico nel tempo sono trascurabili nelle condizioni neurofisiologiche, la legge di Maxwell-Faraday garantisce che il rotore del campo elettrico sia zero:

\begin{equation}
\nabla \times \mathbf{E} = 0
\end{equation}

Questo significa che $\mathbf{E}$ è un \textbf{campo conservativo} e può essere scritto come il gradiente negativo di un potenziale scalare $\phi$:

\begin{equation}
\mathbf{E} = -\nabla\phi
\end{equation}

Sostituendo:

\begin{equation}
\mathbf{v}_s = -Z_s \mu_s \nabla\phi
\end{equation}

\subsubsection{Espressione del Flusso di Deriva}

Il \textbf{flusso molare di deriva} $\mathbf{J}_E^{(s)}$ si ottiene moltiplicando la velocità per la concentrazione locale degli ioni:

\begin{equation}
\mathbf{J}_E^{(s)} = C_s \mathbf{v}_s = -Z_s \mu_s C_s \nabla\phi
\end{equation}

In forma monodimensionale (asse $x$ ortogonale alla membrana):

\begin{equation}
J_E^{(s)} = -Z_s \mu_s C(x) \frac{d\phi}{dx}
\end{equation}

\subsubsection{Collegamento con la Legge di Ohm}

Moltiplicando il flusso molare per $Z_s F$ (dove $F$ è la \textbf{costante di Faraday}, $F \approx 96485\,\text{C/mol}$), si ottiene la \textbf{densità di corrente} $i_s = Z_s F J_E^{(s)}$:

\begin{equation}
i_s = Z_s F \cdot \left(-Z_s \mu_s C_s \frac{d\phi}{dx}\right) = -Z_s^2 F \mu_s C_s \frac{d\phi}{dx}
\end{equation}

Questa è esattamente la \textbf{legge di Ohm} in forma differenziale: la corrente è proporzionale alla differenza di potenziale (o al suo gradiente), e la costante di proporzionalità è $1/R$, dove la resistenza è inversamente proporzionale alla mobilità $\mu_s$.



\section{Il Flusso Diffusivo: Prima Legge di Fick e Relazione di Einstein}

\subsubsection{Prima Legge di Fick}

Il flusso di ioni dovuto al gradiente di concentrazione è descritto dalla \textbf{prima legge di Fick} (enunciata da Adolf Fick nel 1855):

\begin{equation}
\mathbf{J}_D^{(s)} = -D_s \nabla C_s
\end{equation}

In forma monodimensionale:

\begin{equation}
J_D^{(s)} = -D_s \frac{dC_s}{dx}
\end{equation}

dove $D_s\,[\text{m}^2/\text{s}]$ è il \textbf{coefficiente di diffusione} della specie $s$. Il segno negativo indica che il flusso è diretto dalla regione a maggiore concentrazione verso quella a minore concentrazione (in direzione opposta al gradiente).

\subsubsection{Relazione di Einstein (Nernst-Einstein)}

Il coefficiente di diffusione non è indipendente dalla mobilità degli ioni. La \textbf{relazione di Einstein} (o relazione di Nernst-Einstein) mette in relazione $D_s$ con $\mu_s$:

\begin{equation}
\boxed{D_s = \frac{RT}{F} \mu_s}
\end{equation}

dove:
\begin{itemize}
    \item $R = 8.314\,\text{J/(mol K)}$: costante dei gas
    \item $T$: temperatura assoluta in Kelvin
    \item $F = 96485\,\text{C/mol}$: costante di Faraday
    \item $\mu_s$: mobilità della specie $s$
\end{itemize}

Questa relazione ha un contenuto fisico profondo: essa afferma che \textbf{diffusione e mobilità sono due aspetti dello stesso fenomeno microscopico}, ovvero l'agitazione termica delle molecole nel solvente. In particolare:

\begin{itemize}
    \item Maggiore è la temperatura $T$, maggiore è $D_s$ e quindi più intenso è il flusso diffusivo.
    \item Maggiore è la mobilità $\mu_s$ (ovvero, minore è la resistenza del mezzo), maggiore è la facilità con cui gli ioni si muovono sia sotto l'influenza del campo elettrico sia per diffusione.
\end{itemize}

Il flusso diffusivo diventa quindi:

\begin{equation}
J_D^{(s)} = -\frac{RT\mu_s}{F} \frac{dC_s}{dx}
\end{equation}


\section{L'Equazione di Nernst-Planck}

\subsubsection{Derivazione dell'Equazione Completa}

Combinando il flusso di deriva e il flusso diffusivo, si ottiene il \textbf{flusso molare totale} per la specie $s$:

\begin{equation}
\mathbf{J}^{(s)} = \mathbf{J}_E^{(s)} + \mathbf{J}_D^{(s)} = -Z_s \mu_s C_s \nabla\phi - \frac{RT\mu_s}{F} \nabla C_s
\end{equation}

Raccogliendo la mobilità $\mu_s$:

\begin{equation}
\boxed{\mathbf{J}^{(s)} = -\mu_s \left( Z_s C_s \nabla\phi + \frac{RT}{F} \nabla C_s \right)}
\end{equation}

Questa è l'\textbf{equazione di Nernst-Planck}, che descrive il trasporto ionico in presenza sia di un campo elettrico sia di un gradiente di concentrazione. Si tratta di uno strumento fondamentale non solo per la neurofisiologia, ma per tutta la biofisica delle membrane.

\subsubsection{L'Approssimazione Planare della Membrana}

La condizione $A \ll r$ (spessore della membrana molto minore del raggio dell'assone) ci permette di trattare ogni \textit{patch} di membrana come un \textbf{condensatore piano} con facce parallele infinite. Le conseguenze di questa approssimazione sono:

\begin{enumerate}
    \item Il campo elettrico all'interno della membrana è \textbf{sempre perpendicolare} alle facce.
    \item L'unica direzione spaziale rilevante è $x$, l'asse ortogonale alla membrana.
    \item I gradienti nelle direzioni tangenziali alla membrana sono trascurabili.
\end{enumerate}

L'equazione di Nernst-Planck si riduce quindi a una \textbf{forma monodimensionale}:

\begin{equation}
\boxed{J^{(s)} = -\mu_s \left( Z_s C_s(x) \frac{d\phi}{dx} + \frac{RT}{F} \frac{dC_s}{dx} \right)}
\end{equation}

dove $J^{(s)}$ è ora uno scalare (positivo verso l'esterno, negativo verso l'interno) e la derivata è rispetto alla coordinata $x$.



\section{Potenziale di Equilibrio per Singola Specie: Derivazione dell'Equazione di Nernst}

\subsubsection{Ipotesi Semplificative}

Prima di affrontare il caso generale con più specie ioniche, studiamo il caso più semplice: \textbf{una sola specie ionica} (ad esempio K$^+$) in condizioni di \textbf{equilibrio termodinamico}. Le ipotesi sono:

\begin{enumerate}
    \item \textbf{Una sola specie ionica:} ignoriamo per ora la presenza di Na$^+$, Cl$^-$, etc.
    \item \textbf{Equilibrio:} il flusso molare netto è zero, $J^{(K)} = 0$.
    \item \textbf{Concentrazioni ai bordi costanti:} $C(0) = C_{in}$ e $C(A) = C_{out}$ sono fissate (le pompe ioniche mantengono costanti le concentrazioni intracellulare ed extracellulare).
\end{enumerate}



Dalla condizione di equilibrio $J^{(K)} = 0$, l'equazione di Nernst-Planck diventa:

\begin{equation}
0 = -\mu \left( Z C \frac{d\phi}{dx} + \frac{RT}{F} \frac{dC}{dx} \right)
\end{equation}

Poiché $\mu \neq 0$, la parentesi deve essere zero:

\begin{equation}
Z C \frac{d\phi}{dx} + \frac{RT}{F} \frac{dC}{dx} = 0
\end{equation}

Separando le variabili e dividendo per $C$:

\begin{equation}
\frac{ZF}{RT} d\phi = -\frac{dC}{C} = -d(\ln C)
\end{equation}



Integrando entrambi i membri da $x = 0$ a $x = A$:

\begin{equation}
\frac{ZF}{RT} \int_0^A d\phi = -\int_0^A d(\ln C)
\end{equation}

\begin{equation}
\frac{ZF}{RT} \left[\phi(A) - \phi(0)\right] = \ln C(0) - \ln C(A) = \ln \frac{C_{in}}{C_{out}}
\end{equation}

Ricordando che $\phi(0) = V_{in}$, $\phi(A) = V_{out}$, e la convenzione $V_m = V_{in} - V_{out}$:

\begin{equation}
\frac{ZF}{RT}(V_{out} - V_{in}) = \ln \frac{C_{in}}{C_{out}}
\end{equation}

\begin{equation}
-\frac{ZF}{RT} V_m = \ln \frac{C_{in}}{C_{out}}
\end{equation}

Invertendo:

\begin{equation}
\boxed{E_s \equiv V_m^{eq} = \frac{RT}{ZF} \ln \frac{C_{out}}{C_{in}}}
\end{equation}

Questa è la celebre \textbf{equazione di Nernst}, che fornisce il \textbf{potenziale di equilibrio} (o \textbf{potenziale di Nernst}) per la specie ionica $s$. È il valore di $V_m$ al quale il flusso netto di quella specie attraverso la membrana è zero.

\subsection{Il Ruolo delle Pompe Ioniche nel Mantenimento dell'Equilibrio}

Un aspetto importante che emerge da questa analisi è il ruolo delle \textbf{pompe ioniche}. Nel caso di una singola specie ionica in equilibrio, la condizione $J = 0$ è effettivamente soddisfatta e le concentrazioni $C_{in}$ e $C_{out}$ non cambiano nel tempo. In questo caso le pompe ioniche non sarebbero necessarie.

Tuttavia, in presenza di più specie ioniche (come nella situazione fisiologica reale), il sistema non può essere simultaneamente all'equilibrio per tutte le specie, e quindi c'è sempre un flusso netto residuo. Le pompe ioniche intervengono per \textbf{mantenere costanti} le concentrazioni intracellulari ed extracellulari contro questi flussi.

\begin{tcolorbox}[colback=gray!5, colframe=gray!40, arc=2mm, boxrule=0.5pt]
 Le pompe ioniche sono dispositivi attivi: consumano energia metabolica (ATP) per trasportare ioni contro il loro gradiente elettrochimico. Tuttavia  il lavoro che devono svolgere è sorprendentemente piccolo: è per questo motivo che il cervello umano consuma solo circa 20 W (o 50 W se si considera l'intero metabolismo cerebrale). Le pompe lavorano molto vicino all'equilibrio e il numero di cariche che devono essere compensate ad ogni spike è una frazione minuscola del totale.
\end{tcolorbox}

\section{Il Potenziale di Reversione: Numeri Reali}

\subsubsection{Significato Fisico del Potenziale di Nernst}

Il potenziale di Nernst $E_s$ è anche detto \textbf{potenziale di reversione} (\textit{reversal potential}) perché rappresenta il valore di $V_m$ al quale la corrente portata dalla specie $s$ si inverte di segno: per $V_m > E_s$ la corrente fluisce in una direzione, per $V_m < E_s$ nella direzione opposta.

Questo concetto è di grande utilità sperimentale: iniettando una corrente esterna nel neurone e variando il potenziale di membrana, si può determinare il potenziale di reversione di una specifica conduttanza ionica come il valore di $V_m$ al quale quella corrente cambia segno.

\subsubsection{Calcolo dei Potenziali di Nernst per le Principali Specie Ioniche}

La costante $RT/F$ assume un valore determinato dalla temperatura. A $T = 37^\circ\text{C} = 310\,\text{K}$ (temperatura corporea dei mammiferi):

\begin{equation}
\frac{RT}{F} = \frac{8.314 \times 310}{96485} \approx 26.7\,\text{mV}
\end{equation}

Il fattore numerico è quindi $RT/F \approx 26.7\,\text{mV}$, che viene talvolta approssimato a $25\,\text{mV}$ (a $20^\circ\text{C}$) o $27\,\text{mV}$ (a $37^\circ\text{C}$). Il professore utilizza nella sua derivazione la notazione $RT/F \approx 62\,\text{mV}$ per $Z=1$, che deriva dalla formula con il logaritmo in base 10: $2.303 \times RT/F \approx 61.5\,\text{mV}$, arrotondato a $62\,\text{mV}$.

\textbf{Potenziale di Nernst del Potassio (mammifero):}

\begin{equation}
E_K = \frac{RT}{F} \ln \frac{C_K^{out}}{C_K^{in}} = 62\,\text{mV} \cdot \log_{10}\frac{5}{140} \approx -89\,\text{mV}
\end{equation}

\textbf{Potenziale di Nernst del Sodio (mammifero):}

\begin{equation}
E_{Na} = \frac{RT}{F} \ln \frac{C_{Na}^{out}}{C_{Na}^{in}} = 62\,\text{mV} \cdot \log_{10}\frac{142}{10} \approx +71\,\text{mV}
\end{equation}

Il professore cita $\approx +100\,\text{mV}$ come ordine di grandezza per il potenziale di Nernst del sodio.

\textbf{Potenziale di Nernst del Cloruro (mammifero):}

\begin{equation}
E_{Cl} = \frac{RT}{(-1)F} \ln \frac{C_{Cl}^{out}}{C_{Cl}^{in}} = -62\,\text{mV} \cdot \log_{10}\frac{110}{10} \approx -62\,\text{mV}
\end{equation}

Il cloruro, avendo valenza $-1$, ha il potenziale di Nernst con segno opposto rispetto a quanto ci si aspetterebbe dalla sola distribuzione di concentrazione: il risultato è circa $-70\,\text{mV}$, molto vicino al potenziale di Nernst del potassio.

\textbf{Potenziale di Nernst del Calcio (mammifero):}

\begin{equation}
E_{Ca} = \frac{RT}{2F} \ln \frac{C_{Ca}^{out}}{C_{Ca}^{in}} = \frac{62}{2}\,\text{mV} \cdot \log_{10}\frac{2.5}{10^{-4}} \approx +136\,\text{mV}
\end{equation}

Il potenziale di Nernst del calcio è elevatissimo e positivo, a causa della concentrazione intracellulare quasi nulla.


\subsubsection{Il Problema del Compromesso}

L'osservazione cruciale è che \textbf{non esiste un singolo valore di $V_m$ che soddisfi simultaneamente le condizioni di equilibrio per K$^+$, Na$^+$ e Cl$^-$}. In particolare:

\begin{itemize}
    \item $E_K \approx -89\,\text{mV}$: il potenziale di equilibrio del potassio è molto negativo.
    \item $E_{Na} \approx +70\text{ a }+100\,\text{mV}$: il potenziale di equilibrio del sodio è molto positivo.
    \item $E_{Cl} \approx -70\,\text{mV}$: il potenziale di equilibrio del cloruro è vicino a quello del potassio.
\end{itemize}

Il neurone deve trovare un \textbf{compromesso} tra questi potenziali in competizione. Questo compromesso determina il \textbf{potenziale di riposo} del neurone, che è sempre compreso tra i valori estremi dei potenziali di Nernst. Ma questo compromesso implica necessariamente che ci sia sempre un \textbf{flusso netto} di almeno alcune specie ioniche attraverso la membrana, che deve essere compensato dalle pompe ioniche.

\begin{tcolorbox}[colback=gray!5, colframe=gray!40, arc=2mm, boxrule=0.5pt]
 Il fatto che non si possa trovare un potenziale simultaneamente all'equilibrio per il sodio e il potassio è la radice profonda del fatto che il neurone è un sistema \textit{fuori dall'equilibrio}. Questo non è un difetto: è il motivo per cui i neuroni \textit{funzionano}.
\end{tcolorbox}

\section{Struttura delle Pompe Ioniche}

Le \textbf{pompe ioniche} sono proteine transmembrana che trasportano attivamente ioni contro il loro gradiente elettrochimico, consumando energia metabolica sotto forma di ATP. A differenza dei canali ionici (che consentono un flusso passivo secondo il gradiente), le pompe eseguono un lavoro meccanico-chimico per spostare gli ioni "in salita".

L'esempio più importante è la \textbf{Na$^+$/K$^+$-ATPasi} (pompa sodio-potassio):

\begin{itemize}
    \item Espelle \textbf{3 ioni Na$^+$} verso l'esterno per ogni ciclo
    \item Importa \textbf{2 ioni K$^+$} verso l'interno per ogni ciclo
    \item Idrolizza \textbf{1 molecola di ATP} per ogni ciclo
    \item Il trasporto è elettrogenico (carica netta $-1$ per ciclo), contribuendo leggermente alla negatività del potenziale di membrana
\end{itemize}


Un risultato sorprendente che emerge dall'analisi quantitativa è che le pompe ioniche svolgono un lavoro \textbf{molto piccolo} rispetto a quanto si potrebbe intuitivamente pensare. Il motivo è che il numero di ioni che attraversano la membrana durante un evento elettrico (come un potenziale d'azione) è una frazione minuscola del pool ionico totale disponibile.

Questa efficienza è alla base del basso consumo energetico del cervello: nonostante la complessità dei calcoli eseguiti, il cervello umano consuma solo circa \textbf{20 W} di potenza metabolica.

\section{Stima dell'Energia delle Pompe Ioniche: Quante Cariche si Scambiano?}

Per valutare quantitativamente il lavoro delle pompe ioniche, consideriamo un segmento cilindrico di assone con raggio $r$ e lunghezza $l$. Vogliamo calcolare il rapporto tra le cariche scambiate attraverso la membrana durante un potenziale d'azione e le cariche totali contenute nel volume intracellulare.



La \textbf{superficie laterale} del cilindro (la parte di membrana che ci interessa) è:

\begin{equation}
S = 2\pi r l
\end{equation}

Quando il potenziale di membrana cambia di $\Delta V$ (durante uno spike, $\Delta V \approx 100\,\text{mV}$), la variazione di carica sul condensatore-membrana è:

\begin{equation}
\Delta Q = C_m \cdot S \cdot \Delta V = C_m \cdot 2\pi r l \cdot \Delta V
\end{equation}

con $C_m \approx 1\,\mu\text{F/cm}^2 = 10^{-2}\,\text{F/m}^2$ la capacità specifica della membrana.



Il \textbf{volume} del cilindro è:

\begin{equation}
V_{cil} = \pi r^2 l
\end{equation}

Le cariche totali associate agli ioni K$^+$ all'interno (consideriamo solo il potassio come specie dominante, con $Z_K = 1$) sono:

\begin{equation}
Q_{tot} = F \cdot C_K^{in} \cdot \pi r^2 l
\end{equation}

dove $F = 96485\,\text{C/mol}$ è la costante di Faraday.



\begin{equation}
\frac{\Delta Q}{Q_{tot}} = \frac{C_m \cdot 2\pi r l \cdot \Delta V}{F \cdot C_K^{in} \cdot \pi r^2 l} = \frac{2 C_m \Delta V}{F \cdot C_K^{in} \cdot r}
\end{equation}

Sostituendo i valori numerici per l'\textbf{assone di mammifero} ($r \approx 5\,\mu\text{m}$, $C_K^{in} = 140\,\text{mM} = 140\,\text{mol/m}^3$):

\begin{equation}
\frac{\Delta Q}{Q_{tot}} \approx \frac{2 \times 10^{-2} \times 0.1}{96485 \times 140 \times 5 \times 10^{-6}} \approx 10^{-3}
\end{equation}

Per l'\textbf{assone gigante del calamaro} (raggio $r \approx 500\,\mu\text{m}$, molto più grande):

\begin{equation}
\frac{\Delta Q}{Q_{tot}} \approx 10^{-7}
\end{equation}

\subsubsection{Interpretazione Biologica}

Il risultato è notevole: anche durante l'evento più energetico nella vita di un neurone (un potenziale d'azione con $\Delta V \approx 100\,\text{mV}$), \textbf{solo 1/1000 degli ioni} (in un assone di mammifero) devono essere "rimpiazzati" dalle pompe ioniche. Per l'assone gigante del calamaro, la frazione è addirittura di $10^{-7}$.

Questo significa che, a una frequenza di scarica di 1 Hz (un potenziale d'azione al secondo), le pompe ioniche devono compensare solo una variazione relativa dello 0.1\% delle concentrazioni ioniche. L'energia richiesta è quindi minima. Questo è il motivo per cui il cervello umano — nonostante elabori informazioni in modo straordinariamente complesso — consuma solo circa \textbf{20 W} di potenza metabolica netta.


\section{Il Modello di Goldman-Hodgkin-Katz (GHK)}

\subsubsection{Necessità di un Modello Multi-Ionico}

Come abbiamo visto, la realtà fisiologica prevede la presenza contemporanea di più specie ioniche (K$^+$, Na$^+$, Cl$^-$). In questo caso, non si può più parlare di un singolo potenziale di equilibrio: il sistema è in uno \textbf{stato stazionario} (non di equilibrio) in cui i flussi individuali di ciascuna specie sono non nulli, ma la corrente totale si azzera (condizione stazionaria).

Nel 1943, Goldman, e successivamente Hodgkin e Katz nel 1949, derivarono un'equazione per il potenziale di membrana stazionario a partire dall'equazione di Nernst-Planck, sotto tre ipotesi fondamentali.



\textbf{Ipotesi 1 — Campo elettrico costante nella membrana:}\\
In condizioni stazionarie, il campo elettrico all'interno della membrana è uniforme:

\begin{equation}
\frac{d\phi}{dx} = -\frac{V_m}{A} = \text{costante}
\end{equation}

Questo equivale ad assumere un profilo di potenziale \textbf{lineare} all'interno della membrana:

\begin{equation}
\phi(x) = V_{in} - \frac{V_m}{A} x
\end{equation}

Questa ipotesi è giustificata dal fatto che la membrana si comporta come un condensatore piano in condizioni stazionarie.

\textbf{Ipotesi 2 — Indipendenza dei flussi:}\\
I flussi delle diverse specie ioniche sono \textbf{indipendenti}: gli ioni K$^+$, Na$^+$ e Cl$^-$ si muovono senza influenzarsi reciprocamente. Questa è un'approssimazione, non rigorosamente vera, ma eccellente nella maggior parte delle condizioni fisiologiche.

\textbf{Ipotesi 3 — Concentrazioni stazionarie:}\\
Le concentrazioni ioniche ai bordi della membrana sono mantenute costanti dalle pompe ioniche:

\begin{equation}
C_s(0) = C_s^{in} = \text{costante}, \quad C_s(A) = C_s^{out} = \text{costante}
\end{equation}

\subsubsection{Condizione Stazionaria Multi-Ionica}

In condizioni stazionarie, il flusso molare di ogni specie $J_s$ è costante nel tempo e nello spazio (non dipende da $x$), ma è in generale \textbf{non nullo}:

\begin{equation}
\frac{\partial J_s}{\partial x} = 0, \quad J_s = \text{costante} \neq 0
\end{equation}

La condizione di neutralità della corrente totale è:

\begin{equation}
\sum_s Z_s F J_s = 0
\end{equation}

Questa condizione è fondamentalmente diversa da quella di equilibrio per singola specie, in cui si imponeva $J_s = 0$ per ogni specie. Qui ogni specie ha un flusso netto non nullo, ma i contributi alla corrente totale si compensano.


\subsection{Derivazione dell'Equazione GHK per la Corrente}


Partiamo dall'equazione di Nernst-Planck monodimensionale per la specie $s$, con $J_s$ costante:

\begin{equation}
J_s = -\mu_s \left( Z_s C_s \phi' + \frac{RT}{F} C_s' \right)
\end{equation}

dove gli apici denotano derivate rispetto a $x$.

Riorganizzando:

\begin{equation}
C_s' + \frac{Z_s F}{RT} \phi' C_s = -\frac{J_s F}{RT \mu_s}
\end{equation}

Definiamo i parametri:

\begin{equation}
\alpha \equiv \frac{Z_s F}{RT} \phi', \quad \beta \equiv -\frac{J_s F}{RT \mu_s}
\end{equation}

L'equazione diventa una \textbf{EDO lineare del primo ordine} in $C_s(x)$:

\begin{equation}
C_s' + \alpha C_s = \beta
\end{equation}



Introduciamo il \textbf{fattore integrante}:

\begin{equation}
\lambda(x) = \exp\left(\int_0^x \alpha(x')\,dx'\right) = \exp\left(\frac{Z_s F}{RT}\int_0^x \phi'(x')\,dx'\right) = \exp\left(\frac{Z_s F}{RT}(\phi(x) - V_{in})\right)
\end{equation}

La soluzione dell'EDO si scrive come:

\begin{equation}
C_s(x) = \frac{1}{\lambda(x)}\left[C_s(0)\lambda(0) + \beta\int_0^x \lambda(x')\,dx'\right]
\end{equation}

Poiché $\lambda(0) = \exp(0) = 1$, e $C_s(0) = C_{in}$:

\begin{equation}
C_s(x) = \frac{1}{\lambda(x)}\left[C_{in} + \beta\int_0^x \lambda(x')\,dx'\right]
\end{equation}

Con l'ipotesi di campo elettrico costante, il potenziale varia linearmente:

\begin{equation}
\phi(x) = V_{in} - \frac{V_m}{A}x
\end{equation}

Il fattore integrante diventa:

\begin{equation}
\lambda(x) = \exp\left(\frac{Z_s F}{RT}\left(-\frac{V_m}{A}x\right)\right) = \exp\left(-\frac{Z_s F V_m}{RT \cdot A}x\right)
\end{equation}

L'integrale di $\lambda$ da 0 ad $A$:

\begin{equation}
\int_0^A \lambda(x)\,dx = \int_0^A \exp\left(-\frac{Z_s F V_m}{RT A}x\right)dx
\end{equation}

Questo è un integrale di un'esponenziale con esponente lineare in $x$. Il risultato è:

\begin{equation}
\int_0^A \lambda(x)\,dx = \frac{RT A}{Z_s F V_m}\left(1 - e^{-Z_s F V_m/(RT)}\right)
\end{equation}

Definiamo $\gamma \equiv F/(RT)$ (in $\text{mV}^{-1}$), cosicché $1/\gamma = RT/F \approx 26\,\text{mV}$ a $300\,\text{K}$. Allora:

\begin{equation}
\int_0^A \lambda(x)\,dx = \frac{A}{Z_s \gamma V_m}\left(1 - e^{-Z_s \gamma V_m}\right)
\end{equation}

\subsubsection{Espressione per il Flusso Molare}

Applicando la condizione al contorno $C_s(A) = C_{out}$:

\begin{equation}
C_{out} = \frac{1}{\lambda(A)}\left[C_{in} + \beta \int_0^A \lambda(x')\,dx'\right]
\end{equation}

Ricordando che $\lambda(A) = e^{-Z_s \gamma V_m}$ e risolvendo per $\beta = -J_s F/(RT\mu_s)$, si ottiene:

\begin{equation}
J_s = \frac{\mu_s Z_s  V_m}{A} \cdot \frac{C_{in} - C_{out}\,e^{-Z_s \gamma V_m}}{1 - e^{-Z_s \gamma V_m}}
\end{equation}

\subsubsection{Verifica: Recupero dell'Equazione di Nernst}

Per il caso di singola specie all'equilibrio ($J_s = 0$), la condizione $C_{in} - C_{out}\,e^{-Z_s \gamma V_m} = 0$ dà:

\begin{equation}
e^{-Z_s \gamma V_m} = \frac{C_{in}}{C_{out}}
\end{equation}

\begin{equation}
-Z_s \gamma V_m = \ln\frac{C_{in}}{C_{out}}
\end{equation}

\begin{equation}
V_m^{eq} = \frac{1}{Z_s \gamma} \ln\frac{C_{out}}{C_{in}} = \frac{RT}{Z_s F} \ln\frac{C_{out}}{C_{in}}
\end{equation}

Questo è esattamente l'equazione di Nernst. Il modello GHK generalizza l'equazione di Nernst al caso multi-ionico.

\subsubsection{Calcolo Esplicito dell'Integrale di Lambda}

Mostriamo in dettaglio come si calcola l'integrale del fattore integrante all'interno della membrana, un passaggio cruciale della derivazione.

\begin{equation}
\int_0^A \lambda(x)\,dx = \int_0^A \exp\left(\frac{Z_s F}{RT}(\phi(x) - V_{in})\right)dx
\end{equation}

Poiché $\phi(x) - V_{in} = -\frac{V_m}{A}x$, l'integrale diventa:

\begin{equation}
\int_0^A \exp\left(-\frac{Z_s F V_m}{RT A}x\right)dx = \left[\frac{-RT A}{Z_s F V_m}\exp\left(-\frac{Z_s F V_m}{RT A}x\right)\right]_0^A
\end{equation}

\begin{equation}
= \frac{-RT A}{Z_s F V_m}\left(e^{-Z_s \gamma V_m} - 1\right) = \frac{RT A}{Z_s F V_m}\left(1 - e^{-Z_s \gamma V_m}\right) = \frac{A}{Z_s \gamma V_m}\left(1 - e^{-Z_s \gamma V_m}\right)
\end{equation}

\subsubsection{L'Equazione di Corrente GHK}

La \textbf{densità di corrente} per la specie $s$ si ottiene moltiplicando il flusso molare per $Z_s F$:

\begin{equation}
I_s = Z_s F J_s = Z_s^2 \frac{F \mu_s}{A} \cdot \frac{ V_m \left(C_{in} - C_{out}\,e^{-Z_s \gamma V_m}\right)}{1 - e^{-Z_s \gamma V_m}}
\end{equation}

Poiché $Z^2 = 1$ per tutte le specie con valenza $\pm 1$, possiamo scrivere in forma compatta:

\begin{equation}
\boxed{I_s = \frac{F \mu_s}{A} \cdot \frac{ V_m \left(C_s^{in} - C_s^{out}\,e^{-Z_s \gamma V_m}\right)}{1 - e^{-Z_s \gamma V_m}}}
\end{equation}

Questa è l'\textbf{equazione di corrente di Goldman-Hodgkin-Katz} per la specie $s$. Essa descrive come la corrente ionica dipende dal potenziale di membrana $V_m$ e dalle concentrazioni intracellulari ed extracellulari della specie $s$.

Il parametro $\gamma = F/(RT)$ ha le dimensioni di $[\text{mV}^{-1}]$. A temperatura ambiente ($T \approx 300\,\text{K}$):

\begin{equation}
\frac{1}{\gamma} = \frac{RT}{F} \approx 26\,\text{mV}
\end{equation}

a $37^\circ\text{C}$ ($T = 310\,\text{K}$):

\begin{equation}
\frac{1}{\gamma} = \frac{RT}{F} \approx 26.7\,\text{mV}
\end{equation}


\subsection{Il Potenziale di Equilibrio GHK: Il Potenziale di Riposo del Neurone}

\subsubsection{Condizione di Stazionarietà della Corrente Totale}

In condizioni stazionarie, la corrente totale attraverso la membrana deve essere zero:

\begin{equation}
I_{tot} = I_K + I_{Na} + I_{Cl} = 0
\end{equation}

Sostituendo le espressioni GHK per ogni specie, con $Z_K = Z_{Na} = +1$ e $Z_{Cl} = -1$:

Per K$^+$ e Na$^+$ ($Z = +1$):
\begin{equation}
I_s \propto \frac{ V_m (C_s^{in} - C_s^{out}\,e^{-\gamma V_m})}{1 - e^{-\gamma V_m}}
\end{equation}

Per Cl$^-$ ($Z = -1$):
\begin{equation}
I_{Cl} \propto \frac{ V_m (C_{Cl}^{in} - C_{Cl}^{out}\,e^{+\gamma V_m})}{1 - e^{+\gamma V_m}}
\end{equation}

\subsubsection{Algebrizzazione e Soluzione}

Per uniformare le espressioni, moltiplichiamo il numeratore e il denominatore dell'espressione per Cl$^-$ per $e^{+\gamma V_m}$:

\begin{equation}
I_{Cl} \propto \frac{ V_m (C_{Cl}^{in}\,e^{+\gamma V_m} - C_{Cl}^{out})}{e^{+\gamma V_m} - 1}
\end{equation}

Notiamo che $e^{+\gamma V_m} - 1 = -(1 - e^{+\gamma V_m})$ e moltiplichiamo per $e^{-\gamma V_m}$:

\begin{equation}
I_{Cl} \propto \frac{ V_m (C_{Cl}^{out} - C_{Cl}^{in}\,e^{-\gamma V_m})}{1 - e^{-\gamma V_m}}
\end{equation}

Questo ci consente di scrivere tutte e tre le correnti con lo stesso denominatore $1 - e^{-\gamma V_m}$, che è diverso da zero (poiché $V_m \neq 0$ sperimentalmente). Raccogliendo i termini che moltiplicano $e^{-\gamma V_m}$ e quelli indipendenti dall'esponenziale, e imponendo $I_{tot} = 0$:

\begin{equation}
\left(\mu_K C_K^{in} + \mu_{Na} C_{Na}^{in} + \mu_{Cl} C_{Cl}^{out}\right) e^{-\gamma V_m} = \mu_K C_K^{out} + \mu_{Na} C_{Na}^{out} + \mu_{Cl} C_{Cl}^{in}
\end{equation}

Prendendo il logaritmo e invertendo il segno:

\begin{equation}
\boxed{V_m^{rest} = \frac{RT}{F} \ln \frac{\mu_K C_K^{out} + \mu_{Na} C_{Na}^{out} + \mu_{Cl} C_{Cl}^{in}}{\mu_K C_K^{in} + \mu_{Na} C_{Na}^{in} + \mu_{Cl} C_{Cl}^{out}}}
\end{equation}

Questa è la \textbf{equazione di Goldman-Hodgkin-Katz} per il potenziale di riposo della membrana neuronale.

Confrontando con l'equazione di Nernst per singola specie, si notano due differenze fondamentali:

\begin{enumerate}
    \item Al numeratore e al denominatore compaiono le concentrazioni \textbf{pesate dalle mobilità} $\mu_s$ di ciascuna specie.
    \item Per il cloruro ($Z_{Cl} = -1$), le concentrazioni intracellulare ed extracellulare sono \textbf{scambiate} rispetto a K$^+$ e Na$^+$: al numeratore compare $\mu_{Cl} C_{Cl}^{in}$ (anziché $C_{Cl}^{out}$), e al denominatore $\mu_{Cl} C_{Cl}^{out}$ (anziché $C_{Cl}^{in}$). Questo riflette il segno negativo della valenza del cloruro.
\end{enumerate}

Il potenziale di riposo è quindi una \textbf{media logaritmica pesata} dei potenziali di Nernst delle singole specie, con i pesi determinati dalle mobilità.

\subsubsection{Valori Numerici e Interpretazione}

Per i neuroni corticali di mammifero, utilizzando le mobilità relative $\mu_K : \mu_{Na} : \mu_{Cl} = 1.0 : 0.03 : 0.1$ (in unità arbitrarie) e le concentrazioni della tabella in Sezione 4.2, l'equazione GHK dà:

\begin{equation}
V_m^{rest} \approx -72\,\text{mV}
\end{equation}

Questo valore è in accordo con le misure sperimentali del potenziale di riposo nei neuroni corticali di mammifero.

L'interpretazione geometrica è illuminante: $-72\,\text{mV}$ si trova tra i potenziali di Nernst di K$^+$ ($-89\,\text{mV}$) e Cl$^-$ ($\approx -70\,\text{mV}$), ed è molto distante dal potenziale di Nernst di Na$^+$ ($\approx +70\,\text{mV}$). Questo riflette direttamente le mobilità: la mobilità del sodio è circa \textbf{33 volte più piccola} di quella del potassio ($\mu_{Na}/\mu_K = 0.03$), il che significa che i canali del sodio contribuiscono pochissimo al potenziale di riposo.
La bassa mobilità del sodio nel neurone a riposo riflette il fatto che i canali del sodio sono quasi tutti chiusi (nello stato cosiddetto di "inattivazione") al potenziale di riposo. Quando questi canali si aprono — come accade durante il potenziale d'azione — la mobilità del sodio aumenta bruscamente e il potenziale di membrana si sposta rapidamente verso $E_{Na}$.



\newpage 
\section{Sintesi: Dall'Equazione di Nernst alla GHK}


Il percorso seguito in questa lezione ha costruito due risultati principali, che differiscono per il numero di specie ioniche considerate e per la condizione imposta:

\textbf{Caso 1 — Singola specie ionica all'equilibrio:}
\begin{itemize}
    \item Condizione: $J_s = 0$ (flusso netto nullo)
    \item Risultato: \textbf{Equazione di Nernst}
\end{itemize}

\begin{equation}
E_s = \frac{RT}{Z_s F} \ln \frac{C_s^{out}}{C_s^{in}}
\end{equation}

\begin{itemize}
    \item Significato: il potenziale di membrana al quale il flusso di deriva e il flusso diffusivo si compensano esattamente.
    \item Tipo di stato: \textbf{equilibrio termodinamico} (nessun flusso, nessun consumo di energia).
\end{itemize}

\textbf{Caso 2 — Più specie ioniche in condizioni stazionarie:}
\begin{itemize}
    \item Condizione: $\sum_s Z_s F J_s = 0$ (corrente totale nulla, ma flussi individuali non nulli)
    \item Risultato: \textbf{Equazione di Goldman-Hodgkin-Katz}
\end{itemize}

\begin{equation}
V_m^{rest} = \frac{RT}{F} \ln \frac{\mu_K C_K^{out} + \mu_{Na} C_{Na}^{out} + \mu_{Cl} C_{Cl}^{in}}{\mu_K C_K^{in} + \mu_{Na} C_{Na}^{in} + \mu_{Cl} C_{Cl}^{out}}
\end{equation}

\begin{itemize}
    \item Significato: il potenziale di membrana al quale le correnti ioniche delle diverse specie si compensano, ma ogni singola specie ha ancora un flusso netto non nullo.
    \item Tipo di stato: \textbf{stato stazionario} (non equilibrio) — mantenuto attivamente dalle pompe ioniche.
\end{itemize}


La differenza tra i due casi è \textbf{concettualmente profonda}, non solo matematica:

\begin{itemize}
    \item Nel \textbf{caso di una sola specie}, il sistema può effettivamente raggiungere l'equilibrio termodinamico, in cui la corrente è zero e il sistema non ha bisogno di energia esterna per mantenere questo stato.
    \item Nel \textbf{caso multi-ionico}, il sistema è strutturalmente incapace di raggiungere l'equilibrio termodinamico (perché non esiste un $V_m$ che soddisfi simultaneamente le condizioni di Nernst per tutte le specie), e quindi è sempre in uno \textbf{stato stazionario lontano dall'equilibrio}. Questo stato richiede un apporto continuo di energia — fornito dalle pompe ioniche — per essere mantenuto.
\end{itemize}

Il potenziale di riposo di un neurone non è un equilibrio: è uno \textbf{stato attivo} che il neurone mantiene a costo di consumo energetico. Questo ha implicazioni profonde: un neurone privo di ATP (e quindi senza pompe ioniche funzionanti) perde rapidamente i suoi gradienti ionici e il suo potenziale di membrana.



I risultati di questa lezione pongono le basi per i temi delle prossime lezioni:

\begin{itemize}
    \item \textbf{Canali voltaggio-dipendenti:} se la mobilità $\mu_s$ non è costante ma dipende da $V_m$, si ottengono fenomeni non lineari come il potenziale d'azione.
    \item \textbf{Modello di Hodgkin-Huxley:} formalizzazione matematica dei canali voltaggio-dipendenti per Na$^+$ e K$^+$, che spiega la generazione e propagazione del potenziale d'azione.
    \item \textbf{Trasmissione sinaptica:} il ruolo del Ca$^{2+}$ nel rilascio dei neurotrasmettitori e la modulazione della sinapsi.
\end{itemize}


\section*{Tabella Riepilogativa}
\addcontentsline{toc}{section}{Tabella Riepilogativa}

\begin{table}[htbp]
    \centering
    \footnotesize
    \renewcommand{\arraystretch}{1.3}
    \begin{tabularx}{\textwidth}{|l|X|X|}
    \hline
    \textbf{Quantità}  &  \textbf{Equazione} & \textbf{Valore tipico (mammifero, 37$^\circ$C)} \\
    \hline
    Potenziale di membrana  $V_m$  & $V_m = V_{in} - V_{out}$ & Da $-$90 a $-$50 mV \\
    \hline
    Potenziale di Nernst K$^+$  $E_K$ & $\frac{RT}{F}\ln\frac{C_K^{out}}{C_K^{in}}$ & $-$89 mV \\
    \hline
    Potenziale di Nernst Na$^+$  $E_{Na}$  & $\frac{RT}{F}\ln\frac{C_{Na}^{out}}{C_{Na}^{in}}$ & +62 a +100 mV \\
    \hline
    Potenziale di Nernst Cl$^-$  $E_{Cl}$ & $\frac{RT}{(-1)F}\ln\frac{C_{Cl}^{out}}{C_{Cl}^{in}}$ & $-$62 a $-$70 mV \\
    \hline
    Potenziale di Nernst Ca$^{2+}$  $E_{Ca}$ & $\frac{RT}{2F}\ln\frac{C_{Ca}^{out}}{C_{Ca}^{in}}$ & $\approx$ +136 mV \\
    \hline
    Potenziale di riposo (GHK)  $V_m^{rest}$  & $\frac{RT}{F}\ln\frac{\mu_K C_K^{out}+\mu_{Na} C_{Na}^{out}+\mu_{Cl} C_{Cl}^{in}}{\mu_K C_K^{in}+\mu_{Na} C_{Na}^{in}+\mu_{Cl} C_{Cl}^{out}}$ & $-$72 mV \\
    \hline
    Capacità specifica membrana  $C_m$  & — & 1 $\mu$F/cm$^2$ \\
    \hline
    Spessore membrana  $A$  & — & $\approx$ 5 nm \\
    \hline
    Costante termica  $RT/F$  & $R \cdot T / F$ & $\approx$ 26.7 mV (37$^\circ$C) \\
    \hline
    Fattore gamma  $\gamma$ & $\gamma = F/(RT)$ & $\approx$ 37.5 mV$^{-1}$ (37$^\circ$C) \\
    \hline
    Coefficiente di diffusione  $D_s$ &  $D_s = \frac{RT}{F}\mu_s$ & Variabile per specie \\
    \hline
    Flusso molare totale  $J^{(s)}$ & $-\mu_s\left(Z_s C_s \frac{d\phi}{dx}+\frac{RT}{F}\frac{dC_s}{dx}\right)$ & - \\
    \hline
    Corrente GHK  $I_s$ & $\frac{F\mu_s}{A}\frac{\gamma V_m(C_s^{in}-C_s^{out}e^{-Z_s\gamma V_m})}{1-e^{-Z_s\gamma V_m}}$ & - \\
    \hline
    
    Consumo cerebrale  $P$ & — & $\approx$ 20 W \\
    \hline
    Mobilità relative  $\mu_K:\mu_{Na}:\mu_{Cl}$ & — & 1.0 : 0.03 : 0.1 \\
    \hline
    \end{tabularx}
\end{table}


\chapter{Lezione 4}


\textbf{Argomenti:} Equazioni di Goldman-Hodgkin-Katz, approssimazione ohmica, circuito equivalente passivo della membrana, equazione del neurone RC, esperimento di Hodgkin-Huxley, voltage clamp, space clamp, assone gigante di calamaro, correnti ioniche Na$^+$/K$^+$, TTX e TEA, circuito equivalente HH, variabili di gate n, m, h, modello di Markov per i canali ionici

\vspace{0.5cm}
\noindent\rule{\textwidth}{0.4pt}
\vspace{0.5cm}

\section{Riepilogo: Equazioni di Goldman-Hodgkin-Katz}


Il professore riprende dal materiale presentato nella lezione precedente, riepilogando due equazioni fondamentali dell'elettrofisiologia della membrana: l'\textbf{equazione del potenziale di equilibrio di Goldman-Hodgkin-Katz} (GHK) e l'\textbf{equazione della corrente di Goldman-Hodgkin-Katz}.

Entrambe le equazioni descrivono il comportamento elettrico di una membrana biologica in condizioni stazionarie, quando la membrana è permeabile a più specie ioniche simultaneamente.



In condizioni stazionarie, il \textbf{potenziale di equilibrio di membrana} $E_m$ risultante dalla presenza simultanea di più specie ioniche (tipicamente K$^+$, Na$^+$, Cl$^-$) è dato dall'equazione di Goldman-Hodgkin-Katz. Questo potenziale dipende dalla \textbf{temperatura} $T$ e dalla \textbf{mobilità ionica} di ciascuna specie (che è l'inverso della resistenza).

\begin{tcolorbox}[colback=gray!5, colframe=gray!40, arc=2mm, boxrule=0.5pt]
Al numeratore dell'equazione GHK si trovano sempre le concentrazioni esterne, con l'unica eccezione del Cl$^-$, per il quale compare la concentrazione interna. Questo segno invertito riflette la carica negativa dello ione cloruro.
\end{tcolorbox}


\subsubsection{Dati Sperimentali dell'Assone Gigante di Calamaro}

Come riferimento canonico, il professore utilizza i dati dell'\textbf{assone gigante di calamaro} (\textit{Loligo vulgaris}), il sistema modello storico di Hodgkin e Huxley. Le condizioni sperimentali tipiche sono:

\begin{table}[htbp]
    \centering
    \renewcommand{\arraystretch}{1.3}
    \begin{tabularx}{\textwidth}{|l|X|X|X|}
    \hline
    \textbf{Specie ionica} & \textbf{Concentrazione interna} & \textbf{Concentrazione esterna} & \textbf{Potenziale di inversione} \\
    \hline
    K$^+$ & 400 mM & 20 mM & $E_K \approx -76$ mV \\
    Na$^+$ & 50 mM & 440 mM & $E_{Na} \approx +52$ mV \\
    Cl$^-$ & 52 mM & 560 mM & $E_{Cl} \approx -60$ mV \\
    \hline
    \end{tabularx}
\end{table}

\begin{itemize}
    \item \textbf{Temperatura sperimentale:} 6.3 $^\circ$C (condizioni HH storiche)
    \item \textbf{Potenziale di equilibrio di membrana:} $E_m \approx -60$ mV
\end{itemize}

A questo potenziale:
\begin{itemize}
    \item La corrente di K$^+$ è \textbf{uscente} ($E_m > E_K$, quindi la forza motrice è verso l'esterno)
    \item La corrente di Cl$^-$ è anch'essa \textbf{uscente} ($E_m > E_{Cl}$)
    \item La corrente di Na$^+$ è \textbf{entrante} ($E_m \ll E_{Na} = +52$ mV, quindi la forza motrice è verso l'interno)
\end{itemize}


\section{Relazione Corrente-Tensione e Approssimazione Ohmica}

\subsubsection{La Curva I-V dalla Legge GHK}

Il potenziale di membrana di un \textbf{neurone in funzione} fluttua continuamente, a causa dei bombardamenti sinaptici eccitatori e inibitori. Il range operativo fisiologico si estende tipicamente da circa $-100$ mV a $+50$ mV.

La relazione corrente-tensione (curva I-V) ricavata dall'equazione GHK per il potassio è una \textbf{curva esponenziale}. Il punto cruciale è il \textbf{potenziale di inversione} $E_K$: sopra di esso la corrente è uscente (positiva per convenzione), al di sotto è entrante (negativa).

\subsubsection{Linearizzazione di Ohm}

\begin{tcolorbox}[colback=gray!5, colframe=gray!40, arc=2mm, boxrule=0.5pt]
Nell'intervallo operativo del neurone, la curva esponenziale può essere approssimata linearmente senza errore significativo. Questo è il passaggio chiave che permette di costruire il circuito equivalente.
\end{tcolorbox}

La \textbf{linearizzazione ohmica} di Ohm permette di scrivere:

\begin{equation}
I_\alpha \approx G_\alpha (V - E_\alpha)
\end{equation}

dove $G_\alpha$ è la \textbf{conduttanza} della specie ionica $\alpha$ (l'inverso della resistenza, misurata in Siemens) e $E_\alpha$ è il potenziale di inversione (che funge da \textbf{forza elettromotrice} nella notazione del circuito).

Per le due specie ioniche principali:

\begin{equation}
I_{Na} = G_{Na}(V - E_{Na})
\end{equation}

\begin{equation}
I_K = G_K(V - E_K)
\end{equation}

Queste due correnti scorrono \textbf{in parallelo} attraverso la membrana.

\subsubsection{Validità della Linearizzazione e Caso del Ca$^{2+}$}

La linearizzazione ohmica è \textbf{valida per le strutture grandi} del neurone (soma, dendriti, assone), dove le concentrazioni di Ca$^{2+}$ sono trascurabili. In queste strutture, Na$^+$, K$^+$ e Cl$^-$ obbediscono perfettamente all'approssimazione.

Esiste tuttavia un'importante \textbf{eccezione}: le \textbf{spine sinaptiche}, strutture submicrometriche che ricevono l'input sinaptico. In questi compartimenti molto piccoli:

\begin{itemize}
    \item Il volume è ridottissimo, quindi piccoli flussi di Ca$^{2+}$ producono \textbf{grandi variazioni di concentrazione}
    \item Il potenziale di inversione del Ca$^{2+}$ è $E_{Ca} \approx +150$ mV
    \item La curva I-V del Ca$^{2+}$ è quindi \textbf{fortemente non lineare} su tutto l'intervallo operativo
    \item L'approssimazione ohmica \textbf{non è applicabile} al Ca$^{2+}$ nelle spine
\end{itemize}

Per gli scopi di questa lezione — e per costruire il modello di Hodgkin-Huxley — ci si concentra su Na$^+$, K$^+$ e Cl$^-$ nelle strutture grandi, dove la linearizzazione è giustificata.


\section{Circuito Equivalente della Membrana}

\subsubsection{Struttura del Doppio Strato Lipidico}

Consideriamo una \textbf{piccola sezione di membrana}. Il doppio strato lipidico ha due proprietà elettriche fondamentali:

\begin{enumerate}
    \item Le due pareti del doppio strato sono superfici conduttrici: le cariche ioniche in soluzione acquosa si distribuiscono rapidamente sulle superfici, rendendole \textbf{equipotenziali} (approssimazione di conduttori piani).
    \item Le due pareti conduttrici separate da un isolante lipidico costituiscono un \textbf{condensatore piano}.
\end{enumerate}

Questa struttura suggerisce immediatamente un modello circuitale.

\subsubsection{Componenti del Circuito Passivo}

Attraverso una piccola sezione di membrana scorrono tre tipi di corrente:

\begin{enumerate}
    \item \textbf{Corrente del Na$^+$}: $I_{Na} = G_{Na}(V - E_{Na})$, rappresentata da una resistenza (conduttanza $G_{Na}$) in serie con una batteria (forza elettromotrice $E_{Na}$).
    \item \textbf{Corrente del K$^+$}: $I_K = G_K(V - E_K)$, analogamente rappresentata da $G_K$ in serie con $E_K$.
    \item \textbf{Corrente capacitiva}: la membrana immagazzina carica $Q = C_m \cdot V$ (dove $C_m$ è la \textbf{capacitanza di membrana}), e la corrente associata è:
\end{enumerate}

\begin{equation}
I_C = C_m \frac{dV}{dt}
\end{equation}

\subsubsection{Equazione del Circuito Chiuso}

Per il \textbf{principio di conservazione della carica} (legge di Kirchhoff), in un circuito chiuso la somma di tutte le correnti deve essere zero. Includendo anche una possibile corrente esterna $I_e$ iniettata tramite un \textbf{elettrodo intracellulare} (pipetta di vetro):

\begin{equation}
C_m \frac{dV}{dt} + I_{Na} + I_K = 0 \quad \text{(senza corrente esterna)}
\end{equation}

Con corrente esterna per unità di area $I_e/A$ (dove $A$ è la sezione della pipetta):

\begin{equation}
C_m \frac{dV}{dt} + I_{Na} + I_K = \frac{I_e}{A}
\end{equation}


\subsubsection{Derivazione dell'Equazione Completa}

Sostituendo le espressioni ohmiche per le correnti ioniche nell'equazione del circuito, si ottiene:

\begin{equation}
C_m \frac{dV}{dt} = -\sum_\alpha G_\alpha(V - E_\alpha) + \frac{I_e}{A}
\end{equation}

Espandendo la somma e raccogliendo i termini:

\begin{equation}
C_m \frac{dV}{dt} = -\left(\sum_\alpha G_\alpha\right) V + \sum_\alpha G_\alpha E_\alpha + \frac{I_e}{A}
\end{equation}

Definendo la \textbf{conduttanza totale di membrana}:

\begin{equation}
G_m = \sum_\alpha G_\alpha
\end{equation}

l'equazione diventa:

\begin{equation}
C_m \frac{dV}{dt} = -G_m V + \sum_\alpha G_\alpha E_\alpha + \frac{I_e}{A}
\end{equation}

\subsubsection{Il Potenziale di Equilibrio Effettivo}

Si può raccogliere ulteriormente definendo il \textbf{potenziale di equilibrio effettivo} $E_m$:

\begin{equation}
E_m = \frac{\sum_\alpha G_\alpha E_\alpha}{\sum_\alpha G_\alpha} = \frac{\sum_\alpha G_\alpha E_\alpha}{G_m}
\end{equation}

\begin{tcolorbox}[colback=gray!5, colframe=gray!40, arc=2mm, boxrule=0.5pt]
Questo potenziale di equilibrio $E_m$, derivato qui tramite l'approssimazione ohmica, coincide esattamente con il potenziale di Goldman-Hodgkin-Katz! Abbiamo quindi ottenuto lo stesso risultato con due approcci completamente diversi, il che è una conferma della coerenza dei due modelli.
\end{tcolorbox}

L'equazione del circuito passivo assume quindi la forma compatta:

\begin{equation}
C_m \frac{dV}{dt} = -G_m(V - E_m) + \frac{I_e}{A}
\end{equation}

\subsubsection{Il Circuito RC di Membrana}

Definendo la \textbf{costante di tempo di membrana}:

\begin{equation}
\tau_m = C_m R_m = \frac{C_m}{G_m}
\end{equation}

dove $R_m = 1/G_m$ è la \textbf{resistenza totale di membrana}, si ottiene la forma compatta adimensionale:

\begin{equation}
\tau_m \frac{dV}{dt} = E_m - V + R_m \frac{I_e}{A}
\end{equation}

Questa è identica all'\textbf{equazione di un circuito RC} standard con:
\begin{itemize}
    \item Costante di tempo $\tau_m$
    \item Valore di equilibrio $E_m + R_m I_e/A$ (dipendente dalla corrente esterna)
\end{itemize}


\subsubsection{Risposta a un Gradino di Corrente}

Consideriamo il caso in cui a $t = 0$ si applichi un gradino di corrente costante $I_e$ a partire dalla condizione di equilibrio $V(0) = E_m$. La soluzione dell'equazione RC è:

\begin{equation}
V(t) = E_m + R_m \frac{I_e}{A}\left(1 - e^{-t/\tau_m}\right)
\end{equation}

Il sistema presenta quindi:
\begin{itemize}
    \item Una \textbf{crescita esponenziale} dalla condizione iniziale $E_m$
    \item Verso il \textbf{valore asintotico}: $V_\infty = E_m + R_m I_e/A$
    \item Con costante di tempo $\tau_m$
\end{itemize}

Analogamente, alla rimozione della corrente (gradino negativo), si ha un \textbf{decadimento esponenziale} verso $E_m$ con la stessa costante di tempo $\tau_m$.

\subsubsection{Caratterizzazione Sperimentale del Circuito}

Questo semplice protocollo sperimentale (injection di correnti note) permette di \textbf{misurare completamente i parametri} del circuito RC:

\begin{table}[htbp]
    \centering
    \renewcommand{\arraystretch}{1.3}
    \begin{tabularx}{\textwidth}{|l|X|}
    \hline
    \textbf{Misura} & \textbf{Parametro ottenuto} \\
    \hline
    $V_\infty$ vs $I_e = 0$ & Potenziale di equilibrio $E_m$ \\
    Pendenza $dV_\infty / dI_e = R_m/A$ & Resistenza di membrana $R_m$ (e conduttanza $G_m$) \\
    Costante di decadimento & Costante di tempo $\tau_m$ \\
    $\tau_m / R_m$ & Capacitanza di membrana $C_m = \tau_m G_m$ \\
    \hline
    \end{tabularx}
\end{table}

\subsubsection{Validazione Sperimentale Storica}

\begin{tcolorbox}[colback=gray!5, colframe=gray!40, arc=2mm, boxrule=0.5pt]
\textbf{Nota del docente:} Mostro dati sperimentali storici risalenti a più di 50 anni fa, ottenuti sulla \textbf{giunzione neuromuscolare di rana}. Questi esperimenti furono condotti da ricercatori che iniettarono diverse correnti e osservarono esattamente quanto previsto dal circuito RC equivalente: adattamento esponenziale alla corrente, poi decadimento esponenziale verso l'equilibrio al ritiro della corrente. Questo dimostra che l'approssimazione delle proprietà elettriche passive è straordinariamente efficace.
\end{tcolorbox}

Il professore menziona come riferimento bibliografico principale "\textit{Biophysics of Electrical Properties of Neurons}" (Cambridge Press), definendolo "la bibbia in neuroscienze" per questo argomento.

\subsubsection{Limiti del Modello Passivo e Transizione al Modello HH}

\begin{tcolorbox}[colback=gray!5, colframe=gray!40, arc=2mm, boxrule=0.5pt]
\textbf{Nota del docente:} Le proprietà passive non sono però l'unica storia. Il vero potere del neurone è la capacità di \textbf{trasformare nonlinearmente l'informazione sinaptica}. Questo fu scoperto da Hodgkin e Huxley negli anni '50. Dobbiamo estendere il circuito passivo per spiegare questa nonlinearità.
\end{tcolorbox}

Il circuito RC passivo descrive correttamente la risposta subliminale del neurone. Non può però spiegare il \textbf{potenziale d'azione} — lo spike — che è il meccanismo fondamentale di comunicazione neuronale.

\newpage

\section{L'Esperimento di Hodgkin e Huxley}

\subsubsection{La Scelta dell'Assone Gigante di Calamaro}

\textbf{Alan Hodgkin} e \textbf{Andrew Huxley} scelsero come sistema modello l'\textbf{assone gigante di calamaro} (\textit{Loligo vulgaris}). La motivazione era puramente pratica: le fibre nervose giganti di questo mollusco hanno un diametro di circa \textbf{1 mm} (centinaia di volte più grande degli assoni di mammifero), il che rende possibili manipolazioni sperimentali altrimenti impossibili.

Il protocollo sperimentale prevedeva:
\begin{enumerate}
    \item Tagliare un pezzo di assone
    \item Immergerlo in un bagno con \textbf{fluido cerebrospinale artificiale} (ACSF, \textit{Artificial Cerebrospinal Fluid})
    \item Applicare la tecnica dello \textbf{space clamp} e del \textbf{voltage clamp}
\end{enumerate}

\subsubsection{Lo Space Clamp}

Il problema fondamentale nello studio di un assone è che il potenziale d'azione si \textbf{propaga} lungo la fibra, rendendo il potenziale di membrana una funzione sia del tempo che della posizione. Per semplificare drasticamente il problema, HH introdussero la tecnica dello \textbf{space clamp}:

\begin{itemize}
    \item Un \textbf{anodo esterno} che copre l'intera lunghezza del pezzo di assone
    \item Un \textbf{catodo interno} (un filo metallico inserito nel lume dell'assone)
\end{itemize}

Questa configurazione forza l'\textbf{intera membrana} ad avere lo stesso potenziale in ogni istante, eliminando la dipendenza spaziale. Il sistema si riduce così a un problema puramente temporale (equazione differenziale ordinaria, non PDE).

\subsubsection{Il Voltage Clamp}

Il \textbf{voltage clamp} (pinza di voltaggio) è la seconda tecnica fondamentale:

\begin{enumerate}
    \item Due elettrodi sono connessi a un \textbf{amplificatore di feedback}
    \item L'amplificatore confronta continuamente il potenziale misurato $V_{mis}$ con il potenziale desiderato $V_{cmd}$ (imposto dall'esperimentalista)
    \item L'amplificatore inietta la \textbf{corrente necessaria} per mantenere $V = V_{cmd}$
    \item Questa corrente viene misurata con un \textbf{amperometro}
\end{enumerate}

\begin{tcolorbox}[colback=gray!5, colframe=gray!40, arc=2mm, boxrule=0.5pt]
 Con questo setup, fissando $V = -50$ mV, tutta la membrana ha esattamente quel potenziale di membrana in ogni istante. Si può misurare come cambia nel tempo la corrente che l'amplificatore deve iniettare per mantenere il potenziale costante — e questa corrente è esattamente uguale e opposta alla corrente ionica che fluisce attraverso la membrana.
\end{tcolorbox}



Con lo space clamp, il voltage clamp e la \textbf{modellizzazione numerica} (realizzata sui primi mainframe disponibili negli anni '50), Hodgkin e Huxley riuscirono a:
\begin{enumerate}
    \item \textbf{Misurare} le correnti ioniche in funzione del voltaggio e del tempo
    \item \textbf{Isolare} le correnti di Na$^+$ e K$^+$ tramite farmaci specifici
    \item \textbf{Costruire} un modello matematico quantitativo del potenziale d'azione
    \item \textbf{Simulare} numericamente il potenziale d'azione, ottenendo un accordo eccellente con i dati
\end{enumerate}

Per questi risultati, \textbf{Alan Hodgkin} e \textbf{Andrew Huxley} ricevettero il \textbf{Premio Nobel per la Fisiologia e la Medicina nel 1963}, condiviso con \textbf{John Eccles} (che aveva scoperto indipendentemente i meccanismi della trasmissione sinaptica).


\vspace{0.3 cm}


\noindent Senza voltage clamp, se si lascia libero il potenziale, si osserva il \textbf{potenziale d'azione} come uno spike. L'obiettivo era capire \textit{perché} esiste questa amplificazione non lineare. La strategia di Hodgkin e Huxley fu:

\textbf{Fissare} il voltaggio a diversi livelli (voltage clamp) $\rightarrow$ \textbf{misurare} la corrente risultante $\rightarrow$ \textbf{caratterizzare} come la corrente dipende da $V$ e dal tempo $\rightarrow$ \textbf{costruire} un modello.

\subsection{Correnti Non Lineari: Piccolo vs Grande $\Delta$V}

\subsubsection{Esperimento con Piccolo $\Delta$V (+10 mV)}

Partendo dal potenziale di equilibrio $E_m$, si applica un salto di potenziale $\Delta V = +10$ mV. La corrente misurata presenta:

\begin{enumerate}
    \item Un \textbf{overshoot iniziale} (picco istantaneo): è la \textbf{corrente capacitiva} $I_C = C_m \, dV/dt$, proporzionale alla derivata del voltaggio (impulsiva per un gradino ideale)
    \item Poi un \textbf{rilassamento esponenziale} verso il valore asintotico $\Delta V \cdot G_m$: questa è la corrente di leakage passiva, perfettamente descritta dal circuito RC
\end{enumerate}

Questo comportamento è \textbf{esattamente} quanto previsto dal circuito passivo.

\subsubsection{Esperimento con Grande $\Delta$V (+60 mV)}

Applicando invece $\Delta V = +60$ mV (una forte depolarizzazione), il comportamento cambia radicalmente:

\begin{enumerate}
    \item Stesso \textbf{overshoot iniziale} capacitivo
    \item Poi una \textbf{forte corrente entrante} (\textit{inward current}): specie cationiche (Na$^+$) che entrano nella cellula, corrispondente alla depolarizzazione
    \item Poi una \textbf{corrente uscente} (\textit{outward current}): specie cationiche (K$^+$) che escono, corrispondente alla ripolarizzazione
\end{enumerate}

\begin{tcolorbox}[colback=gray!5, colframe=gray!40, arc=2mm, boxrule=0.5pt]
Questa non linearità nasce da qualcosa che il circuito RC passivo non è in grado di descrivere. Il circuito passivo prevede solo un rilassamento esponenziale monotono. Invece vediamo una corrente che cambia segno — prima entrante, poi uscente. Questo è il segnale che ci sono elementi attivi nella membrana.
\end{tcolorbox}


\subsection{Isolamento di $I_{Na}$ e $I_K$: TTX e TEA}


Per identificare i contributi separati delle correnti ioniche, Hodgkin e Huxley (e i loro collaboratori) ricorsero a farmaci specifici che bloccano selettivamente un tipo di canale.

\subsubsection{La Tetrodotossina (TTX)}

La \textbf{tetrodotossina} (TTX) è un potente neurotossico di natura non proteica, prodotto da batteri simbiotici e accumulato in diversi animali marini, il più noto dei quali è il \textbf{pesce palla} (famiglia \textit{Tetraodontidae}).

\begin{itemize}
    \item \textbf{Meccanismo}: blocco specifico dei \textbf{canali ionici voltaggio-dipendenti del Na$^+$} (si inserisce fisicamente nel poro del canale)
    \item \textbf{Effetto fisiologico}: paralisi muscolare (incluso il diaframma) $\rightarrow$ morte per arresto respiratorio
    \item \textbf{Potenza}: circa 1200 volte più tossica del cianuro
\end{itemize}

Con TTX nel bagno di superfusione, la corrente misurata è \textbf{esclusivamente quella del K$^+$} (corrente uscente).

\subsubsection{Il Tetraetilammonium (TEA)}

Il \textbf{tetraetilammonium} (TEA) è un bloccante specifico dei \textbf{canali del K$^+$} voltaggio-dipendenti. Con TEA nel bagno, la corrente misurata è \textbf{esclusivamente quella del Na$^+$}.
La corrente del Na$^+$ isolata è \textbf{non monotona}: presenta prima un picco inward (attivazione rapida dei canali Na$^+$), poi si esaurisce rapidamente (inattivazione).

\subsubsection{Ricostruzione della Corrente Totale}

La corrente totale osservata senza farmaci è la \textbf{somma} delle due correnti isolate:

\begin{equation}
I_{totale} = I_{Na} + I_K
\end{equation}

Questo risultato fondamentale permette di identificare le \textbf{due fasi} del potenziale d'azione:

\begin{enumerate}
    \item \textbf{Fase di salita (spike)}: i \textbf{canali Na$^+$} si aprono $\rightarrow$ corrente entrante $\rightarrow$ \textbf{depolarizzazione} rapida della membrana
    \item \textbf{Fase di recupero (ripolarizzazione)}: i \textbf{canali K$^+$} si aprono (e i canali Na$^+$ si inattivano) $\rightarrow$ corrente uscente $\rightarrow$ \textbf{ripolarizzazione} verso $E_m$
\end{enumerate}



\subsection{Il Circuito Equivalente Hodgkin-Huxley}

\subsubsection{Estensione del Circuito Passivo}

Hodgkin e Huxley riconobbero che il circuito passivo deve essere \textbf{esteso} con canali \textbf{attivi} (voltaggio-dipendenti) per Na$^+$ e K$^+$. La differenza fondamentale rispetto al circuito passivo è che le conduttanze \textbf{non sono più costanti}:

\begin{equation}
G_{Na} = G_{Na}(V, t) \qquad G_K = G_K(V, t)
\end{equation}

Queste conduttanze dipendono sia dal potenziale di membrana $V$ che dal tempo $t$ (tramite le variabili di gate che vedremo nelle prossime sezioni).

\subsubsection{L'Equazione HH Completa}

L'equazione differenziale del modello di Hodgkin-Huxley completo è:

\begin{equation}
C_m \frac{dV}{dt} = -\bar{G}_{Na} m^3 h (V - E_{Na}) - \bar{G}_K n^4 (V - E_K) - G_L (V - E_L) + I_{ext}
\end{equation}

dove:
\begin{itemize}
    \item $\bar{G}_{Na}$ e $\bar{G}_K$ sono le \textbf{conduttanze massime} (valori costanti quando tutti i canali sono aperti)
    \item $m$, $h$, $n$ sono le \textbf{variabili di gate} (comprese tra 0 e 1)
    \item $G_L$ è la \textbf{conduttanza di leakage} (passiva, costante)
    \item $E_L$ è il potenziale di inversione di leakage
    \item $I_{ext}$ è la corrente esterna iniettata
\end{itemize}

Le \textbf{batterie ioniche} hanno polarità opposte:
\begin{itemize}
    \item $E_K \approx -70 \div -90$ mV: tende a iperpolarizzare la membrana
    \item $E_{Na} \approx +50$ mV: tende a depolarizzare la membrana
\end{itemize}

\begin{tcolorbox}[colback=gray!5, colframe=gray!40, arc=2mm, boxrule=0.5pt]

Il circuito equivalente HH è topologicamente identico al circuito passivo, ma con la differenza cruciale che $G_{Na}$ e $G_K$ non sono resistori fissi bensì elementi il cui valore dipende dal voltaggio e dal tempo. Questo rende il sistema non lineare.
\end{tcolorbox}


\subsection{La Conduttanza del Potassio }


Con la TTX nel bagno (per bloccare il Na$^+$), la corrente misurata è solo $I_K$. La conduttanza del potassio in funzione del tempo si ricava come:

\begin{equation}
G_K(t) = \frac{I_K(t)}{V - E_K}
\end{equation}

(il denominatore è la forza motrice, costante grazie al voltage clamp)


I dati sperimentali a diversi valori di $\Delta V$ rivelano tre caratteristiche che \textbf{non possono essere spiegate} dal circuito passivo:

\textbf{Prima differenza — Forma sigmoidale:}\\
L'andamento temporale di $G_K(t)$ non è esponenziale ma \textbf{sigmoidale}. Nel circuito passivo, la risposta sarebbe puramente esponenziale con derivata iniziale diversa da zero. Invece sperimentalmente la derivata iniziale è \textbf{zero}: $G_K$ parte lentamente, poi accelera, poi satura.

\textbf{Seconda differenza — Valore asintotico voltaggio-dipendente:}\\
Il valore asintotico $G_K(\infty)$ dipende fortemente da $V$. Nel circuito passivo, $G_K$ sarebbe costante indipendentemente da $V$.

\textbf{Terza differenza — Costante di tempo voltaggio-dipendente:}\\
Il tempo caratteristico di avvicinamento al valore asintotico dipende da $V$. Ad esempio:
\begin{itemize}
    \item A $\Delta V = 109$ mV: $\tau \approx 1.5$ ms
    \item A $\Delta V$ più piccolo: $\tau \approx 4$ ms
\end{itemize}

Tutti e tre questi comportamenti indicano la presenza di \textbf{canali ionici attivi} con proprietà voltaggio-dipendenti.


\subsection{Il Modello di Markov per i Canali Ionici}

\subsubsection{Struttura Molecolare dei Canali: Due Stati Discreti}

Alla fine degli anni '40, era già noto che i canali ionici sono \textbf{molecole macromolecolari} con almeno due stati conformazionali discreti:
\begin{itemize}
    \item Stato \textbf{aperto} (\textit{open}): il poro è accessibile agli ioni, la corrente può fluire
    \item Stato \textbf{chiuso} (\textit{closed}): il poro è bloccato, nessuna corrente
\end{itemize}

Questa struttura può essere rappresentata con un \textbf{paesaggio di energia a doppio pozzo}: due minimi di energia corrispondenti ai due stati, separati da una \textbf{barriera energetica}.

\subsubsection{Transizioni Stocastiche e Velocità di Kramers}

Le transizioni tra stati possono essere:
\begin{itemize}
    \item \textbf{Spontanee} (stocastiche): fluttuazioni termiche possono fornire l'energia necessaria per superare la barriera
    \item \textbf{Voltaggio-indotte}: se il sensore di voltaggio del canale porta un campo elettrico locale, il potenziale di membrana modifica l'altezza della barriera
\end{itemize}

La \textbf{probabilità di transizione} è descritta dalla \textbf{velocità di Kramers}:

\begin{equation}
\text{rate} \propto e^{-\Delta E / k_B T}
\end{equation}

dove $\Delta E$ è l'altezza della barriera e $k_B T$ è l'energia termica. Se il voltaggio modifica $\Delta E$, allora le rate constants sono \textbf{voltaggio-dipendenti}.

\subsubsection{Equazione Cinetica del Primo Ordine per n}

Hodgkin e Huxley modellarono il canale del K$^+$ come avente una \textbf{variabile di gate} $n$ che rappresenta la \textbf{frazione di canali aperti} in un dato istante.

Le transizioni seguite da un processo di Markov (dipendente solo dallo stato attuale):
\begin{itemize}
    \item Chiuso $\rightarrow$ Aperto con rate $\alpha_n(V)$
    \item Aperto $\rightarrow$ Chiuso con rate $\beta_n(V)$
\end{itemize}

In un intervallo $\Delta t$, la variazione di $n$ è:

\begin{equation}
\Delta n = \alpha_n(1-n)\Delta t - \beta_n \cdot n \cdot \Delta t
\end{equation}

Nel limite $\Delta t \to 0$, si ottiene l'\textbf{equazione cinetica del primo ordine}:

\begin{equation}
\dot{n} = \alpha_n(V)(1 - n) - \beta_n(V) \, n
\end{equation}

Questa è un'equazione di Markov: la dinamica di $n$ dipende solo dal valore attuale di $n$ (non dalla storia passata).


\subsubsection{Riscrittura in Forma Canonica}

L'equazione cinetica per $n$ si può riscrivere nella forma canonica:

\begin{equation}
\dot{n} = \frac{n_\infty(V) - n}{\tau_n(V)}
\end{equation}

dove il \textbf{valore asintotico} e la \textbf{costante di tempo} sono:

\begin{equation}
n_\infty(V) = \frac{\alpha_n(V)}{\alpha_n(V) + \beta_n(V)}, \qquad \tau_n(V) = \frac{1}{\alpha_n(V) + \beta_n(V)}
\end{equation}


\subsubsection{Il Problema della Forma Sigmoidale}

La soluzione di $\dot{n} = (n_\infty - n)/\tau_n$ (con $V$ costante sotto voltage clamp) è un'evoluzione \textbf{esponenziale} da $n(0)$ verso $n_\infty$. Ma la $G_K$ sperimentale ha forma \textbf{sigmoidale}, non esponenziale!

Hodgkin e Huxley proposero una soluzione elegante: la conduttanza del K$^+$ è proporzionale a una \textbf{potenza di n}. Testando $n$, $n^2$, $n^3$, $n^4$, trovarono che:

\begin{equation}
G_K(V, t) = \bar{G}_K \cdot n^4(V, t)
\end{equation}

Il termine $n^4$ (prodotto di quattro variabili esponenziali indipendenti) genera una forma sigmoidale che \textbf{fitta perfettamente} i dati sperimentali.

\begin{tcolorbox}[colback=gray!5, colframe=gray!40, arc=2mm, boxrule=0.5pt]
\textbf{Nota del docente:} Perché n$^4$? Non c'era allora una spiegazione strutturale: era solo il fitting migliore. Decenni dopo, la cristallografia a raggi X dei canali del K$^+$ rivelò che i canali sono effettivamente \textbf{tetrameri} — composti da quattro subunità, ciascuna con una propria variabile di gate. Hodgkin e Huxley avevano indovinato la struttura molecolare solo dall'analisi funzionale!
\end{tcolorbox}



\subsection{Fitting con n$^4$ e Dipendenza da V delle Rate Constants}

Il fitting di $G_K \propto n^4$ viene eseguito per ogni valore di $\Delta V$, ottenendo le coppie $(\alpha_n(V), \beta_n(V))$. In totale si ottengono le \textbf{dipendenze funzionali} di $\alpha_n$ e $\beta_n$ dal voltaggio.


Dai dati sperimentali a diversi $\Delta V$, Hodgkin e Huxley trovarono comportamenti asintotici precisi:

\begin{itemize}
    \item \textbf{Per V molto positivo} (depolarizzazione forte): $\beta_n \to 0$ e $\alpha_n \propto V$ $\rightarrow$ $n_\infty \to 1$ (tutti i canali K$^+$ aperti)
    \item \textbf{Per V molto negativo} (iperpolarizzazione forte): $\alpha_n \to 0$ $\rightarrow$ $n_\infty \to 0$ (tutti i canali K$^+$ chiusi)
\end{itemize}

Le forme funzionali adottate da Hodgkin e Huxley per le rate constants sono combinazioni di termini lineari in $V$ ed esponenziali in $V$. Questa scelta è giustificata fisicamente dalla \textbf{velocità di Kramers}: la barriera energetica dipende linearmente dal voltaggio (attraverso la carica del sensore di voltaggio), quindi il rate di transizione va come $e^{\pm eV/k_BT}$.

Di conseguenza, $n_\infty(V) = \alpha_n(V)/[\alpha_n(V) + \beta_n(V)]$ ha la forma di una \textbf{funzione sigmoidale}:
\begin{itemize}
    \item $n_\infty \to 1$ per $V$ depolarizzato
    \item $n_\infty \to 0$ per $V$ iperpolarizzato
\end{itemize}



\subsubsection{Sintesi del Modello per il Potassio}

L'equazione completa per la conduttanza del potassio è:

\begin{equation}
G_K(V, t) = \bar{G}_K \cdot n^4(V, t)
\end{equation}

con la cinetica governata da:

\begin{equation}
\dot{n} = \alpha_n(V)(1-n) - \beta_n(V) \cdot n
\end{equation}

I comportamenti limite:
\begin{itemize}
    \item $n_\infty \to 1$ per $V$ depolarizzato $\rightarrow$ $G_K \to \bar{G}_K$ (conduttanza massima, tutti i canali aperti)
    \item $n_\infty \to 0$ per $V$ iperpolarizzato $\rightarrow$ $G_K \to 0$ (tutti i canali chiusi, nessuna corrente)
\end{itemize}

Questo meccanismo spiega la fase di \textbf{ripolarizzazione} del potenziale d'azione: la depolarizzazione provoca l'apertura dei canali K$^+$ (con un certo ritardo, per via della cinetica di $n$), che produce una corrente uscente che riporta $V$ verso $E_K$.



\subsection{La Conduttanza del Sodio}

\subsubsection{L'Andamento Non Monotono di $G_{Na}$}

Isolare $G_{Na}$ tramite TEA rivela un comportamento \textbf{qualitativamente diverso} da $G_K$: l'andamento temporale di $G_{Na}$ è \textbf{non monotono}. A grande $\Delta V$:

\begin{enumerate}
    \item \textbf{Fase di attivazione}: $G_{Na}$ cresce rapidamente (in pochi decimi di ms)
    \item \textbf{Fase di inattivazione}: $G_{Na}$ decresce in meno di 1 ms, tornando a zero anche se $V$ è ancora mantenuto a un valore depolarizzato
\end{enumerate}

Questo comportamento non può essere descritto da una sola variabile del primo ordine $\rightarrow$ servono \textbf{due variabili di gate} con dinamiche separate e comportamenti opposti.

\subsubsection{Le Due Variabili m e h}

Hodgkin e Huxley introdussero:

\begin{itemize}
    \item \textbf{m}: variabile di \textbf{attivazione} (apertura del gate)
    \item \textbf{h}: variabile di \textbf{inattivazione} (chiusura/blocco del gate)
\end{itemize}

La conduttanza del Na$^+$ è:

\begin{equation}
G_{Na}(V, t) = \bar{G}_{Na} \cdot m^3(V, t) \cdot h(V, t)
\end{equation}

\begin{tcolorbox}[colback=gray!5, colframe=gray!40, arc=2mm, boxrule=0.5pt]
\textbf{Nota del docente:} La potenza $m^3$ e il fattore $h$ sono trovati per fitting numerico dei dati sperimentali, esattamente come per $n^4$. La struttura molecolare dei canali del Na$^+$ (scoperta molto più tardi) confermò questa intuizione: i canali del Na$^+$ hanno effettivamente tre subunità di attivazione e un dominio di inattivazione separato.
\end{tcolorbox}




\subsubsection{Cinetica di m (Attivazione)}

L'equazione cinetica per la variabile di attivazione $m$ è:

\begin{equation}
\dot{m} = \alpha_m(V)(1-m) - \beta_m(V) \cdot m
\end{equation}

con il valore asintotico e la costante di tempo:

\begin{equation}
m_\infty(V) = \frac{\alpha_m(V)}{\alpha_m(V) + \beta_m(V)}, \qquad \tau_m(V) = \frac{1}{\alpha_m(V) + \beta_m(V)}
\end{equation}

Le rate constants per $m$ hanno una forma funzionale analoga a quelle di $\alpha_n$: il termine lineare domina per $V$ positivo e $\alpha_m \to 0$ per $V \to -\infty$.

Risultato: $m_\infty(V)$ è una \textbf{sigmoide crescente}:
\begin{itemize}
    \item $m_\infty \to 1$ per $V$ depolarizzato (gate di attivazione aperto)
    \item $m_\infty \to 0$ per $V$ iperpolarizzato (gate di attivazione chiuso)
\end{itemize}

La \textbf{costante di tempo} $\tau_m$ è molto piccola rispetto a $\tau_n$ e $\tau_h$: $m$ è la variabile più \textbf{rapida} del modello HH (si attiva quasi istantaneamente).

\subsubsection{Cinetica di h (Inattivazione)}

L'equazione cinetica per la variabile di inattivazione $h$ è:

\begin{equation}
\dot{h} = \alpha_h(V)(1-h) - \beta_h(V) \cdot h
\end{equation}

con:

\begin{equation}
h_\infty(V) = \frac{\alpha_h(V)}{\alpha_h(V) + \beta_h(V)}, \qquad \tau_h(V) = \frac{1}{\alpha_h(V) + \beta_h(V)}
\end{equation}



Le rate constants per $h$ hanno un \textbf{comportamento inverso} rispetto a quelle di $m$:

\begin{itemize}
    \item \textbf{Per V molto positivo} (depolarizzazione): $\alpha_h \to 0$ e $\beta_h$ è grande $\rightarrow$ $h_\infty = \alpha_h/(\alpha_h + \beta_h) \to 0$
    \item \textbf{Per V molto negativo} (iperpolarizzazione): $\beta_h \to 0$ e $\alpha_h$ è grande $\rightarrow$ $h_\infty \to 1$
\end{itemize}

Quindi $h_\infty(V)$ è una \textbf{sigmoide decrescente}:
\begin{itemize}
    \item $h_\infty \to 1$ per $V$ iperpolarizzato: i canali Na$^+$ sono "pronti" ad attivarsi (gate di inattivazione aperto)
    \item $h_\infty \to 0$ per $V$ depolarizzato: i canali Na$^+$ sono inattivati (gate di inattivazione chiuso)
\end{itemize}

\begin{tcolorbox}[colback=gray!5, colframe=gray!40, arc=2mm, boxrule=0.5pt]
\textbf{Nota del docente:} L'interpretazione biologica è precisa: h è la variabile che dice "il canale è disponibile per aprirsi". Quando la membrana è a riposo (iperpolarizzata), h è vicino a 1 — i canali sono "carichi". Dopo un potenziale d'azione (quando V è rimasto a lungo depolarizzato), h crolla verso 0 — i canali sono inattivati e non possono riaprirsi. Questo è il meccanismo del \textbf{periodo refrattario assoluto}.
\end{tcolorbox}

\subsubsection{Generazione del Comportamento Non Monotono di $G_{Na}$}

La combinazione di $m^3$ (sigmoide crescente, rapida) e $h$ (sigmoide decrescente, lenta) crea il comportamento non monotono osservato:

\begin{enumerate}
    \item \textbf{Inizio} (subito dopo depolarizzazione): $m$ cresce rapidamente, $h$ è ancora vicino al suo valore di riposo (alto) $\rightarrow$ $G_{Na} = \bar{G}_{Na} m^3 h$ cresce
    \item \textbf{Picco}: $m$ ha quasi raggiunto $m_\infty$, $h$ sta iniziando a diminuire $\rightarrow$ $G_{Na}$ al massimo
    \item \textbf{Inattivazione}: $m$ è saturo, ma $h$ diminuisce rapidamente $\rightarrow$ $G_{Na}$ crolla
\end{enumerate}



\section*{Riepilogo del Modello di Hodgkin-Huxley}


Il modello di Hodgkin-Huxley è un sistema di \textbf{quattro equazioni differenziali ordinarie} accoppiate:

\textbf{Equazione di membrana:}
\begin{equation}
C_m \frac{dV}{dt} = -\bar{G}_{Na} m^3 h (V - E_{Na}) - \bar{G}_K n^4 (V - E_K) - G_L (V - E_L) + I_{ext}
\end{equation}

\textbf{Cinetica di m (attivazione Na$^+$):}
\begin{equation}
\dot{m} = \alpha_m(V)(1-m) - \beta_m(V) \, m
\end{equation}

\textbf{Cinetica di h (inattivazione Na$^+$):}
\begin{equation}
\dot{h} = \alpha_h(V)(1-h) - \beta_h(V) \, h
\end{equation}

\textbf{Cinetica di n (attivazione K$^+$):}
\begin{equation}
\dot{n} = \alpha_n(V)(1-n) - \beta_n(V) \, n
\end{equation}

\subsubsection{Gerarchia Temporale delle Variabili}

Le tre variabili di gate operano su scale temporali diverse, il che è essenziale per il meccanismo del potenziale d'azione:

\begin{table}[htbp]
    \centering
    \renewcommand{\arraystretch}{1.3}
    \begin{tabularx}{\textwidth}{|l|X|X|X|}
    \hline
    \textbf{Variabile} & \textbf{Tipo} & \textbf{Velocità} & \textbf{Comportamento a V depolarizzato} \\
    \hline
    $m$ & Attivazione Na$^+$ & Rapida ($\tau_m \sim 0.1$ ms) & $m_\infty \to 1$ \\
    $h$ & Inattivazione Na$^+$ & Lenta ($\tau_h \sim 1{-}10$ ms) & $h_\infty \to 0$ \\
    $n$ & Attivazione K$^+$ & Intermedia ($\tau_n \sim 1{-}4$ ms) & $n_\infty \to 1$ \\
    \hline
    \end{tabularx}
\end{table}

\subsubsection{Meccanismo del Potenziale d'Azione (Anteprima)}

Il potenziale d'azione emerge dall'interazione delle variabili come un'\textbf{instabilità autoamplificata}:

\begin{enumerate}
    \item Una perturbazione depolarizzante supera la soglia $\rightarrow$ $m$ si attiva rapidamente
    \item Corrente Na$^+$ entrante $\rightarrow$ ulteriore depolarizzazione (\textbf{feedback positivo})
    \item La depolarizzazione attiva $n$ (K$^+$) più lentamente e inattiva $h$ $\rightarrow$ inizio della ripolarizzazione
    \item $G_K$ cresce, $G_{Na}$ crolla $\rightarrow$ ripolarizzazione rapida (\textbf{feedback negativo})
    \item Iperpolarizzazione transitoria (afterhyperpolarization) $\rightarrow$ periodo refrattario
\end{enumerate}


\newpage
\section*{Tabella Riepilogativa}

\begin{table}[htbp]
    \centering
    \footnotesize
    \renewcommand{\arraystretch}{1.3}
    \begin{tabularx}{\textwidth}{|>{\raggedright\arraybackslash}p{3cm}|>{\raggedright\arraybackslash}X|>{\raggedright\arraybackslash}X|}
    \hline
    \textbf{Simbolo / Termine} & \textbf{Definizione} & \textbf{Unità / Note} \\
    \hline
    $E_{Na}$ & Potenziale di inversione del Na$^+$ & $\approx +52$ mV (assone di calamaro) \\
    \hline
    $E_K$ & Potenziale di inversione del K$^+$ & $\approx -76$ mV (assone di calamaro) \\
    \hline
    $E_{Cl}$ & Potenziale di inversione del Cl$^-$ & $\approx -60$ mV (assone di calamaro) \\
    \hline
    $E_m$ & Potenziale di equilibrio di membrana & $\approx -60$ mV \\
    \hline
    $G_\alpha$ & Conduttanza ionica per la specie $\alpha$ & S/cm$^2$ \\
    \hline
    $G_m$ & Conduttanza totale di membrana & $G_m = \sum_\alpha G_\alpha$ \\
    \hline
    $R_m$ & Resistenza di membrana & $\Omega \cdot$cm$^2$; $R_m = 1/G_m$ \\
    \hline
    $C_m$ & Capacitanza di membrana & F/cm$^2$; tipico $\approx 1 \mu$F/cm$^2$ \\
    \hline
    $\tau_m$ & Costante di tempo di membrana & ms; $\tau_m = C_m/G_m = C_m R_m$ \\
    \hline
    $I_e/A$ & Corrente esterna iniettata per unità di area & A/cm$^2$ \\
    \hline
    $\bar{G}_{Na}$ & Conduttanza massima del Na$^+$ & S/cm$^2$; quando tutti i canali Na$^+$ sono aperti \\
    \hline
    $\bar{G}_K$ & Conduttanza massima del K$^+$ & S/cm$^2$; quando tutti i canali K$^+$ sono aperti \\
    \hline
    $G_L$ & Conduttanza di leakage (passiva) & S/cm$^2$; costante \\
    \hline
    $m$ & Variabile di gate di attivazione del Na$^+$ & $\in [0,1]$; sigmoide crescente in V \\
    \hline
    $h$ & Variabile di gate di inattivazione del Na$^+$ & $\in [0,1]$; sigmoide decrescente in V \\
    \hline
    $n$ & Variabile di gate di attivazione del K$^+$ & $\in [0,1]$; sigmoide crescente in V; potenza 4 \\
    \hline
    $\alpha_x(V)$ & Rate di transizione chiuso$\rightarrow$aperto per la variabile $x$ & ms$^{-1}$; voltaggio-dipendente \\
    \hline
    $\beta_x(V)$ & Rate di transizione aperto$\rightarrow$chiuso per la variabile $x$ & ms$^{-1}$; voltaggio-dipendente \\
    \hline
    $x_\infty(V)$ & Valore asintotico della variabile di gate $x$ a voltaggio $V$ & Adimensionale \\
    \hline
    $\tau_x(V)$ & Costante di tempo della variabile di gate $x$ a voltaggio $V$ & ms \\
    \hline
    Modello di Markov & Transizioni tra stati con rate costanti & Dipende solo dallo stato attuale \\
    \hline
    Velocità di Kramers & Rate di transizione termica su barriera energetica & $\propto e^{-\Delta E/k_BT}$ \\
    \hline
   
    \end{tabularx}
\end{table}

\chapter{Lezione 5}


\textbf{Argomenti:} variabili di gate n, m, h, cinetica del primo ordine, funzioni di transizione $\alpha$ e $\beta$, valori all'equilibrio n$_\infty$/m$_\infty$/h$_\infty$, costanti di tempo $\tau_n$/$\tau_m$/$\tau_h$, simulazione Mathematica, potenziale di riposo HH, correnti ioniche Na$^+$/K$^+$/leakage, sistema a quattro dimensioni, potenziale d'azione come fenomeno stereotipato, soglia di eccitabilità, fasi dello spike (rising/decaying/iperpolarizzazione), periodo refrattario assoluto e relativo, inattivazione del Na$^+$, temperatura e velocità dello spike, propagazione del potenziale d'azione, corrente assiale, unidirezionalità, velocità di propagazione, mielinizzazione, conduzione saltatoria, nodi di Ranvier


\subsubsection{Ricapitolazione dalla Lezione Precedente}

L'assone gigante del calamaro (\textit{Loligo vulgaris}) presenta proprietà di \textbf{permeabilità ionica} al Na$^+$ e al K$^+$ che dipendono dal potenziale di membrana. Grazie agli esperimenti di Hodgkin e Huxley, fu possibile isolare il decorso temporale delle correnti del sodio e del potassio separatamente, rivelando la relazione tra conduttanze e potenziale di membrana.

Per descrivere questo fenomeno, Hodgkin e Huxley ipotizzarono l'esistenza di \textbf{variabili di gate} (gating variables), ovvero variabili interne ai canali ionici che ne regolano l'apertura e la chiusura. Esse sono:

\begin{itemize}
    \item \textbf{$n$}: variabile di gate per il canale del K$^+$ (variabile di attivazione)
    \item \textbf{$m$}: variabile di gate di attivazione per il canale del Na$^+$
    \item \textbf{$h$}: variabile di gate di inattivazione per il canale del Na$^+$
\end{itemize}

Ogni variabile varia tra 0 (canale completamente chiuso) e 1 (canale completamente aperto).



Ciascuna variabile di gate obbedisce a un'\textbf{equazione differenziale del primo ordine}:

\begin{equation}
\tau_x(V) \frac{dx}{dt} = -x + x_\infty(V)
\end{equation}

dove $x \in \{n, m, h\}$, $\tau_x(V)$ è la \textbf{costante di tempo} di rilassamento  e $x_\infty(V)$ è il \textbf{valore asintotico}. Questa equazione descrive il processo di apertura e chiusura dei canali ionici come un sistema a \textbf{due stati} (aperto/chiuso), con:

\begin{itemize}
    \item $\alpha_x(V)$: tasso di transizione verso lo stato aperto (opening rate)
    \item $\beta_x(V)$: tasso di transizione verso lo stato chiuso (closing rate)
\end{itemize}

Le relazioni tra queste quantità e i parametri della cinetica del primo ordine sono:

\begin{equation}
x_\infty(V) = \frac{\alpha_x(V)}{\alpha_x(V) + \beta_x(V)}
\end{equation}

\begin{equation}
\tau_x(V) = \frac{1}{\alpha_x(V) + \beta_x(V)}
\end{equation}

\begin{tcolorbox}[colback=gray!5, colframe=gray!40, arc=2mm, boxrule=0.5pt]
Il rapporto $\alpha_x / (\alpha_x + \beta_x)$ fornisce la frazione di gate aperti all'equilibrio termodinamico.
\end{tcolorbox}


Le funzioni $\alpha$ e $\beta$ per ciascuna variabile di gate hanno un'espressione analitica derivata empiricamente da Hodgkin e Huxley fittando i dati sperimentali di voltage clamp. La struttura generale prevede un \textbf{termine lineare} al numeratore e una \textbf{dipendenza esponenziale} al denominatore:

\begin{equation}
\alpha_n(V) = \frac{0.01(V + 55)}{1 - e^{-(V+55)/10}}
\quad \quad  ,  \quad \quad
\beta_n(V) = 0.125 \, e^{-(V+65)/80}
\end{equation}

\begin{equation}
\alpha_m(V) = \frac{0.1(V + 40)}{1 - e^{-(V+40)/10}}
\quad \quad  ,  \quad \quad
\beta_m(V) = 4 \, e^{-(V+65)/18}
\end{equation}

\begin{equation}
\alpha_h(V) = 0.07 \, e^{-(V+65)/20}
\quad \quad  ,  \quad \quad
\beta_h(V) = \frac{1}{1 + e^{-(V+35)/10}}
\end{equation}


Le funzioni di transizione per le variabili $n$ ed $m$ mostrano un \textbf{comportamento qualitativo simile}: $\alpha$ cresce con il voltaggio (la probabilità di apertura aumenta con la depolarizzazione), mentre $\beta$ decresce. Per la variabile $h$, il comportamento è \textbf{invertito}: quando il voltaggio aumenta, $\alpha_h$ diminuisce e $\beta_h$ aumenta. Questo garantisce che l'inattivazione del sodio proceda con la depolarizzazione.


Le curve $n_\infty(V)$ e $m_\infty(V)$ hanno un andamento \textbf{sigmoidale crescente}: aumentano da 0 a voltaggio molto negativo a 1 a voltaggio molto positivo. La curva $h_\infty(V)$ è invece una \textbf{sigmoide decrescente}: vale circa 1 a valori iperpolarizzati e si riduce verso 0 con la depolarizzazione.

Al potenziale di riposo $V = -65 \; \text{mV}$, i  valori di equilibrio sono:

\begin{table}[htbp]
    \centering
    \renewcommand{\arraystretch}{1.3}
    \begin{tabular}{|c|c|}
    \hline
    \textbf{Variabile} & \textbf{Valore all'equilibrio} \\
    \hline
    $V$ & $-65 \; \text{mV}$ \\
    $n_\infty(-65)$ & $\approx 0.32$ \\
    $m_\infty(-65)$ & $\approx 0.05$ (molto vicino a 0) \\
    $h_\infty(-65)$ & $\approx 0.60$ \\
    \hline
    \end{tabular}
\end{table}

\begin{tcolorbox}[colback=gray!5, colframe=gray!40, arc=2mm, boxrule=0.5pt]
Il fatto che $m_\infty$ sia quasi zero a riposo significa che la conduttanza del sodio è praticamente nulla in assenza di stimolazione: $G_{Na} = \bar{G}_{Na} \cdot m^3 h \approx 0$. D'altra parte $h_\infty \approx 0.6$ — non è né 0 né 1 — e questo implica che una piccola porzione dei gate di inattivazione è già chiusa anche a riposo. Il sistema è quindi "pronto" ad attivarsi, ma non completamente libero dall'inattivazione basale.
\end{tcolorbox}

\subsubsection{Dipendenza dal Voltaggio delle Costanti di Tempo}

Le costanti di tempo di rilassamento delle tre variabili di gate hanno una forma a \textbf{campana} (bell-shaped) in funzione del voltaggio, con un massimo nella regione del potenziale di riposo. A differenza del valore asintotico $x_\infty(V)$, le costanti di tempo non variano monotonicamente: hanno un massimo intorno alla tensione di riposo e si riducono sia a valori molto positivi che a valori molto negativi.

Il risultato più importante è la \textbf{gerarchia delle scale temporali} delle tre variabili:

\begin{equation}
\tau_m(V) \ll \tau_n(V) \approx \tau_h(V)
\end{equation}

Concretamente, al potenziale di riposo e nelle sue vicinanze:
\begin{itemize}
    \item $\tau_m$ è dell'ordine di \textbf{frazioni di millisecondo} (scala rapida)
    \item $\tau_n$ e $\tau_h$ sono dell'ordine di \textbf{qualche millisecondo} (scala lenta)
\end{itemize}



\begin{tcolorbox}[colback=gray!5, colframe=gray!40, arc=2mm, boxrule=0.5pt]
Questa separazione di scala temporale è la chiave di tutta la dinamica del potenziale d'azione. Quando il potenziale cambia bruscamente, $m$ si adatta quasi istantaneamente al nuovo valore di $m_\infty$, mentre $n$ e $h$ rimangono quasi "congelati" al loro valore precedente per diversi millisecondi. Questa differenza di velocità è ciò che genera la cascata non lineare che porta all'esplosione del potenziale d'azione.
\end{tcolorbox}


\section{Correnti Ioniche nel Modello HH}

\subsubsection{Equazione del Modello di Hodgkin-Huxley}

Il cuore del modello è l'equazione per il potenziale di membrana $V(t)$, che bilancia la corrente capacitiva con le correnti ioniche che attraversano la membrana:

\begin{equation}
C_m \frac{dV}{dt} = -G_L(V - E_L) - \bar{G}_{Na} m^3 h (V - E_{Na}) - \bar{G}_K n^4 (V - E_K) + I_{ext}
\end{equation}

Questo sistema è completato dalle tre equazioni differenziali del primo ordine per le variabili di gate:

\begin{equation}
\tau_m(V) \frac{dm}{dt} = m_\infty(V) - m
\end{equation}

\begin{equation}
\tau_n(V) \frac{dn}{dt} = n_\infty(V) - n
\end{equation}

\begin{equation}
\tau_h(V) \frac{dh}{dt} = h_\infty(V) - h
\end{equation}

Il sistema totale è dunque di \textbf{dimensione 4} nello spazio delle fasi: le quattro variabili di stato sono $V$, $n$, $m$, $h$.



\begin{tcolorbox}[colback=gray!5, colframe=gray!40, arc=2mm, boxrule=0.5pt]
I due potenziali di inversione delimitano l'\textbf{intervallo fisiologico operativo} della membrana: $E_K = -77 \; \text{mV}$ rappresenta il limite inferiore (non si può andare più negativi perché la corrente di K$^+$ scompare), mentre $E_{Na} = +50 \; \text{mV}$ rappresenta il limite superiore (sopra tale valore la corrente del Na$^+$ cambia segno e non spinge più la depolarizzazione).
\end{tcolorbox}

\subsubsection{Potenziale di Membrana a Riposo nel Modello HH}


Lo stato di riposo del modello HH corrisponde alla soluzione stazionaria del sistema a quattro dimensioni, ovvero al punto per cui $\dot{V} = 0$, $\dot{m} = 0$, $\dot{n} = 0$, $\dot{h} = 0$. Imponendo queste condizioni, il sistema di quattro equazioni nonlineari da risolvere simultaneamente è:

\begin{equation}
G_L(V - E_L) + \bar{G}_{Na} m^3 h (V - E_{Na}) + \bar{G}_K n^4 (V - E_K) = 0
\end{equation}
\begin{equation}
m = m_\infty(V)
\end{equation}
\begin{equation}
n = n_\infty(V)
\end{equation}
\begin{equation}
h = h_\infty(V)
\end{equation}


Il risultato della ricerca numerica dà:

\begin{equation}
V_{riposo} = -65 \; \text{mV}
\end{equation}

Questo valore coincide esattamente con quello misurato sperimentalmente da Hodgkin e Huxley sull'assone gigante del calamaro. Il professore sottolinea la \textbf{coerenza} tra il modello matematico e l'esperimento: fittando le funzioni $\alpha$ e $\beta$ sui dati di voltage clamp, si riproduce automaticamente il corretto potenziale di riposo.

\begin{tcolorbox}[colback=gray!5, colframe=gray!40, arc=2mm, boxrule=0.5pt]
Il fatto che $m_\infty \approx 0$ a riposo significa che la conduttanza del sodio è quasi nulla. Il potenziale di riposo è quindi determinato prevalentemente dall'equilibrio tra la corrente di leakage e la piccola ma non nulla conduttanza del potassio. Questo è fisicamente corretto: il K$^+$ è lo ione più permeabile a riposo.
\end{tcolorbox}


\section{Iniezione di Corrente a Gradino}

\subsubsection{Setup dell'Esperimento in Silico}

Per esplorare la risposta del modello a una stimolazione, il professore descrive un esperimento di \textbf{iniezione di corrente a gradino} che replica ciò che si fa in laboratorio con una micropipetta. Il protocollo è il seguente:

\begin{enumerate}
    \item Il sistema parte all'equilibrio: $V(0) = -65 \; \text{mV}$
    \item Al tempo $t = 0$, si inietta un impulso di corrente molto breve e intenso che produce una \textbf{depolarizzazione istantanea} di $\Delta V$
    \item Dopo l'impulso, il sistema evolve liberamente 
\end{enumerate}

Questo è fondamentalmente diverso dall'esperimento di voltage clamp originale di Hodgkin e Huxley: qui il potenziale di membrana è libero di variare, e si osserva l'\textbf{evoluzione spontanea} del sistema.

\subsubsection{Risposta a Perturbazioni di Diversa Ampiezza}

Il professore mostra sistematicamente i risultati per tre ampiezze di perturbazione:

\textbf{Perturbazione di 90 mV} ($V(0^+) \approx +25 \; \text{mV}$): Si osserva una risposta rapida e decisa. Il potenziale raggiunge un massimo vicino a $+40 \; \text{mV}$ (leggermente al di sotto di $E_{Na} = +50 \; \text{mV}$), poi cala bruscamente attraversando la linea di equilibrio e raggiungendo una fase di iperpolarizzazione sotto $-65 \; \text{mV}$, per poi tornare lentamente al potenziale di riposo.

\textbf{Perturbazione di 50 mV} ($V(0^+) \approx -15 \; \text{mV}$): Si osserva un comportamento qualitativamente identico. C'è una breve rampa nella fase iniziale, poi una rapida depolarizzazione verso lo stesso valore massimo ($\approx +40 \; \text{mV}$), seguita dalla stessa iperpolarizzazione. Il tempo complessivo e il picco sono quasi indistinguibili dal caso precedente.

\textbf{Perturbazione di 7 mV} ($V(0^+) \approx -58 \; \text{mV}$): Anche in questo caso, dopo un periodo di rampa più lento, si osserva lo stesso potenziale d'azione a tutti gli effetti: stessa ampiezza, stesso picco, stessa iperpolarizzazione.

\begin{tcolorbox}[colback=gray!5, colframe=gray!40, arc=2mm, boxrule=0.5pt]
Questa è la proprietà più importante del potenziale d'azione: è un \textbf{evento stereotipato}. Indipendentemente dall'intensità della perturbazione iniziale (purché superi la soglia), la forma dell'evento è sempre identica. Il neurone non codifica l'informazione nell'ampiezza dello spike, ma nella sua semplice presenza o assenza. Il messaggio trasmesso al neurone postsinaptico è sempre lo stesso.
\end{tcolorbox}

\subsubsection{Il Concetto di Soglia}

Riducendo ulteriormente l'ampiezza della perturbazione, si arriva a un punto critico: per una depolarizzazione di circa \textbf{6 mV} ($V(0^+) \approx -59 \; \text{mV}$), il sistema non produce più alcun potenziale d'azione. Invece di salire verso il picco, il potenziale di membrana torna lentamente all'equilibrio senza alcuna esplosione.

Il confronto tra la traccia a 7 mV (che produce uno spike) e quella a 6 mV (che non lo produce) evidenzia un cambiamento \textbf{qualitativo} del comportamento. Non si tratta di una variazione quantitativa dello stesso fenomeno, ma di un salto di categoria: due traiettorie inizialmente molto simili divergono in modo drammatico.

Questo è il comportamento caratteristico di un \textbf{sistema eccitabile}: esiste un valore soglia della variabile di stato (il potenziale di membrana) al di sopra del quale si innesca una traiettoria stereotipata nello spazio delle fasi (lo spike), e al di sotto del quale il sistema torna semplicemente all'equilibrio.

\subsubsection{Soglia Sperimentale}

Il professore indica che, in base alle simulazioni, la \textbf{soglia di eccitazione} si trova circa a $-58.5 \; \text{mV}$, ovvero una depolarizzazione di circa $6$–$7 \; \text{mV}$ rispetto al potenziale di riposo di $-65 \; \text{mV}$.

\begin{tcolorbox}[colback=gray!5, colframe=gray!40, arc=2mm, boxrule=0.5pt]
 La soglia non è una proprietà fissa del neurone, ma dipende dallo stato delle variabili di gate nel momento della perturbazione. Come vedremo quando discuteremo il periodo refrattario, la stessa perturbazione di 7 mV può non produrre alcuno spike se applicata immediatamente dopo un potenziale d'azione.
\end{tcolorbox}

\section{Dinamica Temporale dello Spike}

\subsubsection{Le Tre Fasi del Potenziale d'Azione}

Un potenziale d'azione prodotto dal modello HH si articola in tre fasi temporalmente distinte, ciascuna governata da meccanismi biofisici specifici:

\begin{enumerate}
    \item \textbf{Fase di salita (rising phase)}: rapido aumento del potenziale da circa $-65 \; \text{mV}$ fino al picco ($\approx +40 \; \text{mV}$), che si completa in circa $1.2 \; \text{ms}$
    \item \textbf{Fase di discesa (decaying phase)}: calo del potenziale dal picco attraverso il valore di riposo, guidato da una corrente netta negativa
    \item \textbf{Fase di iperpolarizzazione}: il potenziale scende al di sotto di $E_m = -65 \; \text{mV}$ e poi torna lentamente al riposo, con una scala temporale di diversi millisecondi
\end{enumerate}

\subsubsection{Fase di Salita: il Ruolo di m e la Reazione a Catena}

All'inizio della perturbazione, il sistema parte con $m \approx 0$ (a riposo). Tuttavia, il salto di voltaggio sposta immediatamente $m_\infty$ a un valore più alto. Poiché $\tau_m$ è molto piccolo, $m$ si adatta quasi istantaneamente al nuovo $m_\infty$.

Si instaura quindi una \textbf{reazione a catena positiva}:

\begin{equation}
V \uparrow \Rightarrow m_\infty \uparrow \Rightarrow m \uparrow \Rightarrow G_{Na} = \bar{G}_{Na} m^3 h \uparrow \Rightarrow I_{Na,in} \uparrow \Rightarrow V \uparrow
\end{equation}

In questa fase, $n$ e $h$ rimangono quasi costanti ai loro valori di riposo ($n \approx 0.32$, $h \approx 0.60$) perché $\tau_n$ e $\tau_h$ sono grandi. La corrente del K$^+$ esiste ma è ancora piccola perché $n$ non ha avuto tempo di cambiare significativamente.

La salita si arresta quando il potenziale si avvicina a $E_{Na} = +50 \; \text{mV}$: a quel punto il fattore $(V - E_{Na})$ tende a zero e la corrente del Na$^+$ scompare non per mancanza di conduttanza (che è anzi massima in quel momento), ma per mancanza di forza motrice.

\begin{tcolorbox}[colback=gray!5, colframe=gray!40, arc=2mm, boxrule=0.5pt]
È fondamentale distinguere tra \textbf{conduttanza} e \textbf{corrente}. Al picco dello spike, la conduttanza del Na$^+$ è al suo valore massimo, ma la corrente del Na$^+$ è quasi zero perché $(V - E_{Na}) \approx 0$. Confondere i due concetti porta a un'errata comprensione del meccanismo di inversione.
\end{tcolorbox}

\subsubsection{Fase di Discesa: il Ruolo di h e n}

Durante la fase di salita, mentre $m$ si adattava rapidamente, $n$ e $h$ stavano evolvendo lentamente verso i loro nuovi valori asintotici che ora corrispondere a voltaggio positivo. A voltaggio positivo:

\begin{itemize}
    \item $n_\infty \approx 1$, quindi $n$ cresce lentamente verso 1 $\rightarrow$ la corrente del K$^+$ \textbf{aumenta} (corrente iperpolarizzante)
    \item $h_\infty \approx 0$, quindi $h$ decresce lentamente verso 0 $\rightarrow$ la conduttanza del Na$^+$ \textbf{diminuisce}
\end{itemize}

La \textbf{fase di discesa} è dunque il risultato combinato di:
\begin{enumerate}
    \item Aumento della conduttanza del K$^+$ (via aumento di $n$)
    \item Diminuzione della conduttanza del Na$^+$ (via diminuzione di $h$, l'inattivazione)
\end{enumerate}

La discesa è più lenta della salita perché la corrente netta negativa (K$^+$ $-$ Na$^+$) rimane relativamente piccola.

\subsubsection{Fase di Iperpolarizzazione: il Ruolo di n}

Nella fase di iperpolarizzazione, $h$ ha raggiunto quasi zero (la conduttanza del Na$^+$ è pressoché nulla) e anche $m$ torna rapidamente verso zero perché $m_\infty \to 0$ a valori negativi. La corrente dominante è quindi esclusivamente quella del K$^+$, che continua a scorrere verso l'esterno perché $n$ è ancora elevato.

Poiché $E_K = -77 \; \text{mV}$, la corrente del K$^+$ spinge il potenziale verso $-77 \; \text{mV}$, \textbf{al di sotto} del potenziale di riposo. La membrana è quindi \textbf{iperpolarizzata} rispetto a $E_m$.

Il ritorno al potenziale di riposo avviene lentamente, man mano che $n$ decresce verso il suo valore di equilibrio $n_\infty(-65) \approx 0.32$ con la costante di tempo $\tau_n$.

\begin{equation}
\underbrace{\text{Fase di salita}}_{\text{dominata da } m} \longrightarrow \underbrace{\text{Fase di discesa}}_{\text{dominata da } m, h, n} \longrightarrow \underbrace{\text{Iperpolarizzazione}}_{\text{dominata da } n}
\end{equation}


\section{Analisi delle Correnti e delle Conduttanze durante lo Spike}

\subsubsection{Evoluzione delle Conduttanze}

Il professore mostra l'andamento temporale delle conduttanze $G_{Na}(t) $ e $G_K(t)$ durante un potenziale d'azione:

\begin{itemize}
    \item $G_{Na}$ ha un picco \textbf{precoce} e molto stretto: sale rapidamente durante la fase di salita, raggiunge il massimo appena prima del picco del potenziale, poi cala rapidamente a causa dell'inattivazione (diminuzione di $h$)
    \item $G_K$ ha un picco \textbf{tardivo} e più largo: cresce lentamente durante la fase di salita, raggiunge il massimo nella fase di discesa, poi decresce lentamente nella fase di iperpolarizzazione
\end{itemize}

La corrente di leakage ha un ruolo del tutto marginale in questo fenomeno altamente nonlineare.

\subsubsection{Corrente Netta e Segno della Derivata del Voltaggio}

La derivata del potenziale di membrana è proporzionale alla corrente netta totale:

\begin{equation}
C_m \dot{V} = I_{Na} + I_K + I_L
\end{equation}

(segni positivi per le correnti entranti, negativi per le uscenti)

\begin{itemize}
    \item \textbf{Fase di salita}: $I_{Na}$ è entrante e grande, $I_K$ è ancora piccola $\rightarrow$ corrente netta positiva $\rightarrow$ $\dot{V} > 0$
    \item \textbf{Apice dello spike}: $I_{Na} \approx 0$ (perché $V \to E_{Na}$), $I_K$ aumenta $\rightarrow$ il bilancio si annulla $\rightarrow$ $\dot{V} = 0$
    \item \textbf{Fase di discesa}: $I_K$ supera $I_{Na}$ $\rightarrow$ corrente netta negativa $\rightarrow$ $\dot{V} < 0$
    \item \textbf{Iperpolarizzazione}: solo $I_K$ significativa, $I_{Na} \approx 0$ $\rightarrow$ $\dot{V} < 0$ fino a $E_K$
\end{itemize}



\section{Periodo Refrattario}

Nella fase di iperpolarizzazione che segue il potenziale d'azione, il sistema è in uno stato molto diverso dall'equilibrio. Il professore dimostra con una simulazione che, se si inietta la stessa perturbazione di 50 mV a $t = 5 \; \text{ms}$ (durante la fase di iperpolarizzazione), \textbf{non si produce alcun secondo potenziale d'azione}, nonostante la perturbazione sia ampiamente sopra-soglia nelle condizioni normali.

Questo periodo di insensibilità agli stimoli è definito \textbf{periodo refrattario}. Durante questo intervallo:
\begin{itemize}
    \item Il neurone è completamente insensibile ai segnali sinaptici in ingresso
    \item Non è un "difetto" del sistema, ma una proprietà funzionale che limita la frequenza massima di firing
\end{itemize}

\begin{tcolorbox}[colback=gray!5, colframe=gray!40, arc=2mm, boxrule=0.5pt] Il periodo refrattario ha un'implicazione profonda per la codifica dell'informazione neurale. Non solo esiste un \textbf{limite massimo alla frequenza di firing} del neurone, ma anche un limite teorico alla velocità con cui il neurone può trasmettere informazioni al neurone postsinaptico. Il neurone non è un trasmettitore ideale: ha una "finestra cieca" di diversi millisecondi dopo ogni spike.
\end{tcolorbox}


La causa del periodo refrattario risiede nella \textbf{stato delle variabili di gate} durante l'iperpolarizzazione:

\begin{enumerate}
    \item \textbf{$m \approx 0$}: il gate di attivazione del Na$^+$ è quasi completamente chiuso perché $m_\infty \to 0$ a potenziali negativi. In assenza di $m$, la conduttanza del Na$^+$ è nulla indipendentemente dal valore di $h$. \textbf{Nessuna reazione a catena del Na$^+$ è possibile}.
    \item \textbf{$h \approx 0$}: il gate di inattivazione del Na$^+$ è quasi completamente chiuso. Anche se il potenziale venisse depolarizzato, i canali del Na$^+$ rimarrebbero inattivati finché $h$ non si recupera.
    \item \textbf{$n$ elevato}: la variabile di gate del K$^+$ è ancora alta, producendo una corrente del K$^+$ iperpolarizzante che contrasta qualsiasi tentativo di depolarizzazione.
\end{enumerate}

Il risultato è che l'unica corrente disponibile è quella del K$^+$, che è \textbf{iperpolarizzante}. Iniettare corrente in questo stato non fa altro che rendere la corrente di K$^+$ ancora più negativa (aumentando $n^4 (V - E_K)$), senza alcuna possibilità di innescare la reazione a catena del Na$^+$.


\subsubsection{Periodo Refrattario Assoluto}

Il \textbf{periodo refrattario assoluto} è l'intervallo di tempo immediatamente successivo a un potenziale d'azione durante il quale \textbf{nessuna} stimolazione, per quanto intensa, è in grado di evocare un secondo spike. Corrisponde alla fase di discesa e alla prima parte dell'iperpolarizzazione.

Il professore dimostra questo con una perturbazione di \textbf{90 mV} (ben al di sopra della soglia normale di circa 7 mV) applicata durante la fase di discesa: anche in questo caso la membrana non genera alcun potenziale d'azione. La corrente netta rimane negativa nonostante la forte depolarizzazione applicata.

Il meccanismo è lo stesso: $m \approx 0$ e i canali del Na$^+$ sono inattivati ($h \approx 0$). Il sistema non è in grado di avviare la reazione a catena del sodio.

\subsubsection{Periodo Refrattario Relativo}

Il \textbf{periodo refrattario relativo} è la fase che segue quella assoluta, durante la quale il neurone può rispondere a stimoli eccezionalmente intensi. In questo periodo, $h$ si sta \textbf{recuperando} verso il suo valore di equilibrio $h_\infty \approx 0.6$, e $n$ sta lentamente diminuendo verso $n_\infty \approx 0.32$.

Il professore mostra che se si applica la stessa perturbazione forte in questa fase più tardiva, si osserva una risposta: il potenziale d'azione viene prodotto, ma con un \textbf{profilo leggermente diverso} dall'ottimale (picco più basso, dinamica alterata), a causa dei valori di $h$ e $n$ ancora anomali rispetto al riposo.

La transizione tra periodo refrattario assoluto e relativo è determinata dal recupero di $h$:

\begin{equation}
h(t) \approx h_\infty - (h_\infty - h_{min}) \cdot e^{-t/\tau_h}
\end{equation}

dove $h_{min}$ è il valore minimo raggiunto durante lo spike e $\tau_h$ è la costante di tempo di recupero.

\begin{tcolorbox}[colback=gray!5, colframe=gray!40, arc=2mm, boxrule=0.5pt]
La distinzione pratica è semplice: durante il periodo assoluto, anche una stimolazione di 90 mV è inefficace. Durante il periodo relativo, una stimolazione normale (7 mV) è ancora insufficiente, ma una stimolazione molto forte (90 mV) può produrre uno spike, seppur distorto.
\end{tcolorbox}



\section{Firing Rate e Codifica della Frequenza}


Poiché il neurone recupera il suo stato normale dopo il periodo refrattario, è possibile elicitare \textbf{treni di potenziali d'azione} iniettando una corrente continua. Il professore discute come il modello HH risponda a una corrente di iniezione costante $I_{ext}$:

\begin{itemize}
    \item Se $I_{ext}$ è sotto-soglia: il potenziale di membrana raggiunge un nuovo stato stazionario leggermente depolarizzato senza produrre spike
    \item Se $I_{ext}$ è sopra-soglia: il neurone genera una serie ripetuta di spike con una certa \textbf{frequenza di firing} (firing rate)
\end{itemize}

\subsubsection{Relazione Frequenza-Corrente (f-I)}

La frequenza di firing $f$ dipende dall'intensità della corrente iniettata $I_{ext}$. Nella forma più semplice:
\begin{itemize}
    \item Sotto la soglia $I_{thr}$: $f = 0$
    \item Sopra la soglia: $f$ aumenta con $I_{ext}$
\end{itemize}

\noindent Il periodo refrattario impone un \textbf{limite massimo alla frequenza} di firing: se il periodo refrattario dura circa $T_{ref}$ millisecondi, la frequenza massima teorica è $f_{max} = 1/T_{ref}$.

\begin{tcolorbox}[colback=gray!5, colframe=gray!40, arc=2mm, boxrule=0.5pt]
 Questo è il fondamento della \textbf{codifica in frequenza} (rate coding): il neurone rappresenta l'intensità di uno stimolo non nell'ampiezza del singolo spike (che è sempre uguale), ma nel \textbf{numero di spike per secondo}. Una sinapsi eccitatore forte produce una frequenza di firing più alta di una sinapsi debole.
\end{tcolorbox}


\section{Effetto della Temperatura sulla Cinetica del Modello HH}


 Le funzioni $\alpha_x$ e $\beta_x$ hanno una dipendenza dalla temperatura descrivibile attraverso l'\textbf{equazione di Arrhenius}:

\begin{equation}
\alpha_x(V, T) \propto e^{-\Delta E / (k_B T)}
\end{equation}

dove $\Delta E$ è l'energia di attivazione per la transizione conformazionale del gate e $k_B T$ è l'energia termica. Aumentare la temperatura abbassa la barriera energetica relativa e aumenta i tassi di transizione.


Poiché $\tau_x = 1/(\alpha_x + \beta_x)$, un aumento di $\alpha_x$ e $\beta_x$ porta a:

\begin{equation}
\tau_x(T_2) < \tau_x(T_1) \quad \text{per } T_2 > T_1
\end{equation}

Le variabili di gate si adattano quindi più velocemente, e l'intero potenziale d'azione si comprime nella scala temporale.

\begin{table}[htbp]
    \centering
    \renewcommand{\arraystretch}{1.3}
    \begin{tabular}{|c|c|}
    \hline
    \textbf{Temperatura} & \textbf{Durata del potenziale d'azione} \\
    \hline
    $6.3 \; ^\circ\text{C}$ (condizioni originali HH) & $\approx 2 \; \text{ms}$ \\
    \hline
    $20 \; ^\circ\text{C}$ & notevolmente più breve \\
    \hline
    $37 \; ^\circ\text{C}$ (mammiferi) & $\approx 1 \; \text{ms}$ \\
    \hline
    \end{tabular}
\end{table}

\section{Propagazione del Potenziale d'Azione: Corrente Assiale}

Nell'esperimento originale di Hodgkin e Huxley, il potenziale era \textbf{uniforme} lungo tutta la lunghezza dell'assone grazie all'elettrodo interno che imponeva il \textbf{space clamp}. In un assone reale, non esiste questa costrizione: il potenziale varia in funzione della posizione $x$ lungo l'assone, e questo genera una \textbf{corrente assiale} $I_{ax}$.


Se il potenziale di membrana varia spazialmente, i gradienti di voltaggio interni generano una corrente assiale che scorre \textbf{longitudinalmente} nel citoplasma. Questa corrente è analoga a quella che scorre in un cavo elettrico:

\begin{equation}
I_{ax} = \frac{1}{r_a} \frac{\partial V}{\partial x}
\end{equation}

dove $r_a$ è la resistenza assiale per unità di lunghezza (resistenza del citoplasma, inversamente proporzionale alla sezione dell'assone).

L'equazione del modello HH si arricchisce di un termine spaziale:

\begin{equation}
C_m \frac{\partial V}{\partial t} = \frac{1}{r_a} \frac{\partial^2 V}{\partial x^2} - I_{Na} - I_K - I_L
\end{equation}



Il professore descrive il meccanismo fisico della propagazione con la seguente immagine: quando un potenziale d'azione è in corso in una posizione $x_0$, le cariche positive (ioni Na$^+$ entrati) si accumulano nella faccia interna della membrana. Queste cariche creano un campo elettrico \textbf{longitudinale} che spinge altre cariche positive verso le regioni adiacenti in avanti ($x > x_0$), \textbf{depolarizzando} così il segmento successivo.

Se la depolarizzazione indotta nel segmento successivo è superiore alla soglia, questo segmento a sua volta genera un potenziale d'azione, che a sua volta depolarizza il segmento ancora più avanti — e così via. Il potenziale d'azione si propaga come un'\textbf{onda solitaria} (soliton) lungo l'assone, con forma invariante nel tempo e nello spazio.

\begin{equation}
I_{ax} = C_m \dot{V} - (I_{Na} + I_K + I_L)
\end{equation}

Il professore mostra uno snapshot del profilo spaziale del potenziale lungo l'assone: la forma è speculare rispetto al profilo temporale, perché la parte "davanti" nella direzione di propagazione (dove lo spike non è ancora arrivato) corrisponde temporalmente al "futuro".



Una delle proprietà più importanti della propagazione è la sua \textbf{unidirezionalità}: il potenziale d'azione si propaga sempre in una sola direzione, dal soma verso le terminazioni assonali (direzione \textbf{ortodromica}), senza tornare indietro.

La spiegazione di questa proprietà è direttamente connessa al periodo refrattario assoluto:

\begin{enumerate}
    \item Un'onda di depolarizzazione si propaga in avanti verso il segmento \textbf{a riposo} ($x > x_0$): questo segmento è eccitabile e risponde generando un potenziale d'azione
    \item La stessa onda genera anche una corrente assiale verso il segmento \textbf{già eccitato} ($x < x_0$): ma questo segmento è nel suo periodo refrattario assoluto ($m \approx 0$, $h \approx 0$, $n$ grande)
    \item Nel segmento refrattario, la corrente assiale produce solo una piccola depolarizzazione che viene annullata dalla grande corrente uscente del K$^+$ — nessun secondo potenziale d'azione è possibile
\end{enumerate}

In sintesi: il segmento posteriore è "esausto" e impossibilitato a rigenerarsi, mentre il segmento anteriore è "fresco" e pienamente eccitabile.

\begin{tcolorbox}[colback=gray!5, colframe=gray!40, arc=2mm, boxrule=0.5pt]
L'unidirezionalità ha un significato funzionale fondamentale per la comunicazione neurale. Se i potenziali d'azione potessero tornare indietro, si genererebbero cicli di rientro che renderebbero impossibile qualsiasi codifica ordinata dell'informazione. La refrattarietà assoluta è quindi un meccanismo che garantisce l'\textbf{integrità direzionale del segnale nervoso}.
\end{tcolorbox}



\subsubsection{Velocità di Propagazione}

La velocità di propagazione del potenziale d'azione dipende dalla rapidità con cui la corrente assiale riesce a depolarizzare il segmento seguente. Questa dipende da due fattori:

\begin{enumerate}
    \item \textbf{L'ampiezza della corrente assiale}: maggiore è la sezione trasversale dell'assone, maggiore è la quantità di citoplasma e quindi maggiore è la corrente che può scorrere
    \item \textbf{La resistenza assiale}: un assone più largo ha \textbf{resistenza assiale più bassa} $r_a \propto 1/A$ (dove $A$ è l'area della sezione), quindi la stessa differenza di voltaggio genera una corrente molto più grande
\end{enumerate}

Il risultato complessivo è che la \textbf{velocità di propagazione} cresce con il \textbf{raggio} dell'assone:

\begin{equation}
v_{prop} \propto \sqrt{r_{assone}}
\end{equation}

Questo spiega perché l'assone gigante del calamaro (diametro $\approx 1 \; \text{mm}$) propaghi i potenziali d'azione molto più velocemente di un assone di neurone mammaliano (diametro $\approx 1 \; \mu\text{m}$).

\begin{tcolorbox}[colback=gray!5, colframe=gray!40, arc=2mm, boxrule=0.5pt]
 Nell'assone gigante del calamaro, la grande sezione serve proprio a ridurre la resistenza assiale e aumentare la velocità di propagazione. Questa è una soluzione evolutiva "ingenua": aumentare il diametro dell'assone fino a raggiungere velocità di conduzione adeguate. I mammiferi hanno adottato una soluzione molto più efficiente: la mielinizzazione.
\end{tcolorbox}



\begin{table}[htbp]
    \centering
    \renewcommand{\arraystretch}{1.3}
    \begin{tabularx}{\textwidth}{|>{\raggedright\arraybackslash}X|c|c|}
    \hline
    \textbf{Tipo di assone} & \textbf{Diametro} & \textbf{Velocità di conduzione} \\
    \hline
    Assone gigante del calamaro (non mielinizzato) & $\approx 1 \; \text{mm}$ & $\approx 25 \; \text{m/s}$ \\
    \hline
    Assone mammaliano non mielinizzato & $\approx 1 \; \mu\text{m}$ & $< 1 \; \text{m/s}$ \\
    \hline
    Assone mammaliano mielinizzato (fibra A$\alpha$) & $\approx 15 \; \mu\text{m}$ & $70$–$120 \; \text{m/s}$ \\
    \hline
    \end{tabularx}
\end{table}


\subsubsection{La Mielina Come Isolante Elettrico}

I neuroni dei mammiferi hanno risolto il problema della lentezza della propagazione in assoni sottili attraverso un meccanismo evolutivo elegante: la \textbf{mielinizzazione}. La \textbf{guaina mielinica} è formata da strati concentrici di membrane lipidiche di cellule di Schwann (nel sistema nervoso periferico) o oligodendrociti (nel sistema nervoso centrale) che avvolgono l'assone.

La mielina agisce come un \textbf{isolante elettrico} che:
\begin{enumerate}
    \item Aumenta enormemente la resistenza transmembrana nei tratti mielinizzati: la corrente non può uscire lateralmente
    \item Riduce la capacità di membrana nei tratti mielinizzati: $C_m$ ridotto significa che servono meno cariche per depolarizzare ogni segmento
\end{enumerate}


La guaina mielinica non è continua: è interrotta a intervalli regolari da piccole finestre di membrana nuda chiamate \textbf{nodi di Ranvier}. Solo in questi nodi:
\begin{itemize}
    \item La membrana è direttamente esposta al liquido extracellulare
    \item Sono presenti canali ionici voltaggio-dipendenti (Na$^+$ e K$^+$) in alta densità
    \item Può avvenire la rigenerazione attiva del potenziale d'azione
\end{itemize}



Il meccanismo risultante è la \textbf{conduzione saltatoria} (\textit{saltatory conduction}). Il potenziale d'azione non si propaga con un fronte continuo, ma "salta" da un nodo di Ranvier al successivo:

\begin{enumerate}
    \item Il potenziale d'azione viene generato nel nodo $n$
    \item La corrente assiale si propaga attraverso il tratto mielinizzato (quasi senza dispersione, grazie all'isolamento) fino al nodo $n+1$
    \item La depolarizzazione al nodo $n+1$ supera la soglia e genera un nuovo potenziale d'azione
    \item Il processo si ripete al nodo $n+2$, e così via
\end{enumerate}

In questo modo la velocità effettiva di propagazione è \textbf{molto superiore} a quella che sarebbe possibile con un'assone non mielinizzato dello stesso diametro. La conduzione saltatoria consente ai neuroni mammaliani di raggiungere velocità di $70$–$120 \; \text{m/s}$ con assoni di soli $15 \; \mu\text{m}$ di diametro, performance altrimenti raggiungibili solo con un assone di $1 \; \text{mm}$ come quello del calamaro.

\chapter{Lezione 6}

\textbf{Argomenti:} Risposta del modello HH a correnti costanti, treni di spike, curva frequenza-corrente (F-I), neuroni di Tipo I e Tipo II, adattamento della frequenza di scarica (SFA), canali HVA calcio-dipendenti, dinamica intracellulare del Ca$^{2+}$, canali K$^+$ Ca$^{2+}$-dipendenti, trasmissione sinaptica chimica, esocitosi, EPSP/IPSP, stochasticità sinaptica, analisi quantale, modello cinetico della conduttanza sinaptica, filtro sinaptico passa-basso

\section{Risposta del Modello HH a Correnti Costanti}

Il professore riprende dal modello di \textbf{Hodgkin-Huxley (HH)}, che descrive l'eccitabilità della membrana neuronale attraverso le variabili di gate $m$, $h$ ed $n$ per i canali Na$^+$ voltaggio-dipendenti (attivazione e inattivazione) e K$^+$ voltaggio-dipendenti (attivazione ritardata). L'obiettivo di questa sezione è comprendere come il neurone risponde a una \textbf{corrente costante a gradino} (step current), che modella in modo semplificato il \textbf{bombardamento sinaptico} continuo che il neurone riceve in vivo dai neuroni della rete.

La corrente iniettata $I_e$ (corrente esterna) simula la somma di tutti gli input sinaptici eccitatori e inibitori che arrivano sul dendrite e sul soma. In laboratorio, questa procedura è realizzata sperimentalmente con l'\textbf{iniezione di corrente intracellulare} attraverso un elettrodo.

\subsubsection{L'Equazione Sub-soglia del Potenziale di Membrana}

Quando la corrente iniettata $I_e$ è sufficientemente piccola, il neurone non supera la soglia di produzione di potenziali d'azione. In questo \textbf{regime lineare sub-soglia}, i canali Na$^+$ e K$^+$ voltaggio-dipendenti non si aprono significativamente: la dinamica del potenziale di membrana $V$ è dominata dal solo canale di perdita (leak) con conduttanza $G_L$ e potenziale di inversione $E_L \approx E_m = -65$ mV (nel modello HH standard).

L'equazione che governa questo regime è:

\begin{equation}
C_m \frac{dV}{dt} = -G_L(V - E_L) + \frac{I_e}{A}
\end{equation}

dove:
\begin{itemize}
    \item $C_m$ è la capacità di membrana per unità di area $[\mu\text{F/cm}^2]$,
    \item $G_L$ è la conduttanza di perdita per unità di area $[\text{mS/cm}^2]$,
    \item $E_L$ è il potenziale di inversione del leak $[\text{mV}]$,
    \item $\frac{I_e}{A}$ è la densità di corrente iniettata $[\mu\text{A/cm}^2]$ (corrente totale divisa per l'area della membrana).
\end{itemize}

\subsubsection{Deviazione Stazionaria dal Potenziale di Riposo}

In condizioni stazionarie ($dV/dt = 0$), il potenziale di membrana raggiunge un nuovo valore di equilibrio $V_{ss}$ spostato rispetto al valore di riposo $E_m = -65$ mV. La \textbf{deviazione stazionaria} $\Delta V$ vale:

\begin{equation}
\Delta V = V_{ss} - E_m = \frac{I_e / A}{G_L}
\end{equation}

Questa è semplicemente la legge di Ohm applicata al circuito RC della membrana. Il neurone si comporta come un filtro passa-basso del primo ordine con costante di tempo $\tau_m = C_m / G_L$ (la \textbf{costante di tempo di membrana}).

\subsubsection{Oscillazioni Smorzate nell'Avvicinamento all'Equilibrio}

Un fenomeno caratteristico e non banale del modello HH è che, durante la transizione verso il nuovo potenziale stazionario, il $V_m$ \textbf{non si avvicina monotonicamente} all'equilibrio, bensì presenta \textbf{oscillazioni smorzate} (damped oscillations). Il sistema HH è, vicino al punto di equilibrio, un oscillatore sottoammortizzato: il potenziale di membrana oscilla attorno a $V_{ss}$ con ampiezza decrescente esponenzialmente prima di stabilizzarsi.

\begin{tcolorbox}[colback=gray!5, colframe=gray!40, arc=2mm, boxrule=0.5pt]
Queste oscillazioni smorzate sono un primo indizio della \textbf{struttura di biforcazione} del sistema HH: il punto fisso è uno spirale stabile (stable focus) e il sistema si trova vicino a una biforcazione di Hopf. Questo è connesso alla classificazione dei neuroni in Tipo I e Tipo II che vedremo tra poco.
\end{tcolorbox}



\section{Treni di Spike}

\subsubsection{Soglia e Produzione di un Singolo Spike}

Quando la corrente iniettata viene raddoppiata a $I_e \approx 4 \, \mu\text{A/cm}^2$, il sistema supera la \textbf{soglia di eccitabilità}: si produce un singolo \textbf{potenziale d'azione} (spike). Il meccanismo è il seguente:

\begin{enumerate}
    \item La corrente capacitiva iniziale carica rapidamente la membrana (fase di ramping di $V_m$).
    \item Quando $V_m$ supera la soglia $V_{th} \approx -58.5$ mV (nel modello HH), i canali Na$^+$ voltaggio-dipendenti si aprono con retroazione positiva (overshoot).
    \item Il potenziale raggiunge il \textbf{potenziale di picco} (peak) vicino al potenziale di inversione del sodio: $E_{Na} \approx +50$ mV.
    \item Segue la \textbf{fase di ripolarizzazione} (apertura dei canali K$^+$ ritardati e inattivazione dei canali Na$^+$) e l'\textbf{iperpolarizzazione} (la cosiddetta \textbf{undershoot} o \textbf{periodo refrattario}).
    \item Il potenziale ritorna lentamente al valore di riposo $E_m$.
\end{enumerate}


\subsubsection{Dal Singolo Spike al Treno di Spike}

A corrente $I_e = 8 \, \mu\text{A/cm}^2$: si produce un \textbf{treno di spike} (spike train). La corrente iniettata è ora sufficiente a mantenere il potenziale di membrana in un regime che supera \textbf{ripetutamente e periodicamente} la soglia, generando una sequenza di potenziali d'azione.

Il meccanismo è il seguente: dopo ogni spike, durante il periodo refrattario il $V_m$ scende sotto soglia (iperpolarizzazione); poi la corrente iniettata fa ricrescere lentamente il $V_m$ (ramping) fino a raggiungere nuovamente la soglia, producendo un nuovo spike. Il processo si ripete indefinitamente a corrente costante.

A $I_e = 16 \, \mu\text{A/cm}^2$: il treno è ancora più denso, perché $dV/dt \propto I_e$ è più elevato $\rightarrow$ il ramping è più veloce $\rightarrow$ la soglia viene raggiunta in minor tempo $\rightarrow$ periodo tra spike più breve.


\subsection{Frequenza di Scarica e Inter-Spike Interval (ISI)}

Durante un treno di spike periodico, il parametro più importante è l'\textbf{Inter-Spike Interval (ISI)}, cioè il tempo tra due spike consecutivi:

\begin{equation}
T_{ISI} = t_{f+1} - t_f
\end{equation}

dove $t_f$ è il tempo del $f$-esimo spike e $t_{f+1}$ è il tempo del successivo.

La \textbf{frequenza di scarica} (firing rate) $\nu$ è l'inverso dell'ISI:

\begin{equation}
\nu = \frac{1}{T_{ISI}}
\end{equation}

Si misura in Hz (scariche al secondo). Per un treno perfettamente periodico, $\nu$ è costante; in presenza di adattamento o rumore, $\nu$ varia nel tempo.



L'ISI dipende direttamente dall'intensità della corrente iniettata:
\begin{itemize}
    \item \textbf{Alta corrente} $\Rightarrow$ ramping più rapido $\Rightarrow$ $T_{ISI}$ più corto $\Rightarrow$ $\nu$ più alta.
    \item \textbf{Bassa corrente (ma sopra soglia)} $\Rightarrow$ ramping lento $\Rightarrow$ $T_{ISI}$ lungo $\Rightarrow$ $\nu$ bassa.
\end{itemize}

La relazione tra $I_e$ e $\nu$ definisce la \textbf{curva F-I} o curva guadagno del neurone.


\subsection{La Curva F-I del Modello HH}

La \textbf{curva frequenza-corrente} (F-I curve, anche chiamata curva guadagno o curva di trasferimento) mappa la corrente iniettata $I_e$ nella frequenza di scarica stazionaria $\nu$:

\begin{equation}
\nu = f(I_e)
\end{equation}

Per il modello HH, la curva F-I ha le seguenti proprietà:
\begin{itemize}
    \item Per $I_e$ piccola (sotto soglia): $\nu = 0$ (nessun firing).
    \item Esiste una \textbf{corrente di reobase} $I_{rh}$ (rheobase current): la corrente minima necessaria per produrre spike. Sotto $I_{rh}$, il sistema HH non oscilla.
    \item A $I_e = I_{rh}$, il firing rate \textbf{salta discontinuamente} a un valore $\nu_{min} > 0$: non si può avere firing a frequenza arbitrariamente bassa.
\end{itemize}

Questo comportamento è la firma dei \textbf{neuroni di Tipo II}.

\subsubsection{Neuroni di Tipo II (Discontinuous F-I Curve)}

Un \textbf{neurone di Tipo II} è caratterizzato da una curva F-I con salto discontinuo alla reobase. Il modello HH è l'esempio prototipico di neurone di Tipo II.

La caratteristica matematica sottostante è che il punto fisso del sistema HH perde stabilità attraverso una \textbf{biforcazione di Hopf subcritica}, nella quale il ciclo limite (oscillazione periodica) appare a frequenza finita. Questo spiega il salto: il neurone passa direttamente dalla quiescenza a una frequenza minima $\nu_{min}$.

\begin{tcolorbox}[colback=gray!5, colframe=gray!40, arc=2mm, boxrule=0.5pt]
La biforcazione di Hopf subcritica implica che il sistema HH è \textbf{bistabile} in un piccolo intervallo di corrente attorno alla reobase: coesistono il punto fisso stabile e il ciclo limite stabile. Questo ha conseguenze importanti per la risposta del neurone a stimoli transitori.
\end{tcolorbox}

\subsubsection{Neuroni di Tipo I (Continuous F-I Curve)}

I \textbf{neuroni di Tipo I} hanno invece una curva F-I \textbf{continua}: la frequenza di firing parte da zero alla reobase e cresce in modo continuo (anche se la derivata $d\nu/dI_e$ può essere discontinua o molto elevata alla soglia). Un neurone di Tipo I può perciò scaricare a frequenze arbitrariamente basse, basta avere una corrente di poco superiore alla reobase.

Il meccanismo sottostante è una \textbf{biforcazione SNIC (Saddle-Node on Invariant Circle)}: il punto fisso collide con un punto di sella su un ciclo invariante. Il ciclo limite nasce con frequenza zero e cresce continuamente.

Esempio biologico: i neuroni della corteccia di tipo RS (Regular Spiking) si comportano approssimativamente come neuroni di Tipo I.

\subsubsection{Classificazione come Prima Tassonomia dei Neuroni}

Questa distinzione Tipo I / Tipo II rappresenta la \textbf{prima classificazione dei neuroni} a seconda del loro \textbf{pattern di eccitabilità}. Fu introdotta da \textbf{Rinzel e Ermentrout} negli anni '80-'90, sulla base dell'analisi delle biforcazioni dei sistemi dinamici neurali. Ha implicazioni importanti per la codifica dell'informazione: un neurone di Tipo I può codificare segnali deboli attraverso la modulazione fine della frequenza; un neurone di Tipo II risponde in modo più tutto-o-nulla.


\section{Dinamica del Calcio}


Negli anni '30, \textbf{Sir Edgar Douglas Adrian} osservò sperimentalmente che i neuroni sensoriali non scaricano a frequenza costante in risposta a uno stimolo costante: la frequenza di scarica \textbf{decresce progressivamente nel tempo}, anche a corrente iniettata costante. Questo fenomeno è chiamato \textbf{adattamento della frequenza di scarica} (Spike Frequency Adaptation, SFA).

Adrian ricevette il Premio Nobel per la Fisiologia o Medicina nel 1932 per le sue scoperte sulla funzione dei neuroni, tra cui questo fenomeno.

Il modello HH standard (con soli canali Na$^+$ e K$^+$ voltaggio-dipendenti) \textbf{non predice} l'SFA: a corrente costante, il modello HH produce un treno di spike a frequenza perfettamente costante. Per spiegare l'SFA è necessario introdurre canali ionici aggiuntivi.

\subsubsection{Canali HVA (High-Voltage Activated) del Calcio}

La chiave dell'SFA è l'esistenza di \textbf{canali ionici Ca$^{2+}$ voltaggio-dipendenti} di tipo \textbf{HVA} (High-Voltage Activated). Questi canali:

\begin{itemize}
    \item Si aprono quando il potenziale di membrana è \textbf{molto elevato}, tipicamente durante la fase ascendente di un potenziale d'azione ($V_m > -20$ mV circa).
    \item Rimangono sostanzialmente chiusi al potenziale di riposo.
    \item Permettono un rapido \textbf{influsso di Ca$^{2+}$} nella cellula durante ogni spike.
\end{itemize}

La concentrazione di calcio extracellulare $[\text{Ca}^{2+}]_o \approx 2$ mM è molto maggiore di quella intracellulare (a riposo: $[\text{Ca}^{2+}]_i \approx 100$ nM $= 10^{-7}$ M), quindi il gradiente elettrochimico favorisce fortemente l'ingresso del calcio.

\begin{tcolorbox}[colback=gray!5, colframe=gray!40, arc=2mm, boxrule=0.5pt]
 I canali HVA del calcio esistono in diversi sottotipi (L, N, P/Q, R), con caratteristiche di cinetica e localizzazione subcellulare diverse. Ai fini del modello di SFA che trattiamo, interessa solo che si aprano durante gli spike e permettano l'ingresso di Ca$^{2+}$.
\end{tcolorbox}


\subsubsection{Equazione Differenziale per la Concentrazione di Calcio}

La concentrazione di calcio intracellulare $[\text{Ca}^{2+}]_i$ (abbreviata $[\text{Ca}]_i$ nelle equazioni) varia dinamicamente durante il treno di spike. La sua evoluzione temporale è descritta da un'equazione del primo ordine:

\begin{equation}
\frac{d[\text{Ca}]_i}{dt} = -\frac{[\text{Ca}]_i}{\tau_{Ca}} + \alpha_{Ca} \sum_{f} \delta(t - t_f)
\end{equation}


\begin{itemize}
    \item $\tau_{Ca} \approx 150$ ms è la \textbf{costante di tempo di rilassamento del Ca$^{2+}$} verso il valore di equilibrio (bassissimo);
    \item $\alpha_{Ca}$ è l'\textbf{ampiezza del salto} di $[\text{Ca}]_i$ per ogni spike;
    \item $t_f$ sono i tempi dei potenziali d'azione;
    \item $\delta(t - t_f)$ è la delta di Dirac, che modella l'apertura istantanea dei canali HVA durante lo spike.
\end{itemize}

Il termine $-[\text{Ca}]_i / \tau_{Ca}$ descrive il rilassamento esponenziale verso la concentrazione basale (effettivamente $\approx 0$ nel modello semplificato), dovuto ai \textbf{trasportatori del calcio} (pompa Ca$^{2+}$-ATPasi, scambiatore Na$^+$/Ca$^{2+}$, tamponamento intracellulare).

\subsubsection{Accumulo di Ca$^{2+}$ durante il Treno di Spike}

Durante un treno di spike con periodo $T_{ISI}$, ogni spike aggiunge un salto $\alpha_{Ca}$ a $[\text{Ca}]_i$. Se il decadimento nell'ISI non è completo (cioè se $T_{ISI} \ll \tau_{Ca}$), la concentrazione di calcio si \textbf{accumula progressivamente}:

\begin{itemize}
    \item Dopo il primo spike: $[\text{Ca}]_i \approx \alpha_{Ca}$.
    \item Dopo il secondo spike: $[\text{Ca}]_i \approx 2\alpha_{Ca} - \alpha_{Ca}(1 - e^{-T_{ISI}/\tau_{Ca}}) = \alpha_{Ca}(1 + e^{-T_{ISI}/\tau_{Ca}})$.
    \item Asintoticamente, $[\text{Ca}]_i$ si stabilizza su un valore medio $\langle[\text{Ca}]_i\rangle = \alpha_{Ca} \tau_{Ca} / T_{ISI} \cdot (1 - e^{-T_{ISI}/\tau_{Ca}})^{-1}$.
\end{itemize}

Poiché $\tau_{Ca} \approx 150$ ms è molto maggiore di un tipico ISI (ad es. 20–50 ms), il Ca$^{2+}$ si accumula efficacemente spike dopo spike.

\subsubsection{Canali K$^+$ Ca$^{2+}$-Dipendenti e la Variabile di Gate $n_\infty$}

Il calcio accumulato attiva una \textbf{corrente di K$^+$ aggiuntiva} attraverso \textbf{canali K$^+$ Ca$^{2+}$-dipendenti} (SKCa: Small-conductance Calcium-activated Potassium channels). Questi canali sono \textbf{diversi} dai canali K$^+$ voltaggio-dipendenti del modello HH standard: si aprono in risposta alla concentrazione intracellulare di Ca$^{2+}$, non al potenziale di membrana.

La variabile di gate $n$ per questi canali ha un'asintoto sigmoidale:

\begin{equation}
n_\infty([\text{Ca}]) = \frac{[\text{Ca}]}{[\text{Ca}] + K_2}
\end{equation}

dove $K_2$ è una costante di dissociazione:
\begin{itemize}
    \item Quando $[\text{Ca}]_i \approx 0$ (a riposo): $n_\infty \approx 0$ $\rightarrow$ canali chiusi.
    \item Quando $[\text{Ca}]_i \gg K_2$: $n_\infty \to 1$ $\rightarrow$ canali completamente aperti.
\end{itemize}

La costante di tempo di questa variabile di gate è $\tau_n \approx 10$ ms, che è molto \textbf{più rapida} di $\tau_{Ca} \approx 150$ ms. Questo significa che $n$ insegue quasi istantaneamente $[\text{Ca}]_i$: si può applicare l'\textbf{approssimazione adiabatica} $n \approx n_\infty([\text{Ca}]_i)$.

\subsubsection{Corrente K$^+$ Ca$^{2+}$-Dipendente}

La corrente $K^+$ Ca$^{2+}$-dipendente è:

\begin{equation}
I_{K,Ca} = G_{K,Ca} \cdot n_\infty([\text{Ca}]_i) \cdot (V - E_K)
\end{equation}

dove $E_K \approx -76$ mV (potenziale di inversione del K$^+$). Poiché durante il normale firing $V > E_K$, questa corrente è \textbf{uscente} (positiva per convenzione elettrofisiologica) e quindi \textbf{iperpolarizzante}: tende ad allontanare il potenziale di membrana dalla soglia.


\subsubsection{Come la Corrente K$^+$ Ca$^{2+}$-Dipendente Allunga l'ISI}

Il meccanismo dell'\textbf{adattamento della frequenza di scarica} è ora chiaro: la corrente $I_{K,Ca}$ (negativa = iperpolarizzante) si \textbf{somma algebricamente} alla corrente iniettata $I_e$, riducendo la corrente netta disponibile per caricare la membrana:

\begin{equation}
I_{netto}(t) = \frac{I_e}{A} - G_{K,Ca} \cdot n_\infty([\text{Ca}]_i(t)) \cdot (V - E_K)
\end{equation}

Poiché $[\text{Ca}]_i$ aumenta nel corso del treno, il termine $I_{K,Ca}$ cresce progressivamente nel tempo. Questo riduce $I_{netto}$, che a sua volta riduce la velocità di ramping del potenziale di membrana tra due spike ($dV/dt = I_{netto}/C_m$).

La conseguenza diretta è che \textbf{l'ISI si allunga progressivamente}: il potenziale ci mette più tempo a raggiungere la soglia, quindi la frequenza di scarica $\nu = 1/T_{ISI}$ \textbf{diminuisce nel tempo}.

\subsubsection{Andamento Temporale dell'Adattamento}

L'evoluzione dell'SFA ha due fasi temporali caratteristiche:

\begin{enumerate}
    \item \textbf{Fase rapida} (scala di $\tau_n \approx 10$ ms): ogni spike produce immediatamente un piccolo incremento di $I_{K,Ca}$, che allunga leggermente l'ISI successivo.
    \item \textbf{Fase lenta} (scala di $\tau_{Ca} \approx 150$ ms): il Ca$^{2+}$ si accumula lentamente, aumentando gradualmente il livello basale di $I_{K,Ca}$, con un ulteriore allungamento progressivo degli ISI.
\end{enumerate}

L'SFA porta il sistema verso un nuovo stato stazionario in cui l'ISI si è stabilizzato su un valore più lungo rispetto all'ISI iniziale. Se la corrente iniettata viene rimossa, $[\text{Ca}]_i$ decade con $\tau_{Ca} \approx 150$ ms e il sistema torna alla condizione di riposo.


\subsubsection{La Frequenza Massima Assoluta}

Esiste un limite fisico superiore alla frequenza di firing, determinato dal \textbf{periodo refrattario assoluto} $\tau_{ref}$. Durante il periodo refrattario, i canali Na$^+$ sono completamente inattivati (variabile $h \approx 0$) e il neurone è assolutamente incapace di produrre un nuovo spike, indipendentemente dall'entità della corrente iniettata.

La \textbf{frequenza massima di scarica} è quindi:

\begin{equation}
\nu_{max} = \frac{1}{\tau_{ref}} \approx \frac{1}{2 \times 10^{-3} \, \text{s}} = 500 \, \text{Hz}
\end{equation}

dove $\tau_{ref} \approx 2$ ms è il periodo refrattario assoluto tipico.

\subsubsection{Rilevanza Biologica}

In vivo, le frequenze di scarica tipiche nella corteccia cerebrale sono dell'ordine di \textbf{poche decine di Hz} (tipicamente 5–80 Hz), ben al di sotto di $\nu_{max} = 500$ Hz. Questo non è casuale: l'SFA è uno dei meccanismi che \textbf{impedisce} ai neuroni di avvicinarsi alla frequenza massima in condizioni fisiologiche normali.

\begin{tcolorbox}[colback=gray!5, colframe=gray!40, arc=2mm, boxrule=0.5pt]
 Se i neuroni corticali potessero scaricare a 500 Hz, il consumo energetico del cervello diventerebbe insostenibile (già a riposo il cervello consuma circa il 20\% dell'energia corporea totale). L'SFA agisce come un meccanismo di \textbf{risparmio energetico}: abbassa automaticamente il firing rate in risposta a stimoli prolungati.
\end{tcolorbox}

L'SFA ha anche un \textbf{significato computazionale}: favorisce la codifica dei \textbf{transitori} degli stimoli (onset) rispetto alla loro componente stazionaria. Un neurone con SFA risponde con un burst iniziale ad alta frequenza e poi adatta la sua scarica, "segnalando" l'inizio di un nuovo stimolo più che la sua durata.


\section{Trasmissione Sinaptica Chimica}

 Le correnti che il neurone riceve e che determinano il suo firing non sono generate direttamente dall'interno della cellula, ma arrivano da altri neuroni attraverso le \textbf{sinapsi}.

Nel sistema nervoso centrale dei mammiferi, la forma di comunicazione interneuronale più comune è la \textbf{trasmissione sinaptica chimica} (in opposizione alle \textbf{sinapsi elettriche} o gap junctions, che sono presenti in strutture più antiche filogeneticamente come il cervelletto e il talamo, ma anche in alcune regioni corticali).

\subsubsection{Morfologia della Sinapsi Chimica}

La struttura morfologica di una sinapsi chimica comprende tre componenti principali:

\begin{enumerate}
    \item \textbf{Bottone presinaptico} (presynaptic terminal o bouton): è il rigonfiamento terminale dell'assone del neurone presinaptico. Contiene migliaia di \textbf{vescicole sinaptiche} — piccole sfere membranose (diametro $\approx$ 40 nm) riempite di molecole di \textbf{neurotrasmettitore}.
    \item \textbf{Fessura sinaptica} (synaptic cleft): spazio extracellulare di circa \textbf{20 nm} che separa la membrana presinaptica da quella postsinaptica. È crucialmente stretto: i neurotrasmettitori possono diffondere attraverso di esso in pochi microsecondi.
    \item \textbf{Membrana postsinaptica}: si trova tipicamente sul \textbf{dendrite} (o sul soma) del neurone postsinaptico. Contiene alta densità di \textbf{recettori ionotropici} (canali ionici ligando-dipendenti) nella \textbf{densità postsinaptica} (PSD, PostSynaptic Density).
\end{enumerate}

\subsubsection{Neuroni Piramidali e Organizzazione Corticale}

I principali neuroni eccitatori della corteccia cerebrale sono i \textbf{neuroni piramidali}. Hanno un \textbf{dendrite apicale} orientato verticalmente verso la superficie corticale e numerosi \textbf{dendriti basali}. Le sinapsi sono distribuite sull'albero dendritico: le sinapsi apicali (provenienti da neuroni di altri strati corticali o aree remote) sono vicine alla superficie, mentre le sinapsi basali ricevono input locali.

\subsubsection{Il Ruolo del Calcio nel Processo Presinaptico}

Quando uno spike raggiunge il bottone presinaptico, la \textbf{depolarizzazione della membrana presinaptica} apre i canali Ca$^{2+}$ \textbf{HVA} presenti nel terminale. Il Ca$^{2+}$ extracellulare entra rapidamente nel bottone presinaptico. Questo evento è \textbf{indispensabile} per la trasmissione: in assenza di Ca$^{2+}$ extracellulare, gli spike presinaptici non producono rilascio di neurotrasmettitori.

Il Ca$^{2+}$ che entra nel bottone presinaptico si lega a proteine sensori del calcio (come la \textbf{sinaptotagmina}), inducendo la fusione delle vescicole sinaptiche con la membrana presinaptica (esocitosi).


\subsection{Meccanismo di Esocitosi}

\subsubsection{Le Proteine SNARE e il Meccanismo di Fusione}

Il meccanismo molecolare dell'esocitosi è mediato dalle \textbf{proteine SNARE} (Soluble NSF Attachment Protein Receptors):

\begin{itemize}
    \item \textbf{Synaptobrevin} (o VAMP): proteina SNARE sulla vescicola (v-SNARE).
    \item \textbf{Syntaxin} e \textbf{SNAP-25}: proteine SNARE sulla membrana presinaptica (t-SNARE).
\end{itemize}

Quando il Ca$^{2+}$ si lega alla \textbf{sinaptotagmina} (il sensore del calcio), si scatena un cambiamento conformazionale che permette alle proteine SNARE di formare un complesso a quattro eliche (SNARE complex). Questo complesso è un motore molecolare che "cuce" insieme la membrana vescicolare e quella presinaptica, permettendo la \textbf{fusione} e l'apertura di un poro di fusione (fusion pore).

\subsubsection{Rilascio dei Neurotrasmettitori nella Fessura}

Una volta aperto il poro di fusione, le molecole di \textbf{neurotrasmettitore} contenute nella vescicola (tipicamente migliaia di molecole per vescicola: ad es. $\sim$5000 molecole di glutammato per vescicola glutamatergica) vengono rilasciate nella fessura sinaptica per \textbf{diffusione}.

Date le dimensioni nanometriche della fessura (20 nm), la diffusione è estremamente rapida: le molecole di neurotrasmettitore raggiungono i recettori postsinaptici in pochi \textbf{microsecondi}.

\subsubsection{Apertura dei Canali Ionici Postsinaptici}

I neurotrasmettitori rilasciati si legano ai \textbf{recettori ionotropici} sulla membrana postsinaptica. I recettori ionotropici sono canali ionici \textbf{ligando-dipendenti}: si aprono quando una o più molecole di neurotrasmettitore si legano al loro sito di riconoscimento. Esempi principali:

\begin{table}[htbp]
    \centering
    \renewcommand{\arraystretch}{1.3}
    \begin{tabularx}{\textwidth}{|X|X|X|X|}
    \hline
    \textbf{Neurotrasmettitore} & \textbf{Recettore} & \textbf{Ione principale} & \textbf{Effetto} \\
    \hline
    Glutammato & AMPA & Na$^+$ & Depolarizzazione (EPSP) \\
    \hline
    Glutammato & NMDA & Na$^+$/Ca$^{2+}$ & Depolarizzazione (EPSP) \\
    \hline
    GABA & GABA$_A$ & Cl$^-$ & Iperpolarizzazione (IPSP) \\
    \hline
    Glicina & Recettore glicina & Cl$^-$ & Iperpolarizzazione (IPSP) \\
    \hline
    \end{tabularx}
\end{table}

Il lavoro sperimentale su questi meccanismi è stato premiato con il \textbf{Nobel per la Fisiologia o Medicina 1963}, assegnato a \textbf{John Eccles} (sinapsi), \textbf{Alan Hodgkin} e \textbf{Andrew Huxley} (potenziale d'azione).


\section{EPSP, IPSP e Chiusura Spontanea dei Canali}

\subsubsection{EPSP (Excitatory PostSynaptic Potential)}

Un \textbf{EPSP} (Potenziale Postsinaptico Eccitatorio) è la variazione di potenziale di membrana indotta dall'apertura di canali ionici eccitatori (tipicamente Na$^+$ o Na$^+$/K$^+$):

\begin{itemize}
    \item I canali AMPA si aprono $\rightarrow$ ingresso di Na$^+$ $\rightarrow$ depolarizzazione della membrana postsinaptica.
    \item Il potenziale di membrana si allontana dal potenziale di riposo verso valori meno negativi (più vicini a $E_{Na} \approx +50$ mV).
    \item L'ampiezza tipica di un singolo EPSP in un neurone corticale di mammifero è di $\sim 0.1$–$1$ mV (quindi molto al di sotto della soglia di firing $\sim 20$ mV sopra il riposo).
\end{itemize}

\subsubsection{IPSP (Inhibitory PostSynaptic Potential)}

Un \textbf{IPSP} (Potenziale Postsinaptico Inibitorio) è la variazione di potenziale opposta:

\begin{itemize}
    \item I canali GABA$_A$ si aprono $\rightarrow$ ingresso di Cl$^-$ (o uscita di K$^+$) $\rightarrow$ iperpolarizzazione della membrana postsinaptica.
    \item Il potenziale di membrana scende verso $E_{Cl}$ o $E_K$, che sono entrambi più negativi del potenziale di riposo.
    \item Questo rende il neurone postsinaptico più difficile da portare a soglia (inibizione).
\end{itemize}

\subsubsection{Decadimento dell'EPSP: Rising e Decaying Phase}

Una volta che i neurotrasmettitori iniziano a diffondere fuori dalla fessura sinaptica, la loro concentrazione nella fessura stessa diminuisce. Di conseguenza i recettori ionotropici postsinaptici \textbf{si chiudono spontaneamente} (per dissociazione del ligando), e la corrente postsinaptica decade.

L'andamento temporale di un EPSP presenta quindi due fasi:
\begin{enumerate}
    \item \textbf{Rising phase} (fase ascendente): durata $\sim 1$–$2$ ms, corrispondente all'apertura dei canali.
    \item \textbf{Decaying phase} (fase discendente): durata variabile (da $\sim 3$ ms per sinapsi veloci AMPA fino a centinaia di ms per sinapsi NMDA/metabotropiche), corrispondente alla chiusura spontanea dei canali.
\end{enumerate}




\subsubsection{L'Esperimento di Stimolazione Ripetuta}

Un esperimento fondamentale per comprendere la trasmissione sinaptica consiste nel stimolare ripetutamente il neurone presinaptico con spike identici (stessa ampiezza, stesso timing) e misurare la risposta postsinaptica. I risultati sono sorprendenti:

\begin{itemize}
    \item A volte si osserva un EPSP \textbf{grande} ($\Delta V \approx 1$–$1.2$ mV).
    \item A volte un EPSP \textbf{piccolo} ($\Delta V \approx 0.4$ mV).
    \item A volte si osserva un \textbf{failure} (nessun EPSP: $\Delta V = 0$).
\end{itemize}

La risposta non è deterministica: è \textbf{stocastica}. Ogni stimolazione ha un esito diverso, anche a parità di stimolo presinaptico. Questo è un dato sperimentale robusto e riproducibile.

\subsubsection{Analisi Quantale: Distribuzione Multimodale di $\Delta$V}

La distribuzione degli $\Delta V$ osservati sperimentalmente presenta caratteristicamente più \textbf{picchi (modes) equidistanti}, separati di circa $\mu_1 \approx 0.4$ mV. Questa struttura a picchi multipli è la firma dell'\textbf{analisi quantale} della trasmissione sinaptica.

L'ipotesi fondamentale è:
\begin{itemize}
    \item Il rilascio del neurotrasmettitore avviene in unità discrete dette \textbf{quanti}: ciascun quanto corrisponde alla fusione di \textbf{una singola vescicola} e produce un EPSP di ampiezza $\mu_1$ (EPSP unitario o mEPSP — miniature EPSP).
    \item Il numero di quanti rilasciati per ogni stimolo sinaptico $k$ è una variabile aleatoria.
\end{itemize}



\section{Modello Statistico: Distribuzione Binomiale + Gaussiana}

\subsubsection{La Distribuzione Binomiale per il Numero di Vescicole}

L'ipotesi del modello quantale è che la sinapsi disponga di $N$ siti di rilascio (release sites), ciascuno con una vescicola pronta al rilascio. Per ogni spike presinaptico, ciascun sito rilascia la sua vescicola con probabilità $p$ (probabilità di rilascio), indipendentemente dagli altri siti. Il numero totale di quanti (vescicole) rilasciati $k$ segue quindi una \textbf{distribuzione binomiale}:

\begin{equation}
P(k) = \binom{N}{k} p^k (1-p)^{N-k}
\end{equation}

I casi speciali:
\begin{itemize}
    \item $k = 0$ (failure): $P(0) = (1-p)^N$.
    \item $k = N$ (rilascio massimale): $P(N) = p^N$.
    \item Valore atteso: $\langle k \rangle = Np$ (\textbf{quantal content}).
\end{itemize}

\subsubsection{La Distribuzione Gaussiana per il $\Delta$V dato k}

Fissato il numero di quanti rilasciati $k$, l'ampiezza dell'EPSP $\Delta V$ non è esattamente $k \mu_1$, ma fluttua attorno a questo valore. Questo è dovuto a:
\begin{itemize}
    \item Variabilità nell'ampiezza del singolo EPSP unitario (fluttuazioni del numero di molecole per vescicola, della conduttanza dei recettori, ecc.).
    \item Rumore di fondo nel potenziale di membrana postsinaptico.
\end{itemize}

Si assume che, dato $k$, la distribuzione di $\Delta V$ sia \textbf{gaussiana}:

\begin{equation}
P(\Delta V \mid k) = \mathcal{N}(k\mu_1, \, \sqrt{k}\,\sigma_1)
\end{equation}

dove:
\begin{itemize}
    \item $\mu_1$ è l'ampiezza media del singolo quanto (mEPSP unitario) $\approx 0.4$ mV,
    \item $\sigma_1$ è la deviazione standard dell'EPSP per un singolo quanto.
\end{itemize}

La varianza del $\Delta V$ per $k$ quanti è $k\sigma_1^2$ (perché la somma di $k$ variabili indipendenti con varianza $\sigma_1^2$ ha varianza $k\sigma_1^2$).

\subsubsection{Distribuzione Complessiva del $\Delta$EPSP}

La distribuzione totale degli EPSP osservata sperimentalmente è la \textbf{somma pesata} delle distribuzioni gaussiane (una per ogni $k$), pesata dalla distribuzione binomiale:

\begin{equation}
P(\Delta V) = \sum_{k=0}^{N} P(k) \cdot \mathcal{N}(k\mu_1, \sqrt{k}\,\sigma_1)
\end{equation}

Questa è una \textbf{miscela di gaussiane} (Gaussian mixture model) con pesi dati dalla distribuzione binomiale. Il primo termine ($k=0$) è un delta di Dirac in $\Delta V = 0$ (i failure).

Il fit di questa distribuzione teorica ai dati sperimentali permette di stimare:
\begin{itemize}
    \item $N$: numero di siti di rilascio.
    \item $p$: probabilità di rilascio per sito.
    \item $\mu_1$: ampiezza del quanto unitario.
    \item $\sigma_1$: variabilità del quanto unitario.
\end{itemize}

\begin{tcolorbox}[colback=gray!5, colframe=gray!40, arc=2mm, boxrule=0.5pt]
Questa distribuzione teorica fitta \textbf{perfettamente} i dati sperimentali. Questo costituisce una delle prove più eleganti della natura quantale della trasmissione sinaptica, introdotta originalmente da \textbf{Bernard Katz} (Nobel 1970).
\end{tcolorbox}

\subsubsection{La Trasmissione Sinaptica come Fonte di Stocasticità}

La trasmissione sinaptica è una delle principali \textbf{fonti di stocasticità} nel sistema nervoso. Con $10^{10}$ neuroni e $\sim 10^4$ sinapsi per neurone, il cervello ha $\sim 10^{14}$ sinapsi, ciascuna con una propria probabilità di rilascio. Eppure questo sistema massiccio stocastico produce comportamenti affidabili e riproducibili, grazie a meccanismi di \textbf{averaging di campo medio} (mean-field averaging): le fluttuazioni di una singola sinapsi sono irrilevanti quando centinaia di sinapsi convergono sullo stesso neurone.


\section{Modello Matematico della Conduttanza Sinaptica}

\subsubsection{La Corrente Sinaptica}

La corrente sinaptica prodotta dall'apertura dei canali postsinaptici è descritta dalla forma ohmica:

\begin{equation}
I_{syn}(t) = G_{syn}(t) \cdot (V - E_{syn})
\end{equation}

dove:
\begin{itemize}
    \item $G_{syn}(t) = G_{max} \cdot R(t)$ è la conduttanza sinaptica variabile nel tempo,
    \item $G_{max}$ è la conduttanza massima (quando tutti i canali sono aperti),
    \item $R(t) \in [0, 1]$ è la \textbf{frazione di canali postsinaptici aperti} in funzione del tempo,
    \item $E_{syn}$ è il potenziale di inversione della corrente sinaptica ($E_{syn} = E_{Na} \approx +50$ mV per sinapsi eccitatorie AMPA; $E_{syn} = E_{Cl} \approx -70$ mV per sinapsi inibitorie GABA$_A$).
\end{itemize}

\subsubsection{Equazione Cinetica per R: Un Processo di Markov}

La variabile $R(t)$ descrive la cinetica di transizione dei canali ionici postsinaptici tra lo stato chiuso (C) e lo stato aperto (O):

\begin{equation}
\text{C} \underset{\beta}{\overset{\alpha N(t)}{\rightleftharpoons}} \text{O}
\end{equation}

La dinamica di $R$ come sistema del primo ordine è:

\begin{equation}
\dot{R} = \alpha N(t)(1-R) - \beta R
\end{equation}

dove:
\begin{itemize}
    \item $N(t)$ è la \textbf{concentrazione di neurotrasmettitore} nella fessura sinaptica in funzione del tempo,
    \item $\alpha$ è il rate di transizione chiuso $\to$ aperto (processo bimolecolare),
    \item $\beta$ è il rate di transizione aperto $\to$ chiuso (decadimento monomolecolare),
    \item $(1-R)$ è la frazione di canali nello stato chiuso.
\end{itemize}



Il profilo temporale della concentrazione di neurotrasmettitore $N(t)$ nella fessura sinaptica in risposta a uno spike presinaptico al tempo $t = 0$ è complesso, ma può essere modellato con una \textbf{funzione alpha}:

\begin{equation}
N(t) \propto N_{max} \frac{t}{\tau^2} \cdot e^{-t/\tau_T}
\end{equation}

Questa funzione ha un'ascesa lineare seguita da un decadimento esponenziale, con un picco al tempo $t_{peak} = \tau_T$.


\subsubsection{L'Approssimazione a Scalino (Heaviside)}

Per semplificare la soluzione dell'equazione per $R$, si approssima il profilo temporale del neurotrasmettitore con una \textbf{funzione a scalino di durata finita} $\Delta$:

\begin{equation}
N(t) \approx N_{max} \cdot [\Theta(t)  \Theta(t - \Delta)]
\end{equation}

dove $N_{max}$ è la concentrazione massima di neurotrasmettitore. Questa approssimazione trasforma l'equazione differenziale per $R$ in un sistema \textbf{a coefficienti costanti a tratti}, risolvibile analiticamente in ciascuna delle due fasi.

\subsubsection{Fase di Rising (0 $\leq$ t $\leq$ $\Delta$)}

Durante la fase di rising, $N = N_{max} = \text{cost}$. L'equazione diventa:

\begin{equation}
\dot{R} = (\alpha N_{max} + \beta)\left(\frac{\alpha N_{max}}{\alpha N_{max} + \beta} - R\right)
\end{equation}

La soluzione è un'esponenziale che si avvicina all'asintoto:

\begin{equation}
R_\infty^{rise} = \frac{\alpha N_{max}}{\alpha N_{max} + \beta} \approx 1 \quad \text{(poiché } \alpha N_{max} \gg \beta \text{)}
\end{equation}

con costante di tempo:

\begin{equation}
\tau_{rise} = \frac{1}{\alpha N_{max} + \beta} \approx \frac{1}{\alpha N_{max}}
\end{equation}

Tipicamente $\tau_{rise} \sim 1$ ms, dell'ordine della durata dello spike presinaptico.

\subsubsection{Fase di Decaying (t > $\Delta$)}

Quando $N = 0$ (il neurotrasmettitore è diffuso fuori dalla fessura), l'equazione si riduce a:

\begin{equation}
\dot{R} = -\beta R
\end{equation}

La soluzione è il puro decadimento esponenziale:

\begin{equation}
R(t) = R(\Delta) \cdot e^{-\beta(t - \Delta)}
\end{equation}

con asintoto $R_\infty^{decay} = 0$ e costante di tempo:

\begin{equation}
\tau_{decay} = \frac{1}{\beta}
\end{equation}

La conduttanza sinaptica decade quindi con la costante di tempo $\tau_{decay}$, che è determinata unicamente dal rate di chiusura spontanea $\beta$ dei canali.

\subsubsection{Rappresentazione Continua con la Funzione Alpha}

Invece della soluzione a tratti, si può usare la soluzione dell'equazione completa con la funzione alpha come forzante. Il risultato è una curva con:
\begin{itemize}
    \item Slope iniziale $\sim 1/\tau_{rise}$ durante la fase ascendente,
    \item Decadimento esponenziale con $\tau_{decay}$ dopo il picco.
\end{itemize}

Questa curva a "campana" asimmetrica è la \textbf{forma d'onda tipica} di un EPSP o di una conduttanza sinaptica.

\begin{tcolorbox}[colback=gray!5, colframe=gray!40, arc=2mm, boxrule=0.5pt]
Nel passaggio tra l'approssimazione a scalino e la funzione alpha, il parametro $\Delta$ viene aggiustato in modo che l'\textbf{integrale della corrente} (cioè la carica totale trasferita) sia conservato. Questo garantisce che le due rappresentazioni siano equivalenti dal punto di vista della fisiologia postsinaptica.
\end{tcolorbox}




\subsubsection{La Carica per Singolo Evento Sinaptico}

La carica totale $Q$ consegnata al neurone postsinaptico durante un singolo evento sinaptico è:

\begin{equation}
Q = \int_0^\infty I_{syn}(t) \, dt = G_{max} \int_0^\infty R(t) \cdot (V - E_{syn}) \, dt
\end{equation}

Poiché l'ampiezza del singolo EPSP è piccola ($\Delta V \approx 0.4$ mV $\ll$ soglia $\approx 20$ mV), il potenziale di membrana postsinaptico $V$ rimane quasi costante durante l'evento sinaptico. Si può quindi estrarre $(V - E_{syn})$ dall'integrale:

\begin{equation}
Q \approx G_{max} (V - E_{syn}) \int_0^\infty R(t) \, dt
\end{equation}

Di conseguenza, \textbf{la conservazione della carica equivale alla conservazione dell'integrale di} $R(t)$. Nel passaggio tra diverse approssimazioni del profilo temporale di $N(t)$, si sceglie $\Delta$ (o l'ampiezza della funzione alpha) in modo che $\int R(t) dt$ rimanga invariato.

\subsubsection{Classificazione delle Sinapsi per la Velocità}

Le sinapsi chimiche si classificano in base alle costanti di tempo $\tau_{rise}$ e $\tau_{decay}$:

\begin{table}[htbp]
    \centering
    \renewcommand{\arraystretch}{1.3}
    \begin{tabularx}{\textwidth}{|l|l|l|X|X|}
    \hline
    \textbf{Tipo} & \textbf{$\tau_{rise}$} & \textbf{$\tau_{decay}$} & \textbf{Esempio biologico} & \textbf{Funzione} \\
    \hline
    \textbf{Sinapsi veloci} (fast) & $\sim 1$ ms & $3$–$5$ ms & AMPA, GABA$_A$ & Trasmissione fedele di strutture temporali fini \\
    \hline
    \textbf{Sinapsi lente} (slow) & $\sim 1$–$5$ ms & $10$–$100$+ ms & NMDA, GABA$_B$ & Integrazione della frequenza media del treno \\
    \hline
    \end{tabularx}
\end{table}

Per le sinapsi veloci: $\tau_{rise} \approx \tau_{decay}/5$ o meno, e $\tau_{decay} \sim$ durata dello spike.
Per le sinapsi lente: $\tau_{decay} \gg \tau_{rise}$, e la decaying phase domina il profilo temporale.



\section{Il Filtro Passa-Basso Sinaptico: Sinapsi Veloci e Lente}

\subsubsection{L'Approssimazione Esponenziale per Sinapsi con $\tau_{rise} \ll \tau_{decay}$}

Quando $\tau_{rise} \ll \tau_{decay}$, la fase di rising è trascurabile rispetto alla fase di decaying. In questo caso, si può approssimare la risposta sinaptica con un singolo esponenziale con un piccolo ritardo temporale $\delta_t$ (che incorpora l'effetto della fase ascendente):

\begin{equation}
R(t) \approx A \cdot e^{-(t )/\tau_{decay}} \cdot \Theta(t)
\end{equation}

dove $t_f$ è il tempo dello spike presinaptico.

Questa è l'equazione di un \textbf{filtro del primo ordine} applicato al treno di spike presinaptico. In termini di risposta all'impulso (risposta impulsiva del sistema), la funzione $R(t)$ è un esponenziale decrescente con costante di tempo $\tau_{decay}$.

\subsubsection{Il Filtro Sinaptico come Filtro Passa-Basso}

In termini di elaborazione del segnale, la sinapsi chimica (con $\tau_{rise} \ll \tau_{decay}$) si comporta come un \textbf{filtro passa-basso del primo ordine} con frequenza di taglio:

\begin{equation}
f_c = \frac{1}{2\pi \tau_{decay}}
\end{equation}

Questo significa che:
\begin{itemize}
    \item Componenti del segnale a frequenza $f \ll f_c$ (variazioni lente): passano attraverso il filtro quasi inalterate.
    \item Componenti del segnale a frequenza $f \gg f_c$ (variazioni rapide): vengono attenuate (filtrate).
\end{itemize}

\subsubsection{Sinapsi Veloci: Trasmissione Fedele delle Strutture Temporali}

Per sinapsi con $\tau_{decay} \sim 4$ ms, la frequenza di taglio è:

\begin{equation}
f_c \approx \frac{1}{2\pi \times 4 \times 10^{-3}} \approx 40 \, \text{Hz}
\end{equation}

Questo significa che variazioni del treno di spike a frequenze fino a $\sim 40$ Hz vengono trasmesse fedelmente alla membrana postsinaptica. Poiché le frequenze di firing tipiche nella corteccia sono nell'ordine di decine di Hz, le \textbf{sinapsi veloci riproducono fedelmente le strutture temporali fini} del treno di spike presinaptico.

Il potenziale postsinaptico $V_m$ varia quindi in modo simile al profilo temporale del treno di spike: ogni spike presinaptico produce un EPSP riconoscibile e separato.

\subsubsection{Sinapsi Lente: Integrazione della Frequenza Media}

Per sinapsi con $\tau_{decay} \sim 64$ ms, la frequenza di taglio è:

\begin{equation}
f_c \approx \frac{1}{2\pi \times 64 \times 10^{-3}} \approx 2.5 \, \text{Hz}
\end{equation}

In questo caso, le \textbf{alte frequenze del segnale vengono filtrate}: la corrente postsinaptica non segue i singoli spike, ma è quasi costante e riflette la \textbf{frequenza media} del treno presinaptico. Il potenziale $V_m$ varia attorno a un valore quasi costante con oscillazioni molto smorzate.

Le \textbf{sinapsi lente integrano la frequenza media} del treno di spike: sono sensori di rate code, non di spike code temporale.

\subsubsection{Significato Computazionale}

La coesistenza sull'albero dendritico di \textbf{sinapsi veloci} (AMPA, GABA$_A$) e \textbf{sinapsi lente} (NMDA, GABA$_B$) ha un profondo significato computazionale:

\begin{itemize}
    \item Le sinapsi veloci trasmettono informazione con granularità temporale fine (coding temporale, sincronizzazione tra neuroni, binding).
    \item Le sinapsi lente calcolano la frequenza media (rate coding), permettendo al neurone di integrare l'attività del neurone presinaptico su finestre temporali più lunghe.
\end{itemize}

Il neurone postsinaptico riceve quindi informazioni su \textbf{molteplici scale temporali} simultaneamente, aumentando la sua capacità computazionale.

\begin{tcolorbox}[colback=gray!5, colframe=gray!40, arc=2mm, boxrule=0.5pt]
 Questo è uno dei motivi biologici per cui sotto l'albero dendritico di ogni neurone corticale coesistono sinapsi veloci e lente. Non è ridondanza, è complessità computazionale distribuita.
\end{tcolorbox}

\subsubsection{Condizione per la Perdita di Informazione}

La condizione limite per la perdita completa dell'informazione temporale è:

\begin{equation}
f_c = \frac{1}{\tau_{decay}} < f_{segnale}
\end{equation}

dove $f_{segnale}$ è la frequenza caratteristica del contenuto informativo del treno di spike presinaptico. Se questa condizione è soddisfatta, la sinapsi filtra completamente le variazioni temporali del segnale, e la corrente postsinaptica diventa proporzionale alla sola frequenza media del firing presinaptico.


\chapter{Lezione 7 }


\textbf{Argomenti:} Modello di conduttanza sinaptica, funzione alpha, cinetica R, recettori ionotropici veloci e lenti, GABA-A e GABA-B, sinapsi inibitorie, IPSP e potenziale di inversione, dipendenza dell'IPSP dal potenziale di membrana, sinapsi eccitatorie AMPA e NMDA, stato bilanciato eccitazione-inibizione, blocco Mg$^{2+}$ del recettore NMDA, memoria di lavoro, plasticità sinaptica a breve termine (STP), depressione sinaptica a breve termine (STD), deplezione delle vescicole presinaptiche, facilitazione sinaptica a breve termine (STF), riduzione dimensionale del modello HH da 4D a 2D, approssimazione adiabatica di m, anticorrelazione h-n, variabile di recupero W, modello FitzHugh-Nagumo, modello Morris-Lecar



\section{Modello di Conduttanza per la Trasmissione Sinaptica}


Il professore riprende dalla lezione precedente il modello della \textbf{trasmissione sinaptica chimica}. Una sinapsi chimica funziona attraverso la seguente catena di eventi: lo \textbf{spike presinaptico} (potenziale d'azione nel terminale presinaptico) provoca l'esocitosi delle \textbf{vescicole sinaptiche}, con rilascio del \textbf{neurotrasmettitore} nello \textbf{spazio sinaptico} (fessura sinaptica). Il neurotrasmettitore diffonde e si lega ai \textbf{recettori postsinaptici}, che sono canali ionici ligando-dipendenti. Quando i canali si aprono, fluiscono ioni attraverso la membrana postsinaptica, generando una variazione del potenziale di membrana: l'\textbf{EPSP} (Excitatory Postsynaptic Potential) per le sinapsi eccitatorie o l'\textbf{IPSP} (Inhibitory Postsynaptic Potential) per quelle inibitorie.

Il neurotrasmettitore delle sinapsi \textbf{eccitatorie} nei neuroni corticali è il \textbf{glutammato}; quello delle sinapsi \textbf{inibitorie} è l'\textbf{acido gamma-aminobutirrico} (\textbf{GABA}). Questi due neurotrasmettitori sono i principali segnali chimici dell'elaborazione corticale nei vertebrati.

\subsubsection{La Corrente Sinaptica: Modello a Conduttanza}

La \textbf{corrente sinaptica} postsinaptica è descritta, nel quadro del modello a conduttanza, da:

\begin{equation}
I_{syn}(t) = G_{max} \cdot R(t) \cdot (V - E_{syn})
\end{equation}

dove:
\begin{itemize}
    \item $G_{max}$ è la conduttanza massima sinaptica $[\text{nS}]$, determinata dal numero di recettori disponibili e dalla loro conduttanza unitaria,
    \item $R(t) \in [0, 1]$ è la \textbf{frazione di recettori aperti} (noti anche come neuroreceptors), adimensionale,
    \item $V$ è il potenziale di membrana postsinaptico $[\text{mV}]$,
    \item $E_{syn}$ è il \textbf{potenziale di inversione} (reversal potential) della sinapsi $[\text{mV}]$, che dipende dagli ioni permeanti il canale.
\end{itemize}

La corrente ha segno: quando $V > E_{syn}$, la corrente è uscente (depolarizzante se $E_{syn}$ è più positivo di $V$, iperpolarizzante se $E_{syn}$ è più negativo di $V$).

\begin{tcolorbox}[colback=gray!5, colframe=gray!40, arc=2mm, boxrule=0.5pt]
La convenzione corretta per la corrente sinaptica nell'equazione del potenziale di membrana è $-G_{syn}(t)(V - E_{syn})$, con il segno negativo davanti. Questo riflette la convenzione in cui le correnti che depolarizzano la membrana contribuiscono positivamente a $dV/dt$.
\end{tcolorbox}

\subsubsection{La Cinetica dei Recettori: Funzione Alpha Generalizzata}

La dinamica della variabile $R(t)$ è governata dall'equazione:

\begin{equation}
\dot{R} = \alpha N(t)(1 - R) - \beta R
\end{equation}


Il termine $\alpha N(t)(1 - R)$ descrive l'\textbf{apertura} dei canali: è proporzionale alla concentrazione di neurotrasmettitore e alla frazione di recettori chiusi. Il termine $-\beta R$ descrive la \textbf{chiusura spontanea} dei canali con costante di tempo $\tau_{decay} = 1/\beta$ (il decadimento esponenziale che avviene quando il neurotrasmettitore si è diffuso via dallo spazio sinaptico).

La risposta $R(t)$ a uno spike presinaptico singolo ha una forma caratteristica: \textbf{salita rapida} (fase di apertura, determinata da $\alpha$ e dalla durata dell'impulso di $[T]$) seguita da \textbf{decadimento esponenziale} (fase di chiusura, con costante $\tau_{decay} = 1/\beta$). Questa forma a campana asimmetrica è la \textbf{funzione alpha} generalizzata.


\subsubsection{Diversità Cinetica nei Recettori per lo Stesso Neurotrasmettitore}

Un aspetto cruciale della trasmissione sinaptica è che per uno stesso \textbf{neurotrasmettitore} possono esistere recettori con cinetiche molto diverse. Il parametro chiave che distingue i diversi tipi di recettori è $\beta$: la costante di chiusura dei canali.

\begin{itemize}
    \item \textbf{Sinapsi veloci} (fast synapses): $\beta$ grande $\Rightarrow$ $\tau_{decay} = 1/\beta$ piccolo $\Rightarrow$ la conduttanza postsinaptica decade rapidamente. Questi recettori codificano eventi sinaptici brevi e precisi.
    \item \textbf{Sinapsi lente} (slow synapses): $\beta$ piccolo $\Rightarrow$ $\tau_{decay} = 1/\beta$ grande $\Rightarrow$ la conduttanza persiste a lungo. Questi recettori producono un'integrazione temporale più ampia dell'input presinaptico.
\end{itemize}

\subsubsection{Conduttanza Sinaptica Totale come Combinazione Lineare}

Quando un neurone postsinaptico esprime sia recettori veloci sia recettori lenti per lo stesso neurotrasmettitore, la \textbf{conduttanza sinaptica totale} è una combinazione lineare delle risposte dei due tipi:

\begin{equation}
G_{syn}(t) = G_{max} \left[ \alpha_f \cdot R_f(t) + (1 - \alpha_f) \cdot R_s(t) \right]
\end{equation}

dove $\alpha_f \in [0,1]$ è la \textbf{frazione di componente veloce} (fast) e $(1 - \alpha_f)$ è la frazione di componente lenta (slow). Ciascuna variabile $R_f(t)$ e $R_s(t)$ segue la cinetica descritta nella sezione precedente, con parametri $\beta_f$ (grande) e $\beta_s$ (piccolo) rispettivamente.


\subsubsection{I Due Tipi di Recettori Inibitori}

Il GABA (acido gamma-aminobutirrico) agisce su due classi principali di recettori:

\begin{itemize}
    \item \textbf{GABA-A}: recettore ionotropico. È un canale ionico ligando-dipendente che, in seguito al legame con il GABA, apre un canale permeabile prevalentemente agli ioni $\text{Cl}^-$ (e in misura minore a $\text{HCO}_3^-$). È una sinapsi \textbf{veloce}: $\tau_{decay} \approx 6 \, \text{ms}$. Il $\tau_{rise} \approx 1 \, \text{ms}$ per entrambi i tipi (limitato dalla durata dello spike presinaptico).
    \item \textbf{GABA-B}: recettore metabotropico. È un recettore accoppiato a proteina G (GPCR). Il GABA-B attiva indirettamente canali per il $\text{K}^+$ (canali GIRK, G-protein-coupled Inwardly-Rectifying Potassium channels) attraverso una cascata di segnale intracellulare. È una sinapsi \textbf{lenta}: $\tau_{decay} \approx 200 \, \text{ms}$.
\end{itemize}

\begin{table}[htbp]
    \centering
    \renewcommand{\arraystretch}{1.3}
    \begin{tabularx}{\textwidth}{|X|X|X|}
    \hline
    \textbf{Proprietà} & \textbf{GABA-A} & \textbf{GABA-B} \\
    \hline
    Tipo di recettore & Ionotropico (canale ligando-dipendente) & Metabotropico (GPCR) \\
    \hline
    Ioni permeanti & $\text{Cl}^-$ (prevalente), $\text{HCO}_3^-$ & $\text{K}^+$ (GIRK) \\
    \hline
    $\tau_{decay}$ & $\approx 6 \, \text{ms}$ & $\approx 200 \, \text{ms}$ \\
    \hline
    Velocità & Veloce & Lenta \\
    \hline
    Meccanismo & Diretto (canale si apre al legame) & Indiretto (cascasca G-proteica) \\
    \hline
    \end{tabularx}
\end{table}




\subsubsection{Meccanismo Ionico dell'IPSP}

I recettori \textbf{GABA-A} aprono canali selettivi per $\text{Cl}^-$ e $\text{K}^+$:
\begin{itemize}
    \item Lo ione $\text{K}^+$ \textbf{fuoriesce} dalla cellula (perdita di cariche positive intracellulari).
    \item Lo ione $\text{Cl}^-$ \textbf{entra} nella cellula (aggiunta di cariche negative intracellulari).
\end{itemize}

Entrambi i movimenti ionici portano a una \textbf{iperpolarizzazione} della membrana postsinaptica: il potenziale di membrana diventa più negativo, allontanandosi dalla soglia di emissione degli spike. Questo è il meccanismo del \textbf{potenziale postsinaptico inibitorio (IPSP)}.

I potenziali di inversione individuali sono:
\begin{equation}
E_K \approx -70 \, \text{mV}, \qquad E_{Cl} \approx -75 \, \text{mV}
\end{equation}

Il potenziale di inversione della sinapsi inibitoria $E_{inh}$ è una media pesata tra $E_K$ e $E_{Cl}$ (pesata sulle rispettive permeabilità del canale):

\begin{equation}
E_{inh} \approx -70 \div -75 \, \text{mV}
\end{equation}

La corrente inibitoria è:

\begin{equation}
I_{inh} = G_{inh}(t)(V - E_{inh}), \qquad E_{inh} \approx -70 \, \text{mV}
\end{equation}


\subsubsection{La Forza Motrice della Corrente Inibitoria}

L'ampiezza dell'IPSP non è fissa: dipende dal \textbf{potenziale di membrana} del neurone postsinaptico al momento dell'arrivo dello spike inibitorio. La corrente inibitoria è:

\begin{equation}
I_{inh}(t) = G_{inh}(t) \cdot (V - E_{inh})
\end{equation}

Il fattore $(V - E_{inh})$ è la \textbf{forza motrice} (driving force). Nel modello HH standard, il potenziale di membrana a riposo è $E_m = -65 \, \text{mV}$, mentre $E_{inh} \approx -70 \, \text{mV}$.

Poiché $E_m = -65 \, \text{mV} > E_{inh} = -70 \, \text{mV}$, in condizioni di riposo la forza motrice è positiva:

\begin{equation}
(V - E_{inh}) = -65 - (-70) = +5 \, \text{mV} > 0
\end{equation}

La corrente inibitoria è quindi una \textbf{corrente uscente} (da dentro a fuori per gli ioni positivi equivalenti), che iperpolarizza la membrana. Si noti che nella convenzione HH il segno della corrente postsinaptica nell'equazione del potenziale è negativo ($-G_{inh}(V - E_{inh})$), il che significa che questa corrente riduce $dV/dt$.

\subsubsection{Saturazione e Sicurezza Biologica}

Si osserva il seguente comportamento cruciale:

\begin{enumerate}
    \item \textbf{Se il neurone è fortemente depolarizzato} ($V \gg E_{inh}$): la forza motrice $(V - E_{inh})$ è grande $\Rightarrow$ la corrente inibitoria è intensa $\Rightarrow$ l'IPSP è grande.
    \item \textbf{Se il neurone è a riposo} ($V \approx E_m = -65 \, \text{mV}$): la forza motrice è piccola (+5 mV) $\Rightarrow$ l'IPSP è piccolo.
    \item \textbf{Se il neurone è già iperpolarizzato a} $V = E_{inh}$: la forza motrice è zero $\Rightarrow$ l'IPSP è zero. La corrente inibitoria non può spingere il potenziale di membrana al di sotto di $E_{inh}$.
\end{enumerate}

Questo costituisce un \textbf{meccanismo di sicurezza biologica}: l'iperpolarizzazione massima raggiungibile per via inibitoria è limitata da $E_{inh} \approx -70/-75 \, \text{mV}$. Il sistema è intrinsecamente saturante: non è possibile iperpolarizzare il neurone indefinitamente tramite stimolazione inibitoria.

\begin{tcolorbox}[colback=gray!5, colframe=gray!40, arc=2mm, boxrule=0.5pt]
Questo comportamento "bounded" dell'IPSP è un esempio di auto-regolazione elettrochimica. Il fatto che l'inibizione sia più forte quando il neurone è più attivo (depolarizzato) crea una retroazione negativa che stabilizza l'attività neuronale.
\end{tcolorbox}



\subsubsection{Recettori per il Glutammato}

Le sinapsi eccitatorie corticali utilizzano il \textbf{glutammato} come neurotrasmettitore. Il glutammato agisce su due principali classi di recettori ionotropici, caratterizzate da cinetiche molto diverse:

\begin{itemize}
    \item \textbf{Recettori AMPA} ($\alpha$-amino-3-idrossi-5-metil-4-isossazolpropionato): recettori ionotropici veloci, con $\tau_{decay} \approx 3 \div 5 \, \text{ms}$. Sono canali permeabili principalmente a $\text{Na}^+$ e $\text{K}^+$, e sono responsabili della trasmissione eccitatoria rapida. Costituiscono la principale corrente portante durante la trasmissione sinaptica basale.
    \item \textbf{Recettori NMDA} ($N$-metil-$D$-aspartato): recettori ionotropici lenti, con $\tau_{decay}$ molto maggiore di $\tau_{rise}$. Sono permeabili non solo a $\text{Na}^+$ e $\text{K}^+$ ma anche, in modo importante, agli ioni $\text{Ca}^{2+}$.
\end{itemize}

\subsubsection{Potenziale di Inversione delle Sinapsi Eccitatorie}

I recettori AMPA e NMDA aprono canali permeabili sia a $\text{Na}^+$ che a $\text{K}^+$. Il potenziale di inversione $E_{exc}$ è una media pesata tra $E_{Na} \approx +50 \, \text{mV}$ e $E_K \approx -70 \, \text{mV}$, con una permeabilità maggiore al $\text{Na}^+$ rispetto al $\text{K}^+$. Il risultato è:

\begin{equation}
E_{exc} \approx 0 \, \text{mV}
\end{equation}

Per la sinapsi eccitatoria si ha quindi:

\begin{equation}
I_{exc} = G_{exc}(t)(V - E_{exc}), \qquad E_{exc} \approx 0 \, \text{mV}
\end{equation}

Analogamente all'IPSP, anche la corrente eccitatoria è auto-limitante: quando il potenziale di membrana si avvicina a $E_{exc} \approx 0 \, \text{mV}$, la forza motrice $(V - E_{exc})$ diminuisce, riducendo l'ampiezza dell'EPSP. Non è possibile depolarizzare il neurone al di sopra di $E_{exc}$ tramite stimolazione eccitatoria glutamatergica da sola.



\section{Lo Stato Bilanciato Eccitazione-Inibizione}

\subsubsection{Bombardamento Sinaptico e Dinamica del Potenziale di Membrana}

In condizioni fisiologiche in vivo, un neurone corticale riceve un \textbf{bombardamento sinaptico} (synaptic bombardment) continuo: migliaia di input sinaptici eccitatori e inibitori al secondo, da centinaia o migliaia di neuroni presinaptici. In questo regime, il potenziale di membrana $V_m$ non è mai costante ma fluttua continuamente.

Il \textbf{quadro dello stato bilanciato} (balanced state) descrive la situazione in cui i contributi eccitatori e inibitori sono numericamente bilanciati. In questo stato, si può analizzare dove si assesta il $V_m$ medio.

\subsubsection{Il Meccanismo dell'Attrattore Bilanciato}

Consideriamo un neurone sotto bombardamento sinaptico stazionario con conduttanze eccitatorie $G_{exc}$ e inibitorie $G_{inh}$ medie non nulle. La forza motrice di ciascuna sinapsi dipende da $V$:

\begin{itemize}
    \item \textbf{Quando $V$ è vicino a $E_{exc} \approx 0 \, \text{mV}$}: la forza motrice eccitatoria $(V - E_{exc}) \approx 0$ è piccola (EPSP debole), ma la forza motrice inibitoria $(V - E_{inh}) = (0 - (-70)) = +70 \, \text{mV}$ è grande (IPSP forte). L'inibizione prevale $\Rightarrow$ il neurone viene spinto verso valori più negativi di $V$.
    \item \textbf{Quando $V$ è vicino a $E_{inh} \approx -70 \, \text{mV}$}: la forza motrice inibitoria $(V - E_{inh}) \approx 0$ è piccola (IPSP debole), ma la forza motrice eccitatoria $(V - E_{exc}) = (-70 - 0) = -70 \, \text{mV}$ è grande in valore assoluto (EPSP forte). L'eccitazione prevale $\Rightarrow$ il neurone viene spinto verso valori meno negativi.
\end{itemize}

Questo crea un \textbf{attrattore stabile} per il potenziale di membrana. Il sistema è auto-regolante: qualunque deviazione dal valore di equilibrio viene corretta dalla retroazione elettrochimica delle correnti sinaptiche conduttanza-based.

\begin{tcolorbox}[colback=gray!5, colframe=gray!40, arc=2mm, boxrule=0.5pt]
 Questo meccanismo di retroazione negativa è automatico e non richiede alcun meccanismo neurale aggiuntivo: emerge direttamente dalle proprietà fisiche della conduzione ionica attraverso i canali selettivi. È una proprietà fondamentale dei neuroni reali operanti sotto bombardamento sinaptico.
\end{tcolorbox}

\subsubsection{Il Potenziale di Equilibrio Bilanciato}

Il valore di equilibrio del potenziale di membrana nello stato bilanciato dipende dal rapporto tra le conduttanze eccitatorie e inibitorie medie. Il $V_m$ è \textbf{sempre compreso} tra i due potenziali di inversione:

\begin{equation}
E_{inh} \lesssim V_m \lesssim E_{exc}, \qquad -70 \, \text{mV} \lesssim V_m \lesssim 0 \, \text{mV}
\end{equation}

In condizioni fisiologiche tipiche, con una predominanza relativa dell'inibizione, il valore medio di $V_m$ si assesta intorno a $-60 \div -70 \, \text{mV}$: ben al di sotto della soglia di emissione degli spike ($V_{th} \approx -55 \div -50 \, \text{mV}$).

\subsubsection{Emissione di Spike per Fluttuazione}

Il fatto che il $V_m$ medio sia al di sotto della soglia implica che gli \textbf{spike vengono emessi solo per fluttuazioni} attorno al valore medio. Le fluttuazioni sono causate dalla \textbf{stocasticità} della trasmissione sinaptica: il numero di vescicole rilasciate per ogni spike presinaptico, il numero di recettori postsinaptici attivati e i tempi di arrivo degli spike presinaptici sono tutti processi stocastici. Il neurone, in questo regime, è quindi un \textbf{rivelatore di fluttuazioni}: la sua frequenza di scarica è determinata dall'ampiezza delle fluttuazioni del $V_m$ rispetto alla soglia.

Questa irregolarità della scarica è pienamente coerente con le osservazioni sperimentali in vivo nei neuroni corticali, che mostrano una distribuzione degli ISI approssimativamente esponenziale (processo di Poisson) con coefficiente di variazione $CV \approx 1$.


\subsection{Il Recettore NMDA}

\subsubsection{Il Blocco Voltaggio-Dipendente da Mg$^{2+}$}

Il recettore \textbf{NMDA} presenta una proprietà di eccezionale importanza funzionale: la sua conduttanza dipende non solo dal legame con il glutammato, ma anche dal \textbf{potenziale di membrana postsinaptico}. Questo lo rende, in un certo senso, sia un recettore chimico che un sensore di voltaggio.

Il meccanismo è il seguente: in condizioni di riposo (bassa attività del neurone postsinaptico, $V_m$ negativo), lo ione $\text{Mg}^{2+}$ penetra nel canale del recettore NMDA e lo \textbf{blocca fisicamente} dall'interno (blocco del canale dal lato intracellulare del vestibolo). Questo blocco da $\text{Mg}^{2+}$ è voltaggio-dipendente:

\begin{itemize}
    \item \textbf{A basso potenziale} ($V_m$ negativo, neurone inattivo): $\text{Mg}^{2+}$ occupa il sito di blocco nel canale $\Rightarrow$ $G_{NMDA} \approx 0$ (anche in presenza di glutammato).
    \item \textbf{Ad alto potenziale} ($V_m$ meno negativo o positivo, neurone attivo): la forza elettrostatica non è sufficiente a trattenere $\text{Mg}^{2+}$ nel canale $\Rightarrow$ $\text{Mg}^{2+}$ viene espulso $\Rightarrow$ il canale si apre $\Rightarrow$ $G_{NMDA} > 0$.
\end{itemize}

Formalmente, la conduttanza del recettore NMDA è modulata da un fattore di dipendenza da $\text{Mg}^{2+}$:

\begin{equation}
G_{NMDA}(V) \propto \frac{1}{1 + [\text{Mg}^{2+}]_o \cdot e^{-V/V_{Mg}}}
\end{equation}

dove $[\text{Mg}^{2+}]_o$ è la concentrazione esterna di $\text{Mg}^{2+}$ (tipicamente $1 \, \text{mM}$) e $V_{Mg} \approx 16.5 \, \text{mV}$ è un parametro caratteristico della dipendenza dal voltaggio.


\subsubsection{Il Recettore NMDA come Rivelatore di Coincidenza}

La dipendenza sia dal glutammato (richiede attività presinaptica) sia dal voltaggio postsinaptico (richiede attività postsinaptica) rende il recettore NMDA un \textbf{rivelatore di coincidenza} (coincidence detector): si attiva solo quando il neurone presinaptico \textbf{e} quello postsinaptico sono attivi simultaneamente. Questa proprietà è alla base del principio hebbiano dell'apprendimento sinaptico associativo.

\subsubsection{Ruolo nella Memoria di Lavoro}

Il professore sottolinea un ruolo funzionale specifico del recettore NMDA: la \textbf{stabilizzazione della memoria di lavoro} (working memory). Consideriamo un neurone che rappresenta un'informazione (ad esempio, l'immagine di una mela nel campo visivo). La corrente NMDA, essendo attiva solo quando il neurone stesso è attivo (depolarizzato), crea una retroazione eccitatoria auto-sostenuta: il neurone che spara mantiene aperto il suo stesso canale NMDA (attraverso la depolarizzazione), il che sostiene ulteriormente la scarica. Questo meccanismo di \textbf{retroazione positiva voltaggio-dipendente} è un substrato biofisico plausibile per il mantenimento dell'attività persistente osservata nella corteccia prefrontale durante i compiti di memoria di lavoro.



\subsection{Plasticità Sinaptica a Breve Termine}

\subsubsection{La Sinapsi come Filtro Non Stazionario}

Fino a questo punto abbiamo trattato la sinapsi come un sistema con parametri fissi: $G_{max}$, $\alpha$, $\beta$. In realtà, la \textbf{plasticità sinaptica a breve termine} (Short-Term Plasticity, \textbf{STP}) descrive il fatto che l'efficacia di una sinapsi è dinamica: cambia in funzione della storia recente dell'attività presinaptica, su scale temporali che vanno da decine di millisecondi a qualche secondo.

Il termine "breve termine" distingue questi fenomeni dalla \textbf{plasticità sinaptica a lungo termine} (Long-Term Potentiation/Depression, LTP/LTD), che opera su scale di ore o giorni e implica modifiche strutturali permanenti della sinapsi.

\subsubsection{Deplezione dei Neuroreceptors Postsinaptici}

Il primo meccanismo di STP che il professore introduce opera \textbf{sul lato postsinaptico} ed è di natura puramente matematica: la conseguenza diretta della cinetica dei recettori.

Il numero totale di \textbf{neuroreceptors} (canali ionici postsinaptici) disponibili è finito. Dopo l'apertura da parte di uno spike, ogni recettore impega un tempo dell'ordine di $\tau_{decay} = 1/\beta$ per chiudersi e tornare disponibile. Se due spike presinaptici arrivano con un \textbf{intervallo inter-spike (ISI)} inferiore a $\tau_{decay}$:

\begin{itemize}
    \item Al momento del secondo spike, una frazione $R(t_2^-)$ di recettori è ancora \textbf{aperta} (in stato conduttivo).
    \item Solo la frazione $1 - R(t_2^-)$ di recettori è \textbf{chiusa} e disponibile per essere riaperta.
    \item Il salto $\Delta R$ prodotto dal secondo spike è proporzionale ai soli recettori chiusi: $\Delta R_2 \propto (1 - R(t_2^-))$, che è minore di $\Delta R_1 \propto (1 - R(t_1^-)) \approx 1$ (per il primo spike).
\end{itemize}

Il secondo EPSP è quindi più piccolo del primo. Questo è il meccanismo \textbf{postsinaptico} della depressione a breve termine.


\subsection{Depressione Sinaptica a Breve Termine (STD)}

\subsubsection{Formalizzazione Matematica}

Si modella il rilascio del neurotrasmettitore come una serie di impulsi di Dirac ai tempi degli spike presinaptici $\{t_k\}$:

\begin{equation}
N(t) = N_{max} \cdot \sum_k \delta(t - t_k)
\end{equation}

dove $N_{max}$ è una costante proporzionale alla quantità di neurotrasmettitore rilasciato per spike.

L'equazione cinetica per $R(t)$ diventa, tra uno spike e il successivo, un decadimento esponenziale:

\begin{equation}
\dot{R} = -\beta R \qquad \text{(tra gli spike)}
\end{equation}

con soluzione $R(t) = R(t_k^+) \cdot e^{-\beta(t - t_k)}$.

Al momento di ogni spike $t_k$, la variabile $R$ riceve un salto impulsivo:

\begin{equation}
\Delta R_k = \alpha N_{max} \cdot (1 - R(t_k^-))
\end{equation}

Il salto è proporzionale alla \textbf{frazione di recettori chiusi} al momento dello spike: $1 - R(t_k^-)$. Se $R(t_k^-)$ è grande (molti recettori ancora aperti dal spike precedente), il salto è piccolo.

\subsubsection{Dinamica di R durante un Treno di Spike}

Consideriamo un treno di spike regolari con ISI costante $T_{ISI}$:

\begin{itemize}
    \item \textbf{Primo spike} ($t_1$): $R(t_1^-) \approx 0$ (sinapsi a riposo). Il salto è massimo: $\Delta R_1 = \alpha N_{max}$.
    \item \textbf{Spike successivi}: $R(t_k^-)$ cresce progressivamente (i recettori non si richiudono completamente tra uno spike e il successivo se ISI $< \tau_{decay}$). Il salto $\Delta R_k = \alpha N_{max}(1 - R(t_k^-))$ diminuisce progressivamente.
    \item \textbf{Stato stazionario}: $R$ raggiunge asintoticamente un valore massimo $R_{max}$ tale che il guadagno (salto) sia esattamente bilanciato dal decadimento nell'ISI. L'EPSP si stabilizza a un valore di plateau ridotto rispetto al primo.
\end{itemize}

Se ISI $\gg \tau_{decay}$: i recettori si richiudono completamente tra uno spike e il successivo $\Rightarrow$ nessuna depressione.\\
Se ISI $\ll \tau_{decay}$: forte depressione; l'EPSP di plateau è molto inferiore al primo EPSP.



\subsection{Meccanismo Presinaptico della Depressione}

Oltre al meccanismo postsinaptico già descritto (deplezione dei neuroreceptors), esiste un meccanismo analogo che opera sul \textbf{lato presinaptico}: la \textbf{deplezione delle vescicole sinaptiche}.

Le \textbf{vescicole sinaptiche} rappresentano il "serbatoio" di neurotrasmettitore del terminale presinaptico. Ogni spike presinaptico causa l'esocitosi di un certo numero di vescicole. Il processo di ripristino del pool di vescicole pronte al rilascio (il cosiddetto \textbf{readily releasable pool}, RRP) richiede un tempo finito: le vescicole devono essere riciclate per endocitosi o sintetizzate ex novo, riempite di neurotrasmettitore e trasportate alla zona attiva (active zone) della membrana presinaptica per il docking.

Se l'attività presinaptica è intensa (alta frequenza di spike), le vescicole vengono consumate più rapidamente di quanto non vengano ricostituiti $\Rightarrow$ il pool si esaurisce $\Rightarrow$ ogni spike successivo rilascia meno neurotrasmettitore $\Rightarrow$ il segnale postsinaptico si riduce.

Al contrario, se il terminale presinaptico è rimasto inattivo per un periodo sufficiente, il pool di vescicole è completamente ricostituito e il primo spike produce il massimo rilascio di neurotrasmettitore.


\subsection{Facilitazione Sinaptica a Breve Termine (STF)}

\subsubsection{Il Meccanismo Opposto alla Depressione}

La \textbf{facilitazione sinaptica a breve termine} (Short-Term Facilitation, \textbf{STF}) è il meccanismo opposto alla STD: invece di ridursi, l'ampiezza dell'EPSP \textbf{aumenta} progressivamente durante un treno di spike. La STF può prevalere sulla STD in certi tipi sinaptici, a seconda delle caratteristiche molecolari del terminale presinaptico.

Il meccanismo biologico della STF è legato all'accumulo di ioni $\text{Ca}^{2+}$ nel terminale presinaptico durante l'attività. Ogni spike presinaptico apre i canali del calcio voltaggio-dipendenti del terminale, causando un influsso di $\text{Ca}^{2+}$. Se gli spike arrivano ravvicinati (ISI breve), il $\text{Ca}^{2+}$ si accumula progressivamente (la pompa del calcio non riesce a espellerlo abbastanza rapidamente). La concentrazione di $\text{Ca}^{2+}$ residua (residual calcium) aumenta, e poiché il rilascio di vescicole è altamente dipendente dal $\text{Ca}^{2+}$ (si stima una dipendenza da $[\text{Ca}^{2+}]^4$), anche un piccolo accumulo di $\text{Ca}^{2+}$ residuo porta a un forte aumento del rilascio $\Rightarrow$ più neurotrasmettitore rilasciato $\Rightarrow$ EPSP più grande.

\subsubsection{Riepilogo: STD vs. STF}

\begin{table}[htbp]
    \centering
    \renewcommand{\arraystretch}{1.3}
    \begin{tabularx}{\textwidth}{|>{\raggedright\arraybackslash}p{3cm}|X|X|}
    \hline
    \textbf{Parametro} & \textbf{STD (Depressione)} & \textbf{STF (Facilitazione)} \\
    \hline
    Effetto sull'EPSP & Riduzione progressiva durante il treno & Aumento progressivo durante il treno \\
    \hline
    Meccanismo principale & Deplezione dei neuroreceptors o delle vescicole & Accumulo di $\text{Ca}^{2+}$ residuo presinaptico \\
    \hline
    Sede principale & Post-sinaptica (recettori) o pre-sinaptica (vescicole) & Pre-sinaptica \\
    \hline
    Recupero & Con $\tau_{decay}$ (chiusura dei canali) o $\tau_{vesicle}$ (ricostituzione) & Con $\tau_{Ca}$ (espulsione del $\text{Ca}^{2+}$ residuo) \\
    \hline
    Effetto sulla codifica & Filtro passa-basso (deprime le alte frequenze) & Filtro passa-alto (potenzia le alte frequenze) \\
    \hline
    \end{tabularx}
\end{table}

\begin{tcolorbox}[colback=gray!5, colframe=gray!40, arc=2mm, boxrule=0.5pt]
In molte sinapsi reali, STD e STF \textbf{coesistono} e interagiscono. Il profilo netto (depressione o facilitazione) dipende dalla frequenza presinaptica, dal tipo sinaptico e dalla temperatura. Alcune sinapsi mostrano facilitazione a basse frequenze e depressione ad alte frequenze (o viceversa). Queste proprietà conferiscono alle sinapsi capacità di filtraggio frequenziale della trasmissione neuronale.
\end{tcolorbox}

\newpage

\section{Transizione al Modello Semplificato}


Il professore introduce la seconda metà della lezione: la \textbf{semplificazione del modello di Hodgkin-Huxley} (HH) da 4 dimensioni a 2. L'obiettivo è rendere il modello trattabile analiticamente tramite l'analisi del \textbf{piano delle fasi} (phase plane analysis), che sarà il tema centrale delle lezioni successive.

Il modello HH completo ha quattro variabili di stato:
\begin{equation}
\{V, m, h, n\}
\end{equation}

L'analisi completa del sistema 4D richiede tecniche avanzate di teoria dei sistemi dinamici (varietà, attrattori, biforcazioni in alta dimensione). La riduzione a 2D consente invece di visualizzare le \textbf{nullcline} nel piano $(V, W)$ e applicare l'analisi grafica del piano delle fasi.



\subsubsection{Le Quattro Variabili del Modello HH e le Loro Scale Temporali}

Il modello HH completo è descritto dal sistema:

\begin{equation}
C_m \frac{dV}{dt} = I_{ext} - \bar{G}_{Na} m^3 h (V - E_{Na}) - \bar{G}_K n^4 (V - E_K) - G_L (V - E_L)
\end{equation}

\begin{equation}
\dot{m} = \frac{m_\infty(V) - m}{\tau_m(V)}, \qquad \dot{h} = \frac{h_\infty(V) - h}{\tau_h(V)}, \qquad \dot{n} = \frac{n_\infty(V) - n}{\tau_n(V)}
\end{equation}

Le scale temporali tipiche delle variabili di gate nel modello HH sono molto diverse tra loro:
\begin{itemize}
    \item $\tau_m(V)$: dell'ordine di \textbf{frazioni di millisecondo} (es. $0.1 \div 0.5 \, \text{ms}$) — scala molto rapida.
    \item $\tau_h(V)$: dell'ordine di \textbf{alcuni millisecondi} (es. $1 \div 10 \, \text{ms}$) — scala lenta.
    \item $\tau_n(V)$: dell'ordine di \textbf{alcuni millisecondi} (es. $1 \div 10 \, \text{ms}$) — scala lenta (simile a $\tau_h$).
\end{itemize}

\subsubsection{Prima Riduzione: Approssimazione Adiabatica di m ($\tau_m \approx 0$)}

Poiché $\tau_m \ll \tau_h, \tau_n, \tau_V$, la variabile di attivazione $m$ segue istantaneamente le variazioni di $V$. Formalmente, nel limite $\tau_m \to 0$:

\begin{equation}
\tau_m \approx 0 \implies m(t) \approx m_\infty(V(t))
\end{equation}

Questa è l'\textbf{approssimazione adiabatica} (o quasi-statica) per la variabile $m$. $m$ non è più una variabile dinamica indipendente ma diventa una \textbf{funzione algebrica} del potenziale di membrana $V$.

Sostituendo $m \to m_\infty(V)$ nel modello HH, il sistema si riduce a:

\begin{equation}
C_m \frac{dV}{dt} = I_{ext} - \bar{G}_{Na} m_\infty^3(V) h (V - E_{Na}) - \bar{G}_K n^4 (V - E_K) - G_L (V - E_L)
\end{equation}

\begin{equation}
\dot{h} = \frac{h_\infty(V) - h}{\tau_h(V)}, \qquad \dot{n} = \frac{n_\infty(V) - n}{\tau_n(V)}
\end{equation}

Il sistema è ora \textbf{tridimensionale}: $\{V, h, n\}$.

\begin{tcolorbox}[colback=gray!5, colframe=gray!40, arc=2mm, boxrule=0.5pt]
 L'approssimazione adiabatica di $m$ è estremamente accurata. La curva $m(t)$ ottenuta dal modello ridotto 3D è praticamente indistinguibile da quella del modello completo 4D durante un potenziale d'azione. La correttezza dell'approssimazione si verifica graficamente: $m(t)$ segue quasi perfettamente $m_\infty(V(t))$.
\end{tcolorbox}


\subsection{Anticorrelazione tra h e n: la Variabile di Recupero W}

\subsubsection{L'Osservazione Empirica di Hodgkin e Huxley}

Hodgkin e Huxley osservarono empiricamente che durante il ciclo del potenziale d'azione, le variabili $h$ (inattivazione del $\text{Na}^+$) e $n$ (attivazione del $\text{K}^+$) si comportano in modo \textbf{anticorrelato}: quando $n$ cresce (apertura progressiva dei canali $\text{K}^+$ durante la ripolarizzazione), $h$ diminuisce (inattivazione progressiva dei canali $\text{Na}^+$), e viceversa.

Più precisamente: se si calcola la combinazione lineare $h + n/2$ (o, equivalentemente, $h + a \cdot n$ per un opportuno coefficiente $a$), questo valore è \textbf{quasi costante} durante il ciclo del potenziale d'azione. Graficamente, nel piano $(n, h)$, la traiettoria del sistema segue approssimativamente una retta.

\subsubsection{La Variabile di Recupero W}

Si introduce la \textbf{variabile di recupero} $W$ come trasformazione lineare di $n$:

\begin{equation}
W = a \cdot n
\end{equation}

con $a \approx 0.78$ (dal fit ai dati originali HH). La variabile $h$ è approssimata come:

\begin{equation}
h \approx b - W, \qquad b \approx 0.78
\end{equation}

o equivalentemente:

\begin{equation}
h \approx b - a \cdot n
\end{equation}

Questa relazione lineare riduce le due variabili indipendenti $\{h, n\}$ a una sola variabile $W$. Il sistema 3D $\{V, h, n\}$ si riduce quindi a un sistema \textbf{bidimensionale} $\{V, W\}$.

\subsubsection{La Dinamica di W}

Poiché $\tau_h(V) \approx \tau_n(V)$ nel modello HH (stesse scale temporali), si può definire una costante di tempo comune $\tau_W(V)$ come media:

\begin{equation}
\tau_W(V) \approx \frac{\tau_h(V) + \tau_n(V)}{2}
\end{equation}

La $\tau_W(V)$ ha forma a campana (bell-shaped): è grande ai valori di $V$ lontani dalla soglia e più piccola attorno alla soglia. L'equazione di evoluzione per $W$ è:

\begin{equation}
\frac{dW}{dt} = \frac{W_\infty(V) - W}{\tau_W(V)}
\end{equation}

dove $W_\infty(V) = a \cdot n_\infty(V)$ è il valore asintotico di $W$ per un dato potenziale $V$.

\begin{tcolorbox}[colback=gray!5, colframe=gray!40, arc=2mm, boxrule=0.5pt]
 La scoperta di Fitzhugh degli anni '60 non fu solo l'anticorrelazione $h$-$n$, ma anche la constatazione che le scale temporali $\tau_h$ e $\tau_n$ sono comparabili, il che giustifica l'introduzione di un'unica variabile $W$ con una singola scala temporale. La riduzione è quindi una semplificazione non solo del numero di variabili, ma anche della struttura temporale del sistema.
\end{tcolorbox}



\subsection{Il Modello HH 2D e le Approssimazioni di Fitzhugh}

\subsubsection{Il Modello HH Ridotto a 2D}

Sostituendo $m = m_\infty(V)$, $n = W/a$ e $h = b - W$ nel modello HH originale, si ottiene il \textbf{modello HH 2D ridotto}:

\begin{equation}
C_m \frac{dV}{dt} = F(V, W)
\end{equation}

\begin{equation}
\frac{dW}{dt} = \frac{W_\infty(V) - W}{\tau_W(V)}
\end{equation}

dove:

\begin{equation}
F(V, W) = I_{ext} - \bar{G}_{Na} \cdot m_\infty^3(V) \cdot (b - W) \cdot (V - E_{Na}) - \bar{G}_K \cdot \left(\frac{W}{a}\right)^4 \cdot (V - E_K) - G_L(V - E_L)
\end{equation}

Questo sistema mantiene la forma funzionale delle variabili di gate originali (sigmoidali) ma è trattabile nel piano delle fasi $(V, W)$.

Il ciclo di un potenziale d'azione nel piano $(V, W)$ è un \textbf{ciclo limite} (limit cycle): una traiettoria chiusa nello spazio delle fasi verso cui convergono le traiettorie vicine. La presenza o assenza di un ciclo limite (e il tipo di biforcazione che lo genera) determina le proprietà di eccitabilità del neurone.



\subsection{Il Modello FitzHugh-Nagumo}

\subsubsection{Le Ulteriori Semplificazioni di Fitzhugh}

\textbf{Richard Fitzhugh}, negli anni '60, identificò non solo la riduzione da 4D a 2D del modello HH, ma propose anche una successiva semplificazione analitica che rende il modello pienamente trattabile algebricamente. Le approssimazioni aggiuntive di Fitzhugh sono:

\begin{enumerate}
    \item \textbf{$W_\infty(V)$ lineare}: la curva sigmoide $W_\infty(V) = a \cdot n_\infty(V)$ viene approssimata con una retta:
\end{enumerate}

\begin{equation}
W_\infty(V) \approx b_1 \cdot V + b_0
\end{equation}

\begin{enumerate}
    \item[2.] \textbf{$\tau_W$ costante} (non dipendente da $V$): la costante di tempo bell-shaped $\tau_W(V)$ viene sostituita con una costante $\tau_W$.
\end{enumerate}

Con queste approssimazioni, l'equazione per $W$ diventa:

\begin{equation}
\tau_W \frac{dW}{dt} = b_1 V + b_0 - W
\end{equation}

Questa equazione è \textbf{lineare} in $W$ e $V$, il che semplifica enormemente l'analisi delle nullcline.

\subsubsection{L'Equazione per V nel Modello FHN}

Per la variabile $V$, Fitzhugh osserva che il termine $n^4 = (W/a)^4 \ll 1$ per $n < 1$ e propone di approssimarlo con un contributo trascurabile. La non-linearità principale del modello è la corrente del sodio, che dipende da $m_\infty^3(V) \cdot (b - W)$. Approssimando $m_\infty^3(V)$ con un polinomio cubico in $V$ (espansione di Taylor della funzione sigmoidale cubicata) e combinando i termini, si ottiene la famosa forma del modello \textbf{FitzHugh-Nagumo (FHN)}:

\begin{equation}
C_m \frac{dV}{dt} \approx V - \frac{V^3}{3} - W + I_{ext}
\end{equation}

\begin{equation}
\tau_W \frac{dW}{dt} = b_1 V + b_0 - W
\end{equation}

Il termine cubico $V - V^3/3$ è la non-linearità essenziale del modello: produce un'isoclina (nullcline per $V$) a forma di cubo nel piano $(V, W)$, con tre rami (due stabili e uno instabile). Questa struttura a N è responsabile dell'eccitabilità del modello.

\subsubsection{Il Contributo di Nagumo: L'Implementazione Circuitale}

\textbf{Jinichi Nagumo} fornì un contributo complementare ma indipendente: la \textbf{realizzazione circuitale} del modello FHN con componenti elettronici:
\begin{itemize}
    \item La non-linearità cubica è implementata con un \textbf{trigger di Schmitt} (Schmitt trigger) o un elemento a resistenza negativa.
    \item La variabile di recupero $W$ è modellata da un circuito \textbf{RL} (resistenza-induttanza).
    \item Il circuito di Nagumo riproduce esattamente le proprietà del modello teorico di Fitzhugh: eccitabilità, ciclo limite, biforcazioni.
\end{itemize}

Il circuito di Nagumo era un importante strumento di simulazione analogica prima dell'era dei computer digitali e dimostrò che le proprietà dell'eccitabilità neuronale sono universali e non specifiche della biologia.



\subsection{Il Modello Morris-Lecar e Anteprima del Piano delle Fasi}

\subsubsection{Il Modello Morris-Lecar (ML)}

\textbf{Cathy Morris e Harold Lecar}, negli anni '80, proposero un modello 2D alternativo al FHN: il \textbf{modello Morris-Lecar (ML)}. A differenza del FHN (che usa polinomi), il modello ML utilizza espressioni \textbf{sigmoidali} basate sulla tangente iperbolica, che approssimano più accuratamente le curve di gate misurate sperimentalmente.

Le funzioni di gate del modello Morris-Lecar sono:

\begin{equation}
m_\infty(V) = \frac{1 + \tanh\!\left(\dfrac{V - V_1}{V_2}\right)}{2}
\end{equation}

\begin{equation}
W_\infty(V) = \frac{1 + \tanh\!\left(\dfrac{V - V_3}{V_4}\right)}{2}
\end{equation}

\begin{equation}
\tau_W(V) = \frac{1}{\cosh\!\left(\dfrac{V - V_3}{2V_4}\right)}
\end{equation}

dove $V_1$, $V_2$, $V_3$, $V_4$ sono parametri liberi del modello, aggiustabili per riprodurre le proprietà di eccitabilità desiderate.

La funzione $\tau_W(V) = 1/\cosh(\cdot)$ è bell-shaped centrata in $V_3$: ha un massimo finito in $V = V_3$ e decade asintoticamente a zero per $|V - V_3| \to \infty$. Questa forma riproduce bene la dipendenza dal voltaggio della costante di tempo delle variabili di gate HH.

\subsubsection{Morris-Lecar vs. FitzHugh-Nagumo: Confronto}

\begin{table}[htbp]
    \centering
    \renewcommand{\arraystretch}{1.3}
    \begin{tabularx}{\textwidth}{|>{\raggedright\arraybackslash}X|>{\raggedright\arraybackslash}p{4cm}|>{\raggedright\arraybackslash}X|}
    \hline
    \textbf{Proprietà} & \textbf{FitzHugh-Nagumo} & \textbf{Morris-Lecar} \\
    \hline
    Anno & 1961 (Fitzhugh), 1962 (Nagumo) & 1981 \\
    \hline
    Forma di $m_\infty(V)$ & Polinomio (cubica) & $\tanh$ (sigmoide) \\
    \hline
    Forma di $W_\infty(V)$ & Lineare ($b_1 V + b_0$) & $\tanh$ (sigmoide) \\
    \hline
    Forma di $\tau_W(V)$ & Costante & $1/\cosh$ (bell-shaped) \\
    \hline
    Accuratezza quantitativa & Bassa (approssimazione grossolana) & Alta (riproduce bene i dati HH) \\
    \hline
    Trattabilità analitica & Alta (nullcline algebriche semplici) & Media (nullcline transcendenti) \\
    \hline
    Uso principale & Analisi qualitativa delle biforcazioni & Modellistica quantitativa \\
    \hline
    \end{tabularx}
\end{table}

\begin{tcolorbox}[colback=gray!5, colframe=gray!40, arc=2mm, boxrule=0.5pt]
\textbf{Nota del docente:} Il modello ML è preferito nelle simulazioni numeriche quando si vuole una buona corrispondenza quantitativa con il modello HH 4D. Il modello FHN è preferito per l'analisi analitica: le sue nullcline polinomiali permettono di calcolare esplicitamente i punti fissi e di tracciare il diagramma di biforcazione in funzione di $I_{ext}$.
\end{tcolorbox}

\subsubsection{Anteprima: il Piano delle Fasi del Modello HH 2D}

Il professore conclude la lezione annunciando che nella prossima lezione si discuterà in dettaglio il \textbf{piano delle fasi del modello HH 2D} (equivalentemente, FHN o ML). Gli strumenti principali saranno:

\begin{itemize}
    \item Le \textbf{nullcline}: le curve nel piano $(V, W)$ su cui rispettivamente $dV/dt = 0$ (nullcline di $V$, la cubica) e $dW/dt = 0$ (nullcline di $W$, la retta per FHN o la sigmoide per ML).
    \item I \textbf{punti fissi}: le intersezioni delle due nullcline, dove entrambe le derivate sono nulle.
    \item La \textbf{stabilità dei punti fissi}: determinata dagli autovalori della matrice Jacobiana del sistema valutata nel punto fisso.
    \item Le \textbf{biforcazioni}: i valori critici del parametro $I_{ext}$ a cui la struttura dei punti fissi cambia qualitativamente (nascita o scomparsa di cicli limite).
\end{itemize}


\chapter{Lezione 8}


\textbf{Argomenti:} riduzione dimensionale del modello HH, modello FitzHugh-Nagumo, modello Morris-Lecar, piano di fase, nullcline, campo vettoriale, punti di equilibrio, analisi lineare e Jacobiano, autovalori, classificazione dei punti fissi (nodo stabile/instabile, fuoco stabile/instabile, sella), eccitabilità di tipo I e tipo II, singolo spike vs. firing ripetitivo, cicli limite, biforcazioni, biforcazione saddle-node su cerchio invariante (SNIC), biforcazione di Hopf, curva f-I, termine esponenziale del campo di forza, modello integrate-and-fire esponenziale


\section{Modelli Bidimensionali}

\subsubsection{Il Problema della Dimensionalità}

Il modello di \textbf{Hodgkin-Huxley} (HH) è un sistema dinamico a \textbf{quattro dimensioni}: le variabili di stato sono il potenziale di membrana $V$ e le tre variabili di gate $m$, $n$, $h$. Sebbene questo modello catturi con accuratezza quantitativa la forma del potenziale d'azione, la sua elevata dimensionalità rende difficile l'analisi geometrica della dinamica. Il piano di fase, strumento potentissimo per la comprensione qualitativa dei sistemi dinamici, è applicabile direttamente solo a sistemi bidimensionali. La lezione di oggi prosegue il filo conduttore della \textbf{riduzione dimensionale} affrontato nella lezione precedente, mostrando come i modelli a due dimensioni siano capaci di riprodurre l'intera ricchezza qualitativa del comportamento neuronale.


La motivazione principale per studiare i modelli bidimensionali non è dunque la convenienza computazionale, bensì la possibilità di ottenere una \textbf{comprensione geometrica} della dinamica che il modello a quattro dimensioni non può offrire. Studiare la struttura del piano di fase permette di rispondere a domande qualitative fondamentali.

\subsubsection{Le Approssimazioni che Portano alla Riduzione}

La derivazione del modello ridotto a due dimensioni dal modello HH si basa su due approssimazioni fondamentali.

\begin{enumerate}
    \item \textbf{La variabile $m$ è ultra-rapida.} La costante di tempo $\tau_m(V)$ è dell'ordine di frazioni di millisecondo, enormemente più piccola di $\tau_n$ e $\tau_h$ (dell'ordine di alcuni millisecondi). Questa separazione di scala temporale giustifica l'approssimazione adiabatica $m \approx m_\infty(V)$: la variabile $m$ segue istantaneamente il potenziale di membrana, eliminando un grado di libertà.
    \item \textbf{Le variabili $n$ e $h$ sono anti-correlate.} Analisi numeriche del modello HH mostrano che durante il potenziale d'azione, la variabile di inattivazione del sodio $h$ e la variabile di attivazione del potassio $n$ evolvono in modo approssimativamente anti-correlato: $h \approx 1 - n$ (con opportuna riscalatura). Questo permette di eliminare $h$ e di introdurre un'unica \textbf{variabile di recupero} $w$, che coincide essenzialmente con la corrente del potassio $I_K$.
\end{enumerate}

Il risultato è un sistema di soli due gradi di libertà:

\begin{equation}
C \frac{dV}{dt} = F(V, w) + I_{\text{ext}}
\end{equation}

\begin{equation}
\frac{dw}{dt} = G(V, w)
\end{equation}

dove $F$ e $G$ sono funzioni non lineari che dipendono dal modello specifico considerato.



\subsection{Il Modello FitzHugh-Nagumo}

Il \textbf{modello di FitzHugh-Nagumo} (FHN) è stato introdotto da Richard FitzHugh nel 1961 (e implementato analogicamente da Nagumo nello stesso anno) come il prototipo minimalista di neurone bidimensionale eccitabile. Le sue equazioni sono:

\begin{equation}
\frac{dV}{dt} = V - \frac{V^3}{3} - w + I_{\text{ext}}
\end{equation}

\begin{equation}
\tau \frac{dw}{dt} = V + a - bw
\end{equation}

dove $a$, $b$, $\tau$ sono parametri del modello con $\tau \gg 1$ (il recupero è più lento dell'eccitazione).

Le caratteristiche salienti di questo modello sono:
\begin{itemize}
    \item Il \textbf{campo di forza per $V$} è un \textbf{polinomio cubico} $V - V^3/3 - w$, che dà alla nullcline del potenziale la sua forma caratteristica a "N rovesciata" con un ramo decrescente e due rami crescenti.
    \item Il \textbf{campo di forza per $w$} è \textbf{lineare} in entrambe le variabili, il che rende la nullcline della variabile di recupero una \textbf{retta} nel piano di fase.
\end{itemize}

\subsubsection{Giustificazione dalla Riduzione HH}

Il professore sottolinea che la scelta di una nullcline cubica per $V$ e lineare per $w$ non è arbitraria: essa è giustificata dalla riduzione del modello HH completo.

Analizzando il modello HH ridotto a due dimensioni, si osserva che:
\begin{itemize}
    \item La nullcline $\dot{V} = 0$ del modello ridotto presenta effettivamente una forma \textbf{non monotona} con un andamento compatibile con un polinomio cubico: una curva con due estremi, una regione di pendenza negativa al centro, e due regioni di pendenza positiva ai lati.
    \item La nullcline $\dot{w} = 0$ del modello ridotto si comporta in modo molto simile a una \textbf{retta} nel range di valori fisiologicamente rilevanti, perché la variabile di gate del potassio $n$ può essere linearizzata in quella regione.
\end{itemize}

Per comprendere quest'ultimo punto nel dettaglio: nella nullcline $\dot{w} = 0$ del FHN, l'equazione è $w = b_0 + b_1 V$ con $b_1 > 0$. Se $b_1$ è positivo, significa che al crescere di $V$ aumenta anche il valore di $w$ sull'equilibrio della variabile di recupero. Questo corrisponde al fatto che un potenziale di membrana più positivo apre progressivamente di più i canali del potassio, alzando il livello di equilibrio della corrente inibitoria. Questa linearizzazione è giustificata dal fatto che nella regione in cui $w$ varia significativamente (ovvero nella regione in cui il potassio si apre e si chiude), la relazione tra la variabile $n$ e il potenziale $V$ è approssimativamente lineare.

\begin{tcolorbox}[colback=gray!5, colframe=gray!40, arc=2mm, boxrule=0.5pt]
\textbf{Nota del docente:} Mostrando le nullcline del modello HH ridotto, si verifica visivamente che la forma cubica per la nullcline di $V$ e quella lineare per $w$ non sono semplici scelte matematiche comode, ma riflettono fedelmente la struttura non lineare del modello biofisico completo. Il FitzHugh-Nagumo non è solo un modello fenomenologico: è una approssimazione controllata del modello di Hodgkin-Huxley.
\end{tcolorbox}

\subsubsection{Parametri del Modello FitzHugh-Nagumo e Regime di Eccitabilità}

I parametri $a$, $b$ e $\tau$ del modello FHN controllano il regime di comportamento:

\begin{itemize}
    \item \textbf{$a$}: sposta il punto fisso della variabile di recupero (corrisponde al livello di corrente di base). Valori tipici: $a \approx 0.7$.
    \item \textbf{$b$}: controlla la pendenza della nullcline lineare di $w$ (la retta nel piano di fase). Valori tipici: $b \approx 0.8$.
    \item \textbf{$\tau$}: scala temporale relativa del recupero (tipicamente $\tau = 12.5$, imponendo che il recupero sia 12.5 volte più lento dell'eccitazione).
\end{itemize}

Variando $I_{\text{ext}}$ e i parametri $a$, $b$, $\tau$ si possono ottenere:
\begin{enumerate}
    \item \textbf{Regime di riposo} (punto fisso stabile): con $I_{\text{ext}}$ basso, il sistema ha un unico punto fisso stabile.
    \item \textbf{Regime eccitabile} (singolo spike): con $I_{\text{ext}}$ di ampiezza sufficiente a superare la separatrice.
    \item \textbf{Regime oscillatorio} (ciclo limite, firing ripetitivo): con $I_{\text{ext}}$ tale da destabilizzare il punto fisso.
\end{enumerate}



\subsection{Il Modello Morris-Lecar}

\subsubsection{Formulazione del Modello Morris-Lecar}

Il \textbf{modello di Morris-Lecar} (ML) fu introdotto da Catherine Morris e Harold Lecar nel 1981 per descrivere le oscillazioni del potenziale nella cellula muscolare del percebes (\textit{Barnacle}). È considerato un modello più vicino alla biofisica del modello HH rispetto al FHN, pur mantenendo la semplicità di due gradi di libertà. Le sue equazioni sono:

\begin{equation}
C \frac{dV}{dt} = -G_{Ca} \, m_\infty(V)(V - E_{Ca}) - G_K \, w(V - E_K) - G_L(V - E_L) + I_{\text{ext}}
\end{equation}

\begin{equation}
\frac{dw}{dt} = \phi \, \frac{w_\infty(V) - w}{\tau_w(V)}
\end{equation}

dove:
\begin{itemize}
    \item $m_\infty(V) = \frac{1}{2}\left[1 + \tanh\!\left(\frac{V - V_1}{V_2}\right)\right]$ è la \textbf{funzione di attivazione del calcio} (istantanea, perché il calcio si attiva molto rapidamente)
    \item $w_\infty(V) = \frac{1}{2}\left[1 + \tanh\!\left(\frac{V - V_3}{V_4}\right)\right]$ è il \textbf{valore di equilibrio della variabile di gate del potassio}
    \item $\tau_w(V) = \frac{1}{\cosh\!\left(\frac{V - V_3}{2V_4}\right)}$ è la \textbf{costante di tempo} voltaggio-dipendente del gate del potassio
    \item $\phi$ è un fattore di scala per la velocità del recupero
    \item $G_{Ca}$, $G_K$, $G_L$ sono le conduttanze massime rispettivamente per calcio, potassio e leakage
    \item $E_{Ca}$, $E_K$, $E_L$ sono i rispettivi potenziali di inversione
\end{itemize}

\subsubsection{Descrizione Fisica delle Variabili del Morris-Lecar}

Il professore spiega in modo approfondito la fisica dietro le equazioni del Morris-Lecar. La variabile $w$ nel modello Morris-Lecar corrisponde alla \textbf{frazione di canali del potassio aperti} (o più precisamente, la probabilità di apertura del gate della conduttanza del potassio), in analogia con la variabile di gate $n$ nel modello di Hodgkin-Huxley. A differenza del FHN, dove $w$ è una variabile puramente fenomenologica, nel Morris-Lecar $w$ ha una diretta interpretazione biofisica.

La funzione $w_\infty(V) = \frac{1}{2}[1 + \tanh((V-V_3)/V_4)]$ è la \textbf{curva di attivazione all'equilibrio} del canale del potassio: per $V \ll V_3$, quasi nessun canale è aperto ($w_\infty \to 0$); per $V \gg V_3$, quasi tutti sono aperti ($w_\infty \to 1$). Il parametro $V_3$ è il voltaggio di mezza-attivazione e $V_4$ controlla la pendenza della transizione.

La costante di tempo $\tau_w(V) = 1/\cosh((V-V_3)/(2V_4))$ ha un massimo in $V = V_3$ e decresce sia a valori più positivi che più negativi. Questo modella il fatto che la velocità di apertura e chiusura dei canali ha un massimo intorno al voltaggio di mezza-attivazione e diminuisce agli estremi.

Il calcio, a differenza del potassio, è assunto avere un'attivazione \textbf{istantanea} ($m \approx m_\infty(V)$, con $m_\infty(V) = \frac{1}{2}[1 + \tanh((V-V_1)/V_2)]$): i canali del Ca$^{2+}$ si aprono e chiudono su una scala temporale molto più rapida di quella del potassio, giustificando questa approssimazione adiabatica analoga all'approssimazione su $m$ nel modello HH.


\begin{table}[htbp]
    \centering
    \renewcommand{\arraystretch}{1.3}
    \begin{tabularx}{\textwidth}{|>{\raggedright\arraybackslash}X|>{\raggedright\arraybackslash}X|>{\raggedright\arraybackslash}X|}
    \hline
    \textbf{Caratteristica} & \textbf{FitzHugh-Nagumo} & \textbf{Morris-Lecar} \\
    \hline
    Funzione di gate & Polinomiale (cubica) & Sigmoidale ($\tanh$) \\
    \hline
    Nullcline $\dot{V}=0$ & Cubica esplicita & Forma a S, sigmoide \\
    \hline
    Nullcline $\dot{w}=0$ & Retta & Curva sigmoidale \\
    \hline
    Scale temporali & Parametro $\tau$ fisso & $\tau_w(V)$ voltaggio-dipendente \\
    \hline
    Amplificazione & Lineare & Guadagno sigmoidale realistico \\
    \hline
    Ispirazione biofisica & Fenomenologica & Canali Ca$^{2+}$ e K$^+$ espliciti \\
    \hline
    \end{tabularx}
\end{table}

\begin{tcolorbox}[colback=gray!5, colframe=gray!40, arc=2mm, boxrule=0.5pt]
La forma sigmoidale di $m_\infty(V)$ e $w_\infty(V)$ nel modello Morris-Lecar è analoga alle funzioni di attivazione che abbiamo introdotto per le variabili di gate del modello HH. Il vantaggio del Morris-Lecar è che, pur rimanendo bidimensionale, le sue non linearità sono biofisicamente più realistiche e possono essere messe in corrispondenza diretta con proprietà misurabili dei canali ionici.
\end{tcolorbox}

\subsection{Analisi del Piano di Fase: Nullcline e Punti di Equilibrio}

\subsubsection{Il Piano di Fase come Spazio degli Stati}

Un sistema bidimensionale della forma

\begin{equation}
\dot{V} = F(V, w)
\end{equation}
\begin{equation}
\dot{w} = G(V, w)
\end{equation}

può essere visualizzato in un \textbf{piano di fase} avente come assi le due variabili di stato $V$ (potenziale di membrana, asse orizzontale) e $w$ (variabile di recupero, asse verticale). Ogni punto del piano di fase rappresenta un possibile \textbf{stato del sistema}, e l'evoluzione temporale è rappresentata da una \textbf{traiettoria} in questo spazio. Una proprietà fondamentale delle traiettorie in un sistema autonomo è che \textbf{non si intersecano mai}: per ogni punto del piano di fase esiste una e una sola traiettoria che vi passa.

\subsubsection{Le Nullcline: Definizione e Significato}

Le \textbf{nullcline} sono curve del piano di fase lungo le quali una delle due componenti del campo vettoriale si annulla:

\begin{itemize}
    \item \textbf{Nullcline di $V$} (nullcline rossa): l'insieme dei punti $(V, w)$ tali che $\dot{V} = F(V, w) = 0$.
    Lungo questa curva, il campo vettoriale ha \textbf{solo componente verticale}: le traiettorie la attraversano in direzione verticale.
    \item \textbf{Nullcline di $w$} (nullcline blu): l'insieme dei punti $(V, w)$ tali che $\dot{w} = G(V, w) = 0$.
    Lungo questa curva, il campo vettoriale ha \textbf{solo componente orizzontale}: le traiettorie la attraversano in direzione orizzontale.
\end{itemize}

I punti di \textbf{intersezione} tra le due nullcline sono i \textbf{punti fissi} (o punti di equilibrio) del sistema: in questi punti entrambe le componenti del campo vettoriale sono nulle e il sistema non evolve.

\subsubsection{Forma della Nullcline di $V$ nel Modello Ridotto HH}

Per il modello HH ridotto, la nullcline $\dot{V} = 0$ si presenta come una \textbf{curva non monotona}: ha una prima regione di pendenza positiva per $V$ molto negativo, poi una regione di pendenza negativa intorno al potenziale di riposo, e infine torna a pendenza positiva per $V$ positivo. Questa forma è reminiscente di un \textbf{polinomio cubico} e giustifica retroattivamente l'assunzione del modello FitzHugh-Nagumo.

La forma non monotona ha una spiegazione fisica precisa. Per impostare $\dot{V} = 0$, la corrente eccitatoria del sodio deve bilanciare la corrente inibitoria del potassio più la leakage. A potenziali molto negativi, la conduttanza del sodio è quasi nulla e serve poco potassio per bilanciarla: la nullcline è bassa. Al crescere di $V$, la conduttanza del sodio cresce esponenzialmente (la sigmoide $m_\infty(V)$ nella sua parte ascendente), per cui serve sempre più potassio per bilanciare $\rightarrow$ la nullcline sale steeply. Poi, quando $V$ è molto positivo, la forza motrice del Na$^+$ si riduce (ci si avvicina al potenziale di inversione $E_{Na}$) e la corrente del sodio satura: la nullcline scende di nuovo. Questo meccanismo produce esattamente il profilo a "N rovesciata" con i tre rami caratteristici.

\subsubsection{Forma della Nullcline di $w$ nel Modello Ridotto HH}

La nullcline $\dot{w} = 0$ nel modello ridotto è approssimativamente \textbf{lineare} nella regione fisiologicamente rilevante. Questo è dovuto al fatto che la variabile di gate del potassio $n$ può essere linearizzata nel range di $V$ in cui si osservano le variazioni di $w$. Anche questa osservazione giustifica l'assunzione del FHN di avere una nullcline lineare per la variabile di recupero.


\subsection{Campo Vettoriale e Traiettorie nel Piano di Fase}

\subsubsection{Direzione del Campo Vettoriale nelle Quattro Regioni}

Le due nullcline dividono il piano di fase in quattro regioni, in ciascuna delle quali il segno di $\dot{V}$ e di $\dot{w}$ è costante. Conoscendo il segno di $F$ e $G$ in ciascuna regione, si può dedurre la direzione qualitativa del campo vettoriale.

\textbf{Ragionamento per $\dot{V}$:}
La variabile di recupero $w$ rappresenta la corrente inibitoria del potassio. Quando $w$ è grande, la corrente inibitoria domina sul campo eccitatorio, abbassando il potenziale di membrana. Pertanto:
\begin{itemize}
    \item \textbf{Al di sopra della nullcline $\dot{V} = 0$}: $w$ è grande, la corrente del K$^+$ domina, quindi $\dot{V} < 0$ (il potenziale diminuisce, il campo vettoriale punta verso sinistra).
    \item \textbf{Al di sotto della nullcline $\dot{V} = 0$}: la corrente eccitatoria del Na$^+$ è dominante, quindi $\dot{V} > 0$ (il potenziale aumenta, il campo vettoriale punta verso destra).
\end{itemize}

\textbf{Ragionamento per $\dot{w}$:}
\begin{itemize}
    \item \textbf{A destra della nullcline $\dot{w} = 0$}: $V$ è grande, i canali del K$^+$ si aprono più intensamente, quindi $\dot{w} > 0$ (la variabile di recupero aumenta, il campo vettoriale punta verso l'alto).
    \item \textbf{A sinistra della nullcline $\dot{w} = 0$}: la corrente del K$^+$ si rilassa, quindi $\dot{w} < 0$ (il campo vettoriale punta verso il basso).
\end{itemize}

\subsubsection{Composizione del Campo nelle Quattro Regioni}

Combinando i segni delle due componenti, si ottiene la direzione del campo vettoriale nelle quattro regioni del piano di fase:

\begin{enumerate}
    \item \textbf{Regione in basso a destra} (sotto nullcline $V$, a destra di nullcline $w$): $\dot{V} > 0$, $\dot{w} > 0$ $\rightarrow$ campo diretto in alto a destra.
    \item \textbf{Regione in alto a destra} (sopra nullcline $V$, a destra di nullcline $w$): $\dot{V} < 0$, $\dot{w} > 0$ $\rightarrow$ campo diretto in alto a sinistra.
    \item \textbf{Regione in alto a sinistra} (sopra nullcline $V$, a sinistra di nullcline $w$): $\dot{V} < 0$, $\dot{w} < 0$ $\rightarrow$ campo diretto in basso a sinistra.
    \item \textbf{Regione in basso a sinistra} (sotto nullcline $V$, a sinistra di nullcline $w$): $\dot{V} > 0$, $\dot{w} < 0$ $\rightarrow$ campo diretto in basso a destra.
\end{enumerate}

Questo schema di campo vettoriale spiega le traiettorie \textbf{orbitali} osservate nel piano di fase: il sistema segue un percorso rotazionale attorno al punto fisso.



\subsection{Descrizione del Potenziale d'Azione nel Piano di Fase}

\subsubsection{Perturbazione Iniziale: Lo Spostamento Orizzontale}

Consideriamo il sistema all'equilibrio nel punto fisso $(V^*, w^*)$, che corrisponde al potenziale di riposo $E_m \approx -65 \; \text{mV}$. Quando si inietta un \textbf{impulso di corrente} di durata brevissima, si produce un improvviso cambiamento $\Delta V$ del potenziale di membrana. Poiché la durata dell'impulso è molto inferiore alle scale temporali della variabile di recupero $w$, quest'ultima non ha il tempo di cambiare: si ha cioè $\Delta w \approx 0$.

Questo implica che la perturbazione corrisponde a uno \textbf{spostamento puramente orizzontale} nel piano di fase: il punto che rappresenta lo stato del sistema si sposta da $(V^*, w^*)$ a $(V^* + \Delta V, w^*)$. La perturbazione è quindi:

\begin{equation}
\Delta V = \frac{Q}{C} = \frac{I_{\text{pulse}} \cdot \Delta t}{C}
\end{equation}

dove $Q$ è la carica iniettata, $C$ la capacità di membrana, $I_{\text{pulse}}$ l'intensità e $\Delta t$ la durata dell'impulso.

\subsubsection{Il Potenziale di Riposo nel Piano di Fase}

Prima di descrivere la traiettoria del potenziale d'azione, vale la pena soffermarsi sullo \textbf{stato di riposo} nel piano di fase. All'equilibrio, il sistema si trova al punto fisso $(V^*, w^*)$, dove $V^* \approx -65 \; \text{mV}$ per il modello HH ridotto. Questo punto corrisponde alla situazione in cui la corrente del sodio, quella del potassio e la corrente di leakage si bilanciano esattamente. La variabile di recupero al riposo è $w^* = w_\infty(V^*)$, ovvero il valore di equilibrio della conduttanza del potassio al potenziale di riposo.

Il professore sottolinea che questo punto fisso è un \textbf{fuoco stabile} in condizioni fisiologiche: se perturbato lievemente, il sistema ruota in spirale attorno a esso e converge. La frequenza di queste oscillazioni smorzate è determinata dalla parte immaginaria degli autovalori della Jacobiana $\omega = \text{Im}(\lambda_\pm)$, e la velocità di convergenza dalla parte reale $|S| = |\text{Re}(\lambda_\pm)|$.

\subsubsection{Le Quattro Fasi del Potenziale d'Azione come Traiettoria}

Se la perturbazione è sufficientemente grande da portare il sistema \textbf{oltre la separatrice} (soglia di eccitabilità), la traiettoria nel piano di fase descrive un'orbita che percorre le quattro fasi del potenziale d'azione:

\textbf{Fase 1 — Fase crescente (rising phase):}\\
Il sistema si trova nella regione in basso a destra del piano di fase ($\dot{V} > 0$, $\dot{w} > 0$). Il potenziale di membrana cresce rapidamente verso valori positivi perché la corrente del sodio (eccitatoria) domina. La traiettoria si dirige verso destra e verso l'alto, fino a quando non incrocia la nullcline $\dot{V} = 0$. All'incrocio, la traiettoria è \textbf{verticale} (solo componente $\dot{w} \neq 0$).

\textbf{Fase 2 — Fase decrescente (decaying phase):}\\
Il sistema ha superato la nullcline di $V$ e si trova ora nella regione in alto a destra ($\dot{V} < 0$, $\dot{w} > 0$). Il potenziale di membrana inizia a decrescere, mentre la variabile di recupero $w$ continua ad aumentare (il potassio si sta aprendo progressivamente). Quando la traiettoria incrocia la nullcline $\dot{w} = 0$, la traiettoria è \textbf{orizzontale}.

\textbf{Fase 3 — Iperpolarizzazione:}\\
Il sistema si trova nella regione in alto a sinistra ($\dot{V} < 0$, $\dot{w} < 0$). Il potenziale di membrana scende al di sotto del potenziale di riposo (iperpolarizzazione) mentre la variabile di recupero inizia a diminuire (i canali del K$^+$ si chiudono). Questa corrisponde alla \textbf{fase iperpolarizzante} successiva all'emissione dello spike.

\textbf{Fase 4 — Recupero dell'equilibrio:}\\
Il sistema si trova nella regione in basso a sinistra ($\dot{V} > 0$, $\dot{w} < 0$). Il potenziale di membrana risale verso il valore di riposo mentre $w$ diminuisce verso il suo valore di equilibrio. La traiettoria a spirale converge verso il punto fisso.

\begin{tcolorbox}[colback=gray!5, colframe=gray!40, arc=2mm, boxrule=0.5pt]
\textbf{Nota del docente:} In questa descrizione bidimensionale, il \textbf{periodo refrattario assoluto} corrisponde alla fase in cui il sistema si trova sulla ramo superiore del percorso, dove la variabile di recupero $w$ è molto elevata (i canali del potassio sono aperti al massimo e impediscono qualsiasi nuova eccitazione). Il \textbf{periodo refrattario relativo} corrisponde alla fase di recupero verso l'equilibrio, in cui il sistema è ancora nel territorio subthreshold ma non è del tutto insensibile a una nuova perturbazione.
\end{tcolorbox}



\subsection{La Separatrice: Soglia di Eccitabilità nel Piano di Fase}

\subsubsection{Definizione della Separatrice}

Non tutte le perturbazioni producono un potenziale d'azione completo. Esiste nel piano di fase una curva speciale, chiamata \textbf{separatrice} (o manifold stabile dell'orbita o threshold manifold), che separa le traiettorie che generano un potenziale d'azione da quelle che ritornano direttamente all'equilibrio.

La separatrice è un \textbf{manifold monodimensionale} (una curva nel piano $V$-$w$) che passa in prossimità del ramo intermedio (quello di pendenza negativa) della nullcline cubica. La sua posizione definisce una \textbf{soglia dinamica}: per ogni valore di $w$, esiste un valore soglia di $V$ tale che:
\begin{itemize}
    \item Se la perturbazione porta il sistema \textbf{a destra della separatrice}, il sistema produrrà un potenziale d'azione percorrendo la grande orbita attorno alle nullcline.
    \item Se la perturbazione porta il sistema \textbf{a sinistra della separatrice}, il sistema tornerà all'equilibrio con una piccola oscillazione smorzata (risposta sub-soglia).
\end{itemize}

\subsubsection{Risposta Sub-Soglia}

Quando l'impulso di corrente è \textbf{inferiore alla soglia}, la perturbazione porta il sistema nella regione immediatamente a sinistra della separatrice. In questo caso:
\begin{itemize}
    \item Il sistema percorre una traiettoria che si avvicina brevemente al ramo intermedio della nullcline cubica, producendo una piccola depolarizzazione transitoria.
    \item Il sistema poi attraversa la nullcline $\dot{V} = 0$ sul ramo sinistro (pendenza positiva) ed entra nella regione in cui $\dot{V} < 0$.
    \item La traiettoria a spirale converge verso il punto fisso stabile, producendo oscillazioni smorzate attorno all'equilibrio.
\end{itemize}

\begin{tcolorbox}[colback=gray!5, colframe=gray!40, arc=2mm, boxrule=0.5pt]
\textbf{Nota del docente:} La soglia di eccitabilità non è una retta verticale nel piano di fase, ma una curva. Questo implica che la "soglia" dipende dallo stato iniziale del sistema: se la variabile di recupero $w$ è già elevata (come durante la fase di recupero dopo uno spike), la soglia per un nuovo spike è più alta. Questo è il meccanismo del \textbf{periodo refrattario relativo} in questa descrizione bidimensionale.
\end{tcolorbox}



\section{Analisi Lineare e Jacobiano}

\subsubsection{Linearizzazione Attorno al Punto Fisso}

Per studiare la \textbf{stabilità locale} del punto fisso $(V^*, w^*)$, si ricorre alla \textbf{linearizzazione} del sistema. Si introducono le variabili perturbative:

\begin{equation}
\tilde{V} = V - V^*, \quad \tilde{w} = w - w^*
\end{equation}

che rappresentano la deviazione dallo stato di equilibrio. Sviluppando in serie di Taylor le funzioni $F$ e $G$ al primo ordine attorno al punto fisso, si ottiene il sistema lineare approssimato:

\begin{equation}
\frac{d}{dt}\begin{pmatrix} \tilde{V} \\ \tilde{w} \end{pmatrix} = J \begin{pmatrix} \tilde{V} \\ \tilde{w} \end{pmatrix}
\end{equation}

dove $J$ è la \textbf{matrice Jacobiana} valutata nel punto fisso:

\begin{equation}
J = \begin{pmatrix} \partial_V F & \partial_w F \\ \partial_V G & \partial_w G \end{pmatrix}\bigg|_{(V^*, w^*)}
\end{equation}

Il professore sottolinea che questa approssimazione lineare descrive solo il comportamento \textbf{locale} nell'intorno del punto fisso: le grandi orbite associate alla produzione del potenziale d'azione richiedono l'analisi non lineare globale.

\subsubsection{Struttura della Jacobiana}

Per i modelli bidimensionali considerati, gli elementi della Jacobiana hanno un'interpretazione geometrica precisa:

\begin{itemize}
    \item $\partial_V F = F_V$: \textbf{pendenza della nullcline $\dot{V} = 0$} nel punto fisso (coefficiente angolare della curva cubica)
    \item $\partial_w F = F_w = -1$ (per FHN, è costante): forza dell'accoppiamento inibitorio di $w$ su $V$
    \item $\partial_V G = G_V$: pendenza della nullcline $\dot{w} = 0$ (coefficiente angolare $m_w$ della retta)
    \item $\partial_w G = G_w$: quanto rapidamente $w$ si rilassa verso il suo valore di equilibrio
\end{itemize}

Il professore introduce due parametri compatti:
\begin{itemize}
    \item $m_V$: coefficiente angolare (pendenza) della nullcline $\dot{V} = 0$ nel punto fisso
    \item $m_w$: coefficiente angolare della nullcline $\dot{w} = 0$
\end{itemize}

Scegliendo opportunamente le unità di tempo in modo che la scala temporale di $V$ sia l'unità ($\tau_V = 1$), e normalizzando la forza inibitoria a $-1$, la Jacobiana assume la forma:

\begin{equation}
J = \begin{pmatrix} m_V & -1 \\ \varepsilon \, m_w & -\varepsilon \end{pmatrix}
\end{equation}

dove $\varepsilon = \tau_V / \tau_w \ll 1$ è il rapporto tra la scala temporale rapida del potenziale ($\tau_V \sim 20 \; \text{ms}$) e quella lenta del recupero ($\tau_w \sim 100 \; \text{ms}$).

\begin{tcolorbox}[colback=gray!5, colframe=gray!40, arc=2mm, boxrule=0.5pt]
\textbf{Nota del docente:} Il parametro $\varepsilon$ quantifica la separazione di scala temporale tra i due gradi di libertà. Il fatto che $\varepsilon \ll 1$ significa che la variabile di recupero $w$ è molto lenta rispetto al potenziale di membrana. Questa separazione di scala è cruciale per l'eccitabilità e per l'esistenza della soglia: senza di essa, il sistema non potrebbe produrre potenziali d'azione.
\end{tcolorbox}



\subsection{Classificazione dei Punti Fissi}

\subsubsection{Il Polinomio Caratteristico e gli Autovalori}

La stabilità del punto fisso è determinata dagli \textbf{autovalori} $\lambda_{\pm}$ della matrice Jacobiana $J$. Per una matrice $2 \times 2$, gli autovalori soddisfano il polinomio caratteristico:

\begin{equation}
\lambda^2 - \text{tr}(J) \cdot \lambda + \det(J) = 0
\end{equation}

dove:
\begin{itemize}
    \item $\text{tr}(J) = \lambda_+ + \lambda_-$ è la \textbf{traccia} della Jacobiana
    \item $\det(J) = \lambda_+ \cdot \lambda_-$ è il \textbf{determinante} della Jacobiana
\end{itemize}

Introducendo la notazione $S = \text{tr}(J)/2$ (metà della traccia) e $D = \det(J)$, la soluzione esplicita per gli autovalori è:

\begin{equation}
\lambda_{\pm} = S \pm \sqrt{S^2 - D}
\end{equation}

\subsubsection{Autovalori Reali vs. Complessi}

Il discriminante $\Delta = S^2 - D$ determina la natura degli autovalori:
\begin{itemize}
    \item Se $S^2 > D$ (cioè $\Delta > 0$): gli autovalori sono \textbf{reali} e distinti.
    \item Se $S^2 = D$ (cioè $\Delta = 0$): autovalori reali coincidenti (caso di frontiera).
    \item Se $S^2 < D$ (cioè $\Delta < 0$): gli autovalori sono \textbf{complessi coniugati} $\lambda_{\pm} = S \pm i\sqrt{D - S^2}$, il che implica una \textbf{componente oscillatoria} nella dinamica locale.
\end{itemize}

Nel piano $(S, D)$ (traccia/2 vs. determinante), questa frontiera è descritta dalla \textbf{parabola} $D = S^2$.

\subsubsection{Relazione tra Autovalori e Traccia/Determinante}

Il professore mostra esplicitamente come dalla conoscenza di $S$ e $D$ si ricavino le proprietà degli autovalori:

\begin{itemize}
    \item \textbf{Parte reale:} $\text{Re}(\lambda_\pm) = S \pm \text{Re}\!\left(\sqrt{S^2 - D}\right)$. Se si è fuori dalla parabola (autovalori reali), entrambi gli autovalori hanno parte reale uguale a $\lambda_\pm = S \pm \sqrt{S^2 - D}$. Se si è dentro la parabola (autovalori complessi coniugati), la parte reale di entrambi è uguale a $S$.
    \item \textbf{Parte immaginaria:} Per autovalori complessi, $\text{Im}(\lambda_\pm) = \pm\sqrt{D - S^2}$. Questa è la \textbf{frequenza angolare} delle oscillazioni smorzate (o crescenti) attorno al punto fisso.
\end{itemize}

Da queste relazioni, il professore deduce le regioni nel piano $(S, D)$:
\begin{itemize}
    \item Per $S < 0$ e fuori dalla parabola ($S^2 > D > 0$): $\lambda_+$ e $\lambda_-$ entrambi reali negativi (nodo stabile).
    \item Per $S > 0$ e fuori dalla parabola ($S^2 > D > 0$): $\lambda_+$ e $\lambda_-$ entrambi reali positivi (nodo instabile).
    \item Per $S < 0$ e dentro la parabola ($D > S^2 > 0$): autovalori complessi con $\text{Re} = S < 0$ (fuoco stabile).
    \item Per $S > 0$ e dentro la parabola: autovalori complessi con $\text{Re} = S > 0$ (fuoco instabile).
    \item Per $D < 0$ (qualunque $S$): autovalori reali di segno opposto (sella).
\end{itemize}

\subsubsection{Rappresentazione nel Piano Traccia-Determinante}

Nel piano $(S, D)$ si possono distinguere le seguenti regioni e i corrispondenti tipi di punti fissi:

\textbf{Nodo stabile} (stable node):
\begin{itemize}
    \item Regione: $D > 0$, $S < 0$, $S^2 > D$ (fuori dalla parabola, semipiano sinistro)
    \item Autovalori: entrambi reali e negativi
    \item Dinamica: il sistema converge al punto fisso come \textbf{sovrapposizione di due esponenziali decrescenti}, senza oscillazioni
\end{itemize}

\textbf{Nodo instabile} (unstable node):
\begin{itemize}
    \item Regione: $D > 0$, $S > 0$, $S^2 > D$ (fuori dalla parabola, semipiano destro)
    \item Autovalori: entrambi reali e positivi
    \item Dinamica: il sistema diverge dal punto fisso come \textbf{sovrapposizione di due esponenziali crescenti}
\end{itemize}

\textbf{Fuoco stabile} (stable focus o stable spiral):
\begin{itemize}
    \item Regione: $D > S^2$ (dentro la parabola), $S < 0$
    \item Autovalori: complessi coniugati con parte reale negativa $\text{Re}(\lambda_\pm) = S < 0$
    \item Dinamica: il sistema \textbf{converge al punto fisso con oscillazioni smorzate} (spirale verso il centro)
\end{itemize}

\textbf{Fuoco instabile} (unstable focus o unstable spiral):
\begin{itemize}
    \item Regione: $D > S^2$ (dentro la parabola), $S > 0$
    \item Autovalori: complessi coniugati con parte reale positiva
    \item Dinamica: il sistema si \textbf{allontana dal punto fisso con oscillazioni crescenti} (spirale che diverge dal centro)
\end{itemize}

\textbf{Centro}:
\begin{itemize}
    \item Regione: $D > 0$, $S = 0$ (asse verticale, dentro la parabola)
    \item Autovalori: immaginari puri $\lambda_\pm = \pm i\omega$
    \item Dinamica: \textbf{orbite chiuse} (caso non genericamente stabile, tipico dei sistemi Hamiltoniani)
\end{itemize}

\textbf{Sella} (saddle point):
\begin{itemize}
    \item Regione: $D < 0$
    \item Autovalori: uno reale positivo e uno reale negativo
    \item Dinamica: il sistema converge lungo una \textbf{direzione instabile} (manifold instabile) e diverge lungo un'altra (\textbf{manifold stabile}). Un punto sella non è mai stabile.
\end{itemize}


\subsection{Scomposizione Spettrale e Soluzione del Sistema Linearizzato}

\subsubsection{Autovettori come Base dello Spazio}

Per risolvere il sistema lineare $\dot{\mathbf{x}} = J \mathbf{x}$ (dove $\mathbf{x} = (\tilde{V}, \tilde{w})^T$), il professore ricorre alla \textbf{scomposizione spettrale} della Jacobiana. Gli autovettori $|\lambda_+\rangle$ e $|\lambda_-\rangle$ (notazione bra-ket della meccanica quantistica, usata per convenienza) formano una base dello spazio bidimensionale se sono \textbf{ortonormali}:

\begin{equation}
\langle \lambda_+ | \lambda_+ \rangle = 1, \quad \langle \lambda_- | \lambda_- \rangle = 1
\end{equation}
\begin{equation}
\langle \lambda_+ | \lambda_- \rangle = 0, \quad \langle \lambda_- | \lambda_+ \rangle = 0
\end{equation}

L'identità si decompone come:

\begin{equation}
\mathbb{1} = |\lambda_+\rangle\langle \lambda_+| + |\lambda_-\rangle\langle \lambda_-|
\end{equation}

e la Jacobiana si decompone come:

\begin{equation}
J = \lambda_+ |\lambda_+\rangle\langle \lambda_+| + \lambda_- |\lambda_-\rangle\langle \lambda_-|
\end{equation}

\subsubsection{Proiezione e Soluzione Temporale}

Definendo le proiezioni del vettore di stato sugli autovettori:

\begin{equation}
a_{\pm}(t) = \langle \lambda_{\pm} | \mathbf{x}(t) \rangle
\end{equation}

il sistema lineare si disaccoppia in due equazioni \textbf{scalari indipendenti}:

\begin{equation}
\dot{a}_{\pm} = \lambda_{\pm} \, a_{\pm}
\end{equation}

ciascuna con soluzione elementare:

\begin{equation}
a_{\pm}(t) = a_{\pm}(0) \, e^{\lambda_{\pm} t}
\end{equation}

Lo stato del sistema al tempo $t$ è quindi:

\begin{equation}
\mathbf{x}(t) = a_+(t) |\lambda_+\rangle + a_-(t) |\lambda_-\rangle = a_+(0) e^{\lambda_+ t} |\lambda_+\rangle + a_-(0) e^{\lambda_- t} |\lambda_-\rangle
\end{equation}

\subsubsection{Interpretazione Fisica}

Se gli autovalori sono complessi coniugati $\lambda_\pm = S \pm i\omega$ (con $\omega = \sqrt{D - S^2}$), l'esponenziale si decompone:

\begin{equation}
e^{\lambda_\pm t} = e^{St} \left[\cos(\omega t) \pm i\sin(\omega t)\right]
\end{equation}

Questo spiega il comportamento oscillatorio smorzato (o divergente) del fuoco stabile (o instabile): la parte reale $S$ determina se le oscillazioni sono smorzate ($S < 0$) o amplificate ($S > 0$), mentre la parte immaginaria $\omega$ determina la \textbf{frequenza di oscillazione}.

\begin{tcolorbox}[colback=gray!5, colframe=gray!40, arc=2mm, boxrule=0.5pt]
\textbf{Nota del docente:} Questa analisi per autovalori è esatta solo per il sistema linearizzato e vale quindi solo localmente, nell'intorno del punto fisso. Le grandi orbite associate ai potenziali d'azione sono fenomeni non lineari globali che non possono essere descritti da questa analisi locale. Tuttavia, l'analisi lineare è fondamentale per capire se il punto fisso è attrattivo o repulsivo, e con quale modalità (con o senza oscillazioni).
\end{tcolorbox}



\subsection{Condizioni per la Stabilità del Punto Fisso}

\subsubsection{Condizioni Necessarie e Sufficienti}

Un punto fisso è \textbf{stabile} (nel senso di Lyapunov, asintoticamente stabile) se e solo se le parti reali di entrambi gli autovalori sono negative:

\begin{equation}
\text{Re}(\lambda_+) < 0 \quad \text{e} \quad \text{Re}(\lambda_-) < 0
\end{equation}

Nelle variabili traccia e determinante, queste condizioni si esprimono come:

\begin{equation}
S = \frac{\text{tr}(J)}{2} < 0 \quad \text{e} \quad D = \det(J) > 0
\end{equation}

La condizione $D > 0$ garantisce che i due autovalori abbiano lo stesso segno reale (non si ha una sella). La condizione $S < 0$ garantisce che la parte reale comune sia negativa.

\subsubsection{Interpretazione Geometrica: Pendenze delle Nullcline}

Nella forma parametrica della Jacobiana introdotta nella Sezione 8.2:

\begin{equation}
J = \begin{pmatrix} m_V & -1 \\ \varepsilon \, m_w & -\varepsilon \end{pmatrix}
\end{equation}

la traccia e il determinante sono:

\begin{equation}
\text{tr}(J) = m_V - \varepsilon
\end{equation}

\begin{equation}
\det(J) = -m_V \varepsilon + \varepsilon \, m_w = \varepsilon(m_w - m_V)
\end{equation}

Le condizioni di stabilità diventano quindi:

\begin{itemize}
    \item $\text{tr}(J) < 0 \Rightarrow m_V < \varepsilon$: la \textbf{pendenza della nullcline $\dot{V}=0$} deve essere inferiore a $\varepsilon = \tau_V/\tau_w$.
    \item $\det(J) > 0 \Rightarrow m_w > m_V$: la \textbf{pendenza della nullcline $\dot{w}=0$} deve essere maggiore di quella della nullcline $\dot{V}=0$.
\end{itemize}

\subsubsection{Significato Pratico per i Modelli Neuronali}

Il professore chiarisce il ruolo numerico del parametro $\varepsilon$. La costante di tempo della membrana (che governa $V$) è tipicamente dell'ordine di $\tau_V \sim 20 \; \text{ms}$, mentre la costante di tempo del recupero (che governa $w$, ossia la cinetica del potassio) è $\tau_w \sim 100 \; \text{ms}$. Quindi:

\begin{equation}
\varepsilon = \frac{\tau_V}{\tau_w} \approx \frac{20}{100} = 0.2
\end{equation}

Poiché $\varepsilon = \tau_V/\tau_w \ll 1$ (tipicamente $\approx 20/100 = 0.2$), la condizione di stabilità sull'elemento diagonale della Jacobiana diventa approssimativamente:

\begin{equation}
m_V < \varepsilon \approx 0 \quad \Rightarrow \quad m_V < 0
\end{equation}

ovvero la \textbf{pendenza della nullcline del potenziale deve essere negativa} nel punto fisso.

Nel piano di fase del modello FHN (o HH ridotto), questo corrisponde al fatto che il punto fisso deve trovarsi sul \textbf{ramo intermedio decrescente} della curva cubica, non sui due rami laterali crescenti. Quando il punto di equilibrio è sul ramo decrescente, la pendenza locale è negativa e il sistema è stabile.

\begin{tcolorbox}[colback=gray!5, colframe=gray!40, arc=2mm, boxrule=0.5pt]
\textbf{Nota del docente:} Ecco il significato fisico di questa condizione: sul ramo decrescente della nullcline cubica, un piccolo aumento del potenziale di membrana fa sì che la forza $F$ diventi negativa (cioè il sistema tende a tornare verso il punto fisso). Sui rami crescenti, invece, l'aumento di $V$ aumenta la forza $F$, producendo un'amplificazione positiva (feedback positivo) che destabilizza il punto fisso.
\end{tcolorbox}



\subsection{Regime di Firing Ripetitivo}

\subsubsection{Effetto di una Corrente Costante Sulle Nullcline}

Quando si inietta una \textbf{corrente costante} $I_{\text{ext}} > 0$ nel neurone, l'equazione per $\dot{V}$ diventa:

\begin{equation}
\dot{V} = F(V, w) + \frac{I_{\text{ext}}}{C}
\end{equation}

Questo aggiunge un \textbf{offset positivo} alla forza che governa il potenziale. La condizione di nullcline $\dot{V} = 0$ diventa $F(V, w) = -I_{\text{ext}}/C$, il che equivale a \textbf{traslare verticalmente verso l'alto} la nullcline cubica (in termini di $w$). In altre parole, per bilanciare la corrente eccitatoria aggiuntiva, è necessaria una corrente inibitoria del potassio $w$ maggiore: la nullcline si sposta verso valori più alti di $w$.

\subsubsection{Spostamento del Punto Fisso con l'Intensità di Corrente}

Al crescere dell'intensità della corrente iniettata $I_{\text{ext}}$, il punto fisso si sposta lungo la curva cubica:
\begin{itemize}
    \item Per \textbf{correnti basse}: il punto fisso rimane sul ramo decrescente della curva cubica, a basso $V$ e basso $w$. Il sistema è stabile e risponde con oscillazioni smorzate.
    \item Per \textbf{correnti moderate}: il punto fisso si sposta verso l'alto lungo la curva cubica e, a un certo valore critico di $I_{\text{ext}}$, raggiunge il ramo decrescente vicino al massimo locale. Qui la pendenza è ancora negativa ma diventa meno negativa (e poi zero nel massimo).
    \item Per \textbf{correnti elevate}: il punto fisso si sposta oltre il massimo della curva cubica, sul ramo ascendente destro. La pendenza della nullcline diventa \textbf{positiva} e il punto fisso diventa \textbf{instabile} (fuoco instabile).
\end{itemize}

\subsubsection{Firing Ripetitivo e Ciclo Limite}

Quando il punto fisso diventa instabile (ramo crescente della cubica), il sistema non può più convergere all'equilibrio. Le traiettorie nel piano di fase non si possono avvolgere in spirale attorno al punto fisso (poiché questo è instabile), ma devono necessariamente percorrere \textbf{orbite grandi}. Se il sistema è dissipativo (non conservativo), queste orbite grandi possono stabilizzarsi su un \textbf{ciclo limite} stabile.

Un \textbf{ciclo limite} è una traiettoria chiusa isolata nel piano di fase: il sistema ripete indefinitamente lo stesso percorso, corrispondendo a un'oscillazione periodica. Nel contesto neuronale, questo corrisponde alla \textbf{generazione di potenziali d'azione ripetitivi} (spike train) in risposta a una corrente sostenuta.

\begin{tcolorbox}[colback=gray!5, colframe=gray!40, arc=2mm, boxrule=0.5pt]
\textbf{Nota del docente:} Il passaggio da un punto fisso stabile a un ciclo limite, al variare del parametro di controllo $I_{\text{ext}}$, è un esempio di \textbf{biforcazione}. Il tipo di biforcazione determina molte proprietà della risposta neuronale, in particolare la relazione frequenza-corrente ($f$-$I$). Distingueremo nel dettaglio due tipi principali di biforcazione nei prossimi paragrafi.
\end{tcolorbox}


\section{Biforcazioni}


Una \textbf{biforcazione} è una modificazione qualitativa della struttura degli attrattori di un sistema dinamico al variare di un parametro di controllo. Nel contesto dei modelli neuronali, il parametro di controllo principale è la corrente iniettata $I_{\text{ext}}$. Al variare di $I_{\text{ext}}$, il sistema può passare da un regime di riposo (punto fisso stabile) a un regime di firing (ciclo limite) attraverso una biforcazione.

Il professore precisa che una biforcazione non è semplicemente un cambiamento quantitativo del comportamento, ma un cambiamento \textbf{qualitativo}: prima della biforcazione il sistema ha un attrattore di un certo tipo (ad esempio un punto fisso), dopo la biforcazione l'attrattore cambia natura (ad esempio diventa un ciclo limite). La corrente $I_c$ alla quale avviene la biforcazione è la \textbf{corrente reobase} o \textbf{corrente critica}: è la minima corrente costante che può indurre il firing ripetitivo.

Le proprietà della biforcazione determinano:
\begin{itemize}
    \item Con quale \textbf{minima frequenza} il neurone può sparare
    \item Se la transizione al firing è \textbf{continua o discontinua} (isteresi)
    \item La forma della \textbf{curva frequenza-corrente} ($f$-$I$)
    \item Se il neurone è sensibile a \textbf{perturbazioni veloci} (oscillazioni) o a \textbf{segnali lenti} (variazione graduale dell'input)
\end{itemize}

Nella classificazione di Hodgkin del 1948, le due classi principali di eccitabilità — Tipo I e Tipo II — corrispondono esattamente alle due biforcazioni che analizzeremo: SNIC e Hopf rispettivamente.

\subsubsection{Biforcazione Saddle-Node su Cerchio Invariante (SNIC)}

La \textbf{biforcazione SNIC} (Saddle-Node on Invariant Circle) è caratteristica dei neuroni di \textbf{Tipo I} (o \textbf{classe 1 di eccitabilità} nella classificazione di Hodgkin). Il suo meccanismo è il seguente:

\textbf{Stato di riposo} (corrente bassa):\\
Il piano di fase presenta un \textbf{punto fisso stabile} (nodo stabile o fuoco stabile) e un \textbf{punto sella} nelle vicinanze. Esiste già nel piano di fase un \textbf{cerchio invariante} (una curva chiusa non attraversata dalle traiettorie) che passa vicino al punto sella e circonda il punto fisso stabile.

\textbf{Al valore critico di corrente} $I_{\text{ext}} = I_{\text{SNIC}}$:\\
Il punto fisso stabile e il punto sella si \textbf{avvicinano e si fondono} in un singolo punto di tipo \textbf{saddle-node}. Questo punto si trova esattamente sul cerchio invariante.

\textbf{Dopo la biforcazione} (corrente elevata):\\
Il saddle-node scompare dal piano di fase, lasciando il cerchio invariante come \textbf{ciclo limite stabile} su cui il sistema percorre orbite chiuse. Il sistema inizia a generare potenziali d'azione ripetitivi.

\textbf{Proprietà chiave della biforcazione SNIC:}
\begin{itemize}
    \item Alla biforcazione, il \textbf{periodo del ciclo limite} $T \to \infty$ (la frequenza tende a zero): il neurone inizia a sparare con frequenza arbitrariamente bassa.
    \item La relazione frequenza-corrente ha andamento $f \propto \sqrt{I_{\text{ext}} - I_{\text{SNIC}}}$ per $I_{\text{ext}}$ appena sopra la soglia.
    \item La transizione è \textbf{continua} (non c'è isteresi): il neurone passa gradualmente dal riposo al firing.
\end{itemize}


\subsection{Biforcazione di Hopf }

La \textbf{biforcazione di Hopf} è la transizione caratteristica dei neuroni di \textbf{Tipo II} (classe 2 di Hodgkin). Il meccanismo è profondamente diverso dalla SNIC:

\textbf{Stato di riposo} (corrente bassa):\\
Il piano di fase ha un \textbf{fuoco stabile} come unico attrattore. Le traiettorie convergono a spirale al punto fisso.

\textbf{Al valore critico di corrente} $I_{\text{ext}} = I_H$:\\
La parte reale degli autovalori complessi del fuoco stabile attraversa zero: $\text{Re}(\lambda_\pm) = S(I_H) = 0$. Il fuoco stabile \textbf{perde la sua stabilità} e diventa un fuoco instabile.

\textbf{Dopo la biforcazione} (corrente superiore alla soglia):\\
Il fuoco instabile è circondato da un \textbf{piccolo ciclo limite stabile} che emerge dal punto fisso. L'ampiezza del ciclo limite cresce con $I_{\text{ext}}$ a partire da zero (biforcazione \textbf{supercritica}).

Nella variante \textbf{sottocritica}, il ciclo limite instabile circonda il fuoco stabile prima della biforcazione, e la transizione è discontinua (con isteresi).

\textbf{Proprietà chiave della biforcazione di Hopf:}
\begin{itemize}
    \item Alla biforcazione, la \textbf{frequenza del ciclo limite} $f \to f_0 \neq 0$: il neurone inizia a sparare con frequenza \textbf{finita} non nulla (la frequenza del ciclo limite al punto di biforcazione è $f_0 = \omega/(2\pi) = \sqrt{D - S^2}/(2\pi)$).
    \item La curva $f$-$I$ ha un \textbf{salto}: da frequenza zero (riposo) il sistema passa bruscamente a una frequenza minima $f_0 > 0$ alla soglia.
    \item Nel caso supercritico, la transizione è \textbf{continua} nell'ampiezza ma con frequenza minima non nulla.
\end{itemize}



\subsection{Confronto Tipo I vs. Tipo II: Curva $f$-$I$}


La \textbf{curva $f$-$I$} (frequenza media di firing in funzione della corrente iniettata) è una delle misure sperimentali più informative sul tipo di eccitabilità di un neurone. Essa riflette direttamente il tipo di biforcazione attraverso cui avviene la transizione riposo $\rightarrow$ firing.

\subsubsection{Tipo I (Biforcazione SNIC)}

Nei neuroni di Tipo I:
\begin{itemize}
    \item La soglia di firing corrisponde alla biforcazione SNIC.
    \item Per correnti appena sopra la soglia $I_c$, la frequenza di firing cresce \textbf{continuamente da zero} secondo la legge:
    \begin{equation}
    f \propto \sqrt{I_{\text{ext}} - I_c} \quad \text{per} \quad I_{\text{ext}} \gtrsim I_c
    \end{equation}
    \item La curva $f$-$I$ è una curva che parte dall'origine (nel piano $I_{\text{ext}}$-$f$) con pendenza infinita e poi si appiattisce.
    \item Il neurone è in grado di generare potenziali d'azione con \textbf{frequenza arbitrariamente bassa}: può codificare segnali deboli con bassa frequenza di firing.
\end{itemize}

\subsubsection{Tipo II (Biforcazione di Hopf)}

Nei neuroni di Tipo II:
\begin{itemize}
    \item La soglia di firing corrisponde alla biforcazione di Hopf.
    \item Per correnti appena sopra la soglia $I_H$, il neurone inizia a sparare con una \textbf{frequenza minima non nulla} $f_0$:
    \begin{equation}
    f(I_H^+) = f_0 > 0
    \end{equation}
    \item La curva $f$-$I$ ha un \textbf{salto discontinuo} alla soglia: non esiste una frequenza di firing intermedia tra 0 e $f_0$.
    \item Il neurone è un \textbf{oscillatore a frequenza preferenziale} e risponde in modo tutto-o-niente alla soglia.
    \item Dopo la transizione, la frequenza cresce con la corrente in modo quasi lineare.
\end{itemize}

\subsubsection{La Curva $f$-$I$ come Strumento Diagnostico Sperimentale}

Prima di presentare la tabella comparativa, il professore spiega come la curva $f$-$I$ venga misurata in laboratorio e come permetta di identificare il tipo di eccitabilità:

\begin{enumerate}
    \item Si prepara una registrazione in patch-clamp in configurazione whole-cell su un neurone isolato (o in slice).
    \item Si iniettano \textbf{gradini di corrente} di ampiezza crescente ($I_1 < I_2 < \ldots < I_n$), ciascuno della durata di $0.5$–$1 \; \text{s}$.
    \item Per ogni gradino, si conta il numero di spike emessi e si calcola la frequenza media $f = N_{\text{spike}}/\Delta t$.
    \item Si costruisce il grafico $f$ vs. $I$.
\end{enumerate}

Un neurone di \textbf{Tipo I} mostrerà:
\begin{itemize}
    \item Nessun firing per $I < I_c$
    \item Per $I$ appena sopra $I_c$: firing molto sporadico (1–2 Hz), con lunghi periodi inter-spike
    \item Crescita quasi continua della frequenza con la corrente
\end{itemize}

Un neurone di \textbf{Tipo II} mostrerà:
\begin{itemize}
    \item Nessun firing per $I < I_H$
    \item Per $I$ appena sopra $I_H$: firing brusco a frequenza $f_0$ (tipicamente $20$–$50 \; \text{Hz}$ per molti neuroni corticali)
    \item Curva $f$-$I$ che parte da $f_0$ e poi cresce con andamento quasi lineare
\end{itemize}

\subsubsection{Implicazioni Funzionali}

\begin{table}[htbp]
    \centering
    \renewcommand{\arraystretch}{1.3}
    \begin{tabularx}{\textwidth}{|>{\raggedright\arraybackslash}p{3cm}|X|X|}
    \hline
    \textbf{Proprietà} & \textbf{Tipo I (SNIC)} & \textbf{Tipo II (Hopf)} \\
    \hline
    Biforcazione & Saddle-Node su cerchio invariante & Hopf (supercritica o sottocritica) \\
    \hline
    Frequenza alla soglia & $f \to 0$ (continua) & $f \to f_0 > 0$ (salto) \\
    \hline
    Modulazione frequenza & Ampia gamma, da 0 in su & Gamma limitata attorno a $f_0$ \\
    \hline
    Sensibilità a correnti deboli & Alta (risponde con $f$ bassa) & Bassa (non risponde sotto soglia) \\
    \hline
    Codifica segnale & Tasso di firing (rate coding) & Presenza/assenza dello spike \\
    \hline
    Esempi biologici & Neuroni piramidali corticali (alcune classi) & Neuroni del neurone B del ganglio viscerale \\
    \hline
    \end{tabularx}
\end{table}

\begin{tcolorbox}[colback=gray!5, colframe=gray!40, arc=2mm, boxrule=0.5pt]
\textbf{Nota del docente:} Questa distinzione tra Tipo I e Tipo II non è solo teorica: è misurabile sperimentalmente iniettando rampe di corrente lentamente crescenti e registrando la frequenza di firing. Un neurone di Tipo I inizierà a sparare con frequenza bassa (intorno a pochi Hz), mentre un neurone di Tipo II inizierà bruscamente con una frequenza di almeno 20-40 Hz, senza frequenze intermedie. Questa classificazione è stata proposta dallo stesso Hodgkin nel 1948, prima ancora della pubblicazione del modello matematico del 1952.
\end{tcolorbox}


\section{Dinamica del Modello HH Ridotto con Corrente Costante}

\subsubsection{Tre Regimi al Crescere della Corrente}

Il professore presenta in dettaglio tre casi distinti per il modello HH ridotto (o equivalentemente, FitzHugh-Nagumo) al crescere dell'intensità della corrente costante iniettata:

\textbf{Caso 1 — Corrente bassa} (sotto-soglia):
\begin{itemize}
    \item La nullcline $\dot{V} = 0$ è traslata verso l'alto di poco.
    \item Il nuovo punto fisso si trova ancora sul ramo decrescente della cubica con pendenza $m_V < 0$.
    \item Il punto fisso è un \textbf{fuoco stabile}: le traiettorie convergono a spirale.
    \item Il sistema risponde con una singola oscillazione smorzata senza produzione di spike.
    \item La costante di tempo di rilassamento è determinata dagli autovalori del fuoco stabile.
\end{itemize}

\textbf{Caso 2 — Corrente intermedia} (vicino alla soglia):\\
Il professore descrive un caso particolarmente interessante: la nullcline è salita abbastanza da portare il punto fisso vicino al massimo della curva cubica, dove la pendenza è quasi nulla. La pendenza $m_V$ è ancora lievemente negativa, ma si avvicina a zero. Il punto fisso è ancora un fuoco stabile, ma la separatrice si è spostata vicino al punto fisso.

In questo regime, dal punto fisso stabile si produce una \textbf{prima eccitazione a spike}: la traiettoria percorre la grande orbita (potenziale d'azione), ma poi — nella fase di recupero — il sistema incontra una separatrice ora diversa dall'iniziale e torna al punto fisso con oscillazioni smorzate. Si produce \textbf{un solo spike} seguito da smorzamento.

\textbf{Caso 3 — Corrente elevata} (sopra la biforcazione):
\begin{itemize}
    \item La nullcline è salita ulteriormente portando il punto fisso sul ramo crescente della cubica, con $m_V > 0$.
    \item Il punto fisso diventa \textbf{instabile} (fuoco instabile se siamo dentro la parabola $D > S^2$, oppure nodo instabile).
    \item Il sistema si trova su un \textbf{ciclo limite}: genera potenziali d'azione ripetitivi indefinitamente.
    \item La frequenza di firing dipende dall'intensità di corrente.
\end{itemize}

\subsubsection{Meccanismo del Singolo Spike con Ritorno Smorzato}

Il professore spiega con attenzione il meccanismo del caso intermedio (un singolo spike seguito da smorzamento), che è rappresentativo dell'\textbf{eccitabilità} in senso stretto:

\begin{enumerate}
    \item Il sistema parte dal punto fisso stabile $(V^*, w^*)$.
    \item La corrente aumenta istantaneamente (step): la nullcline si sposta verso l'alto.
    \item Il nuovo punto fisso (ancora stabile) è a un $w$ più alto e un $V$ più alto, ma vicino alla separatrice spostata.
    \item La prima traiettoria verso il nuovo equilibrio attraversa la separatrice e percorre la grande orbita $\rightarrow$ primo spike.
    \item Dopo il primo spike, il sistema rientra nella fase di recupero con $w$ ancora elevato.
    \item In questa condizione, il sistema è sulla \textbf{sinistra della separatrice del nuovo equilibrio}: le successive traiettorie non sono capaci di generare un altro spike.
    \item Il sistema converge a spirale al nuovo punto fisso: smorzamento oscillatorio.
\end{enumerate}

\begin{tcolorbox}[colback=gray!5, colframe=gray!40, arc=2mm, boxrule=0.5pt]
\textbf{Nota del docente:} Questo comportamento — un singolo spike di transizione seguito dall'adattamento a un nuovo equilibrio — è esattamente ciò che si osserva sperimentalmente in molti neuroni corticali in risposta a un gradino di corrente di intensità moderata. Si chiama \textbf{eccitabilità di Tipo I} nel senso fenomenologico: il neurone risponde all'inizio della stimolazione con un burst iniziale e poi si adatta silenziosamente.
\end{tcolorbox}



\section{Il Firing Ripetitivo}

Il professore spiega in dettaglio perché il firing ripetitivo è possibile anche quando il punto fisso è ancora stabile: la chiave è che la \textbf{separatrice si sposta} al variare della corrente iniettata.

Quando la corrente $I_{\text{ext}}$ aumenta, la nullcline di $V$ si sposta verso l'alto. Questo porta il punto fisso a posizioni diverse lungo la curva cubica. Ma anche la separatrice, che è legata alla geometria delle nullcline, si sposta di conseguenza. In particolare, la separatrice può avvicinarsi così tanto al punto fisso stabile che praticamente qualsiasi perturbazione porta il sistema oltre la soglia.

\subsubsection{Meccanismo del Secondo Spike nel Firing Ripetitivo}

Nel regime di ciclo limite, il meccanismo che porta alla generazione del secondo (e successivi) spike è il seguente:

\begin{enumerate}
    \item Dopo il primo spike, il sistema si trova in fase di recupero con $w$ elevato (iperpolarizzazione).
    \item La corrente costante $I_{\text{ext}}$ continua a "spingere" il sistema: compensa la corrente iperpolarizzante del potassio.
    \item Il punto fisso (istantaneo) del sistema è instabile, quindi il sistema non può convergere verso di esso.
    \item Il sistema è costretto a percorrere nuovamente la grande orbita $\rightarrow$ secondo spike.
    \item Il processo si ripete indefinitamente, generando un treno di spike periodico.
\end{enumerate}

\subsubsection{Periodo del Ciclo Limite e Frequenza di Firing}

La \textbf{frequenza di firing} nel regime di ciclo limite è semplicemente $f = 1/T$, dove $T$ è il periodo dell'orbita chiusa. Il periodo dipende dalla geometria del ciclo limite nel piano di fase e dalla velocità con cui il sistema percorre le diverse parti dell'orbita.

In generale, il sistema trascorre la maggior parte del tempo nella fase di recupero (ramo in basso a sinistra del piano di fase), perché in quella regione il campo vettoriale è più debole (la variabile di recupero si rilassa lentamente). La fase di amplificazione esplosiva (rising phase) è invece molto rapida. Questo spiega perché la durata dello spike (circa $1 \; \text{ms}$) sia molto più breve del periodo inter-spike, che può essere di decine di millisecondi.

\begin{tcolorbox}[colback=gray!5, colframe=gray!40, arc=2mm, boxrule=0.5pt]
\textbf{Nota del docente:} La velocità con cui il sistema percorre il ramo lento del ciclo limite è essenzialmente determinata dalla scala temporale $\tau_w$ della variabile di recupero. Questo è il collo di bottiglia che limita la frequenza massima di firing: non si può sparare più velocemente di quanto $w$ riesca a rilassarsi. Una corrente inibitoria del potassio con $\tau_w$ breve consente frequenze di firing più alte.
\end{tcolorbox}

\subsubsection{Visualizzazione nel Piano di Fase}

Nel piano di fase, il \textbf{ciclo limite} appare come una \textbf{curva chiusa stabile}: le traiettorie che iniziano all'interno del ciclo limite convergono verso di esso dall'interno, e quelle che iniziano all'esterno convergono verso di esso dall'esterno. Questo contrasta con il caso del punto fisso stabile (in cui tutte le traiettorie nelle vicinanze convergono al punto).

Il ciclo limite nel piano di fase del modello FHN/HH ridotto ha la forma caratteristica di una grande orbita che:
\begin{itemize}
    \item Si trova sulla destra (alto $V$) corrispondente alla fase di depolarizzazione
    \item Attraversa la nullcline $\dot{V} = 0$ con traiettoria verticale (massimo del potenziale d'azione)
    \item Si porta in alto a sinistra (alto $w$, basso $V$) corrispondente all'iperpolarizzazione
    \item Attraversa la nullcline $\dot{w} = 0$ con traiettoria orizzontale
    \item Ritorna al punto di partenza passando per la fase di recupero
\end{itemize}


\section{Il Campo di Forza Esponenziale: Verso l'Integrate-and-Fire Esponenziale}

\subsubsection{Il Termine Esponenziale nel Campo di Forza}

Il professore introduce un'analisi quantitativa fondamentale: la struttura del \textbf{campo di forza} $F(V, w)$ nella direzione di $V$, al variare di $w$.

Nel modello HH ridotto, il campo di forza per il potenziale è dominato dalla corrente del sodio attraverso il termine $m_\infty(V)$. Come discusso nelle lezioni precedenti, $m_\infty(V)$ ha una forma sigmoidale. Nella regione in cui si sviluppa la \textbf{fase crescente} del potenziale d'azione (la fase di amplificazione esplosiva), la funzione $m_\infty(V)$ è ancora nella sua parte ascendente ripida, e può essere approssimata dalla sua espansione esponenziale:

\begin{equation}
m_\infty(V) \approx e^{(V - V_T)/\Delta_T}
\end{equation}

dove $V_T$ è la \textbf{soglia di sparo} e $\Delta_T$ è la \textbf{pendenza della soglia} (sharpness of threshold).

\subsubsection{Evidenza Numerica del Comportamento Esponenziale}

Il professore mostra un grafico notevole: per sezioni orizzontali fisse del piano di fase (a $w = \text{costante}$), si traccia $F(V, w)$ in funzione di $V$ su scala \textbf{logaritmica} dell'asse $y$.

L'osservazione chiave è:
\begin{itemize}
    \item Nella regione di potenziale attorno alla soglia, la curva $F$ vs. $V$ è approssimativamente una \textbf{retta su scala logaritmica}.
    \item Una retta su scala logaritmica corrisponde a una \textbf{funzione esponenziale} su scala lineare.
    \item Questo vale sia per il modello HH ridotto che per il modello Morris-Lecar: entrambi mostrano un \textbf{comportamento esponenziale} di $F$ nella regione sub-soglia e peri-soglia.
\end{itemize}

Il professore mostra due sezioni specifiche del campo di forza:

\textbf{Sezione a $w = 0.3$:}\\
- Si traccia $F(V, 0.3)$ vs. $V$ su scala semilogaritmica.\\
- Per potenziali compresi tra circa $-65$ e $-40 \; \text{mV}$, il grafico mostra un andamento in prima approssimazione lineare su scala log (cioè esponenziale su scala lineare).\\
- La pendenza della retta semilog corrisponde alla costante $1/\Delta_T$ del termine esponenziale.

\textbf{Sezione a $w = 0.84$:}\\
- Si traccia $F(V, 0.84)$ vs. $V$ su scala semilogaritmica.\\
- La forma è qualitativamente identica: retta semilog nella regione sub-soglia.\\
- Le due curve (a $w = 0.3$ e $w = 0.84$) sono \textbf{quasi coincidenti} nella regione di amplificazione: il termine esponenziale domina e $w$ non modifica significativamente il comportamento.\\
- Una leggera differenza appare solo a potenziali più positivi, dove la corrente inibitoria $w$ non è più trascurabile rispetto all'amplificazione del sodio.

La conclusione del professore è che \textbf{nella regione peri-soglia e nella fase di amplificazione esplosiva}:

\begin{equation}
F(V, w) \approx -G_L(V - E_L) + G_L \Delta_T \exp\!\left(\frac{V - V_T}{\Delta_T}\right) - w
\end{equation}

dove i parametri $\Delta_T$ e $V_T$ sono determinati dalla forma di $m_\infty(V)$ (la funzione di attivazione del sodio nel modello HH ridotto). Questo è esattamente il campo di forza del modello EIF (con $w$ che rappresenta l'inibitore del potassio).

\subsubsection{Struttura del Modello Integrate-and-Fire Esponenziale}

Questa osservazione motiva l'introduzione di un modello monodimensionale ancora più semplice: il \textbf{modello Integrate-and-Fire Esponenziale} (EIF, Exponential Integrate-and-Fire). In questo modello, si mantiene un solo grado di libertà (il potenziale di membrana $V$) e si aggiunge al termine passivo della membrana una non-linearità esponenziale che modella l'amplificazione del sodio:

\begin{equation}
C \frac{dV}{dt} = -G_L(V - E_L) + G_L \Delta_T \exp\!\left(\frac{V - V_T}{\Delta_T}\right) + I_{\text{ext}}
\end{equation}

L'emissione di uno spike viene definita \textbf{convenzionalmente} quando $V$ raggiunge un valore di soglia $V_{\text{th}} \gg V_T$ (tipicamente $0 \; \text{mV}$), dopodiché il potenziale viene resettato a un valore $V_r$ (tipicamente $-65 \; \text{mV}$).

I parametri del modello EIF sono:
\begin{itemize}
    \item $C$: capacità di membrana
    \item $G_L$: conduttanza di leakage (o conduttanza passiva)
    \item $E_L$: potenziale di inversione del leakage (approssimativamente il potenziale di riposo)
    \item $\Delta_T$: sharpness della soglia (tipicamente $1$–$2 \; \text{mV}$)
    \item $V_T$: potenziale di soglia effettivo (tipicamente $-55 \; \text{mV}$)
\end{itemize}

\begin{tcolorbox}[colback=gray!5, colframe=gray!40, arc=2mm, boxrule=0.5pt]
\textbf{Nota del docente:} Il modello EIF è il più semplice modello monodimensionale in grado di riprodurre in modo realistico la forma della \textbf{rampa suprathreshold} del potenziale d'azione. A differenza del semplice integrate-and-fire lineare (LIF), l'EIF non ha bisogno di definire una soglia fissa: lo spike emerge spontaneamente dalla non-linearità esponenziale quando il potenziale supera $V_T$. Il parametro $\Delta_T$ controlla quanto è "sharp" la soglia: per $\Delta_T \to 0$ si recupera il LIF con soglia fissa; per $\Delta_T$ finito la soglia è "morbida" e il rumore può indurre spike anche sotto $V_T$.
\end{tcolorbox}


\chapter{Lezione 9}


\textbf{Argomenti:} neuroni tipo 1 e tipo 2; biforcazione saddle-node sul ciclo limite (SNIC); biforcazione di Hopf subcritica e supercritica; oscillatore a rilassamento; dinamica veloce/lenta ($\varepsilon \ll 1$); riduzione a modello monodimensionale; modello integrate-and-fire (IF); modello exponential integrate-and-fire (EIF); campo di forza e paesaggio energetico; funzione di Lyapunov; modello leaky integrate-and-fire (LIF); prescrizione di reset e condizioni al contorno di Fokker-Planck

\section{Ripasso: Sistemi Eccitabili, Separatrice e Tipi di Neuroni}

\subsubsection{Il Sistema Eccitabile e la Separatrice}

Il professore apre la lezione ricordando i risultati della lezione precedente sui modelli neuronali bidimensionali. I modelli a due variabili di stato — il potenziale di membrana $V$ e la variabile di recupero $W$ — sono \textbf{sistemi eccitabili} perché esiste una \textbf{separatrice} nel piano di fase che divide due regimi qualitativamente distinti:

\begin{itemize}
    \item Se le condizioni iniziali si trovano \textbf{al di sotto} della separatrice, il campo di forza riporta il sistema al punto di equilibrio senza produrre un potenziale d'azione.
    \item Se le condizioni iniziali si trovano \textbf{al di sopra} della separatrice, la dinamica porta all'emissione di un singolo spike, grazie alle correnti del sodio che rendono il campo di forza $F = \dot{V}$ molto più grande della componente $G = \dot{W}$ (flusso quasi orizzontale nel piano di fase).
\end{itemize}

La fase crescente del potenziale d'azione è dominata dalla \textbf{corrente del sodio} (meccanismo di amplificazione non lineare dell'apertura dei canali $\text{Na}^+$): sezionando orizzontalmente il piano di fase a tre diversi valori della variabile di recupero, si ottengono tre profili del campo di forza $F(V)$ tutti con la medesima componente esponenziale, indipendente da $W$. Questa fase di salita è quindi \textbf{determinata esclusivamente} dalla dinamica del sodio.

\subsubsection{Funzione di Guadagno Corrente-Frequenza (Curva f-I)}

La proprietà più importante che differenzia le famiglie di neuroni è la \textbf{curva di guadagno corrente-frequenza} $f(I)$, dove la frequenza di firing $f$ è definita come

\begin{equation}
f = \frac{1}{\langle T_{\text{ISI}} \rangle}
\end{equation}

con $\langle T_{\text{ISI}} \rangle$ l'intervallo interspike medio. La curva $f(I)$ è monotonamente crescente in quanto iniettare più corrente nel neurone porta a una maggiore produzione di spike per unità di tempo. Le famiglie si distinguono per il comportamento vicino alla corrente di soglia $I_\theta$:

\begin{table}[htbp]
    \centering
    \renewcommand{\arraystretch}{1.3}
    \begin{tabularx}{\textwidth}{|>{\raggedright\arraybackslash}p{2cm}|X|>{\raggedright\arraybackslash}p{4.5cm}|}
    \hline
    \textbf{Tipo} & \textbf{Comportamento di $f(I)$ per $I \to I_\theta^+$} & \textbf{Biforcazione} \\
    \hline
    \textbf{Tipo 1} & $f \to 0$ continuamente (con discontinuità nella derivata) & Saddle-node sul ciclo limite (SNIC) \\
    \hline
    \textbf{Tipo 2} & $f$ salta a un valore minimo non nullo (discontinuità) & Hopf subcritica \\
    \hline
    \end{tabularx}
\end{table}

Il \textbf{tipo 1} è tipico dei neuroni piramidali della corteccia cerebrale dei mammiferi; il \textbf{tipo 2} è invece meglio descritto dal modello di Hodgkin-Huxley nella sua parametrizzazione originale.

\section{Neuroni di Tipo 1}

\subsubsection{Il Modello di Morris-Lecar}

Per descrivere i neuroni di tipo 1 il professore introduce il \textbf{modello di Morris-Lecar} (ML), ricordandone la struttura generale dalla lezione 8:

\begin{equation}
C \frac{dV}{dt} = -G_{Ca}\, m_\infty(V)(V - E_{Ca}) - G_K\, W(V - E_K) - G_L(V - E_L) + I_{\text{ext}}
\end{equation}

\begin{equation}
\frac{dW}{dt} = \phi \frac{W_\infty(V) - W}{\tau_W(V)}
\end{equation}

La funzione $m_\infty(V) = \frac{1}{2}[1 + \tanh((V-V_1)/V_2)]$ ha una forma \textbf{sigmoidale}, a differenza della nullcline cubica del modello FitzHugh-Nagumo. Questo rende le nullcline del Morris-Lecar più simili a sigmoidi che a cubiche, e la nullcline $\dot{W}=0$ è a sua volta una curva non lineare (ramo inferiore di una sigmoide), poiché anche $G$ è funzione di una costante di tempo $\tau_W(V) = 1/\cosh((V-V_3)/(2V_4))$.

\subsubsection{Configurazione di Equilibrio e Spostamento della Nullcline}

Per valori di $I_{\text{ext}}$ bassi o nulli, le due nullcline si intersecano in due punti (e talvolta tre):

\begin{itemize}
    \item Il punto a sinistra, dove la pendenza tangente alla nullcline $\dot{V}=0$ è \textbf{negativa}, corrisponde a un \textbf{fuoco stabile} (attrattore).
    \item Il punto successivo, dove la pendenza è \textbf{positiva}, corrisponde a un \textbf{fuoco instabile} o nodo instabile.
\end{itemize}

Quando si inietta una corrente positiva $I_{\text{ext}} > 0$, la nullcline $\dot{V}=0$ si sposta \textbf{verso l'alto}: occorre una variabile di recupero $W$ più grande per annullare $\dot{V}$, perché la corrente eccitatoria aggiuntiva deve essere bilanciata da un maggiore contributo inibitorio. Aumentando progressivamente $I_{\text{ext}}$, si raggiunge una \textbf{corrente di reobasi} $I_\theta$ al di sopra della quale le due nullcline non si intersecano più, e il neurone entra in un regime di firing ripetitivo.

\subsubsection{La Biforcazione Saddle-Node sul Ciclo Limite (SNIC)}

Questo scenario di transizione è noto come \textbf{biforcazione saddle-node sul ciclo limite} (Saddle-Node on Invariant Circle, SNIC), e deve il suo nome alla presenza di una sella e di un nodo nelle vicinanze del quale le traiettorie effettuano un detour molto lento.

Il meccanismo fisico che genera la curva f-I continua e la frequenza di firing tendente a zero è il seguente: quando la corrente iniettata è appena sopra $I_\theta$, le due nullcline si trovano quasi a toccarsi. In questa condizione, la traiettoria che rappresenta il treno di spike si avvicina alla regione di quasi-coalescenza delle nullcline e vi trascorre un tempo \textbf{arbitrariamente lungo}, poiché in quella regione entrambi i componenti del campo di forza $F$ e $G$ sono quasi nulli. L'intervallo interspike $T_{\text{ISI}}$ può quindi tendere a infinito, rendendo la frequenza di firing $f \to 0^+$ in modo continuo.

\begin{tcolorbox}[colback=gray!5, colframe=gray!40, arc=2mm, boxrule=0.5pt]
\textbf{Nota del docente:} In questo regime, il componente principale del campo di forza vicino alla quasi-coalescenza è quello legato alla variabile di recupero: $F \approx 0$ lungo la nullcline, e la traiettoria scende lungo la nullcline $\dot{V}=0$ con velocità dettata da $G$ e quindi da $\tau_W$. Quanto più le nullcline si avvicinano, tanto più lungo è il tempo necessario per compiere l'orbita completa. È possibile selezionare una corrente tale che le due nullcline siano quasi tangenti: la traiettoria avrebbe bisogno di un tempo infinito per attraversare quella regione. Questo giustifica fisicamente la possibilità di avere una frequenza di firing arbitrariamente bassa.
\end{tcolorbox}

\section{Neuroni di Tipo 2}

\subsubsection{Campo di Forza Più Debole nel Modello Morris-Lecar con Parametri HH}

La transizione ai neuroni di tipo 2 avviene quando si modificano i parametri del modello Morris-Lecar in modo da rendere il campo di forza $F = \dot{V}$ \textbf{più debole}. Fisicamente questo corrisponde a rendere la dinamica della salita del potenziale d'azione più lenta, il che equivale ad aumentare la costante di tempo della corrente eccitatoria (legata ai parametri $\alpha_3$ e $\alpha_4$ della sigmoide di Morris-Lecar).

\begin{tcolorbox}[colback=gray!5, colframe=gray!40, arc=2mm, boxrule=0.5pt]
\textbf{Domanda di uno studente:} Perché il professore menziona il modello di Hodgkin-Huxley per il tipo 2? Il modello Morris-Lecar non è usato per entrambi i tipi?

\textbf{Risposta del professore:} Sì, in entrambi i casi si usa il modello di Morris-Lecar. La differenza è solo nella scelta dei parametri. Scegliendo i parametri $\alpha_3$ e $\alpha_4$ in modo da rendere il modello Morris-Lecar più simile al modello Hodgkin-Huxley originale (dove le scale temporali della corrente eccitatoria sono comparabili a quelle della corrente inibitoria), si ottiene il comportamento di tipo 2. In questo senso si dice che la parametrizzazione tipo HH del Morris-Lecar è associata al tipo 2. Negli script condivisi su Google Classroom sono riportati i valori numerici dei parametri per riprodurre entrambi i comportamenti.
\end{tcolorbox}

\subsubsection{Conseguenza Geometrica: Traiettoria al di qua della Nullcline}
Quando $F$ è relativamente piccolo, le frecce orizzontali nel piano di fase sono più corte. Questo implica che, dopo l'emissione dello spike, la traiettoria \textbf{non attraversa} la nullcline $\dot{V}=0$ verso il basso come nel tipo 1. Al contrario, essa rimane \textbf{al di qua} (a destra) della nullcline, procedendo verso valori crescenti di $V$ prima di voltare. La traiettoria aggira la sella rimanendo sul suo lato destro: da qui il nome \textbf{saddle-node fuori dal ciclo limite}.

\subsubsection{Bistabilità}
Una conseguenza importante di questa geometria è la \textbf{bistabilità}: il sistema può avere simultaneamente un ciclo limite stabile (treno di spike) e un punto fisso stabile (potenziale di riposo), anche \textbf{senza corrente iniettata}. In assenza di corrente, le nullcline presentano ancora due intersezioni (il fuoco stabile $E_M$ e la sella). Dipendentemente dalle condizioni iniziali:
\begin{itemize}
    \item Se il potenziale di membrana iniziale è inferiore alla sella (separatrice), il neurone si rilassa verso $E_M$.
    \item Se il potenziale di membrana iniziale è superiore alla sella, il neurone produce un treno di spike stabile (ciclo limite stabile).
\end{itemize}

\subsubsection{Salto Discontinuo nella Curva f-I}
Questa geometria determina la discontinuità nella curva f-I: poiché le traiettorie passano sul lato destro della sella, la forza del campo in quel punto non è mai trascurabile. Anche quando si aumenta la corrente in modo da far quasi coalescere le nullcline, il tempo necessario per compiere un'orbita rimane \textbf{finito} (e non può tendere a infinito). Ne consegue che la frequenza di firing non può tendere a zero: quando il sistema entra in regime oscillatorio, lo fa con una \textbf{frequenza minima non nulla}, producendo il caratteristico salto discontinuo nella curva f-I.

\section{Analisi degli Autovalori: da Nodo Stabile a Biforcazione di Hopf}

\subsubsection{Variazione degli Autovalori con la Corrente}

Per il modello tipo 2 (parametrizzazione FitzHugh-Nagumo o Morris-Lecar con $F$ debole), si può seguire l'evoluzione degli autovalori $\lambda_\pm$ della matrice Jacobiana al punto fisso al variare della corrente iniettata $I/C_m$.

Ricordando il diagramma usato nella lezione precedente, gli autovalori si classificano in base a due quantità:
\begin{itemize}
    \item La \textbf{semitraccia} $\mu = \text{tr}(J)/2$, che determina il segno della parte reale.
    \item Il \textbf{discriminante} $\Delta = \text{tr}(J)^2/4 - \det(J)$, che separa la regione con autovalori reali ($\Delta > 0$) da quella con autovalori complessi coniugati ($\Delta < 0$, oscillazioni).
\end{itemize}

Al crescere di $I$, il punto fisso percorre il seguente percorso nel piano $(\mu, \Delta)$:

\begin{enumerate}
    \item \textbf{Nodo stabile}: $\lambda_\pm$ reali e negativi ($\Delta > 0$, $\mu < 0$). Il neurone si rilassa senza oscillazioni.
    \item \textbf{Fuoco stabile}: $\lambda_\pm$ complessi coniugati con parte reale negativa ($\Delta < 0$, $\mu < 0$). Il neurone si rilassa con oscillazioni smorzate sempre più lente.
    \item \textbf{Crossing dell'asse immaginario} ($\mu = 0$): la parte reale diventa nulla e poi positiva $\rightarrow$ \textbf{biforcazione di Hopf}.
    \item \textbf{Fuoco instabile}: $\lambda_\pm$ complessi coniugati con parte reale positiva ($\Delta < 0$, $\mu > 0$). Il punto fisso è instabile e compare un ciclo limite.
\end{enumerate}

\begin{tcolorbox}[colback=gray!5, colframe=gray!40, arc=2mm, boxrule=0.5pt]
\textbf{Nota del docente:} Nelle vicinanze della biforcazione di Hopf, la parte immaginaria degli autovalori determina la \textbf{frequenza delle oscillazioni} (smorzate prima, del ciclo limite dopo). Poiché nella biforcazione di Hopf (subcritica) la parte immaginaria è già non nulla prima del crossing, il ciclo limite compare con una \textbf{frequenza di onset non nulla}, producendo il salto discontinuo nella curva f-I caratteristico del tipo 2.
\end{tcolorbox}

\subsubsection{Biforcazione di Hopf Subcritica}

La \textbf{biforcazione di Hopf subcritica} si manifesta con un salto discontinuo dell'ampiezza del ciclo limite: al valore critico di corrente $I_\theta$, il punto fisso stabile perde stabilità e \textbf{immediatamente} compare un ciclo limite di ampiezza finita. Nel diagramma biforcazione (ampiezza massima e minima del potenziale vs $I$):
\begin{itemize}
    \item Per $I < I_\theta$: il potenziale di membrana all'equilibrio è un valore stabile costante (punto fisso).
    \item Appena $I > I_\theta$: compare un'orbita con ampiezza finita, rappresentata da un valore massimo e un valore minimo ben separati.
\end{itemize}

Il comportamento è dunque \textbf{discontinuo}: questa è la firma della Hopf subcritica ed è il meccanismo alla base dei neuroni di tipo 2.

\subsubsection{Biforcazione di Hopf Supercritica}

In alternativa, se la nullcline $\dot{V}=0$ è \textbf{meno ripida} (il parametro $m_V$ della pendenza è maggiore di zero ma non molto grande, ossia il sistema è meno non lineare), si osserva una \textbf{biforcazione di Hopf supercritica}. In questo caso:
\begin{itemize}
    \item La parte reale degli autovalori diventa positiva attraversando il crossing.
    \item La parte immaginaria rimane diversa da zero prima e dopo.
    \item L'ampiezza del ciclo limite cresce in modo \textbf{continuo e dolce} dalla biforcazione.
\end{itemize}

In questa condizione il neurone \textbf{non produce spike}: genera oscillazioni sottosoglia sostenute di ampiezza crescente con la corrente. Questi neuroni si comportano come \textbf{pacemaker di fase} e sono rilevanti in strutture subcorticali, dove un sottoinsieme di neuroni può produrre oscillazioni in grado di sincronizzare l'attività della rete.

\begin{tcolorbox}[colback=gray!5, colframe=gray!40, arc=2mm, boxrule=0.5pt]
\textbf{Nota del docente:} Il punto chiave per la Hopf supercritica è che il raggio del ciclo limite (distanza tra massimo e minimo del potenziale) cresce in modo dolce con la corrente iniettata. Ma sotto queste condizioni il neurone non produce spike: genera solo oscillazioni di ampiezza sub-soglia.
\end{tcolorbox}


\subsubsection{Oscillazioni Sottosoglia come Pacemaker}

I neuroni che sperimentano una biforcazione di Hopf supercritica producono oscillazioni di membrana sottosoglia di ampiezza crescente al crescere della corrente iniettata, senza mai emettere spike. Il professore evidenzia che questo tipo di dinamica è biologicamente rilevante: in strutture \textbf{subcorticali}, alcuni neuroni agiscono come pacemaker producendo oscillazioni ritmiche che inducono una sincronizzazione a livello di popolazione (activity lock in the network). L'ampiezza dell'oscillazione, che cresce dolcemente con $I$, rappresenta una modalità di codifica dell'informazione alternativa al firing.



\subsection{Oscillatore a Rilassamento: Separazione delle Scale Temporali}

\subsubsection{La Condizione $\varepsilon \ll 1$}

Si considera ora una condizione fondamentale che riguarda la separazione delle scale temporali delle due variabili di stato. La costante di tempo del potenziale di membrana è dell'ordine

\begin{equation}
\tau_V \sim 10\text{–}20 \; \text{ms}
\end{equation}

mentre i canali del potassio (corrispondenti alle variabili $N$ e $H$ del modello HH) hanno una cinetica più lenta:

\begin{equation}
\tau_W \sim 150 \; \text{ms}
\end{equation}

Il loro rapporto

\begin{equation}
\varepsilon = \frac{\tau_V}{\tau_W} \ll 1
\end{equation}

è molto minore di 1. Questo parametro $\varepsilon$ è lo stesso che era stato introdotto nelle lezioni precedenti come fattore moltiplicativo davanti al campo di forza $G$ della variabile di recupero:

\begin{equation}
\dot{W} = \varepsilon \cdot G(V, W)
\end{equation}

Poiché $\varepsilon$ è piccolo, il campo di forza verticale $G$ è \textbf{debole} rispetto alla componente orizzontale $F$. Nel piano di fase, il flusso è quindi rappresentato da \textbf{linee quasi orizzontali}.

\subsubsection{Decomposizione del Potenziale d'Azione in Componenti Veloci e Lente}

In questo regime, il potenziale d'azione si decompone in quattro fasi ben distinte, ciascuna con la propria scala temporale:

\begin{enumerate}
    \item \textbf{Fase di salita veloce} (scala $\tau_V$): il neurone si avvicina rapidamente alla nullcline $\dot{V}=0$ da sinistra. Il campo di forza $F$ è positivo e grande: le frecce orizzontali sono lunghe e puntano verso destra.
    \item \textbf{Fase lungo la nullcline} (scala $\tau_W$): una volta raggiunta la nullcline $\dot{V}=0$, la traiettoria la percorre lentamente verso l'alto (verso valori crescenti di $W$). La velocità in questa fase è determinata da $G$ e quindi da $\tau_W$.
    \item \textbf{Fase di discesa veloce} (scala $\tau_V$): quando la variabile di recupero diventa sufficientemente grande da non avere più la nullcline $\dot{V}=0$ disponibile, il campo di forza $F$ è negativo e grande: il sistema collassa rapidamente verso il ramo sinistro della nullcline, producendo la \textbf{fase iperpolarizzante} del potenziale d'azione.
    \item \textbf{Periodo interspike} (scala $\tau_W$): dopo l'iperpolarizzazione, la variabile di recupero diminuisce lentamente (scala $\tau_W$) fino a che la traiettoria non raggiunge il ramo destro della nullcline, producendo il successivo spike.
\end{enumerate}

Questo meccanismo è denominato \textbf{oscillatore a rilassamento} poiché la dinamica si articola in due fasi di rilassamento lente (sulle scale $\tau_W$) intervallate da due transizioni veloci (sulle scale $\tau_V$).

\begin{tcolorbox}[colback=gray!5, colframe=gray!40, arc=2mm, boxrule=0.5pt]
\textbf{Nota del docente:} Il termine "oscillatore a rilassamento" evidenzia che l'intervallo interspike — e quindi la frequenza di firing — è determinata principalmente dalla scala temporale lenta $\tau_W$, ma dipende anche dalla corrente iniettata. Iniettare più corrente accelera la fase di rilassamento lenta, aumentando la frequenza di firing.
\end{tcolorbox}

\subsubsection{Meccanismo a Scalini lungo la Nullcline}

Il professore descrive in dettaglio il meccanismo di progressione lenta lungo la nullcline $\dot{V}=0$. Immaginando di integrare numericamente il sistema con passi discreti:

\begin{itemize}
    \item Quando la traiettoria si trova esattamente sulla nullcline $\dot{V}=0$, la componente orizzontale del campo di forza è nulla: il passo successivo è \textbf{puramente verticale} (determinato da $G$).
    \item Ma dopo il passo verticale, la traiettoria si trova leggermente fuori dalla nullcline: $F \neq 0$ e, poiché $|F|$ è grande, la traiettoria viene \textbf{rapidamente riportata} sulla nullcline (passo quasi orizzontale).
    \item Si ottiene così un percorso a scalini (staircase): un passo verticale piccolo (scala $\tau_W$) seguito da un ritorno quasi istantaneo alla nullcline (scala $\tau_V$).
\end{itemize}

Questo meccanismo si ripete fino a quando la variabile di recupero supera il valore per cui non esiste più alcun punto della nullcline $\dot{V}=0$ con quel valore di $W$: in quel momento il campo di forza $F$ diventa positivo senza possibilità di ritorno alla nullcline, e si innesca la fase di discesa veloce.

\section{Riduzione a una Dimensione: verso il Modello Integrate-and-Fire}

\subsubsection{Regime Sottosoglia con $W \approx W^*$ Costante}

Un esperimento concettuale permette di apprezzare la possibilità di ridurre il sistema da due a una sola dimensione. Si consideri il neurone in una configurazione in cui:
\begin{itemize}
    \item Il potenziale di membrana è \textbf{iperpolarizzato} (ben al di sotto della separatrice).
    \item Non si inietta alcuna corrente, oppure la corrente è insufficiente a produrre spike.
\end{itemize}

In questa condizione, le traiettorie nel piano di fase sono quasi orizzontali: la variabile di recupero $W$ cambia pochissimo durante il moto del potenziale. Si può quindi approssimare $W \approx W^*$, dove $W^*$ è il valore di equilibrio, e l'equazione per $V$ diventa:

\begin{equation}
C_m \frac{dV}{dt} \approx -G_L(V - E_M)
\end{equation}

ovvero una semplice equazione di rilassamento esponenziale verso il potenziale di equilibrio $E_M$. Questa è la dinamica delle \textbf{proprietà elettriche passive} del neurone: la membrana si comporta come un circuito $RC$ con costante di tempo $\tau_m = C_m/G_m$.

\begin{tcolorbox}[colback=gray!5, colframe=gray!40, arc=2mm, boxrule=0.5pt]
\textbf{Nota del docente:} In questo regime possiamo dimenticarci di $W$. La dinamica del potenziale di membrana è ben descritta da quella associata alle proprietà elettro-passive del neurone: la perdita passiva delle correnti (potassio, cloro, sodio in regime di equilibrio), senza alcun contributo della corrente escitatoria del sodio responsabile dello spike.
\end{tcolorbox}

\subsubsection{Risposta a un Impulso di Corrente: la Separatrice come Soglia}

Si inietta un \textbf{impulso di corrente} (corrente di breve durata): questo provoca un salto rapido nel potenziale di membrana, senza modificare sensibilmente la variabile di recupero $W$ (poiché la dinamica di $W$ è lenta).

Due scenari si presentano:
\begin{enumerate}
    \item \textbf{Impulso sub-soglia} ($V$ rimane al di sotto della separatrice): il campo di forza riporta $V$ verso $E_M$ con rilassamento esponenziale di costante di tempo $\tau_m = C_m/G_m$.
    \item \textbf{Impulso sopra-soglia} ($V$ supera la separatrice): la componente esponenziale del sodio domina il campo di forza, producendo l'emissione dello spike.
\end{enumerate}

Questo esperimento dimostra che nel regime sottosoglia possiamo trascurare $W$ e descrivere la dinamica di $V$ con un modello monodimensionale.

\subsubsection{Prescrizione dello Spike: Riduzione a un Grado di Libertà}

La riduzione dimensionale si completa con una \textbf{prescrizione esplicita} per la produzione dello spike: si rimuove il grado di libertà della variabile di recupero e si sostituisce la sua dinamica con una regola semplice. La dinamica dell'oscillatore a rilassamento mostra che, dopo lo spike, la traiettoria percorre la fase di iperpolarizzazione e rientra nella regione dove $W \approx W^*$. In questa regione, il potenziale di membrana si trova a un valore di \textbf{reset} $V_{\text{reset}}$ (tipicamente inferiore a $E_M$).

La prescrizione formale è:

\begin{equation}
\text{se } V(t_k) > V_{\text{soglia}} \Rightarrow \text{spike a } t_k, \quad V(t_k + \tau_0) = V_{\text{reset}}
\end{equation}

dove $\tau_0$ è il \textbf{periodo refrattario assoluto} (la durata dello spike e dell'iperpolarizzazione che segue).

Questa prescrizione definisce la classe dei \textbf{modelli integrate-and-fire (IF)}.

\section{Il Modello Integrate-and-Fire (IF)}

\subsubsection{Equazione Dinamica Sottosoglia}

Il modello integrate-and-fire nella sua formulazione generale è un sistema \textbf{monodimensionale} con la seguente struttura:

\begin{equation}
\tau_m \frac{dV}{dt} = \mathcal{F}(V) + R_m I(t)
\end{equation}

dove:
\begin{itemize}
    \item $\tau_m = C_m / G_m$ è la costante di tempo di membrana (tipicamente 10–20 ms per i neuroni corticali).
    \item $\mathcal{F}(V)$ è il campo di forza sottosoglia (dipende dalla variante del modello).
    \item $R_m = 1/G_m$ è la resistenza di membrana.
    \item $I(t)$ è la corrente iniettata.
\end{itemize}

La \textbf{prescrizione di spike} è:
\begin{itemize}
    \item Quando $V(t) \geq V_{\text{soglia}}$: si registra un'emissione a tempo $t_k$.
    \item Dopo un periodo refrattario $\tau_0$: $V(t_k + \tau_0) = V_{\text{reset}}$.
\end{itemize}

Il nome "integrate-and-fire" deriva dal fatto che:
\begin{itemize}
    \item \textbf{Integrate}: nel regime sottosoglia il potenziale di membrana integra la corrente iniettata (è un \textbf{integratore} del segnale di ingresso).
    \item \textbf{Fire}: quando il potenziale supera la soglia, il neurone emette uno spike.
\end{itemize}

\subsubsection{Il Potenziale di Reset}

Il professore risponde a una domanda degli studenti sul significato del potenziale di reset:

\begin{tcolorbox}[colback=gray!5, colframe=gray!40, arc=2mm, boxrule=0.5pt]
\textbf{Domanda di uno studente:} Dopo lo spike, perché il neurone riprende la dinamica da $V_{\text{reset}}$ e non dal potenziale di equilibrio $E_M$?

\textbf{Risposta del professore:} Il reset non è il potenziale di equilibrio, ma il potenziale che il neurone ha quando rientra nella regione del piano di fase dove $W \approx W^*$, ossia quando la variabile di recupero ha raggiunto il suo valore di quasi-stazionarietà. In questa regione, il potenziale di membrana è tipicamente \textbf{più negativo} di $E_M$ (iperpolarizzato), perché la fase di iperpolarizzazione porta il sistema al di sotto dell'equilibrio. Quindi $V_{\text{reset}} < E_M$. È il prezzo che si paga per dimenticare la variabile di recupero: si introduce una condizione iniziale ben definita con cui il modello monodimensionale riprende la sua evoluzione dopo ogni spike.
\end{tcolorbox}

\section{Il Campo di Forza dell'EIF: Parametri $V_T$ e $\Delta_T$}

\subsubsection{Struttura del Campo di Forza dal Piano di Fase}

Effettuando tagli orizzontali nel piano di fase del modello Morris-Lecar (o del modello HH ridotto) a diversi valori di $W^*$, si ottiene la forma del campo di forza $F(V, W^*)$ come funzione di $V$ a $W$ fissato. Il risultato mostra (in scala logaritmica per la componente positiva):
\begin{itemize}
    \item Una regione centrale con $F < 0$: il campo di forza è negativo, il potenziale è attratto verso l'equilibrio stabile (a sinistra).
    \item Un valore di $V$ dove $F = 0$: l'equilibrio instabile (a destra), corrispondente alla soglia di spike.
    \item Una regione ancora più a destra con $F > 0$, che esplode esponenzialmente: la corrente del sodio domina e si innesca lo spike.
\end{itemize}

Questa struttura si può scomporre come \textbf{sovrapposizione} di due componenti:
\begin{enumerate}
    \item Una componente \textbf{lineare} (leakage): $F_{\text{lin}}(V) = -(V - E_M)/\tau_m$, che descrive il regime passivo e crea l'equilibrio stabile a $E_M$.
    \item Una componente \textbf{esponenziale} (sodio): $F_{\text{exp}}(V) = \frac{\Delta_T}{\tau_m} \exp\!\left(\frac{V - V_T}{\Delta_T}\right)$, che è responsabile dell'esplosione del campo di forza vicino alla soglia.
\end{enumerate}

\subsubsection{Il Modello Exponential Integrate-and-Fire (EIF)}

L'equazione dinamica del modello \textbf{Exponential Integrate-and-Fire (EIF)} è:

\begin{equation}
C_m \frac{dV}{dt} = -G_L(V - E_L) + G_L \Delta_T \exp\!\left(\frac{V - V_T}{\Delta_T}\right) + I
\end{equation}

dove:
\begin{itemize}
    \item $E_L$ è il potenziale di riposo (leakage).
    \item $V_T$ è la \textbf{soglia di iniziazione dello spike} (spike initiation threshold): il valore di potenziale dove il campo di forza $F$ ha il suo \textbf{minimo} (il "saddle" del campo di forza monodimensionale).
    \item $\Delta_T$ è la \textbf{sharpness del threshold} (threshold sharpness): la finestra di tensione entro cui si produce il blow-up esponenziale.
\end{itemize}

Il significato fisico di $\Delta_T$:
\begin{itemize}
    \item Se $V - V_T \gg \Delta_T$: il termine esponenziale è enorme $\rightarrow$ spike inevitabile.
    \item Se $V - V_T \ll -\Delta_T$: il termine esponenziale è trascurabile $\rightarrow$ dinamica puramente passiva.
    \item $\Delta_T$ piccolo: la soglia è "netta" (il blow-up avviene in un intervallo di tensione stretto).
    \item $\Delta_T$ grande: la soglia è "morbida" (il blow-up si distribuisce su un intervallo più ampio).
\end{itemize}

\begin{tcolorbox}[colback=gray!5, colframe=gray!40, arc=2mm, boxrule=0.5pt]
\textbf{Nota del docente:} $V_T$ corrisponde a dove si trova il minimo del campo di forza nel diagramma lineare. In questa regione si ha l'iniziazione dello spike. $\Delta_T$ è la sharpness: se $V - V_T$ è molto maggiore di $\Delta_T$, l'esponenziale esplode. Se $\Delta_T$ è piccolo, il blow-up avviene in una piccola finestra di tensione, rendendola più riproducibile.
\end{tcolorbox}


\section{Corrente di Reobase e Soglia Dipendente dal Protocollo}

\subsubsection{Il Campo di Forza Totale Include la Corrente}

Il professore precisa una domanda emersa durante la lezione: il campo di forza $F$ che si visualizza nel piano di fase è il campo di forza \textbf{totale}, che include già la corrente iniettata divisa per la capacità di membrana:

\begin{equation}
F_{\text{tot}}(V) = F_{\text{intrinseche}}(V, W^*) + \frac{I}{C_m}
\end{equation}

Quando si inietta una corrente positiva $I > 0$, il campo di forza totale si sposta \textbf{verso l'alto} di un valore costante $I/C_m$. Questo sposta il profilo del campo di forza rispetto alla posizione dello zero, avvicinando le radici del campo di forza.

\subsubsection{Due Protocolli di Stimolazione, Due Soglie}

La \textbf{soglia di voltaggio per l'emissione dello spike} dipende dal \textbf{protocollo di stimolazione} adottato:

\textbf{Protocollo 1: impulso di corrente}
\begin{itemize}
    \item La corrente impulsiva cambia rapidamente $V$ senza modificare le nullcline (che dipendono da $I$ costante).
    \item La soglia è il potenziale di sella del piano di fase originale: $V_{\text{soglia, impulso}} = V_{\text{reobase}}$.
    \item Il neurone deve essere spinto fisicamente al di sopra di questa sella dall'impulso.
\end{itemize}

\textbf{Protocollo 2: corrente costante}
\begin{itemize}
    \item La corrente costante sposta il campo di forza totale verso l'alto di $I/C_m$.
    \item La corrente di reobase $I_{\text{reo}}$ è il valore minimo di corrente continua per cui il campo di forza totale è sempre positivo (non ha più radici).
    \item La soglia effettiva di voltaggio in questo caso è $V_T$ (dove il campo di forza ha il minimo), che è \textbf{inferiore} a $V_{\text{reobase}}: V_T < V_{\text{reobase}}$.
\end{itemize}

La corrente di reobase si calcola imponendo che il minimo del campo di forza totale sia zero:

\begin{equation}
I_{\text{reo}} \approx G_L \cdot \Delta_T
\end{equation}

(l'espressione esatta si ricava derivando $F_{\text{tot}}$ e imponendo $F_{\text{tot}} = 0$ e $\partial F_{\text{tot}}/\partial V = 0$ contemporaneamente).

\subsubsection{Analogia del Livello del Mare}

Il professore illustra questo meccanismo con una metafora geometrica: si immagini il campo di forza $F(V)$ come il profilo di una costa con una montagna (il termine esponenziale). Iniettare una corrente positiva equivale ad \textbf{abbassare il livello del mare}. Abbassando il livello del mare:
\begin{itemize}
    \item La "costa" (il valore dove $F_{\text{tot}} = 0$) si sposta verso destra, avvicinandosi alla montagna.
    \item Abbassando ancora il livello del mare, si raggiunge la situazione in cui la costa tocca la base della montagna: questa è la corrente di reobase.
    \item Abbassando ulteriormente, non c'è più costa: il campo di forza è sempre positivo e lo spike è inevitabile.
\end{itemize}



\section{Funzione di Lyapunov e Paesaggio Energetico}

\subsubsection{Dal Campo di Forza all'Energia}

Poiché nel regime sottosoglia la dinamica di $V$ è monodimensionale e il campo di forza $F(V, W^*)$ dipende solo da $V$, esso è integrabile e si può definire una \textbf{funzione di energia totale immagazzinata} (analogia con l'energia potenziale):

\begin{equation}
U(V) = -\int F(V, W^*) \, dV
\end{equation}

La variazione temporale di $U$ è:

\begin{equation}
\frac{dU}{dt} = \frac{\partial U}{\partial V} \dot{V} = -F(V) \cdot \dot{V} = -\dot{V}^2 \leq 0
\end{equation}

Questa quantità è sempre \textbf{non positiva}: $U$ è una funzione monotonamente decrescente nel tempo (eccetto agli equilibri dove $\dot{V} = 0$). Questa proprietà identifica $U$ come una \textbf{funzione di Lyapunov} per il sistema.

\subsubsection{La Funzione di Lyapunov e la Stabilità}

Una funzione di Lyapunov $U(V)$ per un sistema dissipativo ha le seguenti proprietà:
\begin{itemize}
    \item $\dot{U} \leq 0$ per ogni traiettoria del sistema.
    \item $\dot{U} = 0$ se e solo se $\dot{V} = 0$, cioè ai punti di equilibrio.
\end{itemize}

Questo significa che il sistema si muove \textbf{sempre in discesa} lungo il paesaggio energetico definito da $U(V)$: la dinamica equivale a quella di una pallina che rotola giù per una superficie. Gli \textbf{attrattori} corrispondono a \textbf{avvallamenti} (minimi locali di $U$), mentre le \textbf{selle} (punti instabili) corrispondono a \textbf{selle del paesaggio energetico}.

\subsubsection{Paesaggio Energetico del Modello EIF}

Per il modello EIF, il paesaggio energetico $U(V)$ ha la seguente forma:
\begin{itemize}
    \item Un \textbf{avvallamento} (minimo stabile) in corrispondenza di $E_M$: questo è il potenziale di riposo (attrattore).
    \item Una \textbf{sella} in corrispondenza di $V_T$ (circa): questo è il punto di massimo locale di $U$, che corrisponde alla soglia di spike nel paesaggio energetico.
    \item Una \textbf{discesa ripida} (minimo profondo) per $V > V_T$: una volta superata la sella, il sistema cade in una buca praticamente irreversibile, emettendo lo spike.
\end{itemize}

\begin{tcolorbox}[colback=gray!5, colframe=gray!40, arc=2mm, boxrule=0.5pt]
\textbf{Nota del docente:} Il paesaggio energetico rende intuitivo perché lo spike sia difficile da "annullare" una volta innescato: per fermare un neurone che si trova oltre la sella $V_T$, occorrerebbe iniettare una corrente inibitoria molto intensa (che alzi il paesaggio energetico al di là della buca), mentre per innescarlo basta superare la sella. Questa asimmetria è una proprietà fondamentale dell'eccitabilità neuronale.
\end{tcolorbox}

\subsubsection{Effetto della Corrente sul Paesaggio Energetico}

L'iniezione di una corrente costante $I$ sposta il paesaggio energetico abbassando il contributo lineare. Il professore descrive tre regimi:
\begin{itemize}
    \item \textbf{$I = 0$}: un unico attrattore a $E_M$, sella a $V_{\text{reobase}}$, buca profonda per $V > V_{\text{reobase}}$.
    \item \textbf{$0 < I < I_{\text{reo}}$}: l'avvallamento a $E_M$ si sposta verso destra (potenziale di equilibrio $E_M + I/G_L$) e la sella si avvicina: il sistema è metastabile.
    \item \textbf{$I \to I_{\text{reo}}$}: la sella e l'avvallamento si avvicinano finché la sella scompare (coalescenza saddle-node nel paesaggio energetico).
    \item \textbf{$I > I_{\text{reo}}$}: il paesaggio non ha più minimi locali per $V$ finito: lo spike è inevitabile.
\end{itemize}

\subsubsection{Barriera Assorbente e Limite $\Delta_T \to 0$}

Il professore illustra come la variazione del parametro $\Delta_T$ modifichi la forma della barriera nel paesaggio energetico:
\begin{itemize}
    \item \textbf{$\Delta_T$ standard (1.5 mV)}: la sella nel paesaggio energetico è "morbida"; c'è una zona di transizione finita tra il potenziale di riposo e la buca dello spike.
    \item \textbf{$\Delta_T$ ridotto}: la sella diventa più ripida; la barriera assorbente è più netta.
    \item \textbf{$\Delta_T \to 0$}: la componente esponenziale scompare ($G_L \Delta_T \exp((V - V_T)/\Delta_T) \to 0$ per $V < V_T$ e $\to \infty$ per $V > V_T$). La sella diventa una \textbf{barriera assorbente verticale}: tutti i neuroni che raggiungono $V_T$ vengono "assorbiti" verso $+\infty$, ovvero emettono lo spike con certezza.
\end{itemize}

\subsubsection{Il Modello Leaky Integrate-and-Fire (LIF)}

Nel limite $\Delta_T \to 0$, il modello EIF si riduce al \textbf{modello Leaky Integrate-and-Fire (LIF)}, introdotto da \textbf{Louis Lapicque} (matematico francese) all'inizio del XX secolo:

\begin{equation}
\tau_m \frac{dV}{dt} = -(V - E_L) + R_m I(t)
\end{equation}

con la prescrizione:
\begin{itemize}
    \item $V(t) \geq V_{\text{th}} \Rightarrow$ spike a $t$, poi $V \to V_{\text{reset}}$ dopo un periodo refrattario $\tau_0$.
\end{itemize}

Il campo di forza è puramente parabolico (lineare in $V$), con un unico avvallamento in $E_L$ e una \textbf{barriera dura} a $V_{\text{th}}$ (nessuna dinamica intrinseca della soglia: la soglia è una soglia fissa, non una proprietà della dinamica del campo di forza).

Il paesaggio energetico del LIF è quindi:
\begin{equation}
U_{\text{LIF}}(V) = \frac{(V - E_L)^2}{2}
\end{equation}
(parabola) con una barriera assorbente a $V_{\text{th}}$.

\begin{tcolorbox}[colback=gray!5, colframe=gray!40, arc=2mm, boxrule=0.5pt]
\textbf{Nota del docente:} Il LIF è il modello integrate-and-fire più semplice storicamente (non il più semplice in assoluto, poiché ce n'è uno lineare, ma il più usato). È stato inventato da Lapicque circa cento anni fa e permette di evitare la complessità della componente esponenziale mantenendo la proprietà essenziale di integrare il segnale e sparare a soglia.
\end{tcolorbox}


\section{Prescrizione di Reset e Condizioni al Contorno per la Fokker-Planck}

\subsubsection{Il Potenziale di Reset}

La prescrizione di reset completa il modello integrate-and-fire: dopo l'emissione di uno spike al tempo $t_k$, la variabile di membrana riparte da $V_{\text{reset}}$ dopo il periodo refrattario $\tau_0$:

\begin{equation}
V(t_k + \tau_0) = V_{\text{reset}}
\end{equation}

Il potenziale di reset ha un significato preciso in relazione alla dinamica bidimensionale originale: è il valore del potenziale di membrana quando la traiettoria rientra nella regione del piano di fase dove $W \approx W^*$. Tipicamente $V_{\text{reset}} < E_M$ (la fase iperpolarizzante porta il potenziale al di sotto dell'equilibrio di riposo).

\begin{tcolorbox}[colback=gray!5, colframe=gray!40, arc=2mm, boxrule=0.5pt]
\textbf{Nota del docente:} Diciamo che questo è il reset. È il potenziale di membrana che si ha quando la traiettoria rientra nell'intervallo di $W$ dove possiamo approssimare la variabile di recupero come costante e quindi considerare il neurone come descritto da un sistema Hamiltoniano monodimensionale dissipativo. Tipicamente questo valore è inferiore al potenziale di equilibrio $E_M$, quindi $V_{\text{reset}} < E_M$.
\end{tcolorbox}

\subsubsection{Interpretazione come Condizioni al Contorno di Fokker-Planck}

Il professore accenna alla rilevanza della prescrizione di reset per la descrizione di \textbf{popolazioni di neuroni}. Quando si descrive l'evoluzione di una distribuzione di probabilità $P(V, t)$ per una popolazione di neuroni IF soggetti a rumore (equazione di Fokker-Planck), la prescrizione di reset si traduce in condizioni al contorno speciali:

\begin{itemize}
    \item \textbf{Pozzo alla soglia $V_{\text{th}}$}: i neuroni che raggiungono la soglia vengono "assorbiti" (sparano lo spike e scompaiono dall'integrazione). Il flusso di neuroni attraverso la soglia è una sorgente: il \textbf{tasso di firing} della popolazione.
    \item \textbf{Sorgente a $V_{\text{reset}}$}: i neuroni che hanno appena sparato vengono "reiniettati" nel sistema al potenziale di reset con il tasso di firing (come delta di Dirac nel flusso di probabilità).
\end{itemize}

Formalmente:
\begin{equation}
J(V_{\text{th}}, t) = -r(t) \quad \text{(pozzo)}
\end{equation}
\begin{equation}
J(V_{\text{reset}}^+, t) - J(V_{\text{reset}}^-, t) = r(t) \quad \text{(sorgente)}
\end{equation}

dove $J(V, t)$ è la corrente di probabilità e $r(t)$ è il tasso di firing istantaneo della popolazione. Questa struttura è fondamentale per lo studio delle reti di neuroni IF con rumore, che sarà argomento di lezioni future.

\chapter{Lezione 10}


\textbf{Argomenti:} Ripasso EIF e paesaggio energetico; validazione sperimentale EIF (esperimento Gerstner–Markram); estrazione del campo di forza con media condizionata; fit dei parametri EIF su neuroni corticali; predizione del timing degli spike; periodo refrattario relativo e dinamica di $V_T$; modello AdEx (EIF adattivo con recupero); modello QIF e oscillatori di Kuramoto; modello LIF ($\Delta_T \to 0$); tassonomia dei neuroni inibitori (NAC, AC, burst, delay); modello adattivo integrate-and-fire 2D; firing tonico non-adattivo; spike frequency adaptation; burst iniziale + tonico; firing a burst regolare; differenze LIF vs EIF per il burst
\section{Ripasso: Modello EIF e Paesaggio Energetico}

\subsubsection{Campo di Forza dell'EIF}

Il professore apre la lezione con un ripasso della dinamica del modello \textbf{Exponential Integrate-and-Fire} (EIF) discussa nella lezione precedente. Il modello descrive l'evoluzione temporale del potenziale di membrana $V$ tramite il campo di forza monodimensionale:

\begin{equation}
C_m \frac{dV}{dt} = \mathcal{F}(V) + I_{\text{ext}}
\end{equation}

dove il campo di forza è:

\begin{equation}
\mathcal{F}(V) = -G_m(V - E_m) + G_m \Delta_T \exp\!\left(\frac{V - V_T}{\Delta_T}\right)
\end{equation}

Il primo termine è la \textbf{componente passiva lineare} (l'RC circuit che porta il potenziale verso $E_m$); il secondo termine è la \textbf{componente esponenziale} attiva soltanto quando $V$ si avvicina alla soglia $V_T$, con $\Delta_T$ (la sharpness) che determina quanto è appuntita la rampa esponenziale.

\subsubsection{Funzione di Lyapunov e Metafora del Paesaggio Energetico}

Per avere un'intuizione fisica dell'evoluzione del potenziale di membrana, si introduce la \textbf{funzione di Lyapunov} (energia del sistema):

\begin{equation}
U(V) = -\int \mathcal{F}(V)\, dV
\end{equation}

La condizione $dU/dt = -\mathcal{F}(V)^2 \leq 0$ garantisce che il sistema dissipa energia finché non raggiunge un attrattore. Il paesaggio energetico $U(V)$ ha la forma di una \textbf{valle profonda} centrata intorno a $E_m$ (stato di riposo preferenziale) e di una \textbf{sella} in corrispondenza di $V_T$ (l'emissione dello spike corrisponde al superamento della sella con caduta verso $+\infty$).

Quando si inietta una corrente esterna $I_{\text{ext}} > 0$, il paesaggio energetico si \textbf{deforma}: la valle si sposta verso destra (avvicinandosi alla soglia) e la sella si abbassa. Più elevata è la corrente, più il potenziale di equilibrio è vicino alla soglia di emissione.

\begin{tcolorbox}[colback=gray!5, colframe=gray!40, arc=2mm, boxrule=0.5pt]
\textbf{Chiarimento del docente in risposta a una domanda:} La soglia $V_T$ del modello EIF ha un significato diverso a seconda del protocollo di stimolazione. Con una corrente costante, la soglia operativa è quella che determina la corrente minima per produrre un treno di spike (corrente di reobase). Con un impulso molto breve, la soglia effettiva è $V_T$, che coincide con la base della rampa esponenziale. I due valori non coincidono nel modello EIF (a differenza del LIF, come si vedrà più avanti).
\end{tcolorbox}

\section{Validazione Sperimentale: Esperimento Gerstner–Markram}

\subsubsection{Setup Sperimentale}

Il professore descrive un'elegante serie di esperimenti condotti nel \textbf{laboratorio di Wulfram Gerstner} (EPFL, Losanna) in collaborazione con il \textbf{laboratorio di Henry Markram}. Il protocollo sperimentale prevedeva:

\begin{enumerate}
    \item Preparazione di una \textbf{fettina corticale} (\textit{cortical slice}) di cervello di roditore: un segmento di tessuto cerebrale mantenuto in vita in vitro.
    \item Selezione di un \textbf{neurone piramidale} dello strato corticale e inserimento di un elettrodo intracellulare al soma.
    \item Iniezione di una \textbf{corrente rumorosa} $I_{\text{in}}(t)$ attraverso l'elettrodo.
    \item Registrazione simultanea del \textbf{potenziale di membrana} $V(t)$ al soma.
\end{enumerate}

\subsubsection{Perché una Corrente Rumorosa?}

La scelta di usare una corrente stocastica (rumorosa) ha una motivazione fisica precisa: il \textbf{barrage di impulsi sinaptici} che raggiunge il soma attraverso le dendrite può essere efficacemente approssimato da un \textbf{processo stocastico}. Iniettando al soma una corrente rumorosa che simula questo processo stocastico sinaptico, gli sperimentatori avevano accesso contemporaneo a:

\begin{itemize}
    \item $I_{\text{in}}(t)$: l'\textbf{input} al sistema (corrente iniettata, nota).
    \item $V(t)$: l'\textbf{output} del sistema (potenziale di membrana, registrato).
\end{itemize}

Questo permetteva di caratterizzare la relazione ingresso–uscita del neurone biologico.

\subsubsection{Osservazioni Sperimentali}

Dalla registrazione si osservava che:

\begin{itemize}
    \item Il neurone produceva \textbf{spike} in corrispondenza delle fluttuazioni della corrente che avvicinavano il potenziale alla soglia.
    \item Il \textbf{potenziale di membrana trascorreva la maggior parte del tempo al di sotto della soglia}, fluttuando senza emettere spike.
    \item I saltuari attraversamenti della soglia (grazie alle fluttuazioni stocastiche) producevano spike isolati.
    \item Gli \textbf{intervalli interspike (ISI) erano molto variabili}, coerentemente con quanto si osserva \textit{in vivo} in neuroni che ricevono input sinaptici naturalistici.
\end{itemize}



\subsection{Estrazione del Campo di Forza dalla Registrazione}

Il professore illustra la strategia ingegnosa adottata dagli sperimentatori per estrarre il \textbf{campo di forza} $\mathcal{F}(V)$ direttamente dalla registrazione. Sotto l'ipotesi che il neurone sia ben descritto da un circuito RC con un solo grado di libertà (il potenziale di membrana), il bilancio di corrente al soma è:

\begin{equation}
C_m \frac{dV}{dt} = \mathcal{F}(V) + I_{\text{in}}(t)
\end{equation}

La \textbf{corrente transmembrana} $I_m(t)$ (corrente che fluisce attraverso la membrana, escludendo la corrente iniettata) è quindi:

\begin{equation}
I_m(t) = I_{\text{in}}(t) - C_m \frac{dV}{dt}
\end{equation}

Poiché si conoscono sia $I_{\text{in}}(t)$ che $V(t)$ (e quindi $dV/dt$ numericamente), $I_m(t)$ è \textbf{direttamente misurabile}. E dalla definizione del campo di forza si ottiene:

\begin{equation}
\mathcal{F}(V) = -C_m \frac{dV}{dt} + I_{\text{in}} = -I_m(t) \cdot \frac{1}{\text{segno}}
\end{equation}

In termini operativi: $I_m(t)$ è direttamente proporzionale a $-\mathcal{F}(V)$.

\subsubsection{Media Condizionata e Scatter Plot}

Poiché si tratta di un processo stocastico, il campo di forza deterministico si estrae con una \textbf{media condizionata}:

\begin{equation}
\hat{\mathcal{F}}(V_0) = -\frac{1}{C_m} \langle I_m \mid V = V_0 \rangle
\end{equation}

Il procedimento sperimentale era il seguente:

\begin{enumerate}
    \item Per ogni campione temporale, si misuravano la coppia $(V(t), I_m(t))$ e si riportavano su uno \textbf{scatter plot} con $V$ sull'asse $x$ e $I_m$ sull'asse $y$.
    \item Si suddivideva l'asse $x$ in finestre strette di potenziale di membrana.
    \item Per ogni finestra, si calcolava la \textbf{media della distribuzione di} $I_m$ condizionata a quel valore di $V$.
    \item Le medie risultanti (punti rossi con barre di errore = errore standard della media) delineavano la \textbf{curva del campo di forza}.
\end{enumerate}

La forma risultante mostrava:

\begin{itemize}
    \item Un \textbf{ramo lineare} (dinamica passiva sotto-soglia, coerente con il termine $-G_m(V - E_m)$).
    \item Un \textbf{ramo esponenziale} che deviava dalla linearità in prossimità di $V_T$ (amplificazione attiva del sodio).
\end{itemize}



\subsection{Fit dei Parametri EIF su Neuroni Corticali}

Una volta estratta la curva sperimentale del campo di forza, il gruppo Gerstner–Markram procedeva con un \textbf{fit non-lineare} della funzione:

\begin{equation}
\mathcal{F}(V) = -G_m(V - E_m) + G_m \Delta_T \exp\!\left(\frac{V - V_T}{\Delta_T}\right)
\end{equation}

I \textbf{quattro parametri} liberi da determinare erano:

\begin{table}[htbp]
    \centering
    \renewcommand{\arraystretch}{1.3}
    \begin{tabularx}{\textwidth}{|l|X|X|}
    \hline
    \textbf{Parametro} & \textbf{Significato} & \textbf{Valore tipico} \\
    \hline
    $G_m$ & Conduttanza di membrana passiva & Variabile (cfr. $\tau_m$) \\
    \hline
    $E_m$ & Potenziale di equilibrio & $\sim -65$ mV \\
    \hline
    $\Delta_T$ & Sharpness del threshold & $\sim 1$ mV \\
    \hline
    $V_T$ & Soglia di iniziazione dello spike & $\sim -50$ mV \\
    \hline
    \end{tabularx}
\end{table}

\subsubsection{Verifica dell'Esponenziale}

Per verificare che la componente esponenziale fosse davvero di forma $e^{V/\Delta_T}$ (e non, ad esempio, una potenza), si rimuoveva il contributo lineare passivo $(-G_m(V - E_m))$ dalla curva sperimentale e si riportava il residuo $\mathcal{F}(V) - \mathcal{F}_{\text{lin}}(V)$ in \textbf{scala semi-logaritmica} ($y$-axis in scala log, $x$-axis lineare). In questa rappresentazione, una vera esponenziale appare come una \textbf{retta}, con pendenza $1/\Delta_T$. I punti sperimentali in prossimità di $V_T$ seguivano effettivamente una retta, confermando la natura esponenziale dell'amplificazione.

\subsubsection{Distribuzione dei Parametri tra Neuroni}

L'esperimento fu ripetuto su \textbf{decine di neuroni} appartenenti agli stessi strati corticali di diversi topi, ottenendo una \textbf{distribuzione dei parametri} tra neuroni:

\begin{itemize}
    \item $\tau_m = C_m/G_m$: distribuito tra 10 e 40 ms (ampia variabilità biologica).
    \item $E_m$: distribuzione centrata intorno a $-65$ mV.
    \item $V_T - E_m$: variabile (distanza tra soglia ed equilibrio).
    \item $\Delta_T$: $\sim 1$ mV, con variabilità limitata.
\end{itemize}

Il professore sottolinea che questa variabilità inter-neuronale è del tutto attesa in sistemi biologici, così come in qualsiasi sistema fisico. Ciascun neurone ha parametri biologici propri determinati dalla sua storia e dal suo contesto di sviluppo.

\begin{tcolorbox}[colback=gray!5, colframe=gray!40, arc=2mm, boxrule=0.5pt]
\textbf{Domanda di uno studente:} Le distribuzioni sono rispetto a quali condizioni? I parametri sono deterministici?\\
\textbf{Risposta del professore:} I parametri sono deterministici per ogni singolo neurone, ma variano da un neurone all'altro. Le distribuzioni mostrate rappresentano la variabilità tra i diversi neuroni campionati, tutti appartenenti agli stessi strati corticali di diversi topi.
\end{tcolorbox}



\subsubsection{Confronto Tracce: Biologica vs Modello EIF}

Con i parametri estratti dal fit, si testava la \textbf{capacità predittiva} del modello EIF confrontando la \textbf{traccia nera} (potenziale registrato biologicamente) con la \textbf{traccia rossa} (dinamica del neurone EIF fittato, guidato dalla stessa realizzazione di corrente rumorosa iniettata nel neurone reale).

I risultati mostravano:

\begin{itemize}
    \item \textbf{Ottima riproduzione della dinamica sotto-soglia}: le tracce rossa e nera erano spesso indistinguibili nella fase sub-threshold.
    \item \textbf{Buona predizione del timing degli spike}: il modello riproduceva fedelmente la maggior parte degli spike della traccia biologica.
    \item \textbf{Fallimento per doublet e triplet}: in alcune finestre temporali, il neurone reale produceva due o tre spike ravvicinati (doublet/triplet) che il modello 1D non riusciva a riprodurre.
    \item \textbf{Falsi positivi}: in altri casi il modello produceva spike assenti nella traccia biologica.
\end{itemize}

\subsubsection{Elemento Mancante: Effetto dello Spike sul Campo di Forza}

Il professore identifica la causa del fallimento: il modello 1D, per ridurre la dimensionalità da 2D a 1D, aveva \textbf{congelato} la variabile di recupero $W$. Ma la produzione di un potenziale d'azione ha un impatto reale sulla dinamica successiva (correnti di calcio, canali K$^+$ attivati da Ca$^{2+}$) che si manifesta come una \textbf{modifica temporanea del campo di forza} dopo ogni spike.



\subsection{Analisi Spike-Triggered e Dinamica di $V_T$}

\subsubsection{Statistica Spike-Triggered}

Per identificare l'elemento mancante, il gruppo adottò un'analisi \textbf{spike-triggered}: si allineavano tutte le tracce registrate al \textbf{tempo di emissione di uno spike isolato} (non appartenente a burst), e si raccoglievano i campioni di $(V(t), I_m(t))$ in diverse \textbf{finestre temporali} dopo lo spike:

\begin{itemize}
    \item Finestra 1: 5–10 ms dopo lo spike
    \item Finestra 2: 10–20 ms dopo lo spike
    \item Finestra 3: 20–30 ms dopo lo spike
    \item Finestra 4: 30–50 ms dopo lo spike
\end{itemize}

Per ciascuna finestra si ripeteva la procedura della media condizionata e si estraeva il campo di forza $\mathcal{F}(V)$ corrispondente.

\subsubsection{Risultato: Spostamento Transitorio di $V_T$ ed $E_m$}

Il risultato principale era che il campo di forza \textbf{cambia nelle ore successive a uno spike}:

\begin{enumerate}
    \item \textbf{Nelle finestre temporali brevi} (5–10 ms): la soglia $V_T$ era significativamente più alta (verso potenziali più depolarizzati) rispetto al valore stazionario, e anche il punto di intersezione con $\mathcal{F} = 0$ (l'equilibrio $E_m$) era spostato verso destra.
    \item \textbf{Nelle finestre successive} (10–20 ms, 20–30 ms): la soglia $V_T$ si avvicinava progressivamente al valore asintotico $V_T(\infty)$.
    \item \textbf{Dopo $\sim$15 ms}: il campo di forza tornava al suo profilo stazionario.
\end{enumerate}

Questo significava:

\begin{itemize}
    \item Il neurone era \textbf{più difficile da eccitare} subito dopo uno spike (soglia più alta): questo corrisponde al \textbf{periodo refrattario relativo}.
    \item Il potenziale di membrana era \textbf{più depolarizzato} subito dopo uno spike (equilibrio spostato): il potenziale di reset aveva una dipendenza temporale.
\end{itemize}

\subsubsection{Meccanismo Fisico}

L'interpretazione fisica era la seguente: durante l'emissione di uno spike, il potenziale di membrana raggiunge valori molto depolarizzati (0 mV, +50 mV), provocando l'\textbf{apertura di canali Ca$^{2+}$} e un influx di ioni calcio intracellulare. L'aumento della concentrazione di Ca$^{2+}$ attivava correnti del potassio dipendenti dal calcio (canali BK e SK), responsabili di:

\begin{enumerate}
    \item Un'\textbf{iperpolarizzazione} dopo-spike (after-hyperpolarization, AHP).
    \item Un \textbf{innalzamento transitorio della soglia} di emissione (threshold fatigue).
\end{enumerate}

\subsubsection{Dinamica Temporale di $V_T(t)$}

Il professore sintetizza il comportamento di $V_T$ nel tempo:

\begin{equation}
V_T(t) = V_T(\infty) + \sum_k \Delta V_T \cdot \delta(t - t_k^{\text{spike}})
\end{equation}

con una \textbf{dinamica di rilassamento} verso $V_T(\infty)$. In pratica:

\begin{itemize}
    \item Ad ogni spike emesso al tempo $t_k$, $V_T$ \textbf{salta} di un incremento $\Delta V_T > 0$ (soglia più alta).
    \item Tra uno spike e il successivo, $V_T$ \textbf{decade esponenzialmente} verso $V_T(\infty)$.
    \item Se gli spike si susseguono rapidamente (burst), la soglia si accumula progressivamente: più spike vengono emessi, più difficile è emettere il successivo.
\end{itemize}

\subsection{Modello EIF Adattivo: Riproduzione dei Burst}

\subsubsection{Il Modello con Soglia Dinamica}

Introducendo una \textbf{variabile di soglia dinamica} $V_T(t)$ (un grado di libertà aggiuntivo che si comporta come la variabile di recupero $W$ dei modelli bidimensionali), il modello EIF diventava capace di riprodurre le caratteristiche mancanti. La \textbf{traccia verde} nel grafico mostrava il risultato del modello EIF con questo grado di libertà aggiuntivo:

\begin{itemize}
    \item Dopo ogni spike la traccia verde non era più discontinua (il potenziale di reset evolveva gradualmente grazie alla dinamica della soglia).
    \item I \textbf{doublet} venivano riprodotti correttamente.
    \item I \textbf{triplet} erano anch'essi catturati dal modello.
    \item Il potenziale di membrana negli inter-spike interval dei burst era riprodotto in modo affidabile.
\end{itemize}

\subsubsection{Conclusione: i Modelli IF non sono Modelli Giocattolo}

\begin{tcolorbox}[colback=gray!5, colframe=gray!40, arc=2mm, boxrule=0.5pt]
\textbf{Nota del docente:} Questo è per dirvi che siamo partiti dal modello di Hodgkin-Huxley, molto non-lineare e quadridimensionale. Ma semplicemente reintroducendo una variabile di recupero — e vedremo in che modo — in un neurone integrate-and-fire monodimensionale, siamo in grado di essere molto, molto accurati nel riprodurre il potenziale di membrana di neuroni reali. Quindi questi neuroni integrate-and-fire non sono modelli giocattolo, ma strumenti quantitativi precisi per riprodurre ciò che è registrato \textit{in vivo}.
\end{tcolorbox}


\section{Modelli IF come Approssimazioni dell'EIF: QIF e LIF}

\subsubsection{Espansione in Serie di Taylor}

Per ricavare i due modelli IF più utilizzati come approssimazioni dell'EIF, si considera un'espansione in serie di Taylor del campo di forza $\mathcal{F}(V)$ intorno al potenziale di equilibrio $E_m$, limitandosi ai termini fino al \textbf{secondo ordine}:

\begin{equation}
\mathcal{F}(V) \approx c_0 + c_1(V - E_m) + c_2(V - E_m)^2
\end{equation}

Questa forma parabolica dipendente da $V^2$ giustifica il nome \textbf{Quadratic Integrate-and-Fire} (QIF).

\subsubsection{Modello QIF: Campo di Forza e Energia}

Nel modello QIF il campo di forza è parabolico e il corrispondente \textbf{paesaggio energetico} è cubico: ha ancora una \textbf{sella} (che rappresenta la soglia di emissione) e una \textbf{valle} (il punto di riposo), conservando le caratteristiche qualitative essenziali del modello EIF.

\begin{tcolorbox}[colback=gray!5, colframe=gray!40, arc=2mm, boxrule=0.5pt]
\textbf{Domanda di uno studente:} L'espansione è attorno a $V_T$ o attorno a $E_m$?\\
\textbf{Risposta del professore:} In linea di principio si può espandere attorno a qualsiasi punto; l'obiettivo è minimizzare l'errore di approssimazione nell'intervallo di potenziale rilevante (il range sotto-soglia tra $E_m$ e $V_T$). In questa derivazione ho usato l'espansione attorno a $E_m$, ma se si vuole ridurre l'errore vicino alla soglia si può espandere attorno a $V_T$. È una scelta arbitraria che dipende dall'uso che si vuole fare del modello.
\end{tcolorbox}

\subsubsection{Connessione QIF $\leftrightarrow$ Oscillatori di Kuramoto (Theta Neurons)}

Un risultato importante riguarda il legame tra il modello QIF e gli \textbf{oscillatori di Kuramoto}. Introducendo la variabile di fase $\theta$ tramite la trasformazione non-lineare:

\begin{equation}
V = V_T - \cot(\theta/2)
\end{equation}

l'emissione di uno spike ($V \to +\infty$) corrisponde a $\theta \to \pi$, e la fase compie un giro completo sul cerchio unitario. In questa rappresentazione, il modello QIF diventa un \textbf{theta neuron} (neurone di fase), la cui dinamica è:

\begin{equation}
\frac{d\theta}{dt} = (1 - \cos\theta) + (1 + \cos\theta) \cdot I
\end{equation}

che mappa esattamente sulla dinamica di un singolo oscillatore di Kuramoto. Di conseguenza, le \textbf{teorie di campo medio per reti di oscillatori di Kuramoto} si traducono direttamente in descrizioni analitiche per reti di neuroni QIF.

\subsubsection{Modello LIF: Il Limite $\Delta_T \to 0$}

Il \textbf{Leaky Integrate-and-Fire} (LIF) si ottiene dal modello EIF portando $\Delta_T \to 0$. In questo limite:

\begin{itemize}
    \item Il ramo esponenziale scompare e il campo di forza diventa puramente lineare:
    \begin{equation}
    \mathcal{F}_{\text{LIF}}(V) = -G_L(V - E_L)
    \end{equation}
    \item La soglia $V_T$ diventa una \textbf{barriera rigida}: appena $V$ supera $V_T$, viene emesso uno spike; non c'è la possibilità di rientrare sotto-soglia anche in presenza di forte corrente inibitoria.
    \item $V_T$ e $V_{\text{reobase}}$ \textbf{coincidono}: poiché non c'è più il ramo esponenziale che li distingueva, nel modello LIF la soglia per corrente costante e la soglia per impulso breve sono la stessa quantità.
\end{itemize}

\begin{tcolorbox}[colback=gray!5, colframe=gray!40, arc=2mm, boxrule=0.5pt]
\textbf{Nota del docente:} Nel LIF, $V_T$ e $V_{\text{reobase}}$ sono lo stesso numero. Questo risponde alla domanda di prima: nel modello EIF erano diversi (la reobase è il minimo del campo di forza perturbato dalla corrente; $V_T$ è la base dell'esponenziale). Nel limite $\Delta_T \to 0$ questi due valori convergono a una singola soglia.
\end{tcolorbox}

Il LIF è importante per la sua \textbf{trattabilità analitica} nella descrizione di popolazioni di neuroni tramite l'equazione di Fokker-Planck, argomento che sarà affrontato nelle lezioni successive.


\begin{tcolorbox}[colback=gray!5, colframe=gray!40, arc=2mm, boxrule=0.5pt]
\textbf{Domanda:} Quando si espande attorno a $V_T$, è perché si vuole un interruttore (switch) al threshold?\\
\textbf{Risposta del professore:} Sì. Se si espande attorno a $V_T$, e poi si inserisce una soglia rigida a $V_T$, si ha un'approssimazione dell'EIF attorno al punto di emissione. In questo caso la parabola ha il vertice a $V_T$. Nella mia derivazione ho espanso attorno a $E_m$, che è il modo più naturale se si vuole catturare bene la dinamica sottosoglia. La scelta ottimale dipende dall'applicazione: se si vuole una buona approssimazione della rampa pre-spike si espande vicino a $V_T$; se si vuole minimizzare l'errore nell'intero range sottosoglia si espande vicino a $E_m$.
\end{tcolorbox}


\section{Tassonomia dei Pattern di Firing Neuronale}

\subsubsection{Il Paper di Markram e Collaboratori}

Il professore introduce un articolo di riferimento fondamentale in neuroscienza computazionale: la \textbf{review di $\sim$20 anni fa} del laboratorio Markram e collaboratori, che ha tentato una \textbf{tassonomia sistematica} dei neuroni inibitori in base ai loro pattern di firing elettrofisiologici. Lo studio includeva neuroni non solo dalla corteccia cerebrale, ma anche da strutture sottocorticali come \textbf{gangli della base}, \textbf{tronco cerebrale} e \textbf{talamo}, dove si osservano comportamenti più peculiari.

\subsubsection{Protocollo di Stimolazione Standard}

Il protocollo sperimentale era standardizzato:

\begin{enumerate}
    \item Inserire un elettrodo intracellulare nel neurone in studio (nella fettina cerebrale).
    \item Iniettare una \textbf{corrente a gradino} (stepwise current) di intensità costante, per un intervallo di tempo definito.
    \item Ripetere con due (o più) livelli di corrente: una corrente \textbf{bassa} e una corrente \textbf{alta}.
    \item Osservare la risposta del neurone in termini di treni di spike.
\end{enumerate}

\subsubsection{Classe 1: Neuroni NAC (Non-Accommodating)}

I neuroni \textbf{NAC} (Non-Accommodating, non-accomodanti) mostrano un \textbf{firing tonico regolare}:

\begin{itemize}
    \item In risposta alla corrente a gradino, producono un treno di spike con \textbf{intervalli interspike costanti nel tempo} (ISI non cambia dalla prima alla n-esima spike).
    \item Al crescere dell'intensità della corrente iniettata, la frequenza di firing aumenta, ma rimane regolare.
    \item \textbf{Non c'è adattamento della frequenza} nel tempo.
\end{itemize}

\subsubsection{Classe 2: Neuroni AC (Accommodating)}

I neuroni \textbf{AC} (Accommodating, accomodanti) mostrano \textbf{spike frequency adaptation}:

\begin{itemize}
    \item All'inizio della stimolazione: alta frequenza di spike.
    \item Nel tempo: gli ISI diventano \textbf{progressivamente più lunghi} (la frequenza di firing diminuisce).
    \item La frequenza si stabilizza a un valore asintotico (se la stimolazione è sufficientemente lunga).
    \item Questo è il comportamento che abbiamo discusso come \textbf{adattamento della frequenza di spike} nelle lezioni precedenti.
\end{itemize}

\subsubsection{Classi con Burst Iniziale}

Esistono sottoclassi di neuroni che producono un \textbf{burst iniziale} ad alta frequenza all'inizio della stimolazione, seguito da:

\begin{itemize}
    \item \textbf{Burst iniziale + tonico}: burst all'inizio, poi firing tonico regolare.
    \item \textbf{Burst iniziale + adattante}: burst all'inizio, poi firing con spike frequency adaptation.
\end{itemize}

\subsubsection{Classi con Ritardo}

Alcune classi mostrano un \textbf{ritardo} nell'emissione del primo spike:

\begin{itemize}
    \item \textbf{Delayed tonic} (ritardato + tonico): il primo spike arriva con un ritardo rispetto all'inizio della stimolazione; poi firing regolare.
    \item \textbf{Delayed adapting} (ritardato + adattante): ritardo nel primo spike, poi firing adattante.
\end{itemize}

\subsubsection{Regular Bursting}

Una classe particolarmente interessante è quella dei \textbf{neuroni a burst regolare}:

\begin{itemize}
    \item In risposta alla stimolazione, il neurone produce \textbf{burst periodici} separati da \textbf{periodi di silenzio}.
    \item Aumentando l'intensità della corrente, non cambia la frequenza di spike all'interno del burst, ma si riduce la \textbf{distanza temporale tra burst consecutivi} (inter-burst interval).
\end{itemize}

\section{Profilo delle Classi Cellulari Principali}

\subsubsection{Cellule Piramidali Eccitatorie}

Le \textbf{cellule piramidali} della corteccia eccitatoria appartengono generalmente alla classe \textbf{Regular Spiking} (RS):

\begin{itemize}
    \item Mostrano \textbf{spike frequency adaptation} (accomodanti, AC).
    \item Non producono burst.
    \item Firing relativamente basso grazie all'adattamento.
    \item Per questo vengono chiamate "regular spiking" (spiking regolare senza burst) anche se con il tempo la frequenza diminuisce.
\end{itemize}

\begin{tcolorbox}[colback=gray!5, colframe=gray!40, arc=2mm, boxrule=0.5pt]
\textbf{Nota del docente:} I neuroni piramidali eccitatori sono spesso denominati "regular spiking cells" perché il loro firing è regolare, senza burst, anche se hanno spike frequency adaptation. La frequenza non è troppo alta proprio grazie al meccanismo di adattamento.
\end{tcolorbox}

\subsubsection{Interneuroni Inibitori}

Gli \textbf{interneuroni inibitori} mostrano prevalentemente comportamento \textbf{NAC} o fast spiking:

\begin{itemize}
    \item Firing tonico ad \textbf{alta frequenza} (possono produrre decine o centinaia di Hz).
    \item Necessità funzionale: devono poter controllare in modo preciso e rapido l'attività delle cellule piramidali circostanti.
    \item In alcune classi: burst iniziali seguiti da firing tonico.
\end{itemize}

\begin{tcolorbox}[colback=gray!5, colframe=gray!40, arc=2mm, boxrule=0.5pt]
\textbf{Nota del docente:} Gli interneuroni inibitori sono lì per controllare l'attività dei neuroni piramidali in modo che non diventi troppo grande. Quindi devono essere in grado di produrre alta frequenza di spike in modo tonicamente sostenuto.
\end{tcolorbox}

\section{Il Modello Adattivo Integrate-and-Fire (AdIF / AdEx)}

\subsubsection{Equazioni del Modello 2D}

Per riprodurre la ricca varietà di pattern di firing osservati biologicamente, si introduce il \textbf{modello adattivo integrate-and-fire} (AdIF), equivalente al modello \textbf{AdEx} di Brette e Gerstner (2005). Il modello ha due variabili di stato: il potenziale di membrana $V$ e la \textbf{variabile di recupero} (adattamento) $W$:

\begin{equation}
C_m \frac{dV}{dt} = \mathcal{F}(V) + I - W
\end{equation}

\begin{equation}
\tau_W \frac{dW}{dt} = a(V - E_L) - W
\end{equation}

\textbf{Condizioni di reset} ad ogni spike (quando $V$ supera la soglia $V_{\text{th}}$):

\begin{equation}
V \to V_{\text{reset}}
\end{equation}
\begin{equation}
W \to W + b
\end{equation}

I \textbf{parametri chiave} del modello sono:

\begin{table}[htbp]
    \centering
    \renewcommand{\arraystretch}{1.3}
    \begin{tabularx}{\textwidth}{|l|X|}
    \hline
    \textbf{Parametro} & \textbf{Significato} \\
    \hline
    $a$ & Accoppiamento sotto-soglia (subthreshold adaptation) \\
    \hline
    $b$ & Grandezza del salto di $W$ per ogni spike (spike-triggered adaptation) \\
    \hline
    $\tau_W$ & Costante di tempo della variabile di adattamento \\
    \hline
    $V_{\text{reset}}$ & Potenziale di reset post-spike \\
    \hline
    \end{tabularx}
\end{table}

\subsubsection{Meccanismo Fisico del Salto $b$}

Il salto $W \to W + b$ ad ogni spike rappresenta l'\textbf{influx di Ca$^{2+}$ durante lo spike} (attraverso i canali del calcio voltaggio-dipendenti), che attiva le correnti del potassio $K^+$ dipendenti dal calcio (canali SK e BK), aumentando la corrente inibitoria. Nella dinamica del modello, un $W$ più alto corrisponde a una corrente inibitoria più forte che si oppone all'eccitazione, rendendo più difficile la produzione del successivo spike.

\subsubsection{Connessione con i Modelli Molecolari}

Il professore ricorda che la variabile $W$ era stata introdotta nella riduzione del modello FitzHugh-Nagumo come proporzionale alla variabile di gating $n$ (canali K$^+$ del modello HH), con dinamica del primo ordine:

\begin{equation}
\tau_W \frac{dW}{dt} = a(V - E_L) - W
\end{equation}

che rispecchia la cinetica di $n_\infty(V) - n$ (equilibrio dipendente dalla tensione meno valore attuale). Il parametro $a > 0$ introduce un accoppiamento tra $V$ e la nullcline $\dot{W} = 0$, che nel piano di fase si manifesta come una nullcline con pendenza positiva invece che orizzontale.


\subsection{Firing tonico non-adattivo}

Il professore analizza il comportamento del modello 2D nel \textbf{piano di fase} $(V, W)$ per diversi set di parametri. Il primo caso corrisponde al \textbf{firing tonico non-adattivo} (NAC):

\begin{itemize}
    \item $\tau_W = 30$ ms (variabile di recupero relativamente veloce)
    \item $b = 60$ pA (salto grande per ogni spike)
    \item $V_{\text{reset}} = -55$ mV
\end{itemize}

\subsubsection{Analisi della Traiettoria nel Piano di Fase}

Il protocollo di stimolazione è un \textbf{gradino di corrente costante}, che nella rappresentazione del piano di fase sposta verso l'alto la nullcline $\dot{V} = 0$ (la corrente rende il neurone più eccitabile).

\begin{enumerate}
    \item \textbf{Condizione iniziale}: $W(0) = 0$, $V(0) < V_{\text{th}}$. Il sistema si trova sotto la nullcline $\dot{V} = 0$ (campo di forza positivo) e la nullcline $\dot{W} = 0$ orizzontale.
    \item \textbf{Avvicinamento alla soglia}: il campo di forza spinge $V$ verso il threshold $V_{\text{th}}$.
    \item \textbf{Primo spike}: quando $V = V_{\text{th}}$, lo spike viene emesso. Il sistema viene reiniettato in $(V_{\text{reset}},\, W = 0 + b = b)$.
    \item \textbf{Dopo-iperpolarizzazione (AHP)}: partendo da $W = b > 0$, la traiettoria supera la nullcline $\dot{V} = 0$ (la corrente inibitoria $W$ è abbastanza forte da iperpolarizzare il neurone) e poi torna.
    \item \textbf{Secondo spike}: la traiettoria si avvicina nuovamente alla soglia. Questa volta però $W$ non è tornato a 0 (poiché $\tau_W = 30$ ms non è abbastanza lungo): $W$ al momento del secondo spike è diverso da zero, e il jump porta $W$ a un valore ancora più alto.
    \item \textbf{Convergenza a orbita stabile}: dopo alcuni spike, il sistema converge a un ciclo limite in cui l'ISI è costante.
\end{enumerate}

\begin{tcolorbox}[colback=gray!5, colframe=gray!40, arc=2mm, boxrule=0.5pt]
\textbf{Chiarimento del docente a una domanda dello studente:} Non c'è un'orbita nello spazio di fase nel senso tradizionale: c'è un salto verticale in $W$ al momento di ogni spike, che porta il sistema in punti diversi dello spazio di fase ad ogni emissione. Il "detour" visibile nel disegno è artificiale, perché la produzione dello spike non è modellata esplicitamente ma solo con il reset. Il sistema viene reiniettato in posizioni diverse dello spazio di fase per i diversi spike, fino alla convergenza.
\end{tcolorbox}

\subsubsection{ISI Costante $\rightarrow$ Firing Tonico NAC}

La frequenza asintotica di firing $f_\infty = 1/\text{ISI}_\infty$ è costante nel tempo: \textbf{firing tonico non-adattivo} (NAC).

\subsection{Spike Frequency Adaptation nel Piano di Fase}


Per ottenere lo \textbf{spike frequency adaptation}, i parametri chiave da modificare sono:

\begin{itemize}
    \item $\tau_W = 100$ ms (variabile di recupero più lenta, nel precedente $\tau_W = 30$ ms)
    \item $b = 10$ pA (salto più piccolo per ogni spike, nel precedente $b = 60$ pA)
\end{itemize}

\subsubsection{Flusso Quasi-Orizzontale}

Poiché il campo di forza della variabile di recupero è:

\begin{equation}
\frac{dW}{dt} = \frac{a(V - E_L) - W}{\tau_W}
\end{equation}

un $\tau_W$ grande implica che il \textbf{flusso è quasi-orizzontale} nel piano di fase. Le frecce verticali sono molto piccole, e la traiettoria si muove quasi parallelamente all'asse $V$.

\subsubsection{Accumulo di $W$ e Rallentamento del Firing}

In questa configurazione:

\begin{itemize}
    \item Il primo spike produce un salto piccolo in $W$ ($b = 10$ pA).
    \item Nell'interspike interval, $W$ decade molto lentamente (essendo $\tau_W = 100$ ms, paragonabile alla durata dell'ISI).
    \item Al secondo spike, $W$ non è tornato a zero: il salto $b$ porta $W$ a un valore ancora più alto.
    \item Progressivamente, $W$ si accumula verso un valore asintotico $W_\infty$.
\end{itemize}

Il valore asintotico $W_\infty$ determina la \textbf{corrente effettiva netta} percepita dal neurone:

\begin{equation}
I_{\text{netto}} = I - W_\infty < I
\end{equation}

La frequenza di firing asintotica è quindi inferiore a quella iniziale:

\begin{equation}
f_\infty = \frac{1}{\text{ISI}_\infty} \propto I - W_\infty < f(t=0)
\end{equation}

\subsubsection{Interpretazione come Corrente Inibitoria Stazionaria}

Da un punto di vista fisico, $W_\infty$ può essere interpretata come una \textbf{corrente inibitoria equivalente stazionaria} (proporzionale all'attività totale dei canali K$^+$ attivati dallo storico degli spike). Più spike vengono emessi, più forte è questa corrente inibitoria, il che rallenta il firing. Il sistema raggiunge l'equilibrio quando la corrente inibitoria compensa parzialmente la corrente eccitatoria iniettata.

\begin{tcolorbox}[colback=gray!5, colframe=gray!40, arc=2mm, boxrule=0.5pt]
\textbf{Nota del docente:} La pendenza delle traiettorie nel piano di fase nell'interspike interval diventa più e più piccola avvicinandosi alla nullcline $\dot{V} = 0$. Questo spiega fisicamente perché il tempo per raggiungere la soglia (l'ISI) aumenta: la corrente di guida $\mathcal{F}(V)$ si riduce perché il termine $-W$ diventa sempre più grande.
\end{tcolorbox}


\subsection{Burst Iniziale + Tonico nel Piano di Fase}


Per ottenere un \textbf{burst iniziale seguito da firing tonico}, si modificano diversi parametri rispetto ai casi precedenti:

\begin{itemize}
    \item La nullcline $\dot{W} = 0$ diventa \textbf{dipendente da $V$} (slope positivo, $a > 0$), come nel modello FitzHugh-Nagumo.
    \item La \textbf{costante di tempo di membrana} $\tau_m$ viene ridotta (neurone più veloce, più eccitabile).
    \item $V_{\text{reset}}$ viene avvicinato alla soglia $V_{\text{th}}$ (il neurone parte già vicino alla soglia dopo ogni spike).
    \item Corrente iniettata più alta per compensare la maggiore eccitabilità.
    \item $\tau_W$ lungo e $b$ piccolo.
\end{itemize}

\subsubsection{Dinamica del Burst Iniziale}

Nella fase iniziale della stimolazione:

\begin{enumerate}
    \item Il neurone riceve la corrente e, avendo $V_{\text{reset}}$ vicino a $V_{\text{th}}$, emette il \textbf{primo spike} molto rapidamente.
    \item Il salto $b$ è piccolo ($b = 7$ pA): il sistema rientra in una posizione con $W$ quasi nullo.
    \item Il campo di forza per $V$ è forte (il neurone è veloce), quindi il secondo spike viene emesso rapidamente.
    \item Questo produce un \textbf{burst} (ISI brevi) nella fase iniziale.
\end{enumerate}

\subsubsection{Transizione al Firing Tonico}

Man mano che il burst procede, $W$ si accumula. Quando le traiettorie si avvicinano alla nullcline $\dot{V} = 0$, la pendenza del potenziale di membrana nelle fasi di avvicinamento alla soglia diventa \textbf{più lenta}. Dopo il 4$^\circ$–5$^\circ$ spike, l'ISI inizia ad allungarsi e si stabilizza sul valore asintotico (firing tonico).

Il campo di forza $G = \dot{W} \approx 0$ nella regione di lavoro (perché $a(V-E_L) \approx W$, grazie alla nullcline con slope positivo) significa che il flusso nella direzione di $W$ è quasi nullo: le frecce verticali sono quasi assenti, e il sistema scivola quasi orizzontalmente verso la soglia nelle fasi inter-spike.

\subsection{Regular Burst Firing}


Per il \textbf{regular burst firing}, la modifica fondamentale rispetto al caso precedente è la posizione di $V_{\text{reset}}$:

\begin{equation}
V_{\text{reset}} > V_T
\end{equation}

Questo significa che, dopo ogni spike, il potenziale di membrana viene reiniettato \textbf{sul lato destro della nullcline} $\dot{V} = 0$ (dove $\mathcal{F}(V) < 0$, quindi il campo di forza iperpolarizza il neurone).

\subsubsection{Meccanismo del Burst Regolare nel Piano di Fase}

Con questo parametro:

\begin{enumerate}
    \item \textbf{Fase di burst}: il neurone emette spike rapidamente (ISI brevi) perché $V_{\text{reset}}$ è già vicino (o superiore) alla soglia. Le traiettorie sono veloci e producono tanti spike consecutivi.
    \item \textbf{Crossing della nullcline}: ad un certo punto, il salto $b$ porta il sistema sul lato destro della nullcline $\dot{V} = 0$, dove $\mathcal{F}(V) < 0$.
    \item \textbf{Forte iperpolarizzazione}: il campo di forza ora iperpolarizza il neurone. Si ha una discesa lenta del potenziale verso valori iperpolarizzati.
    \item \textbf{Recupero lento}: grazie al lento decadimento di $W$ (con $\tau_W$ lungo), il sistema si avvicina nuovamente alla soglia descrivendo una traiettoria lenta verso il basso.
    \item \textbf{Nuovo burst}: quando il potenziale risale fino alla soglia, comincia un nuovo burst.
\end{enumerate}

Questo ciclo si ripete periodicamente, producendo il \textbf{regular bursting}: burst separati da periodi di silenzio (AHP lenta), con frequenza inter-burst che dipende dall'intensità della corrente iniettata.

\subsubsection{Impossibilità del Regular Bursting nel Modello LIF}

Questo è il punto cruciale che giustifica l'utilizzo del modello EIF rispetto al LIF per descrivere il regular bursting:

\begin{itemize}
    \item \textbf{Condizione necessaria} per il regular bursting: la linea $V = V_{\text{reset}}$ interseca la nullcline $\dot{V} = 0$ alla sua destra ($V_{\text{reset}} > V_T$).
    \item \textbf{Nel modello LIF}, con $\Delta_T = 0$, la nullcline $\dot{V} = 0$ è una retta verticale a $V = E_L$ (il campo di forza è sempre negativo per $V > E_L$). Non esiste una posizione $V_{\text{reset}} > V_T$ che interseca questa nullcline a destra della soglia.
\end{itemize}

In altre parole: \textbf{il modello LIF non può produrre regular bursting}. Se si crede che il LIF sia un modello sufficiente, si deve accettare il fatto che i pattern di regular bursting non possono essere descritti da esso. Al contrario, il modello EIF (con il suo ramo esponenziale e la struttura del piano di fase che ammette $V_{\text{reset}} > V_T$) è necessario per catturare questo comportamento.

\begin{tcolorbox}[colback=gray!5, colframe=gray!40, arc=2mm, boxrule=0.5pt]
\textbf{Nota finale del docente:} Il modello leaky integrate-and-fire ha una nullcline verticale (perché $\Delta_T = 0$), e questa non si interseca mai con la linea $V_{\text{reset}}$ a destra di $V_T$. Quindi se credete che il LIF sia un modello sufficiente, dovete sapere che non sarà mai in grado di descrivere il regular burst firing. Ed è tutto per oggi. Arrivederci alla prossima lezione.
\end{tcolorbox}


\chapter{Lezione 11}


\section{Riepilogo: Neuroni 2D con Focus Stabile e Risonanza}

\subsubsection{Chiusura del Capitolo sui Modelli Neuronali Bidimensionali}

La lezione inizia con la conclusione del percorso dedicato ai \textbf{modelli neuronali bidimensionali}, in particolare al modello integrate-and-fire adattativo esponenziale con variabile di recupero $w$. Il professore fa riferimento alle simulazioni numeriche discusse nelle lezioni precedenti e riepilogo brevemente i due fenomeni fondamentali che questo tipo di neuroni può riprodurre: il \textbf{focus stabile} e la \textbf{risonanza sub-soglia}.

Nella versione del modello con variabile di recupero, il sistema è governato da:

\begin{equation}
C \frac{dV}{dt} = F(V) - w + I_{\text{ext}}
\end{equation}

\begin{equation}
\tau_w \frac{dw}{dt} = a(V - E_L) - w
\end{equation}

dove $F(V)$ contiene il termine esponenziale del modello EIF e $\tau_w$ è la costante di tempo della variabile di recupero. La nullcline $\dot{V} = 0$ presenta la sua caratteristica forma non monotona, mentre la nullcline $\dot{w} = 0$ è una retta di pendenza positiva nel piano di fase $(V, w)$.

\subsubsection{Il Focus Stabile e le Oscillazioni Smorzate Sub-Soglia}

Il caso di interesse è quello in cui il \textbf{punto fisso è un focus stabile}: l'intersezione delle due nullcline cade in una regione in cui la pendenza della nullcline $\dot{V} = 0$ è negativa. Come discusso nelle lezioni precedenti, la condizione necessaria affinché gli autovalori del Jacobiano siano complessi coniugati con parte reale negativa è che il punto fisso si trovi sul ramo decrescente della nullcline del potenziale di membrana. Questa condizione garantisce la presenza di oscillazioni smorzate sub-soglia.

Il profilo di rilassamento da una condizione iniziale è a \textbf{spirale convergente} nel piano di fase: la traiettoria avvolge il punto fisso descrivendo oscillazioni ammortizzate con frequenza di oscillazione $\omega \approx 2\pi / \tau_w$. In termini temporali, si osserva una sequenza di oscillazioni smorzate nel potenziale di membrana, la cui ampiezza decresce esponenzialmente nel tempo con costante di tempo data dalla parte reale degli autovalori.

\subsubsection{Neurone Risonante: Sensibilità alla Frequenza dell'Input}

Il \textbf{neurone risonante} è definito come un neurone che risponde in modo selettivo agli input oscillatori con frequenza prossima alla frequenza di risonanza naturale $\nu_{\text{res}} = 1/(2\pi\tau_w)$. Il professore spiega che se si somministra a questo tipo di neurone una corrente sinusoidale:

\begin{equation}
I_{\text{ext}}(t) = I_0 \sin(2\pi \nu t)
\end{equation}

la risposta del neurone è fortemente amplificata soltanto quando $\nu \approx \nu_{\text{res}}$: a questa frequenza l'energia fornita dall'esterno è in grado di far divergere le oscillazioni sub-soglia e — se l'ampiezza è sufficiente — di portare la traiettoria nel piano di fase oltre la separatrice, inducendo l'emissione di spike.

\begin{tcolorbox}[colback=gray!5, colframe=gray!40, arc=2mm, boxrule=0.5pt]
\textbf{Nota del docente:} Questo è il motivo per cui si parla di neuroni risonanti: si comportano come oscillatori forzati con una frequenza naturale. La risonanza è anche il meccanismo ipotizzato alla base della \textbf{sincronizzazione di rete}: se molti neuroni hanno la stessa frequenza di risonanza, essi risponderanno in modo coerente a stimoli oscillatori del network, e questa coerenza è alla base delle \textbf{oscillazioni cerebrali} (gamma, beta, theta, ecc.) che si ritiene medino la comunicazione tra aree distinte della corteccia.
\end{tcolorbox}

\subsubsection{Il Caso Eccitabile con Punto Fisso Stabile e Singolo Spike}

Il secondo scenario discusso è quello in cui il punto fisso rimane un focus stabile, ma la nullcline $\dot{V} = 0$ è stata spostata verso l'alto da una corrente iniettata di intensità maggiore. In questo caso il punto fisso è ancora stabile, ma esiste una \textbf{separatrice} che separa il bacino di attrazione del punto fisso dal cammino che conduce all'emissione di uno spike.

Partendo da una condizione iniziale che si trova al di là della separatrice — ad esempio, imponendo come condizione iniziale un potenziale di membrana appena sotto la soglia con la variabile di recupero a zero — il campo di forza trasporta il sistema verso il ramo superiore della nullcline $\dot{V} = 0$, producendo \textbf{un singolo spike}. Dopo l'emissione dello spike, il sistema viene resettato al potenziale di reset $V_{\text{res}} = -51 \; \text{mV}$, e da questo punto la traiettoria percorre una spirale smorzata che converge al focus stabile. Il sistema non emetterà altri spike a meno di ulteriori perturbazioni.

Questo comportamento è la firma tipica del \textbf{sistema eccitabile} di tipo II (Hodgkin): la risposta è un singolo grande detour nel piano di fase seguito da un ritorno all'equilibrio.

\newpage

\section{Dinamica Stocastica dei Neuroni}

\subsubsection{Il Problema dell'Input Realistico}

Il professore introduce una \textbf{transizione concettuale fondamentale}: fino a questo punto del corso, i neuroni modello sono stati studiati con correnti iniettate deterministiche (gradini, rampe, sinusoidi). Questa è una semplificazione utile per lo studio delle proprietà intrinseche, ma \textbf{non riflette ciò che accade nel cervello reale}, dove ogni neurone è immerso in una rete e riceve input attraverso migliaia di sinapsi.

Il passo successivo è chiedersi: qual è la natura dell'input che un neurone riceve quando è \textbf{embedded} in una rete reale di neuroni? La risposta, come il professore anticipa subito, è che questo input è di natura \textbf{stocastica}. Il neuroni non solo sono intrinsecamente stocastici per le proprietà dei canali ionici che li compongono, ma sono esposti a un ambiente — la rete — che produce fluttuazioni intrinseche e imprevedibili dell'input.


\subsubsection{L'Esperimento di Shadlen e Newsome: il Paradigma Sperimentale}

Per motivare l'approccio stocastico con dati reali, il professore descrive in dettaglio un esperimento diventato un paradigma delle neuroscienze: l'esperimento di \textbf{Shadlen e Newsome} (Newsome Laboratory, Baylor College of Medicine).

Il soggetto sperimentale è un macaco (\textit{macaca mulatta}) addestrato a eseguire un \textbf{compito di decisione percettiva}. L'animale è posizionato davanti a uno schermo su cui vengono presentati stimoli visivi. Lo stimolo è costituito da un campo di punti in movimento, all'interno di una regione circolare del campo visivo che corrisponde al campo recettivo dei neuroni registrati. Ogni punto si muove in una direzione casuale e, dopo pochi fotogrammi, scompare e viene sostituito da un nuovo punto. La \textbf{coerenza del moto} è il parametro sperimentale chiave: esso indica quale frazione dei punti si muove in una direzione prestabilita (destra o sinistra), mentre i rimanenti si muovono in direzioni casuali.

I compiti sperimentali sono i seguenti:
\begin{enumerate}
    \item L'animale fissa il proprio sguardo su un punto centrale (croce bianca) e non può muovere gli occhi.
    \item L'animale osserva con la visione periferica il campo di punti.
    \item Quando ha accumulato sufficiente evidenza, l'animale sposta la fissazione verso un bersaglio rosso (se percepisce moto verso destra) o verso un bersaglio verde (se percepisce moto verso sinistra).
\end{enumerate}

Questo è un \textbf{compito di decisione percettiva ocuomotoria} (\textit{perceptual decision task}), in cui la percezione di un segnale sensoriale debole deve essere integrata nel tempo per produrre una decisione binaria.

\subsubsection{L'Area MT: Registrazione Elettrofisiologica}

L'area corticale da cui si registra è la \textbf{MT} (area MT, \textit{Middle Temporal}), anche nota come V5. Essa si trova lungo la via visiva dorsale ("dove") che si estende dalla corteccia visiva primaria (V1) verso il lobo parietale. L'area MT è specializzata nel processamento del \textbf{moto visivo}: i suoi neuroni rispondono selettivamente alla direzione e alla velocità del moto degli stimoli visivi nel loro campo recettivo.

Il professore descrive la topografia cerebrale del macaco:
\begin{itemize}
    \item Il \textbf{solco centrale} separa lobo frontale e parietale.
    \item La via visiva dorsale sale dal lobo occipitale verso il parietale, passando per le aree V2, V3, MT (quinta stazione della via visiva dorsale).
    \item Gli elettrodi di tungsteno vengono impiantati nell'area MT per registrare l'attività neuronale \textit{in vivo}.
\end{itemize}

L'esperimento non registra il potenziale di membrana del singolo neurone — operazione impossibile in un animale sveglio e in comportamento. Ciò che viene registrato è il \textbf{segnale extracellulare locale}: il potenziale di campo locale (\textit{local field potential}, LFP) non filtrato, dal quale si identificano gli \textbf{spike} come eventi brevi e di grande ampiezza che superano la soglia di discriminazione. Il registratore memorizza la \textbf{sequenza temporale degli spike}: per ciascun trial $i$ e per ciascun evento $k$, viene registrata la coppia $(i, t_{ik})$, dove $t_{ik}$ è l'istante di emissione del $k$-esimo spike durante il trial $i$.


\subsection{Il Treno di Spike come Processo Stocastico}

\subsubsection{Il Raster Plot e la Variabilità Trial-to-Trial}

La visualizzazione standard della variabilità neuronale è il \textbf{raster plot} (o "dot raster"): ogni riga corrisponde a un trial, ogni punto nero nella riga indica l'istante di emissione di uno spike. Formalmente, se il neurone emette $K$ spike nel trial $j$, la riga $j$ del raster plot contiene $K$ punti alle coordinate temporali $t_{j1}, t_{j2}, \ldots, t_{jK}$.

Il professore mostra e discute in dettaglio la visualizzazione del raster plot dell'esperimento di Shadlen e Newsome. In una finestra temporale ristretta di 100 millisecondi, ingrandita dall'originale, è evidente che:
\begin{enumerate}
    \item Il numero di spike prodotti in quella finestra \textbf{varia da trial a trial}.
    \item Il \textbf{timing} degli spike non è riproducibile: lo spike prodotto al trial $j$ al tempo $t_{jk}$ non compare necessariamente al tempo $t_{(j+1)k}$ nel trial successivo.
    \item Non esiste alcun pattern regolare riconoscibile a occhio nudo.
\end{enumerate}

Questa variabilità è nominata \textbf{variabilità trial-to-trial} (\textit{trial-to-trial variability}). Essa è la firma sperimentale del carattere stocastico dell'attività neuronale.

Il treno di spike per un singolo trial $j$ può essere rappresentato matematicamente come una somma di \textbf{funzioni delta di Dirac}:

\begin{equation}
S_j(t) = \sum_{k} \delta(t - t_{jk})
\end{equation}

Questa rappresentazione cattura in modo completo il contenuto informativo del treno di spike: ogni spike è un evento puntuale che avviene all'istante $t_{jk}$, e la somma di delta descrive la sequenza temporale di tutti gli eventi.

\subsubsection{Il Peristimulus Time Histogram (PSTH)}

Per estrarre l'informazione sistematica dalla variabilità, si calcola il \textbf{PSTH} (\textit{Peristimulus Time Histogram}): si divide l'asse temporale in bin di durata $\Delta t$ (tipicamente 5 ms) e per ciascun bin si conta il numero di spike prodotti attraverso tutti i trial $N$, dividendo poi per $N \cdot \Delta t$ per ottenere una stima della \textbf{frequenza di firing istantanea} $\nu(t)$.

Formalmente, la frequenza di firing istantanea come media sul numero di trial $N$ è:

\begin{equation}
\nu(t) = \frac{1}{N} \sum_{j=1}^{N} \int_t^{t+\Delta t} S_j(t') \, dt' \cdot \frac{1}{\Delta t}
\end{equation}

oppure, in forma compatta:

\begin{equation}
\nu(t) \, \Delta t \approx \frac{1}{N} \sum_{j=1}^{N} \int_t^{t+\Delta t} S_j(t') \, dt'
\end{equation}

Il professore sottolinea un punto cruciale: la frequenza di firing istantanea così calcolata ha una \textbf{doppia interpretazione} che sarà di fondamentale importanza per tutta la lezione:

\begin{enumerate}
    \item \textbf{Interpretazione temporale:} è la frequenza media con cui un singolo neurone emette spike, mediata su molti trial identici.
    \item \textbf{Interpretazione spaziale (di popolazione):} è equivalente alla frequenza di firing di una \textbf{popolazione} di $N$ neuroni identici esposti allo stesso stimolo — ciascun trial è una realizzazione stocastica indipendente dello stesso processo. Se i $N$ neuroni sono statisticamente equivalenti e indipendenti, la media sull'ensemble di trial corrisponde alla media sull'ensemble di neuroni.
\end{enumerate}



\subsection{Frequenza di Firing Istantanea e Stima dal Raster Plot}

\subsubsection{Numero di Spike in una Finestra Temporale}

Il professore formalizza la costruzione del PSTH in modo rigoroso. Si definisce la variabile aleatoria:

\begin{equation}
n_j(t, \Delta t) = \int_t^{t+\Delta t} S_j(t') \, dt'
\end{equation}

ovvero il numero di spike prodotti dal neurone nel trial $j$ durante la finestra $[t, t+\Delta t]$. Questa è una variabile intera non negativa, tipicamente 0 o 1 per $\Delta t$ molto piccolo.

La stima empirica del tasso di firing istantaneo è:

\begin{equation}
\hat{\nu}(t) = \frac{1}{N \cdot \Delta t} \sum_{j=1}^{N} n_j(t, \Delta t)
\end{equation}

Il valore atteso di $n_j(t, \Delta t)$ su $N$ trial equivalenti (ipotesi di ergodicity o stazionarietà dello stato sperimentale) è:

\begin{equation}
\langle n(t, \Delta t) \rangle = \nu(t) \cdot \Delta t
\end{equation}

Il professore introduce la notazione $\nu(t)$ per il tasso di firing che può essere dipendente dal tempo — una scelta che anticipa il concetto di \textbf{processo di Poisson inomogeneo}.

\subsubsection{Varianza del Numero di Spike e Fluttuazioni di Dimensione Finita}

Una questione importante, con implicazioni dirette per le reti di neuroni, è la \textbf{varianza} di $\hat{\nu}(t)$. Il professore la calcola esplicitamente:

\begin{equation}
\text{Var}\left[\hat{\nu}(t)\right] = \text{Var}\left[\frac{1}{N\Delta t}\sum_{j=1}^N n_j\right] = \frac{1}{N^2 \Delta t^2} \sum_{j=1}^N \text{Var}[n_j]
\end{equation}

Assumendo che la produzione degli spike in trial diversi sia \textbf{statisticamente indipendente} (ipotesi di processi indipendenti), si può portare l'operatore varianza dentro la somma. Per una variabile binaria (0 o 1 nella finestra $\Delta t$), la varianza di $n_j$ è:

\begin{equation}
\text{Var}[n_j] = p(1-p) \approx p = \nu \cdot \Delta t
\end{equation}

dove $p = \nu \cdot \Delta t$ è la probabilità di emettere uno spike nella finestra e si usa l'approssimazione valida per $\Delta t$ piccolo (probabilità bassa). Quindi:

\begin{equation}
\text{Var}\left[\hat{\nu}(t)\right] \approx \frac{1}{N^2 \Delta t^2} \cdot N \cdot \nu \cdot \Delta t = \frac{\nu}{N \cdot \Delta t}
\end{equation}

e la deviazione standard del tasso stimato scala come:

\begin{equation}
\sigma_{\hat{\nu}} \sim \sqrt{\frac{\nu}{N \cdot \Delta t}}
\end{equation}

Questo risultato ha una conseguenza importante: le \textbf{fluttuazioni di dimensione finita} nel tasso di firing sono dell'ordine di $\sqrt{\nu / N}$ e decrescono come $1/\sqrt{N}$ all'aumentare del numero di neuroni (o di trial) considerati. In una rete di dimensione finita, queste fluttuazioni costituiscono una \textbf{sorgente intrinseca di rumore} che il neurone postsinaptico riceve come segnale di input.



\subsection{Tasso di Firing come Probabilità di Emissione dello Spike}

\subsubsection{Interpretazione Probabilistica del Tasso di Firing}

Il professore dimostra con rigore che il tasso di firing istantaneo $\nu(t)$ può essere interpretato come la \textbf{probabilità di emettere uno spike per unità di tempo}. Formalmente:

\begin{equation}
\langle n(t, \Delta t) \rangle = \nu(t) \cdot \Delta t = P(\text{uno spike in } [t, t+\Delta t])
\end{equation}

per $\Delta t \to 0$ (nell'approssimazione in cui la probabilità di due o più spike nella finestra è trascurabile, $P(n \geq 2) = O(\Delta t^2)$).

Questa interpretazione è fondamentale perché connette la grandezza osservabile sperimentalmente (il conteggio degli spike nel raster plot) con la struttura probabilistica sottostante del processo stocastico. Non si sta facendo un'assunzione ad hoc: si sta semplicemente riconoscendo che la media empirica su $N$ ripetizioni di un esperimento identico è un estimatore della probabilità dell'evento contato.

La probabilità $P(\text{spike in } \Delta t) = \nu(t) \Delta t$ è la \textbf{definizione operativa} del tasso di Poisson locale: il processo puntuale è completamente caratterizzato dal suo tasso istantaneo $\nu(t)$.

\subsubsection{Il Treno di Spike come Processo Puntuale Stocastico}

Il professore introduce la nozione formale di \textbf{processo puntuale stocastico} (\textit{stochastic point process}). Un processo puntuale è una sequenza di eventi (i tempi degli spike $\{t_k\}$) che occorrono su una linea temporale, caratterizzata statisticamente dalla distribuzione delle distanze tra eventi consecutivi (la \textbf{distribuzione degli intervalli inter-spike}) e dalla distribuzione del numero di eventi in intervalli temporali finiti.

La proprietà chiave richiesta per un processo di Poisson omogeneo è la \textbf{proprietà di Markov locale}: la probabilità di emettere uno spike nell'intervallo $[t, t+\Delta t]$ dipende soltanto dal tasso corrente $\nu$ e non dalla storia passata del neurone (assenza di memoria). Questa proprietà, denominata \textbf{proprietà di assenza di memoria} o \textbf{proprietà di rinnovamento}, è equivalente alla condizione che gli \textbf{intervalli inter-spike} siano \textbf{variabili aleatorie indipendenti e identicamente distribuite} (\textit{i.i.d.}).

Il processo che ha questa proprietà e tasso costante nel tempo è il \textbf{processo di Poisson omogeneo}, che sarà trattato in dettaglio nella sezione successiva.


\subsection{Il Processo di Poisson Omogeneo: Definizione e Proprietà}

\subsubsection{Definizione Assiomatica}

Il professore introduce il \textbf{processo di Poisson omogeneo} (o stazionario) con tasso $\nu$ tramite le seguenti proprietà fondamentali:

\begin{enumerate}
    \item \textbf{Stazionarietà:} la probabilità di osservare $n$ spike nell'intervallo $[t, t+T]$ dipende soltanto da $T$, non da $t$.
    \item \textbf{Indipendenza degli incrementi:} il numero di spike in intervalli temporali disgiunti è statisticamente indipendente.
    \item \textbf{Rarità degli eventi:} nell'approssimazione $\Delta t \to 0$, la probabilità di due o più spike in $[t, t+\Delta t]$ è trascurabile: $P(n \geq 2 \text{ in } \Delta t) = O(\Delta t^2)$.
    \item \textbf{Linearità del tasso:} $P(\text{esattamente uno spike in } [t, t+\Delta t]) = \nu \Delta t + O(\Delta t^2)$.
\end{enumerate}

Queste quattro proprietà sono le condizioni che definiscono univocamente il processo di Poisson. Dal punto di vista neuroscientistico, l'ipotesi che il treno di spike sia un processo di Poisson corrisponde al dire che il neurone \textbf{non ha memoria}: la probabilità di emettere il prossimo spike non dipende da quando è stato emesso l'ultimo spike.

\subsubsection{Il Contatore del Processo: $N(t)$}

Si definisce il processo di conteggio:

\begin{equation}
N(t) = \int_0^t S(t') \, dt'
\end{equation}

che conta il numero cumulativo di spike prodotti dall'istante $t=0$ fino al tempo $t$. Il processo $N(t)$ è un processo stocastico a valori interi non decrescenti, che aumenta di 1 ogni volta che si verifica uno spike. Il suo grafico — una funzione a gradini — è la rappresentazione temporale del processo puntuale.

Per un processo di Poisson omogeneo, il valor medio di $N(t)$ cresce linearmente nel tempo:

\begin{equation}
\langle N(t) \rangle = \nu \cdot t
\end{equation}

Il professore sottolinea che $N(t)$ è dipendente dal tempo, mentre il tasso $\nu$ è costante: la confusione tra queste due quantità è una possibile fonte di errore, come emerge da un'interazione con uno studente durante la lezione.


\subsubsection{Derivazione della Distribuzione}

Il professore deriva la probabilità di osservare esattamente $n$ spike nell'intervallo temporale $[0, T]$ per un processo di Poisson di tasso $\nu$.

Si divide l'intervallo $[0, T]$ in $M$ sottofinestre di durata $\Delta t = T/M$. Per $M \to \infty$ (e quindi $\Delta t \to 0$), ciascuna sottofinestra contiene al più uno spike con probabilità $p = \nu \Delta t$, oppure nessuno con probabilità $1-p$. Grazie all'indipendenza degli incrementi, il numero totale di spike in $[0, T]$ ha distribuzione \textbf{binomiale} $B(M, p)$, che nel limite $M \to \infty$ con $Mp = \nu T$ costante converge alla \textbf{distribuzione di Poisson}:

\begin{equation}
P(N(T) = n) = \frac{(\nu T)^n}{n!} e^{-\nu T}, \qquad n = 0, 1, 2, \ldots
\end{equation}

Si introduce la notazione $\lambda = \nu T$ (il numero medio di spike nell'intervallo $T$), cosicché:

\begin{equation}
P_\lambda(n) = \frac{\lambda^n}{n!} e^{-\lambda}
\end{equation}

Questa è la \textbf{distribuzione di Poisson} con parametro $\lambda$.

\subsubsection{Valore Atteso della Distribuzione di Poisson}

Il professore calcola il \textbf{valor medio} della distribuzione di Poisson:

\begin{equation}
\langle n \rangle = \sum_{n=0}^{\infty} n \cdot \frac{\lambda^n}{n!} e^{-\lambda}
\end{equation}

Si osserva che il termine $n=0$ non contribuisce. Per $n \geq 1$:

\begin{equation}
\langle n \rangle = e^{-\lambda} \sum_{n=1}^{\infty} \frac{\lambda^n}{(n-1)!} = e^{-\lambda} \cdot \lambda \sum_{n=1}^{\infty} \frac{\lambda^{n-1}}{(n-1)!}
\end{equation}

Eseguendo il cambio di indice $m = n-1$:

\begin{equation}
\langle n \rangle = e^{-\lambda} \cdot \lambda \sum_{m=0}^{\infty} \frac{\lambda^m}{m!} = e^{-\lambda} \cdot \lambda \cdot e^{\lambda} = \lambda = \nu T
\end{equation}

Il valore atteso del numero di spike nell'intervallo $T$ è quindi $\langle n \rangle = \nu T$, come atteso.

\subsubsection{Varianza della Distribuzione di Poisson e il Fano Factor}

Per calcolare la \textbf{varianza}, si usa la relazione $\text{Var}[n] = \langle n^2 \rangle - \langle n \rangle^2$. Il calcolo di $\langle n^2 \rangle$ (omesso nella derivazione completa ma ricavabile con gli stessi metodi) fornisce:

\begin{equation}
\langle n^2 \rangle = \lambda^2 + \lambda
\end{equation}

Quindi:

\begin{equation}
\text{Var}[n] = \langle n^2 \rangle - \langle n \rangle^2 = (\lambda^2 + \lambda) - \lambda^2 = \lambda
\end{equation}

Il risultato fondamentale è che \textbf{varianza e media coincidono} per la distribuzione di Poisson:

\begin{equation}
\text{Var}[n] = \langle n \rangle = \lambda = \nu T
\end{equation}

Questa proprietà è la firma caratteristica del processo di Poisson e permette di definire il \textbf{Fano factor} (o \textit{fattore di Fano}):

\begin{equation}
F = \frac{\text{Var}[N(T)]}{\langle N(T) \rangle} = \frac{\lambda}{\lambda} = 1
\end{equation}

Il Fano factor è \textbf{uguale a 1} per il processo di Poisson omogeneo, indipendentemente da $\nu$ e $T$. Questo costituisce un test empirico per verificare se un treno di spike reale è approssimabile da un processo di Poisson: si misura la coppia $(\langle n \rangle, \text{Var}[n])$ su molti trial/finestre temporali e si verifica se i punti si distribuiscono lungo la bisettrice del piano $\text{Var}[n]$ vs $\langle n \rangle$ (retta con pendenza 1 passante per l'origine).



\subsection{Intervallo Inter-Spike (ISI): Distribuzione Esponenziale}

\subsubsection{Definizione dell'ISI e il Processo di Rinnovamento}

L'\textbf{intervallo inter-spike} (\textit{Inter-Spike Interval}, ISI) è definito come la durata temporale tra due spike consecutivi:

\begin{equation}
\tau_k = t_{k+1} - t_k
\end{equation}

dove $t_k$ è il tempo del $k$-esimo spike. Per il processo di Poisson omogeneo, gli ISI sono variabili aleatorie \textbf{i.i.d.} (indipendenti e identicamente distribuite), il che definisce una classe di processi chiamati \textbf{processi di rinnovamento} (\textit{renewal processes}).

La condizione di rinnovamento ha un significato fisico preciso: il neurone "dimentica" tutto ciò che è accaduto prima dell'ultimo spike. In altre parole, il momento in cui avverrà il prossimo spike dipende soltanto dal tasso $\nu$ corrente e non dall'intera storia del neurone. Questo è esattamente la condizione di assenza di memoria (\textit{memoryless}) che caratterizza il processo di Poisson.

\subsubsection{Derivazione della Distribuzione Esponenziale}

Il professore mostra come derivare la distribuzione degli ISI a partire dalla definizione del processo di Poisson. La probabilità di \textbf{non} osservare spike nell'intervallo $[0, t]$ — ovvero che $\tau_k > t$ — corrisponde a $N(t) = 0$:

\begin{equation}
P(\tau > t) = P(N(t) = 0) = \frac{(\nu t)^0}{0!} e^{-\nu t} = e^{-\nu t}
\end{equation}

Quindi la \textbf{funzione di sopravvivenza} (complemento della CDF) degli ISI è:

\begin{equation}
P(\tau > t) = e^{-\nu t}
\end{equation}

La distribuzione di probabilità degli ISI è la \textbf{distribuzione esponenziale}. La densità di probabilità (PDF) si ottiene derivando rispetto a $t$:

\begin{equation}
\rho(\tau) = -\frac{d}{dt} P(\tau > t) = \nu \, e^{-\nu \tau}, \qquad \tau \geq 0
\end{equation}

Questa derivazione è chiara nel ragionamento fisico: se $\tau > t$ è la probabilità di non aver visto spike in $[0,t]$ — che sappiamo essere $e^{-\nu t}$ — allora la probabilità che il primo spike cada nell'intervallino $[t, t+d\tau]$ è:

\begin{equation}
P(\tau \in [t, t+d\tau]) = \underbrace{e^{-\nu t}}_{\text{nessuno spike in }[0,t]} \cdot \underbrace{\nu \, d\tau}_{\text{spike in }[t, t+d\tau]}
\end{equation}

Il professore illustra questa derivazione passo per passo sulla lavagna, ricollegandola alla formula:

\begin{equation}
P(\tau > t) = 1 - \int_0^t \rho(\tau') \, d\tau'
\end{equation}

da cui, derivando, si ottiene $\rho(t) = \nu e^{-\nu t}$.

\subsubsection{Momenti della Distribuzione Esponenziale}

Dalla distribuzione esponenziale $\rho(\tau) = \nu e^{-\nu \tau}$ si ricavano i momenti:

\textbf{Valore atteso (ISI medio):}

\begin{equation}
\langle \tau \rangle = \int_0^\infty \tau \, \nu e^{-\nu \tau} \, d\tau = \frac{1}{\nu}
\end{equation}

Il risultato $\langle \tau \rangle = 1/\nu$ è il \textbf{reciproco del tasso di firing} e non sorprende: se il neurone spara $\nu$ spike al secondo, l'intervallo medio tra spike consecutivi è $1/\nu$ secondi. Questo era già stato introdotto nella lezione sulla curva di risposta ingresso-uscita, dove si era definito il tasso di firing come l'inverso dell'ISI medio.

\textbf{Secondo momento:}

\begin{equation}
\langle \tau^2 \rangle = \int_0^\infty \tau^2 \, \nu e^{-\nu \tau} \, d\tau = \frac{2}{\nu^2}
\end{equation}

\textbf{Varianza:}

\begin{equation}
\text{Var}[\tau] = \langle \tau^2 \rangle - \langle \tau \rangle^2 = \frac{2}{\nu^2} - \frac{1}{\nu^2} = \frac{1}{\nu^2}
\end{equation}

\textbf{Deviazione standard:}

\begin{equation}
\sigma_\tau = \sqrt{\text{Var}[\tau]} = \frac{1}{\nu} = \langle \tau \rangle
\end{equation}

Il risultato notevole è che la \textbf{deviazione standard coincide con la media}: $\sigma_\tau = \langle \tau \rangle = 1/\nu$. Questo è una proprietà esclusiva della distribuzione esponenziale.



\subsubsection{Definizione del CV e Confronto con il Fano Factor}

Il professore introduce il \textbf{coefficiente di variazione} (\textit{coefficient of variation}, CV) degli ISI, definito come:

\begin{equation}
\text{CV} = \frac{\sigma_\tau}{\langle \tau \rangle} = \frac{\sqrt{\text{Var}[\tau]}}{\langle \tau \rangle}
\end{equation}

Il CV misura la \textbf{variabilità relativa} degli ISI rispetto alla loro media. Esso è adimensionale e permette di confrontare la regolarità del treno di spike tra neuroni con tassi di firing diversi.

Il professore sottolinea la differenza tra CV e Fano factor:
\begin{itemize}
    \item Il \textbf{Fano factor} $F = \text{Var}[N(T)] / \langle N(T) \rangle$ misura la variabilità del numero di spike \textbf{in una finestra temporale} $T$ (spazio di conteggio).
    \item Il \textbf{CV} misura la variabilità degli ISI nel \textbf{dominio temporale} (durate tra eventi consecutivi).
\end{itemize}

Per un processo di Poisson omogeneo, usando $\sigma_\tau = 1/\nu$ e $\langle \tau \rangle = 1/\nu$:

\begin{equation}
\text{CV}_{\text{Poisson}} = \frac{1/\nu}{1/\nu} = 1
\end{equation}

Il processo di Poisson ha dunque $\text{CV} = 1$. Questo è il \textbf{test del CV}: se un neurone reale ha $\text{CV} = 1$, il suo treno di spike è compatibile con un processo di Poisson omogeneo.


\begin{table}[htbp]
    \centering
    \renewcommand{\arraystretch}{1.3}
    \begin{tabularx}{\textwidth}{|c|c|X|}
    \hline
    \textbf{Valore del CV} & \textbf{Tipo di processo} & \textbf{Significato fisico} \\
    \hline
    $\text{CV} \ll 1$ & Molto regolare & Gli spike sono quasi periodici (es. pacemaker) \\
    \hline
    $\text{CV} = 1$ & Poisson & Indipendenza totale tra ISI successivi \\
    \hline
    $\text{CV} > 1$ & Super-Poisson & Bursting, clustering temporale degli spike \\
    \hline
    \end{tabularx}
\end{table}


Il valore $\text{CV} = 1$ separa i regimi regolare e irregolare ed è la firma del processo di Poisson.


\subsubsection{Limite del Processo Omogeneo}

Il processo di Poisson omogeneo è inadeguato per descrivere l'attività di neuroni reali in risposta a stimoli che variano nel tempo. La risposta di un neurone visivo a una barra in movimento, ad esempio, ha un tasso di firing che cambia dinamicamente con la posizione e la velocità dello stimolo. Allo stesso modo, nel paradigma di Shadlen e Newsome, il tasso di firing del neurone MT varia durante il corso del trial.

Per gestire questa non-stazionarietà, si introduce il \textbf{processo di Poisson inomogeneo} (o non-stazionario), in cui il tasso $\nu(t)$ è una funzione deterministica del tempo.


Il Fano factor rimane $F = 1$ anche per il processo inomogeneo, a condizione di considerare finestre temporali sufficientemente piccole da poter approssimare $\nu(t) \approx \text{const}$ all'interno della finestra.


\section{Il Test Empirico}

\subsubsection{L'Esperimento di Tolhurst}

Il professore descrive l'esperimento di \textbf{Tolhurst et al.} sulla corteccia visiva primaria (V1), uno dei lavori fondativi nello studio della variabilità neuronale corticale. In questo esperimento si misurava la risposta di neuroni V1 del gatto a barre in movimento, calcolando il Fano factor della distribuzione dei conteggi degli spike su molte ripetizioni dello stesso stimolo.

Il risultato fu che il Fano factor era tipicamente compreso tra 1 e 2, quindi \textbf{maggiore di 1}, indicando una variabilità superiore a quella attesa per un processo di Poisson. Questo suggeriva che il processo stocastico sottostante fosse più complesso di un semplice Poisson omogeneo.

\subsubsection{L'Esperimento di Softky e Koch: Il Test del CV}

Un contributo fondamentale fu quello di \textbf{Softky e Koch} (1993), che analizzarono i dati raccolti da Shadlen e Newsome sull'area MT del macaco e da Van Essen e collaboratori sulla corteccia visiva, per testare l'ipotesi del processo di Poisson usando il \textbf{coefficiente di variazione degli ISI}.

Il problema metodologico principale che Softky e Koch dovettero affrontare è la \textbf{non-stazionarietà del processo}: il PSTH mostra chiaramente che $\nu(t)$ varia nel corso del trial, e questo da solo può gonfiare il CV anche se il processo locale è Poisson. Per eliminare questo artefatto, la procedura adottata fu la seguente:

\begin{enumerate}
    \item Si calcola il PSTH (il profilo temporale del tasso di firing) con bin di durata $\Delta t_{\text{bin}}$.
    \item Si classificano i bin in \textbf{classi di frequenza}: per esempio, tutti i bin in cui il tasso stimato ricade nell'intervallo $[20, 30]$ Hz formano una classe, quelli nell'intervallo $[30, 40]$ Hz un'altra classe, ecc.
    \item Per \textbf{ciascuna classe}, si raccolgono gli ISI degli spike prodotti in quei bin e si costruisce l'istogramma degli ISI per quella classe.
    \item Da questo istogramma si calcolano $\langle \tau \rangle$ e $\sigma_\tau$, e quindi il CV.
\end{enumerate}

In questo modo, all'interno di ciascuna classe, il tasso è approssimativamente costante, e il CV stimato non è inflazionato dalla non-stazionarietà del processo.

\subsubsection{Risultati: Scatter Plot CV vs. ISI Medio}

Il professore descrive in dettaglio i risultati dello scatter plot di Softky e Koch, in cui sull'asse $x$ si pone $\langle \tau \rangle$ (l'ISI medio per quella classe) e sull'asse $y$ il CV. Ogni punto rappresenta una classe di frequenza per un determinato neurone.

I risultati mostrano:

\begin{enumerate}
    \item \textbf{Per ISI brevi (alto tasso di firing, $\langle \tau \rangle$ piccolo):} il CV è significativamente \textbf{inferiore a 1}. Questo indica che quando il neurone spara a frequenza elevata, gli ISI sono più regolari di quanto atteso per un Poisson. La spiegazione è il \textbf{periodo refrattario assoluto} $\tau_0$: non è fisicamente possibile produrre spike con ISI inferiore a $\tau_0 \approx 2-5$ ms. Per tassi elevati, l'ISI medio si avvicina a $\tau_0$, la distribuzione degli ISI è compressa e stretta intorno alla media, e il CV scende. In formula: quando $\langle \tau \rangle \approx \tau_0$, la distribuzione degli ISI è quasi unimodale con varianza piccola, e $\text{CV} \ll 1$.
    \item \textbf{Per ISI lunghi (basso tasso di firing, $\langle \tau \rangle$ grande):} il CV si avvicina a \textbf{1}, cioè al valore atteso per un Poisson. Questo è coerente con l'idea che per tassi bassi il periodo refrattario è ininfluente ($\langle \tau \rangle \gg \tau_0$) e la distribuzione degli ISI si avvicina alla esponenziale.
    \item \textbf{Differenza tra aree:} I neuroni dell'area MT (Shadlen-Newsome) mostrano un andamento con CV vicino a 1 per ISI lunghi, più coerente con il Poisson. I neuroni di V1 (Van Essen) mostrano maggiore variabilità.
\end{enumerate}

\begin{tcolorbox}[colback=gray!5, colframe=gray!40, arc=2mm, boxrule=0.5pt]
\textbf{Nota del docente:} Questo risultato è fondamentale. Ci dice che a basso tasso di firing il processo di Poisson è una buona approssimazione, ma ad alto tasso di firing il periodo refrattario introduce regolarità nel treno di spike. Questo non è un fallimento della modellistica Poisson: è semplicemente il segnale che dobbiamo aggiungere al modello la nozione di periodo refrattario.
\end{tcolorbox}




\subsubsection{La Distribuzione ISI dei Neuroni Corticali}

Tornando all'esperimento di Shadlen e Newsome, il professore discute la distribuzione degli ISI misurata per i neuroni dell'area MT. Contrariamente a quanto osservato per la giunzione neuromuscolare di Kuffler-Katz, la distribuzione degli ISI dei neuroni corticali \textbf{non è esponenziale}: ha invece una forma a campana con:
\begin{itemize}
    \item Un \textbf{dip} vicino a $\tau = 0$ (rarità degli ISI brevissimi, dovuta al periodo refrattario assoluto).
    \item Un picco a un valore positivo di $\tau$.
    \item Una coda lunga (decadimento lento per ISI grandi).
\end{itemize}

Questa forma è incompatibile con la distribuzione esponenziale di Poisson. La presenza del dip per $\tau \to 0$ indica che il neurone \textbf{non può fisicamente} produrre spike consecutivi separati di meno del periodo refrattario assoluto $\tau_0$.

\subsubsection{Il Fano Factor nei Neuroni Corticali}

Il professore mostra il grafico dello scatter plot varianza vs. media (test del Fano factor) per i neuroni di Shadlen e Newsome. I punti sono distribuiti intorno alla retta di Poisson (pendenza 1, Fano factor = 1) per bassi tassi di firing, ma si discostano sistematicamente sopra questa retta per alti tassi di firing.

Questo comportamento è interpretato come segue:
\begin{itemize}
    \item A \textbf{basso tasso di firing:} i neuroni sono ben approssimati da un processo di Poisson.
    \item Ad \textbf{alto tasso di firing:} la variabilità è maggiore del Poisson. Questo è dovuto alla non-stazionarietà del processo: le finestre temporali ad alto tasso corrispondono a periodi di forte eccitazione dove il tasso stesso varia rapidamente, gonfiando la varianza del conteggio.
\end{itemize}

La deviazione dal Poisson è quindi in larga misura un \textbf{artefatto di non-stazionarietà}, e non una proprietà intrinseca del generatore di spike.


\section{Sorgenti di Stocasticità: Canali, Sinapsi, Stato della Rete}

\subsubsection{Le Tre Sorgenti Principali}

Il professore enumera e discute in modo sistematico le \textbf{sorgenti di stocasticità} che operano a diversi livelli biologici di organizzazione:

\textbf{1. Stocasticità dei canali ionici (\textit{channel noise}):} Per un singolo canale, il processo di apertura e chiusura è un processo di Markov a due stati.

\textbf{2. Stocasticità della trasmissione sinaptica:} il rilascio di neurotrasmettitore dalla terminazione presinaptica è intrinsecamente stocastico. 

\textbf{3. Stocasticità dello stato della rete (\textit{network state noise}):}
Anche se i canali e le sinapsi fossero completamente deterministici, i neuroni in una rete riceverebbero un input fluttuante a causa della variabilità nell'attività degli altri neuroni della rete. 

\subsubsection{Il Rumore come Risorsa Computazionale}

Il professore introduce il concetto controintuitivo ma fondamentale che il \textbf{rumore non è soltanto un disturbo da minimizzare}, ma può avere un \textbf{ruolo costruttivo nella computazione neuronale}. Questo verrà approfondito nelle lezioni successive, ma il professore ne anticipa l'intuizione.

Il meccanismo fondamentale è il \textbf{risonanza stocastica} (\textit{stochastic resonance}): in alcuni regimi, l'aggiunta di rumore a un sistema non lineare può \textbf{aumentare} la capacità del sistema di rilevare segnali deboli. Questo fenomeno, originariamente studiato per spiegare la sensibilità straordinaria di certi animali a stimoli periodici deboli (come il ritmo delle ere glaciali nel caso dell'asteroide mammifero), è stato successivamente identificato anche nei sistemi neurali.

Il professore anticipa la serie di esperimenti di Segundo, Mainen e Sejnowski che dimostrano questo principio in modo diretto.


\subsubsection{L'Esperimento di Mainen e Sejnowski}

Il professore descrive l'esperimento cruciale di \textbf{Mainen e Sejnowski} (1995, \textit{Science}), che ha definitivamente mostrato il ruolo costruttivo del rumore nella precisione temporale degli spike.

Il protocollo:
\begin{enumerate}
    \item \textbf{Stimolo deterministico (corrente costante):} si inietta una corrente costante $I_0$ per circa 1 secondo. Il raster plot mostra spike molto irregolari da trial a trial, con jitter crescente nel tempo.
    \item \textbf{Stimolo stocastico (corrente fluttuante):} si inietta una corrente fluttuante $I(t) = \mu + \sigma \xi(t)$, dove $\xi(t)$ è rumore colorato (filtrato passa-basso), con la stessa realizzazione ripetuta identicamente da trial a trial. Il raster plot mostra spike con \textbf{alta precisione temporale}: gli spike tendono a occorrere agli stessi istanti da trial a trial, con jitter di pochi millisecondi.
\end{enumerate}

Il risultato è paradossale ma robusto: \textbf{l'input rumoroso produce spike più precisi} dell'input deterministico. La spiegazione fisica è la seguente:
\begin{itemize}
    \item Con input costante, il neurone vive in una condizione di \textbf{firing tonicamente regolare} in cui si accumulano le fluttuazioni intrinseche del sistema (da canali ionici, etc.), che amplificano il jitter nel tempo.
    \item Con input fluttuante, le grandi escursioni positive di $I(t)$ forniscono \textbf{impulsi forti} che resettano il sistema a ogni spike, eliminando l'accumulo di jitter. Ogni spike è "ancorato" a un picco del segnale di input, il che produce alta affidabilità temporale.
\end{itemize}

Il professore quantifica la \textbf{reliability} come il rapporto segnale-rumore del PSTH:

\begin{equation}
\text{Reliability} = \frac{\nu_{\text{peak}}}{\sigma_{\nu}}
\end{equation}

dove $\nu_{\text{peak}}$ è il valore di picco del PSTH e $\sigma_{\nu}$ è la sua variabilità standard. Mainen e Sejnowski trovarono che:
\begin{itemize}
    \item Per $\sigma_I = 0$ (corrente costante): la reliability è bassa.
    \item La reliability aumenta \textbf{monotonicamente} con la deviazione standard $\sigma_I$ dell'input.
\end{itemize}

Questo risultato suggerisce che i neuroni corticali siano \textbf{progettati evolutivamente} per operare con input fluttuanti naturalistici, e che la presenza di rumore ne aumenti l'efficacia nel trasdurre segnali.



\subsubsection{L'Esperimento di Shadlen e Newsome con Diversa Coerenza}

Il professore ritorna all'esperimento di Shadlen e Newsome per illustrare la \textbf{dipendenza dall'attività della rete}. Nei dati di questo esperimento, il professore separa le prove (trials) in base alla \textbf{coerenza} del moto dei punti:

\begin{itemize}
    \item \textbf{Coerenza 0\% (rumore puro):} i punti si muovono tutti in direzioni casuali. La risposta dell'area MT è irregolare ma con variabilità relativamente bassa del tasso di firing — simile al caso dell'input costante nell'esperimento di Mainen-Sejnowski.
    \item \textbf{Coerenza 50\%:} metà dei punti si muovono in una direzione definita. Intuitivamente ci si aspetterebbe una risposta più deterministica. Invece, il raster plot mostra una \textbf{maggiore irregolarità} del treno di spike.
    \item \textbf{Coerenza 100\% (segnale deterministico):} tutti i punti si muovono nella stessa direzione. Il compito è facilissimo — l'animale sa già dove muovere gli occhi. Eppure il treno di spike dei neuroni MT è \textbf{molto irregolare}, con CV vicino a 1.
\end{itemize}

\subsubsection{La Spiegazione: Input Top-Down dalla Corteccia Frontale}

Il professore spiega questa contraddizione apparente con la nozione di \textbf{input top-down}: il cervello non è un sistema feedforward a una via, ma una rete di reti. L'area MT non riceve solo input dalla corteccia visiva (via sensoriale ascendente), ma anche da aree frontali (corteccia prefrontale, area FEF, corteccia motoria oculare) attraverso proiezioni \textbf{discendenti} (\textit{top-down}).

Il ragionamento è il seguente:
\begin{itemize}
    \item A \textbf{coerenza 0\%}, il segnale sensoriale è ambiguo. La corteccia frontale non ha ancora preso una decisione e non invia input top-down a MT. La risposta di MT è guidata solo dal segnale sensoriale ascendente e ha variabilità moderata.
    \item A \textbf{coerenza 100\%}, la decisione è già stata presa quasi istantaneamente. I neuroni della corteccia prefrontale che codificano la decisione "vai a destra" si attivano e inviano una \textbf{scarica irregolare sostenuta} verso MT. Questo input top-down aggiuntivo, irregolare e ad alta conduttanza, rende il regime operativo di MT più simile al regime di alta conduttanza, aumentando la variabilità del treno di spike.
\end{itemize}

Questa interpretazione è coerente con il modello di \textbf{active inference} e con la visione del cervello come sistema generativo: le aree superiori non aspettano passivamente l'input sensoriale, ma inviano proiezioni discendenti che modulano attivamente l'elaborazione nelle aree sensoriali.

\begin{tcolorbox}[colback=gray!5, colframe=gray!40, arc=2mm, boxrule=0.5pt]
\textbf{Nota del docente:} Questo risultato è anche una splendida dimostrazione di come la variabilità neuronale non sia un semplice "rumore" da studiare in isolamento, ma dipenda dallo stato globale del network. Un neurone in isolamento e lo stesso neurone in un network in comportamento sono sistemi radicalmente diversi.
\end{tcolorbox}


\subsubsection{L'Esperimento di Greenwald: Optical Imaging e LFP}

Il professore conclude la lezione con un esperimento di grande impatto visivo, svolto nel laboratorio di \textbf{Greenwald} alla fine degli anni '90 e pubblicato su \textit{Science}. Questo esperimento utilizza contemporaneamente l'\textbf{imaging a voltaggio-sensibile} (\textit{voltage-sensitive dye imaging}, VSDI) e la registrazione del \textbf{potenziale di campo locale} (LFP) per caratterizzare la variabilità trial-to-trial nella corteccia visiva del macaco.

Il paradigma sperimentale: un macaco osserva una barra in movimento (stimolo per V1), e ogni trial viene presentato nelle stesse condizioni fisiche (stessa barra, stessa direzione, stessa velocità). Nonostante la ripetibilità dello stimolo, la risposta corticale è \textbf{drasticamente diversa tra trial diversi}.

Il VSDI mostra la depolarizzazione media dei dendriti dei neuroni in una regione della corteccia visiva. In due trial consecutivi del tutto identici (trial A e trial B), la risposta VSD è completamente diversa:
\begin{itemize}
    \item Nel trial A: ampia attivazione corticale, sostenuta nel tempo.
    \item Nel trial B: attivazione minima, che si estingue rapidamente.
\end{itemize}

\subsubsection{Correlazione tra LFP, VSD e Treno di Spike}

Nella stessa area corticale studiata con VSDI, viene inserito un elettrodo di tungsteno che registra simultaneamente:
\begin{itemize}
    \item Il \textbf{LFP} (potenziale di campo locale, componente a bassa frequenza): riflette l'attività sinaptica media dei neuroni intorno all'elettrodo.
    \item Il \textbf{treno di spike} (\textit{multi-unit activity}): il segnale LFP filtrato a alta frequenza permette di identificare gli spike delle unità neurali adiacenti.
\end{itemize}

Il confronto tra trial A e trial B mostra:
\begin{itemize}
    \item LFP trial A: oscillazioni ampie e prolungate.
    \item LFP trial B: oscillazioni molto più deboli.
    \item Spike trial A: attività sostenuta con tasso elevato.
    \item Spike trial B: attività sporadica con tasso basso.
\end{itemize}

Questo nonostante lo stimolo sia identico!

\subsubsection{La Media Non Cattura la Variabilità Trial-to-Trial}

Il punto cruciale che il professore enfatizza è che la \textbf{media} dei trial (il PSTH e il potenziale LFP medio) può mascherare la variabilità:

\begin{equation}
\langle S(t) \rangle = \frac{1}{N} \sum_{j=1}^N S_j(t)
\end{equation}

In media, il segnale può sembrare modulato in modo coerente dallo stimolo. Ma questa media nasconde il fatto che trial diversi sono in stati radicalmente diversi: alcuni sono altamente reattivi, altri praticamente silenti. La media non cattura questa \textbf{variabilità dello stato globale} del cervello.

Il professore illustra questo punto mostrando che la media del VSD e del LFP su molti trial produce una curva "media" che è reminiscente del segnale di un singolo trial attivo, ma che non riflette affatto la risposta del singolo trial inattivo.

\subsubsection{La Quarta Sorgente di Stocasticità: Lo Stato del Network}

Questo esperimento introduce quella che il professore chiama la \textbf{quarta sorgente di stocasticità}: non i canali, non le sinapsi, non le fluttuazioni di dimensione finita, ma lo \textbf{stato globale della rete}. Il cervello in ogni momento si trova in uno stato complesso determinato da:
\begin{itemize}
    \item Il livello di arousal (sonno, veglia, attenzione).
    \item Le aspettative dell'animale rispetto allo stimolo.
    \item L'engagement nel compito ("voglio la ricompensa" vs. "sono annoiato").
    \item Lo stato delle proiezioni top-down da corteccia frontale, parietale, ecc.
\end{itemize}

In condizioni di alta attenzione e alta motivazione, le proiezioni top-down sono attive e il regime di alta conduttanza è più marcato, producendo una risposta più variabile. In condizioni di disattenzione, le proiezioni top-down sono silenti e il regime è più simile al firing regolare.


\chapter{Lezione 12}

\section{Dal Modello a Conduttanza al Modello a Corrente}

\subsubsection{Il Problema della Non-Linearità}

In questa lezione il professore presenta una \textbf{semplificazione fondamentale} che sarà alla base di tutte le analisi delle lezioni successive: il passaggio dal \textbf{modello sinaptico a conduttanza} (\textit{conductance-based}) al \textbf{modello sinaptico a corrente} (\textit{current-based}). L'obiettivo è eliminare una non-linearità moltiplicativa che renderebbe l'analisi analitica della dinamica di rete notevolmente più complessa, senza perdere le proprietà qualitative essenziali del sistema.

Il punto di partenza è il modello \textbf{leaky integrate-and-fire a conduttanza}, in cui la corrente sinaptica che fluisce attraverso la membrana del neurone post-sinaptico non è semplicemente proporzionale ai neurotrasmettitori rilasciati, ma dipende in modo moltiplicativo dalla differenza tra il \textbf{potenziale di membrana} $V$ e il \textbf{potenziale di inversione} (o reversal potential) $E_\alpha$ della sinapsi:

\begin{equation}
I_{\text{syn}} = \sum_\alpha \sum_j g_{\alpha j}(t)\,(V - E_\alpha)
\end{equation}

Questa dipendenza da $V$ rende l'input \textbf{dipendente dallo stato} (\textit{state-dependent}): l'efficacia della sinapsi nel depolarizzare o iperpolarizzare il neurone post-sinaptico cambia istante per istante con il potenziale di membrana. In termini dinamici, il sistema ha almeno tre gradi di libertà ($V$, $G_E$, $G_I$) e i termini $G_E \cdot V$ e $G_I \cdot V$ sono \textbf{termini non-lineari} (prodotto di due variabili di stato). Questa non-linearità complica enormemente sia la simulazione numerica su larga scala sia, soprattutto, qualsiasi trattamento analitico delle popolazioni neurali.



\subsubsection{Il Neurone Post-Sinaptico e il suo Albero Dendritico}

Si considera un neurone corticale immerso in una rete. Il neurone è rappresentato con il suo \textbf{soma} (in azzurro nel disegno del professore), il suo \textbf{assone} che si prolunga verso destra e il suo \textbf{albero dendritico} con numerosi rami che si estendono dal soma. Lungo l'albero dendritico sono distribuite molte sinapsi: alcune di tipo \textbf{glutamatergico} (eccitatorie, indicate in rosso) e alcune di tipo \textbf{GABAergico} (inibitorie, indicate in verde).

\begin{tcolorbox}[colback=gray!5, colframe=gray!40, arc=2mm, boxrule=0.5pt]
\textbf{Nota del docente:} "In questo schema uso il colore verde per le sinapsi inibitorie e il rosso per le eccitatorie. Attenzione: questa è una scelta di rappresentazione, non c'è nulla di 'buono' o 'cattivo' nei due tipi di sinapsi. Entrambe sono fondamentali per il funzionamento della rete."
\end{tcolorbox}

\subsubsection{L'Etichetta $\alpha$: Classe di Appartenenza del Neurone Pre-Sinaptico}

Ogni neurone pre-sinaptico è identificato da \textbf{due etichette}. La prima è $\alpha$, che indica a quale \textbf{classe di popolazione} appartiene il neurone pre-sinaptico:

\begin{equation}
\alpha \in \{E,\, I\}
\end{equation}

\begin{itemize}
    \item $\alpha = E$: il neurone pre-sinaptico è \textbf{eccitatorio} (glutamatergico, indicato in rosso).
    \item $\alpha = I$: il neurone pre-sinaptico è \textbf{inibitorio} (GABAergico, indicato in verde).
\end{itemize}

\subsubsection{L'Etichetta $j$: Indice all'Interno della Classe}

La seconda etichetta è $j$, un \textbf{indice intero} che identifica il singolo neurone all'interno della classe $\alpha$:

\begin{equation}
j \in \{1, 2, \ldots, K_\alpha\}
\end{equation}

Il numero di neuroni pre-sinaptici in ciascuna classe, $K_\alpha$, dipende dalla classe: si usa $K_E$ per i pre-sinaptici eccitatori e $K_I$ per quelli inibitori. Il numero totale di contatti pre-sinaptici sul neurone studiato è:

\begin{equation}
K = K_E + K_I
\end{equation}

\subsubsection{I Numeri Biologici: Connettività Corticale}

I numeri tipici per un neurone corticale di mammifero sono fondamentali per giustificare le approssimazioni successive. Il professore li presenta sistematicamente:

% Ricordati di includere questi pacchetti nel preambolo del tuo documento:
% \usepackage{amsmath}
% \usepackage{nicefrac} 

\begin{table}[htbp]
    \centering
    \renewcommand{\arraystretch}{1.3}
    % Ho rimosso \textwidth e impostato l'allineamento a sinistra (l) e centrato (c)
    \begin{tabular}{|l|c|}
    \hline
    \textbf{Quantità} & \textbf{Valore tipico} \\
    \hline
    Numero totale di contatti pre-sinaptici $K$ & $\approx 20\,000 - 30\,000$ \\
    \hline
    Frazione di sinapsi eccitatorie & $\approx 80\%$  \\
    \hline
    Frazione di sinapsi inibitorie & $\approx 20\%$  \\
    \hline
    Tasso di firing medio (background) & $\approx 1\text{ Hz}$ \\
    \hline
    Tasso di firing di background totale & $\approx 20\text{ kHz}$ \\
    \hline
    \end{tabular}
    \caption{Caratteristiche tipiche dell'input di background per un neurone corticale.}
    \label{tab:neurone_background}
\end{table}

La \textbf{asimmetria 80/20} tra sinapsi eccitatorie e inibitorie è biologicamente fondamentale: siccome i neuroni inibitori sono in numero inferiore ma devono controllare efficacemente l'eccitazione, ciascuna sinapsi inibitoria deve avere una \textbf{forza 5 volte maggiore} rispetto alle eccitatorie. Questa compensazione garantisce il bilanciamento eccitatório/inibitório (\textit{E/I balance}) che caratterizza lo stato fisiologico della corteccia.


\subsubsection{Definizione Formale del Treno di Spike}

Il \textbf{treno di spike} prodotto dal neurone pre-sinaptico $(\alpha, j)$ è definito come una somma di \textbf{funzioni delta di Dirac} centrate sugli istanti di emissione degli spike:

\begin{equation}
S_{\alpha j}(t) = \sum_k \delta\!\left(t - t_{\alpha j}^{(k)}\right)
\end{equation}

dove $t_{\alpha j}^{(k)}$ è il tempo del $k$-esimo spike emesso dal neurone $(\alpha, j)$. Questo è lo stesso formalismo introdotto nella lezione precedente per descrivere i processi di Poisson.

\subsubsection{Il Ritardo Assonale}

Il professore precisa che nel contesto di rete, il tempo $t_{\alpha j}^{(k)}$ che compare nel treno di spike include già il \textbf{ritardo assonale} (\textit{axonal delay}) nella trasmissione del segnale dal neurone pre-sinaptico al neurone post-sinaptico. Formalmente:

\begin{equation}
t_{\alpha j}^{(k)} = t_{\alpha j}^{(k),\text{emissione}} + \Delta_{\alpha j}
\end{equation}

dove $t_{\alpha j}^{(k),\text{emissione}}$ è il vero istante di emissione dello spike nel soma del neurone $(\alpha, j)$ e $\Delta_{\alpha j}$ è il ritardo assonale, che dipende sia dal neurone pre-sinaptico sia (in generale) dal neurone post-sinaptico. In questa trattazione, poiché ci concentriamo su un unico neurone post-sinaptico, l'indice post-sinaptico non viene esplicitato, ma la sua esistenza è implicita.


\subsubsection{Definizione del Tasso di Firing Medio}

Il \textbf{tasso di firing medio} (o frequenza di firing media) del neurone $(\alpha, j)$ è definito come la media temporale del numero di spike prodotti per unità di tempo:

\begin{equation}
\bar{\nu}_{\alpha j} = \frac{1}{T} \int_0^T S_{\alpha j}(t)\, dt
\end{equation}

dove $T$ è un intervallo di osservazione sufficientemente lungo. La condizione su $T$ richiede che sia:
\begin{itemize}
    \item Molto più grande dei tempi caratteristici del sistema: il \textbf{periodo refrattario assoluto} $\tau_{\text{ref}} \approx 1-2\,\text{ms}$, la \textbf{costante di tempo di membrana} $\tau_m \approx 10-20\,\text{ms}$.
    \item Molto più grande dell'\textbf{intervallo inter-spike medio} $\langle \text{ISI} \rangle \approx 1/\bar{\nu}_{\alpha j}$.
\end{itemize}

Quest'ultima condizione implica $T \gg 100\,\text{ms}$ (poiché $\bar{\nu} \approx 1\,\text{Hz}$). In pratica, le stime sperimentali del tasso di firing vengono eseguite su finestre temporali dell'ordine di \textbf{minuti o decine di secondi}, ovvero dello stesso ordine della durata di un determinato stato cerebrale.

\subsubsection{Il Tasso di Background: $\sim$1 Hz}

Dai dati sperimentali emerge un risultato sorprendente e molto robusto: \textbf{in condizioni di riposo}, i neuroni corticali producono spike a un tasso medio di circa:

\begin{equation}
\bar{\nu}_{\text{bg}} \approx 1\,\text{Hz}
\end{equation}

Questo valore è notevolmente conservato tra diverse specie, diverse aree corticali e diversi stati cerebrali. Il professore menziona in particolare un \textbf{esperimento chiave}: un gruppo dello \textbf{NYU} (New York University) ha registrato l'attività di un grande numero di neuroni corticali in topi sia durante il sonno sia durante la veglia. Sorprendentemente, il tasso di firing medio temporale era quasi identico nelle due condizioni (circa 1 Hz), nonostante i pattern di attività fossero dramaticamente diversi:
\begin{itemize}
    \item \textbf{Durante il sonno} (onde lente): alternanza quasi-periodica tra stati di quiescenza e stati attivi (\textbf{slow wave activity}).
    \item \textbf{Durante la veglia}: attività \textbf{asincrona} e irregolare.
\end{itemize}

\subsubsection{Il Bombardamento Sinaptico: $\sim$20 kHz}

Avendo stabilito che $\bar{\nu}_{\alpha j} \approx 1\,\text{Hz}$ per ogni neurone pre-sinaptico, si può calcolare il \textbf{tasso totale di spike ricevuti} dal neurone post-sinaptico:

\begin{equation}
\nu_{\text{tot}} = \sum_\alpha \sum_{j=1}^{K_\alpha} \bar{\nu}_{\alpha j} \approx \bar{\nu}_{\text{bg}} \cdot K
\end{equation}

Con $\bar{\nu}_{\text{bg}} = 1\,\text{Hz}$ e $K = 20\,000$:

\begin{equation}
\nu_{\text{tot}} \approx 1\,\text{Hz} \times 20\,000 = 20\,\text{kHz}
\end{equation}

Questo è un numero \textbf{enorme}: il neurone riceve mediamente uno spike ogni:

\begin{equation}
\langle \Delta t_{\text{input}} \rangle = \frac{1}{\nu_{\text{tot}}} = \frac{1}{20\,000\,\text{Hz}} = 0.05\,\text{ms}
\end{equation}

Questo intervallo medio di $0.05\,\text{ms}$ è \textbf{molto più piccolo} sia del periodo refrattario ($\tau_{\text{ref}} \approx 1-2\,\text{ms}$) sia della costante di tempo di membrana ($\tau_m \approx 20\,\text{ms}$). Questo fatto è cruciale: il neurone, agendo come un \textbf{filtro passa-basso} con costante di tempo $\tau_m$, media su moltissimi eventi in ogni finestra temporale di ampiezza $\tau_m$. In una finestra di durata $\tau_m \approx 20\,\text{ms}$, il numero atteso di spike ricevuti è:

\begin{equation}
N_{\text{spikes}} \approx \nu_{\text{tot}} \cdot \tau_m \approx 20\,000\,\text{Hz} \times 0.02\,\text{s} = 400\,\text{spike}
\end{equation}

Il neurone media dunque circa \textbf{400 eventi sinaptic} in ogni costante di tempo di membrana.


\subsubsection{Separazione tra Input di Background e Input Selettivo}

Il tasso di firing totale ricevuto dal neurone non è sempre costante a $20\,\text{kHz}$. Il professore introduce una distinzione fondamentale tra due componenti dell'input:

\begin{enumerate}
    \item \textbf{Input di background}: l'attività spontanea della rete, che rimane circa costante a $\approx 20\,\text{kHz}$.
    \item \textbf{Input selettivo} (o perturbazione): una componente aggiuntiva prodotta da un sottoinsieme di neuroni pre-sinaptici che, in risposta a uno stimolo o a un compito cognitivo, aumentano il loro tasso di firing al di sopra del background.
\end{enumerate}

Formalmente, il tasso di firing istantaneo ricevuto è scritto come:

\begin{equation}
\nu_{\text{tot}}(t) = \nu_{\text{bg}} + \delta\nu(t)
\end{equation}

dove:
\begin{itemize}
    \item $\nu_{\text{bg}} \approx 20\,\text{kHz}$ è il \textbf{tasso di background} (fisso, determinato da $K \cdot \bar{\nu}_\text{bg}$).
    \item $\delta\nu(t)$ è la \textbf{perturbazione temporale} dovuta all'attività selettiva.
\end{itemize}

\subsubsection{Stima della Perturbazione Selettiva}

Il professore stima la perturbazione con un calcolo esplicito. Siano:
\begin{itemize}
    \item $f \approx 1\% = 0.01$: la \textbf{frazione di codifica} (\textit{coding fraction}) — la frazione dei neuroni pre-sinaptici che sono selettivamente attivi.
    \item $\bar{\nu}_{\text{sel}} \approx 10\,\text{Hz}$: il tasso di firing dei neuroni selettivi (al di sopra del background; il massimo fisiologico è $\bar{\nu}_{\text{max}} \approx 1/\tau_{\text{ref}} \approx 100-200\,\text{Hz}$).
\end{itemize}

Il contributo selettivo al tasso di input è:

\begin{equation}
\delta\nu_{\text{sel}} = f \cdot K \cdot \bar{\nu}_{\text{sel}} = 0.01 \times 20\,000 \times 10\,\text{Hz} = 2\,\text{kHz}
\end{equation}

La perturbazione relativa rispetto al background è:

\begin{equation}
\frac{\delta\nu_{\text{sel}}}{\nu_{\text{bg}}} = \frac{2\,\text{kHz}}{20\,\text{kHz}} = 10\%
\end{equation}

Questo è un \textbf{numero piccolo}: la perturbazione selettiva rappresenta circa il 10\% del background. Questo fatto giustifica un \textbf{approccio perturbativo}: si può linearizzare la dinamica attorno al punto di lavoro determinato dal background, trattando la perturbazione selettiva come un piccolo disturbo.


\subsection{Il Modello Conduttanza-Based: Equazione del Potenziale di Membrana}

\subsubsection{Scelta del Riferimento di Potenziale}

Prima di scrivere le equazioni, si sceglie come \textbf{riferimento di potenziale} il potenziale di equilibrio della membrana passiva $E_m = -65\,\text{mV}$ (il valore tipico del modello di Hodgkin-Huxley). Ponendo $E_m = 0$ per semplificare la notazione, i potenziali di inversione si ridefiniscono come:

\begin{table}[htbp]
    \centering
    \renewcommand{\arraystretch}{1.3}
    \begin{tabularx}{\textwidth}{|X|c|c|}
    \hline
    \textbf{Sinapsi} & \textbf{$E_\alpha$ originale} & \textbf{$E_\alpha$ nel nuovo riferimento} \\
    \hline
    Eccitatoria (AMPA/NMDA) & $0\,\text{mV}$ & $+65\,\text{mV}$ \\
    \hline
    Inibitoria (GABA) & $-75\,\text{mV}$ & $-10\,\text{mV}$ \\
    \hline
    Membrana passiva & $-65\,\text{mV}$ & $0\,\text{mV}$ \\
    \hline
    \end{tabularx}
\end{table}

Con questo riferimento, $V$ rappresenta il potenziale di membrana misurato rispetto all'equilibrio passivo.

\subsubsection{L'Equazione del Potenziale di Membrana Conduttanza-Based}

L'equazione che governa la dinamica sub-soglia del potenziale di membrana del neurone post-sinaptico nel modello conduttanza-based leaky integrate-and-fire è:

\begin{equation}
C_m \frac{dV}{dt} = -g_m V - \sum_{\alpha \in \{E,I\}} \sum_{j=1}^{K_\alpha} g_{\alpha j}(t)\,(V - E_\alpha)
\end{equation}

dove:
\begin{itemize}
    \item $C_m$: \textbf{capacità di membrana} (in $\text{F}$ o $\text{nF}$).
    \item $g_m$: \textbf{conduttanza di fuga} (leak conductance), responsabile del ritorno al potenziale di riposo.
    \item $g_{\alpha j}(t)$: \textbf{conduttanza sinaptica} del contatto $(\alpha, j)$ al tempo $t$ (in Siemens o nS).
    \item $E_\alpha$: potenziale di inversione della sinapsi di tipo $\alpha$ (nel riferimento scelto sopra).
\end{itemize}

Il termine $-g_m V$ è il termine di rilassamento passivo che riporta $V \to 0$ in assenza di input. Il termine $-g_{\alpha j}(t)\,(V - E_\alpha)$ rappresenta la corrente sinaptica di tipo $\alpha$ generata dalla sinapsi $j$-esima: essa è positiva (corrente depolarizzante) per le sinapsi eccitatorie (poiché $V < E_E = +65\,\text{mV}$ quasi sempre) e negativa (corrente iperpolarizzante) per le sinapsi inibitorie (poiché $V > E_I = -10\,\text{mV}$ quasi sempre in condizioni normali).



\subsubsection{La Dinamica del Primo Ordine per la Conduttanza}

Come stabilito nelle lezioni precedenti (in particolare nella lezione dedicata alla trasmissione sinaptica chimica), la conduttanza sinaptica $g_{\alpha j}(t)$ evolve secondo una \textbf{cinetica del primo ordine} guidata dall'arrivo degli spike pre-sinaptici. La semplificazione standard (che approssima la più realistica "funzione alpha" con un'esponenziale semplice) dà:

\begin{equation}
\frac{d g_{\alpha j}}{dt} = -\frac{g_{\alpha j}}{\tau_\alpha} + G_{\alpha j} \cdot S_{\alpha j}(t)
\end{equation}

dove:
\begin{itemize}
    \item $\tau_\alpha$: \textbf{costante di tempo sinaptica} per le sinapsi di tipo $\alpha$ (in ms).
    \item $G_{\alpha j}$: \textbf{conduttanza massima} (o "jump") associata all'arrivo di un singolo spike pre-sinaptico al contatto $(\alpha, j)$.
    \item $S_{\alpha j}(t) = \sum_k \delta(t - t_{\alpha j}^{(k)})$: il treno di spike pre-sinaptico.
\end{itemize}

I valori tipici delle costanti di tempo sinaptiche sono:

\begin{table}[htbp]
    \centering
    \renewcommand{\arraystretch}{1.3}
    \begin{tabularx}{\textwidth}{|X|l|l|}
    \hline
    \textbf{Tipo di sinapsi} & \textbf{Recettore} & \textbf{$\tau_\alpha$} \\
    \hline
    Eccitatoria rapida & AMPA & $2-5\,\text{ms}$ \\
    \hline
    Eccitatoria lenta & NMDA & $50-100\,\text{ms}$ \\
    \hline
    Inibitoria rapida & GABA$_A$ & $5-10\,\text{ms}$ \\
    \hline
    Inibitoria lenta & GABA$_B$ & $60-100\,\text{ms}$ \\
    \hline
    \end{tabularx}
\end{table}

\subsubsection{Soluzione in Assenza di Spike: Rilassamento Esponenziale}

In assenza di spike pre-sinaptici ($S_{\alpha j} = 0$), l'equazione si riduce a:

\begin{equation}
\frac{d g_{\alpha j}}{dt} = -\frac{g_{\alpha j}}{\tau_\alpha}
\end{equation}

con soluzione $g_{\alpha j}(t) = g_{\alpha j}(0)\,e^{-t/\tau_\alpha}$. La conduttanza rilassa esponenzialmente a zero con costante di tempo $\tau_\alpha$.

\subsubsection{L'Effetto di un Singolo Spike Pre-Sinaptico}

All'arrivo di un singolo spike pre-sinaptico all'istante $t^{(k)}$, la conduttanza riceve un \textbf{salto istantaneo} di ampiezza $G_{\alpha j}$:

\begin{equation}
g_{\alpha j}(t^{(k)+}) = g_{\alpha j}(t^{(k)-}) + G_{\alpha j}
\end{equation}

Seguito da un rilassamento esponenziale fino all'eventuale spike successivo. Questa è la ragione per cui questa cinetica è talvolta chiamata il "modello di conduttanza a funzione alpha": la risposta a uno spike singolo è un'esponenziale decrescente. La risposta in termine di \textbf{potenziale post-sinaptico eccitatorio} (EPSP) avrà una forma a campana asimmetrica (rapida salita, lenta discesa) determinata dalla convoluzione della risposta di conduttanza con il kernel RC della membrana.

\begin{tcolorbox}[colback=gray!5, colframe=gray!40, arc=2mm, boxrule=0.5pt]
\textbf{Nota del docente:} "Il modello vero dovrebbe avere una funzione alpha — rapida salita, lenta discesa — come forma della conduttanza. Qui stiamo semplificando ulteriormente usando un'esponenziale. L'importante è che la costante di tempo $\tau_\alpha$ e la quantità di carica trasmessa siano corretti."
\end{tcolorbox}

\subsubsection{Definizione della Conduttanza Totale per Popolazione}

Poiché l'equazione del potenziale di membrana dipende dalla \textbf{somma} sulle conduttanze sinaptiche di tutti i neuroni pre-sinaptici della stessa classe, è conveniente definire la \textbf{conduttanza sinaptica totale} per ciascuna classe:

\begin{equation}
G_\alpha(t) = \sum_{j=1}^{K_\alpha} g_{\alpha j}(t)
\end{equation}

Questa operazione è legittima perché l'equazione per la corrente sinaptica è \textbf{lineare} nelle conduttanze. La corrente sinaptica totale si riscrive come:

\begin{equation}
I_{\text{syn}} = -\sum_{\alpha \in \{E,I\}} G_\alpha(t)\,(V - E_\alpha)
\end{equation}

\subsubsection{Equazione per la Conduttanza Totale}

Poiché ogni singola conduttanza $g_{\alpha j}$ soddisfa l'equazione del primo ordine, per linearità la conduttanza totale $G_\alpha(t)$ soddisfa:

\begin{equation}
\frac{d G_\alpha}{dt} = -\frac{G_\alpha}{\tau_\alpha} + \sum_{j=1}^{K_\alpha} G_{\alpha j} \cdot S_{\alpha j}(t)
\end{equation}

Il sistema è ora ridotto a \textbf{tre gradi di libertà}: $V$, $G_E$, $G_I$. Il termine sorgente è la sovrapposizione di tutti i treni di spike pre-sinaptici, pesati per le rispettive conduttanze massime $G_{\alpha j}$.


\subsection{Decomposizione Media + Fluttuazione}

\subsubsection{La Decomposizione Perturbativa}

Il passo cruciale verso la linearizzazione è decomporre ogni variabile dinamica in una \textbf{componente media} (di stato stazionario) e una \textbf{fluttuazione temporale}:

\begin{equation}
G_\alpha(t) = \bar{G}_\alpha + \delta G_\alpha(t), \qquad \langle \delta G_\alpha \rangle_t = 0
\end{equation}

\begin{equation}
V(t) = \bar{V} + \delta V(t), \qquad \langle \delta V \rangle_t = 0
\end{equation}

dove la media temporale $\langle \cdot \rangle_t = \frac{1}{T}\int_0^T \cdot\, dt$ è presa su un intervallo $T$ sufficientemente lungo.

Analogamente, il treno di spike di ogni neurone pre-sinaptico viene decomposto:

\begin{equation}
S_{\alpha j}(t) = \bar{\nu}_{\alpha j} + \delta S_{\alpha j}(t), \qquad \langle \delta S_{\alpha j} \rangle_t = 0
\end{equation}

dove $\bar{\nu}_{\alpha j}$ è il tasso di firing medio e $\delta S_{\alpha j}$ è la componente fluttuante (media nulla).

La decomposizione è consistente se valgono le seguenti relazioni:

\begin{equation}
\frac{1}{T}\int_0^T G_\alpha(t)\, dt = \bar{G}_\alpha, \qquad \frac{1}{T}\int_0^T \delta G_\alpha(t)\, dt = 0
\end{equation}

Queste condizioni sono soddisfatte per definizione.


\subsubsection{Applicazione dell'Operatore di Media Temporale}

Per calcolare $\bar{G}_\alpha$, si applica la media temporale a entrambi i membri dell'equazione della conduttanza totale. Allo stato stazionario, la derivata temporale della media è nulla ($\langle \dot{G}_\alpha \rangle_t = 0$). L'equazione diventa:

\begin{equation}
0 = -\frac{\bar{G}_\alpha}{\tau_\alpha} + \sum_{j=1}^{K_\alpha} G_{\alpha j}\,\bar{\nu}_{\alpha j}
\end{equation}

\subsubsection{Il Risultato: Conduttanza di Stato Stazionario}

Risolvendo per $\bar{G}_\alpha$:

\begin{equation}
\boxed{\bar{G}_\alpha = \tau_\alpha \sum_{j=1}^{K_\alpha} G_{\alpha j}\,\bar{\nu}_{\alpha j}}
\end{equation}

Questo risultato ha una interpretazione fisica immediata: la conduttanza media è il prodotto della \textbf{costante di tempo sinaptica} per la \textbf{somma di tutti i contributi alla conduttanza prodotti da spike per unità di tempo}. Ogni neurone pre-sinaptico $j$ contribuisce con $G_{\alpha j} \cdot \bar{\nu}_{\alpha j}$ (conduttanza massima per salto moltiplicata per il numero di salti per secondo). La costante di tempo $\tau_\alpha$ appare perché ogni salto "vive" per un tempo $\tau_\alpha$ prima di rilassare.

Poiché $G_{\alpha j} > 0$ e $\bar{\nu}_{\alpha j} > 0$, si ha sempre $\bar{G}_\alpha > 0$: la conduttanza media è sempre positiva, come atteso fisicamente.

\subsubsection{Il Ruolo delle Fluttuazioni}

La componente fluttuante $\delta G_\alpha(t)$ è associata all'\textbf{attività selettiva} della rete: i neuroni che aumentano il loro tasso di firing in risposta a uno stimolo producono fluttuazioni di $G_\alpha$ attorno alla media. Come stimato nella sezione precedente, queste fluttuazioni sono di ordine \textbf{10\%} rispetto alla media, ovvero $|\delta G_\alpha| \approx 0.1\,\bar{G}_\alpha$.

\subsection{Linearizzazione del Potenziale di Membrana}

\subsubsection{Sostituzione della Decomposizione nell'Equazione del Moto}

Si sostituisce $G_\alpha = \bar{G}_\alpha + \delta G_\alpha$ e $V = \bar{V} + \delta V$ nell'equazione del moto per il potenziale di membrana:

\begin{equation}
C_m \frac{d(\bar{V} + \delta V)}{dt} = -g_m(\bar{V} + \delta V) - \sum_\alpha (\bar{G}_\alpha + \delta G_\alpha)\bigl((\bar{V} + \delta V) - E_\alpha\bigr)
\end{equation}

Espandendo il prodotto e raccogliendo per ordine perturbativo:

\begin{equation}
C_m \frac{d \delta V}{dt} = -g_m(\bar{V} + \delta V) - \sum_\alpha \bar{G}_\alpha(\bar{V} + \delta V - E_\alpha) - \sum_\alpha \delta G_\alpha(\bar{V} + \delta V - E_\alpha)
\end{equation}

\subsubsection{Identificazione dei Termini}

I termini dell'espansione si classificano come:
\begin{itemize}
    \item \textbf{Ordine zero} (solo quantità medie): $-g_m \bar{V} - \sum_\alpha \bar{G}_\alpha(\bar{V} - E_\alpha)$
    \item \textbf{Ordine lineare} in $\delta V$: $-g_m \delta V - \sum_\alpha \bar{G}_\alpha \delta V$
    \item \textbf{Ordine lineare} in $\delta G_\alpha$: $-\sum_\alpha \delta G_\alpha(\bar{V} - E_\alpha)$
    \item \textbf{Ordine quadratico} in $\delta G_\alpha \cdot \delta V$: questo termine viene \textbf{trascurato} nell'approssimazione perturbativa.
\end{itemize}

L'ordine quadratico $\delta G_\alpha \cdot \delta V$ è il prodotto di due quantità dell'ordine del 10\%: il suo contributo relativo è dell'ordine dell'\textbf{1\%} e può essere trascurato nell'approssimazione lineare.

\subsubsection{Equazione per il Valore Medio $\bar{V}$}

Applicando la media temporale all'equazione (dove $\langle d\delta V/dt \rangle = 0$ allo stato stazionario), si ottiene un'equazione per il potenziale medio $\bar{V}$:

\begin{equation}
0 = -g_m \bar{V} - \sum_\alpha \bar{G}_\alpha(\bar{V} - E_\alpha)
\end{equation}

Raccogliendo i termini proporzionali a $\bar{V}$:

\begin{equation}
0 = -\bar{V}\left(g_m + \sum_\alpha \bar{G}_\alpha\right) + \sum_\alpha \bar{G}_\alpha E_\alpha
\end{equation}

Si definisce la \textbf{conduttanza effettiva totale}:

\begin{equation}
\boxed{g_m' \equiv g_m + \sum_\alpha \bar{G}_\alpha = g_m + \bar{G}_E + \bar{G}_I}
\end{equation}

Con questa definizione, il potenziale medio di stato stazionario è:

\begin{equation}
\bar{V} = \frac{\sum_\alpha \bar{G}_\alpha E_\alpha}{g_m'} = \frac{\bar{G}_E E_E + \bar{G}_I E_I}{g_m'}
\end{equation}

Questo è un potenziale di equilibrio \textbf{effettivo}, che dipende dal bilanciamento tra eccitazione e inibizione: $\bar{G}_E E_E = \bar{G}_E \times (+65\,\text{mV})$ (positivo) vs. $\bar{G}_I E_I = \bar{G}_I \times (-10\,\text{mV})$ (negativo). L'equilibrio $\bar{V}$ può essere sopra o sotto il potenziale di riposo $V = 0$ a seconda della prevalenza relativa di eccitazione o inibizione. Si può definire un \textbf{potenziale di equilibrio effettivo}:

\begin{equation}
\boxed{E_m' = V \frac{\sum_\alpha \bar{G}_\alpha }{g_m'}}
\end{equation}


\subsection{Passaggio al Modello a Corrente}

\subsubsection{Equazione Linearizzata per la Fluttuazione del Potenziale}

Dopo aver separato il termine di ordine zero (che determina $\bar{V}$), l'equazione per la fluttuazione del potenziale $\delta V$ diventa:

\begin{equation}
C_m \frac{d\delta V}{dt} = -g_m' \,\delta V - \sum_\alpha \delta G_\alpha (\bar{V} - E_\alpha)
\end{equation}

Questa equazione ha una struttura fondamentale: la fluttuazione del potenziale di membrana è guidata da due termini:
\begin{enumerate}
    \item Un termine di rilassamento con conduttanza effettiva $g_m'$.
    \item Un termine sorgente proporzionale alle \textbf{fluttuazioni di conduttanza} $\delta G_\alpha$ pesate dalla \textbf{driving force di stato stazionario} $(\bar{V} - E_\alpha)$.
\end{enumerate}

In questa equazione la \textbf{non-linearità è scomparsa}: il potenziale $\delta V$ compare solo linearmente, senza più termini del tipo $G_\alpha \cdot V$.

\subsubsection{Riscrittura in Termini di V Totale}

Sommando l'equazione per $\bar{V}$ (ordine zero) e l'equazione per $\delta V$ (ordine lineare), e ricordando che $V = \bar{V} + \delta V$, si ottiene l'equazione per il potenziale totale:

\begin{equation}
C_m \frac{dV}{dt} = -g_m'(V - E_m') - \sum_\alpha \delta G_\alpha(\bar{V} - E_\alpha)
\end{equation}

Questa si può riscrivere in modo più compatto. Si noti che:
\begin{equation}
\delta G_\alpha (\bar{V} - E_\alpha) = (G_\alpha - \bar{G}_\alpha)(\bar{V} - E_\alpha)
\end{equation}

Quindi l'input sinaptico è proporzionale a $(G_\alpha - \bar{G}_\alpha)$ moltiplicato per una costante $(\bar{V} - E_\alpha)$. Questo è equivalente a un \textbf{modello a corrente} in cui l'efficacia sinaptica è riscalata.

\subsubsection{La Costante di Tempo Effettiva}

La costante di tempo di rilassamento del potenziale è ora:

\begin{equation}
\tau_m' = \frac{C_m}{g_m'}
\end{equation}

Sostituendo la definizione di $g_m'$:

\begin{equation}
\tau_m' = \frac{C_m}{g_m + \bar{G}_E + \bar{G}_I}
\end{equation}

Dividendo numeratore e denominatore per $g_m$:

\begin{equation}
\boxed{\tau_m' = \frac{\tau_m}{1 + \tau_m \dfrac{\bar{G}_E + \bar{G}_I}{C_m}}}
\end{equation}

dove $\tau_m = C_m/g_m$ è la costante di tempo passiva originale. Poiché tutte le conduttanze sono positive, si ha sempre $g_m' > g_m$, da cui:

\begin{equation}
\tau_m' < \tau_m
\end{equation}

\textbf{Risultato cruciale}: il bombardamento sinaptico \textbf{accorcia la costante di tempo di membrana}. Maggiore è l'attività della rete (maggiori sono $\bar{G}_E$ e $\bar{G}_I$), più piccola è $\tau_m'$ e più veloce è la risposta del neurone. Questo fenomeno è noto come \textbf{regime di alta conduttanza} (\textit{high-conductance state}) ed è una caratteristica fondamentale dei neuroni corticali in vivo.



\subsubsection{Introduzione della Corrente Sinaptica Linearizzata}

Il passo finale è introdurre come nuova variabile di stato la \textbf{corrente sinaptica linearizzata} per ciascuna classe $\alpha$:

\begin{equation}
I_\alpha(t) \equiv -G_\alpha(t)\,(\bar{V} - E_\alpha)
\end{equation}

Con questa definizione, l'equazione del potenziale di membrana assume la forma:

\begin{equation}
C_m \frac{dV}{dt} = -g_m' V + g_m' E_m' + \sum_\alpha I_\alpha(t)
\end{equation}

o, equivalentemente:

\begin{equation}
\tau_m' \frac{dV}{dt} = -(V - E_m') + \frac{\tau_m'}{C_m}\sum_\alpha I_\alpha(t)
\end{equation}

Questa è la forma standard del \textbf{leaky integrate-and-fire a corrente} (\textit{current-based LIF}).

\subsubsection{Equazione per la Corrente Sinaptica}

L'equazione del moto per la corrente sinaptica $I_\alpha(t)$ si ottiene moltiplicando l'equazione di $G_\alpha$ per $(\bar{V} - E_\alpha)$:

\begin{equation}
\frac{d I_\alpha}{dt} = -\frac{I_\alpha}{\tau_\alpha} + \sum_{j=1}^{K_\alpha} J_{\alpha j}\, S_{\alpha j}(t)
\end{equation}

dove si è definita l'\textbf{efficacia sinaptica} (synaptic efficacy):

\begin{equation}
\boxed{J_{\alpha j} = -G_{\alpha j}\,(\bar{V} - E_\alpha) = G_{\alpha j}\,(E_\alpha - \bar{V})}
\end{equation}

Il sistema completo di equazioni del modello current-based LIF è dunque:

\begin{equation}
\begin{cases}
\tau_m' \dfrac{dV}{dt} = -(V - E_m') + \dfrac{\tau_m'}{C_m}(I_E + I_I) \\[8pt]
\dfrac{d I_E}{dt} = -\dfrac{I_E}{\tau_E} + \displaystyle\sum_{j=1}^{K_E} J_{Ej}\, S_{Ej}(t) \\[8pt]
\dfrac{d I_I}{dt} = -\dfrac{I_I}{\tau_I} + \displaystyle\sum_{j=1}^{K_I} J_{Ij}\, S_{Ij}(t)
\end{cases}
\end{equation}

Le variabili di stato sono ora \textbf{tre}: $V$, $I_E$, $I_I$ — e il sistema è \textbf{lineare} nella dinamica di $I_E$ e $I_I$.

\subsubsection{Segno dell'Efficacia Sinaptica}

Il segno di $J_{\alpha j}$ determina se la sinapsi è \textbf{eccitatoria} o \textbf{inibitoria}:

\begin{itemize}
    \item \textbf{Sinapsi eccitatoria} ($\alpha = E$): $E_E = +65\,\text{mV}$, $\bar{V} \in [-10, 20]\,\text{mV}$ tipicamente, quindi $E_E - \bar{V} > 0$, e $G_{Ej} > 0$. Dunque $J_{Ej} > 0$: la corrente eccita il neurone post-sinaptico.
    \item \textbf{Sinapsi inibitoria} ($\alpha = I$): $E_I = -10\,\text{mV}$, e il potenziale medio $\bar{V}$ è in genere vicino o superiore a $E_I$, quindi $E_I - \bar{V} \leq 0$, e $G_{Ij} > 0$. Dunque $J_{Ij} \leq 0$: la corrente inibisce il neurone post-sinaptico.
\end{itemize}

\begin{tcolorbox}[colback=gray!5, colframe=gray!40, arc=2mm, boxrule=0.5pt]
\textbf{Nota del docente:} "La linearizzazione ci ha portato a una riscalatura dell'efficacia sinaptica: $J_{\alpha j} = G_{\alpha j} (E_\alpha - \bar{V})$. Il fattore $(\bar{V} - E_\alpha)$ è una costante che assorbe la 'driving force' media. Questo numero è misurabile sperimentalmente: è semplicemente l'ampiezza del salto di potenziale post-sinaptico quando arriva uno spike isolato (EPSP o IPSP)."
\end{tcolorbox}


\section{Riepilogo}

La lezione ha sviluppato una \textbf{gerarchia di semplificazioni} successive, ciascuna con le proprie ipotesi e limiti di validità:

\begin{table}[htbp]
    \centering
    \renewcommand{\arraystretch}{1.3}
    \begin{tabularx}{\textwidth}{|l|l|X|X|}
    \hline
    \textbf{Livello} & \textbf{Modello} & \textbf{Variabili di stato} & \textbf{Non-linearità} \\
    \hline
    1 & Conduttanza-based completo & $V, g_{\alpha j}$ (una per sinapsi, $\sim$20 000) & $g_{\alpha j} \cdot V$ \\
    \hline
    2 & Conduttanza totale & $V, G_E, G_I$ & $G_\alpha \cdot V$ \\
    \hline
    3 & Current-based linearizzato & $V, I_E, I_I$ & nessuna (lineare) \\
    \hline
    4 & White-noise LIF & $V$ & nessuna (lineare stocastico) \\
    \hline
    \end{tabularx}
\end{table}

Un neurone corticale è immerso in un bombardamento sinaptico massiccio ($\sim 20\,\text{kHz}$) proveniente da migliaia di neuroni pre-sinaptici. In questo contesto:

\begin{enumerate}
    \item La \textbf{conduttanza sinaptica media} modifica le proprietà passive del neurone: la costante di tempo si riduce ($\tau_m' < \tau_m$) e il potenziale di equilibrio si sposta ($E_m' \neq 0$).
    \item La \textbf{non-linearità} del modello conduttanza-based scompare nella linearizzazione perturbativa: il sistema si riduce a un \textbf{current-based LIF} con parametri dipendenti dall'attività di rete.
    \item Le \textbf{fluttuazioni dell'input} sono ben descritte da un \textbf{processo di Ornstein-Uhlenbeck}: un processo stazionario gaussiano di Markov con media $\mu_\alpha \tau_\alpha$ e varianza $\sigma_\alpha^2 \tau_\alpha/2$.
    \item Nel regime bilanciato della corteccia, la media netta dell'input è quasi zero (bilancio E/I) e le fluttuazioni dominano, producendo un potenziale di membrana irregolare con CV $\approx 1$.
\end{enumerate}

Queste sono le fondamenta su cui si costruirà la \textbf{teoria di campo medio} (\textit{mean-field theory}) per le reti neurali nelle lezioni successive.

\chapter{Lezione 13}


\section{Riepilogo}

\subsubsection{Il Punto di Partenza: Dinamica Sub-Soglia del Neurone LIF}

Questa lezione è la continuazione diretta della lezione precedente e costituisce il nucleo tecnico della derivazione dell'approssimazione diffusiva, che trasforma il modello discreto di salti sinaptici in un'equazione stocastica continua (l'equazione di Langevin), e introduce successivamente l'equazione di Fokker-Planck come strumento per descrivere la densità di probabilità del potenziale di membrana a livello di popolazione.

Il professore inizia richiamando il risultato principale della lezione 12. Il punto di partenza è un \textbf{neurone integrate-and-fire current-based} (LIF), il cui potenziale di membrana $V(t)$ evolve, al di sotto della soglia, secondo la dinamica di rilassamento:

\begin{equation}
\tau_m \frac{dV}{dt} = -V + I_E(t) + I_I(t)
\end{equation}

dove $\tau_m \approx 10$–$20\,\text{ms}$ è la \textbf{costante di tempo di membrana}. Le variabili $I_E(t)$ e $I_I(t)$ sono le correnti sinaptiche eccitatoria e inibitoria che confluiscono nel soma del neurone attraverso l'albero dendritico.

Le correnti sinaptiche, a loro volta, soddisfano ciascuna un'equazione del primo ordine guidata dall'arrivo degli spike pre-sinaptici:

\begin{equation}
\tau_\alpha \frac{dI_\alpha}{dt} = -I_\alpha + \tau_\alpha \sum_{j=1}^{K_\alpha} J_{\alpha j}\, S_{\alpha j}(t)
\end{equation}

dove $\alpha \in \{E, I\}$ identifica la classe (eccitatoria o inibitoria), $K_\alpha$ è il numero di contatti pre-sinaptici della classe $\alpha$, $J_{\alpha j}$ è l'\textbf{efficacia sinaptica} del contatto $j$-esimo della classe $\alpha$, e $S_{\alpha j}(t) = \sum_k \delta(t - t_{\alpha j}^{(k)})$ è il \textbf{treno di spike} del neurone pre-sinaptico $(\alpha, j)$.

Il sistema completo $(V, I_E, I_I)$ è dunque \textbf{tridimensionale}. La struttura accoppiata si traduce in una gerarchia: $V$ è forzato dalle correnti, le correnti a loro volta sono forzate dai treni di spike pre-sinaptici.

\subsubsection{Classificazione delle Sinapsi: Veloci e Lente}

Il professore ricorda che le sinapsi eccitatorie e inibitorie si dividono in due grandi categorie in base alle loro costanti di tempo di decadimento:

\begin{itemize}
    \item $\tau_\alpha^{\text{fast}} \approx 3\,\text{ms}$ (AMPA, GABA$_A$): \textbf{molto più piccolo} di $\tau_m \approx 10$–$20\,\text{ms}$.
    \item $\tau_\alpha^{\text{slow}} \approx 70$–$150\,\text{ms}$ (NMDA, GABA$_B$): \textbf{molto più grande} di $\tau_m$.
\end{itemize}

Questa separazione di scale temporali è fondamentale per giustificare la semplificazione successiva.

\section{Il Modello LIF a Salti:}

\subsubsection{Il Trattamento delle Sinapsi Lente: Quasi-Adiabaticità}

Le sinapsi \textbf{lente} (NMDA e GABA$_B$), avendo $\tau_\alpha \gg \tau_m$, filtrano completamente le fluttuazioni rapide del treno di spike. Il professore argomenta che in questo regime si può adottare una \textbf{approssimazione quasi-adiabatica}: la corrente lenta $I_\alpha^{\text{slow}}$ varia così lentamente rispetto alla dinamica di $V$ che può essere trattata come una corrente esterna lentamente variante, $I_{\alpha 0}(t)$, senza alcuna struttura di fluttuazioni su scale temporali $\sim \tau_m$. Questa componente viene quindi inclusa come un input esterno aggiuntivo ma, per la trattazione principale di questa lezione, viene messa da parte. La complessità aggiuntiva delle sinapsi lente sarà reintrodotta successivamente, quando si discuterà il limite di campo medio esteso.


\subsubsection{Il Limite Rapido: $\tau_\alpha \to 0$ e l'Approssimazione Delta}

Per le sinapsi \textbf{veloci} ($\tau_\alpha \ll \tau_m$), il professore mostra come si ottiene una drastica semplificazione. Si parte dall'equazione della corrente sinaptica veloce:

\begin{equation}
\tau_\alpha \frac{dI_\alpha}{dt} = -I_\alpha + \tau_\alpha \sum_{j=1}^{K_\alpha} J_{\alpha j}\, S_{\alpha j}(t)
\end{equation}

In questo regime, la corrente $I_\alpha$ raggiunge istantaneamente (su scale $\ll \tau_m$) il suo valore di stato stazionario forzato dal treno di spike. Matematicamente, si impone:

\begin{equation}
\tau_\alpha \frac{dI_\alpha}{dt} \approx 0
\end{equation}

che restituisce:

\begin{equation}
I_\alpha(t) \approx \tau_\alpha \sum_{j=1}^{K_\alpha} J_{\alpha j}\, S_{\alpha j}(t)
\end{equation}

Questa espressione mostra che la corrente diventa \textbf{proporzionale al treno di spike} in tempo reale. Sostituendo nella dinamica di $V$, si ottiene un modello \textbf{unidimensionale} in cui il potenziale di membrana è direttamente forzato dai treni di spike:

\begin{equation}
\tau_m \frac{dV}{dt} = -V + \sum_{\alpha \in \{E,I\}} \tau_\alpha \sum_{j=1}^{K_\alpha} J_{\alpha j}\, S_{\alpha j}(t)
\end{equation}

L'effetto fisico di questa approssimazione è il seguente. Quando arriva uno spike pre-sinaptico al tempo $t^{(k)}$, la corrente sinaptica produce un \textbf{salto istantaneo} nella corrente, seguito da un rilassamento esponenziale con scala $\tau_\alpha$. Nell'approssimazione $\tau_\alpha \to 0$ a \textbf{carica costante trasmessa}, questo diventa un \textbf{impulso di Dirac} $\delta(t - t^{(k)})$ moltiplicato per $\tau_\alpha J_{\alpha j}$, ovvero l'impulso trasferisce una quantità di carica $q_{\alpha j} = \tau_\alpha J_{\alpha j}$ in modo istantaneo. Il potenziale di membrana riceve quindi un \textbf{salto discontinuo} di ampiezza:

\begin{equation}
\Delta V = \frac{q_{\alpha j}}{C_m} \equiv J_{\alpha j}^{(V)}
\end{equation}

dove $J_{\alpha j}^{(V)}$ è l'\textbf{efficacia sinaptica ridefinita a livello di potenziale di membrana} (con dimensioni di Volt). Il professore chiarisce:

\begin{tcolorbox}[colback=gray!5, colframe=gray!40, arc=2mm, boxrule=0.5pt]
\textbf{Nota del docente:} "Quello che stiamo facendo è spostare il 'salto' dal livello della corrente al livello del potenziale di membrana. Nella corrente, il salto era $J_{\alpha j}$ e aveva dimensioni di corrente; ora il salto è $J_{\alpha j}^{(V)}$ e ha dimensioni di volt. Questo è il prezzo dell'approssimazione $\tau_\alpha \to 0$: i salti si spostano verso 'a monte'."
\end{tcolorbox}

\subsubsection{ l'Equazione Ridotta}

Con questa approssimazione, la dinamica diventa la seguente equazione scalare:

\begin{equation}
\tau_m \frac{dV}{dt} = -V + \sum_{\alpha \in \{E,I\}} J_\alpha \sum_{j=1}^{K_\alpha} S_{\alpha j}(t)
\end{equation}

dove si è introdotta l'efficacia sinaptica \textbf{media} (sotto l'assunzione di omogeneità della popolazione):

\begin{equation}
J_\alpha = \frac{1}{K_\alpha} \sum_{j=1}^{K_\alpha} J_{\alpha j}^{(V)}
\end{equation}

con la convenzione che $J_E > 0$ (sinapsi eccitatorie) e $J_I < 0$ (sinapsi inibitorie, poiché il GABA iperpolarizza).

Questo modello è noto nella letteratura come \textbf{modello di Stein} (introdotto negli anni '60 del XX secolo). Esso include esplicitamente i salti discontinui del potenziale di membrana ad ogni arrivo di spike, con l'ampiezza del salto determinata dall'efficacia sinaptica.

\section{Il Processo di Conteggio: Spike Presinaptici come Processo a Gradini}

\subsubsection{Definizione del Processo di Conteggio}

Per caratterizzare statisticamente l'input al neurone, il professore introduce formalmente il \textbf{processo di conteggio} $N_\alpha(t)$. Questo è il numero cumulativo di spike ricevuti dalla popolazione $\alpha$ nell'intervallo temporale $[0, t]$:

\begin{equation}
N_\alpha(t) = \sum_{j=1}^{K_\alpha} \sum_k \Theta(t - t_{\alpha j}^{(k)})
\end{equation}

dove $\Theta$ è la funzione di Heaviside. Il suo \textbf{incremento infinitesimale} è:

\begin{equation}
dN_\alpha(t) = N_\alpha(t + dt) - N_\alpha(t) = \sum_{j=1}^{K_\alpha} \int_t^{t+dt} S_{\alpha j}(t')\, dt'
\end{equation}

Questo incremento conta il numero di spike ricevuti dall'intera popolazione $\alpha$ nella finestra $[t, t+dt]$. È un processo a \textbf{valori interi non negativi}: $dN_\alpha \in \{0, 1, 2, \ldots\}$.

\subsubsection{Struttura a Gradini del Processo di Conteggio}

La realizzazione tipica di $N_\alpha(t)$ è un \textbf{processo a gradini}: piatto per lunghi tratti, con salti unitari in corrispondenza di ogni spike ricevuto. Ciascun salto si verifica in un istante che dipende dalla storia dell'attività di tutta la rete. Il professore descrive due esempi:

\begin{itemize}
    \item Il processo $N_E(t)$, associato agli spike eccitatori, ha salti verso l'alto (contributi positivi a $V$).
    \item Il processo $N_I(t)$, associato agli spike inibitori, ha salti verso il basso (contributi negativi a $V$).
\end{itemize}

L'equazione del modello di Stein può ora essere riscritta in forma integrale come:

\begin{equation}
V(t + dt) - V(t) = -\frac{V(t)}{\tau_m}\, dt + J_E\, dN_E(t) + J_I\, dN_I(t)
\end{equation}


\section{Regime di Sparsità: Condizioni per l'Approssimazione Diffusiva}

\subsubsection{I Numeri Biologici Cruciali}

Il professore introduce ora i numeri biologici che giustificano l'approssimazione che seguirà. I dati fondamentali sono:

\begin{table}[htbp]
    \centering
    \renewcommand{\arraystretch}{1.3}
    \begin{tabularx}{\textwidth}{|X|c|}
    \hline
    \textbf{Quantità} & \textbf{Valore tipico} \\
    \hline
    Numero totale di contatti pre-sinaptici $K = K_E + K_I$ & $\approx 20\,000$–$30\,000$ \\
    \hline
    Frazione eccitatoria ($K_E / K$) & $\approx 80\%$ \\
    \hline
    Frazione inibitoria ($K_I / K$) & $\approx 20\%$ \\
    \hline
    Tasso di firing medio di background $\bar{\nu}_\alpha$ & $\approx 1\,\text{Hz}$ \\
    \hline
    Input totale: $K \cdot \bar{\nu}$ & $\approx 20\,\text{kHz}$ \\
    \hline
    Finestra temporale di integrazione $\tau_m$ & $\approx 20\,\text{ms}$ \\
    \hline
    \end{tabularx}
\end{table}

\subsubsection{Sparsità del Singolo Neurone Pre-Sinaptico}

La probabilità che un singolo neurone pre-sinaptico emetta uno spike nella finestra $[t, t + \tau_m]$ è:

\begin{equation}
p_{\text{spike}} = \bar{\nu}_\alpha \cdot \tau_m \approx 1\,\text{Hz} \times 0.02\,\text{s} = 0.02
\end{equation}

Questo numero è \textbf{piccolo}: ciascun neurone pre-sinaptico emette, in media, uno spike ogni $1/\bar{\nu} = 1\,\text{s}$, che è molto maggiore di $\tau_m = 20\,\text{ms}$. Il treno di spike del singolo neurone è quindi \textbf{sparso} nella finestra di integrazione del neurone post-sinaptico. Il professore formalizza questa condizione:

\begin{equation}
\bar{\nu}_{\alpha j} \cdot \tau_m \ll 1 \qquad \forall\, (\alpha, j)
\end{equation}

\subsubsection{La Somma di Molti Processi Sparsi e Indipendenti}

Il processo $N_\alpha(t)$ è la somma di $K_\alpha \approx 10\,000$–$20\,000$ processi di punto individuali, ciascuno sparso. Il professore evidenzia che siamo in presenza di \textbf{due condizioni simultanee}:

\begin{enumerate}
    \item \textbf{Numero grande}: $K_\alpha \to \infty$ (limite termodinamico della connettività).
    \item \textbf{Sparsità}: il singolo processo di punto ha probabilità di emissione $\to 0$ in ogni finestra infinitesima.
\end{enumerate}

Questo è esattamente il regime in cui si applica un \textbf{teorema limite per processi di punto}, analogo al Teorema del Limite Centrale per variabili casuali. Nella letteratura matematica della teoria dei processi stocastici, questo risultato è il \textbf{Teorema di Griglioni} (Griglioni limit theorem) per la superposizione di processi di punto.



\subsection{Il Teorema di Griglioni: Convergenza a un Processo di Poisson}

Il professore descrive il \textbf{Teorema di Griglioni} (o Griglioni's limit theorem), che è il teorema di convergenza per la superposizione di processi di punto. Il teorema afferma quanto segue:

\textbf{Teorema (Griglioni):} Consideriamo una famiglia di $n$ processi di punto, con $n \to \infty$. Ciascun processo è un \textbf{processo di rinnovamento} (renewal process) — ovvero gli interspike interval sono indipendenti e identicamente distribuiti, ma con distribuzione arbitraria (non necessariamente esponenziale). I processi sono tra loro \textbf{indipendenti}. Si richiede che siano \textbf{infinitesimali} nel seguente senso: la probabilità che un singolo processo emetta almeno uno spike nell'intervallo $[0, t]$ tende a zero quando $n \to \infty$, mentre la somma di queste probabilità rimane finita e pari a $\Lambda(0, t)$. Allora la superposizione dei $n$ processi converge \textbf{in distribuzione} a un \textbf{processo di Poisson non omogeneo} con funzione di intensità $\lambda(t)$.

La funzione di intensità $\lambda(t)$ è semplicemente la somma dei tassi di firing istantanei:

\begin{equation}
\lambda(t) = \sum_{j=1}^{K_\alpha} \nu_{\alpha j}(t)
\end{equation}

Per i treni di spike delle popolazioni eccitatoria e inibitoria, sotto l'assunzione di stazionarietà $\nu_{\alpha j}(t) = \bar{\nu}_\alpha$ per ogni $j$, si ottiene:

\begin{equation}
\lambda_\alpha = K_\alpha \bar{\nu}_\alpha
\end{equation}

\subsubsection{Il Processo di Poisson Non Omogeneo}

Per trattare il caso generale (treni di spike non stazionari), il professore introduce formalmente il \textbf{processo di Poisson non omogeneo} con intensità $\lambda(t)$.

La \textbf{funzione di intensità cumulata} (o leading function) è:

\begin{equation}
\Lambda(s, t) = \int_s^t \lambda(u)\, du
\end{equation}

La probabilità di osservare esattamente $k$ spike nell'intervallo $[s, t]$ è data dalla distribuzione di Poisson con parametro $\Lambda(s, t)$:

\begin{equation}
P(N(s,t) = k) = \frac{[\Lambda(s,t)]^k}{k!}\, e^{-\Lambda(s,t)}
\end{equation}

Il significato di $\Lambda(s, t)$ è intuitivo: è il \textbf{numero medio atteso di spike} prodotti dall'intera popolazione nell'intervallo $[s, t]$, tenendo conto di una possibile variazione temporale del tasso di firing.

\subsubsection{La Convergenza al Processo di Poisson: Argomento Fisico}

Il professore presenta anche una lettura fisica della convergenza. Immaginiamo di passare da una famiglia di $n_1$ neuroni a una famiglia di $n_2 > n_1$ neuroni. Vogliamo mantenere costante il numero totale di spike prodotti per unità di tempo ($K_\alpha \bar{\nu}_\alpha = \text{cost}$). Per farlo, mentre aumentiamo il numero di neuroni, dobbiamo \textbf{scalare in basso} il loro tasso di firing individuale:

\begin{equation}
\nu_{\alpha j}(n) = \frac{K_\alpha \bar{\nu}_\alpha}{n}
\end{equation}

Così, nel limite $n \to \infty$, ogni singolo neurone diventa infinitamente sparso ($\nu_{\alpha j} \to 0$), ma la loro somma produce un flusso costante di spike. Il professore mostra che in questo limite la statistica dell'interspike interval individuale \textbf{non conta}: qualunque sia la distribuzione degli ISI del singolo neurone (esponenziale, log-normale, gamma, \dots), la superposizione converge sempre a un processo di Poisson. Questo è esattamente l'analogo, per i processi di punto, del Teorema del Limite Centrale per variabili casuali.


\section{Momenti Infinitesimali}

\subsubsection{I Momenti del Processo di Conteggio di Poisson}

Stabilito che $N_\alpha(t)$ è (con buona approssimazione) un processo di Poisson con intensità $\lambda_\alpha(t) = K_\alpha \nu_\alpha(t)$, il professore vuole ora caratterizzare statisticamente i momenti del processo di conteggio nell'intervallo infinitesimale $[t, t + dt]$.

Si definisce:

\begin{equation}
\mathbb{E}[dN_\alpha(t)^n]
\end{equation}

come l'$n$-esimo momento del numero di spike ricevuti in $[t, t+dt]$. Per un processo di Poisson con parametro $\Lambda = \lambda_\alpha(t)\, dt$, si ha:

\begin{equation}
\mathbb{E}[dN_\alpha^n] = \sum_{k=0}^{\infty} k^n \frac{\Lambda^k}{k!} e^{-\Lambda}
\end{equation}

Per calcolare i momenti a questo livello infinitesimale, è fondamentale considerare solo i termini \textbf{al primo ordine in $dt$}, poiché i termini $O(dt^2)$ vanno a zero nel limite $dt \to 0$ più rapidamente del necessario. Il professore calcola esplicitamente:

\begin{itemize}
    \item Termine $k=0$: contribuisce $0$ a qualsiasi momento.
    \item Termine $k=1$: contribuisce $1^n \cdot \Lambda \cdot e^{-\Lambda} \approx \Lambda = \lambda_\alpha(t)\, dt$ per $\Lambda \to 0$.
    \item Termine $k=2$: contribuisce $2^n \cdot \frac{\Lambda^2}{2} \cdot e^{-\Lambda} \approx O(dt^2)$.
\end{itemize}

Il risultato fondamentale è:

\begin{equation}
\mathbb{E}[dN_\alpha^n] = \lambda_\alpha(t)\, dt + O(dt^2) = K_\alpha \nu_\alpha(t)\, dt + O(dt^2)
\end{equation}

\textbf{Tutti i momenti dell'incremento di Poisson, a leading order in $dt$, sono uguali al primo momento.} Questa è la proprietà caratteristica del processo di Poisson.

\subsubsection{La Carica Sinaptica Consegnata in dt}

L'incremento del potenziale di membrana dovuto agli spike della classe $\alpha$ nella finestra $[t, t+dt]$ è:

\begin{equation}
dV_\alpha = J_\alpha\, dN_\alpha(t)
\end{equation}

La \textbf{carica sinaptica totale} consegnata in $[t, t+dt]$ è $Q_\alpha = J_E\, dN_E + J_I\, dN_I$. Per caratterizzare la distribuzione di $Q$, il professore calcola i momenti:

\begin{equation}
\mathbb{E}[Q^n] = \mathbb{E}[(J_E\, dN_E + J_I\, dN_I)^n]
\end{equation}

Poiché $N_E$ e $N_I$ sono \textbf{indipendenti} (spike eccitatori e inibitori sono processi distinti), si ha:

\begin{equation}
\mathbb{E}[Q^n] = \sum_{\alpha} J_\alpha^n \cdot \mathbb{E}[dN_\alpha^n] + \text{termini misti} = \sum_\alpha J_\alpha^n K_\alpha \nu_\alpha\, dt + O(dt^2)
\end{equation}

I termini misti coinvolgono prodotti del tipo $dN_E \cdot dN_I$, che sono $O(dt^2)$ e spariscono nel limite $dt \to 0$.



\subsection{Derivazione della Media Infinitesimale}

Il professore definisce il \textbf{momento infinitesimale di primo ordine} (o \textbf{drift infinitesimale}) come il limite del rapporto tra il valor medio dell'incremento di $V$ e l'ampiezza temporale $dt$:

\begin{equation}
\mu(V, t) \equiv \lim_{dt \to 0} \frac{\mathbb{E}[dV \mid V(t) = V]}{dt}
\end{equation}

L'incremento del potenziale di membrana in $dt$ è:

\begin{equation}
dV = -\frac{V}{\tau_m}\, dt + J_E\, dN_E + J_I\, dN_I
\end{equation}

Prendendo il valor medio e usando $\mathbb{E}[dN_\alpha] = K_\alpha \nu_\alpha\, dt$:

\begin{equation}
\mathbb{E}[dV] = -\frac{V}{\tau_m}\, dt + J_E\, K_E \nu_E\, dt + J_I\, K_I \nu_I\, dt
\end{equation}

Dividendo per $dt$ e passando al limite:

\begin{equation}
\mu(V, t) = -\frac{V}{\tau_m} + J_E K_E \nu_E(t) + J_I K_I \nu_I(t)
\end{equation}

Il termine $\mu$ è la somma di tre contributi:
\begin{enumerate}
    \item \textbf{Termine di rilassamento}: $-V/\tau_m$ che riporta il potenziale verso zero.
    \item \textbf{Input eccitatorio medio}: $J_E K_E \nu_E$, sempre positivo.
    \item \textbf{Input inibitorio medio}: $J_I K_I \nu_I$, sempre negativo (poiché $J_I < 0$).
\end{enumerate}

Definendo la \textbf{corrente di drift media} $\mu_0$:

\begin{equation}
\mu_0(t) \equiv J_E K_E \nu_E(t) + J_I K_I \nu_I(t)
\end{equation}

si può riscrivere in forma compatta:

\begin{equation}
\mu(V, t) = -\frac{V}{\tau_m} + \mu_0(t)
\end{equation}

In regime stazionario ($\nu_\alpha = \text{cost}$), la media infinitesimale $\mu$ è una funzione lineare di $V$ con una componente costante $\mu_0$. Il \textbf{punto fisso} del drift (dove $\mu = 0$) è il potenziale di equilibrio effettivo:

\begin{equation}
V^* = \tau_m \mu_0 = \tau_m (J_E K_E \nu_E + J_I K_I \nu_I)
\end{equation}

Questo è il potenziale medio attorno al quale il sistema fluttua sotto l'azione del rumore sinaptico.

\subsubsection{Interpretazione Fisica della Media}

Il professore si sofferma sull'interpretazione fisica. Il termine $\mu_0$ è l'\textbf{input sinaptico medio}: è la corrente che un neurone "vedrebbe" se ricevesse spike a un tasso costante pari alla media, senza fluttuazioni. In assenza di fluttuazioni, il potenziale si stabilizzerebbe in $V^* = \tau_m \mu_0$.

\begin{tcolorbox}[colback=gray!5, colframe=gray!40, arc=2mm, boxrule=0.5pt]
\textbf{Nota del docente:} "Se non ci fosse rumore, il potenziale di membrana si stabilizzerebbe al potenziale effettivo $V^*$ e il neurone emetterebbe uno spike solo se $V^*$ superasse la soglia $V_{\text{thr}}$. La presenza del rumore è ciò che rende la dinamica ricca: anche se $V^* < V_{\text{thr}}$, il neurone può comunque sparare grazie alle fluttuazioni. Ed è questo il punto centrale del modello che stiamo costruendo."
\end{tcolorbox}


\subsection{Derivazione della Varianza Infinitesimale}



Il \textbf{momento infinitesimale di secondo ordine} (o \textbf{coefficiente di diffusione infinitesimale}) è definito come:

\begin{equation}
\sigma^2(V, t) \equiv \lim_{dt \to 0} \frac{\mathbb{E}[(dV)^2 \mid V(t) = V]}{dt}
\end{equation}

Per calcolarlo, si scrive:

\begin{equation}
(dV)^2 = \left(-\frac{V}{\tau_m}\, dt + J_E\, dN_E + J_I\, dN_I\right)^2
\end{equation}

Espandendo il quadrato e tenendo solo i termini $O(dt)$ (poiché il termine $(-V/\tau_m\, dt)^2 = O(dt^2)$ e i termini incrociati tra il drift e $dN_\alpha$ sono $O(dt^{3/2})$ in senso stocastico, e i termini $dN_E \cdot dN_I$ sono $O(dt^2)$ per l'indipendenza), si ottiene:

\begin{equation}
\mathbb{E}[(dV)^2] \approx J_E^2\, \mathbb{E}[dN_E^2] + J_I^2\, \mathbb{E}[dN_I^2]
\end{equation}

Usando il risultato fondamentale $\mathbb{E}[dN_\alpha^2] = K_\alpha \nu_\alpha\, dt + O(dt^2)$:

\begin{equation}
\mathbb{E}[(dV)^2] = J_E^2 K_E \nu_E\, dt + J_I^2 K_I \nu_I\, dt + O(dt^2)
\end{equation}

Dividendo per $dt$ e passando al limite:

\begin{equation}
\sigma^2(t) = J_E^2 K_E \nu_E(t) + J_I^2 K_I \nu_I(t)
\end{equation}

Poiché $J_I^2 = J_I^2 > 0$ (il quadrato è sempre positivo), entrambi i termini contribuiscono positivamente alla varianza. La varianza infinitesimale $\sigma^2$ è \textbf{indipendente da $V$}: la diffusione è \textbf{omogenea} nello spazio del potenziale di membrana. Questo è una conseguenza del fatto che l'ampiezza dei salti sinaptici ($J_E$, $J_I$) non dipende dal valore corrente di $V$.

\subsubsection{Il Significato dei Momenti Centrali di Ordine Superiore}

Il professore mostra anche che i \textbf{momenti centrali di ordine superiore} ($n \geq 3$) sono anch'essi proporzionali a $dt$:

\begin{equation}
m_\alpha^{(n)} \equiv \mathbb{E}[(dN_\alpha - \mathbb{E}[dN_\alpha])^n] = J_\alpha^n K_\alpha \nu_\alpha\, dt + O(dt^2)
\end{equation}

Tuttavia, il punto cruciale è che nel \textbf{limite diffusivo} (che verrà introdotto formalmente nella sezione successiva), $J_\alpha \to 0$ con $K_\alpha \to \infty$. I momenti di ordine $n \geq 3$ scalano come $J_\alpha^n \cdot K_\alpha \propto J_\alpha^{n-2} \cdot J_\alpha^2 K_\alpha \to 0$ per $n \geq 3$ (poiché $J_\alpha \to 0$ più velocemente di quanto $J_\alpha^2 K_\alpha$ cresca). Di conseguenza, nel limite diffusivo, \textbf{solo i primi due momenti infinitesimali sopravvivono}: la media $\mu$ e la varianza $\sigma^2$. Questa è la condizione che caratterizza i \textbf{processi di diffusione gaussiana}.


\section{Il Limite Diffusivo: Condizioni di Scaling}

\subsubsection{La Tensione tra Media e Varianza}

Il professore affronta un problema sottile. Nel limite biologico, $J_\alpha \to 0$ (i singoli salti sinaptici sono piccoli) e $K_\alpha \to \infty$ (il numero di contatti è grande). In questo doppio limite, la varianza infinitesimale scala come:

\begin{equation}
\sigma^2 \propto J_\alpha^2 \cdot K_\alpha \nu_\alpha
\end{equation}

affinché $\sigma^2$ rimanga \textbf{finita}, occorre che:

\begin{equation}
J_\alpha^2 K_\alpha \nu_\alpha = \text{cost} \qquad \Rightarrow \qquad J_\alpha \sim \frac{1}{\sqrt{K_\alpha}}
\end{equation}

Questo è il \textbf{limite diffusivo} (o \textbf{limite idrodinamico}): $J_\alpha$ deve scalare come $1/\sqrt{K_\alpha}$.

\subsubsection{Il Problema della Media nel Limite Diffusivo}

Tuttavia, se si inserisce questo scaling nella media infinitesimale:

\begin{equation}
\mu_0 = J_E K_E \nu_E + J_I K_I \nu_I \propto \frac{1}{\sqrt{K_\alpha}} \cdot K_\alpha \propto \sqrt{K_\alpha} \to \infty
\end{equation}

La media diverge! Questo è un problema: nel limite $K_\alpha \to \infty$, se si mantiene solo lo scaling $J_\alpha \sim 1/\sqrt{K_\alpha}$, la media diventa infinita mentre la varianza rimane finita. Il potenziale di membrana verrebbe immediatamente spinto a valori enormi (o enormemente negativi), e il neurone non potrebbe funzionare in modo biologicamente plausibile.

\subsubsection{Il Bilanciamento E/I come Soluzione}

Il professore mostra che la \textbf{soluzione biologica} a questo problema è il \textbf{bilanciamento eccitatorio/inibitorio} (\textit{E/I balance}). Se il contributo eccitatorio e quello inibitorio alla media quasi si cancellano, si può avere una media finita anche nel limite $K_\alpha \to \infty$.

Formalmente, si richiede che la media infinitesimale rimanga $O(1)$ anche quando $K_\alpha \to \infty$. Questo richiede:

\begin{equation}
J_E K_E \nu_E + J_I K_I \nu_I = O(1)
\end{equation}

con $J_\alpha \sim 1/\sqrt{K_\alpha}$. Questo significa che i due termini, ciascuno $O(\sqrt{K_\alpha})$, devono cancellarsi quasi perfettamente. Il residuo finito è la media effettiva $\mu_0$.

Il professore dimostra che questa condizione di cancellazione è equivalente a richiedere il seguente \textbf{rapporto tra efficacie sinaptiche}:

\begin{equation}
\frac{|J_I|}{J_E} \approx \frac{K_E}{K_I} \approx \frac{80\%}{20\%} = 4
\end{equation}

ovvero, le sinapsi inibitorie devono essere circa \textbf{4 volte più forti} delle eccitatorie per compensare la loro minore numerosità. Questo è esattamente il \textbf{bilanciamento E/I} documentato dalla biologia: i neuroni inibitori, pur essendo il 20\% dei neuroni totali, controllano efficacemente l'eccitazione grazie a sinapsi più potenti.

\section{L'Equazione di Langevin per il Potenziale di Membrana}

\subsubsection{Dalla Dinamica a Salti all'Equazione di Langevin}

Una volta stabilito il limite diffusivo, il professore mostra che la dinamica del potenziale di membrana converge a una \textbf{equazione differenziale stocastica} (SDE) della forma di \textbf{Langevin}. Questa equazione descrive l'evoluzione di $V$ come la somma di un termine deterministico (il drift) e un termine stocastico gaussiano (il rumore):

\begin{equation}
\tau_m \frac{dV}{dt} = -V + \mu_0 + \sigma_V \tau_m\, \xi(t)
\end{equation}

dove:
\begin{itemize}
    \item $\mu_0 = \tau_m(J_E K_E \nu_E + J_I K_I \nu_I)$ è il \textbf{potenziale di deriva effettivo} (il contributo medio dell'input sinaptico, già moltiplicato per $\tau_m$);
    \item $\sigma_V^2 = \tau_m (J_E^2 K_E \nu_E + J_I^2 K_I \nu_I)$ è la \textbf{varianza del potenziale} (varianza del rumore sinaptico, moltiplicata per $\tau_m$);
    \item $\xi(t)$ è un \textbf{rumore bianco gaussiano} con proprietà:
\end{itemize}

\begin{equation}
\langle \xi(t) \rangle = 0, \qquad \langle \xi(t)\, \xi(t') \rangle = \delta(t - t')
\end{equation}

L'equazione di Langevin può anche essere riscritta in forma più compatta come:

\begin{equation}
\tau_m \frac{dV}{dt} = -(V - \mu_0) + \sigma_V \sqrt{2\tau_m}\, \eta(t)
\end{equation}

dove si è ridefinito il rumore per rendere esplicita la convenzione di normalizzazione.

\subsubsection{Origine Fisica del Rumore Bianco}

Il professore spiega perché il rumore è \textbf{bianco} (ovvero, la funzione di correlazione è una delta di Dirac). Questo è la conseguenza diretta del Teorema di Griglioni: la superposizione di $K_\alpha \to \infty$ processi di Poisson sparsi e indipendenti ha correlazioni temporali che vanno a zero nell'intervallo temporale elementare. In altre parole, i singoli incrementi $dN_\alpha(t)$ e $dN_\alpha(t')$ per $t \neq t'$ sono \textbf{stocasticamente indipendenti}, e questa indipendenza si traduce nella struttura $\delta$-correlata del rumore bianco.

\begin{tcolorbox}[colback=gray!5, colframe=gray!40, arc=2mm, boxrule=0.5pt]
\textbf{Nota del docente:} "Il rumore bianco è un'idealizzazione. Nella realtà, se le sinapsi non sono istantanee, il rumore ha una correlazione temporale finita, dell'ordine di $\tau_\alpha$. Ma poiché $\tau_\alpha \ll \tau_m$, questa correlazione è trascurabile sulla scala temporale che ci interessa. L'approssimazione di rumore bianco è dunque ottima per le sinapsi veloci."
\end{tcolorbox}



\subsection{Il Processo di Wiener come Limite del Rumore di Spike}

\subsubsection{Definizione del Processo di Wiener}

Il professore introduce formalmente il \textbf{processo di Wiener} $W(t)$ (o moto browniano standard), che è il processo stocastico di riferimento nell'equazione di Langevin. Esso è definito dalle proprietà:

\begin{enumerate}
    \item $W(0) = 0$ quasi certamente.
    \item $W(t) - W(s) \sim \mathcal{N}(0, t-s)$ per ogni $t > s \geq 0$ (gli incrementi sono gaussiani con varianza $t-s$).
    \item Gli incrementi su intervalli disgiunti sono \textbf{indipendenti}: $W(t_2) - W(t_1)$ e $W(t_4) - W(t_3)$ sono indipendenti se $[t_1, t_2]$ e $[t_3, t_4]$ non si sovrappongono.
    \item Le traiettorie sono continue quasi certamente, ma non differenziabili in nessun punto.
\end{enumerate}

L'equazione di Langevin, scritta in termini del processo di Wiener, diventa:

\begin{equation}
dV = \left(-\frac{V - \mu_0}{\tau_m}\right) dt + \frac{\sigma_V}{\sqrt{\tau_m}}\, dW(t)
\end{equation}

dove $dW(t) = W(t + dt) - W(t) \sim \mathcal{N}(0, dt)$ è l'\textbf{incremento di Wiener}.

\subsubsection{Perché le Traiettorie non sono Differenziabili}

Il professore menziona una proprietà importante: le traiettorie del processo di Wiener sono \textbf{continue ma non differenziabili} in nessun punto. Questo spiega perché il "rumore bianco" $\xi(t) = dW/dt$ non esiste come funzione ordinaria, ma solo come \textbf{distribuzione}. L'equazione di Langevin è dunque un'equazione differenziale stocastica nel senso di Ito o Stratonovich, non un'equazione differenziale ordinaria.


\subsection{Il Processo di Ornstein-Uhlenbeck: LIF Stocastico}

\subsubsection{L'Equazione di Langevin come Processo OU}

L'equazione di Langevin del potenziale di membrana:

\begin{equation}
\tau_m \frac{dV}{dt} = -(V - \mu_0) + \sigma_V \sqrt{2\tau_m}\, \xi(t)
\end{equation}

è l'equazione del celebre \textbf{processo di Ornstein-Uhlenbeck} (OU). Questo processo, introdotto nel 1930 da Ornstein e Uhlenbeck per descrivere la velocità di una particella browniana nel fluido, ha proprietà analitiche molto ricche che lo rendono il modello stocastico di riferimento in molti campi della fisica, dell'ingegneria e delle neuroscienze.

Le proprietà principali del processo OU sono:

\begin{enumerate}
    \item \textbf{Mean-reverting}: il processo tende a tornare verso la media $\mu_0$ con una velocità proporzionale alla distanza $V - \mu_0$. La costante di tempo del ritorno è $\tau_m$.
    \item \textbf{Stazionarietà}: a tempi lunghi ($t \gg \tau_m$), il processo raggiunge una \textbf{distribuzione stazionaria} che è \textbf{gaussiana} con media $\mu_0$ e varianza $\sigma_V^2/2$.
    \item \textbf{Continuità}: le traiettorie sono continue ma non differenziabili.
    \item \textbf{Correlazione temporale}: la funzione di autocorrelazione è:
\end{enumerate}

\begin{equation}
\langle (V(t) - \mu_0)(V(t') - \mu_0) \rangle = \frac{\sigma_V^2}{2}\, e^{-|t-t'|/\tau_m}
\end{equation}

un decadimento esponenziale con costante di tempo $\tau_m$.

\subsubsection{Il LIF Stocastico come Processo OU con Reset}

Il \textbf{neurone LIF stocastico} non è un puro processo OU, perché presenta le condizioni di \textbf{soglia} e \textbf{reset}:
\begin{itemize}
    \item \textbf{Condizione di soglia}: quando $V(t)$ raggiunge $V_{\text{thr}}$, il neurone emette uno spike.
    \item \textbf{Condizione di reset}: immediatamente dopo l'emissione, il potenziale viene riportato al valore di reset $V_r$, ovvero $V(t^+) = V_r$.
\end{itemize}

Il processo OU con soglia e reset è tecnicamente un \textbf{processo OU assorbente-e-rigenerante} (\textit{killing-and-regeneration}): all'assorbimento alla soglia (corrispondente all'emissione dello spike), il processo viene rigenerato nel punto di reset.

\vspace{0.5cm}
\noindent\rule{\textwidth}{0.4pt}
\vspace{0.5cm}

\section{Densità di Probabilità del Potenziale di Membrana: Equazione di Fokker-Planck}

\subsubsection{Dal Singolo Neurone alla Distribuzione di Popolazione}

Fino ad ora si è studiata la dinamica di un \textbf{singolo neurone stocastico}. Il professore introduce ora una prospettiva completamente diversa: la descrizione della \textbf{distribuzione di probabilità} del potenziale di membrana $p(V, t)$. Questa è la probabilità di trovare il potenziale di membrana nel valore $V$ all'istante $t$.

Questa distribuzione ha due interpretazioni equivalenti:
\begin{enumerate}
    \item \textbf{Interpretazione singolo neurone}: $p(V, t)\, dV$ è la probabilità che, al tempo $t$, il potenziale del singolo neurone si trovi nell'intervallo $[V, V + dV]$.
    \item \textbf{Interpretazione di popolazione}: se si considera un insieme di $N \to \infty$ neuroni identici ma soggetti a realizzazioni indipendenti del rumore, $p(V, t)$ è la \textbf{densità di frequenza} del potenziale nell'insieme (ensemble).
\end{enumerate}

\subsubsection{Derivazione dell'Equazione di Fokker-Planck}

L'equazione di evoluzione per $p(V, t)$, associata all'equazione di Langevin, è la \textbf{equazione di Fokker-Planck} (FP). Essa è un'equazione alle derivate parziali di tipo parabolico. Per il processo OU del modello LIF, l'equazione di Fokker-Planck ha la forma:

\begin{equation}
\frac{\partial p(V, t)}{\partial t} = -\frac{\partial}{\partial V}\left[A(V)\, p(V, t)\right] + \frac{1}{2}\frac{\partial^2}{\partial V^2}\left[B(V)\, p(V, t)\right]
\end{equation}

dove:
\begin{itemize}
    \item $A(V) = \mu(V)$ è il \textbf{coefficiente di drift} (momento infinitesimale di primo ordine): $A(V) = (-V + \mu_0)/\tau_m$.
    \item $B(V) = \sigma^2$ è il \textbf{coefficiente di diffusione} (momento infinitesimale di secondo ordine): $B(V) = \sigma_V^2/\tau_m$.
\end{itemize}

Sostituendo esplicitamente:

\begin{equation}
\frac{\partial p}{\partial t} = \frac{\partial}{\partial V}\left[\frac{V - \mu_0}{\tau_m}\, p\right] + \frac{\sigma_V^2}{2\tau_m}\frac{\partial^2 p}{\partial V^2}
\end{equation}

\subsubsection{Forma della Corrente di Probabilità}

È conveniente riscrivere l'equazione di Fokker-Planck in forma di \textbf{equazione di continuità}:

\begin{equation}
\frac{\partial p}{\partial t} = -\frac{\partial J(V, t)}{\partial V}
\end{equation}

dove $J(V, t)$ è la \textbf{corrente di probabilità} (o flusso di probabilità):

\begin{equation}
J(V, t) = -\frac{V - \mu_0}{\tau_m}\, p(V, t) - \frac{\sigma_V^2}{2\tau_m}\frac{\partial p}{\partial V}
\end{equation}

Il primo termine è la \textbf{corrente di drift} (o corrente di avvezione), che trasporta la probabilità nella direzione del campo di forza deterministico. Il secondo termine è la \textbf{corrente di diffusione}, che tende a uniformare la distribuzione di probabilità (dispersione gaussiana). L'equazione di continuità esprime la \textbf{conservazione della probabilità}: la variazione di $p$ in un punto è uguale all'opposto della divergenza del flusso di probabilità in quel punto.



\section{Condizioni al Contorno: Soglia e Reset nel Modello LIF}

\subsubsection{Il Dominio di Definizione della Distribuzione}

Nel modello LIF, il potenziale di membrana $V$ è confinato nell'intervallo:

\begin{equation}
V \in (-\infty, V_{\text{thr}})
\end{equation}

con $V_{\text{thr}}$ la \textbf{soglia di emissione} (tipicamente $V_{\text{thr}} \approx 20\,\text{mV}$ sopra il potenziale di riposo). Al di sotto del potenziale di reset $V_r$, il potenziale può tecnicamente scendere per effetto dell'inibizione, ma in pratica si considera un limite inferiore $V_r \ll V_{\text{thr}}$.

\subsubsection{Condizione alla Soglia: Assorbimento}

Alla soglia $V = V_{\text{thr}}$, il potenziale viene assorbito e lo spike viene emesso. Questo si traduce nella condizione al contorno:

\begin{equation}
p(V_{\text{thr}}, t) = 0
\end{equation}

Questa è una \textbf{condizione di Dirichlet}: la densità di probabilità si annulla alla soglia. Il significato fisico è che ogni volta che una traiettoria tocca $V_{\text{thr}}$, viene immediatamente rimossa dal dominio (emissione dello spike) e riapparirà al reset.

\begin{tcolorbox}[colback=gray!5, colframe=gray!40, arc=2mm, boxrule=0.5pt]
\textbf{Nota del docente:} "Questa è la condizione al contorno \textbf{assorbente} alla soglia. La probabilità di trovare il neurone esattamente in $V_{\text{thr}}$ è zero, perché ogni volta che $V$ raggiunge $V_{\text{thr}}$, il neurone spara e $V$ viene riportato a $V_r$."
\end{tcolorbox}

\subsubsection{Condizione al Reset: Iniezione di Probabilità}

Al reset $V = V_r$, ogni volta che uno spike viene emesso, il potenziale viene riportato a $V_r$. Questo equivale a un'\textbf{iniezione di densità di probabilità} in $V = V_r$ al tasso di firing $\nu_{\text{out}}$. La condizione al contorno al reset si scrive come una \textbf{condizione di flusso}:

\begin{equation}
J(V_r^+, t) - J(V_r^-, t) = \nu_{\text{out}}(t)
\end{equation}

ovvero, il flusso di probabilità ha un salto in $V_r$ pari al tasso di firing: ogni traiettoria che arriva a $V_{\text{thr}}$ (e genera un flusso $J(V_{\text{thr}}, t)$) viene reinserita a $V_r$ con lo stesso flusso.

\subsubsection{Condizione alle Grandi Distanze}

Per $V \to -\infty$, si richiede che la distribuzione vada a zero:

\begin{equation}
p(V, t) \to 0 \quad \text{per } V \to -\infty
\end{equation}

e che il flusso di probabilità si annulli:

\begin{equation}
J(V, t) \to 0 \quad \text{per } V \to -\infty
\end{equation}

Queste condizioni garantiscono la normalizzazione: $\int_{-\infty}^{V_{\text{thr}}} p(V, t)\, dV = 1$.

\vspace{0.5cm}
\noindent\rule{\textwidth}{0.4pt}
\vspace{0.5cm}

\section{Soluzione Stazionaria: Distribuzione di Stato Stazionario}

\subsubsection{Il Regime Stazionario}

Il professore si concentra sulla soluzione in regime stazionario, in cui la distribuzione $p(V, t) = p_0(V)$ non varia nel tempo:

\begin{equation}
\frac{\partial p}{\partial t} = 0 \qquad \Rightarrow \qquad \frac{\partial J}{\partial V} = 0
\end{equation}

L'equazione di continuità in regime stazionario dice dunque che la corrente di probabilità è \textbf{indipendente da $V$} nel bulk (lontano da $V_r$):

\begin{equation}
J(V) = \text{cost} = J_0 \quad \text{per } V \neq V_r
\end{equation}

\subsubsection{La Forma della Corrente Stazionaria}

La corrente stazionaria è:

\begin{equation}
J_0 = -\frac{V - \mu_0}{\tau_m} p_0(V) - \frac{\sigma_V^2}{2\tau_m} \frac{dp_0}{dV}
\end{equation}

Questa è un'equazione differenziale del primo ordine per $p_0(V)$, che può essere riscritta come:

\begin{equation}
\frac{dp_0}{dV} + \frac{2(V - \mu_0)}{\sigma_V^2}\, p_0 = -\frac{2\tau_m J_0}{\sigma_V^2}
\end{equation}

L'equazione è un'equazione lineare del primo ordine non omogenea, con soluzione:

\begin{equation}
p_0(V) = \frac{2\tau_m J_0}{\sigma_V^2}\, e^{-\frac{(V - \mu_0)^2}{\sigma_V^2}} \int_{V}^{V_{\text{thr}}} e^{\frac{(u - \mu_0)^2}{\sigma_V^2}}\, du
\end{equation}

Questa soluzione si ottiene moltiplicando per il fattore integrante $e^{(V-\mu_0)^2/\sigma_V^2}$ e integrando da $V$ a $V_{\text{thr}}$ con la condizione al contorno $p_0(V_{\text{thr}}) = 0$.

\subsubsection{Il Significato della Costante $J_0$}

La costante $J_0$ ha un preciso significato fisico: è il \textbf{flusso di probabilità uniforme} nel bulk, ovvero il tasso con cui le traiettorie "scorrono" da destra a sinistra (o da sinistra a destra) nel dominio. Per le condizioni al contorno del modello LIF, si può mostrare che $J_0 = \nu_{\text{out}}$: il \textbf{tasso di firing in uscita} è esattamente uguale al flusso di probabilità stazionario. La costante $J_0$ viene determinata imponendo la normalizzazione della distribuzione:

\begin{equation}
\int_{-\infty}^{V_{\text{thr}}} p_0(V)\, dV = 1
\end{equation}

\vspace{0.5cm}
\noindent\rule{\textwidth}{0.4pt}
\vspace{0.5cm}

\section{Il Tasso di Firing $\nu_{\text{out}}$ come Flusso di Probabilità alla Soglia}

\subsubsection{Relazione tra Flusso e Tasso di Firing}

Il professore deriva il \textbf{tasso di firing in uscita} $\nu_{\text{out}}$ dalla soluzione dell'equazione di Fokker-Planck. In regime stazionario, il flusso di probabilità è costante nel bulk:

\begin{equation}
J(V) = J_0 = \nu_{\text{out}} \quad \text{per qualsiasi } V \in (V_r, V_{\text{thr}})
\end{equation}

Alla soglia, la condizione $p_0(V_{\text{thr}}) = 0$ e il flusso diretto verso la soglia determinano il tasso di firing:

\begin{equation}
\nu_{\text{out}} = J_0 = \left[\int_{V_r}^{V_{\text{thr}}} \frac{2\tau_m}{\sigma_V^2}\, e^{\frac{(u - \mu_0)^2}{\sigma_V^2}} \int_{u}^{V_{\text{thr}}} e^{-\frac{(w - \mu_0)^2}{\sigma_V^2}}\, dw\, du\right]^{-1}
\end{equation}

Questa formula, per quanto complessa, è \textbf{esplicita}: dati $\mu_0$, $\sigma_V^2$, $\tau_m$, $V_r$, $V_{\text{thr}}$, il tasso di firing in uscita $\nu_{\text{out}}$ è completamente determinato.

\subsubsection{Forma Compatta con la Funzione di Errore}

L'integrale può essere espresso in termini della \textbf{funzione di errore complementare} $\text{erfc}(x) = 1 - \text{erf}(x)$. Definendo le variabili adimensionali:

\begin{equation}
u_{\text{thr}} = \frac{V_{\text{thr}} - \mu_0}{\sigma_V / \sqrt{2}}, \qquad u_r = \frac{V_r - \mu_0}{\sigma_V / \sqrt{2}}
\end{equation}

il tasso di firing si scrive nella forma nota:

\begin{equation}
\nu_{\text{out}}^{-1} = \tau_{\text{ref}} + \tau_m \sqrt{\pi} \int_{u_r}^{u_{\text{thr}}} e^{x^2}\left(1 + \text{erf}(x)\right)\, dx
\end{equation}

dove si è incluso anche il \textbf{periodo refrattario assoluto} $\tau_{\text{ref}}$ (la durata minima tra due spike consecutivi), che può essere aggiunto come una costante additiva al tempo medio di primo passaggio.

Il professore sottolinea che questa formula è \textbf{esatta} nel regime diffusivo (LIF stocastico con rumore bianco gaussiano). Non è un'approssimazione: è la soluzione analitica completa del problema del primo passaggio per il processo OU con soglia e reset.

\begin{tcolorbox}[colback=gray!5, colframe=gray!40, arc=2mm, boxrule=0.5pt]
\textbf{Nota del docente:} "Questa formula è uno dei risultati più belli della teoria dei neuroni stocastici. Vi dà il tasso di firing in funzione di soli due numeri — la media $\mu_0$ e la varianza $\sigma^2$ dell'input — più le costanti biologiche $\tau_m$, $V_{\text{thr}}$, $V_r$. È una formula chiusa, analitica. Il neurone è stato risolto."
\end{tcolorbox}

\vspace{0.5cm}
\noindent\rule{\textwidth}{0.4pt}
\vspace{0.5cm}

\section{La Funzione di Trasferimento: $\nu_{\text{out}} = \Phi(\mu_0, \sigma_V^2)$}

\subsubsection{Definizione della Funzione di Trasferimento}

Il risultato fondamentale dell'equazione di Fokker-Planck è che il \textbf{tasso di firing in uscita} $\nu_{\text{out}}$ è una funzione deterministica dei parametri statistici dell'input sinaptico:

\begin{equation}
\nu_{\text{out}} = \Phi(\mu_0,\, \sigma_V^2)
\end{equation}

Questa funzione $\Phi$ è la \textbf{funzione di trasferimento} del neurone LIF stocastico. Il professore analizza le proprietà di $\Phi$:

\subsubsection{Regime Sub-Soglia: $\mu_0 < V_{\text{thr}}$}

Quando l'input medio è sotto soglia ($\mu_0 < V_{\text{thr}}$), il neurone può comunque sparare grazie alle \textbf{fluttuazioni stocastiche}. In questo regime, il tasso di firing cresce con la varianza $\sigma_V^2$: un rumore più grande spinge il potenziale oltre la soglia più frequentemente. La funzione $\Phi$ è una funzione \textbf{crescente} di $\sigma_V^2$ in questo regime.

Intuitivamente, le fluttuazioni sinaptiche svolgono qui un ruolo paradossale: anche se in media il potenziale non raggiunge la soglia, le code della distribuzione gaussiana la raggiungono con probabilità finita. Maggiore è il "rumore", maggiore è la probabilità di emissione. Questo è un esempio di \textbf{stochastic resonance} o, più precisamente, di \textbf{noise-driven firing}.

\subsubsection{Regime Sopra-Soglia: $\mu_0 > V_{\text{thr}}$}

Quando l'input medio è sopra soglia ($\mu_0 > V_{\text{thr}}$), il neurone spara anche in assenza di rumore. In questo regime, le fluttuazioni hanno un effetto ambivalente:
\begin{itemize}
    \item Se il rumore è grande, le traiettorie raggiungono la soglia prima che il drift deterministico le porti lì, \textbf{aumentando} il tasso di firing.
    \item Ma le fluttuazioni possono anche portare il potenziale verso valori bassi, \textbf{ritardando} l'emissione.
\end{itemize}

Il risultato netto è che in questo regime $\Phi$ è una funzione \textbf{decrescente} di $\sigma_V^2$: il rumore riduce il tasso di firing rispetto al caso deterministico.

\subsubsection{Andamento Sigmoide della Funzione di Trasferimento}

La funzione $\Phi(\mu_0, \sigma_V^2 = \text{cost})$ come funzione di $\mu_0$ ha un andamento \textbf{sigmoidale} (a "S"): cresce monotonicamente da 0 a $1/\tau_{\text{ref}}$ (il tasso di firing massimo determinato dal periodo refrattario), con un punto di inflessione vicino a $\mu_0 \approx V_{\text{thr}}$.

Questo andamento sigmoidale è notevolmente simile alla \textbf{funzione di attivazione sigmoide} usata nelle reti neurali artificiali, suggerendo che il comportamento collettivo di un neurone biologico stocastico è ben catturato da questa funzione matematica. Il professore sottolinea che, al contrario delle reti artificiali, qui la sigmoide non è assunta \textit{a priori} ma è \textbf{derivata rigorosamente} dai principi biofisici.

\begin{tcolorbox}[colback=gray!5, colframe=gray!40, arc=2mm, boxrule=0.5pt]
\textbf{Nota del docente:} "Guardate cosa abbiamo ottenuto: partendo dall'equazione microscopica di un neurone biologico con rumore stocastico, abbiamo derivato una funzione di trasferimento sigmoidale. Questa è la giustificazione biofisica delle funzioni di attivazione che si usano nelle reti neurali artificiali. Non è un'assunzione: è un risultato."
\end{tcolorbox}

\vspace{0.5cm}
\noindent\rule{\textwidth}{0.4pt}
\vspace{0.5cm}

\section{Teoria di Campo Medio: Consistenza Self-Consistent}

\subsubsection{Il Problema della Consistenza}

Fino a ora, si è trattato un singolo neurone immerso in un campo sinaptico esterno, con parametri $\mu_0$ e $\sigma_V^2$ assegnati come input. Nella realtà, però, il neurone fa parte di una rete: la sua attività contribuisce al campo sinaptico degli altri neuroni, che a loro volta influenzano il neurone in esame. Questa \textbf{retroazione} (feedback) deve essere trattata in modo consistente.

La \textbf{teoria di campo medio} (mean field theory) fornisce il quadro formale per farlo. L'idea fondamentale è che, in una rete sufficientemente grande e sufficientemente connessa (il limite termodinamico della connettività), l'input ricevuto da ciascun neurone è determinato dall'\textbf{attività media} dell'intera rete, con piccole fluttuazioni trascurabili. Ogni neurone è dunque equivalente, e si può risolvere un \textbf{problema self-consistent} per l'unico neurone rappresentativo.

\subsubsection{La Condizione di Self-Consistenza}

Il professore formalizza la condizione di campo medio. Si considerino due popolazioni di neuroni: una eccitatoria (E) con tasso di firing $\nu_E$, e una inibitoria (I) con tasso di firing $\nu_I$. I parametri dell'input sinaptico ricevuto dal neurone rappresentativo (in entrambe le popolazioni) sono:

\begin{equation}
\mu_0 = J_E K_E \nu_E + J_I K_I \nu_I
\end{equation}

\begin{equation}
\sigma_V^2 = J_E^2 K_E \nu_E + J_I^2 K_I \nu_I
\end{equation}

La funzione di trasferimento determina il tasso di firing di ciascuna popolazione:

\begin{equation}
\nu_E = \Phi(\mu_0^E,\, \sigma_{V,E}^2)
\end{equation}

\begin{equation}
\nu_I = \Phi(\mu_0^I,\, \sigma_{V,I}^2)
\end{equation}

dove $\mu_0^{E,I}$ e $\sigma_{V,E,I}^2$ sono i parametri dell'input calcolati per un neurone della popolazione $E$ o $I$ rispettivamente.

La \textbf{condizione di self-consistenza} è che i tassi di firing scelti per calcolare $\mu_0$ e $\sigma_V^2$ siano gli stessi che risultano dalla funzione di trasferimento. Si tratta dunque di trovare il \textbf{punto fisso} del sistema:

\begin{equation}
\nu_E^* = \Phi\!\left(J_E K_{EE} \nu_E^* + J_I K_{EI} \nu_I^*,\; J_E^2 K_{EE} \nu_E^* + J_I^2 K_{EI} \nu_I^*\right)
\end{equation}

\begin{equation}
\nu_I^* = \Phi\!\left(J_E K_{IE} \nu_E^* + J_I K_{II} \nu_I^*,\; J_E^2 K_{IE} \nu_E^* + J_I^2 K_{II} \nu_I^*\right)
\end{equation}

dove $K_{\beta\alpha}$ è il numero di contatti sinaptici che un neurone della popolazione $\beta$ riceve dalla popolazione $\alpha$.

\subsubsection{Soluzione Grafica del Punto Fisso}

Il professore illustra la soluzione della condizione di self-consistenza in modo grafico, nel caso semplificato di una sola popolazione (rete eccitatoria omogenea). In questo caso:

\begin{equation}
\nu^* = \Phi\!\left(\mu_0(\nu^*),\, \sigma^2(\nu^*)\right)
\end{equation}

Si tratta di trovare l'\textbf{intersezione} tra la curva $\nu_{\text{out}} = \Phi(\mu_0(\nu), \sigma^2(\nu))$ (funzione di trasferimento, curva sigmoidale) e la retta $\nu_{\text{out}} = \nu$ (condizione di consistenza, retta bisettrice del quadrante I).

Le possibili soluzioni dipendono dai parametri della rete:
\begin{enumerate}
    \item \textbf{Un solo punto fisso} (soluzione unica): la rete ha un unico stato di attività stabile.
    \item \textbf{Tre punti fissi} (bistabilità): coesistono uno stato di bassa attività, uno stato di alta attività, e un terzo punto fisso instabile. Questo è il substrato di neurobiologico per la \textbf{working memory} e altri fenomeni di attività persistente.
    \item \textbf{Nessuna soluzione} (impossibile per una $\Phi$ normalizzata): non fisicamente rilevante.
\end{enumerate}

\begin{tcolorbox}[colback=gray!5, colframe=gray!40, arc=2mm, boxrule=0.5pt]
\textbf{Nota del docente:} "La bistabilità è di enorme importanza per le neuroscienze cognitive. Se la rete ha due stati stabili, può mantenere una traccia di memoria anche in assenza dello stimolo che l'ha generata. Questo è il meccanismo proposto per la memoria di lavoro. E tutto ciò emerge naturalmente dalla condizione self-consistent, senza alcuna assunzione aggiuntiva."
\end{tcolorbox}

\subsubsection{Proprietà del Punto Fisso: Stabilità}

Un punto fisso $\nu^*$ è \textbf{stabile} se piccole perturbazioni $\delta\nu$ intorno ad esso decadono nel tempo. La condizione di stabilità è:

\begin{equation}
\frac{d\Phi}{d\nu}\bigg|_{\nu = \nu^*} < 1
\end{equation}

ovvero, il guadagno effettivo della rete in quel punto deve essere inferiore a 1. Un guadagno uguale a 1 corrisponde a una \textbf{biforcazione} (il punto fisso perde la stabilità), mentre un guadagno maggiore di 1 indica instabilità e transizione verso un altro stato.

\subsubsection{La Mappa dei Regimi di Rete (Diagramma di Brunel)}

La teoria di campo medio permette di costruire una \textbf{mappa dei regimi di attività} della rete nel piano dei parametri. Nel lavoro seminale di Brunel (2000), questa mappa è costruita nel piano ($g$, $\nu_{\text{ext}}/\nu_{\text{thr}}$), dove:

\begin{itemize}
    \item $g = |J_I|/J_E$ è il \textbf{rapporto di efficacia inibitoria/eccitatoria} (misura della forza relativa dell'inibizione).
    \item $\nu_{\text{ext}}/\nu_{\text{thr}}$ è il tasso di input esterno normalizzato alla soglia.
\end{itemize}

Questa mappa rivela quattro regimi fondamentali della rete, ciascuno con caratteristiche distintive:

\begin{enumerate}
    \item \textbf{Regime Sincrono e Regolare (SR — Synchronous Regular)}: la rete oscilla in modo sincrono e regolare. Tutti i neuroni sparano allo stesso ritmo. Questo regime corrisponde a situazioni patologiche (epilessia) o a certi ritmi oscillatori cerebrali.
    \item \textbf{Regime Sincrono e Irregolare (SI — Synchronous Irregular)}: la rete oscilla collettivamente (il firing è sincrono a livello di popolazione) ma i singoli neuroni sparano in modo irregolare (ISI variabili). Questo regime è osservato in alcune condizioni fisiologiche, come i ritmi gamma nella corteccia.
    \item \textbf{Regime Asincrono e Irregolare (AI — Asynchronous Irregular)}: il regime fisiologico normale della corteccia in veglia. I neuroni sparano in modo indipendente (asincrono) e irregolare. Il firing rate è basso ($\approx 1$–$5\,\text{Hz}$), il coefficiente di variazione degli ISI è vicino a 1 (come in un processo di Poisson). Questo regime è stabilizzato dall'inibizione: richiede $g > 4$ (inibizione dominante).
    \item \textbf{Regime Fast-Oscillating (FO — Fast Oscillatory)}: oscillazioni rapide della rete, tipicamente nella banda gamma ($40$–$80\,\text{Hz}$). Si verifica al confine tra il regime AI e il regime SI.
\end{enumerate}

Il regime \textbf{AI} è quello biologicamente più rilevante per la corteccia in condizioni normali. La teoria di campo medio mostra che esso è stabilizzato da una \textbf{forte inibizione} ($g \approx 4$–$6$) in combinazione con un \textbf{input esterno sufficientemente grande} ($\nu_{\text{ext}} \sim \nu_{\text{thr}}$). Questo è consistente con la condizione di bilanciamento E/I derivata nel limite diffusivo.

\begin{tcolorbox}[colback=gray!5, colframe=gray!40, arc=2mm, boxrule=0.5pt]
\textbf{Nota del docente:} "Il diagramma di Brunel è uno dei risultati più belli della neuroscienze computazionali: con una manciata di parametri, cattura l'intera fenomenologia dell'attività corticale spontanea. Ai regimi corrispondono stati cerebrali diversi: il regime AI è la veglia normale, il regime SR è simile all'epilessia, il regime oscillatorio è associato a certi stati del sonno o dell'attenzione."
\end{tcolorbox}

\vspace{0.5cm}
\noindent\rule{\textwidth}{0.4pt}
\vspace{0.5cm}

\section{Esempi Numerici e Stima dei Parametri}

\subsubsection{Parametri Tipici del Modello}

Il professore fornisce i valori tipici dei parametri del modello LIF diffusivo per una rete corticale:

\begin{table}[htbp]
    \centering
    \renewcommand{\arraystretch}{1.3}
    \begin{tabularx}{\textwidth}{|X|c|c|}
    \hline
    \textbf{Parametro} & \textbf{Simbolo} & \textbf{Valore tipico} \\
    \hline
    Costante di tempo di membrana & $\tau_m$ & $20\,\text{ms}$ \\
    \hline
    Soglia di emissione & $V_{\text{thr}}$ & $20\,\text{mV}$ \\
    \hline
    Potenziale di reset & $V_r$ & $0\,\text{mV}$ \\
    \hline
    Periodo refrattario assoluto & $\tau_{\text{ref}}$ & $2\,\text{ms}$ \\
    \hline
    Efficacia eccitatoria & $J_E$ & $0.1$–$0.2\,\text{mV}$ \\
    \hline
    Efficacia inibitoria & $|J_I|$ & $0.4$–$0.8\,\text{mV}$ \\
    \hline
    Numero di contatti eccitatori & $K_E$ & $\approx 16\,000$ \\
    \hline
    Numero di contatti inibitori & $K_I$ & $\approx 4\,000$ \\
    \hline
    Tasso di firing di background & $\bar{\nu}$ & $1$–$5\,\text{Hz}$ \\
    \hline
    \end{tabularx}
\end{table}

\subsubsection{Stima di $\mu_0$ e $\sigma^2$}

Con questi parametri tipici, il professore calcola l'ordine di grandezza di $\mu_0$ e $\sigma_V^2$:

\begin{equation}
\mu_0 = J_E K_E \nu_E + J_I K_I \nu_I
\end{equation}
\begin{equation}
\approx 0.1\,\text{mV} \times 16\,000 \times 5\,\text{Hz} + (-0.4\,\text{mV}) \times 4\,000 \times 5\,\text{Hz}
\end{equation}
\begin{equation}
= 8\,\text{mV/s} \times 20\,\text{ms} - 8\,\text{mV/s} \times 20\,\text{ms} = \text{(quasi cancellazione)}
\end{equation}

Il valore risultante di $\mu_0$ è piccolo (dell'ordine di $1$–$5\,\text{mV}$, sotto soglia), evidenziando che la rete opera nel regime di \textbf{bilanciamento eccitatorio/inibitorio}.

Per la varianza:

\begin{equation}
\sigma_V^2 = J_E^2 K_E \nu_E + J_I^2 K_I \nu_I
\end{equation}
\begin{equation}
\approx (0.1)^2 \times 16\,000 \times 5 + (0.4)^2 \times 4\,000 \times 5
\end{equation}
\begin{equation}
= 0.01 \times 80\,000 + 0.16 \times 20\,000 = 800 + 3200 = 4000\,\text{mV}^2/\text{s}
\end{equation}
\begin{equation}
\Rightarrow \sigma_V = \sqrt{4000 \times 0.02\,\text{s}} \approx 9\,\text{mV}
\end{equation}

Il valore $\sigma_V \approx 9\,\text{mV}$ è dello stesso ordine di grandezza del "gap" tra la media $\mu_0$ e la soglia $V_{\text{thr}} \approx 20\,\text{mV}$, il che indica che le fluttuazioni svolgono un ruolo fondamentale nel portare il potenziale alla soglia. Il regime è genuinamente diffusivo.

\vspace{0.5cm}
\noindent\rule{\textwidth}{0.4pt}
\vspace{0.5cm}

\section{Domande degli Studenti e Approfondimenti}

\subsubsection{La Questione della Dipendenza Temporale del Tasso}

Durante la lezione uno studente pone la seguente domanda: «\textit{Il modello di Stein che hai descritto usa un processo di Poisson omogeneo in tempo per i treni di spike presinaptici. Ma noi stessi abbiamo detto che la densità di spike sperimentale è non stazionaria. Come si giustifica l'uso di un tasso costante $\nu_\alpha$?}»

Il professore risponde che la questione ha due livelli:

\begin{enumerate}
    \item \textbf{Livello del singolo neurone (approssimazione Poissoniana)}: l'approssimazione che i singoli neuroni presinaptici siano processi di Poisson è una semplificazione. I neuroni reali hanno distribuzioni di ISI non-esponenziali (a causa del periodo refrattario, dell'adattamento, etc.). Tuttavia, il Teorema di Griglioni mostra che questa assunzione non è necessaria: anche se i singoli neuroni hanno ISI non-esponenziali, la superposizione converge comunque a un processo di Poisson (non omogeneo in tempo).
    \item \textbf{Livello della rete (non-stazionarietà)}: la non-stazionarietà sperimentale della densità di spike può essere incorporata usando un processo di Poisson non omogeneo, cioè con intensità $\lambda_\alpha(t) = K_\alpha \nu_\alpha(t)$ che varia nel tempo. Il modello di campo medio stazionario che deriviamo è valido nell'approssimazione di \textbf{lentezza}: si assume che $\nu_\alpha(t)$ vari su scale temporali molto più lente di $\tau_m$, così che il sistema raggiunga quasi-stazionariamente lo stato di campo medio corrispondente al valore istantaneo di $\nu_\alpha(t)$.
\end{enumerate}

\begin{tcolorbox}[colback=gray!5, colframe=gray!40, arc=2mm, boxrule=0.5pt]
\textbf{Nota del docente:} "In pratica, consideriamo sempre il caso stazionario qui. Quando si vuole studiare la risposta dinamica della rete a stimoli temporali, si fa un'analisi di risposta lineare o non lineare attorno al punto fisso stazionario. Questo è esattamente ciò che fa la teoria della risposta lineare che vedrete nelle prossime lezioni."
\end{tcolorbox}

\subsubsection{La Questione dell'Indipendenza dei Neuroni Presinaptici}

Un'altra domanda riguarda l'assunzione di \textbf{indipendenza} dei neuroni presinaptici richiesta dal Teorema di Griglioni:

\textit{«Ma i neuroni presinaptici non sono indipendenti: appartengono alla stessa rete e sono quindi correlati tra loro. Come si giustifica l'assunzione di indipendenza?»}

Il professore risponde in modo dettagliato:

\begin{itemize}
    \item In una rete di $N$ neuroni con connettività sparsa (ogni neurone riceve da una frazione $\sim K/N$ dei neuroni), nel limite $N \to \infty$ con $K$ fissato o crescente più lentamente di $N$, le correlazioni tra neuroni diversi tendono a zero come $\sim 1/N$. Questa è la giustificazione del campo medio: in reti grandi, ogni neurone vede i suoi presinaptici come essenzialmente indipendenti.
    \item La situazione è più complessa nelle \textbf{reti piccole} o nelle \textbf{reti con correlazioni forti} (ad esempio, reti con plasticità sinaptica specifica). In questi casi l'approccio di campo medio deve essere corretto includendo termini di correlazione.
\end{itemize}

\begin{tcolorbox}[colback=gray!5, colframe=gray!40, arc=2mm, boxrule=0.5pt]
\textbf{Nota del docente:} "La giustificazione rigorosa dell'indipendenza viene dalla teoria di campo medio stessa: mostreremo che la soluzione di campo medio è autoconsistente, ovvero le correlazioni tra neuroni che derivano dalla soluzione di campo medio sono effettivamente trascurabili per reti grandi. È una verifica a posteriori della validità dell'approssimazione."
\end{tcolorbox}

\vspace{0.5cm}
\noindent\rule{\textwidth}{0.4pt}
\vspace{0.5cm}

\section{Riepilogo: Dal Microscopico al Macroscopico}

\subsubsection{La Catena Logica della Derivazione}

Il professore chiude la lezione tracciando la catena logica completa che ha portato dal modello microscopico alla teoria di campo medio:

\textbf{Passo 1} — Modello LIF a sinapsi veloci ($\delta$-approssimazione, $\tau_\alpha \to 0$):
\begin{equation}
\tau_m \frac{dV}{dt} = -V + J_E \sum_j S_{Ej}(t) + J_I \sum_j S_{Ij}(t)
\end{equation}

\textbf{Passo 2} — I treni di spike presinaptici, somma di $K_\alpha \gg 1$ processi sparsi indipendenti, convergono al processo di Poisson per il Teorema di Griglioni.

\textbf{Passo 3} — Nel limite diffusivo ($J_\alpha \to 0$, $K_\alpha \to \infty$, $J_\alpha^2 K_\alpha = \text{cost}$) con bilanciamento E/I, i momenti infinitesimali del processo di conteggio danno:
\begin{equation}
\mu = -\frac{V}{\tau_m} + \frac{\mu_0}{\tau_m}, \qquad \sigma^2 = \frac{\sigma_V^2}{\tau_m}
\end{equation}

\textbf{Passo 4} — L'equazione di Langevin per il singolo neurone:
\begin{equation}
\tau_m \frac{dV}{dt} = -(V - \mu_0) + \sigma_V \sqrt{2\tau_m}\, \xi(t)
\end{equation}
(processo di Ornstein-Uhlenbeck)

\textbf{Passo 5} — L'equazione di Fokker-Planck per la densità di probabilità:
\begin{equation}
\frac{\partial p}{\partial t} = \frac{\partial}{\partial V}\left[\frac{V - \mu_0}{\tau_m}\, p\right] + \frac{\sigma_V^2}{2\tau_m}\frac{\partial^2 p}{\partial V^2}
\end{equation}

\textbf{Passo 6} — Soluzione stazionaria con condizioni al contorno (assorbente a $V_{\text{thr}}$, reset a $V_r$) $\rightarrow$ \textbf{funzione di trasferimento} $\nu_{\text{out}} = \Phi(\mu_0, \sigma_V^2)$.

\textbf{Passo 7} — \textbf{Condizione self-consistent} di campo medio:
\begin{equation}
\nu_\alpha^* = \Phi\!\left(\mu_0(\nu_E^*, \nu_I^*),\, \sigma_V^2(\nu_E^*, \nu_I^*)\right)
\end{equation}

\subsubsection{Connessione con le Reti Neurali Artificiali}

Il professore apre una breve digressione sul collegamento tra i risultati derivati e le reti neurali artificiali, evidenziando come i concetti biofisici rigorosi derivati giustifichino le scelte euristiche delle reti artificiali.

Nella rete artificiale più semplice (percettrone, rete feed-forward), ogni unità computazionale esegue la trasformazione:

\begin{equation}
a_i = f\!\left(\sum_j w_{ij}\, x_j - \theta_i\right)
\end{equation}

dove $f$ è una funzione di attivazione (tipicamente sigmoide o ReLU), $w_{ij}$ sono i pesi sinaptici e $\theta_i$ è la soglia. La scelta della funzione sigmoide come funzione di attivazione è solitamente \textbf{motivata euristicamente} (differenziabilità, range $[0,1]$, saturazione).

La teoria sviluppata in questa lezione fornisce una \textbf{giustificazione biofisica rigorosa}: la funzione di trasferimento $\Phi(\mu_0, \sigma^2)$ del LIF stocastico è effettivamente una funzione sigmoidale di $\mu_0$ (per $\sigma^2$ fissato). I pesi $w_{ij}$ corrispondono alle efficacie sinaptiche $J_{\alpha j}$, la soglia $\theta_i$ corrisponde a $V_{\text{thr}}$. Il rumore sinaptico $\sigma^2$, che nelle reti biologiche è una proprietà emergente del bombardamento sinaptico, svolge il ruolo di un \textbf{parametro di temperatura} che controlla la "morbidezza" della non-linearità: per $\sigma^2 \to 0$ la funzione di trasferimento converge alla funzione di Heaviside (comportamento deterministico), mentre per $\sigma^2 \to \infty$ la sigmoide si appiattisce.

\begin{tcolorbox}[colback=gray!5, colframe=gray!40, arc=2mm, boxrule=0.5pt]
\textbf{Nota del docente:} "Questo è uno dei punti più belli del corso: la biofisica e le reti artificiali si incontrano qui. Non perché qualcuno abbia copiato la biologia nel disegnare le reti artificiali — in realtà la maggior parte delle scelte di design delle reti artificiali è stata fatta senza questa giustificazione. Ma il fatto che ci sia una convergenza è indicativo di qualcosa di più profondo: il bilanciamento tra eccitazione, inibizione e rumore impone una struttura sulla funzione di trasferimento che ha la forma sigmoidale."
\end{tcolorbox}

\subsubsection{Il Significato Fisico dell'Intera Costruzione}

Il professore sottolinea il salto di scala concettuale compiuto:

\begin{itemize}
    \item Si è partiti dalla descrizione del \textbf{singolo neurone} con input sinaptico a salti discreti.
    \item Si è costruita la \textbf{statistica dell'input} per mezzo del Teorema di Griglioni.
    \item Si è derivata l'\textbf{equazione di Langevin} per la dinamica stocastica del potenziale.
    \item Si è ottenuta l'\textbf{equazione di Fokker-Planck} per la distribuzione della popolazione.
    \item Si è ricavata la \textbf{funzione di trasferimento} che connette input e output del singolo neurone.
    \item Infine, si è costruita la \textbf{teoria di campo medio} che descrive l'attività collettiva della rete come soluzione self-consistent.
\end{itemize}

Questo percorso è l'equivalente neuronale di quanto in fisica statistica si fa passando dalla meccanica molecolare alla termodinamica: la complessità microscopica (miliardi di spike) si condensa in pochi parametri macroscopici ($\mu_0$, $\sigma_V^2$, $\nu_{\text{out}}$).

\chapter{Lezione 17}


\section{Organizzazione Strutturale e Funzionale delle Reti Cerebrali}

\subsubsection{Obiettivi della Lezione}

Uno degli obiettivi fondamentali delle neuroscienze computazionali è definire in modo rigoroso che cosa si intende per \textbf{rete neurale}, partendo da un punto di vista biologico e identificando le approssimazioni che permettono di costruire modelli matematici trattabili pur rimanendo fedeli alla realtà sperimentale. A tale scopo è necessario prima comprendere come le reti neurali siano organizzate nella corteccia cerebrale.

Come introdotto nelle prime lezioni del corso, esistono diversi approcci sperimentali per investigare l'organizzazione di una rete cerebrale. Questi approcci si suddividono tradizionalmente in due grandi categorie:

\begin{itemize}
    \item \textbf{Connettività strutturale}: descrive l'esistenza di connessioni fisiche (fibre assonali) tra aree cerebrali.
    \item \textbf{Connettività funzionale}: descrive le correlazioni statistiche nell'attività di regioni diverse del cervello.
\end{itemize}

È importante tenere presente che la presenza di una connessione strutturale non implica necessariamente un'interazione funzionale rilevante in ogni istante. Viceversa, due aree possono mostrare attività correlata anche in assenza di una connessione diretta, grazie a percorsi polisinaptici che connettono indirettamente ogni area corticale a ogni altra.

\subsubsection{La Corteccia come Oggetto di Studio}

La trattazione si concentra sulle \textbf{reti corticali}, ovvero le reti formate dai neuroni della corteccia cerebrale. La corteccia è la struttura laminare che riveste la superficie dei due emisferi cerebrali ed è responsabile delle funzioni cognitive superiori: percezione, movimento volontario, linguaggio, memoria, pianificazione.


\subsection{Diffusion Tensor Imaging (DTI)}

\subsubsection{Il Problema della Connettività a Lunga Distanza}

Per investigare l'esistenza di connessioni a lunga distanza tra aree corticali distinte si ricorre a tecniche di risonanza magnetica. In particolare, il \textbf{Diffusion Tensor Imaging (DTI)} sfrutta le proprietà di diffusione delle molecole d'acqua nei tessuti cerebrali per rivelare la presenza e l'orientamento dei fasci di fibre assonali.

Il \textbf{globo bianco} (white matter) è la porzione più interna del cervello, dove sono concentrate le fibre mielinizzate. La \textbf{mielina} — prodotta dalle cellule di Schwann nel sistema nervoso periferico e dagli oligodendrociti nel sistema nervoso centrale — funge da guaina isolante attorno agli assoni, permettendo la propagazione rapida dei potenziali d'azione su lunghe distanze. Le fibre assonali all'interno del white matter sono organizzate in \textbf{fasci} (bundles) orientati coerentemente.

\subsubsection{Principio Fisico del DTI}

Il DTI misura l'\textbf{anisotropia della diffusione} dei protoni dell'acqua. In un'area priva di strutture orientate, l'acqua diffonde isotropicamente in tutte le direzioni. All'interno di un assone mielinizzato, invece, l'acqua è confinata a diffondere lungo la direzione del fascio (il termine "pipe" è usato dal docente come analogia): si ottiene quindi un \textbf{tensore di diffusione} $\sigma$ con un'eigenvalue dominante corrispondente alla direzione del fascio.

Il volume cerebrale viene suddiviso in \textbf{voxel} tridimensionali. Per ciascun voxel attraversato da una fibra mielinizzata, il tensore di diffusione risulta anisotropico. Procedendo da un voxel sorgente e seguendo la direzione preferenziale del tensore voxel per voxel, si ricostruisce il percorso del fascio: questa procedura prende il nome di \textbf{trattografia} (tractography).

\begin{tcolorbox}[colback=gray!5, colframe=gray!40, arc=2mm, boxrule=0.5pt]
\textbf{Nota del docente:} L'analogia corretta è quella di una serie di tubi contenenti acqua: i protoni dell'acqua non possono muoversi liberamente in ogni direzione, ma sono vincolati all'interno del "tubo" assonale, quindi il loro moto preferenziale è parallelo all'asse del fascio.
\end{tcolorbox}

\subsubsection{Limitazioni della Trattografia DTI}

La trattografia presenta alcune limitazioni intrinseche:

\begin{enumerate}
    \item \textbf{Risoluzione spaziale}: la dimensione del voxel limita la capacità di risolvere fasci sottili.
    \item \textbf{Fibre incrociate}: quando due fasci di fibre si incrociano nello stesso voxel, il tensore di diffusione risultante è la sovrapposizione di due contributi, rendendo difficile disambiguare le due direzioni.
    \item \textbf{Directionality}: il DTI non fornisce informazioni sulla direzione di propagazione del segnale (da pre- a post-sinaptico), né sull'intensità funzionale della connessione.
\end{enumerate}


\subsubsection{Dal Voxel al Grafo di Connettività}

Partendo da un voxel sorgente posizionato sulla superficie corticale (o prossimo ad essa), la trattografia permette di identificare tutti i voxel destinazione raggiunti dalle fibre originate da quella sorgente. Campionando casualmente molti voxel sorgente sulla superficie corticale, si costruisce una \textbf{matrice di connettività strutturale} (o \textbf{connettoma strutturale}). Questa matrice può essere rappresentata come un \textbf{grafo} orientato i cui nodi sono le aree corticali macroscopiche e i cui archi sono i fasci di fibre.

\subsubsection{Limitazioni del Connettoma Strutturale}

Il connettoma strutturale ottenuto tramite DTI fornisce informazioni qualitative: è possibile distinguere la presenza di un fascio numeroso (molte fibre) da uno esiguo (poche fibre), ma con risoluzione limitata. Soprattutto, non è possibile ricavare:

\begin{itemize}
    \item Il \textbf{numero di sinapsi} stabilite tra le fibre della connessione
    \item La \textbf{direzione funzionale} della comunicazione (il DTI ricostruisce traiettorie geometriche, non circuiti orientati)
    \item La \textbf{forza sinaptica} della connessione
\end{itemize}

Queste limitazioni rendono necessario integrare il connettoma strutturale con misure funzionali.

\subsubsection{La Matrice di Adiacenza del Connettoma}

La rappresentazione matematica del connettoma strutturale è una \textbf{matrice di adiacenza binaria} (o pesata) $A_{ij}$ di dimensione $M \times M$, dove $M$ è il numero di aree corticali macroscopiche considerate:

\begin{equation}
A_{ij} = \begin{cases} 1 & \text{se esiste almeno un fascio di fibre da area } j \text{ ad area } i \\ 0 & \text{altrimenti} \end{cases}
\end{equation}

Nel caso pesato, $A_{ij}$ rappresenta il numero di fibre o l'indice di anisotropia frazionale (FA) medio lungo il tratto. La matrice $A$ è tipicamente:

\begin{itemize}
    \item \textbf{Sparsa}: la maggior parte delle coppie di aree non è direttamente connessa
    \item \textbf{Non simmetrica in senso stretto}: la direzione nel DTI non è determinabile, ma si può fare un'assunzione di simmetria come prima approssimazione
    \item \textbf{Gerarchica}: esistono hubs corticali (aree molto connesse) e nodi periferici poco connessi, in accordo con una struttura "small-world"
\end{itemize}



\subsection{Connettività Funzionale}

\subsubsection{Il Segnale BOLD}

La \textbf{risonanza magnetica funzionale} (fMRI) misura l'attività neurale indirettamente attraverso il \textbf{segnale BOLD} (Blood-Oxygen-Level-Dependent). Il segnale BOLD riflette il consumo metabolico di un volume corticale: in risposta a un aumento dell'attività neurale, il flusso sanguigno locale aumenta (risposta emodinamica), e questo modifica il rapporto tra emoglobina ossigenata e deossigenata, che hanno proprietà magnetiche differenti.



\subsubsection{Misura della Connettività Funzionale}

Dividendo la corteccia cerebrale in voxel, si registrano le serie temporali del segnale BOLD da ciascun voxel. La \textbf{connettività funzionale} viene tipicamente quantificata tramite la \textbf{correlazione incrociata} (cross-correlation)  tra le serie temporali di coppie di voxel:

\begin{equation}
C_{ij} = \frac{\langle R_i(t)\, R_j(t) \rangle_t}{\sigma_i \sigma_j}
\end{equation}

dove $R_i(t)$ è il segnale BOLD del voxel $i$ e $\langle \cdot \rangle_t$ denota la media temporale.

\subsubsection{Misurazioni a Riposo e Non-Stazionarietà}

Le misure di connettività funzionale vengono tipicamente eseguite su soggetti in \textbf{stato di riposo} (resting state), in sessioni della durata di diversi minuti. In questo stato, il cervello non è garantito di trovarsi sempre nello stesso stato funzionale: può vagare tra diversi stati di attività (non-stazionarietà), il che influenza le misure di correlazione.

\subsection{Confronto tra Connettività Strutturale e Funzionale}

\subsubsection{Il Repository CoCoMac}

Un caso di studio standard nella letteratura di neuroscienze sistemiche è il \textbf{connettoma del macaco}, disponibile nel repository pubblico \textbf{CoCoMac} (Collation of Connectivity data on the Macaque brain). Questo dataset fornisce sia la matrice di connettività strutturale (da DTI e tracing anatomico) sia misure di connettività funzionale (da fMRI a riposo).

\subsubsection{Rappresentazione Inflata della Corteccia del Macaco}

La corteccia del macaco viene spesso rappresentata in forma \textbf{inflata}: la superficie corticale viene "sgonfiata" per rendere visibili le aree nascoste nelle fessure (sulci). In questa rappresentazione, le diverse \textbf{aree di Brodmann} (aree citoarchitettoniche) sono contrassegnate da colori diversi:

\begin{itemize}
    \item Polo occipitale (lobo occipitale): aree visive primarie e secondarie
    \item Lobo parietale: integrazione sensorimotoria
    \item Lobo temporale: processamento uditivo e del linguaggio
    \item Lobo frontale: pianificazione motoria e funzioni cognitive superiori; in particolare, la \textbf{corteccia premotoria} (area di Brodmann 6) gestisce la pianificazione del movimento
\end{itemize}

\subsubsection{Confronto tra le Due Matrici}

Il confronto tra la matrice strutturale e quella funzionale rivela un risultato fondamentale:

\begin{enumerate}
    \item \textbf{Corrispondenza parziale}: le coppie di aree con forte connettività strutturale tendono a mostrare elevata correlazione funzionale.
    \item \textbf{Connettività funzionale senza connessione strutturale diretta}: è possibile osservare correlazione (anche anti-correlazione) tra aree non direttamente connesse strutturalmente, a causa di percorsi polisinaptici.
    \item \textbf{Anti-correlazione}: alcune coppie di aree mostrano attività anti-correlata, fenomeno interpretabile come inibizione indiretta o competizione tra sottoreti funzionali.
\end{enumerate}

\begin{tcolorbox}[colback=gray!5, colframe=gray!40, arc=2mm, boxrule=0.5pt]
\textbf{Nota del docente:} Il cervello è un sistema integrato: a uno o due passi sinaptici, tutte le popolazioni corticali risultano indirettamente connesse l'una all'altra. Questo spiega perché possano emergere correlazioni funzionali anche in assenza di connessione strutturale diretta.
\end{tcolorbox}

\subsection{Microscala: Organizzazione in Strati della Colonna Corticale}

\subsubsection{Dalla Macroscala alla Microscala}

Dopo aver caratterizzato la rete corticale alla macroscala (aree corticali intere), è possibile scendere alla \textbf{microscala} di una singola colonna corticale. La colonna corticale è l'unità funzionale verticale della corteccia: un cilindro di tessuto corticale di circa 0.5 mm di diametro e spessore pari all'intera corteccia (circa 2–4 mm), organizzato in sei strati orizzontali (layers 1–6) con caratteristiche citoarchitettoniche distinte.

\subsubsection{Preparazione Istologica: Sezioni Corticali}

In neurofisiologia, un approccio classico consiste nel sacrificare l'animale, estrarre il cervello e ottenere sezioni istologiche dello spessore di circa 0.5 mm. Le sezioni vengono colorate con metodi specifici (ad esempio il \textbf{metodo Golgi} con inchiostro di osmio-bicromato di potassio) che impregnano selettivamente alcuni neuroni, rendendoli visibili come strutture nere su fondo chiaro.

Questa procedura permette di:
\begin{itemize}
    \item Ricostruire la \textbf{morfologia dettagliata} dei neuroni (soma, dendriti, assone)
    \item Identificare la \textbf{posizione laminare} dei neuroni
    \item Visualizzare la \textbf{distribuzione spaziale} delle diverse classi cellulari
\end{itemize}

\subsubsection{I Neuroni Piramidali dello Strato 5}

I \textbf{neuroni piramidali} sono i neuroni eccitatori principali della corteccia. Il loro soma ha una caratteristica forma piramidale, con un lungo apice dendritico che risale verso gli strati superficiali e ramificazioni basali estese orizzontalmente. I neuroni piramidali dello \textbf{strato 5} sono spesso descritti come il "motore" della colonna corticale: inviano proiezioni a lunga distanza verso il tronco encefalico, il midollo spinale e altre aree corticali.


\subsection{Mappatura della Connettività Intra-Colonnare}

\subsubsection{L'Esperimento Multi-Patch del Laboratorio Petersen}

Un esperimento seminale per caratterizzare la connettività intra-colonnare è stato condotto presso il \textbf{laboratorio Petersen} (EPFL, Losanna). In questo esperimento, vengono patchati simultaneamente i soma di \textbf{sei neuroni piramidali} identificati visivamente in una sezione corticale, ciascuno con una pipetta di patch-clamp indipendente. La tecnica di \textbf{patch-clamp whole-cell} consente di misurare il potenziale di membrana $V_i(t)$ di ciascuno dei sei neuroni in modo simultaneo e continuo.

\begin{tcolorbox}[colback=gray!5, colframe=gray!40, arc=2mm, boxrule=0.5pt]
\textbf{Nota del docente:} L'esperimento viene eseguito mantenendo il tessuto in silenzio (senza attività di background), per massimizzare il rapporto segnale-rumore nella misura dei potenziali post-sinaptici.
\end{tcolorbox}

\subsubsection{Protocollo di Stimolazione e Registrazione}

Il protocollo sperimentale è il seguente:

\begin{enumerate}
    \item Si inietta un \textbf{impulso di corrente} nel neurone pre-sinaptico $i$ (il "sender"), sufficiente a evocare un singolo potenziale d'azione.
    \item Si osservano i potenziali di membrana dei cinque neuroni rimanenti ($j = 1, \ldots, 5$, $j \neq i$) per rilevare eventuali \textbf{potenziali post-sinaptici eccitatori} (EPSP) con il ritardo atteso dalla trasmissione sinaptica.
    \item Se in un neurone $j$ si osserva una deflessione depolarizzante dopo il picco di $i$, si conclude che esiste una sinapsi eccitatoria da $i$ a $j$.
\end{enumerate}

Ripetendo questa procedura per ogni coppia ordinata di neuroni, si ricostruisce la \textbf{matrice di connettività sinaptica} del piccolo circuito.

Poiché si lavora con neuroni piramidali (neuroni eccitatori), tutti gli EPSP osservati sono depolarizzanti (positivi). La mappa di connettività prodotta permette di quantificare, per ogni coppia di strati corticali, la \textbf{probabilità di connessione} e la \textbf{dimensione dell'EPSP}.

\subsubsection{Scala dei Potenziali Post-Sinaptici}

Un punto tecnico importante riguarda la scala di misura: un \textbf{potenziale d'azione} ha ampiezza di circa 100 mV e durata di 1–2 ms. Un \textbf{EPSP} prodotto da una singola sinapsi ha ampiezza tipica di 0.1–1 mV e durata di 10–50 ms. Queste scale temporali e di ampiezza differiscono di due ordini di grandezza, il che richiede un amplificatore patch-clamp con elevata risoluzione in tensione.

Il ritardo sinaptico tipico tra il picco del potenziale d'azione pre-sinaptico e l'inizio dell'EPSP post-sinaptico è di circa \textbf{1–3 ms}, composto da:
\begin{itemize}
    \item \textbf{Ritardo di conduzione} lungo l'assone
    \item \textbf{Ritardo di trasmissione sinaptica} (esocitosi delle vescicole, diffusione del neurotrasmettitore, apertura dei recettori)
\end{itemize}

La presenza di questo ritardo con il timing corretto rispetto alla stimolazione è uno dei criteri per identificare un EPSP monosinaptico genuino.



\subsubsection{Matrice della Probabilità di Connessione}

Raggruppando i dati di connettività per strato di appartenenza del neurone pre- e post-sinaptico, si ottiene una \textbf{matrice di probabilità di connessione} $P(l_{\rm pre} \to l_{\rm post})$, dove gli indici $l_{\rm pre}$ e $l_{\rm post}$ identificano lo strato di origine e destinazione.

Questa matrice è \textbf{non uniforme}: alcuni strati sono preferenzialmente connessi tra loro, mentre altre combinazioni sono quasi assenti. In particolare, si osservano connessioni significative tra:

\begin{itemize}
    \item Strato 4 $\to$ Strato 5 (via di processamento principale)
    \item Strato 2/3 $\to$ Strato 5 (input superficiali allo "strato motore")
    \item Strato 2/3 $\to$ Strato 2/3 (ricorrenza all'interno degli strati superficiali)
    \item Strato 5 $\to$ Strato 6 (proiezione verso gli strati profondi)
\end{itemize}

Le aree \textbf{grigie} nella matrice corrispondono a combinazioni per le quali non sono stati trovati neuroni pre-sinaptici nell'esperimento; le aree \textbf{nere} indicano assenza di connessione.

\subsubsection{Matrice dell'Efficacia Sinaptica Media}

In aggiunta alla matrice delle probabilità, l'esperimento misura la \textbf{dimensione media dell'EPSP} per ciascuna combinazione di strati, fornendo una matrice dell'efficacia sinaptica $\langle J(l_{\rm pre} \to l_{\rm post}) \rangle$. Questa matrice è più sparsa: la colorazione rossa indica un valore medio di circa \textbf{0.5 mV}.

\subsubsection{Matrice della Forza di Connessione}

La \textbf{forza di connessione} tra due strati è determinata dal prodotto tra probabilità di connessione ed efficacia sinaptica media, moltiplicato per il numero di sinapsi:

\begin{equation}
W(l_{\rm pre} \to l_{\rm post}) = P_{\rm conn} \cdot \langle J \rangle \cdot K
\end{equation}

Questa quantità rappresenta la vera misura della capacità di uno strato di influenzare un altro. La matrice risultante mostra chiaramente che:

\begin{itemize}
    \item Lo \textbf{strato 5} integra input da tutti gli strati (soprattutto 2/3 e 4)
    \item Esiste una forte \textbf{ricorrenza} all'interno di ogni strato
    \item Lo strato 5b riceve input aspecifici da strati multipli, coerentemente con il suo ruolo integrativo
\end{itemize}

\subsubsection{Confronto tra Valore Teorico e Valore Sperimentale}

Il valore teorico dell'\textbf{EPSP} (potenziale post-sinaptico eccitatorio) atteso dal modello integra-e-spara è:

\begin{equation}
J_{\rm teorico} \approx \frac{V_{\rm thr}}{100} \approx \frac{20\,\text{mV}}{100} = 0.2\,\text{mV}
\end{equation}

dove il fattore 100 rappresenta il numero approssimativo di spike pre-sinaptici necessari a portare la membrana da $V_{\rm rest}$ a $V_{\rm thr}$.

Il valore \textbf{sperimentale} misurato nell'esperimento multi-patch è di circa \textbf{0.5 mV}, leggermente superiore.


La discrepanza è spiegabile considerando la \textbf{forza motrice della sinapsi eccitatoria} (driving force):

\begin{equation}
\Delta V_{\rm EPSP} = g_{\rm syn} \cdot (V - E_{\rm syn})
\end{equation}

dove $E_{\rm syn} \approx 0$ mV è il \textbf{potenziale di inversione} per le sinapsi eccitatorie (canali AMPA/NMDA, con corrente cationica mista).

\textbf{A riposo:} $V \approx -76$ mV, quindi la forza motrice è $(-76 - 0) = -76$ mV (grande valore assoluto), che produce un EPSP ampio.

\textbf{Durante l'attività normale:} Il bombardamento sinaptico da neuroni di background porta il potenziale di membrana più vicino alla soglia, $V \approx -55$ mV, riducendo la forza motrice a $(-55 - 0) = -55$ mV. Il valore $J = 0.2$ mV del modello teorico si riferisce a questa condizione.

Il rapporto di scala è quindi:

\begin{equation}
\frac{J_{\rm sperimentale}}{J_{\rm teorico}} \approx \frac{76}{55} \approx 1.4
\end{equation}

consistente con il rapporto osservato $0.5/0.35 \approx 1.4$.

\begin{tcolorbox}[colback=gray!5, colframe=gray!40, arc=2mm, boxrule=0.5pt]
\textbf{Nota del docente:} Questa non è una contraddizione, ma una conferma: il valore $J \approx V_{\rm thr}/100$ è valido per neuroni bombardati da input presinaptici (condizione fisiologica), non per neuroni quiescenti come in questo esperimento.
\end{tcolorbox}

\begin{equation}
\boxed{J_{\rm sperimentale} \approx 0.5\,\text{mV} \quad \text{vs.} \quad J_{\rm modello} \approx 0.2\,\text{mV}}
\end{equation}


\subsection{Mesoscala: Connettività Inter-Colonnare e Approccio Optogenetico}

Dopo aver caratterizzato la microscala (colonna singola), si passa alla \textbf{mesoscala}: la connettività tra colonne corticali adiacenti. Le colonne corticali possono essere immaginate come esagoni che pavimentano la superficie corticale; ognuno di questi esagoni si estende verticalmente attraverso tutti e sei gli strati.

L'obiettivo è descrivere come colonne adiacenti comunicano tra loro attraverso la \textbf{sostanza grigia} (gray matter), a distanza corta, distinta dalle connessioni a lunga distanza attraverso il white matter.

\subsubsection{Esperimento Optogenetico su Sezione di Topo}

L'approccio sperimentale per mappare la connettività inter-colonnare sfrutta la \textbf{optogenetica}. Un virus viene utilizzato per infettare i neuroni corticali di un topo e indurli a esprimere la \textbf{channelrhodopsina-2 (ChR2)}, un canale ionico fotosensibile (opsina). La ChR2 si apre in risposta alla luce blu (lunghezza d'onda $\sim$470 nm), iniettando corrente depolarizzante e inducendo la scarica del neurone.

Il protocollo sperimentale è:

\begin{enumerate}
    \item Si prepara una sezione corticale sottile dallo stesso topo
    \item I bordi tra i diversi strati sono visualizzati con linee tratteggiate (istologia)
    \item Si punta un \textbf{laser} in modo sistematico su diversi punti della sezione
    \item Si registra il potenziale di membrana di un \textbf{singolo neurone piramidale di strato 5b} con una pipetta di patch-clamp
\end{enumerate}

Ogni volta che il laser illumina un pixel della sezione, i neuroni presenti in quel pixel sparano. Se il neurone di strato 5b registrato riceve un EPSP in risposta a quella stimolazione, si conclude che esiste una connessione sinaptica tra i neuroni del pixel illuminato e il neurone bersaglio.

\begin{tcolorbox}[colback=gray!5, colframe=gray!40, arc=2mm, boxrule=0.5pt]
\textbf{Nota del docente:} Il neurone di strato 5b funge da "antenna": misurando le variazioni del suo potenziale di membrana in risposta alla stimolazione sistematica di tutti i punti della sezione, si mappa l'intera distribuzione spaziale dei neuroni pre-sinaptici.
\end{tcolorbox}


\subsubsection{Mappa della Probabilità di Connessione in Funzione della Distanza}

Il risultato dell'esperimento optogenetico è una mappa bidimensionale della probabilità di connessione: per ogni pixel della sezione, si conosce sia la distanza orizzontale dal neurone bersaglio, sia lo strato di appartenenza dei neuroni stimolati.

Tracciando la probabilità di connessione in funzione della distanza $d$ per ciascun strato, si ottiene un andamento che decade monotonicamente:

\begin{equation}
P_{\rm conn}(d) \propto e^{-d/\lambda_l}
\end{equation}

dove $\lambda_l$ è la \textbf{costante di decadimento spaziale} per lo strato $l$.

\subsubsection{Scale Spaziali $\lambda$ per i Diversi Strati}

I valori sperimentali delle costanti di decadimento spaziale sono:

\begin{table}[htbp]
    \centering
    \renewcommand{\arraystretch}{1.3}
    \begin{tabularx}{\textwidth}{|c|X|}
    \hline
    \textbf{Strato} & \textbf{Costante $\lambda$} \\
    \hline
    Strato 2/3 & $\lambda_{2/3} \approx 350\,\mu\text{m}$ \\
    \hline
    Strato 4 & $\lambda_4 \approx 260\,\mu\text{m}$ \\
    \hline
    Strato 5 & $\lambda_5 \lesssim \lambda_4$ (valore più basso) \\
    \hline
    Strato 6 & $\lambda_6 \gtrsim \lambda_{2/3}$ (valore elevato) \\
    \hline
    \end{tabularx}
\end{table}

Il pattern che emerge è:

\begin{itemize}
    \item Gli \textbf{strati superficiali} (2/3) e gli \textbf{strati profondi} (6) hanno la maggiore portata spaziale nel gray matter
    \item Lo \textbf{strato 4} e in particolare lo \textbf{strato 5} hanno portata spaziale più ridotta
    \item Questo significa che i neuroni di strato 5 sono relativamente \textbf{isolati} rispetto alle colonne adiacenti nel gray matter
\end{itemize}

\subsubsection{Interpretazione Funzionale}

I neuroni di strato 5 ricevono input dalle colonne adiacenti \textbf{indirettamente}, attraverso gli strati 2/3 e 6, che fungono da \textbf{strati di input} per l'integrazione mesoscalare. Lo strato 5, essendo il "motore" della colonna, integra segnali provenienti da un vasto territorio corticale senza essere direttamente accoppiato a breve distanza con le colonne vicine.

\subsection{Due Vie di Comunicazione Corticale}

Un risultato recente e importante nella neuroscienze sistemiche riguarda l'\textbf{asimmetria direzionale} nella comunicazione tra aree corticali distanti. Le due direzioni principali di comunicazione — dal polo posteriore (aree sensoriali) al polo anteriore (aree frontali) e viceversa — utilizzano preferenzialmente \textbf{vie laminari diverse}.

\subsubsection{Via Ascendente (Bottom-Up): Informazioni Sensoriali}

La via ascendente trasporta informazioni sensoriali dalle aree posteriori (occipitale, parietale, temporale) verso la corteccia frontale. Questa via è preferenzialmente mediata dagli \textbf{strati intra-granulari} e profondi (strati 4, 5, 6). L'informazione risale attraverso:

\begin{equation}
\text{Area sensoriale} \xrightarrow{\text{strati 4–6}} \text{Aree associative} \xrightarrow{} \text{Corteccia frontale}
\end{equation}

\subsubsection{Via Discendente (Top-Down): Controllo Frontale}

La via discendente trasporta segnali di controllo dalla corteccia frontale verso le aree posteriori, implementando funzioni di \textbf{attenzione selettiva}, \textbf{previsione} (predictive coding) e modulazione top-down. Questa via è preferenzialmente mediata dagli \textbf{strati superficiali} (strati 2/3):

\begin{equation}
\text{Corteccia frontale} \xrightarrow{\text{strati 2/3}} \text{Aree associative} \xrightarrow{} \text{Aree sensoriali}
\end{equation}

\newpage
\section{Reti di Reti e Popolazioni Omogenee}

\subsubsection{Il Problema della Scala}

Le osservazioni sperimentali descritte nelle sezioni precedenti forniscono una descrizione dettagliata e multi-scala dell'organizzazione corticale. Tuttavia, costruire un modello computazionale che includa tutti questi dettagli è proibitivo dal punto di vista computazionale. È necessario identificare le \textbf{approssimazioni biologicamente plausibili} che permettono di ridurre la complessità mantenendo i meccanismi essenziali.

L'approccio proposto è quello delle \textbf{reti di reti}: si considera la corteccia non come un insieme di neuroni individuali, ma come un insieme di \textbf{popolazioni} di neuroni. Ogni popolazione corrisponde a un gruppo di neuroni omogenei (stesso tipo, stesso strato, stessa area corticale) connessi tra loro secondo statistiche precise.

A livello della microscala (singola colonna corticale, singolo strato), i neuroni che si trovano nella stessa posizione laminare condividono:

\begin{itemize}
    \item La stessa morfologia
    \item Gli stessi parametri biofisici
    \item La stessa organizzazione di connettività
    \item Le stesse sorgenti di input
\end{itemize}

Questo giustifica l'introduzione del concetto di \textbf{popolazione omogenea di neuroni}.



\subsection{Definizione di Popolazione Omogenea}

Una \textbf{popolazione omogenea di neuroni} è un insieme di $N$ neuroni che condividono:

\begin{enumerate}
    \item Gli stessi parametri microscopici neuronali:
    \begin{itemize}
        \item $\tau_m$: costante di tempo della membrana
        \item $V_{\rm thr}$: potenziale di soglia
        \item $V_{\rm reset}$: potenziale di reset post-spike
        \item $V_{\rm rest}$: potenziale di riposo (o potenziale di equilibrio in assenza di input)
    \end{itemize}
    \item Lo stesso input statistico esterno:
    \begin{itemize}
        \item $\mu_{\rm ext}$: media della corrente di input proveniente da aree esterne alla popolazione
        \item $\sigma_{\rm ext}$: deviazione standard della corrente di input esterna
    \end{itemize}
    \item La stessa probabilità di connessione interna:
    \begin{itemize}
        \item $\varepsilon$: probabilità che esista una sinapsi da qualsiasi neurone $j$ al neurone $i$ all'interno della stessa popolazione ($\varepsilon \approx 20\%$)
    \end{itemize}
\end{enumerate}

\subsubsection{Input Esterno come Variabile Statistica}

L'input esterno alla popolazione proviene da:
\begin{itemize}
    \item Colonne corticali distanti (via white matter)
    \item Strati diversi della stessa colonna
    \item Input talamici e subcorticali
\end{itemize}

Tutti questi contributi vengono parametrizzati in modo sintetico dalla coppia $(\mu_{\rm ext}, \sigma_{\rm ext})$, ovvero come \textbf{media e varianza di un processo stocastico} che rappresenta l'input aggregato. Questo è giustificato dal fatto che, trattandosi di molte sorgenti indipendenti, il teorema centrale del limite garantisce la gaussianità dell'input totale.

\begin{tcolorbox}[colback=gray!5, colframe=gray!40, arc=2mm, boxrule=0.5pt]
\textbf{Nota del docente:} L'assunzione di popolazione omogenea è un'approssimazione: neuroni dello stesso strato non sono identici in natura. Tuttavia, i parametri biofisici variano entro un range relativamente stretto, e questa variabilità può essere incorporata come rumore aggiuntivo nel modello.
\end{tcolorbox}


\subsection{Distribuzione Binomiale del Numero di Contatti Sinaptici}


All'interno di una popolazione omogenea di $N$ neuroni con probabilità di connessione $\varepsilon$, il \textbf{numero di contatti sinaptici} $K$ che il neurone $i$ riceve da tutti gli altri $N-1$ neuroni della popolazione segue una \textbf{distribuzione binomiale}:

\begin{equation}
\boxed{P(K = k) = \binom{N}{k}\,\varepsilon^k\,(1-\varepsilon)^{N-k}}
\end{equation}


Il valore atteso è:

\begin{equation}
\langle K \rangle = \varepsilon\, N \approx 10^4
\end{equation}

Considerando $\varepsilon \approx 0.20$ e $N \approx 5 \times 10^4$ neuroni nella colonna (ordine di grandezza), si ottiene un numero di sinapsi per neurone dell'ordine di $10^4$, in accordo con i dati sperimentali.


La varianza della distribuzione binomiale è:

\begin{equation}
\text{Var}(K) = N\,\varepsilon\,(1-\varepsilon) = \langle K \rangle\,(1-\varepsilon)
\end{equation}

Il \textbf{coefficiente di variazione} (CV) del numero di sinapsi è:

\begin{equation}
\text{CV}(K) = \frac{\sqrt{\text{Var}(K)}}{\langle K \rangle} = \frac{\sqrt{N\,\varepsilon\,(1-\varepsilon)}}{N\,\varepsilon} = \sqrt{\frac{1-\varepsilon}{\varepsilon\,N}} = \sqrt{\frac{1-\varepsilon}{\langle K \rangle}}
\end{equation}

Sostituendo $\varepsilon \approx 0.20$ e $\langle K \rangle \approx 10^4$:

\begin{equation}
\text{CV}(K) = \sqrt{\frac{0.8}{10^4}} = \sqrt{8 \times 10^{-5}} \approx \frac{0.894}{100} \approx 0.9\%
\end{equation}

\begin{tcolorbox}[colback=gray!5, colframe=gray!40, arc=2mm, boxrule=0.5pt]
\textbf{Nota del docente:} Sebbene la probabilità di connessione sia solo del 20\%, grazie al gran numero di neuroni presenti nella popolazione il numero di sinapsi per neurone è quasi costante: la variabilità è solo dell'1\%. Questo è una manifestazione del \textbf{teorema centrale del limite} applicato alla distribuzione binomiale: con $N$ grande e $\varepsilon$ non troppo piccolo, la distribuzione binomiale si approssima a una gaussiana con piccola varianza relativa.
\end{tcolorbox}

\subsection{Distribuzione Statistica dell'Efficacia Sinaptica}

La forza della connessione tra due neuroni è determinata dall'\textbf{efficacia sinaptica} $w_{ij}$. Ci sono due fonti di variabilità:

\begin{enumerate}
    \item \textbf{Presenza/assenza della sinapsi}: con probabilità $1-\varepsilon$ non esiste alcuna sinapsi ($w_{ij} = 0$), con probabilità $\varepsilon$ esiste una sinapsi con efficacia non nulla.
    \item \textbf{Variabilità dell'efficacia tra sinapsi}: quando la sinapsi esiste, la sua efficacia segue una distribuzione gaussiana con media $J$ e deviazione standard $\Delta J$.
\end{enumerate}


La distribuzione di probabilità è data da:

\begin{equation}
\boxed{P(J_{ij} \in [w;w+dw]) = \left[ (1-\varepsilon)\,\delta(w) + \frac{\varepsilon}{\sqrt{2\pi}\,\Delta J}\,\exp\!\left(-\frac{(w - J)^2}{2\Delta J^2}\right)\right]dw}
\end{equation}

dove:
\begin{itemize}
    \item $(1-\varepsilon)\,\delta(w)$: contributo delle sinapsi assenti 
    \item Il termine gaussiano: distribuzione dell'efficacia delle sinapsi presenti, con media $J$ e varianza $\Delta J^2$
    \item La gaussiana è normalizzata a $\varepsilon$: $\int_{-\infty}^{+\infty} \frac{\varepsilon}{\sqrt{2\pi}\Delta J} e^{-(w-J)^2/2\Delta J^2} dw = \varepsilon$
\end{itemize}

\begin{tcolorbox}[colback=gray!5, colframe=gray!40, arc=2mm, boxrule=0.5pt]
\textbf{Nota del docente:} Il dominio dell'integrazione si estende da $-\infty$ a $+\infty$ per includere sia sinapsi eccitatorie ($J > 0$) sia inibitorie ($J < 0$). 
\end{tcolorbox}

La densità di probabilità è:
\begin{equation}
    \rho(w)=(1-\varepsilon)\,\delta(w) + \frac{\varepsilon}{\sqrt{2\pi}\,\Delta J}\,\exp\!\left(-\frac{(w - J)^2}{2\Delta J^2}\right)
\end{equation}

\subsection{Momenti della Distribuzione di w}

\subsubsection{Calcolo della Media $\langle w \rangle$}

Il \textbf{valore atteso} dell'efficacia sinaptica si calcola integrando $w\,\rho_j(w)$:

\begin{equation}
\langle w \rangle = \int_{-\infty}^{+\infty} w\,\rho_j(w)\,dw
\end{equation}

Il contributo del termine $\delta(w)$ è nullo (poiché $w\,\delta(w) = 0$ per $w \to 0$). Il contributo del termine gaussiano è:

\begin{equation}
\langle w \rangle = \int_{-\infty}^{+\infty} w \cdot \frac{\varepsilon}{\sqrt{2\pi}\,\Delta J}\,e^{-(w-J)^2/2\Delta J^2}\,dw = \varepsilon\,J
\end{equation}

poiché il valore atteso di una gaussiana con media $J$ normalizzata a $\varepsilon$ è $\varepsilon \cdot J$.

\begin{equation}
\boxed{\langle w \rangle = \varepsilon\, J}
\end{equation}

\subsubsection{Calcolo del Secondo Momento $\langle w^2 \rangle$}

Il \textbf{secondo momento} si calcola analogamente:

\begin{equation}
\langle w^2 \rangle = \int_{-\infty}^{+\infty} w^2\,\rho_j(w)\,dw = \int_{-\infty}^{+\infty} w^2 \cdot \frac{\varepsilon}{\sqrt{2\pi}\,\Delta J}\,e^{-(w-J)^2/2\Delta J^2}\,dw
\end{equation}

Usando la relazione tra momento secondo, varianza e quadrato della media per una distribuzione gaussiana:

\begin{equation}
\text{Var}(w) = \langle w^2 \rangle_{\rm gauss} - \langle w \rangle_{\rm gauss}^2 = \Delta J^2
\end{equation}

quindi $\langle w^2 \rangle_{\rm gauss} = \Delta J^2 + J^2$, e tenendo conto della normalizzazione a $\varepsilon$:

\begin{equation}
\boxed{\langle w^2 \rangle = \varepsilon\left(\Delta J^2 + J^2\right)}
\end{equation}

\subsubsection{Significato dei Due Momenti}

I due momenti $\langle w \rangle $ e $\langle w^2 \rangle$ sono le quantità fondamentali che caratterizzano la distribuzione statistica della connettività. Nella teoria di campo medio, questi momenti determinano rispettivamente:

\begin{itemize}
    \item La \textbf{corrente media} (drift) ricevuta da un neurone dalla popolazione
    \item La \textbf{varianza della corrente} (diffusione) ricevuta da un neurone dalla popolazione
\end{itemize}

Questi due parametri sono i coefficienti $\mu$ e $\sigma^2$ dell'equazione di Langevin che governa la dinamica del potenziale di membrana.

    

\chapter{Lezione 18}



\section{Verso l’Approssimazione di Campo Medio}

\subsubsection{Contesto}
Il punto di partenza è la considerazione che i neuroni non operano in isolamento, ma sono immersi in una rete in cui ogni neurone $i$ riceve una \textbf{corrente sinaptica} $I_i(t)$ generata dall'attività collettiva dei neuroni presinaptici che lo raggiungono.

Come discusso nelle lezioni precedenti, quando il numero di contatti presinaptici è sufficientemente elevato (dell'ordine di migliaia), la corrente sinaptica può essere descritta in maniera eccellente come un \textbf{processo di diffusione}, caratterizzato dai suoi momenti infinitesimali. Si considera una rete \textbf{omogenea}, ovvero una rete in cui tutti i neuroni condividono gli stessi parametri biofisici e in cui la statistica delle connessioni sinaptiche è identica per ogni coppia di neuroni. Questa assunzione di omogeneità è il punto di partenza per costruire una teoria di campo medio della rete.

\subsubsection{Momenti Infinitesimali della Corrente Sinaptica}
Per un processo di diffusione, i momenti infinitesimali si definiscono come i limiti per $dt \to 0$ dei cumulanti dell'incremento del processo. Per il neurone $i$:

\begin{equation}
\mu_i(t) = \lim_{dt \to 0} \frac{\mathbb{E}[dI_i]}{dt}
\end{equation}

\begin{equation}
\sigma_i^2(t) = \lim_{dt \to 0} \frac{\mathbb{V}[dI_i]}{dt}
\end{equation}

dove $dI_i$ è l'incremento della corrente sinaptica in un intervallo infinitesimo $dt$. Il parametro $\mu_i$ rappresenta la \textbf{media infinitesimale} (o drift) e $\sigma_i^2$ la \textbf{varianza infinitesimale} (o coefficiente di diffusione) della corrente sinaptica ricevuta dal neurone $i$.

\subsection{ Assunzione di Campo Medio}
La media infinitesimale si calcola come somma dei contributi di tutti i neuroni presinaptici $j$:

\begin{equation}
\mu_i(t) = \sum_{j=1}^{N} J_{ij} \, \nu_j(t)
\end{equation}

dove $J_{ij}$ è l'\textbf{efficacia sinaptica} (o peso sinaptico) dalla connessione $j \to i$, e $\nu_j(t)$ è la probabilità di emissione di uno spike da parte del neurone $j$ per unità di tempo, ovvero la sua \textbf{frequenza di scarica istantanea}. I treni di spike presinaptici sono trattati come processi di Poisson, il che è giustificato dal \textbf{teorema di Grigelionis} quando si sovrappongono contributi di un numero elevato di neuroni, anche se i singoli neuroni non emettono secondo un processo strettamente poissoniano.

L'\textbf{assunzione di campo medio} (mean field approximation) consiste nel considerare che le proprietà statistiche della scarica del neurone presinaptico $j$ siano solo debolmente influenzate dall'identità dello specifico contatto sinaptico $J_{ij}$. Sotto questa assunzione, l'emissione di uno spike da parte del neurone $i$ dovuta al neurone $j$ è trascurabile. Questo permette di considerare i pesi sinaptici $J_{ij}$ e gli incrementi degli spike $dN_j(t)$ come variabili aleatorie effettivamente indipendenti.

Questa approssimazione è giustificata da due argomenti:

\begin{enumerate}
    \item \textbf{Argomento di scala}: l'efficacia sinaptica di una singola sinapsi è:
    \begin{equation}
    J \approx \frac{V_{\rm thr}}{100} \approx \frac{20\,\text{mV}}{100} = 0.2\,\text{mV}
    \end{equation}
    Questo è un contributo trascurabile rispetto alla soglia: ogni singola sinapsi contribuisce per circa 1\% a un singolo spike.
    \item \textbf{Argomento di frequenza}: i neuroni corticali sparano a circa 1–10 Hz. La probabilità che il singolo spike prodotto dal neurone $j$ in un secondo sia la causa determinante dello spike del neurone $i$ è infinitesimale.
\end{enumerate}

Combinando i due argomenti: il contributo di una singola connessione alla decisione di sparare del neurone $i$ è dell'ordine di $J/V_{\rm thr} \cdot \nu_j \cdot \Delta t \ll 1$, rendendo trascurabili le correlazioni indotte dalla connettività.


\subsubsection{Perdita dell'Identità del Singolo Neurone in MFA}

Un'implicazione concettuale profonda dell'approssimazione di campo medio è la \textbf{perdita dell'identità} del singolo neurone presinaptivo. Nella descrizione esatta, il neurone $i$ riceve spike specificamente dal neurone $j$ tramite una sinapsi con efficacia $w_{ij}$ ben definita: la storia passata di $j$ influenza direttamente $i$.

Nell'approssimazione di campo medio, questa correlazione viene negletta: si assume che ogni spike ricevuto provenga da un neurone "tipico" della popolazione, estratto casualmente dalla distribuzione $\rho_j(w)$. I contributi di spike successivi sono trattati come \textbf{eventi indipendenti} (processo di Poisson).

Il risultato è che la corrente sinaptica totale diventa un processo stocastico caratterizzato solo dai suoi due momenti ($\mu$ e $\sigma^2$), e non dalla struttura di correlazione tra spike di neuroni distinti. Questo è un'approssimazione \textbf{valida nel limite di grande $N$ e piccolo $J/V_{\rm thr}$}, ma diventa meno accurata quando:
\begin{itemize}
    \item La rete è piccola ($N \lesssim 10^3$)
    \item Le sinapsi sono forti ($J \gtrsim V_{\rm thr}/10$)
    \item Esistono strutture di correlazione significative (es. oscillazioni sincrone)
\end{itemize}



\section{Distribuzione dei Coefficienti Sinaptici }

Per procedere nel calcolo del valor medio di $\mu_i$ sull'intera popolazione neuronale (ovvero al variare dell'indice $i$), è necessario conoscere la distribuzione statistica dei pesi sinaptici $J_{ij}$. 

\begin{equation}
Pr\{J_{ij} \in (w, w+dw)\} = \rho(w)dw = \left[ (1 - \epsilon)\delta(w) + \frac{\epsilon}{\sqrt{2\pi}\Delta J}e^{-(w-J)^2/(2\Delta J^2)} \right] dw
\end{equation}

ovvero:
\begin{equation}
\rho(w) = (1 - \epsilon)\delta(w) + \epsilon \mathcal{N}_{J, \Delta J}(w)
\end{equation}

dove:
\begin{itemize}
    \item $\epsilon = K/N$ è la \textbf{probabilità di connessione} tra due neuroni arbitrari della rete; $K$ è il numero medio di contatti presinaptici per neurone e $N$ è la dimensione totale della popolazione. Un valore tipico è $\epsilon \sim 0.2$, ovvero circa il $20\%$ dei neuroni della rete è presinapticamente connesso a un dato neurone postsinaptico.
    \item Con probabilità $1 - \epsilon$ i due neuroni non sono connessi, e il contributo sinaptico è nullo: $J_{ij} = 0$ (termine $\delta(w)$).
    \item Con probabilità $\epsilon$ i due neuroni sono connessi, e in tal caso l'efficacia sinaptica non nulla è estratta da una \textbf{gaussiana} di media $J$ e deviazione standard $\Delta J$: $\mathcal{N}_{J, \Delta J}(w)$.
\end{itemize}



\subsubsection{Momenti di $J_{ij}$}

Il calcolo del primo momento di $J_{ij}$ rispetto alla distribuzione introdotta è diretto. Il contributo del termine $\delta(w)$ è nullo, e quello del termine gaussiano fornisce:

\begin{equation}
\mathbb{E}[J_{ij}] = \int_{-\infty}^{+\infty} w \, \rho(w) \, dw = \epsilon \int_{-\infty}^{+\infty} w \mathcal{N}_{J, \Delta J}(w) dw = \epsilon J
\end{equation}

Il secondo momento si calcola analogamente, ricordando che per la variabile aleatoria associata all'efficacia sinaptica $w \sim \mathcal{N}(J, \Delta J^2)$:

\begin{equation}
\mathbb{E}[w^2] = \mathbb{V}[w] + (\mathbb{E}[w])^2 = \Delta J^2 + J^2 = J^2 \left( 1 + \left(\frac{\Delta J}{J}\right)^2 \right)
\end{equation}

Si ottiene quindi per l'intera distribuzione $\rho(w)$:

\begin{equation}
\mathbb{E}[J_{ij}^2] = \int_{-\infty}^{+\infty} w^2 \rho(w) dw = \epsilon \mathbb{E}[w^2] = \epsilon \left( \Delta J^2 + J^2 \right) = \epsilon J^2 \left( 1 + \left(\frac{\Delta J}{J}\right)^2 \right)
\end{equation}

dove si è evidenziata la varianza relativa dell'efficacia sinaptica $\left(\frac{\Delta J}{J}\right)^2$. Il docente osserva che in fisiologia è ragionevole assumere $\Delta J \sim \frac{1}{4}J \Rightarrow \left(\frac{\Delta J}{J}\right)^2 \approx 0.0625 \ll 1$.

La varianza si ricava dalla definizione standard:

\begin{equation}
\mathbb{V}[J_{ij}] = \mathbb{E}[J_{ij}^2] - (\mathbb{E}[J_{ij}])^2 = \epsilon J^2 \left( 1 + \left(\frac{\Delta J}{J}\right)^2 \right) - \epsilon^2 J^2 = \epsilon J^2 \left[ 1 - \epsilon + \left(\frac{\Delta J}{J}\right)^2 \right]
\end{equation}

\section{Media Infinitesimale}
\subsection{Valore Atteso di $\mu_i$ nella Popolazione}

Si calcola ora il valor medio di $\mu_i$ al variare dell'indice postsinaptico $i$, ovvero mediando sulla statistica (spaziale) della popolazione:

\begin{equation}
\mathbb{E}_i[\mu_i] = \mathbb{E}_i\left[ \sum_{j=1}^{N} J_{ij} \nu_j(t) \right]
\end{equation}

Trattando $J_{ij}$ come variabile aleatoria indipendente dalle frequenze $\nu_j$, l'operatore di media si applica solo ai pesi:

\begin{equation}
\mathbb{E}_i[\mu_i] = \sum_{j=1}^{N} \mathbb{E}_i[J_{ij}] \nu_j(t) = \epsilon J \sum_{j=1}^{N} \nu_j(t)
\end{equation}

Definendo la \textbf{frequenza di scarica media della popolazione}:

\begin{equation}
\nu(t) := \frac{1}{N} \sum_{j=1}^{N} \nu_j(t)
\end{equation}

Riconoscendo che $\epsilon N = K$ (numero medio di contatti presinaptici), si ottiene:

\begin{equation}
\mathbb{E}_i[\mu_i] = \epsilon N J \nu(t) = K J \nu(t)
\end{equation}

Questo risultato è coerente con l'approssimazione diffusiva: la corrente media ricevuta da ogni neurone dipende solo dall'attività media della popolazione. Questo è il \textbf{campo medio} (mean field) della rete.

\subsection{Varianza di $\mu_i$}

Per quantificare quanto bene l'approssimazione $\mu_i \approx \mathbb{E}_i[\mu_i] = \mu$ rappresenti ogni singolo neurone, calcoliamo la varianza di $\mu_i$ (utilizzando il teorema della varianza di combinazioni lineari di variabili aleatorie indipendenti $J_{ij}$):

\begin{equation}
\mathbb{V}_i[\mu_i] = \mathbb{V}_i\left[ \sum_{j=1}^{N} J_{ij} \nu_j(t) \right] = \sum_{j=1}^{N} \nu_j^2(t) \cdot \mathbb{V}_i[J_{ij}]
\end{equation}

Analizzando la somma dei quadrati delle frequenze e riscrivendola come valore atteso sulla popolazione:

\begin{equation}
\sum_{j=1}^{N} \nu_j^2(t) = N \frac{1}{N} \sum_{j=1}^{N} \nu_j^2(t) = N \mathbb{E}_j[\nu_j^2] = N \left( \mathbb{V}_j[\nu_j] + \mathbb{E}_j[\nu_j]^2 \right) = N(\Delta \nu^2 + \nu^2) = N \nu^2 \left( 1 + \left(\frac{\Delta \nu}{\nu}\right)^2 \right)
\end{equation}

Sostituendo l'espressione della varianza dei pesi e la sommatoria, otteniamo:

\begin{equation}
\mathbb{V}_i[\mu_i] = \epsilon J^2 \left[ 1 - \epsilon + \left(\frac{\Delta J}{J}\right)^2 \right] N \nu^2 \left( 1 + \left(\frac{\Delta \nu}{\nu}\right)^2 \right)
\end{equation}

Sapendo che $1 - \epsilon \approx 1$ e riordinando i termini con $\epsilon N = K$:

\begin{equation}
\mathbb{V}_i[\mu_i] \simeq K J^2 \nu^2 \left[ 1 + \left(\frac{\Delta J}{J}\right)^2 \right] \left[ 1 + \left(\frac{\Delta \nu}{\nu}\right)^2 \right]
\end{equation}

\subsection{Rapporto Segnale-Rumore}
Il \textbf{rapporto segnale-rumore} (SNR) delle fluttuazioni di $\mu_i$ tra i vari neuroni è:

\begin{equation}
\text{SNR}^{-1} = \frac{\sqrt{\mathbb{V}_i[\mu_i]}}{|\mathbb{E}_i[\mu_i]|} = \frac{J \sqrt{K} \nu \sqrt{1 + \left(\frac{\Delta J}{J}\right)^2} \sqrt{1 + \left(\frac{\Delta \nu}{\nu}\right)^2}}{K J \nu} \sim \frac{1}{\sqrt{K}}
\end{equation}

Essendo il numero di connessioni $K \sim 10^4$ nella corteccia, la variabilità relativa tra neuroni è estremamente piccola ($\sim 1/\sqrt{10000} = 1\%$). Di conseguenza, è un'eccellente approssimazione sostituire i vari $\mu_i$ con un singolo valor medio condiviso: $\mu_i \simeq \mathbb{E}_i[\mu_i] = \mu$.

\section{Varianza Infinitesimale}


Procedendo analogamente per la varianza infinitesimale (il coefficiente di diffusione del singolo neurone):

\begin{equation}
\sigma_i^2 = \lim_{dt \to 0} \frac{\mathbb{V}[dI_i]}{dt} = \sum_{j=1}^{N} J_{ij}^2 \nu_j(t)
\end{equation}

Il suo valore atteso sulla popolazione è:

\begin{equation}
\mathbb{E}_i[\sigma_i^2] = \sum_{j=1}^{N} \mathbb{E}_i[J_{ij}^2] \nu_j(t) = \epsilon J^2 \left( 1 + \left(\frac{\Delta J}{J}\right)^2 \right) \sum_{j=1}^{N} \nu_j(t)
\end{equation}

Riconoscendo che $\sum \nu_j(t) = N \nu(t)$ e $\epsilon N = K$:

\begin{equation}
\mathbb{E}_i[\sigma_i^2] = K J^2 \left( 1 + \left(\frac{\Delta J}{J}\right)^2 \right) \nu(t) = \sigma^2(t)
\end{equation}

Questa formula generalizza quella derivata nell'approssimazione diffusiva pura; ponendo $\Delta J = 0$ si riottiene l'espressione esatta del caso omogeneo assoluto.

\subsection{Rete Bilanciata e Finitezza della Varianza}
La condizione che $\sigma^2$ sia finita nel limite $N \to \infty$ a $K$ costante è garantita dall'ipotesi di \textbf{rete bilanciata} (balanced network). In una rete bilanciata, il contributo eccitatorio e quello inibitorio al drift medio tendono a cancellarsi parzialmente, garantendo che le fluttuazioni (la varianza infinitesimale) rimangano finite e che il regime della rete sia asincrono e irregolare. Come per $\mu_i$, anche le fluttuazioni relative $\mathbb{V}_i[\sigma_i^2]$ scalano come $1/K$, giustificando l'impiego del valore $\sigma^2$ per tutti i neuroni.

\section{Approssimazione di Campo Medio Estesa (Extended MFA)}

Avendo stabilito che $\mu_i \sim \mu$ e $\sigma_i \sim \sigma$, la dinamica del potenziale di membrana per ciascun neurone $i$ può essere descritta dall'equazione stocastica differenziale sotto l'\textbf{approssimazione di campo medio estesa}:

\begin{equation}
dV_i \simeq \left( -\frac{V_i}{\tau_m} + \mu \right)dt + \sigma dW_i(t)
\end{equation}


Tutti i neuroni seguono la stessa dinamica integrata e di decadimento guidata dal drift comune $\mu$ e dal termine di rumore $\sigma$. I processi stocastici di Wiener $W_i$ presentano la seguente statistica:

\begin{equation}
\mathbb{E}[dW_i(t) dW_j(t')] = \delta_{ij} \delta(t - t') dt
\end{equation}

Da questa relazione si evince che:
\begin{itemize}
    \item Non c'è correlazione nel rumore tra neuroni differenti ($\delta_{ij}$).
    \item Non c'è correlazione temporale nel rumore per lo stesso neurone (rumore bianco gaussiano standard, $\delta(t-t')$).
\end{itemize}


Nell'MFA estesa, ogni neurone del network è visto come una \textbf{realizzazione indipendente dello stesso processo stocastico}. La dipendenza da $t$ dei termini $\mu$ e $\sigma^2$ è puramente \textbf{parametrica} attraverso l'attività media collettiva $\nu(t)$. Questo fornisce una cornice teorica formidabile per studiare il sistema anche in condizioni non stazionarie.

\subsection{Rete di Reti}

Una rete neuronale corticale (es. un layer) non è un sistema isolato. Interagisce con altri layer o riceve afferenze dall'esterno. I momenti infinitesimali ricevono così componenti addizionali non derivanti dall'attività ricorrente:

\begin{equation}
\mu(\nu, t) = K J \nu(t) + \mu_{ext}(t)
\end{equation}

\begin{equation}
\sigma^2(\nu, t) = K J^2 \left( 1 + \left(\frac{\Delta J}{J}\right)^2 \right) \nu(t) + \sigma_{ext}^2(t)
\end{equation}

Il termine $\mu_{ext}(t)$ include influenze come la \textbf{stimolazione transcranica a corrente continua} (tDCS). Un campo elettrico applicato polarizza l'eccitabilità neuronale dipendente dalla morfologia dendritica rispetto alle linee di campo. Mentre correnti estreme causano shock, campi moderati sono usati in clinica per aumentare la plasticità, modulando $\mu_{ext}(t)$. Negli scenari comuni, queste variazioni esterne vengono assunte come costanti a tratti (stepwise changes).

\section{Equazione di Auto-Consistenza per la Frequenza di Scarica}

Il comportamento stazionario collettivo si trova chiudendo il loop: input dipendente dall'output e viceversa. Definiamo la funzione di guadagno $\Phi(\mu, \sigma)$ del processo stocastico:

\begin{equation}
\nu = \frac{1}{\text{ISI}} = \Phi(\mu(\nu), \sigma(\nu))
\end{equation}

dove l'equazione di \textbf{auto-consistenza} $\nu^* = \Phi(\nu^*)$ determina i punti fissi della rete.


Per il modello LIF con rumore gaussiano e frequenza asintotica, il rate (inverso del tempo di primo passaggio) è dato dall'equazione di Ricciardi-Siegert:

\begin{equation}
\Phi(\mu, \sigma) = \left[ \tau_0 + \tau \sqrt{\pi} \int_{\frac{V_{res} - \mu\tau}{\sigma\sqrt{\tau}}}^{\frac{V_{thr} - \mu\tau}{\sigma\sqrt{\tau}}} e^{z^2} (1 + \text{erf}(z)) dz \right]^{-1}
\end{equation}

dove $\tau_0$ è il periodo refrattario, che impone la frequenza limite massima del network.


La funzione di trasferimento esibisce un comportamento \textbf{sigmoidale} dipendente dalla varianza:
\begin{itemize}
    \item \textbf{Regime Supra-Soglia:} Dominato dal parametro $\mu$, dove $V_{THR} \approx 20\text{mV} < \mu\tau$. La frequenza cresce in base al drift medio.
    \item \textbf{Regime Sotto-Soglia:} L'elemento trainante è la varianza $\sigma^2$. Il sistema spara solo grazie al rumore stocastico superando la barriera. Aumentando $\sigma$, gli spike persistono anche con $\mu\tau \ll V_{THR}$.
\end{itemize}

\subsection{Analisi Geometrica dell'Equazione di Auto-Consistenza}

Esplicitando $\nu$ dall'equazione della corrente:

\begin{equation}
\mu = K J \nu + \mu_{ext} \implies K J \nu = \mu - \mu_{ext}
\end{equation}

Portando al piano $\mu\tau$ come variabile indipendente $x$:

\begin{equation}
\nu(x=\mu\tau) = \frac{1}{K J \tau} x - \frac{\mu_{ext}}{K J \tau}
\end{equation}

Quest'equazione di retta descrive l'input richiesto. L'intersezione tra questa linea retta (la cui pendenza è $\propto 1/J$) e la curva sigmoidale $\Phi(\mu, \sigma)$ produce i punti fissi stazionari.

Modificando i pesi $J$, l'inclinazione della retta varia, aprendo scenari di rottura di simmetria (simili ai sistemi ferromagnetici) in cui emergono molteplici punti fissi (tipicamente fino a tre incroci): due nodi stabili e un punto a sella instabile.
La condizione matematica per la stabilità locale del punto fisso all'intersezione $\nu^*$ è data dalla pendenza del guadagno rispetto all'input:

\begin{equation}
\left. \frac{d\Phi}{d\nu} \right|_{\nu^*} < 1
\end{equation}

Questa bistabilità dinamica (attività alta vs. silente) è centrale per modellare stati mentali discreti come la memoria di lavoro.

\section{Sinapsi Lente: NMDA e GABA-B}

Rispetto al paradigma esposto:
\begin{itemize}
    \item \textbf{Sinapsi Veloci:} AMPA e $\text{GABA}_A$, con costanti $\tau_F \sim 3\text{ ms}$. Possono essere trattate come spike $\delta$-Dirac impulsivi.
    \item \textbf{Sinapsi Lente:} NMDA e $\text{GABA}_B$, con $\tau_S \sim 70-100\text{ ms}$. La cinetica prolungata funge da memoria e inerzia temporale per il sistema.
\end{itemize}


Il contributo lento si modella aggiungendo una corrente $\mathcal{I}$ regolata dal processo di Ornstein-Uhlenbeck guidata dal network:

\begin{equation}
d\mathcal{I} = \left(-\frac{\mathcal{I}}{\tau_S} + J_S K \nu\right)dt + \sqrt{J_S^2 K \nu} dW
\end{equation}

I valori asintotici di media e deviazione standard sono:

\begin{equation}
\mathbb{E}[\mathcal{I}_{ss}] = K J_S \nu \tau_S, \quad \sigma_{\mathcal{I}} = \sqrt{\frac{K J_S^2 \nu \tau_S}{2}}
\end{equation}

\subsection{Rapporto SNR: Sinapsi Veloci vs. Lente}

Valutando il rapporto segnale-rumore (fluttuazioni relative) per le correnti lente:

\begin{equation}
\text{SNR}_{S}^{-1} = \frac{\sigma_{\mathcal{I}}}{\mathbb{E}[\mathcal{I}_{ss}]} = \frac{\sqrt{K J_S^2 \nu \tau_S / 2}}{K J_S \nu \tau_S} \propto \frac{1}{\sqrt{K \nu \tau_S}}
\end{equation}

Il confronto con le fluttuazioni derivate dalle sinapsi veloci restituisce:

\begin{equation}
\frac{\text{SNR}_{S}^{-1}}{\text{SNR}_{F}^{-1}} = \frac{\sqrt{K \nu \tau_F}}{\sqrt{K \nu \tau_S}} = \sqrt{\frac{\tau_F}{\tau_S}} \sim \sqrt{\frac{3}{100}} \approx \frac{1}{5.8}
\end{equation}

Le correnti lente possiedono fluttuazioni quasi 6 volte minori, fornendo una giustificazione rigorosa per l'\textbf{approssimazione quasi-adiabatica}: all'interno del limitato orizzonte dell'ISI, le sinapsi lente forniscono un offset quasi interamente deterministico.

\subsection{Dinamica della Rete con la Corrente Sinaptica Lenta}

Esplicitando l'apporto lento come dinamica addizionale ma approssimata deterministica:

\begin{equation}
\tau_m dV_i = \left(-V_i + \mu_{\text{fast}} + \mathcal{I}\right)dt + \sigma_{\text{fast}} dW_i
\end{equation}

Sotto condizioni di quasi equilibrio l'evoluzione di network forma il sistema:

\begin{equation}
\begin{cases}
\nu(\mathcal{I}) = \Phi(\mu_{\text{fast}} + \mathcal{I}, \sigma_{\text{fast}}^2) = \Phi(K J \nu(\mathcal{I}) + \mathcal{I}, J^2 K \nu(\mathcal{I})) \\
\tau_S \dot{\mathcal{I}} = -\mathcal{I} + \tau_S K J_S \nu(\mathcal{I})
\end{cases}
\end{equation}

L'equazione di stazionarietà per il punto fisso $\mathcal{I}^*$ risulta essere:

\begin{equation}
\mathcal{I}^* = \tau_S K J_S \nu(\mathcal{I}^*)
\end{equation}

\subsubsection{Analisi dei Punti Fissi e Linearizzazione}

Sottoponendo il sistema a una perturbazione locale:

\begin{equation}
\mathcal{I}(t) = \mathcal{I}^* + \delta\mathcal{I}(t)
\end{equation}

Sviluppiamo la linearizzazione dell'equazione dinamica:

\begin{equation}
\tau_S \delta\dot{\mathcal{I}} = -\delta\mathcal{I} + \tau_S K J_S \nu'(\mathcal{I}^*) \delta\mathcal{I} \implies \delta\dot{\mathcal{I}} = \frac{1}{\tau_S} \left( -1 + \tau_S K J_S \nu'(\mathcal{I}^*) \right) \delta\mathcal{I}
\end{equation}

Per l'estinzione asintotica delle perturbazioni ($\delta\dot{\mathcal{I}} / \delta\mathcal{I} < 0$), si ha:

\begin{equation}
-1 + \tau_S K J_S \nu'(\mathcal{I}^*) < 0 \implies \tau_S K J_S \nu'(\mathcal{I}^*) < 1
\end{equation}

\subsubsection{Condizione di Stabilità e Funzione Guadagno Effettiva}

La derivata totale $\frac{d\nu}{d\mathcal{I}}$ all'equilibrio si disaccoppia secondo la regola della catena nelle derivate parziali di media e inviluppo di varianza:

\begin{equation}
\frac{d\nu}{d\mathcal{I}}\bigg|_{\mathcal{I}^*} = \frac{\partial \Phi}{\partial \mu}\bigg|_{\mathcal{I}^*} \left[ K J \nu'(\mathcal{I}^*) + 1 \right] + \frac{\partial \Phi}{\partial \sigma^2}\bigg|_{\mathcal{I}^*} \left[ J^2 K \nu'(\mathcal{I}^*) \right]
\end{equation}

Per facilitare l'analisi si definisce la \textbf{Funzione Guadagno Effettiva}:

\begin{equation}
\tilde{\Phi}(\nu) = \Phi(K J \nu + \tau_S J_S K \nu, J^2 K \nu)
\end{equation}

la cui differenziazione totale rispetto a $\nu$ all'equilibrio fornisce:

\begin{equation}
\frac{d\tilde{\Phi}}{d\nu}\bigg|_{\nu^*} = \frac{\partial \Phi}{\partial \mu}\bigg|_{\nu^*} (K J + \tau_S J_S K) + \frac{\partial \Phi}{\partial \sigma^2}\bigg|_{\nu^*} (J^2 K)
\end{equation}

Con passaggi elementari, si dimostra che la precedente restrizione sulla sensibilità $\tau_S K J_S \nu'(\mathcal{I}^*) < 1$ è rigorosamente \textbf{equivalente a}:

\begin{equation}
\tilde{\Phi}'(\nu^*) < 1
\end{equation}

Questa versione ricattura il postulato geometrico espresso precedentemente (la retta tangente alla funzione di guadagno deve avere pendenza minore della bisettrice), integrando formalmente la dicotomia tra feedback veloci e inerzia stocastica lenta.



Il risultato stabilisce che lo stato di un network (bistabile o meno) riposa essenzialmente sulla forma di $\tilde{\Phi}$. Sapendo che la frequenza incrementa sia all'aumentare dell'input deterministico ($\frac{\partial\Phi}{\partial\mu} > 0$) sia della varianza ($\frac{\partial\Phi}{\partial\sigma^2} > 0$), la condizione:

\begin{equation}
\tilde{\Phi}'(\nu^*) < 1
\end{equation}
è costantemente minacciata dall'auto-amplificazione dell'eccitazione. 

Per questo l'\textbf{inibizione} è vitale: operando sottrattivamente nel dominio dei tassi, riduce le pendenze effettive e permette alla rete di sostenere tassi attivi moderati anziché collassare sul silenzio totale o scivolare patologicamente in saturazione neurodinamica.


\chapter{Lezione 19}

\section{Recap}

Il punto di partenza per l'analisi delle reti corticali è definire le condizioni sotto le quali un'estesa rete di neuroni può essere considerata \textbf{omogenea}. In una rete omogenea, tutti i neuroni condividono le stesse proprietà statistiche e dinamiche. Più formalmente, si richiede che:
\begin{enumerate}
    \item Tutti i neuroni abbiano gli stessi parametri identici: $\tau_{m,i} = \tau_m$, $V_{thr,i} = V_{thr}$, $V_{res,i} = V_{res}$ per ogni $i \in [1, N]$.
    \item Tutti i neuroni ricevano input esterni con le stesse proprietà statistiche: $\mathbb{E}[I_{ext,i}] = \mu_{ext}$ e $Var[I_{ext,i}] = \sigma_{ext}^2$.
    \item La connettività della rete sia statisticamente omogenea, ovvero la probabilità di connessione dipenda al più da regole spaziali condivise ma non da eterogeneità individuali: $Pr\{J_{ij} \in (x, x+dx)\} \equiv \rho_{ij}(x)dx = \rho(x)dx$.
\end{enumerate}
Se queste condizioni sono soddisfatte (se non lo sono, la rete è \textit{eterogenea}), l'approssimazione di campo medio è pienamente applicabile.

\subsubsection{Approssimazione di Campo Medio Estesa}

Sotto queste ipotesi, la dinamica del potenziale di membrana del singolo neurone è descritta da un'equazione stocastica (equazione di Langevin):

\begin{equation}
dV_i \simeq \left(-\frac{V_i}{\tau} + \mu\right)dt + \sigma \, dW_i
\end{equation}

dove $dW_i$ è un processo di Wiener tale che $\mathbb{E}[dW_i(t)dW_j(t+\tau)] = \delta_{ij}\delta(\tau)dt$. La media $\mu$ e la varianza $\sigma^2$ dell'input sinaptico dipendono dal tasso di firing di rete $\nu(t)$:

\begin{equation}
\begin{cases}
\mu(t) = JK\nu(t) = \mu(\nu) \\
\sigma^2(t) = J^2\left(1 + \frac{\Delta J^2}{J^2}\right)K\nu = \sigma^2(\nu)
\end{cases}
\end{equation}

\section{Condizione di Stabilità con Correnti Veloci e Lente}

\subsubsection{Dinamica della Corrente Lenta}

Consideriamo una rete in cui sono presenti sia correnti sinaptiche veloci (che contribuiscono in modo quasi istantaneo) sia correnti sinaptiche lente (es. recettori NMDA). La frequenza di firing della rete $\nu$ è determinata sia dall'input ricorrente rapido sia dalla componente lenta $\mathcal{I}$:

\begin{equation}
\nu(\mathcal{I}) = \Phi\!\left(KJ\nu(\mathcal{I}) + \mathcal{I},\, J^2K\nu(\mathcal{I})\right)
\end{equation}

La corrente lenta evolve nel tempo avvicinandosi al valore di stato stazionario dettato dal tasso di firing, con costante di tempo $\tau_S$:

\begin{equation}
\dot{\mathcal{I}} = -\frac{\mathcal{I}}{\tau_S} + J_S K \nu(\mathcal{I})
\end{equation}

Un \textbf{punto fisso} del sistema si ha quando la corrente lenta bilancia esattamente l'input ricorrente che essa stessa genera ($\dot{\mathcal{I}} = 0$), da cui la condizione di auto-consistenza:

\begin{equation}
\tau_S J_S K \nu(\mathcal{I}^*) = \mathcal{I}^*
\end{equation}

\subsubsection{Perturbazione e Condizione di Stabilità}

Per verificare se questo punto fisso è stabile, si studia l'evoluzione di una piccola perturbazione $\delta\mathcal{I}$ attorno all'equilibrio, ponendo $\mathcal{I}(t) = \mathcal{I}^* + \delta\mathcal{I}(t)$:

\begin{equation}
\delta\dot{\mathcal{I}} = -\frac{\mathcal{I}^* + \delta\mathcal{I}}{\tau_S} + J_S K \left[\nu(\mathcal{I}^*) + \nu'(\mathcal{I}^*)\delta\mathcal{I} + \mathcal{O}(\delta\mathcal{I}^2)\right]
\end{equation}

Tenendo conto della condizione di equilibrio, l'equazione linearizzata diventa:

\begin{equation}
\delta\dot{\mathcal{I}} = -\frac{\delta\mathcal{I}}{\tau_S} + J_S K \nu'(\mathcal{I}^*) \delta\mathcal{I}
\end{equation}

La stabilità richiede che la perturbazione decada nel tempo ($\partial I > 0 \Rightarrow \delta\dot{I} < 0$). Moltiplicando per $\tau_S$ e isolando il termine dipendente da $\delta\mathcal{I}$, otteniamo:

\begin{equation}
\tau_S \frac{\delta\dot{\mathcal{I}}}{\delta\mathcal{I}} = -1 + \tau_S J_S K \nu'(\mathcal{I}^*) < 0
\end{equation}

da cui deriva la fondamentale \textbf{condizione di stabilità}:

\begin{equation}
\tau_S J_S K \nu'(\mathcal{I}^*) < 1
\end{equation}

\subsubsection{La Funzione Guadagno Effettiva $\tilde{\Phi}$}

Per un sistema in cui l'input $\mathcal{I}$ è dominato dall'attività di rete, si può introdurre una \textbf{funzione guadagno effettiva} $\tilde{\Phi}(\nu)$, che esprime la dipendenza totale della frequenza di firing considerando sia il contributo della media $\mu$ sia quello della varianza $\sigma^2$:

\begin{equation}
\tilde{\Phi}'(\nu) = \frac{\partial\Phi}{\partial\mu} \frac{d\mu}{d\nu} + \frac{\partial\Phi}{\partial\sigma^2} \frac{d\sigma^2}{d\nu}
\end{equation}

Sviluppando le derivate:

\begin{equation}
\tilde{\Phi}'(\nu) = (\partial_\mu\Phi) \cdot \left[KJ + \tau_S J_S K\right] + (\partial_{\sigma^2}\Phi) \cdot \left[J^2 K\right]
\end{equation}

La stabilità del punto fisso dipende dunque dalla pendenza della funzione guadagno effettiva: il sistema è stabile se $\tilde{\Phi}'(\nu^*) < 1$. Una pendenza ripida (maggiore di 1, che in un grafico $\Phi(\nu)$ vs $\nu$ corrisponde a una retta più ripida di $45^\circ$) porta inevitabilmente a un'amplificazione a valanga delle perturbazioni (instabilità).


\section{Approssimazione della Funzione Guadagno nel Regime Noise-Dominated}

\subsubsection{Espressione Integrale Esatta}

Per analizzare concretamente la stabilità, il docente riprende la \textbf{funzione guadagno di Ricciardi-Siegert}, derivata risolvendo l'equazione di Fokker-Planck stazionaria per il neurone LIF sottoposto a corrente stocastica. Essa fornisce il tempo medio di primo passaggio $T_1$ dalla barriera di reset $V_r$ alla soglia $V_\theta$, e dunque la frequenza di scarica stazionaria $\nu = 1/T_1$:

\begin{equation}
\Phi(\mu, \sigma^2) = \left[\tau_{\text{ref}} + \tau\sqrt{\pi} \int_{\alpha_r}^{\alpha_\theta} e^{z^2}\!\left(1 + \text{erf}(z)\right) dz \right]^{-1}
\end{equation}

I limiti di integrazione adimensionali sono:

\begin{equation}
\alpha_r = \frac{V_r - \mu\tau}{\sigma\sqrt{\tau}}, \qquad \alpha_\theta = \frac{V_\theta - \mu\tau}{\sigma\sqrt{\tau}} \equiv x_T
\end{equation}

Questi quantificano la distanza tra il potenziale di reset (o di soglia) e il \textbf{valore asintotico del potenziale di membrana} $V_\infty = \mu\tau$ (la posizione verso cui il potenziale tenderebbe in assenza di barriere), normalizzata per l'ampiezza delle fluttuazioni $\sigma\sqrt{\tau}$.


La frequenza di scarica dipende in modo non banale dalla posizione di $V_\infty$ rispetto a $V_r$ e $V_\theta$: più $V_\infty = \mu\tau$ si avvicina a $V_\theta$ (al crescere di $\mu$), più la frequenza di scarica cresce rapidamente.



\subsubsection{Il Regime Noise-Dominated}

Nel \textbf{regime dominato dal rumore}, il valore asintotico del potenziale è significativamente inferiore alla soglia: $\mu\tau \ll V_\theta$. Gli spike vengono emessi esclusivamente grazie alle fluttuazioni casuali che spingono occasionalmente $V$ oltre $V_\theta$. In questo regime $\sigma$ è grande, e il parametro adimensionale:

\begin{equation}
x_T = \frac{V_\theta - \mu\tau}{\sigma\sqrt{\tau}}
\end{equation}

è \textbf{grande e positivo} ($x_T \ll 1$). L'equazione di Ricciardi-Siegert può essere approssimata utilizzando la funzione di errore complementare (erfc):

\begin{equation}
\Phi \approx \frac{1}{\tau} \int_{\frac{V_\theta - \mu\tau}{\sigma\sqrt{\tau}}}^\infty \frac{e^{-z^2/2}}{\sqrt{2\pi}} dz = \frac{1}{2\tau}\,\text{erfc}\!\left(\frac{x_T}{\sqrt{2}}\right)
\end{equation}

\subsubsection{Sviluppo Asintotico per $x_T \ll 1$}

Poiché siamo nel regime noise-dominated, $x_T \ll 1$. La funzione erfc si sviluppa in serie di Taylor asintotica come:

\begin{equation}
\text{erfc}\!\left(\frac{x_T}{\sqrt{2}}\right) \approx \frac{e^{-x_T^2/2}}{\sqrt{\pi}\, x_T/\sqrt{2}}
\end{equation}

Sostituendo, la frequenza di scarica diventa:

\begin{equation}
\Phi \simeq \frac{e^{-x_T^2 / 2}}{\sqrt{2\pi} \tau x_T} \sim \frac{\sigma\sqrt{\tau}}{\tau(V_\theta - \mu\tau)} e^{-x_T^2 / 2}
\end{equation}

\subsubsection{I Due Meccanismi di Regolazione della Frequenza di Scarica}

Questa formula approssimata (iperbolica) sintetizza i due \textbf{meccanismi fondamentali} di regolazione della frequenza di scarica:
\begin{enumerate}
    \item \textbf{Ampiezza delle fluttuazioni ($\sigma$)}: al numeratore, aumentare $\sigma$ incrementa la probabilità che le fluttuazioni raggiungano la soglia.
    \item \textbf{Corrente media ($\mu$)}: al denominatore, aumentare $\mu$ avvicina $V_\infty = \mu\tau$ a $V_\theta$, facendo diminuire il denominatore e facendo crescere la frequenza di scarica molto rapidamente.
\end{enumerate}

Questa espressione mostra come la funzione guadagno cresca lentamente a bassi $\nu$ e poi \textbf{esploda rapidamente} quando $\mu\tau \to V_\theta$.

\section{Rete Puramente Eccitatoria}

\subsubsection{Il Modello di Amit-Brunel}

Il docente introduce l'analisi sistematica condotta da \textbf{Daniel Amit e Nicolas Brunel nel 1997} (Cerebral Cortex). Il primo passo è verificare se una \textbf{rete omogenea puramente eccitatoria} — $N$ neuroni LIF eccitatori con $J > 0$ e $K_E$ contatti presinaptici — sia in grado di esprimere uno stato stabile a frequenza moderata. I momenti della rete sono:

\begin{equation}
\mu = K_E J \nu + \mu_{\text{ext}}, \qquad \sigma^2 = K_E J^2\!\left(1 + \frac{\Delta J^2}{J^2}\right)\!\nu + \sigma^2_{\text{ext}}
\end{equation}

\subsubsection{Conseguenze della Pendenza $\Phi' > 1$}

Per una rete puramente eccitatoria, la funzione $\Phi(\nu)$ cresce sempre più rapidamente per l'effetto di feedback positivo. Amit e Brunel mostrarono che \textbf{tutti i punti di intersezione} (tranne le soluzioni banali) presentano una pendenza della funzione guadagno $\tilde{\Phi}' > 1$.

Il fatto che la pendenza alle intersezioni sia maggiore di 1 implica che tutti i punti fissi della rete puramente eccitatoria sono \textbf{instabili}. La dinamica della rete è quindi attratta verso uno dei due estremi:
\begin{enumerate}
    \item \textbf{Silenzio totale} ($\nu = 0$): il punto fisso a $\nu = 0$ con pendenza zero è l'unico punto stabile, ma la corteccia cerebrale non tace mai completamente in condizioni fisiologiche.
    \item \textbf{Esplosione dell'attività}: la rete è attratta verso saturazione ($\nu \approx 1/\tau_{\text{ref}}$).
\end{enumerate}

\begin{tcolorbox}[colback=gray!5, colframe=gray!40, arc=2mm, boxrule=0.5pt]
\textbf{Nota del docente:} Questo risultato ha una solida base intuitiva: l'eccitazione è un meccanismo di feedback positivo. Senza un contrappeso, una perturbazione positiva si amplifica indefinitamente, portando all'esplosione dell'attività. L'inibizione è il meccanismo naturale per bilanciare questo feedback.
\end{tcolorbox}

\section{Il Circuito Canonico Corticale: Tre Popolazioni Interagenti}

\subsubsection{Struttura del Microcircuito Canonico}

Per produrre uno stato asincrono spontaneo stabile, è necessario introdurre neuroni inibitori. Il \textbf{microcircuito canonico corticale} è composto da tre popolazioni distinte:

\begin{enumerate}
    \item \textbf{Popolazione E (eccitatoria)}: neuroni locali, con sinapsi glutamatergiche ($J_{EE} > 0$);
    \item \textbf{Popolazione I (inibitoria)}: interneuroni locali, con sinapsi GABAergiche ($J_{EI} < 0$);
    \item \textbf{Popolazione di Background (X)}: altre aree corticali; modellate come sorgente esterna.
\end{enumerate}

La convenzione adottata è: \textbf{indice sinistro = postsinaptico}, \textbf{indice destro = presinaptico}. Ad esempio, $J_{EI}$ è l'efficacia sinaptica dalla popolazione I (presinaptica) alla popolazione E (postsinaptica).

\subsubsection{Momenti per la Popolazione Multipla}

Generalizzando l'approssimazione diffusiva, i momenti per i neuroni eccitatori sono:

\begin{equation}
\mu_E = K_{EE} J_{EE} \nu_E - K_{IE}\,|J_{EI}|\,\nu_I + \mu_{\text{ext}}
\end{equation}

\begin{equation}
\sigma_E^2 = K_{EE} J_{EE}^2\!\left(1 + \frac{\Delta_{EE}^2}{J_{EE}^2}\right)\!\nu_E + K_{IE} J_{EI}^2\!\left(1 + \frac{\Delta_{EI}^2}{J_{EI}^2}\right)\! \nu_I + \sigma^2_{\text{ext}}
\end{equation}

Il contributo inibitorio è \textbf{sottratto} in $\mu_E$ (le efficacie inibitorie sono negative) e \textbf{sommato} in $\sigma^2_E$ (la varianza è sempre positiva). 

Il sistema da risolvere per lo stato stazionario diventa un sistema di equazioni accoppiate:
\begin{equation}
\nu_E = \Phi_E\!\left(\mu_E(\nu_E, \nu_I, \nu_X),\, \sigma_E^2(\nu_E, \nu_I, \nu_X)\right)
\end{equation}
\begin{equation}
\nu_I = \Phi_I\!\left(\mu_I(\nu_E, \nu_I, \nu_X),\, \sigma_I^2(\nu_E, \nu_I, \nu_X)\right)
\end{equation}


\subsection{Semplificazione di Amit-Brunel}

\subsubsection{La Condizione di Simmetria}

Per illustrare il ruolo dell'inibizione in modo analitico, Amit e Brunel impongono che le due popolazioni E e I ricevano la \textbf{stessa media} $\mu$ e la \textbf{stessa varianza} $\sigma^2$. In questo modo $\nu_E = \nu_I \equiv \nu$, riducendo il sistema a \textbf{una sola equazione}.

Questa condizione si realizza scegliendo:
\begin{itemize}
    \item Stessa costante di tempo $\tau$ per tutti i neuroni;
    \item Efficacie eccitatorie uguali indipendentemente dal postsinaptico: $J_{EE} = J_{IE} = J > 0$;
    \item Efficacie inibitorie proporzionali: $J_{EI} = J_{II} = -gJ$ con $g > 0$.
\end{itemize}

\subsubsection{Il Parametro $g$}

Il parametro $g$ misura la \textbf{forza relativa dell'inibizione rispetto all'eccitazione}: $g = |J_{EI}/J_{EE}|$. Con questa parametrizzazione, e definendo $K_E = \epsilon_E N_E$ e $K_I = \epsilon_I N_I$, i momenti condivisi diventano:

\begin{equation}
\mu = \mu_{\text{ext}} + J(K_E - g K_I)\,\nu
\end{equation}

\begin{equation}
\sigma^2 = \sigma_{\text{ext}}^2 + J^2(K_E + g^2 K_I)\nu
\end{equation}

\subsubsection{Condizione Analitica di Stabilità della Rete E-I}

La condizione di stabilità richiede $\tilde{\Phi}'(\nu^*) < 1$. Applicando la regola della catena all'approssimazione di diffusione (trascurando i derivati rispetto a $\sigma^2$):

\begin{equation}
\tilde{\Phi}'(\nu) \approx \frac{\partial \Phi}{\partial \mu}\bigg|_{\nu^*} \cdot \frac{d\mu}{d\nu} = \partial_\mu\Phi \cdot J(K_E - gK_I)
\end{equation}

\textbf{All'aumentare di $g$}, il fattore $(K_E - gK_I)$ diminuisce, riducendo la pendenza effettiva $\tilde{\Phi}'$. La condizione di stabilità può essere soddisfatta aumentando sufficientemente $g$. L'inibizione stabilizza lo stato asincrono riducendo il guadagno effettivo della rete rispetto alle fluttuazioni dell'input.

\begin{tcolorbox}[colback=gray!5, colframe=gray!40, arc=2mm, boxrule=0.5pt]
\textbf{Nota del docente:} Il docente mostra l'evoluzione della funzione guadagno al crescere di $g$. Per $g=3$ e $g=5$, la pendenza nel punto di intersezione è ancora $> 1$. Per $g=5.5$, l'inibizione diventa sufficiente: $\tilde{\Phi}' < 1$, e il punto fisso diventa \textbf{stabile}. Questo dimostra analiticamente che una forte inibizione è necessaria e sufficiente per stabilizzare lo stato asincrono spontaneo.
\end{tcolorbox}


\subsection{Validazione Numerica}

\subsubsection{Irregolarità del Singolo Neurone}

Simulazioni di rete con parametri biologicamente realistici ($N_E = 4000$, $N_I = 1000$) validano il campo medio. Nel regime asincrono, i singoli neuroni esibiscono dinamiche fortemente irregolari. Il potenziale di membrana fluttua costantemente al di sotto di $V_{thr}$, producendo spike solo in coincidenza di ampie fluttuazioni stocastiche. 

Per neuroni a basso firing rate (es. $\sim$5 Hz), la distribuzione dell'Inter-Spike Interval (ISI) decade in maniera esponenziale, tipica di un \textbf{processo di Poisson}, con Coefficiente di Variazione $CV \approx 1$. A frequenze più elevate (es. $\sim$20 Hz), la regolarità aumenta a causa del periodo refrattario, abbassando il $CV$.

\subsubsection{Disordine Quenched e Distribuzione Log-Normale dei Firing Rate}

Il modello a campo medio predice una frequenza singola per l'intera rete. Tuttavia, la rete strutturale possiede \textbf{disordine quenched}: il numero effettivo di connessioni in ingresso $K_i$ varia da neurone a neurone secondo una distribuzione binomiale, ben approssimata da una Gaussiana:

\begin{equation}
\begin{cases}
\mathbb{E}[K] = \epsilon N \\
Var(K) = N\epsilon(1-\epsilon)
\end{cases}
\end{equation}

Poiché il drift $\mu_i$ del singolo neurone è proporzionale al suo numero di contatti $K_i$, $\mu_i$ sarà distribuito in maniera normale nella popolazione. Inserendo questa variabile $\mu_i$ normalmente distribuita nella funzione guadagno $\Phi(\mu)$ — che nella regione sub-soglia è approssimabile con un esponenziale — si produce una \textbf{distribuzione Log-Normale} per i tassi di firing $\nu_i$ della popolazione.

\subsubsection{Cross-Correlazioni e Ritardi di Trasmissione}

Il bilanciamento dinamico garantisce asincronia a livello globale, ma imprime una fine struttura temporale misurabile attraverso la \textbf{cross-correlazione} $C_{\alpha\beta}(\tau)$:

\begin{equation}
C_{\alpha\beta}(\tau) = \frac{\langle \nu_\alpha(t) \nu_\beta(t+\tau) \rangle}{\langle \nu_\alpha(t) \rangle \langle \nu_\beta(t) \rangle}, \qquad \alpha, \beta \in \{E, I\}
\end{equation}

I grafici di cross-correlazione esibiscono oscillazioni smorzate (ripples) attorno allo zero. La periodicità dei ripples è guidata dal \textbf{ritardo di trasmissione assonale} medio $\langle \Delta_{ij} \rangle$. Incrementando tale ritardo, le oscillazioni si allargano nel tempo.
Inoltre, l'asimmetria nel cross-correlogramma $C_{EI}$ indica una precisa causalità temporale in cui l'inibizione "insegue" l'eccitazione con un breve scarto, prevenendo risposte a valanga e riaffermando l'importanza del bilanciamento stretto E-I.

\newpage

\section{Densità di Popolazione}

\subsubsection{Dinamica Transitoria e l'Equazione di Langevin}

L'approccio di campo medio fin qui sviluppato descrive solo i punti fissi stazionari. Per descrivere la \textbf{dinamica transitoria} della rete è necessario il \textbf{population density approach}.
Il punto di partenza è l'equazione di Langevin per il neurone LIF:

\begin{equation}
\tau \frac{dV}{dt} = -(V - V_{\text{eq}}) + \mu(t)\tau + \sigma(t)\sqrt{\tau}\,\eta(t)
\end{equation}

dove $\eta(t)$ è un rumore bianco gaussiano. Questo descrive un processo di Ornstein-Uhlenbeck con una \textbf{barriera assorbente} a $V = V_\theta$ e una condizione di \textbf{reset} a $V_r$.

\subsubsection{L'Equazione di Fokker-Planck}

Nell'approssimazione di campo medio, la distribuzione statistica del potenziale $V$ al tempo $t$ è descritta dalla \textbf{densità di probabilità} $P(V, t)$. Applicando il calcolo stocastico, l'evoluzione di $P(V,t)$ è data dall'\textbf{equazione di Fokker-Planck} (FP):

\begin{equation}
\frac{\partial P(V,t)}{\partial t} = -\frac{\partial}{\partial V}\!\left[\left(F(V) + \mu(t)\right) P(V,t)\right] + \frac{\sigma^2(t)}{2}\frac{\partial^2 P(V,t)}{\partial V^2}
\end{equation}

dove $F(V) = -(V - V_{\text{eq}})/\tau$ è il drift deterministico e $\sigma^2(t)/2$ è il coefficiente di diffusione. 


L'equazione di Fokker-Planck può essere riscritta come equazione di continuità:

\begin{equation}
\frac{\partial P}{\partial t} = -\frac{\partial J_{\text{FP}}}{\partial V}
\end{equation}

dove $J_{\text{FP}}(V, t)$ è il \textbf{flusso di probabilità}:

\begin{equation}
J_{\text{FP}}(V, t) = \left(F(V) + \mu(t)\right) P(V, t) - \frac{\sigma^2(t)}{2}\frac{\partial P(V,t)}{\partial V}
\end{equation}

Il \textbf{firing rate istantaneo} della popolazione è il flusso di realizzazioni che attraversano la barriera assorbente alla soglia $V_\theta$:

\begin{equation}
\nu(t) = J_{\text{FP}}(V_\theta, t)
\end{equation}

\subsubsection{Condizioni al Contorno e Re-Iniezione}

Alla \textbf{barriera assorbente} $V = V_\theta$, la condizione al contorno è $P(V_\theta, t) = 0$. Inserendo questo nel flusso di probabilità, la formula del firing rate si semplifica a:

\begin{equation}
\nu(t) = -\frac{\sigma^2(t)}{2}\,\frac{\partial P(V,t)}{\partial V}\bigg|_{V = V_\theta}
\end{equation}

Infine, per preservare la normalizzazione della densità ($\int P(V,t)dV = 1$), ogni volta che una realizzazione viene assorbita a $V_\theta$, deve essere \textbf{re-iniettata} a $V = V_r$. Si aggiunge un termine sorgente all'equazione di Fokker-Planck:

\begin{equation}
\frac{\partial P}{\partial t} = -\frac{\partial J_{\text{FP}}}{\partial V} + \nu(t)\,\delta(V - V_r)
\end{equation}


\chapter{Lezione 21}




\section{Introduzione alle Funzioni Cognitive: Memoria di Lavoro}

\subsubsection{Definizione di Memoria di Lavoro}

La \textbf{memoria di lavoro} (\textit{working memory}, WM) è una funzione cognitiva che consente di mantenere attiva un'informazione rilevante per un breve periodo di tempo, necessario per risolvere un problema imminente. È distinta dalla memoria a lungo termine per:
\begin{itemize}
    \item \textbf{Durata}: secondi o frazioni di minuto (non ore o anni).
    \item \textbf{Capacità}: limitata (circa $7\pm2$ elementi, secondo Miller).
    \item \textbf{Goal-directed}: si mantiene solo l'informazione rilevante al compito corrente.
\end{itemize}

La WM è espressa da tutti i mammiferi (e da molti altri animali, compresi insetti e rettili), suggerendo che sia un meccanismo evolutivamente conservato.

\begin{tcolorbox}[colback=gray!5, colframe=gray!40, arc=2mm, boxrule=0.5pt]
\textbf{Nota del docente:} L'analogia quotidiana: un amico chiama e dice "arrivo in cinque minuti". Nei cinque minuti seguenti si mantiene attiva questa informazione, e quando il campanello suona si sa che è l'amico. Questo è un esempio di WM \textit{goal-directed}: non si tiene a mente tutto ciò che accade (la televisione in sottofondo, la finestra aperta), ma solo l'informazione rilevante per rispondere all'amico.
\end{tcolorbox}


\subsection{Il Task Delayed Match-to-Sample Spaziale}

\subsubsection{Struttura del Task}

Il professore descrive il task sperimentale classico usato per studiare la WM nei primati: il \textbf{delayed match-to-sample} (DMS) con cue spaziale, introdotto da \textbf{Patricia Goldman-Rakic} alla fine degli anni '80 con esperimenti sul macaco.

Il protocollo è il seguente (per il singolo trial):
\begin{enumerate}
    \item \textbf{Punto di fissazione}: lo schermo mostra solo un punto centrale. L'animale deve fissarlo per tutta la durata del trial.
    \item \textbf{Presentazione del cue} ($\sim$500 ms): appare brevemente un \textbf{secondo punto} (cue periferico) in una delle 8 posizioni angolari (es. 0$^\circ$, 45$^\circ$, 90$^\circ$, \dots, 315$^\circ$).
    \item \textbf{Periodo di delay} (2–8 secondi): lo schermo mostra solo il punto di fissazione. La posizione del cue \textit{non è più visibile}.
    \item \textbf{Segnale "go"}: il punto di fissazione sparisce. L'animale deve spostare lo sguardo verso la posizione \textit{ricordata} del cue periferico.
\end{enumerate}

Se risponde correttamente riceve ricompensa (succo di frutta). Il task richiede che l'animale mantenga in mente la posizione del cue per l'intera durata del delay.

\subsubsection{Attività dei Neuroni Prefrontali durante il Delay}

I ricercatori, usando elettrodi impiantati nella \textbf{corteccia prefrontale dorsolaterale} (dlPFC) del macaco, registrano l'attività di singoli neuroni durante il task. La dlPFC è situata nella parte anteriore del cervello, nel lobo frontale.

Risultato principale: esiste una classe di neuroni che \textbf{aumenta il firing rate durante il delay} in modo selettivo per una specifica posizione del cue. Questi neuroni:
\begin{itemize}
    \item Hanno bassa attività basale (in assenza di cue).
    \item Si attivano fortemente durante la \textbf{presentazione} del cue e mantengono un'elevata attività durante tutto il \textbf{delay}.
    \item La selettività è \textbf{spaziale}: un neurone risponde fortemente solo quando il cue è in una certa direzione angolare ("posizione preferita"), e ha risposta bassa per le altre direzioni.
\end{itemize}



\subsubsection{Raster Plot e Spike Density}

Il professore mostra i dati sperimentali per un neurone registrato nel dlPFC durante il task DMS. La visualizzazione usa:
\begin{itemize}
    \item \textbf{Raster plot}: ogni riga è un trial; ogni barretta è uno spike. I trial sono raggruppati per posizione del cue (8 direzioni).
    \item \textbf{Spike density} (firing rate istantaneo): integrazione del raster in piccole finestre temporali.
\end{itemize}

Osservazioni per questo neurone:
\begin{itemize}
    \item \textbf{Cue a 45$^\circ$} (non preferita): bassa attività durante tutta la prova, incluso il delay.
    \item \textbf{Cue a 135$^\circ$} (intermedia): attività moderata durante la presentazione, ma silence durante il delay.
    \item \textbf{Cue a 270$^\circ$} (posizione preferita): forte attivazione durante il cue E \textbf{mantenimento dell'attività durante tutto il delay}. Il firing rate sale da $\sim$5 Hz basale a $\sim$30–40 Hz durante il delay, per una durata di diversi secondi.
\end{itemize}

\begin{tcolorbox}[colback=gray!5, colframe=gray!40, arc=2mm, boxrule=0.5pt]
\textbf{Nota del docente:} Questo neurone è un classico \textbf{WM neuron}: mantiene un'attività elevata e persistente durante l'assenza dello stimolo, "tenendo in mente" la posizione del cue. Questa attività persistente è il correlato neurale della memoria di lavoro.
\end{tcolorbox}



\subsubsection{WM per Oggetti, non solo per Posizioni}

Goldman-Rakic mostrò la selettività spaziale. Ma la WM non è limitata alla posizione: si può ricordare un oggetto, un suono, un odore. \textbf{Earl Miller} e \textbf{Jonathan Reinert} eseguirono un esperimento con un task DMS per \textbf{oggetti}:

\begin{enumerate}
    \item L'animale vede un \textbf{sample object} (es. una campana).
    \item Dopo un delay di qualche secondo, vengono presentati due oggetti: il sample e un distrattore (es. una cassetta delle lettere).
    \item L'animale deve scegliere (con lo sguardo o con una leva) l'oggetto uguale al sample.
\end{enumerate}

Risultati: esistono neuroni nella PFC che mantengono un'attività elevata durante il delay \textbf{selettivamente per l'oggetto da ricordare}. Un neurone può rispondere fortemente quando l'animale ricorda la "campana" e debolmente quando ricorda la "cassetta delle lettere".

\subsubsection{Codifica Distribuita: Non Un Singolo Neurone}

Osservazione cruciale: lo \textbf{stesso} neurone può avere un'attività durante il delay superiore alla basale anche per l'oggetto \textit{non preferito}, sebbene più bassa che per l'oggetto preferito. Questo suggerisce che la rappresentazione della WM non è affidata a un singolo neurone ma è \textbf{distribuita} su una popolazione di neuroni, ciascuno con diversi livelli di selettività.

Questo è un \textbf{fenomeno di rete}: non è la proprietà di un singolo neurone, ma l'orchestrazione concertata di molti neuroni. Implica che per modellare correttamente la WM è necessario considerare \textbf{reti di neuroni} con interazioni.


\subsection{Robustezza della WM}

\subsubsection{Il Task Multi-Oggetto con Distrattori}

Per studiare la \textbf{resilienza} della WM, il professore descrive un esperimento in cui il task DMS è reso più complesso: durante il delay vengono presentati \textbf{oggetti distrattori} che l'animale deve ignorare. La sequenza è:
\begin{itemize}
    \item Sample: farfalla.
    \item Delay 1 $\rightarrow$ Distrattore: ombrello $\rightarrow$ Delay 2 $\rightarrow$ Distrattore: leone $\rightarrow$ Delay 3 $\rightarrow$ Distrattore: lampada $\rightarrow$ Delay 4 $\rightarrow$ Distrattore: occhiali $\rightarrow$ Delay 5 $\rightarrow$ Test: farfalla (e un altro oggetto).
\end{itemize}

L'animale deve ricordare la farfalla per tutto il periodo, nonostante i distrattori.

\subsubsection{Resilienza del Segnale di WM}

I neuroni nella PFC che codificano la farfalla:
\begin{itemize}
    \item Mostrano alta attività durante il primo delay (ricordando la farfalla).
    \item Quando arriva l'ombrello: l'attività si modifica (risposta al distrattore), ma poi \textbf{ritorna} al livello precedente durante il delay successivo.
    \item Stesso comportamento per leone, lampada, occhiali.
    \item Alla presentazione della farfalla (matching): forte risposta $\rightarrow$ l'animale risponde correttamente.
\end{itemize}

Questo dimostra che il segnale di WM è \textbf{robusto alle perturbazioni}: la rete è capace di mantenere l'informazione nonostante l'irruzione di distrattori. Il professore sottolinea che questo è un fenomeno di rete: la resilienza emerge dalle proprietà dinamiche dei circuiti prefrontali, che verranno studiate nella prossima lezione.


\subsection{Schizofrenia e Deficit di Memoria di Lavoro}

\subsubsection{Il Deficit Comportamentale}

I pazienti affetti da \textbf{schizofrenia} mostrano deficit marcati nella WM: tendono a "perdere il filo", a cambiare frequentemente l'oggetto della loro attenzione, e a fallire nei task DMS con delay più lunghi. Le loro prestazioni sono significativamente inferiori ai soggetti normali nell'esecuzione di questi task.

\subsubsection{Differenze nell'Attività Cerebrale (fMRI BOLD)}

Confrontando pazienti schizofrenici e soggetti normali tramite \textbf{fMRI} (imaging a risonanza magnetica funzionale, segnale BOLD), durante l'esecuzione di un task DMS:
\begin{itemize}
    \item In entrambi i gruppi, la \textbf{corteccia prefrontale dorsolaterale} (DLPFC) si attiva durante il delay.
    \item Nei soggetti \textbf{normali}: forte attivazione della DLPFC bilaterale, ben visibile come blob giallo nell'immagine fMRI.
    \item Nei pazienti \textbf{schizofrenici}: l'attivazione è presente ma \textbf{meno intensa} e meno estesa (globo giallo più piccolo, significativo solo in un emisfero).
\end{itemize}

Quando si normalizza l'attività per il delay corto (task facile) e si confronta con il delay lungo (task difficile):
\begin{itemize}
    \item Soggetti normali: \textbf{incremento} dell'attività BOLD con il delay lungo.
    \item Pazienti schizofrenici: incremento \textbf{inferiore}, compatibile con un'incapacità di reclutare efficacemente il circuito prefrontale per mantenere la WM per periodi più lunghi.
\end{itemize}




\subsubsection{Riduzione delle Sinapsi nelle Dendriti della PFC}

Studi post-mortem sui cervelli di pazienti schizofrenici rivelano una \textbf{riduzione del numero di spine dendritiche} (bottoni sinaptici) sui neuroni della corteccia prefrontale. Le spine sono le strutture morfologiche su cui si formano le sinapsi eccitatorie; una riduzione delle spine implica una riduzione della \textbf{connettività eccitatoria} della rete prefrontale.

\subsubsection{Gerarchia Corticale delle Sinapsi Glutammatergiche e GABAergiche}

Un'analisi statistica recente ha misurato l'\textbf{efficacia sinaptica} (il parametro $J$ nel nostro modello) sia delle sinapsi \textbf{glutammatergiche} (eccitatorie) che \textbf{GABAergiche} (inibitorie) lungo la \textbf{gerarchia corticale}: da V1/V2 (corteccia visiva primaria, back del cervello) fino alla PFC (top della gerarchia).

Risultati nei \textbf{soggetti normali}:
\begin{itemize}
    \item \textbf{Glutammatergiche (eccitatorie)}: $J_E$ \textbf{aumenta} dal basso verso l'alto della gerarchia. Le sinapsi eccitatorie sono più forti nella PFC che in V1.
    \item \textbf{GABAergiche (inibitorie)}: $J_I$ \textbf{diminuisce} spostandosi verso la PFC. Meno inibizione in PFC, più in V1.
\end{itemize}

La PFC è quindi caratterizzata da \textbf{forte eccitazione e debole inibizione} relativa, rendendola un circuito altamente \textbf{eccitabile} e capace di generare e mantenere stati di alta attività (necessari per la WM). Al contrario, V1 ha forte inibizione per seguire fedelmente gli input sensoriali senza amplificazione.

\subsubsection{La Schizofrenia come Rottura del Bilancio E/I}

Nei pazienti \textbf{schizofrenici}, questa modulazione gerarchica è \textbf{assente o ridotta}: i $J_E$ non mostrano il graduale aumento verso la PFC. Il bilancio eccitazione-inibizione nella DLPFC è sbilanciato verso l'inibizione (o insufficientemente eccitatorio), impedendo alla rete di sostenere gli stati di alta attività necessari per la WM.

\begin{tcolorbox}[colback=gray!5, colframe=gray!40, arc=2mm, boxrule=0.5pt]
\textbf{Nota del docente:} Questo ci dice qualcosa di fondamentale: per avere WM robusta, serve un eccesso di eccitazione ricorrente nella PFC rispetto all'inibizione locale. Quando questo balance si rompe (schizofrenia), la WM fallisce. Questo sarà il punto di partenza per i modelli computazionali della WM nella prossima lezione.
\end{tcolorbox}


\subsection{Plasticità Sinaptica a Lungo Termine}

Per esprimere la WM, il cervello deve non solo mantenere un'attività persistente (problema dinamico), ma anche \textbf{scegliere} quale informazione mantenere (apprendimento selettivo). Questo richiede che le sinapsi siano \textbf{modificabili dall'esperienza}, ovvero che la rete possa "imparare" quali pattern di attività sono rilevanti da mantenere.

Il meccanismo biologico che permette questo è la \textbf{plasticità sinaptica a lungo termine} (\textit{long-term synaptic plasticity}): l'efficacia di una sinapsi (il parametro $J$ del nostro modello) può cambiare in modo permanente (su scale di ore, giorni, anni) in risposta a pattern specifici di attività neuronale.

\subsubsection{L'Ippocampo come Sito di Studio}

La plasticità sinaptica è stata storicamente studiata nell'\textbf{ippocampo}, struttura sottocorticale a forma di cavalluccio marino, responsabile della codifica della \textbf{memoria episodica} (la memoria degli eventi della vita quotidiana). L'ippocampo funge da interfaccia tra la memoria a breve termine (WM) e la memoria a lungo termine (corticale): gli episodi memorizzati nell'ippocampo vengono progressivamente consolidati nella corteccia.

L'ippocampo è organizzato in subregioni distinte: \textbf{CA1}, \textbf{CA2}, \textbf{CA3} e il \textbf{giro dentato} (DG). Queste aree sono interconnesse in modo specifico, con proiezioni denominate \textbf{collaterali di Schaffer}: gli assoni dei neuroni di CA3 proiettano verso le cellule piramidali di CA1.


\subsubsection{Protocollo Sperimentale}

I neurofisiologi hanno compiuto il seguente esperimento:
\begin{enumerate}
    \item Posizionano un \textbf{elettrodo di stimolazione} sulle \textbf{collaterali di Schaffer} (gli assoni CA3 $\rightarrow$ CA1).
    \item Registrano con una \textbf{pipetta} il potenziale di membrana di un neurone postsinaptico di CA1.
    \item Prima della stimolazione ("pre-tetanica"), misurano l'\textbf{EPSP} (\textit{Excitatory Post-Synaptic Potential}): la forma d'onda del potenziale postsinaptico eccitatorio in risposta a uno spike presinaptico.
    \item Applicano la \textbf{stimolazione tetanica}: quattro treni di spike a \textbf{100 Hz}, ciascuno della durata di $\sim$1 s. Questa stimolazione ad alta frequenza mima l'attività neuronale intensa durante un evento di apprendimento.
    \item Misurano l'EPSP nuovamente a vari istanti dopo la stimolazione: 10 min, 30 min, 1 ora, 2 ore, 24 ore, \dots
\end{enumerate}

\subsubsection{Il Risultato: Long-Term Potentiation (LTP)}

Il risultato è sorprendente: l'EPSP registrato dopo la stimolazione tetanica è \textbf{significativamente più grande} di quello pre-tetanico. E questa differenza:
\begin{itemize}
    \item \textbf{Non decade} nel tempo: rimane stabile per ore e ore dopo la stimolazione.
    \item Rappresenta quindi un \textbf{cambiamento permanente} dell'efficacia sinaptica: la sinapsi è stata \textbf{potenziata}.
\end{itemize}

Formalmente: l'efficacia sinaptica dopo la stimolazione è $J' = J + \Delta J$, con $\Delta J > 0$. Questo fenomeno è chiamato \textbf{Long-Term Potentiation} (LTP), termine introdotto da Bliss e Lømo nel 1973.

La condizione necessaria per indurre LTP: la stimolazione tetanica deve essere \textbf{prolungata} (quattro treni, non uno solo). Una singola stimolazione tetanica non è sufficiente: si osserva solo un transitorio che poi decade. Sono necessarie più sessioni di stimolazione intensa per un cambiamento duraturo.


\subsubsection{Long-Term Depression (LTD)}

Complementare all'LTP esiste anche un fenomeno opposto: la \textbf{Long-Term Depression} (LTD). Se la stimolazione consiste in una frequenza bassa (es. 1 Hz per molti minuti, invece di 100 Hz per 1 s), la sinapsi si \textbf{deprime}: l'EPSP post-stimolazione è più piccolo di quello pre-stimolazione, e il cambiamento è duraturo ($J' = J - \Delta J$, con $\Delta J > 0$).

LTD rappresenta quindi il meccanismo di \textbf{indebolimento} sinaptico a lungo termine, in complemento all'LTP che potenzia le sinapsi.

\subsection{La Regola di Hebb}

\textbf{Donald Hebb} nel 1949, nel suo libro "The Organization of Behavior", postulò il principio oggi noto come \textbf{regola di Hebb} (o \textit{hebbiana}):

\begin{tcolorbox}[colback=white, colframe=gray!40, arc=2mm, boxrule=0.5pt]
\textit{"When an axon of cell A is near enough to excite cell B, and repeatedly or persistently takes part in firing it, some growth process or metabolic change takes place in one or both cells such that A's efficiency, as one of the cells firing B, is increased."}
\end{tcolorbox}

In termini moderni: \textbf{se il neurone presinaptico e quello postsinaptico sono entrambi attivi nello stesso momento, la sinapsi tra loro si rinforza} (LTP). Se uno è attivo e l'altro no (attività anticorrelata), la sinapsi si indebolisce (LTD).

\subsubsection{La Matrice di Hebb: Regola $\Delta J$}

Il professore formalizza la regola di Hebb in forma di matrice $\Delta J$, considerando quattro scenari per una sinapsi tra neurone presinaptico e neurone postsinaptico:

\begin{table}[htbp]
    \centering
    \renewcommand{\arraystretch}{1.3}
    \begin{tabularx}{\textwidth}{|X|X|c|X|}
    \hline
    \textbf{Presinaptico $\nu_{\rm pre}$} & \textbf{Postsinaptico $\nu_{\rm post}$} & \textbf{$\Delta J$} & \textbf{Tipo} \\
    \hline
    Basso & Basso & $0$ & Nessun cambiamento \\
    \hline
    Basso & Alto & $< 0$ & LTD (depressione) \\
    \hline
    Alto & Basso & $< 0$ & LTD (depressione) \\
    \hline
    Alto & Alto & $> 0$ & \textbf{LTP (potenziamento)} \\
    \hline
    \end{tabularx}
\end{table}

La regola può essere espressa come:

\begin{equation}
\Delta J \propto (\nu_{\rm pre} - \bar{\nu})(\nu_{\rm post} - \bar{\nu})
\end{equation}

dove $\bar{\nu}$ è un firing rate di riferimento (soglia). Quando entrambi i neuroni sono correlati positivamente nella loro attività (entrambi sopra soglia), si ha LTP. Quando le loro attività sono anticorrelate (uno sopra e l'altro sotto soglia), si ha LTD.

\begin{tcolorbox}[colback=gray!5, colframe=gray!40, arc=2mm, boxrule=0.5pt]
\textbf{Nota del docente:} Questo è il principio "neurons that fire together, wire together" (Hebb, 1949). La regola di Hebb è alla base di tutti i modelli di apprendimento nella rete neurale artificiale e biologica. Nella prossima lezione vedremo come questa regola, applicata a una rete LIF, determina la struttura delle connessioni sinaptiche necessaria per esprimere la WM.
\end{tcolorbox}


\chapter{Lezione 22}


\section{La Regola di Hebb}

\subsubsection{Contesto e Motivazione}

Nella lezione precedente il docente ha introdotto il concetto di \textbf{plasticità sinaptica} (\textit{synaptic plasticity}) e di \textbf{apprendimento ambientale} (\textit{ambient learning}), in vista della costruzione di un modello di \textbf{memoria di lavoro} (\textit{working memory}). La memoria di lavoro è definita come una memoria a breve termine orientata all'obiettivo (\textit{goal-directed short-term memory}): essa consente al soggetto di mantenere attivamente un'informazione sensoriale durante il periodo in cui questa non è più disponibile nell'ambiente, per poi utilizzarla a scopo comportamentale.

L'obiettivo è comprendere se una rete composta da neuroni eccitatori e inibitori possa — attraverso la plasticità — apprendere come risolvere un task di memoria di lavoro, introducendo dinamicamente una macchina capace di codificare un'informazione attiva.

\subsubsection{La Regola di Hebb (1949)}

Considerando una coppia di neuroni — uno \textbf{pre-sinaptico} e uno \textbf{post-sinaptico} connessi da una sinapsi — la \textbf{regola di Hebb}, formulata da Donald Hebb negli anni '50 del Novecento, stabilisce come l'efficacia sinaptica $J$ si modifichi in funzione dell'attività correlata delle due cellule. In questa formulazione binaria si distinguono quattro scenari in base al firing rate (alto $\uparrow$, o basso $\downarrow$):

\begin{table}[htbp]
    \centering
    \renewcommand{\arraystretch}{1.3}
    \begin{tabularx}{\textwidth}{|X|X|X|}
    \hline
    \textbf{Pre-sinaptico} & \textbf{Post-sinaptico} & \textbf{Variazione di $J$} \\
    \hline
    $\uparrow$ (alto) & $\uparrow$ (alto) & \textbf{LTP}: $\Delta J > 0$ — potenziamento sinaptico \\
    \hline
    $\uparrow$ (alto) & $\downarrow$ (basso) & \textbf{LTD}: $\Delta J < 0$ — depressione sinaptica \\
    \hline
    $\downarrow$ (basso) & $\uparrow$ (alto) & \textbf{LTD}: $\Delta J < 0$ — depressione sinaptica \\
    \hline
    $\downarrow$ (basso) & $\downarrow$ (basso) & Nessuna variazione: $\Delta J = 0$ \\
    \hline
    \end{tabularx}
\end{table}

In sintesi: quando l'attività dei due neuroni è \textbf{positivamente correlata} (entrambi attivi) si ha \textbf{potenziamento a lungo termine} (LTP, \textit{Long-Term Potentiation}); quando essa è \textbf{anti-correlata} (uno attivo, l'altro inattivo — indipendentemente da quale sia il pre- e quale il post-sinaptico), si ha \textbf{depressione a lungo termine} (LTD, \textit{Long-Term Depression}); quando entrambi sono inattivi, l'efficacia sinaptica rimane invariata.

\begin{tcolorbox}[colback=gray!5, colframe=gray!40, arc=2mm, boxrule=0.5pt]
\textbf{Nota del docente:} La regola di Hebb costituisce la base microscopica per capire come una rete di neuroni eccitatori e inibitori possa apprendere a risolvere un task di memoria di lavoro, inducendo dinamicamente una macchina che codifica un'informazione attiva durante il periodo di delay. 
\end{tcolorbox}

\section{Il Task di Memoria di Lavoro}

\subsubsection{Il Compito Oculomotorio a Risposta Differita}

Il paradigma sperimentale di riferimento è il \textbf{compito oculomotorio a risposta differita} (\textit{oculomotor delayed-response task}). L'animale è posizionato davanti a uno schermo su cui compare un punto centrale di fissazione (\textit{central cue}). Ad un certo istante compare brevemente un \textbf{bersaglio periferico} (\textit{peripheral target}) in una di otto possibili posizioni sullo schermo. Il bersaglio scompare, inizia il \textbf{periodo di delay} (\textit{delay period}), al termine del quale l'animale deve orientare lo sguardo verso la posizione del bersaglio presentato, anche se questo non è più visibile. L'informazione deve dunque essere mantenuta attivamente in memoria durante tutto il delay.

La struttura temporale del trial è:
\begin{enumerate}
    \item \textbf{Onset del cue centrale} — l'animale si fissa sul centro dello schermo
    \item \textbf{Comparsa del bersaglio periferico} — stimolo sensoriale in una delle 8 posizioni
    \item \textbf{Scomparsa del bersaglio} — inizio del periodo di delay
    \item \textbf{Go signal} (scomparsa del cue centrale) — l'animale deve rispondere saccadando verso la posizione ricordata
\end{enumerate}

L'animale deve essere \textbf{istruito} a svolgere questo task: inizialmente non conosce la regola, e i neuroni selettivi per una data posizione non mantengono l'attività durante il delay. È questo processo di istruzione che corrisponde all'apprendimento modellato con la plasticità sinaptica.

\subsubsection{Neuroni Foreground e Background}

All'interno della rete corticale, solo una piccola \textbf{frazione $F$ dei neuroni eccitatori} mostra selettività per il bersaglio presentato: questi sono i \textbf{neuroni foreground} (o neuroni selettivi). La rimanente maggioranza costituisce i \textbf{neuroni background} (o neuroni non selettivi).

Le evidenze neurofisiologiche indicano che tipicamente:
\begin{equation}
F \approx 1\% \quad \text{(frazione di neuroni eccitatori selettivi per un dato stimolo)}
\end{equation}

Per una rete di $N \approx 10\,000$ neuroni totali — che rappresenta una piccola patch della corteccia cerebrale — la struttura demografica è:
\begin{itemize}
    \item \textbf{Neuroni eccitatori totali:} $N_E \approx 0.8 N = 8\,000$
    \item \textbf{Neuroni inibitori totali:} $N_I \approx 0.2 N = 2\,000$
    \item \textbf{Neuroni foreground:} $F \cdot N_E \approx 80$
    \item \textbf{Neuroni background:} $(1-F) \cdot N_E \approx 7\,920$
\end{itemize}

\begin{tcolorbox}[colback=gray!5, colframe=gray!40, arc=2mm, boxrule=0.5pt]
\textbf{Nota del docente:} Il valore $F \approx 1\%$ riflette la sparsa selettività osservata nelle registrazioni elettrofisiologiche della corteccia prefrontale di primati durante task di working memory. È importante sottolineare che questa è una semplificazione: in generale, durante l'apprendimento la composizione delle popolazioni può cambiare (alcuni neuroni inizialmente attivi vengono "silenziati" perché non appartengono alla rete deputata a rispondere a quello specifico stimolo). Qui si assume per semplicità che la stessa famiglia di neuroni rimanga selettiva per tutta la durata dell'esperimento.
\end{tcolorbox}



\subsubsection{Attività Prima del Training}

Prima che l'animale sia addestrato a risolvere il task, i neuroni foreground presentano il seguente comportamento durante un trial:

\begin{itemize}
    \item \textbf{Stato spontaneo} (prima dello stimolo): firing rate di fondo, $\nu_{\rm sp} \approx 1\ \mathrm{Hz}$
    \item \textbf{Risposta allo stimolo}: quando il bersaglio preferito appare, i neuroni foreground sparano a $\nu_F \approx 10\ \mathrm{Hz}$ (\textit{selectivity firing rate})
    \item \textbf{Dopo lo stimolo}: ritorno immediato allo stato spontaneo — \textbf{nessuna attività persistente durante il delay}
\end{itemize}

I neuroni background e inibitori rimangono sempre alla frequenza basale di $\approx 1\ \mathrm{Hz}$.

\subsubsection{La Matrice Sinaptica Prima del Training}

La matrice sinaptica $N \times N$ può essere partizionata in blocchi in base al tipo (eccitatorio/inibitorio) del neurone pre-sinaptico e post-sinaptico. I quattro blocchi fondamentali sono:

\begin{table}[htbp]
    \centering
    \renewcommand{\arraystretch}{1.3}
    \begin{tabularx}{\textwidth}{|X|c|c|c|}
    \hline
    \textbf{Connessione} & \textbf{Simbolo} & \textbf{Tipo} & \textbf{Segno} \\
    \hline
    Eccitatorio $\rightarrow$ Eccitatorio & $J_{EE}$ & Eccitatoria & $+$ \\
    \hline
    Eccitatorio $\rightarrow$ Inibitorio & $J_{EI}$ & Eccitatoria verso I & $+$ (ma $J_{EI} < 0$ in corrente) \\
    \hline
    Inibitorio $\rightarrow$ Eccitatorio & $J_{IE}$ & Inibitoria verso E & $-$ \\
    \hline
    Inibitorio $\rightarrow$ Inibitorio & $J_{II}$ & Inibitoria & $-$ \\
    \hline
    \end{tabularx}
\end{table}


Prima dell'apprendimento la matrice è omogenea: tutti i neuroni eccitatori ricevono lo stesso $J_{EE}$ da tutti gli altri neuroni eccitatori, indipendentemente dalla loro selettività.

\begin{tcolorbox}[colback=gray!5, colframe=gray!40, arc=2mm, boxrule=0.5pt]
\textbf{Nota del docente:} Per convenzione, le sinapsi inibitorie vengono indicate con un cerchio sull'assonema (anziché la freccia utilizzata per le eccitatorie). Questa distinzione morfologica ha un corrispettivo funzionale nel segno negativo di $J_{IE}$ e $J_{II}$.
\end{tcolorbox}


\subsection{Modifiche Hebbiane della Matrice Sinaptica}

\subsubsection{Analisi Delle Correlazioni Durante il Trial}

Durante i trial di training, si analizza l'attività di coppie di neuroni nelle diverse sottopopolazioni per ricavare il segno della modifica hebbiana:

\textbf{Connessioni Foreground $\rightarrow$ Foreground (F$\rightarrow$F):}\\
Il pre-sinaptico foreground è attivo $\uparrow$ (risponde allo stimolo) e il post-sinaptico foreground è attivo $\uparrow$ (è anch'esso selettivo per quello stimolo) $\rightarrow$ attività \textbf{positivamente correlata} $\rightarrow$ \textbf{LTP} $\rightarrow$ l'efficacia sinaptica aumenta:
\begin{equation}
J_+ > J_{EE}
\end{equation}

\textbf{Connessioni Foreground $\rightarrow$ Background (F$\rightarrow$B):}\\
Il pre-sinaptico foreground è attivo $\uparrow$ (risponde allo stimolo) e il post-sinaptico background è inattivo $\downarrow$ (non selettivo, rimane a $\nu_{\rm sp}$) $\rightarrow$ attività \textbf{anti-correlata} $\rightarrow$ \textbf{LTD} $\rightarrow$ l'efficacia sinaptica si riduce:
\begin{equation}
J_- < J_{EE}
\end{equation}

\textbf{Connessioni Background $\rightarrow$ Foreground (B$\rightarrow$F):}\\
Il pre-sinaptico background è inattivo $\downarrow$ e il post-sinaptico foreground è attivo $\uparrow$ $\rightarrow$ attività \textbf{anti-correlata} $\rightarrow$ \textbf{LTD} $\rightarrow$ l'efficacia sinaptica si riduce allo stesso $J_-$.

\textbf{Connessioni Background $\rightarrow$ Background (B$\rightarrow$B):}\\
Entrambi inattivi $\downarrow$ durante la presentazione dello stimolo $\rightarrow$ \textbf{nessuna variazione}:
\begin{equation}
J_{BB} = J_{EE}
\end{equation}

\subsubsection{Struttura della Matrice Sinaptica Post-Apprendimento}

Dopo l'addestramento, la parte eccitatoria della matrice sinaptica acquisisce una struttura a blocchi con tre regioni distinte:

\begin{itemize}
    \item \textbf{Blocco F$\rightarrow$F:} efficacia $J_+$ (potenziato — LTP)
    \item \textbf{Blocco F$\rightarrow$B e B$\rightarrow$F:} efficacia $J_-$ (depresso — LTD)
    \item \textbf{Blocco B$\rightarrow$B:} efficacia $J_{EE}$ (invariato)
\end{itemize}

\begin{tcolorbox}[colback=gray!5, colframe=gray!40, arc=2mm, boxrule=0.5pt]
\textbf{Nota del docente:} In questo schema le sinapsi inibitorie ($J_{EI}, J_{IE}, J_{II}$) sono mantenute costanti. Questa è una semplificazione: in generale esiste plasticità anche nelle sinapsi inibitorie, ma per rendere il modello analiticamente trattabile si assume invarianza delle connessioni che coinvolgono interneuroni.
\end{tcolorbox}

La matrice con questa struttura a blocchi definisce tre \textbf{popolazioni omogenee}: all'interno di ciascuna popolazione, ogni neurone riceve statisticamente lo stesso insieme di input sinaptici (albero dendritico statisticamente equivalente). Questo è il prerequisito per l'applicazione dell'\textbf{approssimazione di campo medio} (\textit{mean-field approximation}).

\section{Equazioni di Campo Medio per le Tre Popolazioni}

\subsubsection{Corrente Media e Varianza per i Neuroni Foreground}

La corrente sinaptica media $\mu_F$ ricevuta da un neurone foreground è la somma lineare dei contributi di tutte le popolazioni pre-sinaptiche, ciascuno ponderato per l'efficacia sinaptica media e per la frequenza di arrivo degli spike:

\begin{equation}
\mu_F = J_+ K_E F \,\nu_F + J_- K_E (1-F)\,\nu_B + J_{EI} K_I \,\nu_I
\end{equation}

dove:
\begin{itemize}
    \item $K_E$ = numero medio di contatti sinaptici eccitatorio$\rightarrow$eccitatorio per neurone
    \item $K_I$ = numero medio di contatti sinaptici inibitorio$\rightarrow$eccitatorio per neurone
    \item $\nu_F, \nu_B, \nu_I$ = firing rate medi delle tre popolazioni
    \item Il primo termine: auto-eccitazione potenziata ($J_+$) dai $K_E F$ neuroni foreground pre-sinaptici
    \item Il secondo termine: eccitazione depressa ($J_-$) dai $K_E(1-F)$ neuroni background
    \item Il terzo termine: inibizione ($J_{EI} < 0$) dai $K_I$ neuroni inibitori
\end{itemize}

La \textbf{varianza infinitesimale} della corrente (proveniente dal rumore di shot degli input di Poisson) è:

\begin{equation}
\sigma_F^2 = J_+^2 (1+\delta_+^2)\, K_E F\, \nu_F + J_-^2 (1+\delta_-^2)\, K_E (1-F)\, \nu_B + J_{EI}^2 (1+\delta_{EI}^2)\, K_I\, \nu_I
\end{equation}

dove $\delta_+^2, \delta_-^2, \delta_{EI}^2$ sono le varianze relative delle distribuzioni delle efficacie sinaptiche nei rispettivi blocchi. La varianza è sempre proporzionale al quadrato dell'efficacia sinaptica e al numero di spike ricevuti per unità di tempo.

\subsubsection{Corrente Media e Varianza per i Neuroni Background}

Analogamente, la corrente media per un neurone background è:

\begin{equation}
\mu_B = J_- \, K_E \, F \, \nu_F + J_{EE} \, K_E \, (1-F) \, \nu_B + J_{EI} \, K_I \, \nu_I
\end{equation}

Il confronto chiave con $\mu_F$ è nel primo termine: i neuroni foreground contribuiscono con $J_-$ (depresso) ai background, anziché con $J_+$ (potenziato) come per i foreground stessi. Il secondo termine è $J_{EE}$ (non modificato), anziché $J_-$.

La varianza $\sigma_B^2$ ha la stessa struttura formale di $\sigma_F^2$, con $J_+$ sostituito da $J_-$ nel primo termine e $J_{EE}$ nel secondo:

\begin{equation}
\sigma_B^2 = J_-^2 (1+\delta_-^2)\, K_E F\, \nu_F + J_{EE}^2 (1+\delta_{EE}^2)\, K_E (1-F)\, \nu_B + J_{EI}^2 (1+\delta_{EI}^2)\, K_I\, \nu_I
\end{equation}

\subsubsection{Corrente Media per i Neuroni Inibitori}

I neuroni inibitori ricevono input eccitatorio da tutti i neuroni eccitatori (foreground e background) attraverso le sinapsi $J_{EI}$ e si auto-inibiscono tramite $J_{II}$:

\begin{equation}
\mu_I = J_{EI} \, K_E \left[ F \, \nu_F + (1-F) \, \nu_B \right] + J_{II} \, K_I \, \nu_I
\end{equation}

Il termine tra parentesi quadre è il \textbf{firing rate medio pesato} dell'intera popolazione eccitatoria. Poiché $F \approx 1\%$, anche un grande aumento di $\nu_F$ produce una variazione minima di questo termine: \textbf{l'inibizione è approssimativamente invariata}. Questo è un punto fondamentale che il docente sottolinea esplicitamente: grazie alla piccola dimensione della popolazione foreground rispetto alla rete totale, i neuroni inibitori non riescono a "distinguere" se il segnale eccitatorio proviene da neuroni selettivi o meno.

\begin{tcolorbox}[colback=gray!5, colframe=gray!40, arc=2mm, boxrule=0.5pt]
\textbf{Nota del docente:} L'inibizione è simile alla media pesata del firing rate sull'intera rete eccitatoria. Poiché $F$ è piccolo, l'inibizione varia poco — e questa variazione si scarica uniformemente su foreground e background. Questo meccanismo non distrugge la selettività; al contrario, la rende possibile attraverso la differenza di input eccitatorio tra le due popolazioni.
\end{tcolorbox}




\subsubsection{Corrente Media come Funzione Lineare del Firing Rate}

Tutte le correnti medie $\mu_\alpha$ e le varianze $\sigma^2_\alpha$ (con $\alpha \in \{F, B, I\}$) sono funzioni lineari dei tre firing rate:

\begin{equation}
\mu_\alpha = \mu_\alpha(\nu_F, \nu_B, \nu_I), \qquad \sigma^2_\alpha = \sigma^2_\alpha(\nu_F, \nu_B, \nu_I)
\end{equation}

\subsubsection{Differenza Cruciale tra Corrente Foreground e Background}

Il confronto diretto tra $\mu_F$ e $\mu_B$ rivela la \textbf{differenza fondamentale} che genera la selettività:

\begin{itemize}
    \item $\mu_F$ contiene $J_+ K_E F \nu_F$ (auto-eccitazione \textbf{potenziata})
    \item $\mu_B$ contiene $J_- K_E F \nu_F$ (eccitazione \textbf{depressa} dai foreground)
\end{itemize}

Poiché $J_+ > J_{EE} > J_-$, a parità di $\nu_F$ i neuroni foreground ricevono un input eccitatorio molto maggiore rispetto ai background. Il termine dal background verso i foreground ($J_- K_E(1-F)\nu_B$) è depresso rispetto al termine omologa nei background ($J_{EE}K_E(1-F)\nu_B$).

\subsubsection{Invarianza Approssimata dell'Inibizione}

Poiché $F \approx 1\%$, la corrente inibitoria:

\begin{equation}
\mu_I \propto F\,\nu_F + (1-F)\,\nu_B \approx \nu_B \quad \text{(per } F \ll 1\text{)}
\end{equation}

varia pochissimo al variare di $\nu_F$. Di conseguenza, anche l'inibizione che si riversa sui foreground e sui background è \textbf{praticamente uguale} nelle due popolazioni. Il vantaggio dei foreground rispetto ai background è quindi esclusivamente di natura eccitatoria ($J_+$ vs $J_-$), non inibitoria.

\begin{tcolorbox}[colback=gray!5, colframe=gray!40, arc=2mm, boxrule=0.5pt]
\textbf{Nota del docente:} Ricapitolando: se si aumenta l'input ai foreground tramite uno stimolo esterno, i foreground aumentano la loro attività grazie a $J_+$, mentre i background ricevono da questi un input depresso ($J_-$). L'inibizione aumenta leggermente per tutti, ma in egual misura. Il risultato netto è che i foreground salgono e i background scendono — esattamente il pattern di attività osservato sperimentalmente.
\end{tcolorbox}


\subsection{Equazioni di Auto-Consistenza}


Le \textbf{equazioni di auto-consistenza del campo medio} sono tre equazioni accoppiate, una per ciascuna popolazione:

\begin{equation}
\nu_\alpha = \Phi\!\left(\mu_\alpha(\nu_F, \nu_B, \nu_I),\; \sigma^2_\alpha(\nu_F, \nu_B, \nu_I)\right), \qquad \alpha \in \{F, B, I\}
\end{equation}

dove $\Phi(\mu, \sigma^2)$ è la \textbf{funzione di trasferimento} (\textit{transfer function} o \textit{gain function}) del neurone LIF — la stessa già derivata nelle lezioni precedenti tramite l'equazione di Fokker-Planck e il problema del tempo di primo passaggio. Essa fornisce il firing rate stazionario di un neurone LIF in funzione della media $\mu$ e della varianza $\sigma^2$ dell'input corrente.

La soluzione del sistema fornisce i \textbf{punti fissi} della rete, cioè gli stati stazionari di attività collettiva.

\subsubsection{La Dinamica di Rilassamento al Punto Fisso}

La dinamica che porta il sistema al punto fisso è descritta dalle equazioni del firing rate:

\begin{equation}
\tau_\alpha \frac{d\nu_\alpha}{dt} = \Phi(\mu_\alpha, \sigma^2_\alpha) - \nu_\alpha
\end{equation}

dove $\tau_\alpha$ è la costante di rilassamento del firing rate per la popolazione $\alpha$. Allo stato stazionario si ha $d\nu_\alpha/dt = 0$, che fornisce l'equazione di auto-consistenza.

\subsubsection{Riduzione a Equazione Scalare in $\nu_F$}

Il sistema di tre equazioni accoppiate può essere ridotto a una \textbf{singola equazione scalare} per $\nu_F$. A partire dalle equazioni per background e inibitori, fissato $\nu_F$, si possono trovare le soluzioni stazionarie:

\begin{equation}
\nu_B^* = \nu_B^*(\nu_F), \qquad \nu_I^* = \nu_I^*(\nu_F)
\end{equation}

che esprimono il firing rate delle due popolazioni "passive" come funzione della variabile di interesse $\nu_F$. Sostituendo nella prima equazione si ottiene la \textbf{funzione di guadagno effettiva}:

\begin{equation}
\nu_F = \Phi_{\rm eff}(\nu_F) \equiv \Phi\!\left(\mu_F\bigl(\nu_F,\, \nu_B^*(\nu_F),\, \nu_I^*(\nu_F)\bigr),\; \sigma_F^2\bigl(\nu_F,\, \nu_B^*(\nu_F),\, \nu_I^*(\nu_F)\bigr)\right)
\end{equation}

Le soluzioni di questa equazione scalare (intersezioni tra la curva $\Phi_{\rm eff}(\nu_F)$ e la retta identità $\nu_F$) sono i punti fissi del sistema.

\begin{tcolorbox}[colback=gray!5, colframe=gray!40, arc=2mm, boxrule=0.5pt]
\textbf{Nota del docente:} Questa riduzione dimensionale è possibile perché, nella fase stazionaria, i firing rate di background e inibitori si adattano "adiabaticamente" al valore di $\nu_F$, che è la variabile lenta — quella che porta la memoria selettiva. Si può in tal modo costruire un'\textbf{equazione efficace monodimensionale} che cattura tutta la fisica essenziale del sistema.
\end{tcolorbox}


\subsection{Analisi Grafica dei Punti Fissi con Stimolo Esterno}

\subsubsection{La Funzione di Guadagno Effettiva}

Nel grafico $(\nu_F, \Phi_{\rm eff}(\nu_F))$, la funzione di guadagno è una curva monotonicamente crescente e concava (a forma di funzione di risposta sigmoidale). I punti fissi si trovano alle intersezioni con la bisettrice $\nu_F = \nu_F$.

Prima dell'apprendimento (o con $J_+ \approx J_{EE}$), la curva interseca la bisettrice in un \textbf{unico punto}: il punto fisso spontaneo.

\subsubsection{Effetto dell'Input Sensoriale Esterno}

In presenza del bersaglio periferico sullo schermo, i neuroni foreground ricevono un \textbf{input eccitatorio aggiuntivo} $\nu_{\rm stim} > 0$ — assente per i neuroni background. Questo termina si aggiunge alla corrente media:

\begin{equation}
\mu_F \to \mu_F + \nu_{\rm stim}
\end{equation}

Graficamente, aggiungere $\nu_{\rm stim}$ equivale a traslare la funzione di guadagno effettiva verso \textbf{sinistra} sull'asse $\nu_F$ (il punto in cui la corrente netta è zero si sposta a sinistra). L'intersezione si sposta a un valore maggiore di $\nu_F$: i foreground aumentano il loro firing rate.

\subsubsection{Risposta delle Tre Popolazioni}

Quando $\nu_{\rm stim} \neq 0$:

\begin{enumerate}
    \item \textbf{Neuroni foreground:} $\nu_F$ aumenta perché ricevono $\nu_{\rm stim}$ amplificato da $J_+$
    \item \textbf{Neuroni inibitori:} $\nu_I^*(\nu_F)$ aumenta lievemente, perché la media $F\nu_F + (1-F)\nu_B$ cresce leggermente
    \item \textbf{Neuroni background:} $\nu_B^*(\nu_F)$ \textbf{diminuisce}, perché l'aumento dell'inibizione si riversa su di loro senza che abbiano il vantaggio dell'auto-eccitazione potenziata $J_+$
\end{enumerate}

Questo realizza il pattern osservato sperimentalmente: i neuroni selettivi aumentano, i non selettivi vengono soppressi, l'inibizione globale sale lievemente.

\subsection{Evoluzione Temporale dell'Apprendimento e Vincolo Omeostatico}

\subsubsection{Dinamica di $J_+$ e $J_-$ Durante il Training}

Sull'asse del tempo dell'apprendimento (scala temporale dei trial, molto più lenta della dinamica neuronale), le efficacie sinaptiche evolvono monotonicamente:
\begin{itemize}
    \item $J_+$ cresce progressivamente (LTP cumulativo F$\rightarrow$F ad ogni trial)
    \item $J_-$ decresce progressivamente (LTD cumulativo F$\rightarrow$B e B$\rightarrow$F ad ogni trial)
\end{itemize}

Il docente menziona che esistono processi omeostatici che regolano questa evoluzione. Il vincolo imposto è che il \textbf{firing rate medio dell'intera popolazione eccitatoria} rimanga costante:

\begin{equation}
\langle \nu_E \rangle = F \, \nu_F + (1-F) \, \nu_B = \nu_{E0} = \text{const}
\end{equation}

\subsubsection{Derivazione Analitica del Vincolo su $J_\pm$}

Linearizzando intorno alla soluzione omogenea (pre-apprendimento) $\nu_F = \nu_B = \nu_{E0}$, si scrive:
\begin{equation}
\nu_F = \nu_{E0} + \delta\nu_F, \qquad \nu_B = \nu_{E0} + \delta\nu_B
\end{equation}

Imponendo il vincolo al primo ordine in $\delta\nu_F, \delta\nu_B$:
\begin{equation}
F \, \delta\nu_F + (1-F) \, \delta\nu_B = 0 \implies \delta\nu_B = -\frac{F}{1-F}\,\delta\nu_F
\end{equation}

Il calcolo di Taylor delle equazioni di campo medio esprime $\delta\nu_F$ e $\delta\nu_B$ in funzione di $\delta J_+$ e $\delta J_-$. Sostituendo si ottiene un'equazione lineare che lega i cambiamenti delle efficacie sinaptiche, definendo il \textbf{percorso di apprendimento} nel piano $(J_+, J_-)$.

\begin{tcolorbox}[colback=gray!5, colframe=gray!40, arc=2mm, boxrule=0.5pt]
\textbf{Nota del docente:} Ci sono due scale temporali distinte che coesistono nel sistema: la scala temporale della dinamica del firing rate (millisecondi) e la scala temporale dell'apprendimento sinaptico (secondi–minuti per trial). Questa separazione di scale è ciò che rende valida l'analisi adiabatica: per ogni valore di $J_+$ e $J_-$, il firing rate raggiunge il suo punto fisso prima che le sinapsi cambino apprezzabilmente.
\end{tcolorbox}


\section{Deformazione della Funzione di Guadagno con l'Apprendimento}

\subsubsection{La Pendenza della Funzione di Guadagno È Proporzionale a $J_+$}

Espandendo al primo ordine intorno al punto fisso $\nu_F^*$, la derivata della funzione di guadagno effettiva è:

\begin{equation}
\left.\frac{\partial \Phi_{\rm eff}}{\partial \nu_F}\right|_{\nu_F^*} = \frac{\partial \Phi}{\partial \mu}\cdot \frac{\partial \mu_F}{\partial \nu_F} + \frac{\partial \Phi}{\partial \sigma^2}\cdot \frac{\partial \sigma^2_F}{\partial \nu_F}
\end{equation}

Sostituendo le espressioni di $\mu_F$ e $\sigma^2_F$:
\begin{equation}
\frac{\partial \mu_F}{\partial \nu_F} \approx J_+ K_E F, \qquad \frac{\partial \sigma^2_F}{\partial \nu_F} \approx J_+^2 (1+\delta_+^2) K_E F
\end{equation}

si ottiene che la pendenza è \textbf{proporzionale a $J_+$}:

\begin{equation}
\left.\frac{\partial \Phi_{\rm eff}}{\partial \nu_F}\right|_{\nu_F^*} \propto J_+
\end{equation}

(più precisamente, contiene un termine lineare in $J_+$ che proviene da $\partial_\mu\Phi$ e un termine in $J_+$ che proviene da $\partial_{\sigma^2}\Phi \cdot J_+^2$, dove $\partial_{\sigma^2}\Phi$ è sempre positivo per neuroni con rumore).

\subsubsection{Compressione dell'Asse e Aumento della Non-Linearità}

Poiché la funzione di guadagno effettiva può essere scritta nella forma:
\begin{equation}
\Phi_{\rm eff}(\nu_F) \approx f(J_+ \cdot \nu_F + C)
\end{equation}

dove $f$ è la funzione di risposta LIF e $C$ è una costante, aumentare $J_+$ equivale a \textbf{comprimere l'asse} $\nu_F$. Graficamente: la stessa variazione di $\nu_F$ produce una variazione maggiore nella curva $\Phi_{\rm eff}$. La curva diventa via via più ripida e più non-lineare.

Partendo da una curva quasi lineare (pre-apprendimento, $J_+ \approx J_{EE}$), all'aumentare di $J_+$ la curva si deforma progressivamente verso una \textbf{forma a S} (\textit{sigmoid}) sempre più pronunciata.

\subsubsection{Comparsa di Tre Punti Fissi}

Quando $J_+$ supera il valore critico $J_+^c$, la curva $\Phi_{\rm eff}(\nu_F)$ diventa abbastanza ripida da intersecare la retta identità $\nu_F$ in \textbf{tre punti distinti}:

\begin{itemize}
    \item \textbf{Punto fisso inferiore} ($\nu_F = \nu_{\rm sp}$): \textbf{stabile} — corrisponde allo stato spontaneo
    \item \textbf{Punto fisso intermedio} ($\nu_F = \nu_{\rm saddle}$): \textbf{instabile} — è il punto sella (o punto di separatrice)
    \item \textbf{Punto fisso superiore} ($\nu_F = \nu_{\rm sel}$): \textbf{stabile} — corrisponde al nuovo stato selettivo
\end{itemize}

\begin{tcolorbox}[colback=gray!5, colframe=gray!40, arc=2mm, boxrule=0.5pt]
\textbf{Nota del docente:} La stabilità dei punti fissi si determina dalla condizione sulla pendenza: un punto fisso è stabile se in quel punto la pendenza di $\Phi_{\rm eff}(\nu_F)$ è inferiore a 1 (pendenza della bisettrice); è instabile se la pendenza è superiore a 1. Questo è il criterio classico per la stabilità lineare di un punto fisso in una mappa scalare.
\end{tcolorbox}


\section{Funzione di Lyapunov e Paesaggio Energetico}

\subsubsection{L'Equazione del Firing Rate come Sistema Dissipativo}

La dinamica del firing rate della popolazione foreground è un sistema dissipativo governato dall'equazione:

\begin{equation}
\tau_F \frac{d\nu_F}{dt} = \Phi_{\rm eff}(\nu_F) - \nu_F \equiv \mathcal{F}(\nu_F)
\end{equation}

Il termine a destra, $\mathcal{F}(\nu_F) = \Phi_{\rm eff}(\nu_F) - \nu_F$, è il \textbf{campo di forza} (\textit{force field}): determina il verso e l'intensità della deriva del firing rate. I punti fissi si trovano dove $\mathcal{F}(\nu_F^*) = 0$.

\begin{itemize}
    \item Dove $\mathcal{F} > 0$ (cioè $\Phi_{\rm eff} > \nu_F$): il firing rate cresce
    \item Dove $\mathcal{F} < 0$ (cioè $\Phi_{\rm eff} < \nu_F$): il firing rate decresce
\end{itemize}

\subsubsection{Costruzione della Funzione di Lyapunov}

Si cerca una \textbf{funzione di Lyapunov} $E(\nu_F)$ che soddisfi:

\begin{equation}
\frac{dE}{dt} \leq 0 \quad \forall t \quad (\text{non-crescente lungo le traiettorie})
\end{equation}

Poiché:
\begin{equation}
\frac{dE}{dt} = \frac{\partial E}{\partial \nu_F} \cdot \frac{d\nu_F}{dt} = \frac{\partial E}{\partial \nu_F} \cdot \frac{\mathcal{F}(\nu_F)}{\tau_F}
\end{equation}

per garantire $dE/dt \leq 0$, si impone:

\begin{equation}
\frac{\partial E}{\partial \nu_F} = -\mathcal{F}(\nu_F) = \nu_F - \Phi_{\rm eff}(\nu_F)
\end{equation}

In questo modo:
\begin{equation}
\frac{dE}{dt} = \left(\nu_F - \Phi_{\rm eff}(\nu_F)\right) \cdot \frac{\Phi_{\rm eff}(\nu_F) - \nu_F}{\tau_F} = -\frac{\left[\Phi_{\rm eff}(\nu_F) - \nu_F\right]^2}{\tau_F} \leq 0
\end{equation}

La funzione di Lyapunov è quindi l'integrale del campo di forza cambiato di segno:

\begin{equation}
\boxed{E(\nu_F) = -\int_0^{\nu_F} \left[\Phi_{\rm eff}(x) - x\right] dx}
\end{equation}

\subsubsection{Interpretazione come Paesaggio Energetico}

$E(\nu_F)$ si comporta come un \textbf{paesaggio energetico} (\textit{energy landscape}): il sistema evolve come una pallina che rotola seguendo il gradiente discendente di questo paesaggio. I \textbf{minimi locali} corrispondono agli stati stazionari stabili (attrattori); i \textbf{massimi locali} corrispondono agli stati instabili (selle).

\textbf{Caso pre-apprendimento (un solo punto fisso):}\\
$\mathcal{F}(\nu_F)$ ha un solo zero: il campo di forza è negativo a destra e positivo a sinistra del punto fisso. $E(\nu_F)$ ha una \textbf{singola buca} — un unico attrattore. Qualunque condizione iniziale porta il sistema verso $\nu_{\rm sp}$.

\textbf{Caso post-apprendimento (tre punti fissi):}\\
$\mathcal{F}(\nu_F)$ ha tre zeri. $E(\nu_F)$ ha \textbf{due buche separate da una barriera} — il sistema è \textbf{bistabile}: la pallina può trovarsi nel minimo sinistro ($\nu_{\rm sp}$, stato spontaneo) o nel minimo destro ($\nu_{\rm sel}$, stato selettivo), a seconda delle condizioni iniziali.

\begin{tcolorbox}[colback=gray!5, colframe=gray!40, arc=2mm, boxrule=0.5pt]
\textbf{Nota del docente:} Il sistema è dissipativo: l'energia $E$ può solo decrescere o rimanere costante (nei punti fissi). Questo significa che in assenza di perturbazioni esterne, il sistema "rotola verso il basso" nel paesaggio energetico. Non può mai aumentare la propria energia spontaneamente — è un sistema termodinamicamente ben definito.
\end{tcolorbox}


\subsection{Diagramma di Biforcazione}

\subsubsection{Struttura del Diagramma}

Al variare di $J_+$ (parametro d'ordine dell'apprendimento), il sistema attraversa una \textbf{biforcazione sella-nodo} (\textit{saddle-node} o \textit{fold bifurcation}): il numero di punti fissi stazionari cambia discontinuamente.

Il diagramma di biforcazione mostra il firing rate stazionario $\nu_F^*$ in funzione di $J_+$:

\begin{itemize}
    \item \textbf{Per $J_+ < J_+^c$:} esiste un \textbf{unico punto fisso stabile} a basso firing rate: lo stato spontaneo $\nu_{\rm sp}$
    \item \textbf{Per $J_+ = J_+^c$:} \textbf{punto critico di biforcazione} — il secondo e terzo punto fisso emergono simultaneamente (saddle-node)
    \item \textbf{Per $J_+ > J_+^c$:} esistono \textbf{tre punti fissi}: due stabili ($\nu_{\rm sp}$ e $\nu_{\rm sel}$) e uno instabile ($\nu_{\rm saddle}$) — regime bistabile
    \item \textbf{Per $J_+$ molto grande:} lo stato spontaneo e il punto sella si avvicinano e si annullano (seconda biforcazione); rimane solo $\nu_{\rm sel}$
\end{itemize}

Il \textbf{regime di interesse biologico} è quello bistabile ($J_+ > J_+^c$ ma non troppo grande), in cui i due attrattori coesistono.

\subsubsection{Comportamento delle Altre Popolazioni}

Il diagramma di biforcazione completo mostra anche:
\begin{itemize}
    \item $\nu_B^*(J_+)$: il background si abbassa nel regime bistabile quando il foreground è nello stato selettivo
    \item $\nu_I^*(J_+)$: l'inibitore sale leggermente nello stato selettivo
    \item Al crescere di $J_+$, il breaking of symmetry tra foreground e background diventa sempre più pronunciato
\end{itemize}

\begin{tcolorbox}[colback=gray!5, colframe=gray!40, arc=2mm, boxrule=0.5pt]
\textbf{Nota del docente:} È importante mantenere il bilanciamento delle risorse sinaptiche: l'aumento di $J_+$ è accoppiato alla diminuzione di $J_-$ in modo tale da preservare il firing rate medio eccitatorie. Il parametro d'ordine "fisico" del sistema è dunque la differenza $J_+ - J_-$, che cresce monotonicamente con l'apprendimento fino alla biforcazione.
\end{tcolorbox}



\subsection{Dinamica del Task: Stimolazione, Delay e Memoria}

\subsubsection{Le Fasi del Trial nel Regime Bistabile}

Con la rete nel regime bistabile ($J_+ > J_+^c$), si può analizzare la dinamica di $\nu_F(t)$ durante il trial di working memory interpretandola come moto di una pallina nel paesaggio energetico $E(\nu_F)$:

\textbf{Fase 1 — Stato spontaneo pre-stimolo:}\\
La rete è nell'attrattore spontaneo: $\nu_F = \nu_{\rm sp}$ (pallina nel minimo sinistro). Il paesaggio ha due buche separate da una barriera.

\textbf{Fase 2 — Presentazione del bersaglio preferito ($\nu_{\rm stim} \neq 0$):}\\
L'input sensoriale aggiuntivo modifica la corrente dei foreground:
\begin{equation}
\mu_F \to \mu_F + \nu_{\rm stim}
\end{equation}

Questo equivale a uno spostamento della funzione di guadagno verso sinistra, che trasforma il paesaggio energetico: la buca sinistra si innalza, la barriera si abbassa, e per un $\nu_{\rm stim}$ sufficientemente forte \textbf{la buca sinistra scompare}. Non esiste più un attrattore spontaneo: rimane solo la buca destra. La pallina, che si trovava nella buca sinistra, ora segue il gradiente discendente verso la buca destra, raggiungendo $\nu_F = \nu_{\rm sel}$ (firing rate selettivo elevato).

\textbf{Fase 3 — Rimozione dello stimolo (inizio del delay):}\\
Quando il bersaglio scompare dallo schermo: $\nu_{\rm stim} \to 0$. Il paesaggio energetico torna alla configurazione bistabile con due buche. La pallina si trova ora nella buca destra — \textbf{l'attrattore selettivo} — con $\nu_F = \nu_{\rm sel}$. \textbf{In assenza di stimolo, il sistema rimane nell'attrattore selettivo.} Questo è il substrato neuronale della \textbf{memoria di lavoro}: l'informazione "quale bersaglio è stato presentato" è codificata nell'elevato firing rate persistente dei neuroni foreground durante il delay.

\subsubsection{Meccanismo di Encoding e Mantenimento}

Il passaggio dalla buca spontanea a quella selettiva realizza la fase di \textbf{encoding} della memoria. Il mantenimento nella buca selettiva durante il delay realizza la fase di \textbf{retention}. Non sono necessari ulteriori meccanismi o input esterni per mantenere la memoria: la rete è autonomamente stabile nell'attrattore selettivo grazie alla struttura ricorrente $J_+$.

\begin{equation}
\nu_F(t) = \nu_{\rm sel} \gg \nu_{\rm sp} \quad \text{durante il delay}
\end{equation}

\begin{equation}
\nu_B(t) = \nu_B^*(\nu_{\rm sel}) < \nu_{\rm sp} \quad \text{(background silenziato durante il delay)}
\end{equation}



\section{Lettura della Memoria, Stimoli Non-Preferiti e Selettività}

\subsubsection{Stato Selettivo: Firma di Attività delle Tre Popolazioni}

Nello stato selettivo (durante il delay), il pattern di attività della rete è:

\begin{itemize}
    \item \textbf{Neuroni foreground:} $\nu_F = \nu_{\rm sel} \gg \nu_{\rm sp}$ — firing rate elevato e persistente (stato "selettivo" o "up state")
    \item \textbf{Neuroni inibitori:} $\nu_I$ leggermente aumentato rispetto allo spontaneo (per il lieve aumento di $\bar{\nu}_E$)
    \item \textbf{Neuroni background:} $\nu_B < \nu_{\rm sp}$ — \textbf{silenziati} dalla maggiore inibizione globale (soppressione del background)
\end{itemize}

Questa firma di attività può essere \textbf{letta da altre aree corticali} per decodificare quale stimolo è stato presentato e mantenuto in working memory.

\subsubsection{Stimolo Non-Preferito: Assenza di Memoria}

Quando viene presentato un bersaglio in una posizione non preferita dalla popolazione foreground in esame, l'input sensoriale che i foreground ricevono è assente o insufficiente. Lo shift della funzione di guadagno è troppo piccolo per eliminare la buca sinistra: il paesaggio energetico rimane bistabile anche durante la stimolazione. La pallina risale verso la barriera ma \textbf{non la supera}. Quando lo stimolo viene rimosso, la rete torna allo stato spontaneo:

\begin{equation}
\nu_F \to \nu_{\rm sp} \quad \text{(nessuna memoria)}
\end{equation}

Questo garantisce che la working memory sia \textbf{selettiva per il contenuto}: solo lo stimolo "preferito" dalla popolazione foreground è capace di spingere il sistema nell'attrattore selettivo.

\subsubsection{Molteplici Popolazioni e Capacità di Memoria}

In una rete corticale realistica esistono \textbf{molte popolazioni foreground distinte}, ciascuna selettiva per uno dei possibili bersagli (es. le 8 posizioni del task oculomotorio). Ciascuna popolazione ha il proprio attrattore selettivo. Quando viene presentato un bersaglio:

\begin{itemize}
    \item Solo la popolazione foreground selettiva per quel bersaglio viene spinta nel proprio attrattore selettivo
    \item Tutte le altre popolazioni rimangono nell'attrattore spontaneo
\end{itemize}

La working memory è dunque implementata come un \textbf{codice di popolazione}: la posizione del bersaglio ricordato è codificata nell'identità della popolazione foreground attiva.

\begin{tcolorbox}[colback=gray!5, colframe=gray!40, arc=2mm, boxrule=0.5pt]
\textbf{Nota del docente:} Questo tipo di attrattore, che viene implementato davvero nel cervello, è stato introdotto per la prima volta da \textbf{Daniel Amit} e \textbf{Nicola Brunel} negli anni '90 del secolo scorso. Il paper è stato caricato sulla piattaforma del corso. Qui ne viene descritta la dinamica qualitativa; grazie alla conoscenza della funzione di guadagno e dei punti fissi, è possibile risolvere numericamente le equazioni di auto-consistenza e ottenere tutti i risultati quantitativi.
\end{tcolorbox}


\section{Il Modello Amit–Brunel e il Rumore Neuronale}

\subsubsection{Il Diagramma di Biforcazione Numerico}

La soluzione numerica delle equazioni di auto-consistenza per il modello Amit–Brunel produce il diagramma di biforcazione completo, con i parametri realistici dei neuroni LIF. Il docente mostra questo diagramma e descrive le sue caratteristiche:

\begin{itemize}
    \item La curva rossa mostra il firing rate dei foreground ($\nu_F^*$) in funzione di $J_+$: si vede chiaramente il regime monostabile iniziale, la comparsa del secondo ramo (stato selettivo) alla biforcazione, e la coesistenza dei due rami stabili nel regime bistabile
    \item La curva del background ($\nu_B^*$, blu) decresce nello stato selettivo: i background sono inibiti
    \item La curva degli inibitori ($\nu_I^*$) sale leggermente nello stato selettivo
\end{itemize}

Questi risultati confermano quantitativamente tutte le previsioni qualitative del modello analitico discusso in questa lezione.

\subsubsection{Effetti del Rumore in Reti a $N$ Finito}

In una rete di dimensione \textbf{finita} $N$, il firing rate non è deterministico ma \textbf{fluttua stocasticamente} attorno al valore medio previsto dal campo medio. Queste fluttuazioni hanno un'ampiezza che scala come $1/\sqrt{N}$ e agiscono come una \textbf{temperatura efficace} nel paesaggio energetico.

L'analogia con i sistemi fisici ferromagnetici è puntuale: come la temperatura termica permette a un magnete di perdere l'ordine passando dalla fase ferromagnetica a quella paramagnetica, il \textbf{rumore neuronale} consente transizioni spontanee da un attrattore all'altro. Il rate di tali transizioni (di tipo Kramers) è proporzionale a:

\begin{equation}
r_{\rm escape} \propto e^{-\Delta E / T_{\rm eff}}
\end{equation}

dove $\Delta E$ è l'altezza della barriera energetica e $T_{\rm eff} \propto 1/N$ è la "temperatura" efficace della rete.

\subsubsection{Implicazioni per la Robustezza della Memoria}

Quando $J_+$ è appena al di sopra di $J_+^c$, la buca selettiva è poco profonda e la barriera è bassa. In questo caso:

\begin{itemize}
    \item Il sistema, spinto nello stato selettivo dallo stimolo, può \textbf{uscire spontaneamente} dall'attrattore selettivo per effetto delle fluttuazioni
    \item Il firing rate ritorna verso $\nu_{\rm sp}$ anche senza la presentazione di un nuovo stimolo
\end{itemize}

Questo fenomeno modella il \textbf{decadimento della memoria di lavoro} osservato sperimentalmente: la memoria non è indefinitamente stabile ma decade su una scala temporale determinata dalla profondità della buca (cioè da $J_+$) e dall'intensità del rumore (cioè da $N$). Reti con $J_+$ più grande (apprendimento più consolidato) hanno barriere più alte e memorie più robuste; reti con $N$ più grande hanno meno rumore e memorie più durature.

\begin{tcolorbox}[colback=gray!5, colframe=gray!40, arc=2mm, boxrule=0.5pt]
\textbf{Nota del docente:} Questa analogia con i sistemi ferromagnetici è profonda e non accidentale: la matematica dei sistemi di spin con interazioni, sviluppata da Hopfield, Amit e collaboratori, è alla base di tutta la teoria delle reti neurali attrattorali. Il rumore termico nei magneti e il rumore di Poisson nel firing rate dei neuroni giocano lo stesso ruolo formale. Un file matematico con la derivazione esatta della funzione di guadagno e il codice numerico verrà caricato sulla piattaforma del corso.
\end{tcolorbox}

\end{document}
